\documentclass[11pt,a4paper,twoside,openright]{book}

% ============================================================
% ENCODING & FONTS
% ============================================================
\usepackage[utf8]{inputenc}
\usepackage[T1]{fontenc}
\usepackage{lmodern}

% ============================================================
% MATHEMATICS
% ============================================================
\usepackage{amsmath,amssymb,amsthm,mathtools}
\usepackage{thmtools}
\usepackage{stmaryrd}        % \llbracket, \rrbracket
\usepackage{bm}              % bold math

% ============================================================
% GRAPHICS & DIAGRAMS
% ============================================================
\usepackage{tikz}
\usetikzlibrary{
  arrows.meta, positioning, shapes.geometric, fit, calc,
  backgrounds, decorations.pathmorphing, decorations.markings,
  trees, matrix, chains, patterns, shadows
}
\usepackage{pgfplots}
\pgfplotsset{compat=1.18}

% ============================================================
% TABLES
% ============================================================
\usepackage{booktabs}
\usepackage{longtable}
\usepackage{tabularx}
\usepackage{multirow}
\usepackage{array}

% ============================================================
% LAYOUT
% ============================================================
\usepackage[margin=1in]{geometry}
\usepackage{enumitem}
\usepackage{multicol}
\usepackage{caption}
\usepackage{float}
\usepackage{graphicx}
\usepackage{titlesec}
\usepackage{fancyhdr}
\usepackage{emptypage}       % blank pages truly blank

% ============================================================
% HYPERLINKS & CROSS-REFERENCES
% ============================================================
\usepackage[
  colorlinks=true,
  linkcolor=defblue,
  citecolor=thmred,
  urlcolor=edgeorange,
  bookmarksdepth=3
]{hyperref}
\usepackage{cleveref}

% ============================================================
% COLORS
% ============================================================
\usepackage{xcolor}
\definecolor{defblue}{RGB}{31,78,121}
\definecolor{thmred}{RGB}{153,0,0}
\definecolor{sepgreen}{RGB}{0,100,0}
\definecolor{layergray}{RGB}{245,245,245}
\definecolor{edgeorange}{RGB}{204,102,0}
\definecolor{notegray}{RGB}{100,100,100}
\definecolor{boundbrown}{RGB}{139,90,43}
\definecolor{analogypurple}{RGB}{102,0,153}
\definecolor{goldcolor}{RGB}{184,134,11}
\definecolor{onlineblue}{RGB}{0,51,153}

% ============================================================
% COLORED BOXES
% ============================================================
\usepackage{tcolorbox}
\tcbuselibrary{theorems,skins,breakable}

\newtcolorbox{separation}{
  colback=sepgreen!5, colframe=sepgreen!50!black,
  title=Separation Result, fonttitle=\bfseries\small,
  breakable
}
\newtcolorbox{traversal}{
  colback=defblue!5, colframe=defblue!50!black,
  title=Graph Traversal, fonttitle=\bfseries\small,
  breakable
}
\newtcolorbox{obstbox}{
  colback=edgeorange!5, colframe=edgeorange!50!black,
  title=Obstruction, fonttitle=\bfseries\small,
  breakable
}
\newtcolorbox{historical}{
  colback=goldcolor!5, colframe=goldcolor!50!black,
  title=Historical Note, fonttitle=\bfseries\small,
  breakable
}
\newtcolorbox{computational}{
  colback=layergray, colframe=notegray!50!black,
  title=Computational Illustration, fonttitle=\bfseries\small,
  breakable
}
\newtcolorbox{graphnode}[1][]{
  colback=defblue!3, colframe=defblue!30!black,
  fonttitle=\bfseries\small\sffamily,
  title={Graph Node}, #1,
  breakable
}

% ============================================================
% THEOREM ENVIRONMENTS
% ============================================================
\theoremstyle{definition}
\newtheorem{defn}{Definition}[chapter]
\newtheorem{example}[defn]{Example}
\newtheorem{notation}[defn]{Notation}
\newtheorem{convention}[defn]{Convention}

\theoremstyle{plain}
\newtheorem{thm}[defn]{Theorem}
\newtheorem{prop}[defn]{Proposition}
\newtheorem{cor}[defn]{Corollary}
\newtheorem{lemma}[defn]{Lemma}
\newtheorem{claim}[defn]{Claim}

\theoremstyle{remark}
\newtheorem{remark}[defn]{Remark}
\newtheorem{openprob}[defn]{Open Problem}
\newtheorem{warning}[defn]{Warning}

% ============================================================
% PAGE STYLE
% ============================================================
\pagestyle{fancy}
\fancyhf{}
\fancyhead[LE]{\small\textit{A Textbook of Formal Learning Theory}}
\fancyhead[RO]{\small\textit{\leftmark}}
\fancyfoot[C]{\thepage}
\renewcommand{\headrulewidth}{0.4pt}

% ============================================================
% CHAPTER TITLE FORMAT
% ============================================================
\titleformat{\chapter}[display]
  {\normalfont\huge\bfseries\color{defblue}}
  {\chaptertitlename\ \thechapter}{20pt}{\Huge}
\titlespacing*{\chapter}{0pt}{50pt}{40pt}

% ============================================================
% MACROS
% ============================================================
% Dimensions
\newcommand{\VC}{\ensuremath{\mathrm{VCdim}}}
\newcommand{\Ldim}{\ensuremath{\mathrm{Ldim}}}
\newcommand{\DSdim}{\ensuremath{\mathrm{DSdim}}}
\newcommand{\Ndim}{\ensuremath{d_{\mathrm{N}}}}
\newcommand{\Pdim}{\ensuremath{\mathrm{Pdim}}}
\newcommand{\fatshat}{\ensuremath{\mathrm{fat}}}
\newcommand{\SQdim}{\ensuremath{\mathrm{SQdim}}}

% Complexities
\newcommand{\Rad}{\ensuremath{\widehat{\mathcal{R}}}}
\newcommand{\KL}{\ensuremath{\mathrm{KL}}}
\newcommand{\ERM}{\ensuremath{\mathrm{ERM}}}
\newcommand{\SOA}{\ensuremath{\mathrm{SOA}}}

% Learning criteria
\newcommand{\Ex}{\ensuremath{\mathbf{Ex}}}
\newcommand{\BC}{\ensuremath{\mathbf{BC}}}
\newcommand{\FIN}{\ensuremath{\mathbf{FIN}}}

% Probability and expectation
\newcommand{\E}{\ensuremath{\mathbb{E}}}
\renewcommand{\Pr}{\ensuremath{\mathbb{P}}}
\newcommand{\ind}{\ensuremath{\mathbf{1}}}
\newcommand{\R}{\ensuremath{\mathbb{R}}}
\newcommand{\N}{\ensuremath{\mathbb{N}}}
\newcommand{\Z}{\ensuremath{\mathbb{Z}}}

% Risk
\newcommand{\risk}{\ensuremath{R}}
\newcommand{\emprisk}{\ensuremath{\hat{R}}}

% Graph relations
\newcommand{\rel}[1]{\texttt{\small #1}}
\newcommand{\nde}[1]{\textsf{\small #1}}
\newcommand{\edge}[3]{\nde{#1}\ensuremath{\;\xrightarrow{\text{\rel{#2}}}\;}\nde{#3}}

% Growth function
\newcommand{\growth}{\ensuremath{\Pi}}

% ============================================================
% TITLE
% ============================================================
\title{%
  \textbf{\Huge A Textbook of Formal Learning Theory}\\[1em]
  \Large Six Paradigms, Their Characterization Theorems,\\
  and the Separations Between Them\\[1em]
  \large Across Six Decades of Computational and Statistical Learning Theory
}
\author{%
  Dhruv Gupta\\
  Zetesis Labs, ARTPARK @ IISc Bangalore\\
  \texttt{dhruv@zetesis.org}
}
\date{March 2026}

% ============================================================
\begin{document}
% ============================================================

\frontmatter
\maketitle

% ============================================================
% EPIGRAPH
% ============================================================
\thispagestyle{empty}
\vspace*{\fill}
\begin{center}
\begin{minipage}{0.7\textwidth}
\centering\itshape\large
``The question is not which functions can be learned,\\
but which functions can be learned \emph{given what you already know}.''
\end{minipage}
\end{center}
\vspace*{\fill}
\newpage

% ============================================================
% PREFACE
% ============================================================
\chapter*{Preface}
\addcontentsline{toc}{chapter}{Preface}

This book presents formal learning theory as a single, navigable mathematical structure.  It treats six paradigms---PAC, online, Gold-style, exact, universal, and multiclass learning---within a unified framework of typed relations, and devotes equal attention to the separations between them.

The need for such a book arises from a structural gap. The field's foundational results are distributed across several research traditions that developed largely in parallel:

\begin{itemize}[leftmargin=*]
\item \textbf{PAC learning} (Valiant, 1984) and its combinatorial characterization through VC dimension (Vapnik--Chervonenkis, 1971; Blumer et al., 1989).
\item \textbf{Online learning} (Littlestone, 1988) and its characterization through the Littlestone dimension.
\item \textbf{Gold-style identification in the limit} (Gold, 1967) and its mind-change hierarchy (Freivalds--Smith, Case--Smith).
\item \textbf{Universal learning} (Bousquet et al., 2021) and the trichotomy theorem.
\item \textbf{Multiclass learning} and its resolution through the DS dimension (Brukhim et al., 2022).
\item \textbf{Generalization bounds} from five distinct traditions: VC/Rademacher, PAC-Bayes, algorithmic stability, information-theoretic, and margin-based.
\end{itemize}

No single textbook covers all these traditions with their precise interconnections and, critically, their precise \emph{non}-connections. Shalev-Shwartz and Ben-David~\cite{ShalevShwartzBenDavid2014} treat PAC and online learning comprehensively but omit Gold-style identification, universal learning, and information-theoretic bounds. Jain et al.~\cite{JainOshersonRoyerSharma1999} cover inductive inference thoroughly but not PAC theory. The generalization-bound literature is scattered.

This book fills the gap by treating all traditions within a single framework organized around typed relations between concepts.  The key structural innovation is not any new theorem, but a \emph{negative layer}: the book devotes equal attention to what \emph{does not} hold---separation results with explicit witnesses, cross-paradigm analogies with their precise obstructions, and the boundaries where one framework's characterization theorem fails to generalize to another.

\subsection*{How to Read This Book}

The book is organized in five parts:

\begin{description}[leftmargin=*]
\item[Part I: Foundations] (Chapters 1--3) introduces the base types---domains, concepts, hypothesis spaces, data presentations, and the automata-theoretic vocabulary underlying identification in the limit.
\item[Part II: Paradigms] (Chapters 4--9) presents the six major learning paradigms, each with its characterization theorem proved in full: PAC learning (VC dimension), online learning (Littlestone dimension), Gold-style identification (mind-change ordinals), exact learning (query models), and universal learning (the trichotomy).
\item[Part III: Complexity Measures] (Chapters 10--13) develops the 33 complexity measures in the graph, organized by what they measure: combinatorial dimensions, sample complexity and compression, generalization bounds, and mind-change ordinals.
\item[Part IV: The Negative Layer] (Chapters 14--15) presents all separation results with full proofs and all analogy-obstruction pairs with analysis.
\item[Part V: Extensions] (Chapters 16--18) covers computational hardness, extensions beyond binary classification, and application-layer concepts.
\end{description}

Each chapter includes TikZ diagrams showing the subgraph relevant to that chapter's concepts and computational illustrations where appropriate.  Cross-references to the companion knowledge graph are given in the appendices.

\subsection*{Prerequisites}

The reader should be comfortable with:
\begin{itemize}[leftmargin=*]
\item Basic probability (expectation, concentration inequalities, Markov/Chebyshev/Hoeffding bounds).
\item Basic combinatorics (growth functions, pigeonhole principle).
\item Computability theory at the level of decidability, Turing machines, and the halting problem (for Gold-style chapters).
\item Mathematical maturity sufficient for $\varepsilon$-$\delta$ arguments.
\end{itemize}

Background in machine learning is helpful but not required; the book is self-contained from the definitions upward.

\subsection*{Companion Materials}

The book is accompanied by a machine-readable knowledge graph and a benchmark suite, available at \url{https://github.com/Zetetic-Dhruv/Formal-Learning-Theory}.  The graph encodes 142 concepts connected by 260 typed edges across 13 relation types; it is described in detail in Appendix~A (edge inventory) and Appendix~C (validation).  The repository also contains a complete bibliography and scripts for validating the graph against the textbook's content.

\subsection*{Notation}

A notation index appears in Appendix~D. We highlight the most frequently used symbols here:

\begin{center}\small
\begin{tabular}{ll@{\qquad}ll}
\toprule
$\mathcal{C}$ & concept class & $\VC(\mathcal{H})$ & VC dimension \\
$\mathcal{H}$ & hypothesis space & $\Ldim(\mathcal{H})$ & Littlestone dimension \\
$D$ & distribution on $X \times Y$ & $\DSdim(\mathcal{H})$ & DS dimension \\
$S = \{(x_i, y_i)\}_{i=1}^m$ & training sample & $\Ndim(\mathcal{H})$ & Natarajan dimension \\
$\risk(h) = \E_{(x,y) \sim D}[\ell(h(x),y)]$ & true risk & $\Pdim(\mathcal{F})$ & pseudodimension \\
$\emprisk(h) = \frac{1}{m}\sum \ell(h(x_i),y_i)$ & empirical risk & $\fatshat_\gamma(\mathcal{F})$ & fat-shattering dim \\
$\varepsilon$ & accuracy parameter & $\Rad_n(\mathcal{H})$ & Rademacher complexity \\
$\delta$ & confidence parameter & $\growth_\mathcal{H}(n)$ & growth function \\
\bottomrule
\end{tabular}
\end{center}

\vspace{1em}
\hfill \textit{Dhruv Gupta, Bangalore, March 2026}

\tableofcontents

\listoffigures
\addcontentsline{toc}{chapter}{List of Figures}

\listoftables
\addcontentsline{toc}{chapter}{List of Tables}

% ============================================================
\mainmatter
% ============================================================

% Chapter 1: The Objects of Learning
% Rewritten against the design-lab Lean 4 kernel (FLT_Proofs.*):
% typed definitions, full proofs isomorphic to the verified arguments,
% audit-driven corrections, and a closing set of open problems.
% concept_class = Profile G (hub node, parallax),
% inductive_bias = Profile D (negative result centerpiece),
% MDL/MML/Bayesian = Profile E (analogy triangle obstruction analysis)

\chapter{The Objects of Learning}
\label{ch:objects}

Every learning problem begins with the same question: given data drawn from an unknown source, find a rule that predicts well on future data. This chapter introduces the mathematical vocabulary for making that question precise.

The vocabulary divides into three groups of unequal importance. The first group, domain, label, concept, hypothesis, is atomic: these are the sets and functions from which everything else is built. They require only brief definitions. The second group, concept class, hypothesis space, version space, inductive bias, carries genuine mathematical content, because the relationships between these objects determine what is learnable; we prove the load-bearing facts in full. The third group, Bayesian inference, minimum description length, minimum message length, appears to be vocabulary but is actually a network of near-miss analogies, each encoding a different answer to the same question about model selection.

Every typed object in this chapter has a machine-checked counterpart, named where one exists.

\section{Atomic Vocabulary}
\label{sec:atoms}

The following objects are the type-level substrate of learning theory. They are presented together because none of them, individually, contains any mathematical surprise.

\begin{defn}[Domain, Label, Concept]
\label{def:domain-label-concept}
A \emph{domain} $X$ is a set whose elements are called \emph{instances}. A \emph{label set} $Y$ is a set of possible outputs; in binary classification, $Y = \{0, 1\}$. A \emph{concept} is a function $c : X \to Y$.\formal{FLT_Proofs.Basic.Concept}

Equivalently, when $Y = \{0,1\}$, a concept $c$ can be identified with the subset $\{x \in X : c(x) = 1\} \subseteq X$. Both perspectives, function and set, are used throughout the literature. We adopt the function view as primary, matching the typed definition $\leanname{Concept X Y} := X \to Y$, and switch to the set view when it simplifies combinatorial arguments (as in shattering, \cref{def:shatter-vc}).
\end{defn}

\begin{defn}[Hypothesis, Target Concept, Proper/Improper]
\label{def:hypothesis-target}
A \emph{hypothesis} $h : X \to Y$ is a candidate prediction rule that a learning algorithm might output; as a type it is a concept by another name, $\leanname{Hypothesis X Y} := \leanname{Concept X Y}$, and the distinction is semantic: a hypothesis is the learner's guess, a concept is ground truth. The \emph{target concept} $c^\ast \in \mathcal{C}$ is the specific concept that generated the training data, an element of the subtype $\{c : c \in \mathcal{C}\}$. Learning is \emph{proper} when the learner's hypothesis space equals the concept class, $\mathcal{H} = \mathcal{C}$, and \emph{improper} when it may use a larger space $\mathcal{H} \supseteq \mathcal{C}$.\formal{FLT_Proofs.Basic.IsProper}
\end{defn}

\begin{defn}[Realizable and Agnostic]
\label{def:realizable-objects}
The learning problem $(\mathcal{C}, \mathcal{H})$ is \emph{realizable} when every target is available to the learner, $\mathcal{C} \subseteq \mathcal{H}$.\formal{FLT_Proofs.Basic.Realizable} It is \emph{agnostic} when no such containment is assumed; the agnostic setting is precisely the \emph{absence} of a realizability hypothesis, and the learner competes against the best hypothesis in $\mathcal{H}$, whatever its error, since a perfect one may be unavailable.
\end{defn}

An \emph{alphabet} $\Sigma$ is a finite symbol set; it appears only in the formal language chapters (\Cref{ch:automata,ch:gold}) and plays no role in statistical learning theory.

\medskip

These definitions are unremarkable individually. The mathematical content begins when we ask how they compose.

% ================================================================
% CONCEPT CLASS: Hub Node, Profile G (Framework Bridge)
% Centerpiece: the node everything measures
% ================================================================

\section{The Concept Class: Where All Paradigms Converge}
\label{sec:concept-class}

\begin{defn}[Concept Class]
\label{def:concept-class}
A \emph{concept class} $\mathcal{C} \subseteq Y^X$ is a collection of concepts over a common domain; as a type, $\leanname{ConceptClass X Y} := \mathrm{Set}\,(\leanname{Concept X Y})$.\formal{FLT_Proofs.Basic.ConceptClass} The class is the primary object of study: every major question in learning theory is a question about concept classes.
\end{defn}

This definition is simple. Its importance is not. The concept class is the single object on which the chapter's deepest theorem turns: a family of apparently different demands one might place on it (learnability from few examples, finite shattering, bounded compression, vanishing correlation with noise, sub-exponential growth, finite sample size) are secretly one condition, and the class is the vertex at which they coincide (Figure~\ref{fig:concept-class-hub}).\formal{FLT_Proofs.Theorem.PAC.fundamental_theorem}

\begin{figure}[ht]
\centering
\begin{tikzpicture}[
    pivot/.style={draw, very thick, rounded corners, fill=defblue!20, font=\footnotesize\bfseries, minimum width=2.1cm, minimum height=0.95cm, align=center},
    face/.style={draw, thick, rounded corners, fill=defblue!8, font=\scriptsize, minimum width=1.85cm, minimum height=0.74cm, align=center},
    onlinenode/.style={draw, thick, rounded corners, fill=onlineblue!15, font=\scriptsize, minimum width=1.45cm, minimum height=0.62cm, align=center},
    frontier/.style={draw, dashed, rounded corners, fill=edgeorange!7, font=\scriptsize, minimum width=1.6cm, minimum height=0.58cm, align=center},
    chipgood/.style={draw, circle, thick, fill=sepgreen!25, font=\tiny, inner sep=0.4pt, minimum size=0.56cm},
    chipbad/.style={draw, circle, thick, fill=thmred!18, font=\tiny, inner sep=0.4pt, minimum size=0.56cm},
    equiv/.style={{Stealth[length=1.5mm]}-{Stealth[length=1.5mm]}, semithick, defblue},
    strictbridge/.style={-{Stealth[length=1.8mm]}, thick, onlineblue!75!black},
    fedge/.style={-{Stealth[length=1.3mm]}, dashed, edgeorange!85!black},
    elab/.style={font=\tiny, text=notegray, inner sep=1.1pt, fill=white},
    flab/.style={font=\tiny\itshape, text=edgeorange!65!black, inner sep=1pt}
]
% measurability gate
\draw[dashed, thick, notegray] (0,-0.1) ellipse (3.15cm and 2.75cm);
\node[notegray, font=\tiny] at (0,-2.35) {measurability gate};
% pivot: the loop-closing measure on the carrier
\node[pivot] (vc) at (0,0) {$\mathcal{C}$\\$\VC(\mathcal{C})<\infty$};
% verified equivalence faces (conjuncts of the Fundamental Theorem)
\node[face] (pac)  at (-2.05,-1.5) {PAC\\learnable};
\node[face] (rad)  at (-2.5, 0.95) {Rademacher\\$\to 0$};
\node[face] (grow) at (0, 2.2)     {sub-exp.\\growth};
\node[face] (comp) at (2.5, 0.95)  {compression\\(with info)};
\node[face] (samp) at (2.05,-1.5)  {finite sample\\complexity};
% two-way machine-checked equivalences
\draw[equiv] (vc) -- node[elab]{Fund.\ Thm} (pac);
\draw[equiv] (vc) -- (rad);
\draw[equiv] (vc) -- node[elab]{Sauer--Shelah} (grow);
\draw[equiv] (vc) -- node[elab]{Moran--Yehud.} (comp);
\draw[equiv] (vc) -- node[elab]{$\Theta(\VC)$} (samp);
\draw[equiv] (pac) -- (rad);
% nested verified loop: online learning
\node[onlinenode] (onl) at (-2.55,-3.05) {online\\learnable};
\node[onlinenode] (ld)  at (-0.65,-3.05) {Littlestone\\dim};
\draw[equiv, onlineblue!70!black] (onl) -- (ld);
\draw[strictbridge] (onl) to[bend left=14] node[elab, pos=0.58]{strict $\not\Leftarrow$} (pac);
% attention chips straddling the gate
\node[chipgood] (soft) at (1.95,1.6) {soft};
\node[chipbad]  (hard) at (3.0,2.45) {hard};
\node[font=\tiny, text=sepgreen!55!black, anchor=east] at (1.58,1.6) {softmax};
\node[font=\tiny, text=thmred!70!black, anchor=west] at (3.34,2.45) {argmax};
% generality frontier: literature / open / known-to-break
\node[frontier] (ds)    at (0, 3.35)    {multiclass / DS};
\node[frontier] (noinf) at (4.05, 0.95) {compression,\\no side info};
\node[frontier] (agn)   at (-4.5,-1.5)  {agnostic PAC};
\node[frontier, draw=thmred!70!black, fill=thmred!5] (part) at (-4.5, 0.5) {partial\\concepts};
\draw[fedge] (ds)    -- node[flab,right]{lit.\ '22} (grow);
\draw[fedge] (noinf) -- node[flab,above]{open} (comp);
\draw[fedge] (agn)   -- node[flab,below]{lit.} (pac);
\draw[fedge, thmred!70!black] (part) -- node[flab,above,text=thmred!75!black]{VC$\,\not\Leftrightarrow\,$PAC} (pac);
% legend
\node[font=\tiny, anchor=west, text=notegray] at (-4.5,-3.7) {%
  \textcolor{defblue}{\rule[0.45ex]{0.5cm}{0.8pt}}~verified equivalence \quad
  \textcolor{edgeorange!85!black}{\rule[0.45ex]{0.15cm}{0.8pt}\,\rule[0.45ex]{0.15cm}{0.8pt}}~literature / open \quad
  \textcolor{onlineblue!75!black}{\rule[0.45ex]{0.5cm}{0.8pt}}~strict one-way};
\end{tikzpicture}
\caption[The concept class as the vertex of an equivalence]{\textbf{The concept class as the vertex of an equivalence.} Several demands on a class of concepts (PAC learnability, finite shattering, bounded compression, vanishing Rademacher correlation, sub-exponential growth, finite sample complexity) are one property wearing several faces, and the class is the object on which they coincide. A tighter wheel for online learning sits inside, joined by a one-way bridge since the online demand implies the statistical one but not conversely; the whole rests on a measurability gate, and a dashed frontier marks coincidences expected but not yet verified, with the edge for partial concepts known to break.}
\label{fig:concept-class-hub}
\end{figure}

Each measure asks a different question about the same object:

\begin{itemize}[leftmargin=*,itemsep=2pt]
\item \textbf{VC dimension} asks: how many points can $\mathcal{C}$ shatter? (Determines PAC learnability.)
\item \textbf{Littlestone dimension} asks: how deep a mistake tree can $\mathcal{C}$ support? (Determines online learnability.)
\item \textbf{DS dimension} asks: what pseudo-cubes exist in $\mathcal{C}$'s one-inclusion hypergraph? (Determines multiclass learnability.)
\item \textbf{Teaching dimension} asks: how many examples suffice for a teacher to identify each $c \in \mathcal{C}$? (Measures cooperative communication.)
\item \textbf{Rademacher complexity} asks: how well can $\mathcal{C}$ correlate with random noise? (Controls generalization bounds.)
\end{itemize}

These 13 measures are not independent: many are related by \rel{upper\_bounds}, \rel{lower\_bounds}, and \rel{characterizes} edges in the graph. But none is reducible to another; each captures an aspect of $\mathcal{C}$'s complexity the others miss. That a single set of functions supports 13 inequivalent complexity measures is one of the structural surprises of the field.

A single theorem is what installs the concept class as the hub, in the place a paradigm or a measure might naively be expected to occupy. To state it we need the first of the measures.

\begin{defn}[Shattering and VC dimension]
\label{def:shatter-vc}
A finite set $S \subseteq X$ is \emph{shattered} by $\mathcal{C}$ if every labeling of $S$ is realized by some concept in $\mathcal{C}$:
\[
\text{$S$ shattered} \iff \forall\, f : S \to \{0,1\},\ \exists\, c \in \mathcal{C},\ \forall\, x \in S,\ c(x) = f(x).
\]\formal{FLT_Proofs.Complexity.VCDimension.Shatters}
The \emph{VC dimension} $\VC(\mathcal{C}) \in \N \cup \{\infty\}$ is the supremum of the cardinalities of shattered sets:
\[
\VC(\mathcal{C}) = \sup\{\,|S| : S \subseteq X \text{ finite and shattered by } \mathcal{C}\,\}.
\]\formal{FLT_Proofs.Complexity.VCDimension.VCDim}
\end{defn}

The full theory of VC dimension is the subject of \Cref{ch:dimensions}; here we need only the definition and one elementary closure property, used repeatedly below: a subset of a shattered set is shattered. If $T \subseteq S$ and $S$ is shattered, then any labeling of $T$ extends to a labeling of $S$, which some concept realizes; restricting to $T$ proves $T$ shattered.\formal{FLT_Proofs.Complexity.Littlestone.Shatters.subset}

\begin{thm}[VC characterization of PAC learnability]
\label{thm:vc-char}
For a concept class $\mathcal{C}$ over a measurable domain $X$ (subject to the mild measurability conditions of \Cref{ch:generalization}),
\[
\mathcal{C} \text{ is PAC learnable} \iff \VC(\mathcal{C}) < \infty.
\]\formal{FLT_Proofs.Theorem.PAC.vc_characterization}
\end{thm}

The proof occupies \Cref{ch:pac} (sufficiency, via empirical risk minimization and uniform convergence) and \Cref{ch:dimensions} (necessity, via a hard distribution on a shattered set); it is the reason the concept class is the hub of \cref{fig:concept-class-hub}. The class, through a single combinatorial invariant of itself, decides its own learnability. The equivalence extends to a five-way characterization: PAC learnability, finite VC dimension, the existence of a bounded sample compression scheme, uniform convergence of Rademacher complexity, and a finite sample-complexity bound are all equivalent.\formal{FLT_Proofs.Theorem.PAC.fundamental_theorem} We meet the compression and Rademacher faces in \Cref{ch:compression,ch:generalization}.

\begin{remark}[The equivalence is asymmetric]
\label{rmk:vc-asymmetry}
\Cref{thm:vc-char} is an ``iff'', but its two halves are proved by opposite kinds of argument. The sufficiency direction ($\VC < \infty \Rightarrow$ learnable) is constructive: it exhibits an algorithm, empirical risk minimization, and bounds its error. The necessity direction (learnable $\Rightarrow \VC < \infty$) is non-constructive: from an infinite VC dimension it produces a distribution on which every learner fails, without exhibiting any particular hard instance the learner must confront in advance. The two directions cross the boundary between combinatorics and measure theory in opposite directions, and whether the necessity direction admits a constructive proof is open (\cref{op:asymmetry}).
\end{remark}

We can now compute the invariant on the smallest interesting classes, and in doing so witness a separation that recurs throughout the book.

\begin{example}[Four concept classes of increasing complexity]
\label{ex:concept-classes}
Fix $X = \R$.
\begin{enumerate}[label=(\alph*),leftmargin=*,itemsep=3pt]
\item \textbf{Thresholds:} $\mathcal{C}_1 = \{x \mapsto \ind[x \geq \theta] : \theta \in \R\}$. Here $\VC(\mathcal{C}_1) = 1$ and $\Ldim(\mathcal{C}_1) = \infty$ (\cref{prop:threshold-vc,prop:threshold-ldim}).
\item \textbf{Intervals:} $\mathcal{C}_2 = \{x \mapsto \ind[a \leq x \leq b] : a \leq b \in \R\}$. Here $\VC(\mathcal{C}_2) = 2$: the two-point set $\{0,1\}$ is shattered (the four labelings are realized by intervals disjoint from both, containing only $0$, only $1$, or both), while no three points $x_1 < x_2 < x_3$ are, because the labeling $(1,0,1)$ would require an interval containing $x_1$ and $x_3$ but not $x_2$.
\item \textbf{Unions of $k$ intervals:} $\mathcal{C}_3^{(k)}$, with $\VC(\mathcal{C}_3^{(k)}) = 2k$ by the same argument applied $k$ times; the detailed counting is deferred to \Cref{ch:dimensions}.
\item \textbf{All measurable functions:} $\mathcal{C}_4 = \{0,1\}^X$, with $\VC(\mathcal{C}_4) = \infty$. By \cref{thm:vc-char} it is not PAC learnable; by the same token it is not online learnable nor identifiable in the limit. The No-Free-Lunch theorem (\cref{thm:nfl}) applies to it in full force.
\end{enumerate}
The thresholds deserve attention: finite VC dimension but infinite Littlestone dimension. This single class witnesses the separation between PAC and online learning (\Cref{ch:separations}). We prove both values now, because the proofs are the prototypes for the shattering and adversary arguments used throughout.
\end{example}

\begin{prop}[Thresholds have VC dimension one]
\label{prop:threshold-vc}
The class $\mathcal{C}_1 = \{x \mapsto \ind[x \geq \theta] : \theta \in \R\}$ has $\VC(\mathcal{C}_1) = 1$.
\end{prop}

\begin{proof}
\emph{At least one.} The singleton $\{x_0\}$ is shattered: the labeling $x_0 \mapsto 1$ is realized by any $\theta \leq x_0$, and $x_0 \mapsto 0$ by any $\theta > x_0$.

\emph{Not two.} Let $x_1 < x_2$. The labeling $f(x_1) = 1$, $f(x_2) = 0$ asks for a threshold $\theta$ with $x_1 \geq \theta$ and $x_2 < \theta$, hence $x_2 < \theta \leq x_1$, contradicting $x_1 < x_2$. So no two-point set is shattered. By \cref{def:shatter-vc}, $\VC(\mathcal{C}_1) = 1$.
\end{proof}

\begin{prop}[Thresholds have infinite Littlestone dimension]
\label{prop:threshold-ldim}
The class $\mathcal{C}_1$ admits, against any online learner, a realizable sequence forcing an unbounded number of mistakes; hence $\Ldim(\mathcal{C}_1) = \infty$.
\end{prop}

\begin{proof}
We exhibit an adversary that forces a mistake every round while keeping the transcript consistent with some threshold. The adversary maintains an open interval $(a,b)$ of still-admissible thresholds, initialized to $(0,1)$. At each round it presents the midpoint $x = (a+b)/2$ and reads the learner's prediction $\hat{y} = \ind[x \geq \hat\theta]$. It then declares the true label to be $1 - \hat{y}$ and updates:
\begin{itemize}[nosep]
\item if $\hat{y} = 1$, declare label $0$ (meaning $x < \theta$, i.e.\ $\theta > x$) and set $(a,b) \leftarrow (x, b)$;
\item if $\hat{y} = 0$, declare label $1$ (meaning $x \geq \theta$, i.e.\ $\theta \leq x$) and set $(a,b) \leftarrow (a, x)$.
\end{itemize}
After the update the interval $(a,b)$ remains nonempty (the midpoint strictly separates it), and \emph{every} $\theta$ in it is consistent with the entire transcript so far: each declared label $\ind[x_i \geq \theta]$ is correct for such $\theta$ by construction. Thus the sequence is realizable by the class at every finite stage, yet the learner errs at every round. Since the number of rounds is unbounded, the realizable mistake bound is infinite, and therefore $\Ldim(\mathcal{C}_1) = \infty$.\formal{FLT_Proofs.Theorem.Separation.adversary_from_shatters}
\end{proof}

\Cref{prop:threshold-vc,prop:threshold-ldim} together discharge the claim of \cref{ex:concept-classes}(a): the same class is trivially PAC learnable (one example pins the threshold up to $\varepsilon$) yet admits no finite online mistake bound. Because \cref{prop:threshold-ldim} forces unbounded mistakes against every learner, the gap is a genuine separation; we return to it in \Cref{ch:separations}.

% ================================================================
% HYPOTHESIS SPACE: Profile B-light
% Centerpiece: proper vs improper tension
% ================================================================

\section{Hypothesis Spaces and the Proper--Improper Divide}
\label{sec:hypothesis-space}

\begin{defn}[Hypothesis Space]
\label{def:hypothesis-space}
A \emph{hypothesis space} $\mathcal{H} \subseteq Y^X$ is the set of functions available to the learner as candidate outputs. As a type it is structurally identical to a concept class, $\leanname{HypothesisSpace X Y} := \mathrm{Set}\,(\leanname{Concept X Y})$, and semantically distinct: $\mathcal{C}$ is the ground-truth collection, $\mathcal{H}$ is what the learner may return. When $\mathcal{H} = \mathcal{C}$, learning is \emph{proper}; when $\mathcal{H} \supsetneq \mathcal{C}$, learning is \emph{improper}.
\end{defn}

The proper--improper distinction carries real weight. Its computational consequences are among the sharpest in the theory:

\begin{enumerate}[leftmargin=*,itemsep=3pt]
\item \textbf{Improper learning can be exponentially cheaper.} There exist concept classes where proper learning requires exponentially more samples or computation than improper learning~\cite{PittValiant1988}. The graph records this as a \rel{does\_not\_imply} edge from \nde{pac\_learning} to \nde{exact\_learning} (\Cref{ch:separations}).

\item \textbf{Computational hardness often depends on properness.} Cryptographic hardness results (Kearns--Valiant, 1994) typically establish hardness of \emph{proper} learning; improper learners may circumvent the barrier by using richer representations~\cite{KearnsValiant1994}.

\item \textbf{The characterization is representation-independent.} \Cref{thm:vc-char} holds regardless of whether learning is proper or improper. What decides learnability is $\VC(\mathcal{C})$, the VC dimension of the \emph{target} class; properties of the learner's hypothesis space $\mathcal{H}$ leave the verdict untouched. The learner's freedom to range outside $\mathcal{C}$ moves the cost while leaving possibility fixed.
\end{enumerate}

The simplest object built from $\mathcal{H}$ and data together is the set of survivors.

\begin{defn}[Version Space]
\label{def:version-space}
Given a hypothesis space $\mathcal{H}$ and a sample $S = \{(x_i, y_i)\}_{i=1}^m$, the \emph{version space} is
\[
\mathrm{VS}_{\mathcal{H},S} = \{h \in \mathcal{H} : h(x_i) = y_i \text{ for all } i\},
\]
the set of hypotheses consistent with all observed data.\formal{FLT_Proofs.Complexity.GameInfra.versionSpace} The verified development carries a bundled $\leanname{VersionSpace}$ structure packaging $\mathcal{H}$, the data, and the consistency obligations.\formal{FLT_Proofs.Basic.VersionSpace}
\end{defn}

Version-space elimination (returning any element of $\mathrm{VS}_{\mathcal{H},S}$) is among the most natural learning algorithms. Its basic behavior is captured by four propositions, each a one-line consequence of the definition and each a prototype for the convergence arguments of \Cref{ch:gold}.

\begin{prop}[Soundness, persistence, and monotonicity of the version space]
\label{prop:version-space}
Fix $\mathcal{H}$, a target $c \in \mathcal{H}$, and a sample $S$ labeled by $c$. Then
\begin{enumerate}[label=(\roman*),leftmargin=*,itemsep=2pt]
\item $\mathrm{VS}_{\mathcal{H},S} \subseteq \mathcal{H}$;\formal{FLT_Proofs.Complexity.GameInfra.versionSpace_subset}
\item $c \in \mathrm{VS}_{\mathcal{H},S}$;\formal{FLT_Proofs.Complexity.GameInfra.target_in_versionSpace}
\item $\mathrm{VS}_{\mathcal{H},\,S \cup \{(x,y)\}} \subseteq \mathrm{VS}_{\mathcal{H},S}$ for any further labeled example $(x,y)$;\formal{FLT_Proofs.Complexity.GameInfra.versionSpace_append}
\item $\Ldim(\mathrm{VS}_{\mathcal{H},S}) \leq \Ldim(\mathcal{H})$.\formal{FLT_Proofs.Complexity.GameInfra.ldim_versionSpace_le}
\end{enumerate}
\end{prop}

\begin{proof}
(i) An element of $\mathrm{VS}_{\mathcal{H},S}$ is by definition a member of $\mathcal{H}$ satisfying the consistency constraints; forgetting the constraints gives membership in $\mathcal{H}$.

(ii) The target $c$ lies in $\mathcal{H}$ by hypothesis, and since $S$ is labeled by $c$ we have $c(x_i) = y_i$ for every $i$; so $c$ meets both defining conditions and lies in the version space. In particular the version space is never empty in the realizable case.

(iii) Adding a labeled example $(x,y)$ adjoins the constraint $h(x) = y$ to the defining conjunction. Any $h$ satisfying the larger conjunction satisfies the smaller one, so the consistent set can only shrink.

(iv) By (i), $\mathrm{VS}_{\mathcal{H},S} \subseteq \mathcal{H}$. The Littlestone dimension is monotone under inclusion: a mistake tree shattered by a subclass is shattered, unchanged, by the superclass (each labeling realized inside the subclass is realized inside the superclass). Taking the supremum of shattered-tree depths over the smaller class gives a value no larger than over the larger one.
\end{proof}

Item (iii) is the engine of identification in the limit: as data accumulates, the version space contracts monotonically, and the question of \Cref{ch:gold} is whether it contracts to the target. Item (iv) says that this contraction never makes the residual problem harder in the online sense, a fact the Standard Optimal Algorithm exploits to achieve the optimal mistake bound.

The trouble with version-space elimination is computational. For most hypothesis classes of interest, the version space cannot be efficiently represented, enumerated, or sampled from. Two dual invariants quantify aspects of this difficulty. The \emph{eluder dimension}~\cite{RussoVanRoy2013} measures how long an adversary can keep producing examples on which consistent hypotheses still disagree, controlling exploration--exploitation tradeoffs. The \emph{star number} measures a complementary, purely combinatorial quantity.

\begin{defn}[Star number]
\label{def:star-number}
The \emph{star number} of $\mathcal{C}$ is the largest $d$ for which there is a set $S = \{x_1,\dots,x_d\}$ such that each point can be isolated: for every $x \in S$ there is a concept $c_x \in \mathcal{C}$ with $c_x(x) = 1$ and $c_x(y) = 0$ for all other $y \in S$.\formal{FLT_Proofs.Complexity.VCDimension.StarNumber}
\end{defn}

The star number is never smaller than the VC dimension, and the proof is immediate from shattering.

\begin{prop}[Star number dominates VC dimension]
\label{prop:star-ge-vc}
For every concept class, $\VC(\mathcal{C}) \leq \mathrm{star}(\mathcal{C})$.
\end{prop}

\begin{proof}
Let $S$ be a shattered set of size $d = \VC(\mathcal{C})$ (if $\VC(\mathcal{C}) = \infty$, apply the argument to shattered sets of every size). For each $x \in S$ the single labeling ``$x \mapsto 1$, all other points of $S$ $\mapsto 0$'' is one of the labelings realized by shattering, so there is a concept isolating $x$ within $S$. Hence $S$ witnesses the star condition, and $\mathrm{star}(\mathcal{C}) \geq |S| = d$.
\end{proof}

\begin{warning}[The reverse inequality fails, and so does the folklore equivalence]
\label{rmk:star-not-equiv}
It is tempting to believe that finiteness of the star number is \emph{equivalent} to finiteness of the VC dimension. It is not. Consider the class of singletons on an infinite domain, $\mathcal{C} = \{\ind[\cdot = a] : a \in X\}$. Its VC dimension is $1$: a single point $\{a\}$ is shattered ($\ind[\cdot=a]$ labels it $1$, any other singleton labels it $0$), but no pair $\{a,b\}$ is, because labeling both points $1$ would require a single concept true at $a$ and $b$, which no singleton is. Yet its star number is infinite: for \emph{any} finite set $S$ and any $x \in S$, the concept $\ind[\cdot = x]$ isolates $x$ within $S$, so every finite $S$ is a star set. Thus $\VC = 1$ while $\mathrm{star} = \infty$. Finiteness of the star number is therefore strictly stronger than finiteness of the VC dimension; the implication runs one way only (\cref{prop:star-ge-vc}). The star number governs the label complexity of \emph{active} learning, where this gap is exactly the point~\cite{HannekeYang2015}; we return to it in \Cref{ch:dimensions}. Identifying which finite-VC classes have finite star number is open (\cref{op:star}).
\end{warning}

% ================================================================
% INDUCTIVE BIAS: Profile D (Negative Result)
% Centerpiece: NFL theorem makes this NECESSARY
% ================================================================

\section{Inductive Bias and the No-Free-Lunch Barrier}
\label{sec:inductive-bias}

\begin{defn}[Inductive Bias]
\label{def:inductive-bias}
The \emph{inductive bias} of a learner is the set of assumptions, explicit or implicit, that it uses to generalize from training data to unseen instances.
\end{defn}

This informal definition is the one usually given, and it suggests that inductive bias is not a single mathematical object. That suggestion is half right. The \emph{content} of a bias does change across paradigms, as the table below records. But the \emph{form} is uniform, and can be typed: a bias is a hypothesis space together with a real-valued preference ranking on concepts, structured so that nothing outside the space is preferred to anything inside it.

\begin{defn}[Inductive bias, typed]
\label{def:inductive-bias-typed}
An \emph{inductive bias} over $X, Y$ is a triple $(\mathcal{H}, \mathrm{pref}, \pi)$ where $\mathcal{H} \subseteq Y^X$ is a hypothesis space, $\mathrm{pref} : Y^X \to \R$ is a preference score (lower is more preferred), and $\pi$ is the constraint that every concept outside $\mathcal{H}$ is scored at least as high as every concept inside it: $h \notin \mathcal{H} \Rightarrow \mathrm{pref}(h') \leq \mathrm{pref}(h)$ for all $h' \in \mathcal{H}$.\formal{FLT_Proofs.Basic.InductiveBias}
\end{defn}

This single structure subsumes the standard readings. Reading $\mathrm{pref}(h) = -\log P(h)$ recovers the Bayesian prior; reading $\mathrm{pref}(h)$ as description length recovers minimum description length; reading it as the index of the first class in a nested hierarchy containing $h$ recovers structural risk minimization. The differences between these readings are real (\cref{sec:model-selection}), but they are differences of \emph{which} preference score, inside one type.

\begin{center}\small
\begin{tabularx}{\textwidth}{@{}l l X@{}}
\toprule
\textbf{Paradigm} & \textbf{Inductive bias formalized as} & \textbf{Graph edge} \\
\midrule
PAC & Restriction of $\mathcal{H}$ to finite VC dimension & \rel{used\_in\_proof} by \nde{nfl\_theorem} \\
Bayesian & Prior $P(h)$ over hypotheses & \rel{analogy} from \nde{bayesian\_inference} \\
SRM & Nested class hierarchy $\mathcal{H}_1 \subset \mathcal{H}_2 \subset \cdots$ & \rel{analogy} from \nde{srm} \\
Computational & Cryptographic assumption restricting efficient search & \rel{requires\_assumption} by \nde{computational\_hardness} \\
\bottomrule
\end{tabularx}
\end{center}

Inductive bias is a \emph{mathematical necessity}, forced on every learner, and the reason is the No-Free-Lunch theorem. We prove it through its combinatorial engine, then read off the classical statement and the necessity of bias as corollaries.

\begin{lemma}[Shattering forces a hard concept]
\label{lem:nfl-engine}
Let $T \subseteq X$ be shattered by $\mathcal{C}$ with $|T| > 2m$. For any $m$ instances $x_1,\dots,x_m$ and any hypothesis $h : X \to \{0,1\}$, there is a concept $c \in \mathcal{C}$ such that
\begin{enumerate}[label=(\roman*),leftmargin=*,nosep]
\item $c$ agrees with $h$ on every sample point that lies in $T$: $c(x_i) = h(x_i)$ whenever $x_i \in T$; and
\item $c$ disagrees with $h$ on more than a quarter of $T$: $\bigl|\{x \in T : c(x) \neq h(x)\}\bigr| > |T|/4$.
\end{enumerate}\formal{FLT_Proofs.Complexity.Generalization.nfl_per_sample_shattered}
\end{lemma}

\begin{proof}
Define an adversarial labeling $f : T \to \{0,1\}$ that copies $h$ on the seen points and flips it on the rest:
\[
f(x) = \begin{cases} h(x) & \text{if } x \in \{x_1,\dots,x_m\},\\ 1 - h(x) & \text{otherwise.}\end{cases}
\]
Because $T$ is shattered, some $c \in \mathcal{C}$ realizes $f$ on $T$, i.e.\ $c(x) = f(x)$ for all $x \in T$. On a seen point $x_i \in T$ we have $c(x_i) = f(x_i) = h(x_i)$, giving (i). On an unseen point $x \in T$ we have $c(x) = f(x) = 1 - h(x) \neq h(x)$, so $c$ disagrees with $h$ on the entire unseen part of $T$. The unseen part has size at least $|T| - m$, since the $m$ sample points can cover at most $m$ elements of $T$. Finally, $|T| > 2m$ gives $|T| - m > |T|/2 > |T|/4$, so the disagreement count exceeds $|T|/4$, which is (ii).
\end{proof}

\begin{thm}[No Free Lunch {\cite{Wolpert1996}}]
\label{thm:nfl}
Let $A$ be any learning algorithm that observes $m$ labeled examples, and suppose $\mathcal{C}$ shatters some set $T$ with $|T| \geq 2m$. Then there is a target concept $c^\ast \in \mathcal{C}$ and the uniform distribution $D$ on $T$ such that
\[
\E_{S \sim D^m}\bigl[\risk_D(A(S))\bigr] \geq \tfrac{1}{4},
\]
where $S$ is labeled by $c^\ast$. In particular, no learner does better than this on its worst target: averaging cannot rescue what shattering has spoiled.
\end{thm}

\begin{proof}
Take $|T| = 2m$ and let $D$ be uniform on $T$. Shattering realizes all $2^{|T|}$ labelings of $T$ inside $\mathcal{C}$; let $c$ range uniformly over these $2^{2m}$ concepts. It suffices to bound the average risk over $c$ from below by $1/4$, since the maximum over targets is at least the average. Writing the risk as an average of disagreement indicators over $T$ and exchanging the (finite) expectations,
\[
\E_{c}\,\E_{S \sim D^m}\bigl[\risk_D(A(S))\bigr]
= \frac{1}{|T|}\sum_{x \in T} \E_{S}\,\E_{c}\bigl[\ind[A(S)(x) \neq c(x)]\bigr].
\]
Fix a draw $S$ of $m$ points and a point $x \in T$ not among them. The learner's output $A(S)$ depends on $c$ only through the labels at the seen points, so $A(S)(x)$ is independent of $c(x)$; and $c(x)$ is a uniform bit. Hence the inner expectation is exactly $\tfrac{1}{2}$ for every unseen $x$. The unseen points number at least $|T| - m = m = |T|/2$, so the average is at least $\tfrac{1}{|T|}\cdot \tfrac{|T|}{2}\cdot\tfrac{1}{2} = \tfrac14$. Some target $c^\ast$ therefore attains $\E_S[\risk_D(A(S))] \geq \tfrac14$. (The per-sample form of this argument, which builds the hard target adversarially in place of averaging over targets, is \cref{lem:nfl-engine}.)
\end{proof}

\begin{cor}[Inductive bias is mandatory]
\label{cor:bias-mandatory}
A concept class of infinite VC dimension is not PAC learnable; in particular the class $\{0,1\}^X$ of all binary functions on an infinite domain is not learnable. Consequently every learnable class is a proper subset of $\{0,1\}^X$, and the choice of which subset is exactly the inductive bias.
\end{cor}

\begin{proof}
If $\VC(\mathcal{C}) = \infty$, then $\mathcal{C}$ shatters sets of every size, so for each sample budget $m$ it shatters some $T$ with $|T| \geq 2m$. \Cref{thm:nfl} then supplies, for every learner, a target and distribution on which the expected error is at least $1/4$; taking $\varepsilon, \delta < 1/4$ shows the PAC criterion cannot be met. Hence infinite VC dimension precludes PAC learning, the necessity half of \cref{thm:vc-char}, here proved directly.\formal{FLT_Proofs.Complexity.Generalization.vcdim_infinite_not_pac} Since $\{0,1\}^X$ has infinite VC dimension, it is unlearnable, so any learnable class omits some functions; the omission \emph{is} the assumption that those functions are not the target.
\end{proof}

\begin{remark}[No-Free-Lunch as a boundary of askability]
\label{rmk:nfl-pl}
\Cref{cor:bias-mandatory} is more than a negative result; it relocates inductive bias from the engineering of a learner to the precondition of the question. Before a class is fixed, ``which functions are efficiently learnable?'' has no nontrivial answer, because none are. Inductive bias is what makes the learnability question askable at all. The graph encodes this as a \rel{requires\_assumption} dependency from \nde{inductive\_bias} into \nde{computational\_hardness}: the hardness theory presupposes the restriction that No-Free-Lunch shows to be unavoidable.
\end{remark}

% ================================================================
% MODEL SELECTION: Profile E (Analogy Triangle)
% Centerpiece: MDL/MML/Bayesian near-miss network
% ================================================================

\section{Three Faces of Model Selection}
\label{sec:model-selection}

Bayesian inference, minimum description length, and minimum message length are often described informally as ``the same idea.'' They are not. Each formalizes a different answer to the question ``which hypothesis should I prefer?''; the differences between them are more instructive than any individual definition. Each is, in the sense of \cref{def:inductive-bias-typed}, a choice of preference score; the analogies among them are analogies among scores, and the obstructions are precisely where the scores fail to agree.

\subsection{The Definitions}

\begin{defn}[Bayesian Inference]
\label{def:bayesian}
Given a prior $P(h)$ over $\mathcal{H}$ and a likelihood $P(S \mid h)$, the posterior is $P(h \mid S) \propto P(S \mid h) \cdot P(h)$. The Bayesian learner outputs $\hat{h} = \arg\max_{h} P(h \mid S)$ (MAP) or the full posterior $P(\cdot \mid S)$. As a type, a Bayesian learner bundles a prior with the posterior-update rule and converts to an ordinary batch learner by MAP read-off.\formal{FLT_Proofs.Learner.Bayesian.BayesianLearner}
\end{defn}

\begin{defn}[Minimum Description Length {\cite{Rissanen1978}}]
\label{def:mdl}
Select $h$ minimizing $L(h) + L(S \mid h)$, where $L(\cdot)$ denotes description length under a fixed prefix-free code. The first term penalizes complexity; the second penalizes misfit.
\end{defn}

\begin{defn}[Minimum Message Length {\cite{WallaceBoulton1968}}]
\label{def:mml}
Select $h$ minimizing the total message length $\ell(h) + \ell(S \mid h)$ required to transmit the hypothesis and then the data, under a two-part coding scheme. Wallace--Freeman (1987): $\ell(h) = -\log P(h) + \frac{1}{2}\log |F(h)| + \text{const}$, where $F(h)$ is the Fisher information.
\end{defn}

\subsection{The Analogy Network}

These three methods are usually treated as one idea. The deeper structure, made precise by the chapter's compression characterization, is that they are approximations to a single machine-checked equivalence: a class is learnable exactly when its data compress to a bounded summary that reconstructs every consistent labelling (Figure~\ref{fig:model-selection-triangle}).\formal{FLT_Proofs.Theorem.PAC.fundamental_vc_compression} The pairwise analogies among the methods remain genuine but unproven, each blocked by a different obstruction.

\begin{figure}[ht]
\centering
\begin{tikzpicture}[
    spine/.style={draw, very thick, rounded corners, fill=defblue!14, font=\small, minimum width=2.3cm, minimum height=0.85cm, align=center},
    sub/.style={draw, thick, rounded corners, fill=defblue!5, font=\scriptsize, minimum width=2.1cm, minimum height=0.7cm, align=center},
    bridge/.style={draw, thick, rounded corners, fill=sepgreen!12, font=\scriptsize, minimum width=2.0cm, minimum height=0.7cm, align=center},
    method/.style={draw, thick, rounded corners, fill=analogypurple!10, font=\small\bfseries, minimum width=1.7cm, minimum height=0.7cm, align=center},
    vedge/.style={-{Stealth[length=1.7mm]}, semithick, defblue},
    centerpiece/.style={{Stealth[length=2.2mm]}-{Stealth[length=2.2mm]}, line width=1.3pt, defblue},
    gedge/.style={-{Stealth[length=1.7mm]}, semithick, sepgreen!70!black},
    dedge/.style={-{Stealth[length=1.5mm]}, dashed, edgeorange!85!black},
    obs/.style={dashed, analogypurple!55, semithick},
    openedge/.style={-{Stealth[length=1.5mm]}, dashed, thmred!75!black},
    elab/.style={font=\tiny, text=notegray, inner sep=1.3pt, fill=white},
    olab/.style={font=\tiny, text=analogypurple!70!black, inner sep=1.5pt, fill=white, align=center}
]
% spine: the machine-checked equivalence
\node[spine] (vc)   at (-2.3,0) {finite VC\\dimension};
\node[spine] (comp) at ( 2.3,0) {bounded compression\\(with side info)};
\draw[centerpiece] (vc) -- node[elab,above=1pt]{learnable $\Leftrightarrow$ compressible} node[elab,below=1pt]{Moran--Yehudayoff} (comp);
% side-information axis: proven vs open asymmetry
\node[sub] (noinf) at (-2.3,-2.0) {compression\\(no side info)};
\draw[vedge]    ([xshift=-4pt]noinf.north) -- node[elab,left]{VC $\le 2(k{+}1)^2$} ([xshift=-4pt]vc.south);
\draw[openedge] ([xshift=4pt]vc.south) -- node[elab,right]{$O(d)$? open} ([xshift=4pt]noinf.north);
% PAC-Bayes bridge (the one proven cross-edge)
\node[bridge] (pb) at (0,1.7) {prior-weighted\\complexity};
\draw[gedge] (pb) -- node[elab,right=1pt]{PAC-Bayes} (vc);
% model-selection principles as approximating instances
\node[method] (bayes) at (-3.3,3.5) {Bayesian};
\node[method] (mdl)   at ( 0,  3.5) {MDL};
\node[method] (mml)   at ( 3.3,3.5) {MML};
\draw[vedge] (bayes) -- node[elab,left=1pt]{MAP $\to$ batch} (pb);
\draw[dedge] (mdl) -- node[elab,right]{description length} (comp);
\draw[dedge] (mml) to[bend left=8] node[elab,right]{two-part code} (comp);
% obstruction arcs: no equivalence theorem exists
\draw[obs] (bayes) -- node[olab]{type\\mismatch} (mdl);
\draw[obs] (mdl) -- node[olab]{missing\\witness} (mml);
\draw[obs] (bayes) to[bend left=18] node[olab]{missing equivalence witness} (mml);
\node[font=\scriptsize\itshape, text=notegray] at (0,5.0) {model-selection principles (each a choice of inductive bias)};
% legend
\node[font=\tiny, anchor=west, text=notegray] at (-4.6,-2.95) {%
  \textcolor{defblue}{\rule[0.45ex]{0.5cm}{1pt}}~machine-checked \quad
  \textcolor{edgeorange!85!black}{\rule[0.45ex]{0.15cm}{0.8pt}\,\rule[0.45ex]{0.15cm}{0.8pt}}~literature \quad
  \textcolor{thmred!75!black}{\rule[0.45ex]{0.15cm}{0.8pt}\,\rule[0.45ex]{0.15cm}{0.8pt}}~open};
\end{tikzpicture}
\caption[The compression core of model selection]{\textbf{The compression core of model selection.} Bayesian inference, minimum description length, and minimum message length are usually treated as one idea. The figure separates them around the fact the chapter proves: a class is learnable exactly when its data admit a bounded summary reconstructing every consistent labelling. That equivalence is the spine; the three principles hang off it as approximations, with every link among them unproven. Whether the auxiliary tag can be removed while keeping the summary short is open (\cref{op:compression-objects}).}
\label{fig:model-selection-triangle}
\end{figure}

\subsection{Why They Are Not the Same}

\begin{obstbox}
\textbf{Bayesian $\leftrightarrow$ MML:} obstruction type \emph{missing\_equivalence\_witness}. MML's two-part code corresponds to Bayesian MAP estimation when the Fisher information correction term is constant. But for models with varying curvature in parameter space, MML departs from Bayesian posteriors. No theorem establishes their equivalence beyond the constant-Fisher case.
\end{obstbox}

\begin{obstbox}
\textbf{MML $\leftrightarrow$ MDL:} obstruction type \emph{missing\_equivalence\_witness}. Both minimize a two-part code length. But MML uses Shannon-optimal coding relative to a prior (subjective), while MDL uses a universal code based on Kolmogorov complexity (objective). In the limit of infinite data, both converge to the true model under standard regularity conditions, but their finite-sample behavior differs.
\end{obstbox}

\begin{obstbox}
\textbf{MDL $\leftrightarrow$ Bayesian:} obstruction type \emph{type\_mismatch}. MDL uses description lengths (integers under a code); Bayesian inference uses probability densities (reals under a measure). The connection $L(h) = -\log P(h)$ identifies description lengths with log-probabilities, but this identification breaks down when the code is not prefix-free or the prior is improper.
\end{obstbox}

That three different formalizations of ``prefer the simpler hypothesis'' are \emph{not} equivalent, and that the reasons for non-equivalence are themselves different (missing witness, missing witness, and type mismatch), is a structural feature of learning theory that no single definition can convey. It is visible only in the graph.

One bridge across this network is an actual theorem. The preference score of a Bayesian or MDL learner controls generalization through the PAC-Bayes bound: for any prior fixed before seeing the data and any posterior, the expected risk is bounded by the empirical risk plus a complexity term governed by the $\KL$ divergence from prior to posterior~\cite{McAllester1999}, proved here for finite hypothesis spaces.\formal{FLT_Proofs.Theorem.PACBayes.pac_bayes_finite} The obstruction boxes concern the three \emph{procedures}, which are not interchangeable; PAC-Bayes concerns the \emph{quantity} they share, a prior-weighted complexity, and shows it bounds error.

\subsection{Compression, Occam, and the Fundamental Theorem}

The MDL idea (prefer the hypothesis that compresses the data) connects directly to \cref{thm:vc-char}. A \emph{sample compression scheme} of size $k$ reconstructs any consistent labeling from a subsequence of at most $k$ of its own examples. Compression and finite VC dimension are equivalent.

\begin{thm}[Compression characterizes finite VC dimension {\cite{MoranYehudayoff2016}}]
\label{thm:compression-char}
A concept class has finite VC dimension if and only if it admits a sample compression scheme of bounded size (with side information):
\[
\VC(\mathcal{C}) < \infty \iff \exists\, k,\ \mathcal{C} \text{ has a size-}k\text{ compression scheme with side information.}
\]\formal{FLT_Proofs.Theorem.PAC.fundamental_vc_compression}
\end{thm}

We prove this in \Cref{ch:compression}; here it closes the circle. Through \cref{thm:vc-char,thm:compression-char}, ``prefer the hypothesis that compresses'' and ``prefer the learnable hypothesis'' are the same condition on the class. The qualification ``with side information'' carries weight: whether every class of VC dimension $d$ admits a compression scheme of size $O(d)$ \emph{without} side information is the Littlestone--Warmuth conjecture, still open (\cref{op:compression-objects}).\formal{FLT_Proofs.Complexity.Structures.Open_NoInfoCompressionStrengthening}

\subsection{Structural Risk Minimization}

\begin{defn}[Structural Risk Minimization {\cite{Vapnik1998}}]
\label{def:srm}
Given a nested sequence of hypothesis classes $\mathcal{H}_1 \subset \mathcal{H}_2 \subset \cdots$ with $\VC(\mathcal{H}_k) = d_k$, select $k^\ast$ minimizing $\emprisk_{\mathcal{H}_k}(h) + \mathrm{penalty}(d_k, m, \delta)$.
\end{defn}

SRM is the operational realization of inductive bias within VC theory: instead of choosing a single $\mathcal{H}$, choose a hierarchy indexed by complexity, and let the data select the level. In the language of \cref{def:inductive-bias-typed}, SRM's preference score is the index $k$ of the first class containing $h$; the nesting guarantees the monotonicity condition $\pi$. The graph records an \rel{analogy} edge from \nde{srm} to \nde{inductive\_bias} with a \emph{type\_mismatch} obstruction, because SRM constructs a bias by structural nesting while ``inductive bias'' as a concept is broader than any one construction.

\begin{defn}[Algorithmic Probability {\cite{Solomonoff1964}}]
\label{def:alg-prob}
The \emph{Solomonoff prior} assigns to each string $x$ the probability $M(x) = \sum_{p : U(p)=x} 2^{-|p|}$, summing over all programs $p$ that produce $x$ on a universal Turing machine $U$. This is the ``prior that assigns higher probability to simpler explanations'' in the most literal possible sense.
\end{defn}

Algorithmic probability connects to Kolmogorov complexity through the coding theorem, $K(x) = -\log M(x) + O(1)$, and provides the information-theoretic foundation for MDL. Its detailed treatment belongs to \Cref{ch:compression}, where it interacts with compression schemes and Occam's razor; the explicit constant in the coding theorem is itself a target for formalization (\cref{op:solomonoff}).

\section{What This Chapter Established}
\label{sec:ch1-summary}

The base vocabulary of learning theory is small: domain, label, concept, concept class, hypothesis space, version space. Each is a type, and each typed object has a machine-checked counterpart. At this level alone, four structural features are already visible.

\begin{enumerate}[leftmargin=*,itemsep=3pt]
\item The concept class is a hub node: the same object is measured by 13 inequivalent complexity measures, and one of them, the VC dimension, decides learnability outright (\cref{thm:vc-char}). The multiplicity earns its keep: each of the 13 reflects a genuinely different aspect of a single object.

\item Inductive bias is not optional. The No-Free-Lunch theorem (\cref{thm:nfl}), proved here through its shattering engine, converts bias from a design choice into a precondition for the learnability question to be askable at all (\cref{cor:bias-mandatory}).

\item Negative results are content, and rigor catches folklore. The asserted equivalence between finiteness of the star number and of the VC dimension is false, with singletons as an explicit witness (\cref{rmk:star-not-equiv}); the true relationship is a one-sided inequality (\cref{prop:star-ge-vc}).

\item The three standard model-selection principles, Bayesian, MDL, and MML, form a triangle of near-miss analogies whose obstructions differ pair by pair, while a genuine theorem (PAC-Bayes) and a genuine equivalence (compression $\leftrightarrow$ finite VC, \cref{thm:compression-char}) sit nearby and are often confused with them.
\end{enumerate}

These patterns recur throughout the book: hub nodes, negative results as content, analogy obstructions, and the discipline of separating what is proved from what is merely said.

% ================================================================
% TRANSFORMER CONNECTION (retrofit): attention as a concept class,
% well-behavedness, and the soft/hard = Borel/analytic boundary.
% ================================================================

\section{The Objects of a Transformer}
\label{sec:ch1-transformer}

The vocabulary of this chapter is not confined to the classical examples. A single softmax-attention head, read as a binary predictor, is a concept in the sense of \cref{def:domain-label-concept}, a function $X \to \{0,1\}$, and letting its parameters vary produces a concept class in the sense of \cref{def:concept-class}. The question returns with a sharp answer for the central object of modern machine learning: \emph{is the attention class a legitimate concept class, one to which the VC characterization applies?} Answering it requires exactly the measurability that \cref{thm:vc-char} deferred to \Cref{ch:generalization}; this section makes those ``mild measurability conditions'' concrete and shows them non-vacuous.

\begin{defn}[The softmax-attention concept class]
\label{def:attention-concept-class}
Fix an embedding dimension $d$ and a context of $n_K$ key--value pairs.  Take the domain to be the query space $X = \R^d$.  For keys $K = (k_1,\dots,k_{n_K})$ and values $V = (v_1,\dots,v_{n_K})$, the head maps a query $x$ to the readout $r_{K,V}(x) = \sum_{i=1}^{n_K} \mathrm{softmax}(\langle x, k_i\rangle)_i\, v_i$, and thresholding at zero gives a concept
\[
  c_{K,V}(x) \;=\; \ind[\,r_{K,V}(x) > 0\,] \;:\; \R^d \to \{0,1\}.
\]
The \emph{softmax-attention concept class} is $\mathcal{A}_{d,n_K} = \{\, c_{K,V} : K, V \,\}$, the range of the parametrized family.\formal{TLT.Bridge.softAttentionConcept}
\end{defn}

So far this is only a particular concept class among many. Its first non-trivial property is that it is a \emph{well-behaved} object: the measurability hypothesis that \cref{thm:vc-char} assumes is satisfied.

\begin{prop}[Attention is a well-behaved object]
\label{prop:attention-wellbehaved}
For every $d$ and $n_K$, the class $\mathcal{A}_{d,n_K}$ is a well-behaved VC-measurability target: the bad events over which the uniform-convergence argument of \Cref{ch:generalization} quantifies are measurable, so $\mathcal{A}_{d,n_K}$ is a bona-fide concept class for the VC characterization.\formal{TLT.Bridge.softAttention_wellBehaved}
\end{prop}

\begin{proof}
A parametrized concept family is a well-behaved VC-measurability target as soon as the evaluation map $(K,V,x) \mapsto c_{K,V}(x)$ is jointly measurable.  Here that map factors through the readout $(K,V,x) \mapsto r_{K,V}(x)$, which is a finite sum of products of softmax weights and values; each softmax weight is a continuous function of its arguments, so the readout is jointly continuous and therefore Borel measurable.  The concept is the indicator of the open condition $\{r_{K,V}(x) > 0\}$, whose preimage is the measurable set $\{(K,V,x) : 0 < r_{K,V}(x)\}$.  Joint measurability follows, and with it well-behavedness.  The base case is the same statement for the simplest concept class of all: the singleton-indicator class on a measurable set is well-behaved,\formal{TLT.Tame.singletonClassOn_wellBehavedVCMeasTarget} and attention's well-behavedness is the parametrized lift of that fact.
\end{proof}

The proof used continuity of the softmax in an essential way.  Far from a cosmetic change, removing it carries the attention class clear across a boundary in descriptive set theory.

\begin{prop}[Softmax is the boundary of measurability]
\label{prop:soft-hard-borel}
Assume the existence of an analytic set of reals that is not Borel (a standard fact of classical descriptive set theory).  Then softmax attention is well-behaved at every $(d,n_K)$, \emph{and} there is a hard-attention witness, a continuous score map and two measurable experts, whose argmax-routed head has a bad event that is not measurable.\formal{TLT.Bridge.soft_vs_hard_attention_measurabilitySeparation}
\end{prop}

\begin{proof}[Proof idea]
The first conjunct is \cref{prop:attention-wellbehaved}.  For the second, replace the softmax by a hard $\arg\max$ over the scores.  The set of queries routed to a given key is then a level set of an $\arg\max$, and one can arrange, using a continuous reduction from an analytic non-Borel set, that the resulting bad event is the continuous preimage of that analytic set, hence analytic but not Borel, hence not measurable.  The two routes to attention thus sit on opposite sides of the Borel/analytic line: softmax stays measurable; $\arg\max$ can fall out of the measurable world entirely.
\end{proof}

\begin{remark}[Where the deferred conditions bite]
\label{rmk:attention-measurability}
\Cref{prop:soft-hard-borel} is the concrete content behind the parenthetical of \cref{thm:vc-char}.  The ``mild measurability conditions'' are not automatically met by every natural hypothesis class: hard attention violates them, and the violation is genuine: an analytic-not-Borel bad event over which no uniform-convergence statement can even be phrased.  The softmax is precisely the smoothing that restores measurability, and so the choice between soft and hard attention is, at the level of \emph{objects}, the choice of whether the learner's hypotheses form a class to which the theory of this book applies at all.
\end{remark}

\section{Exercises and Open Problems}
\label{sec:ch1-open}

The following thirteen problems are research-grade, each anchored to a concept introduced above and tagged by its epistemic character: a \emph{known unknown} (\textsc{ku}) is an articulated question awaiting an answer; an \emph{unknown known} (\textsc{uk}) is a regularity in the development or literature not yet absorbed into a clean statement.

\begin{openprob}[Sample compression conjecture, \textsc{ku}]
\label{op:compression-objects}
\Cref{thm:compression-char} characterizes finite VC dimension by compression \emph{with} side information. Prove or refute the Littlestone--Warmuth conjecture: every class of VC dimension $d$ admits a sample compression scheme of size $O(d)$ \emph{without} side information. The gap is isolated in the verified development as \leanname{FLT_Proofs.Complexity.Structures.Open_NoInfoCompressionStrengthening}; the best known unconditional bound is exponential in $d$, and $O(d)$ would tighten every compression-based generalization bound in \Cref{ch:compression}.
\end{openprob}

\begin{openprob}[The optimal No-Free-Lunch constant, \textsc{ku}]
\label{op:nfl-constant}
\Cref{lem:nfl-engine} forces disagreement on more than a quarter of a shattered set, and \cref{thm:nfl} yields expected error $\geq 1/4$. Determine the exact minimax expected error as a function of the ratio $m/|T|$ of sample size to shattered-set size, and identify the largest constant replacing $1/4$ in the averaged bound. The verified engine proves the factor-$4$ form; is $1/2 - o(1)$ attainable as $m/|T| \to 0$?
\end{openprob}

\begin{openprob}[Which finite-VC classes have finite star number? \textsc{ku}]
\label{op:star}
\Cref{prop:star-ge-vc} gives $\VC \leq \mathrm{star}$, and \cref{rmk:star-not-equiv} shows the gap can be infinite (singletons: $\VC = 1$, $\mathrm{star} = \infty$). Give a structural characterization (ideally a forbidden-subclass criterion) of the concept classes with finite VC dimension and finite star number. Because the star number controls active-learning label complexity~\cite{HannekeYang2015}, such a characterization would delimit exactly the classes for which active learning yields an exponential saving.
\end{openprob}

\begin{openprob}[An inductive-bias dimension, \textsc{uk}]
\label{op:bias-dimension}
\Cref{cor:bias-mandatory} makes bias mandatory but offers no measure of how much bias a class encodes. Construct a complexity measure $b(\mathcal{C})$, an ``inductive-bias dimension'', that is monotone under inclusion, well-behaved under products, and lower-bounds achievable sample complexity uniformly over distributions, so that the No-Free-Lunch barrier acquires a quantitative form to match its present qualitative one. The typed bias of \cref{def:inductive-bias-typed} (a preference score with a monotonicity constraint) is the natural carrier for such a measure.
\end{openprob}

\begin{openprob}[The proper/improper gap, \textsc{ku}]
\label{op:proper-improper}
It is known that proper learning can be super-polynomially costlier than improper learning~\cite{PittValiant1988}. Exhibit a single combinatorial invariant of $\mathcal{C}$ that controls the proper-versus-improper sample (or time) complexity ratio, and characterize the classes for which the ratio is bounded. The qualitative separation is settled; the separating dimension is not.
\end{openprob}

\begin{openprob}[Model-selection equivalence witnesses, \textsc{ku}]
\label{op:model-selection}
For each edge of the triangle in \cref{fig:model-selection-triangle}, either supply the missing equivalence witness or upgrade the obstruction into a theorem of non-equivalence. Concretely: characterize the parametric model families on which two-part MML coincides with Bayesian MAP beyond the constant-Fisher-information case, and decide whether the residual discrepancy can be absorbed into the side information of a compression scheme.
\end{openprob}

\begin{openprob}[Efficient version-space representation, \textsc{ku}]
\label{op:version-space}
\Cref{prop:version-space} shows the version space shrinks monotonically and never increases Littlestone dimension, yet for most classes it cannot be efficiently represented or sampled. Characterize the classes whose version space admits a polynomial-size representation under arbitrary consistent data, and relate this property to the eluder dimension~\cite{RussoVanRoy2013} and the star number (\cref{def:star-number}).
\end{openprob}

\begin{openprob}[The coding-theorem constant, \textsc{uk}]
\label{op:solomonoff}
The coding theorem $K(x) = -\log M(x) + O(1)$ (\cref{def:alg-prob}) hides an additive constant depending on the universal machine. Formalize the theorem in a proof assistant with the constant made explicit as a function of the reference machine's description, closing the gap between the algorithmic-probability and Kolmogorov-complexity foundations of MDL. No verified companion currently exists for this result.
\end{openprob}

\begin{openprob}[Constructive necessity in the fundamental theorem, \textsc{uk}]
\label{op:asymmetry}
\Cref{thm:vc-char} is proved with a constructive sufficiency direction (empirical risk minimization) and a non-constructive necessity direction (\cref{rmk:vc-asymmetry}). Give a constructive necessity proof: an explicit algorithm that, from oracle access to a class of infinite VC dimension, produces the hard distribution witnessing non-learnability with a stated dependence on the sample budget. This would remove the asymmetry noted at \leanname{FLT_Proofs.Theorem.PAC.vc_characterization} and upgrade the existential lower bounds to effective ones.
\end{openprob}

\begin{openprob}[A functorial account of the thirteen measures, \textsc{ku}]
\label{op:dimension-multiplicity}
\Cref{fig:concept-class-hub} shows 13 complexity measures targeting the concept class. Determine whether they organize into a single partial order, or a functor from the category of concept classes (with label-preserving maps) to $\N \cup \{\infty\}$, and classify the pairs that admit two-sided polynomial comparisons against the pairs that are provably incomparable. A clean answer would explain why no measure subsumes the others while many pairs are nonetheless linked.
\end{openprob}

\begin{openprob}[The VC dimension of attention, \textsc{ku}]
\label{op:attention-vc-objects}
\Cref{prop:attention-wellbehaved} shows the softmax-attention class $\mathcal{A}_{d,n_K}$ is a well-behaved object, so \cref{thm:vc-char} applies once its VC dimension is known to be finite.  Is $\VC(\mathcal{A}_{d,n_K})$ finite for every $d, n_K$, and what is its exact growth in $d$ and $n_K$?  A Pfaffian-function argument suggests finiteness, but no machine-checked bound exists; closing it would place attention definitively on the learnable side of the hub diagram (\cref{fig:concept-class-hub}) and is the verified counterpart to the well-behavedness handle \leanname{softAttention_wellBehaved}.
\end{openprob}

\begin{openprob}[A legitimate object for hard attention, \textsc{uk}]
\label{op:hard-attention-object}
\Cref{prop:soft-hard-borel} shows that argmax (hard) attention can fall outside the measurable concept classes, with an analytic-not-Borel bad event (\leanname{soft_vs_hard_attention_measurabilitySeparation}).  Identify a refined class of objects, an o-minimal or otherwise tame family of hypothesis classes, that contains hard attention and to which a uniform-convergence theory still applies, so that hard attention re-enters the theory as a legitimate learning object and sheds its status as a measurability pathology.
\end{openprob}

\begin{openprob}[Attention among the thirteen measures, \textsc{ku}]
\label{op:attention-hub}
\Cref{fig:concept-class-hub} measures the concept class with thirteen instruments.  Compute, for the softmax-attention class, the values of as many of them as possible (VC dimension, Littlestone dimension, Rademacher complexity, compression size) and determine which instruments agree and which diverge on this single modern object.  The well-behavedness of \cref{prop:attention-wellbehaved} guarantees the measures are defined; whether attention is a place where they coincide or a place where they pull apart is open.
\end{openprob}

\subsection*{Counterfactual exercises}

Each exercise perturbs the wheel of \Cref{fig:concept-class-hub} or the compression spine of \Cref{fig:model-selection-triangle} and asks where the perturbed claim lands. State for each whether the listed conditions still coincide, and name the result that settles it.

\begin{enumerate}[leftmargin=*,itemsep=3pt]
\item \textbf{Infinite label space.} Let labels range over many classes, so the hub measure becomes the DS dimension and the criterion multiclass PAC learnability. Show that finiteness of the Natarajan dimension no longer suffices, and that the characterizing quantity is the DS dimension: a class is multiclass PAC learnable iff its DS dimension is finite (Brukhim et al., 2022).
\item \textbf{Online learnability.} This is the nested inner wheel (\leanname{littlestone_characterization}), joined to the outer one by a one-way bridge \leanname{online_strictly_stronger_pac}: the online demand implies the statistical one and not conversely, with threshold rules standing at the gap.
\item \textbf{Label noise.} Drop realizability and measure success against the best rule in the class. \leanname{AgnosticPACLearnable} is defined but carries no equivalence theorem, so the classical agnostic characterization (Vapnik 1998) enters by citation.
\item \textbf{Partial concepts.} Allow partially-defined concepts. Here the central edge \emph{breaks}: there are partial classes learnable with infinite VC dimension (Alon--Hanneke--Holzman--Moran 2021), so the VC-to-learnability face fails outright.
\item \textbf{Drop measurability.} Remove the well-behavedness side-condition. A hard argmax-routed attention head has a failure event that escapes measurability (\leanname{soft_vs_hard_attention_measurabilitySeparation}), so uniform convergence cannot even be stated.
\item \textbf{Constructive necessity.} Ask the necessity direction to exhibit an explicit hard distribution. The sufficiency direction builds an empirical risk minimizer; the necessity direction is existential, and whether it can be made constructive is open (\cref{op:asymmetry}).
\item \textbf{Non-computable prior.} Replace the prior by one that is not lower-semicomputable. This breaks the Bayesian approximation edge (\leanname{bayesianToBatch}) while leaving the compression spine intact, since the spine never mentions a prior.
\item \textbf{Compression without side information.} Demand a summary of size linear in the dimension with no tag. The reverse direction is verified (\leanname{compression_bounds_vcdim}); the forward strengthening is the open Littlestone--Warmuth conjecture (\cref{op:compression-objects}).
\item \textbf{PAC-Bayes bridge.} \leanname{pac_bayes_finite} bounds the Gibbs risk by empirical risk plus a prior-weighted complexity term, so the quantity the three model-selection principles share is itself a generalization theorem.
\end{enumerate}

% Chapter 2: Data Presentations and Oracles
% Rewritten against the design-lab Lean 4 kernel (FLT_Proofs.*):
% typed data models, full proofs isomorphic to the verified arguments
% (Text<Informant, VCdim<=Ldim, uniform convergence => PAC), and open problems.
%   time_index, counterexample, advice = Profile A (vocabulary, minimal)
%   noisy_input = Profile A-plus (brief + one theorem reference)
%   iid_sample = Profile G (Framework Bridge: i.i.d. is THE statistical assumption)
%   text_presentation + informant_presentation = Profile D (Separation: text < informant)
%   membership_oracle + equivalence_oracle = Profile G (Framework Bridge: queries vs examples)
%   data_stream = Profile G (Framework Bridge: adversarial vs distributional)

\chapter{Data Presentations and Oracles}
\label{ch:data}

Chapter~\ref{ch:objects} established \emph{what} is learned: concept classes, hypotheses, the objects on which complexity measures act. This chapter asks the question that comes before: \emph{how does data reach the learner?}

The answer cuts deep. A learner given i.i.d.\ samples from an unknown distribution operates in a setting where empirical frequencies converge reliably to true probabilities; the VC dimension decides whether learning is possible. Give the same class to a learner facing an adversarially chosen sequence and the VC dimension becomes silent: thresholds on $\R$ have VC dimension one, yet no algorithm bounds its mistakes against an adversary. Supply positive-only data and the class of all finite languages collapses into the unlearnable. Switch to interactive queries and a class that resists passive learning (deterministic finite automata, under cryptographic hardness) becomes learnable in polynomial query time. Same class, same learner, four outcomes: what changed is the data model.

The data model is the axis along which the field divides into four distinct theories, each with its own characterization theorem, its own complexity measure, and its own impossibility result. Same class, same learner, four outcomes: the variable is the data model.

We begin with vocabulary shared across all paradigms (\S\ref{sec:data-vocab}), then develop each data model in turn, organized by the mathematical content each one carries. Each data model is a type in the companion Lean~4 development, and the three results we prove in full (the text/informant separation, the inequality between the VC and Littlestone dimensions, and the passage from uniform convergence to PAC learnability) are mirrored there.

% ================================================================
% SECTION 1: Shared Vocabulary (Profile A concepts, ~12 lines total)
% ================================================================

\section{Shared Vocabulary}
\label{sec:data-vocab}

Before each paradigm goes its own way, it needs a few shared tools: terms for the mechanical skeleton of any learner-environment interaction. The definitions below are brief because their content is thin; they do real work later as types in the formal development and as building blocks for the richer stream definitions that follow.

A \emph{time index} $t \in \N$ is the discrete step indexing rounds of interaction between learner and environment. A \emph{counterexample} to a hypothesis $h$ with respect to target $c$ is a point $x \in X$ where $h(x) \neq c(x)$, witnessing disagreement; as a type it is the subtype $\{x : h(x) \neq c(x)\}$.\formal{FLT_Proofs.Data.Counterexample} \emph{Advice} is any additional information given to a learner beyond its primary data source; its formal role is explored in the advice reductions of \Cref{ch:exact}.

The most general carrier for sequential data is a \emph{data stream}: a function from time steps to labeled examples, $\text{observe} : \N \to X \times Y$.\formal{FLT_Proofs.Data.DataStream} Texts, informants, and noisy inputs are all data streams with extra structure, defined below as extensions of this type.

\begin{defn}[Noisy Input]
\label{def:noisy-input}
In the \emph{classification noise model} at rate $\eta < 1/2$, each label $y_i$ in the training sample is independently flipped with probability $\eta$: the learner receives $(x_i, y_i \oplus z_i)$ where $z_i \sim \mathrm{Bernoulli}(\eta)$. As a type, a noisy stream extends a data stream with a true-label map and a noise rate bounded by $1/2$.\formal{FLT_Proofs.Data.NoisyDataStream}
\end{defn}

Classification noise does not change the landscape of what is learnable.

\begin{thm}[Noise robustness of PAC learning {\cite{AngluinLaird1988}}]
\label{thm:noise-robust}
For every fixed rate $\eta < 1/2$, a concept class is PAC learnable from clean data if and only if it is PAC learnable under classification noise at rate $\eta$. The sample complexity increases by a factor $O\!\left(1/(1-2\eta)^2\right)$, but the family of learnable classes is identical.
\end{thm}

The forward direction is the substance: a clean learner is recovered from a noisy one by estimating each label's majority vote, the variance of which is controlled by $(1-2\eta)^{-2}$. This robustness is special to the i.i.d.\ model; it fails for adversarial noise in the online setting, which requires different techniques (\cref{op:noise}).

\medskip

With these shared pieces in hand, the four data models can be developed in turn. Each receives treatment proportional to its mathematical content.

% ================================================================
% SECTION 2: The i.i.d. Assumption (Profile G: Framework Bridge)
% ================================================================

\section{The Statistical Data Model: I.I.D.\ Samples}
\label{sec:iid}

\begin{defn}[I.I.D.\ Sample]
\label{def:iid-sample}
An \emph{i.i.d.\ sample} of size $m$ from a distribution $D$ over $X \times Y$ is a sequence $S = ((x_1, y_1), \ldots, (x_m, y_m))$ with each $(x_i, y_i)$ drawn independently from $D$. As a type it bundles a probability measure $D$ on $X \times Y$, a size $m$, and a function $\text{Fin } m \to X \times Y$; the learner knows $m$ but learns about $D$ only through $S$.\formal{FLT_Proofs.Data.IIDSample}
\end{defn}

The i.i.d.\ assumption is the most consequential modeling choice in the field. Everything that makes statistical learning theory work, PAC learnability, uniform convergence, the VC characterization, all of it rests on this single hypothesis about how data arrives. The payoff is uniform convergence; the price is three specific fragilities, each one a door into a different paradigm, and that is why the field has several theories.

\subsection{What i.i.d.\ buys: uniform convergence}

The payoff of the i.i.d.\ assumption is \emph{uniform convergence}: empirical frequencies converge to true probabilities simultaneously across the entire hypothesis class. The word ``uniform'' carries weight here. The required sample size must be fixed before the distribution is chosen and before the target is revealed, which means the same $m_0$ works for every distribution and every target at once. This single-$m_0$ uniformity is what lets a learner pick a hypothesis without knowing the distribution it will be tested against.

\begin{defn}[Uniform convergence]
\label{def:uc}
A hypothesis space $\mathcal{H}$ has \emph{uniform convergence} if for every $\varepsilon, \delta > 0$ there is a sample size $m_0$ such that for \emph{every} distribution $D$ and \emph{every} target $c$, an i.i.d.\ sample of size $m \geq m_0$ satisfies, with probability at least $1 - \delta$,
\[
\sup_{h \in \mathcal{H}} \bigl|\emprisk(h) - \risk(h)\bigr| \leq \varepsilon.
\]\formal{FLT_Proofs.Complexity.Generalization.HasUniformConvergence}
\end{defn}

\begin{thm}[Uniform convergence for VC classes, {\cite{VapnikChervonenkis1971,Valiant1984}}]
\label{thm:uc-vc}
If $\VC(\mathcal{H}) = d < \infty$, then $\mathcal{H}$ has uniform convergence, with $m_0 = O\!\left(\frac{d + \log(1/\delta)}{\varepsilon^2}\right)$ sufficing.\formal{FLT_Proofs.Complexity.Symmetrization.vcdim_finite_imp_uc'}
\end{thm}

The proof, deferred to \Cref{ch:pac}, uses symmetrization and the growth-function bound, both of which exploit i.i.d.\ draws essentially. The payoff is provable here in full: uniform convergence makes empirical risk minimization succeed.

\begin{thm}[Uniform convergence yields PAC learnability]
\label{thm:uc-pac}
If $\mathcal{H}$ is nonempty and has uniform convergence, then $\mathcal{H}$ is PAC learnable, and a consistent learner witnesses it.\formal{FLT_Proofs.Complexity.Generalization.uc_imp_pac}
\end{thm}

\begin{proof}
Define the \emph{consistent learner} $L$: on a sample $S$, output any $h \in \mathcal{H}$ that agrees with every labeled point of $S$ if such an $h$ exists, and otherwise output a fixed $h_0 \in \mathcal{H}$ (available since $\mathcal{H} \neq \varnothing$). Given $\varepsilon, \delta > 0$, let $m_0$ be the sample size from uniform convergence at $(\varepsilon, \delta)$, and set the learner's sample budget to $m_0$.

Fix a target $c \in \mathcal{H}$ and a distribution $D$ realizing it (the labels are $c$'s). The target $c$ is itself consistent with every sample, so a consistent hypothesis always exists, and $L(S)$ has empirical error $\emprisk(L(S)) = 0$. On the uniform-convergence event (which has probability at least $1 - \delta$ for $m \geq m_0$), every hypothesis, $L(S)$ among them, satisfies $|\emprisk - \risk| \leq \varepsilon$. Therefore
\[
\risk(L(S)) \leq \emprisk(L(S)) + \varepsilon = 0 + \varepsilon = \varepsilon
\]
on that event, so $\Pr_{S \sim D^{m}}[\risk(L(S)) \leq \varepsilon] \geq 1 - \delta$. This is exactly the PAC guarantee.
\end{proof}

\Cref{thm:uc-vc,thm:uc-pac} compose into the sufficiency half of the fundamental theorem (\cref{thm:vc-char}): finite VC dimension gives uniform convergence, which gives PAC learnability through ERM. This chain underlies the whole statistical theory, and every link uses the i.i.d.\ assumption.

\subsection{What i.i.d.\ costs: three fragilities}
\label{subsec:iid-costs}

Independence and identical distribution buy uniform convergence. The price is three specific breakdowns, each one a door into a different paradigm.

\begin{enumerate}[leftmargin=*,itemsep=4pt]
\item \textbf{No adversarial robustness.} If an adversary chooses the data sequence, independence fails and uniform convergence no longer implies sequential learnability. The obstruction is a theorem, with thresholds as its witness: there are classes with uniform convergence that are not online learnable (\cref{prop:uc-not-online}). The VC dimension cedes jurisdiction to the Littlestone dimension (\S\ref{sec:streams}).

\item \textbf{No distribution-free sequential guarantees.} The i.i.d.\ model is batch: the learner receives all $m$ examples at once. A sequential version must either fix the distribution across time (the online-with-distributional hybrid, which preserves PAC-style results) or drop distributional assumptions entirely (online learning proper).

\item \textbf{No interaction.} The learner is a passive recipient. It cannot choose which points to query, negotiate with a teacher, or request counterexamples. This passivity is the gap that Angluin's query model fills (\S\ref{sec:queries}).
\end{enumerate}

\begin{prop}[Uniform convergence does not imply online learnability]
\label{prop:uc-not-online}
There is a concept class with uniform convergence (hence finite VC dimension) that is not online learnable. Thresholds on $\R$ are such a class.\formal{FLT_Proofs.Complexity.GeneralizationResults.uc_does_not_imply_online}
\end{prop}

\begin{proof}
Thresholds have $\VC = 1$ (\cref{prop:threshold-vc}), so by \cref{thm:uc-vc} they have uniform convergence. But $\Ldim = \infty$ (\cref{prop:threshold-ldim}): the bisection adversary forces an unbounded number of mistakes against any online learner. By the online characterization (\cref{thm:ldim-char}) infinite Littlestone dimension means not online learnable. Hence uniform convergence does not transport to the adversarial setting.
\end{proof}

\begin{remark}[The VC dimension's jurisdiction]
\label{rmk:vc-jurisdiction}
\Cref{prop:uc-not-online} sharpens a point that is easy to state loosely. The VC dimension characterizes PAC learnability \emph{because of}, and only under, the i.i.d.\ assumption. The same class can have $\VC(\mathcal{C}) = 1$ (PAC learnable from i.i.d.\ data) and $\Ldim(\mathcal{C}) = \infty$ (not online learnable). A complexity measure is not a property of the class alone; it is a property of the class \emph{together with} a data model.
\end{remark}

% ================================================================
% SECTION 3: Text vs Informant (Profile D: Separation Result)
% ================================================================

\section{Enumerative Data: Text and Informant}
\label{sec:text-informant}

Remove the distribution. Remove the sample. Instead, let data arrive as an endless enumeration, one element at a time, with an adversary controlling the order and permitted to withhold any element for as long as it chooses. The learner must eventually commit to a correct hypothesis and stop changing it. This is Gold's model of identification in the limit~\cite{Gold1967}, and it strips away the statistical scaffolding entirely.

Within this stripped setting, a difference that looks quantitative turns out to be qualitative. A \emph{text} shows the learner only positive evidence; an \emph{informant} adds negative evidence as well. One might expect that the second model is just a richer version of the first. The separation proved below shows it is a different paradigm: there is a class identifiable from informant that cannot be identified from any text.

\begin{defn}[Text Presentation]
\label{def:text}
A \emph{text} for a language $L \subseteq \N$ is a data stream $t_1, t_2, \ldots$ of elements of $L \cup \{\#\}$ ($\#$ a pause symbol) in which every element of $L$ appears at least once. As a type, a text for $c : X \to \{0,1\}$ extends a data stream with three conditions: every emitted element is positive, every emitted element satisfies $c$, and every positive point eventually appears.\formal{FLT_Proofs.Data.TextPresentation} A class $\mathcal{C}$ is \emph{identifiable in the limit from text} if some learner, on every $C \in \mathcal{C}$ and every text for $C$, eventually stabilizes on a single correct hypothesis.\formal{FLT_Proofs.Criterion.Gold.EXLearnable}
\end{defn}

\begin{defn}[Informant Presentation]
\label{def:informant}
An \emph{informant} for $L$ is a stream of pairs $(x_i, \ell_i)$ with $\ell_i = \ind[x_i \in L]$ in which every natural number eventually appears. As a type it extends a data stream with two conditions: labels are correct, and every point of $X$ eventually appears.\formal{FLT_Proofs.Data.InformantPresentation}
\end{defn}

A text reports what is in the language; an informant reports membership for every point. The informant gives more, and what it gives is not just quantity: the following theorem shows the gap is qualitative.

\begin{thm}[Text $<$ Informant {\cite{Gold1967}}]
\label{thm:text-informant}
There is a concept class identifiable in the limit from informant but not from text. The class
\[
\mathcal{C}_{\mathrm{fin}} = \{L \subseteq \N : L \text{ finite}\} \cup \{\N\}
\]
of all finite languages together with the full language witnesses the separation.
\end{thm}

\begin{proof}
\emph{From informant: identifiable.} Use the learner that conjectures $\N$ until it sees a negative example, and thereafter conjectures the finite set of positive points seen so far. If the target is $\N$, the informant (being correct) never produces a negative label, so the learner conjectures $\N$ at every step and has converged. If the target is a finite $L$, then $\N \setminus L$ is nonempty and the informant is exhaustive, so a negative example appears at some finite time; from then on the learner conjectures ``positives seen so far,'' and since $L$ is finite, all of its elements appear after finitely many steps, after which the conjecture equals $L$ permanently. Either way the learner stabilizes on a correct hypothesis.

\emph{From text: not identifiable.} Suppose, for contradiction, that a learner $M$ identifies $\mathcal{C}_{\mathrm{fin}}$ from every text. We construct a single text on which $M$ changes its conjecture infinitely often, contradicting convergence. Build the text in stages, maintaining a finite prefix $\sigma$.

Begin enumerating $0, 1, 2, \ldots$, a valid text for $\N$. Since $\N \in \mathcal{C}_{\mathrm{fin}}$ and this is a legitimate text for it, $M$ must eventually conjecture a correct hypothesis for $\N$; let stage $1$ end at the first such time, with prefix $\sigma_1 = \langle 0,1,\ldots,k_1\rangle$.

The prefix $\sigma_1$ is also a legitimate beginning of a text for the finite set $F_1 = \{0,\ldots,k_1\}$. Extend the text by repeating elements of $F_1$ (or padding with $\#$): this makes it a valid text for $F_1 \in \mathcal{C}_{\mathrm{fin}}$, on which $M$ must converge to a hypothesis for $F_1$. Since $F_1 \neq \N$, $M$ must abandon its $\N$-conjecture and conjecture $F_1$ at some later time, ending stage $2$ with prefix $\sigma_2$.

Now resume enumerating fresh natural numbers $k_1+1, k_1+2, \ldots$, again building toward a text for $\N$; $M$ must return to an $\N$-conjecture, ending stage $3$. Iterating, each stage forces $M$ to switch between an $\N$-hypothesis and a finite-set hypothesis. The limiting sequence is a genuine text (every element of its underlying language appears), yet $M$ never stabilizes. This contradicts identification, so no learner succeeds on $\mathcal{C}_{\mathrm{fin}}$ from text.
\end{proof}

The witness $\mathcal{C}_{\mathrm{fin}}$ is instructive: an extremely natural class (all finite sets plus the universal set) already separates the two data models. The obstruction is the impossibility of distinguishing ``all of $\N$'' from ``a large finite set'' using positive data alone.

\begin{remark}[The structural role of negative data]
\label{rmk:negative-data}
The separation shows that negative data does a qualitatively different job: it provides \emph{exclusion}. A text reports what is in $L$; an informant also reports what is \emph{not}, and that single extra power, the ability to say a point is absent, is the whole difference between the two models. Exclusion lets the learner confirm finiteness, which no amount of positive data can do, since any finite prefix of a text for $\N$ is also a prefix of a text for a finite superset. The same asymmetry between inclusion and exclusion recurs in query learning (\S\ref{sec:queries}), where equivalence queries deliver exclusion as counterexamples.
\end{remark}

\begin{historical}
Gold's 1967 paper~\cite{Gold1967} introduced both data models and proved the separation for language identification. It predates PAC learning by 17 years and online learning by 21, making it one of the earliest impossibility results in computational learning theory. The framework remains the standard model for inductive inference and underlies the mind-change hierarchy of Freivalds--Smith and Case--Smith (\Cref{ch:ordinals}), with its criterion hierarchy of finite, explanatory, behaviorally correct, vacillatory, and anomalous identification.\formal{FLT_Proofs.Criterion.Gold.BCLearnable}
\end{historical}

% ================================================================
% SECTION 4: Queries (Profile G: Framework Bridge)
% ================================================================

\section{Query Access: Membership and Equivalence Oracles}
\label{sec:queries}

Both models so far, statistical and enumerative, leave the learner as a passive recipient of data that arrives on someone else's schedule. Angluin's query model~\cite{Angluin1988} breaks that constraint: it grants oracle access, letting the learner ask the questions it needs rather than wait for relevant data to arrive.

\begin{defn}[Membership Oracle]
\label{def:mq}
A \emph{membership oracle} $\mathrm{MQ}$ for a target $c$ answers ``what is $c(x)$?'' for any $x$ chosen by the learner. As a type it bundles a truthful response function with the target it answers about.\formal{FLT_Proofs.Data.MembershipOracle}
\end{defn}

\begin{defn}[Equivalence Oracle]
\label{def:eq}
An \emph{equivalence oracle} $\mathrm{EQ}$ for $c$ answers ``does $h = c$?'': it returns success if $h = c$, and otherwise a counterexample $x$ with $h(x) \neq c(x)$. As a type its query map returns an optional counterexample, with proofs that ``none'' certifies equality and ``some $x$'' is a genuine disagreement.\formal{FLT_Proofs.Data.EquivalenceOracle}
\end{defn}

\begin{defn}[Exact Learning]
\label{def:exact-learning}
A class $\mathcal{C}$ is \emph{exactly learnable} from membership and equivalence queries if some algorithm, for every target $c \in \mathcal{C}$, outputs an $h$ with $h = c$ after a number of queries polynomial in the representation size of $c$ and the size of counterexamples received.\formal{FLT_Proofs.Criterion.PAC.ExactLearnable}
\end{defn}

What these oracles do becomes clear only against what passive data cannot.

\subsection{The power of queries: DFAs as witness}

The comparison is sharpest for DFAs. Passive, distributional data cannot efficiently learn them under standard cryptographic assumptions. With both oracles, Angluin's algorithm learns every $n$-state automaton in polynomial time. The class did not change; the data model did.

\begin{thm}[Angluin's $L^\ast$ algorithm {\cite{Angluin1988}}]
\label{thm:lstar}
The class of $n$-state DFAs over an alphabet $\Sigma$ is exactly learnable using $\mathrm{MQ} + \mathrm{EQ}$ in time polynomial in $n$ and $|\Sigma|$, with at most $n - 1$ equivalence queries and $O(n|\Sigma| m)$ membership queries, where $m$ bounds the length of counterexamples.
\end{thm}

The algorithm builds an observation table encoding the Myhill--Nerode classes of the target, filling it via membership queries and repairing it via equivalence queries; it is developed in full in \Cref{ch:exact}. Contrast the passive setting.

\begin{prop}[DFAs are not PAC-learnable under cryptographic assumptions {\cite{KearnsValiant1994}}]
\label{prop:dfa-pac-hard}
Under standard cryptographic assumptions, DFAs are not efficiently PAC-learnable from random examples alone.
\end{prop}

\begin{remark}[Queries as paradigm shift]
\label{rmk:queries-paradigm}
The DFA example witnesses a genuine paradigm separation: the same class is computationally intractable under passive, distributional data and efficiently learnable under interactive, worst-case queries. Membership queries let the learner \emph{explore} the target's structure; equivalence queries deliver \emph{directed} negative information at the learner's current frontier. The separation runs in both directions: some classes are PAC-learnable but not efficiently exactly learnable (equivalence queries can be costly to simulate), and DFAs go the other way.
\end{remark}

\subsection{Neither oracle suffices alone}

The combination $\mathrm{MQ} + \mathrm{EQ}$ is necessary because each oracle has a specific weakness the other supplies. Membership queries explore structure but cannot certify global correctness; equivalence queries certify correctness but cannot probe structure efficiently. The two propositions below make each weakness precise.

\begin{prop}[Membership queries alone cannot verify]
\label{prop:mq-alone}
Let $X$ be infinite and let $\mathcal{C}$ contain two concepts $c_1 \neq c_2$. A learner restricted to membership queries cannot exactly identify the target after any finite number of queries, unless its queries already pin down the target on all of $X$.
\end{prop}

\begin{proof}
After $q$ membership queries the learner has observed the target on a finite set $Q$ of points. If $c_1$ and $c_2$ agree on $Q$ but differ at some $x \notin Q$ (which can be arranged whenever the queried points fail to separate the class), then the transcript is identical whether the target is $c_1$ or $c_2$, so any fixed output is wrong for one of them. Membership queries supply local values but no certificate of global correctness; only an equivalence query can certify $h = c$.
\end{proof}

\begin{prop}[Equivalence queries alone learn finite classes, but slowly]
\label{prop:eq-alone}
A finite class $\mathcal{C} = \{c_1, \ldots, c_N\}$ is exactly learnable from equivalence queries alone using at most $N - 1$ queries; in the worst case $N - 1$ are necessary, which is exponential in the representation size.
\end{prop}
\begin{proof}
Maintain the version space of candidates consistent with all counterexamples received, initialized to $\mathcal{C}$. Query any surviving candidate. A success ends the process; a counterexample $x$ eliminates at least the queried candidate (it disagrees with the target at $x$) and is consistent only with the candidates that survive. Each query removes at least one candidate, so after at most $N - 1$ queries a single candidate remains and is correct. An adversary answering against the largest surviving set forces one elimination per query, giving the matching lower bound. Equivalence queries certify globally but carry no mechanism to probe structure, so the count is governed by $|\mathcal{C}|$ rather than by representation size.
\end{proof}

% ================================================================
% SECTION 5: Adversarial Streams (Profile G: Framework Bridge)
% ================================================================

\section{Adversarial Streams and the Online Data Model}
\label{sec:streams}

\begin{defn}[Data Stream (Online Protocol)]
\label{def:data-stream}
In the \emph{online protocol}, data arrives as an adversarially chosen stream. At each round $t = 1, 2, \ldots$: the adversary selects $x_t \in X$ (possibly depending on the learner's past predictions); the learner predicts $\hat{y}_t$; the true label $y_t = c(x_t)$ is revealed. The learner minimizes the total mistake count $|\{t : \hat{y}_t \neq y_t\}|$, and there is \emph{no distributional assumption}. A class is \emph{online learnable} if some learner has a finite worst-case mistake bound over all targets in the class and all adversarial sequences.\formal{FLT_Proofs.Learner.Core.OnlineLearner}
\end{defn}

Strip away everything the i.i.d.\ model provides: the distribution, the independence, the convergence of frequencies. What remains is a sequential game: the adversary places instances one by one, adapting to the learner's past predictions; the learner must predict each label before seeing it; the running score is mistakes. No distributional guarantee softens the adversary's choices. This austerity is not just philosophical severity; it forces a different complexity measure, one the VC dimension cannot see.

\begin{defn}[Mistake tree and Littlestone dimension]
\label{def:mistake-tree-data}
A \emph{mistake tree} of depth $d$ is a complete binary tree whose internal nodes are labeled by instances $x \in X$; the inductive type is a leaf or a branch carrying an instance and two subtrees.\formal{FLT_Proofs.Complexity.GameInfra.LTree} The tree is \emph{shattered} by $\mathcal{C}$ if, recursively, the root instance $x$ admits a concept in $\mathcal{C}$ labeling it $1$ whose restriction shatters the left subtree and a concept labeling it $0$ whose restriction shatters the right subtree, so that every root-to-leaf path is realized by some concept forcing a mistake at each node.\formal{FLT_Proofs.Complexity.GameInfra.LTree.isShattered} The \emph{Littlestone dimension} $\Ldim(\mathcal{C})$ is the greatest depth of a shattered mistake tree.\formal{FLT_Proofs.Complexity.GameInfra.LittlestoneDim}
\end{defn}

\begin{warning}[Two shattering conventions]
\label{rmk:branchwise}
There are two readings of ``shattered.'' The \emph{branch-wise} reading asks only that each node admit some concept of the right label, with no consistency required along a path; the \emph{path-consistent} reading (above, and used in the characterization theorem) restricts the class at each branch and demands one concept realize an entire path. The branch-wise reading over-counts: the class $\{x \mapsto 1, x \mapsto 0\}$ of two constant concepts has branch-wise dimension $\infty$ yet is online learnable with a single mistake. The path-consistent dimension is the correct invariant; we use it throughout.
\end{warning}

The Littlestone dimension is to online learning what the VC dimension is to PAC learning, and the two are ordered.

\begin{prop}[The VC dimension lower-bounds the Littlestone dimension]
\label{prop:vc-le-ldim}
For every concept class, $\VC(\mathcal{C}) \leq \Ldim(\mathcal{C})$.\formal{FLT_Proofs.Complexity.Littlestone.BranchWiseLittlestoneDim_ge_VCDim}
\end{prop}

\begin{proof}
Let $S = \{x_1, \ldots, x_d\}$ be shattered by $\mathcal{C}$ in the VC sense, with $d = \VC(\mathcal{C})$. Build the complete binary mistake tree $T$ of depth $d$ that queries $x_1$ at the root and, along every branch, queries $x_2, \ldots, x_d$ in order (the same instance at every node of a given level). Each root-to-leaf path corresponds to a labeling $b \in \{0,1\}^d$ of $S$. Because $S$ is shattered, some concept $c \in \mathcal{C}$ realizes exactly the labeling $b$ on $S$, and this single concept forces the prescribed mistake at every node of the path; this is precisely path-consistent shattering. Hence $T$ is a shattered mistake tree of depth $d$, so $\Ldim(\mathcal{C}) \geq d = \VC(\mathcal{C})$.
\end{proof}

The inequality can be arbitrarily loose. Thresholds on $\R$ have $\VC = 1$ but $\Ldim = \infty$ (\cref{prop:threshold-vc,prop:threshold-ldim}), so finite VC dimension does not bound the Littlestone dimension, which is the formal content of fragility~(1) above. What the Littlestone dimension buys in return is a clean characterization and an algorithm that meets it.

\begin{thm}[Littlestone characterization {\cite{Littlestone1988}}]
\label{thm:ldim-char}
A concept class is online learnable if and only if $\Ldim(\mathcal{C}) < \infty$, and the optimal worst-case mistake bound equals $\Ldim(\mathcal{C})$, achieved by the Standard Optimal Algorithm.\formal{FLT_Proofs.Theorem.Online.littlestone_characterization}
\end{thm}

The Standard Optimal Algorithm predicts, at each point, the label whose surviving version space has the larger Littlestone dimension; each mistake strictly decreases that dimension, so at most $\Ldim(\mathcal{C})$ mistakes occur. The characterization and the matching bound are developed in \Cref{ch:online}.\formal{FLT_Proofs.Theorem.Online.optimal_mistake_bound_eq_ldim} Online learnability is strictly stronger than PAC learnability: every online-learnable class is PAC learnable, and thresholds show the converse fails.\formal{FLT_Proofs.Theorem.Separation.online_strictly_stronger_pac}

\begin{remark}[What adversarial data buys: robustness]
\label{rmk:online-robustness}
A mistake-bounded online learner guarantees its bound against \emph{any} sequence: worst-case, non-stationary, distribution-shifting. The bound $\Ldim(\mathcal{C})$ holds regardless of how data is generated. This unconditional guarantee is what makes online learning the right framework when distributional assumptions are unjustified, and it is exactly the strength that the i.i.d.\ model trades away for the smaller invariant $\VC \leq \Ldim$.
\end{remark}

% ================================================================
% SECTION 6: Taxonomy (all four models compared)
% ================================================================

\section{Taxonomy: Four Data Models, Four Theories}
\label{sec:taxonomy}

The four data models have now been developed individually. Placed side by side, they form a definite structure: each is characterized by a distinct combinatorial dimension, the paradigms are ordered by strict containment in some directions and separated by verified impossibility in others, and every claim in Figure~\ref{fig:data-taxonomy} is backed by a machine-checked proof.\formal{FLT_Proofs.Theorem.Separation.online_strictly_stronger_pac}

\begin{figure}[ht]
\centering
\begin{tikzpicture}[
    cell/.style={draw, thick, rounded corners, minimum width=2.9cm, minimum height=1.05cm, font=\scriptsize, align=center},
    contain/.style={-{Stealth[length=2mm]}, semithick, defblue},
    strict/.style={-{Stealth[length=2mm]}, semithick, thmred},
    lit/.style={{Stealth[length=1.5mm]}-{Stealth[length=1.5mm]}, dashed, notegray, semithick},
    litd/.style={-{Stealth[length=1.5mm]}, dashed, notegray, semithick},
    elab/.style={font=\tiny, text=notegray, inner sep=1.3pt, fill=white, align=center},
    slab/.style={font=\tiny, text=thmred!80!black, inner sep=1.3pt, fill=white, align=center}
]
% compound nodes: data model + its characterizing dimension
\node[cell, fill=onlineblue!12]         (str)  at (0, 2.85)   {\textbf{adversarial stream}\\Littlestone dim};
\node[cell, fill=defblue!12]            (iid)  at (0, 0.7)    {\textbf{i.i.d.\ sample}\\VC dim};
\node[cell, fill=sepgreen!12]           (text) at (-3.55,-1.5){\textbf{enumerative text}\\mind-change ordinal};
\node[cell, fill=edgeorange!12, dashed] (mq)   at (3.55,-1.5) {\textbf{MQ\,+\,EQ}\\query complexity};
\node[cell, fill=layergray, dashed, minimum height=0.8cm] (ctx) at (0,-3.35) {\textbf{in-context (attention)}};
% spine: verified containment + strict separation (online vs PAC)
\draw[contain] ([xshift=-7pt]str.south) -- node[elab,left]{online $\subseteq$ PAC} ([xshift=-7pt]iid.north);
\draw[strict]  ([xshift=7pt]iid.north) -- node[slab,right]{$\not\Rightarrow$ thresholds:\\VC$=1$, Ldim$=\infty$} ([xshift=7pt]str.south);
% Gold vs PAC (verified strict)
\draw[strict] (text) -- node[slab,below right,pos=0.6]{$\not\Rightarrow$ finite support} (iid);
% literature / open edges
\draw[litd] (mq) -- node[elab,above right,pos=0.5]{exact $\not\Rightarrow$ PAC\\(DFAs; crypto, lit.)\\reverse open: \cref{op:reverse-separation}} (iid);
\draw[litd] (ctx) -- node[elab,right]{measurability\\side-condition} (iid);
\node[font=\tiny, text=sepgreen!55!black, anchor=north, inner sep=1pt] at (text.south) {\leanname{gold_theorem}: text obstruction};
% legend
\node[font=\tiny, anchor=west, text=notegray] at (-5.1,-4.35) {%
  \textcolor{defblue}{\rule[0.45ex]{0.5cm}{0.9pt}}~containment (verified) \quad
  \textcolor{thmred}{\rule[0.45ex]{0.5cm}{0.9pt}}~$\not\Rightarrow$ strict separation (verified) \quad
  \textcolor{notegray}{\rule[0.45ex]{0.15cm}{0.8pt}\,\rule[0.45ex]{0.15cm}{0.8pt}}~literature / open};
\end{tikzpicture}
\caption[The four data models, their characterizing dimensions, and the proven separations]{Each node pairs a data model with the combinatorial quantity that decides learnability within it: the VC dimension for i.i.d.\ samples, the Littlestone dimension for adversarial streams, the mind-change ordinal for enumerative texts, query complexity for oracle access. The edges prove what the nodes cannot say alone. The red arrows are verified separations: the same class, thresholds on $\R$, has VC dimension one yet Littlestone dimension infinite, so finite VC dimension does not carry over to the adversarial setting. This is the chapter's argument in diagram form: a complexity measure is not a property of the class in isolation; it is a property of the class under a specific data model.}
\label{fig:data-taxonomy}
\end{figure}

\Cref{tab:data-comparison} displays the correspondence between data model, complexity measure, characterization theorem, and impossibility result.

\begin{table}[ht]
\centering
\small
\caption{Each row is a complete paradigm: the data model fixes the complexity measure, which fixes the characterization theorem, which fixes the impossibility result. Change the data model and every column changes with it.}
\label{tab:data-comparison}
\begin{tabularx}{\textwidth}{>{\bfseries\small}l X X X X}
\toprule
& \textbf{Data Model} & \textbf{Complexity Measure} & \textbf{Characterization} & \textbf{Key Impossibility} \\
\midrule
PAC &
I.I.D.\ sample from unknown $D$ &
$\VC(\mathcal{C})$ &
$\mathcal{C}$ is PAC-learnable iff $\VC(\mathcal{C}) < \infty$ \cite{Valiant1984} &
NFL: $\{0,1\}^X$ not learnable \\[6pt]

Gold &
Text (positive) or informant (positive + negative) &
Finite thickness, mind-change ordinals &
Angluin's condition; telltale sets \cite{Gold1967} &
$\mathcal{C}_{\mathrm{fin}} \cup \{\N\}$: text $<$ informant \\[6pt]

Exact &
MQ + EQ (interactive) &
Query complexity (polynomial in representation) &
$L^\ast$ for DFAs; class-specific \cite{Angluin1988} &
MQ alone or EQ alone insufficient \\[6pt]

Online &
Adversarial stream (no distribution) &
$\Ldim(\mathcal{C})$ &
$\mathcal{C}$ is online-learnable iff $\Ldim(\mathcal{C}) < \infty$ \cite{Littlestone1988} &
$\VC < \infty \not\Rightarrow \Ldim < \infty$ \\
\bottomrule
\end{tabularx}
\end{table}

\section{Cross-Paradigm Separations}
\label{sec:data-separations}

The three separations below are the chapter's principal content. Each one names a class that is learnable under one data model and not another, together with the explicit mechanism blocking transfer. Each has a verified proof in the Lean~4 development, developed in full in \Cref{ch:separations}.

\begin{center}\small
\begin{tabular}{llll}
\toprule
\textbf{Separation} & \textbf{Witness} & \textbf{Mechanism} & \textbf{Reference} \\
\midrule
Text $<$ Informant & $\mathcal{C}_{\mathrm{fin}} \cup \{\N\}$ & Diagonalization & \cite{Gold1967} \\[3pt]
PAC $\not\Rightarrow$ Online & Thresholds on $\R$ & $\VC = 1$, $\Ldim = \infty$ & \cite{Littlestone1988} \\[3pt]
Passive $\not\Rightarrow$ Exact & DFAs & Crypto.\ hardness of PAC; & \cite{Angluin1988}, \\
(and vice versa) & & $L^\ast$ in poly queries & \cite{KearnsValiant1994} \\
\bottomrule
\end{tabular}
\end{center}

Each separation has the same logical structure: a witness class learnable under one data model but not another, together with an explicit obstruction. The PAC/online separation is verified directly (some class is PAC learnable but not online learnable\formal{FLT_Proofs.Theorem.Separation.pac_not_implies_online}) and extends to a three-way separation of PAC, online, and Gold identification.\formal{FLT_Proofs.Theorem.Separation.online_pac_gold_separation} The witnesses (thresholds, finite sets, finite automata) are among the simplest classes in the theory. Their simplicity is the signature of a genuine paradigm boundary: even the simplest classes distinguish the paradigms.

\section{What This Chapter Established}
\label{sec:ch2-summary}

The chapter's argument is one sentence: the data model determines the theory. More precisely:

\begin{enumerate}[leftmargin=*,itemsep=3pt]
\item \textbf{Four data models yield four paradigms.} I.i.d.\ samples give PAC learning; enumerative texts give Gold identification; query oracles give exact learning; adversarial streams give online learning. These are different frameworks with different characterization theorems, complexity measures, and impossibility results, and each data model is a distinct type.

\item \textbf{The separations are witnessed by simple classes.} Thresholds separate PAC from online (\cref{prop:vc-le-ldim} and the threshold gap); $\mathcal{C}_{\mathrm{fin}} \cup \{\N\}$ separates text from informant (\cref{thm:text-informant}); DFAs separate passive from interactive. The simplicity of the witnesses shows the boundaries are fundamental.

\item \textbf{The complexity measure is a joint property of class and data model.} The VC dimension measures $\mathcal{C}$ against i.i.d.\ data; the Littlestone dimension measures it against adversarial data; $\VC \leq \Ldim$ orders them, and the gap is unbounded. Neither is ``the'' complexity of $\mathcal{C}$; each is the complexity of the class under a presentation. The same object has different complexities under different data models; this is a central structural insight of the field, and the reason uniform convergence does not transport to the adversarial setting (\cref{prop:uc-not-online}).
\end{enumerate}

Chapters~\ref{ch:pac}--\ref{ch:universal} develop each paradigm in full. The present chapter has identified the axis along which they separate.

% ================================================================
% TRANSFORMER CONNECTION (retrofit): attention as a measurable batch
% learner over i.i.d. data; routing as a measurable partition;
% the float32 forward-pass error of the readout.
% ================================================================

\section{The Data of a Transformer}
\label{sec:ch2-transformer}

The consistent learner of \cref{thm:uc-pac} was defined inline: ``output any $h \in \mathcal{H}$ agreeing with $S$,'' with no $\sigma$-algebra on $\mathcal{H}$ and no requirement that $S \mapsto L(S)$ be measurable. For the simplest classes that omission is harmless. For a parametric attention model it is not free, and making it explicit reveals two axes the classical model suppresses: the \emph{measurability} of the data-to-hypothesis map, and the \emph{finite precision} in which data is actually read by a real machine.

\begin{prop}[Attention routing is a measurable partition]
\label{prop:attention-measurable-routing}
The data-dependent routing of a multi-head attention model (the partition of the query space according to which key each query attends to) is measurable, and the hypothesis class it induces is a well-behaved VC-measurability target.\formal{TLT_Proofs.Routing.MeasurableScoreRouting.multiHeadArgmax_wellBehaved}  The same guarantee holds for a finite-cell router.\formal{TLT_Proofs.Tame.FiniteCellRouter.finiteCellRouter_wellBehaved}  Consequently the attention learner is a \emph{measurable} batch learner in the sense of the condition \leanname{LearnEvalMeasurable}\formal{FLT_Proofs.Theorem.Separation.LearnEvalMeasurable} that \cref{thm:uc-pac} left implicit.
\end{prop}

\begin{proof}
Fix the parameters and a query~$x$.  The routing decision is the $\arg\max$ over keys of the scores $\langle x, k_i\rangle$, so the set of queries routed to key~$i$ is the intersection of the half-spaces $\{x : \langle x, k_i - k_j\rangle \geq 0\}$ over $j \neq i$, which is a Borel set.  The routing map $(\text{parameters}, x) \mapsto (\text{winning key})$ is therefore measurable, and the head's prediction is a measurable function of parameters and query.  A parametrized family of concepts with this joint measurability is by definition a well-behaved VC-measurability target, and the learner's selection $S \mapsto L(S)$ inherits measurability from it.  This is the condition \leanname{LearnEvalMeasurable}: for every sample size~$m$, the evaluation $(S, x) \mapsto L(S)(x)$ is measurable, which is what the uniform-convergence argument of \cref{thm:uc-pac} requires.
\end{proof}

The second axis is invisible in the real-valued data model of \cref{def:iid-sample} and unavoidable in practice: the learner does not read real numbers.

\begin{prop}[The forward pass reads data in finite precision]
\label{prop:attention-fp32}
On real hardware the attention readout is a \texttt{float32} fold-sum of softmax-weighted values, the idealized real-valued sum being unavailable to a finite-precision machine.  The executed result differs from the ideal by a bounded amount: for a fold of $m$ summands whose absolute values total at most~$S$,
\[
  \bigl|\,\mathrm{fp32fold}(v) - \textstyle\sum_i v_i\,\bigr| \;\leq\; 2^{-24}\,(m+1)\,S,
\]\formal{TLT_Proofs.Bridge.Fp32.ClampedReductionRounding.fp32Foldl_error_le_of_sum_bound}
and the executed IEEE-754 fold obeys an analogous accumulated-error envelope.\formal{TLT_Proofs.Bridge.Fp32.SequentialSummationBackwardError.ie32_foldl_envelope}
\end{prop}

\begin{proof}[Proof idea]
Each addition in the fold commits a relative rounding error of at most the unit roundoff $2^{-24}$; propagating these through the $m$ summations and bounding by the total magnitude~$S$ gives the stated linear-in-$m$ envelope by backward-error analysis.  What matters is that such a constant exists at all: the map from data to prediction that actually runs is the idealized one perturbed by a bounded amount.
\end{proof}

\begin{remark}[A third coordinate of the data model]
\label{rmk:attention-data-model}
\Cref{rmk:vc-jurisdiction} held that a complexity measure is a property of the class \emph{together with} a data model.  Attention adds a coordinate that the four classical models share but never expose: the \emph{precision} in which the data is represented.  The i.i.d.\ jurisdiction is where the measurable routing of \cref{prop:attention-measurable-routing} lets the VC machinery apply at all; the float32 envelope of \cref{prop:attention-fp32} is the cost of executing that machinery on data a real machine can hold.  The classical theory idealizes both axes away: it assumes the learner measurable and the arithmetic exact, and a transformer is where neither assumption is free.
\end{remark}

\section{Exercises and Open Problems}
\label{sec:ch2-open}

These problems sit at the edges where the chapter's argument is not yet complete. Each is anchored to a specific construct above. The tags signal what kind of incompleteness you face: a \emph{known unknown} (\textsc{ku}) is a question whose answer is open; an \emph{unknown known} (\textsc{uk}) is a result the literature already contains, but one not yet formalized or absorbed into a clean general statement in this framework. In both cases the question is real, the difficulty is genuine, and the answer, if you find one, changes something.

\begin{openprob}[Formalizing noise robustness, \textsc{uk}]
\label{op:noise}
Formalize \cref{thm:noise-robust} (Angluin--Laird), the equivalence ``PAC under classification noise at rate $\eta < 1/2$ $\Leftrightarrow$ clean PAC'' with the explicit $O(1/(1-2\eta)^2)$ overhead, and identify the exact threshold behavior as $\eta \to 1/2$.
\end{openprob}

\begin{openprob}[The optimal uniform-convergence constant, \textsc{ku}]
\label{op:uc-constant}
\Cref{thm:uc-vc} gives $m_0 = O((d + \log(1/\delta))/\varepsilon^2)$. Determine the exact leading constant for VC classes: symmetrization, chaining, and the Dudley entropy integral give different bounds, and which is tight for general classes of dimension $d$ is open.
\end{openprob}

\begin{openprob}[A characterization of identifiability from text, \textsc{uk}]
\label{op:text-char}
\Cref{thm:text-informant} shows $\mathcal{C}_{\mathrm{fin}} \cup \{\N\}$ is not text-identifiable. Angluin's telltale condition characterizes which classes are: a class is text-identifiable iff every language has a finite ``telltale'' subset distinguishing it from its sublanguages in the class. Formalize this characterization on top of the typed criterion hierarchy (\leanname{FLT_Proofs.Criterion.Gold.EXLearnable}).
\end{openprob}

\begin{openprob}[Closing the VC--Littlestone gap, \textsc{ku}]
\label{op:vc-ldim-gap}
\Cref{prop:vc-le-ldim} gives $\VC \leq \Ldim$, and thresholds make the gap infinite. Characterize the classes for which $\Ldim$ is finite given finite $\VC$ (the Littlestone classes), and give a quantitative bound on $\Ldim$ in terms of $\VC$ and a second structural parameter for the cases where one exists. When is $\Ldim = \VC$ exactly?
\end{openprob}

\begin{openprob}[Query complexity of exact learning, \textsc{ku}]
\label{op:query-complexity}
For the DFA learner of \cref{thm:lstar}, the equivalence-query count $n-1$ is optimal but the membership-query count $O(n|\Sigma| m)$ is not known to be. Determine the optimal joint $\mathrm{MQ} + \mathrm{EQ}$ complexity of exactly learning $n$-state DFAs, and characterize the classes for which $\mathrm{EQ}$ alone (no membership queries) suffices in a polynomial number of queries.
\end{openprob}

\begin{openprob}[The minimally adequate teacher, \textsc{uk}]
\label{op:mat}
\Cref{prop:mq-alone,prop:eq-alone} show each oracle alone is insufficient. Characterize exactly which concept classes become efficiently exactly learnable when the two oracles are combined (the ``minimally adequate teacher'' question) and relate this to the VC and Littlestone dimensions, since neither passive invariant governs query learning.
\end{openprob}

\begin{openprob}[An invariant for query learning, \textsc{ku}]
\label{op:query-dimension}
PAC learning has the VC dimension and online learning has the Littlestone dimension, each characterizing its paradigm. Exact learning has no single characterizing invariant. Construct a combinatorial dimension that characterizes polynomial-query exact learnability the way $\VC$ and $\Ldim$ characterize their paradigms, or prove that no such single invariant exists.
\end{openprob}

\begin{openprob}[Distribution shift and online-to-batch, \textsc{ku}]
\label{op:online-batch}
\Cref{rmk:online-robustness} notes that online bounds hold under distribution shift while i.i.d.\ guarantees do not. Quantify the degradation: given a class of Littlestone dimension $\ell$ and a sequence whose distribution drifts with total variation budget $B$, give matched upper and lower bounds on the achievable risk, interpolating between the i.i.d.\ ($B = 0$) and adversarial regimes.
\end{openprob}

\begin{openprob}[A unified data-model complexity, \textsc{uk}]
\label{op:unified-measure}
The four paradigms use four invariants (VC, mind-change ordinal, query complexity, Littlestone). Determine whether they are instances of a single complexity functor on (class, data-model) pairs, recovering each invariant by specializing the data model, and classify which pairs of data models admit a transfer theorem versus a provable separation (\cref{tab:data-comparison}).
\end{openprob}

\begin{openprob}[Verifying the three-way separation constructively, \textsc{uk}]
\label{op:three-way}
The PAC/online/Gold separation is verified (\leanname{FLT_Proofs.Theorem.Separation.online_pac_gold_separation}), but the witnesses are produced existentially. Give a constructive, effective version: a single family of classes, indexed by a parameter, that realizes each pairwise separation with explicit witnesses and explicit obstructions, so that the entire separation lattice of \cref{tab:data-comparison} is exhibited by one parameterized construction rather than three independent ones.
\end{openprob}

\begin{openprob}[The measurable attention learner, \textsc{uk}]
\label{op:attention-measurable}
\Cref{prop:attention-measurable-routing} shows the attention learner is measurable, the condition \leanname{LearnEvalMeasurable} that \cref{thm:uc-pac} left implicit, with the routing well-behavedness as witness (\leanname{multiHeadArgmax_wellBehaved}).  Assemble these into a single statement: a parametric softmax-attention learner over i.i.d.\ data is PAC, with the measurability discharged rather than assumed.
\end{openprob}

\begin{openprob}[Does PAC survive finite precision? \textsc{ku}]
\label{op:attention-fp32-pac}
\Cref{prop:attention-fp32} bounds the gap between the executed float32 attention learner and its idealized real-valued counterpart (\leanname{fp32Foldl_error_le_of_sum_bound}).  Does the PAC guarantee transfer to the executed learner, and at what cost: is there a sample-complexity penalty that scales with the precision (the unit roundoff $2^{-24}$), so that learnability degrades gracefully as arithmetic coarsens?
\end{openprob}

\begin{openprob}[Attention under the other data models, \textsc{ku}]
\label{op:attention-data-models}
The mapping of \cref{sec:ch2-transformer} is i.i.d.\ (batch).  Where does attention sit under the other three data models of this chapter: is there an online (streaming) attention learner with a finite mistake bound, or a query-driven (active) one, and does the measurable-routing well-behavedness (\leanname{finiteCellRouter_wellBehaved}) transfer to those presentations, or does the data model change attention's learnability as it changes thresholds'?
\end{openprob}

\subsection*{Counterfactual exercises}

The chapter's argument is that the data model determines everything downstream: the complexity measure, the characterization, the impossibility. The exercises below press on that claim directly. Each one takes a single feature of a data model, corrupts it, weakens it, or swaps it for an adversarial counterpart, and asks what breaks. The answer in each case either confirms the paradigm boundary or locates the exact condition that holds the theory together.

\begin{enumerate}[leftmargin=*,itemsep=3pt]
\item \textbf{Noisy oracle.} Corrupt labels at a rate below one half. Show that the VC dimension is invariant under the noise, so PAC learnability is preserved (\cref{thm:noise-robust}).
\item \textbf{Semi-supervised text.} Add a few labels to a positive-only text. Pin down the condition on the added labels under which the class becomes identifiable, that is, when they supply the exclusion information the finite-support witness needs.
\item \textbf{Stream with bounded memory.} Cap the learner's state. Explain why this can only raise the mistake bound above the Littlestone dimension, and why no characterizing dimension is known for the memory-bounded model.
\item \textbf{Distribution shift.} Let the data be non-stationary. Show that the Littlestone mistake bound covers a drifting distribution verbatim, since it holds against any sequence (\leanname{optimal_mistake_bound_eq_ldim}).
\item \textbf{Partial labels.} Let the oracle abstain. Identify the abstention rate at which the surviving labels recover enough exclusion to defeat the Gold diagonalization (\leanname{gold_theorem}), tipping the model from text-like to informant-like.
\item \textbf{Adversarial versus stochastic order.} Switch the sample order from random to adversarial. Stochastic order makes VC govern, adversarial order makes Littlestone govern; name the witness (thresholds, VC one, Littlestone infinite) that forces the gap (\leanname{online_strictly_stronger_pac}).
\item \textbf{Privacy constraint.} Require a differentially private learner. Argue that under pure privacy the governing quantity is a representation dimension distinct from VC, so the i.i.d.\ characterization no longer applies.
\item \textbf{In-context presentation.} Let the data arrive as a prompt consumed by an attention learner. A hard attention head's failure event can be analytic while failing to be Borel (\leanname{attention_nonBorel_badEvent}); exhibit the measurability side-condition under which the VC characterization still applies.
\end{enumerate}

% Chapter 3: Automata, Languages, and Computability
% Rewritten against the design-lab Lean 4 kernel (FLT_Proofs.Computation, FLT_Proofs.Complexity.MindChange):
% every model typed against the verified computability substrate; full proofs for the
% pumping separation, the subset construction, and the limit-lemma direction.
%   dfa, nfa, pushdown_automaton, turing_machine, cfg = Profile A (vocabulary)
%   regular_language, cfl, re_set, partial_computable = Profile A (one-line)
%   godel_numbering = Profile A-plus (brief but load-bearing for Gold)
%   restricts/extends_grammar edge pair = Profile C (Revelation: conceptual centerpiece)
%   limiting_recursion + delta_02_class = Profile G (Framework Bridge to Gold)

\chapter{Automata, Languages, and Computability}
\label{ch:automata}

This chapter assembles the automata-theoretic and computability-theoretic vocabulary required by Gold's model of identification in the limit (\Cref{ch:gold}). Much of the material is standard; the reader who needs the full classical development should consult Hopcroft and Ullman~\cite{HopcroftUllman1979} or Sipser~\cite{Sipser2013}. Every model defined here is typed in the companion Lean~4 development's computability substrate, and three results---the pumping separation of regular from context-free languages, the subset construction, and one direction of the limit lemma---are proved in full.

Two things here are not standard vocabulary. First, the distinction between \emph{restricts} and \emph{extends\_grammar} edges---two ways one computational model can generalize another, one conservative and one genuinely new---is illustrated through the NFA/DFA and PDA/DFA pairs (\S\ref{sec:restricts-extends}); the distinction pervades the entire knowledge graph. Second, limiting recursion connects computability to learning by placing Gold's identification in the limit at a precise level of the arithmetical hierarchy (\S\ref{sec:bridge}).

% ================================================================
% SECTION 1: The Chomsky Hierarchy
% ================================================================

\section{The Chomsky Hierarchy}
\label{sec:chomsky}

The Chomsky hierarchy classifies formal languages by the computational resources needed to recognize them. Figure~\ref{fig:chomsky} displays the containment structure; the definitions follow. A \emph{formal language} over an alphabet $\Sigma$ is a set of finite words, $L \subseteq \Sigma^\ast$.\formal{FLT_Proofs.Computation.FormalLanguage}

\begin{figure}[ht]
\centering
\begin{tikzpicture}[
    every node/.style={font=\small},
    level/.style={draw, thick, ellipse, minimum height=1.2cm, align=center}
]
% Nested ovals, outermost first
\node[level, minimum width=10cm, minimum height=5.8cm, fill=layergray!60] (re) at (0,0) {};
\node[font=\scriptsize, anchor=north] at (re.north) {Recursively enumerable (Type 0)};

\node[level, minimum width=7.8cm, minimum height=4.2cm, fill=defblue!6] (rec) at (0,-0.3) {};
\node[font=\scriptsize, anchor=north] at (rec.north) {Recursive (decidable)};

\node[level, minimum width=5.6cm, minimum height=2.8cm, fill=edgeorange!8] (cfl) at (0,-0.5) {};
\node[font=\scriptsize, anchor=north] at (cfl.north) {Context-free (Type 2)};

\node[level, minimum width=3.2cm, minimum height=1.4cm, fill=sepgreen!10] (reg) at (0,-0.6) {};
\node[font=\scriptsize] at (reg.center) {Regular (Type 3)};

% Recognizer labels on the right
\node[font=\scriptsize\itshape, notegray, anchor=west] at (4.2,2.2) {Turing machine};
\node[font=\scriptsize\itshape, notegray, anchor=west] at (3.2,0.8) {TM (halting)};
\node[font=\scriptsize\itshape, notegray, anchor=west] at (2.2,-0.1) {Pushdown automaton};
\node[font=\scriptsize\itshape, notegray, anchor=west] at (1.2,-0.6) {DFA / NFA};
\end{tikzpicture}
\caption{The Chomsky hierarchy. Each level is strictly contained in the next. The recognizer type is indicated at right. All containments are proper.}
\label{fig:chomsky}
\end{figure}

\subsection{Recognizers}

\begin{defn}[Deterministic Finite Automaton]
\label{def:dfa}
A \emph{DFA} is a 5-tuple $M = (Q, \Sigma, \delta, q_0, F)$ where $Q$ is a finite state set, $\Sigma$ a finite input alphabet, $\delta : Q \times \Sigma \to Q$ a total transition function, $q_0 \in Q$ the start state, and $F \subseteq Q$ the accepting states. $M$ accepts $w \in \Sigma^\ast$ if the unique run on $w$ ends in $F$; the accepted language is $\{w : \hat\delta(q_0, w) \in F\}$, where $\hat\delta$ folds $\delta$ along $w$.\formal{FLT_Proofs.Computation.DFA'}
\end{defn}

\begin{defn}[Nondeterministic Finite Automaton]
\label{def:nfa}
An \emph{NFA} is a 5-tuple identical to a DFA except that $\delta : Q \times (\Sigma \cup \{\varepsilon\}) \to \mathcal{P}(Q)$ maps each state--symbol pair to a \emph{set} of successors and $\varepsilon$-transitions are permitted; $M$ accepts $w$ if \emph{some} run ends in $F$.\formal{FLT_Proofs.Computation.NFA'}
\end{defn}

\begin{defn}[Pushdown Automaton]
\label{def:pda}
A \emph{pushdown automaton} (PDA) is a 7-tuple $M = (Q, \Sigma, \Gamma, \delta, q_0, Z_0, F)$ where $\Gamma$ is a finite stack alphabet, $Z_0 \in \Gamma$ the initial stack symbol, and $\delta : Q \times (\Sigma \cup \{\varepsilon\}) \times \Gamma \to \mathcal{P}(Q \times \Gamma^\ast)$ maps each state--input--stack-top triple to a set of (state, stack-replacement) pairs. The PDA has an unbounded LIFO stack and accepts by final state or empty stack.\formal{FLT_Proofs.Computation.PDA}
\end{defn}

\begin{defn}[Turing Machine]
\label{def:tm}
A \emph{Turing machine} is a 7-tuple $M = (Q, \Sigma, \Gamma, \delta, q_0, q_{\mathrm{acc}}, q_{\mathrm{rej}})$ where $\Gamma \supseteq \Sigma$ is the tape alphabet (with a blank), and $\delta : Q \times \Gamma \to Q \times \Gamma \times \{L, R\}$ is a partial transition function specifying the next state, the symbol to write, and head movement; $M$ accepts $w$ if it halts in $q_{\mathrm{acc}}$.\formal{FLT_Proofs.Computation.TuringMachine}
\end{defn}

\begin{defn}[Context-Free Grammar]
\label{def:cfg}
A \emph{context-free grammar} (CFG) is a 4-tuple $G = (V, \Sigma, R, S)$ with $V$ a finite variable set, $\Sigma$ the terminal alphabet, $R \subseteq V \times (V \cup \Sigma)^\ast$ a finite rule set, and $S \in V$ the start variable; $G$ generates $L(G) = \{w \in \Sigma^\ast : S \Rightarrow^\ast w\}$, the reflexive--transitive closure of one-step derivation.\formal{FLT_Proofs.Computation.CFG}
\end{defn}

\subsection{Language classes}

A language $L \subseteq \Sigma^\ast$ is \emph{regular} if accepted by some DFA;\formal{FLT_Proofs.Computation.IsRegularLang} \emph{context-free} if generated by some CFG (equivalently, accepted by some PDA);\formal{FLT_Proofs.Computation.IsContextFree} \emph{recursively enumerable} (r.e.) if accepted by some Turing machine;\formal{FLT_Proofs.Computation.RESet} and \emph{recursive} (decidable) if accepted by a Turing machine that halts on every input. A function $f : \N \rightharpoonup \N$ is \emph{partial computable} if some Turing machine outputs $f(n)$ when $f(n)$ is defined and diverges otherwise.\formal{FLT_Proofs.Computation.PartialComputable}

\subsection{Separating the levels by pumping}

The containments of \cref{fig:chomsky} are proper, and the regular/context-free gap drives the edge distinction of \S\ref{sec:restricts-extends}. We prove it from a single combinatorial lemma, whose only input is that a DFA has finitely many states.

\begin{lemma}[Pumping lemma for regular languages {\cite{HopcroftUllman1979}}]
\label{lemma:pumping}
For every regular $L$ there is a length $p$ such that every $w \in L$ with $|w| \geq p$ factors as $w = xyz$ with $|xy| \leq p$, $|y| \geq 1$, and $xy^iz \in L$ for all $i \geq 0$.
\end{lemma}

\begin{proof}
Let $M$ be a DFA accepting $L$, with $p = |Q|$ states. Take $w \in L$ with $|w| \geq p$. The run of $M$ on $w$ passes through the states $r_0, r_1, \ldots, r_{|w|}$ with $r_0 = q_0$; among the first $p+1$ of these, two coincide by the pigeonhole principle, say $r_j = r_k$ with $j < k \leq p$. Set $x = w_{1\ldots j}$, $y = w_{j+1\ldots k}$, $z = w_{k+1\ldots |w|}$. Then $|xy| = k \leq p$ and $|y| = k - j \geq 1$. Reading $y$ takes state $r_j$ back to itself, so reading $y$ any number of times leaves the run in $r_j$; hence $xy^iz$ drives $M$ from $q_0$ to the same accepting state $r_{|w|}$ for every $i \geq 0$, and $xy^iz \in L$.
\end{proof}

\begin{prop}[The language $\{a^n b^n\}$ is context-free but not regular]
\label{prop:anbn}
$L = \{a^n b^n : n \geq 0\}$ is context-free and not regular; hence $\mathrm{REG} \subsetneq \mathrm{CFL}$.
\end{prop}

\begin{proof}
$L$ is generated by the CFG $S \to aSb \mid \varepsilon$ (equivalently, accepted by a PDA that pushes on each $a$ and pops on each $b$), so $L$ is context-free. Suppose $L$ were regular, with pumping length $p$ from \cref{lemma:pumping}. Apply the lemma to $w = a^p b^p \in L$, with $|w| = 2p \geq p$. The factorization $w = xyz$ has $|xy| \leq p$, so $xy$ lies within the initial block of $a$'s and $y = a^k$ for some $k \geq 1$. Then $xy^2z = a^{p+k} b^p$, which has more $a$'s than $b$'s and so is not in $L$, contradicting the lemma. Therefore $L$ is not regular, and since $L \in \mathrm{CFL}$, the containment $\mathrm{REG} \subseteq \mathrm{CFL}$ is proper.
\end{proof}

\subsection{Closure properties}

The classes differ sharply in closure. Table~\ref{tab:closure} records the facts used in later chapters; the proofs are standard~\cite{HopcroftUllman1979}, and we give the two that recur. Regular languages are closed under union by the product construction; context-free languages are \emph{not} closed under intersection.

\begin{prop}[Two closure facts]
\label{prop:closure}
(i) If $L_1, L_2$ are regular then so is $L_1 \cup L_2$. (ii) Context-free languages are not closed under intersection.
\end{prop}

\begin{proof}
(i) Let $M_1 = (Q_1,\Sigma,\delta_1,s_1,F_1)$ and $M_2 = (Q_2,\Sigma,\delta_2,s_2,F_2)$ accept $L_1, L_2$. The product DFA on state set $Q_1 \times Q_2$ with transition $\delta((q_1,q_2),a) = (\delta_1(q_1,a),\delta_2(q_2,a))$, start $(s_1,s_2)$, and accepting set $\{(q_1,q_2): q_1 \in F_1 \text{ or } q_2 \in F_2\}$ simulates both machines in parallel and accepts exactly $L_1 \cup L_2$.

(ii) The languages $L_1 = \{a^n b^n c^m : n,m \geq 0\}$ and $L_2 = \{a^m b^n c^n : n,m \geq 0\}$ are each context-free (a single matched block, the other free), but $L_1 \cap L_2 = \{a^n b^n c^n : n \geq 0\}$, which is not context-free---the pumping lemma for context-free languages forces any pumped factor to disturb the three-way balance. Hence the context-free languages are not closed under intersection.
\end{proof}

\begin{table}[ht]
\centering
\small
\caption{Closure properties of the Chomsky hierarchy classes. $\checkmark$ = closed, $\times$ = not closed. Union for regular and intersection for context-free are proved in \cref{prop:closure}; the remaining entries are standard~\cite{HopcroftUllman1979}.}
\label{tab:closure}
\begin{tabular}{lccccc}
\toprule
\textbf{Class} & \textbf{Union} & \textbf{Intersection} & \textbf{Complement} & \textbf{Concat.} & \textbf{Kleene $*$} \\
\midrule
Regular       & $\checkmark$ & $\checkmark$ & $\checkmark$ & $\checkmark$ & $\checkmark$ \\
Context-free  & $\checkmark$ & $\times$     & $\times$     & $\checkmark$ & $\checkmark$ \\
Recursive     & $\checkmark$ & $\checkmark$ & $\checkmark$ & $\checkmark$ & $\checkmark$ \\
R.E.          & $\checkmark$ & $\checkmark$ & $\times$     & $\checkmark$ & $\checkmark$ \\
\bottomrule
\end{tabular}
\end{table}

% ================================================================
% SECTION 2: restricts vs extends_grammar (Profile C — Revelation)
% ================================================================

\section{Two Kinds of Generalization: Restricts vs.\ Extends}
\label{sec:restricts-extends}

The Chomsky hierarchy is a containment chain, but its edges are not all of one type. Moving between levels sometimes removes a constraint (a more flexible model with no new primitive) and sometimes adds a genuinely new resource (a model with a mechanism it lacked). This distinction---conservative versus generative generalization---is a recurring structural motif of the knowledge graph, and the automata hierarchy provides its first concrete instance.

\begin{defn}[Restricts]
\label{def:restricts}
Model $B$ \rel{restricts} model $A$ if $B$ is obtained from $A$ by \emph{removing a constraint}: every $A$ is already a $B$ (after trivial embedding), and the two recognize the same class of languages. The generalization is conservative---it adds representational flexibility without expanding expressive power.
\end{defn}

\begin{defn}[Extends Grammar]
\label{def:extends-grammar}
Model $B$ \rel{extends\_grammar} model $A$ if $B$ is obtained by \emph{adding a new computational primitive}: $B$ recognizes a strictly larger class. The generalization is generative---it requires a fundamentally new resource.
\end{defn}

The DFA is the anchor for both edges. Figure~\ref{fig:restricts-extends} displays the structure, and the two propositions below supply the rigorous content of each edge.

\begin{figure}[ht]
\centering
\begin{tikzpicture}[
    box/.style={draw, thick, rounded corners, minimum width=2.4cm, minimum height=0.9cm, font=\small\bfseries, align=center},
    arr/.style={-{Stealth[length=2.5mm]}, thick},
    lbl/.style={font=\scriptsize, fill=white, inner sep=2pt, align=center}
]
\node[box, fill=sepgreen!12] (dfa) at (0,0) {DFA};
\node[box, fill=sepgreen!12] (nfa) at (5,1.5) {NFA};
\node[box, fill=edgeorange!12] (pda) at (5,-1.5) {PDA};

\draw[arr, sepgreen!70!black] (nfa) -- node[lbl, above left] {\rel{restricts}\\{\scriptsize\itshape removes determinism}\\{\scriptsize\itshape same power}} (dfa);
\draw[arr, edgeorange!80!black] (pda) -- node[lbl, below left] {\rel{extends\_grammar}\\{\scriptsize\itshape adds stack}\\{\scriptsize\itshape strictly more power}} (dfa);

% Annotations
\node[font=\scriptsize, text=sepgreen!70!black, anchor=west] at (6, 1.5) {Regular languages};
\node[font=\scriptsize, text=edgeorange!80!black, anchor=west] at (6, -1.5) {Context-free languages};
\end{tikzpicture}
\caption{Two generalizations of the DFA. The NFA removes the determinism constraint (conservative: same expressive power, \cref{thm:nfa-dfa}). The PDA adds a stack (generative: strictly more power, \cref{prop:anbn}). Both arrows point toward the DFA to indicate that the DFA is the more constrained model.}
\label{fig:restricts-extends}
\end{figure}

\paragraph{The NFA--DFA edge: conservative generalization.}
An NFA relaxes the DFA by allowing several successors per transition and $\varepsilon$-moves, removing the determinism constraint. The subset construction shows this buys no expressive power.

\begin{thm}[Subset construction {\cite{RabinScott1959}}]
\label{thm:nfa-dfa}
Every NFA with $n$ states has an equivalent DFA with at most $2^n$ states; hence $\mathrm{REG}_{\mathrm{NFA}} = \mathrm{REG}_{\mathrm{DFA}}$.
\end{thm}

\begin{proof}
Let $N = (Q, \Sigma, \Delta, S_0, F)$ be an NFA (with $S_0$ the set of start states after $\varepsilon$-closure). Define a DFA $D$ on the power set $2^Q$ by $\delta(T, a) = \bigcup_{q \in T} \Delta(q, a)$ (each set extended by $\varepsilon$-closure), start state $S_0$, and accepting sets $\{T : T \cap F \neq \varnothing\}$. By induction on $|w|$, the state $\hat\delta(S_0, w)$ of $D$ after reading $w$ is exactly the set of states $N$ can occupy after $w$: the base case is $S_0$, and the inductive step is the definition of $\delta$. Thus $D$ accepts $w$ iff some run of $N$ ends in $F$ iff $N$ accepts $w$, so $L(D) = L(N)$. The DFA has $|2^Q| = 2^n$ states.
\end{proof}

This is a \rel{restricts} edge: the NFA needs no new primitive, only a relaxed structural requirement. The cost is representational---an exponential blow-up in states---not expressive.

\paragraph{The PDA--DFA edge: generative generalization.}
A PDA adds an unbounded stack: not the removal of a constraint but the introduction of a new primitive, the ability to store and retrieve unbounded LIFO information. By \cref{prop:anbn}, the consequence is a strict increase in power---$\{a^n b^n\}$ is context-free but not regular, so $\mathrm{REG} \subsetneq \mathrm{CFL}$. This is an \rel{extends\_grammar} edge, and the gap is witnessed by the failure of the pumping lemma, a property that holds for every regular language and fails for $\{a^n b^n\}$.

\paragraph{Why this distinction matters.}
When $B$ \rel{restricts} $A$, the theory of $A$ transfers intact---same characterization theorems, same closure properties, same complexity measures. When $B$ \rel{extends\_grammar} $A$, the old theory typically breaks: new characterization theorems are needed, and closure properties change (\cref{tab:closure}: context-free languages lose closure under intersection and complement, \cref{prop:closure}). The distinction reappears in the learning setting. Relaxing a convergence requirement (behaviorally correct learning relaxing syntactic convergence) is conservative; adding a resource (a teacher who supplies counterexamples) is generative. Recognizing which is at work predicts whether existing results transfer.

% ================================================================
% SECTION 3: Computability Foundations
% ================================================================

\section{Computability Foundations for Learning}
\label{sec:computability}

Gold's model requires learners to be effective procedures that output hypotheses from an effective enumeration. This section fixes the computability vocabulary; each notion is typed in the verified development's computability substrate.

\begin{defn}[G\"odel Numbering]
\label{def:godel-numbering}
A \emph{G\"odel numbering} of a countable set is an effective bijection with $\N$.\formal{FLT_Proofs.Computation.GodelNumbering} We write $\varphi_e$ for the partial computable function with index $e$, and $W_e = \mathrm{dom}(\varphi_e)$ for the $e$-th r.e.\ set; the sequences $\varphi_0, \varphi_1, \ldots$ and $W_0, W_1, \ldots$ are universal effective enumerations of all partial computable functions and all r.e.\ sets.
\end{defn}

In Gold's framework a learner reading a text for $L$ outputs at each step a number $e_t$, interpreted through the G\"odel numbering as the conjecture $L = W_{e_t}$. Identification in the limit means the sequence $e_0, e_1, \ldots$ stabilizes on an index $e^\ast$ with $W_{e^\ast} = L$. The learner is itself a program: a computable map $M : \Sigma^\ast \to \N$ from finite data sequences to hypothesis indices, and the admissible learners are exactly the total computable such maps. This effectivity requirement separates Gold's model from a purely set-theoretic notion of convergence.

% ================================================================
% SECTION 4: The Learning-Theoretic Bridge
% ================================================================

\section{The Learning-Theoretic Bridge: Limiting Recursion and $\Delta^0_2$}
\label{sec:bridge}

The connection between computability and learning is an equivalence that identifies Gold's criterion with a precise level of the arithmetical hierarchy.

\begin{defn}[Limiting Recursion]
\label{def:limiting-recursion}
A function $f : \N \to \N$ is \emph{limiting recursive} if there is a total computable $\hat{f} : \N \times \N \to \N$ with $f(x) = \lim_{s \to \infty} \hat{f}(x, s)$ for every $x$: the approximation $\hat{f}(x,0), \hat{f}(x,1), \ldots$ stabilizes on $f(x)$, though the stabilization time is not itself computable.\formal{FLT_Proofs.Computation.LimitingRecursive}
\end{defn}

\begin{defn}[$\Delta^0_2$ Class]
\label{def:delta-02}
A set $A \subseteq \N$ is \emph{$\Delta^0_2$} if both $A$ and $\bar{A}$ are $\Sigma^0_2$, that is, of the form $\{x : \exists y\, \forall z\, R(x,y,z)\}$ for a computable $R$.\formal{FLT_Proofs.Computation.Delta02Class}
\end{defn}

The two notions coincide, and one direction is a clean unfolding of the definitions; this is the half we prove. The converse is Shoenfield's limit lemma.

\begin{prop}[Limiting recursion gives $\Delta^0_2$]
\label{prop:limit-delta}
If the characteristic function $\chi_A$ of $A \subseteq \N$ is limiting recursive, then $A$ is $\Delta^0_2$.
\end{prop}

\begin{proof}
Let $\hat{\chi} : \N \times \N \to \{0,1\}$ be a total computable approximation with $\chi_A(x) = \lim_s \hat{\chi}(x,s)$ for all $x$. Because the limit exists, the approximation is eventually constant at each $x$, so
\[
x \in A \iff \exists s_0\, \forall s \geq s_0\ \hat{\chi}(x,s) = 1, \qquad
x \notin A \iff \exists s_0\, \forall s \geq s_0\ \hat{\chi}(x,s) = 0 .
\]
The right-hand side of the first equivalence is $\Sigma^0_2$ (an $\exists\forall$ over the computable predicate $\hat{\chi}(x,s)=1$), and the second exhibits $\bar{A}$ as $\Sigma^0_2$ likewise. Both $A$ and $\bar A$ are $\Sigma^0_2$, so $A$ is $\Delta^0_2$.
\end{proof}

\begin{thm}[Shoenfield's limit lemma {\cite{Shoenfield1959}}]
\label{thm:shoenfield}
$A \subseteq \N$ is $\Delta^0_2$ if and only if $\chi_A$ is limiting recursive, equivalently iff $A$ is computable from an oracle for the halting problem.
\end{thm}

The forward direction is \cref{prop:limit-delta}; the converse extracts a computable mind-changing approximation from a $\Delta^0_2$ definition, and the oracle reformulation is the relativization. The lemma is the hinge of the following bridge.

\begin{thm}[Limiting recursion and learnability {\cite{Gold1965,Putnam1965,Shoenfield1959}}]
\label{thm:limit-bridge}
A family of r.e.\ sets $\mathcal{L} = \{L_i\}_{i \in I}$ is identifiable in the limit (Gold) if and only if the index function $i \mapsto e_i$, returning a correct hypothesis index for each $L_i$, is limiting recursive; the classes identifiable in the limit are exactly those whose index sets are $\Delta^0_2$.
\end{thm}

On the computability side, $\Delta^0_2$ is the level computable with a halting oracle, equivalently the sets decidable by a procedure allowed finitely many mind changes (\cref{thm:shoenfield}). On the learning side, identification in the limit is exactly the requirement that the conjecture sequence stabilize, which is exactly limiting recursion. The learning half of the bridge is verified: explanatory (Gold) identifiability is equivalent to the existence of a learner converging with finitely many conjecture changes.\formal{FLT_Proofs.Theorem.Gold.mind_change_characterization} The number of changes is itself a complexity measure, recorded as an ordinal that is finite exactly when the learner converges correctly.\formal{FLT_Proofs.Complexity.MindChange.MindChangeOrdinal}

The equivalence settles three points used in \Cref{ch:gold}.

\begin{enumerate}[leftmargin=*,itemsep=3pt]
\item \textbf{Upper bound on learnability.} No class whose index set lies above $\Delta^0_2$ in the arithmetical hierarchy is identifiable in the limit. This bounds Gold-style learning from above.

\item \textbf{Mind changes as approximation stages.} The number of $s$ with $\hat{f}(x,s) \neq \hat{f}(x,s+1)$ is the number of mind changes the learner makes before converging. The mind-change ordinals of \Cref{ch:ordinals} refine this into a transfinite hierarchy.

\item \textbf{Identification in the limit is canonical.} Gold's criterion is the natural first level above the computable: computability-in-the-limit, equivalently $\Delta^0_2$-computability. It is not one convergence requirement among many but the $\Delta^0_2$ one.
\end{enumerate}

\section{What This Chapter Established}
\label{sec:ch3-summary}

The chapter fixed the definitions---DFA, NFA, PDA, Turing machine, CFG, G\"odel numbering, partial computable function---needed for the identification-in-the-limit chapters, each typed against the verified computability substrate, and proved the load-bearing facts in full. Two structural ideas carry forward.

\begin{enumerate}[leftmargin=*,itemsep=3pt]
\item \textbf{The \rel{restricts}/\rel{extends\_grammar} distinction.} The NFA--DFA pair (conservative, same power, \cref{thm:nfa-dfa}) and the PDA--DFA pair (generative, strictly more power, \cref{prop:anbn}) are the first concrete instance of an edge-type distinction that recurs throughout the graph. Whether a generalization is conservative or generative determines whether existing theory transfers or must be rebuilt.

\item \textbf{The $\Delta^0_2$ bridge.} Identification in the limit is a precise computability concept: computation in the limit, equivalent to $\Delta^0_2$-computability (\cref{thm:shoenfield,thm:limit-bridge}). The equivalence anchors Gold's model in the arithmetical hierarchy and grounds the mind-change hierarchy of \Cref{ch:ordinals}.
\end{enumerate}

\section{Exercises and Open Problems}
\label{sec:ch3-open}

Each problem is anchored to a concept above and tagged: a \emph{known unknown} (\textsc{ku}) is an articulated open question; an \emph{unknown known} (\textsc{uk}) is a result present in the literature or the verified development but not yet absorbed into a clean general statement. Verified handles are named where they exist.

\begin{openprob}[Formalizing the limit-lemma bridge, \textsc{uk}]
\label{op:bridge-formal}
\Cref{prop:limit-delta} proves one direction. The computability substrate types both sides (\leanname{FLT_Proofs.Computation.LimitingRecursive}, \leanname{FLT_Proofs.Computation.Delta02Class}) and verifies the learning half (\leanname{FLT_Proofs.Theorem.Gold.mind_change_characterization}). Formalize the converse of Shoenfield's limit lemma and assemble the full bridge of \cref{thm:limit-bridge} ($\mathrm{Ex}$-identifiable $\Leftrightarrow$ $\Delta^0_2$ index set) inside the development.
\end{openprob}

\begin{openprob}[Optimal NFA determinization blow-up, \textsc{ku}]
\label{op:subset-tight}
\Cref{thm:nfa-dfa} gives a DFA with at most $2^n$ states. Exhibit, for each $n$, an $n$-state NFA whose minimal equivalent DFA has exactly $2^n$ states (Moore, Meyer--Fischer), and formalize the matching lower bound on top of \leanname{FLT_Proofs.Computation.NFA'} and \leanname{FLT_Proofs.Computation.DFA'}, so that the conservative-generalization cost is pinned exactly.
\end{openprob}

\begin{openprob}[Myhill--Nerode and the exact regularity criterion, \textsc{uk}]
\label{op:myhill-nerode}
The pumping lemma (\cref{lemma:pumping}) is necessary but not sufficient for regularity. Formalize the Myhill--Nerode theorem---$L$ is regular iff its right-congruence has finite index, which also yields the unique minimal DFA---on top of \leanname{FLT_Proofs.Computation.IsRegularLang}. This is the structural foundation that the $L^\ast$ learner of \Cref{ch:exact} exploits.
\end{openprob}

\begin{openprob}[Arithmetical complexity of the criterion hierarchy, \textsc{ku}]
\label{op:criterion-hierarchy}
\Cref{thm:limit-bridge} places explanatory identification at $\Delta^0_2$. Locate the other Gold criteria---behaviorally correct, vacillatory, anomalous, finite---in the arithmetical hierarchy, and determine which strictly exceed $\Delta^0_2$. The criteria are typed (\leanname{FLT_Proofs.Criterion.Gold.BCLearnable} and siblings) but their hierarchy positions are not yet pinned.
\end{openprob}

\begin{openprob}[A functorial account of restricts vs.\ extends, \textsc{uk}]
\label{op:restricts-functor}
\Cref{def:restricts,def:extends-grammar} distinguish conservative from generative generalization informally. Make it precise: define a category of computational models and a notion of generalization morphism such that \rel{restricts} morphisms provably transport characterization theorems and \rel{extends\_grammar} morphisms provably need not, and decide whether the classification of a given pair is itself decidable.
\end{openprob}

\begin{openprob}[Learnability of context-free languages, \textsc{ku}]
\label{op:cfl-learnable}
$\{a^n b^n\}$ (\cref{prop:anbn}) is context-free. Characterize which subclasses of the context-free languages are identifiable in the limit from text and from informant, extending the text/informant separation of \Cref{ch:data} up one level of the Chomsky hierarchy, and relate the answer to the $\Delta^0_2$ ceiling of \cref{thm:limit-bridge}.
\end{openprob}

\begin{openprob}[The coding theorem for the computability substrate, \textsc{uk}]
\label{op:coding-theorem}
The substrate types Kolmogorov complexity and algorithmic probability (\leanname{FLT_Proofs.Computation.KolmogorovComplexity}, \leanname{FLT_Proofs.Computation.AlgorithmicProbability}). Formalize the coding theorem $K(x) = -\log \mathrm{P}(x) + O(1)$ relating them, with the constant explicit in the reference machine---the same gap left open for model selection in \Cref{ch:objects} (Problem in \S\ref{sec:ch1-open}).
\end{openprob}

\begin{openprob}[The transfinite mind-change hierarchy, \textsc{uk}]
\label{op:mind-change-ordinals}
The mind-change count is typed as an ordinal (\leanname{FLT_Proofs.Complexity.MindChange.MindChangeOrdinal}), finite exactly when the learner converges correctly. Formalize the Freivalds--Smith transfinite mind-change hierarchy---identification with at most $\alpha$ mind changes for a countable ordinal $\alpha$---and prove the strictness of the resulting hierarchy, which \Cref{ch:ordinals} develops informally.
\end{openprob}

\begin{openprob}[Closure of the recursive languages under learning operations, \textsc{ku}]
\label{op:recursive-closure}
\Cref{prop:closure} and \cref{tab:closure} record language-class closure. Determine the closure of the \emph{identifiable-in-the-limit} classes under union, intersection, and complement: if $\mathcal{C}_1$ and $\mathcal{C}_2$ are each identifiable from text, when is $\mathcal{C}_1 \cup \mathcal{C}_2$, and does the answer mirror the regular/context-free split of \cref{tab:closure}?
\end{openprob}

\begin{openprob}[Effective separation witnesses, \textsc{uk}]
\label{op:effective-witness}
\Cref{prop:anbn} witnesses $\mathrm{REG} \subsetneq \mathrm{CFL}$ by a single language. Give a uniform, effective procedure that, from a description of any two adjacent Chomsky levels, outputs a witness language separating them together with its non-membership proof, turning the proper-containment chain of \cref{fig:chomsky} into one constructive theorem rather than four independent separations.
\end{openprob}

% Chapter 4: Learners and Teachers
% Rewritten against the design-lab Lean 4 kernel (FLT_Proofs.Learner.*) and the
% transformers library (TLT_Proofs): typed learner abstraction, the Gibbs variational
% principle in full, and a modern attention-learner instance of MeasurableBatchLearner.
%   14 x Profile A (vocabulary: learner types and teacher types - TABLE, not definitions)
%   bayesian_learner = Profile E (Obstruction: posterior ≠ single hypothesis)
%   gibbs_posterior = Profile C (Revelation: the Bayesian-PAC bridge)
%   meta_learner = Profile F (Open Frontier: learning to learn)
%   team_learner = Profile C (Revelation: teams > individuals)
%   optimal_teacher = Profile G (Framework Bridge → teaching_dimension)
%   minimally_adequate_teacher = Profile G (Framework Bridge → L*)

\chapter{Learners and Teachers}
\label{ch:agents}

Chapters~\ref{ch:objects} and~\ref{ch:data} established what is learned (concept classes) and how data arrives (i.i.d.\ samples, texts, queries, streams). This chapter introduces the agents that do the learning and, in some settings, the agents that provide the data.

Much of the vocabulary here is taxonomic: a learner is a function from data to hypotheses, and the field has named a dozen variants distinguished by which constraints they satisfy. These are best presented as a table. The chapter's mathematical content lies in three places: the type mismatch between Bayesian posteriors and PAC hypotheses, resolved by the Gibbs posterior, whose optimality we prove; the power of team learning; and the teacher models that connect to the teaching dimension and to Angluin's $L^\ast$. We close the opening section with a modern instance, a finite-head attention learner, that satisfies the measurability condition the PAC theory requires.

% ================================================================
% SECTION 1: The Learner Abstraction
% ================================================================

\section{The Learner Abstraction}
\label{sec:learner}

\begin{defn}[Learner]
\label{def:learner}
A \emph{learner} is a function
\[
L \colon \bigcup_{m=0}^{\infty} (X \times Y)^m \;\to\; \mathcal{H}
\]
mapping a finite data sequence to a hypothesis: having observed $S_t = ((x_1, y_1), \ldots, (x_t, y_t))$, it outputs $h_t = L(S_t) \in \mathcal{H}$. The verified development splits this single abstraction into the three paradigm-specific carriers: a \emph{batch learner} maps a whole sample to a hypothesis,\formal{FLT_Proofs.Learner.Core.BatchLearner} an \emph{online learner} carries a state updated example by example,\formal{FLT_Proofs.Learner.Core.OnlineLearner} and a \emph{Gold learner} conjectures from the data seen so far.\formal{FLT_Proofs.Learner.Core.GoldLearner}
\end{defn}

This definition is minimal. It says nothing about how $h_t$ is selected: whether the learner stores the full history or only its previous guess, whether it uses randomness, whether it may choose which $x$ to query. Each constraint yields a named variant (Figure~\ref{fig:learner-taxonomy} and Table~\ref{tab:learner-types}). One further condition is invisible at the level of types but indispensable for the statistical theory.

\begin{defn}[Measurable batch learner]
\label{def:measurable-learner}
A batch learner $L$ over a measurable domain is a \emph{measurable batch learner} if its evaluation map is jointly measurable: for every sample size $m$, the map $(S, x) \mapsto L(S)(x)$ from $(X \times Y)^m \times X$ to $Y$ is measurable.\formal{FLT_Proofs.Learner.Core.MeasurableBatchLearner}
\end{defn}

This is the minimal regularity that makes the PAC success event $\{S : \Pr_x[L(S)(x) \neq c(x)] \leq \varepsilon\}$ a measurable set, so that ``with probability $1 - \delta$'' is well-defined; without it the PAC criterion of \Cref{ch:pac} cannot even be stated~\cite{KrappWirth2024}. Every theorem in the statistical chapters carries this hypothesis silently, and the learners that matter in practice are exactly the ones for which measurability is not obvious.

\subsection{A modern instance: the attention learner}
\label{sec:attention-learner}

A transformer's attention layer is a learner in the sense of \cref{def:learner}: its parameters are fit to data, and at inference it routes each input through one of several ``heads.'' Whether this learner is admissible for the statistical theory is exactly the question of \cref{def:measurable-learner}, and the answer is yes.

\begin{prop}[A finite-head attention learner is a measurable batch learner]
\label{prop:attention-mbl}
Fix $k \geq 1$. A \emph{$k$-head attention learner} is given by a standard-Borel parameter space $\Theta$, measurable score functions $s_\theta : X \to \R^k$, measurable experts $e_\theta : X \to Y^k$, and a measurable parameter-selection rule $\theta(\cdot)$ from samples; its prediction routes $x$ to the expert of highest score, $L(S)(x) = e_{\theta(S)}\bigl(x\bigr)_{\arg\max_i s_{\theta(S)}(x)_i}$. Then $L$ is a measurable batch learner, so every PAC theorem of \Cref{ch:pac} applies to it.\formal{TLT_Proofs.Learner.MeasurableParametricAttentionLearner.ParametricFiniteHeadAttentionLearner.instMBL'}
\end{prop}

\begin{proof}
The argmax router partitions the input space into the cells $\{x : i = \arg\max_j s_\theta(x)_j\}$, each a finite intersection of the measurable score-comparison sets $\{x : s_\theta(x)_i \geq s_\theta(x)_j\}$; hence the routing map $(\theta, x) \mapsto \mathrm{route}_\theta(x)$ is measurable.\formal{TLT_Proofs.Routing.MeasurableScoreRouting.FiniteScoreRouterCode.route_measurable} Composing the measurable router with the measurable experts gives a jointly measurable evaluation $(\theta, x) \mapsto e_\theta(x)_{\mathrm{route}_\theta(x)}$. Precomposing with the measurable selection $S \mapsto \theta(S)$ in the first coordinate preserves joint measurability, so $(S,x) \mapsto L(S)(x)$ is measurable, which is exactly \cref{def:measurable-learner}. The whole construction is a parametric learner family, and every parametric family with a measurable evaluation and selection is a measurable batch learner.\formal{TLT_Proofs.Learner.MeasurableParametricAttentionLearner.ParametricLearnerFamily.instMeasurableBatchLearner}
\end{proof}

The same holds for the softmax-top-1 and multi-head argmax routers, and for the scaled dot-product attention of a full transformer.\formal{TLT_Proofs.Routing.ScaledDotProductScoreRouter.attentionRouting_wellBehaved} The abstract learner of \cref{def:learner} therefore covers the dominant modern architectures, and places them in the measurable sub-class on which the PAC apparatus operates. We return to the measurability boundary, where attention routing \emph{fails} to be Borel, in \Cref{ch:separations}.

% ================================================================
% SECTION 2: Taxonomy Table
% ================================================================

\section{Taxonomy of Learner Types}
\label{sec:learner-taxonomy}

\begin{figure}[ht]
\centering
\begin{tikzpicture}[
    font=\small,
    root/.style={draw, very thick, rounded corners, fill=defblue!15, font=\small\bfseries,
                 minimum width=2.4cm, minimum height=0.6cm, align=center},
    pred/.style={draw, semithick, rounded corners, fill=layergray, font=\small,
                 minimum width=2.5cm, minimum height=0.7cm, align=center},
    sink/.style={draw, thick, rounded corners, align=center, font=\scriptsize,
                 text width=3.3cm, inner sep=4pt},
    gedge/.style={-{Stealth[length=2mm]}, sepgreen!60!black, semithick},
    oedge/.style={-{Stealth[length=2mm]}, edgeorange!80!black, semithick, dashed},
    meet/.style={notegray, dotted, thick},
    elbl/.style={font=\scriptsize\itshape, fill=white, inner sep=1pt, text=black!65}
]
% --- root anchor (held-constant Will) ---
\node[root] (root) at (0,3.0) {Learner\\[-1pt]{\scriptsize\mdseries one learner, many constraints}};
% --- predicate spine (verified typed predicates, solid-outlined nodes) ---
\node[pred] (cons)  at (0,1.55) {Consistent\\[-2pt]{\tiny\ttfamily IsConsistent}};
\node[pred] (consv) at (0,0.45) {Conservative\\[-2pt]{\tiny\ttfamily IsConservative}};
\node[pred] (sd)    at (0,-0.65) {Set-driven\\[-2pt]{\tiny\ttfamily IsSetDriven}};
\node[pred] (it)    at (0,-1.75) {Iterative\\[-2pt]{\tiny\ttfamily IsIterative}};
% --- faint lattice-meet connectors (constraints on one learner) ---
\draw[meet] (root) -- (cons);
\draw[meet] (cons) -- (consv);
\draw[meet] (consv) -- (sd);
\draw[meet] (sd) -- (it);
% --- regime sinks ---
\node[sink, sepgreen!60!black, fill=sepgreen!8] (pac) at (-4.7,0.1)
  {\textbf{In PAC: constraints are FREE}\\[2pt]learnable by a consistent proper ERM
   (\leanname{ermLearn}); i.i.d.\ makes order irrelevant\\[2pt]
   verified: \leanname{erm_consistent_realizable}};
\node[sink, edgeorange!80!black, fill=edgeorange!8, dashed] (gold) at (4.7,-0.55)
  {\textbf{In Gold/online: constraints BITE}\\[2pt]conservative from informant, not text;
   set-driven forces unbounded memory on iterative\\[2pt]
   literature: Gold; Angluin; K\"otzing--Palenta};
% --- green solid edges: free in PAC ---
\draw[gedge] (cons) -- (pac);
\draw[gedge] (sd)   -- node[elbl, above, pos=0.45]{i.i.d.} (pac);
\draw[gedge] (it)   -- (pac);
% --- orange dashed edges: biting in Gold ---
\draw[oedge] (consv) -- (gold);
\draw[oedge] (sd)    -- node[elbl, above, pos=0.45]{memory} (gold);
\draw[oedge] (it)    -- (gold);
% --- legend ---
\node[font=\scriptsize, text=black!60, align=center] at (0,-3.05)
  {Solid green: free in PAC (kernel-verified).\quad Dashed orange: biting in Gold (literature).\quad
   Boxes: verified predicates.};
\end{tikzpicture}
\caption[Lattice of learner constraints: free in PAC, biting in Gold]{The learner taxonomy as a
lattice of typed constraints on the base learner of \cref{def:learner}. The same four predicates read
two ways: in PAC they are free, since every PAC-learnable class is learnable by a consistent, proper
empirical-risk minimizer (\leanname{erm_consistent_realizable}) and the i.i.d.\ assumption makes order
irrelevant; in Gold's framework they bite, separating conservative from text-only learners and forcing
unbounded memory on iterative ones. The named data-access and randomness variants are in
\cref{tab:learner-types}.}
\label{fig:learner-taxonomy}
\end{figure}

The named variants are constraints on \cref{def:learner}, and each is a typed predicate on the learner in the verified development: consistency,\formal{FLT_Proofs.Learner.Properties.IsConsistent} conservativeness,\formal{FLT_Proofs.Learner.Properties.IsConservative} set-drivenness,\formal{FLT_Proofs.Learner.Properties.IsSetDriven} and iterativeness.\formal{FLT_Proofs.Learner.Properties.IsIterative} Table~\ref{tab:learner-types} gives the precise constraints.

\begin{table}[ht]
\centering
\small
\caption{Named learner types. Each row gives the defining constraint and the chapter where the concept plays its primary role. These are not independent definitions; a single learner may satisfy several constraints simultaneously (e.g., a conservative, set-driven, passive learner).}
\label{tab:learner-types}
\begin{tabularx}{\textwidth}{>{\bfseries\small}l X l}
\toprule
\textbf{Name} & \textbf{Defining Constraint} & \textbf{Primary Role} \\
\midrule
Passive &
  Receives data without choosing which instances to observe. The standard assumption in PAC and Gold settings. &
  Ch.~\ref{ch:pac}, \ref{ch:gold} \\[4pt]

Active &
  Chooses which instances to query. Performance measured by \emph{label complexity}: the number of labels requested, not the number of unlabeled instances seen. &
  Ch.~\ref{ch:exact} \\[4pt]

Conservative &
  Changes hypothesis only when forced by inconsistency: if $h_t$ is consistent with $S_{t+1}$, then $h_{t+1} = h_t$. &
  Ch.~\ref{ch:gold} \\[4pt]

Consistent &
  Always outputs a hypothesis consistent with all data seen so far: $h_t(x_i) = y_i$ for all $i \leq t$. &
  Ch.~\ref{ch:pac}, \ref{ch:gold} \\[4pt]

Set-driven &
  Output depends only on the \emph{set} $\{(x_i, y_i)\}$ of data, not on the order of presentation. Formally, $L(S) = L(S')$ whenever $S$ and $S'$ are permutations of each other. &
  Ch.~\ref{ch:gold} \\[4pt]

Iterative &
  Hypothesis at step $t$ depends only on $h_{t-1}$ and the new datum $(x_t, y_t)$, not on the full history $S_t$. Memory is bounded: the learner's state is a single hypothesis. &
  Ch.~\ref{ch:gold}, \ref{ch:online} \\[4pt]

Probabilistic &
  May use internal randomness in selecting hypotheses. The success criterion (PAC, identification, mistake bound) is then required to hold with high probability over both the data and the learner's coin flips. &
  Ch.~\ref{ch:pac} \\[4pt]

With advice &
  Receives $b$ bits of advice in addition to data. The advice may encode information about the target class or the distribution, and $b$ parameterizes the amount of side information. &
  Ch.~\ref{ch:exact} \\[4pt]

\bottomrule
\end{tabularx}
\end{table}

These distinctions generate strict separations in learning power within Gold's framework: there are classes identifiable by a conservative learner from informant but not from text, and classes identifiable by a set-driven learner that require unbounded memory for an iterative learner (\Cref{ch:gold,ch:separations}).

In the PAC setting many of the distinctions collapse, and the collapse has the status of a proved theorem: every PAC-learnable class is learnable by a consistent, proper learner (the empirical-risk minimizer, which returns any hypothesis consistent with the sample and succeeds once uniform convergence holds (\cref{thm:uc-pac})).\formal{FLT_Proofs.Complexity.Generalization.ermLearn} Because the i.i.d.\ assumption makes example order irrelevant, the set-driven/order-sensitive distinction also vanishes. The taxonomy has its sharpest teeth in the Gold and online paradigms, where the learner's memory and update policy genuinely affect what is learnable.

% ================================================================
% SECTION 3: The Bayesian–PAC Bridge
% ================================================================

\section{The Bayesian--PAC Bridge}
\label{sec:bayesian-pac}

The learner types in Table~\ref{tab:learner-types} all output a single hypothesis $h \in \mathcal{H}$. The Bayesian learner does not.

\begin{defn}[Bayesian Learner]
\label{def:bayesian-learner}
A \emph{Bayesian learner} maintains a prior $\pi$ over $\mathcal{H}$ and, given data $S$, outputs the posterior
\[
Q(h \mid S) \;=\; \frac{\pi(h) \cdot P(S \mid h)}{\sum_{h' \in \mathcal{H}} \pi(h') \cdot P(S \mid h')},
\]
a whole distribution over hypotheses, with every $h$ retained at its posterior weight.\formal{FLT_Proofs.Learner.Bayesian.BayesianLearner}
\end{defn}

\subsection{The type mismatch}

PAC learning requires a learner that outputs a single $h \in \mathcal{H}$ with $\risk_D(h) \leq \risk_D(h^\ast) + \varepsilon$ at confidence $1 - \delta$. The Bayesian learner outputs a probability measure on $\mathcal{H}$. These are different types:
\[
\text{PAC learner:} \quad (X \times Y)^m \to \mathcal{H},
\qquad
\text{Bayesian learner:} \quad (X \times Y)^m \to \Delta(\mathcal{H}).
\]
The naive resolution, outputting the MAP hypothesis $\hat{h} = \arg\max_h Q(h \mid S)$, discards most of the posterior and yields suboptimal generalization. The verified development records both the MAP read-off and the conversion of a Bayesian learner to an ordinary batch learner.\formal{FLT_Proofs.Bridge.bayesianToBatch} The sharper resolution is the Gibbs posterior.

\begin{obstbox}
\textbf{Bayesian $\to$ PAC type mismatch.} The Bayesian learner's output type ($\Delta(\mathcal{H})$) does not match the PAC learner's output type ($\mathcal{H}$). The PAC-Bayes theorem (\Cref{ch:generalization}) bounds the risk of a \emph{distribution} over hypotheses; extracting a single hypothesis with the bound intact requires an explicit randomized construction: the Gibbs posterior.
\end{obstbox}

\subsection{The Gibbs posterior and its optimality}

\begin{defn}[Gibbs Posterior]
\label{def:gibbs-posterior}
Given a prior $\pi$ over $\mathcal{H}$, the empirical risk $\emprisk_S$, and an inverse temperature $\lambda > 0$, the \emph{Gibbs posterior} is
\[
Q_\lambda(h) \;\propto\; \pi(h) \cdot \exp\!\bigl(-\lambda \,\emprisk_S(h)\bigr).
\]
The \emph{Gibbs classifier} draws $h \sim Q_\lambda$ and predicts with $h$; its expected risk is $\E_{h \sim Q_\lambda}[\risk_D(h)]$.\formal{FLT_Proofs.Learner.Bayesian.GibbsPosterior}
\end{defn}

The PAC-Bayes theorem~\cite{McAllester1999} states that for any prior $\pi$ and any posterior $Q$, with probability $1 - \delta$ over $S \sim D^m$,
\[
\E_{h \sim Q}[\risk_D(h)] \;\leq\; \E_{h \sim Q}[\emprisk_S(h)] + \sqrt{\frac{\KL(Q \,\|\, \pi) + \ln(m/\delta)}{2m}},
\]
proved for finite hypothesis spaces in the verified development.\formal{FLT_Proofs.Theorem.PACBayes.pac_bayes_finite} The bound holds for \emph{every} $Q$, and so leaves open which $Q$ the learner should use. The Gibbs posterior is the answer, by a variational argument.

\begin{prop}[The Gibbs posterior minimizes the PAC-Bayes objective {\cite{Catoni2007}}]
\label{prop:gibbs-variational}
Let $\mathcal{H}$ be finite and $\lambda > 0$. Among all distributions $Q$ on $\mathcal{H}$, the objective
\[
F_\lambda(Q) \;=\; \lambda\,\E_{h \sim Q}[\emprisk_S(h)] + \KL(Q \,\|\, \pi)
\]
is minimized uniquely by the Gibbs posterior $Q_\lambda$, with minimum value $-\ln Z_\lambda$ where $Z_\lambda = \sum_{h} \pi(h)\,e^{-\lambda \emprisk_S(h)}$.
\end{prop}

\begin{proof}
Write $Q_\lambda(h) = \pi(h)\,e^{-\lambda \emprisk_S(h)} / Z_\lambda$. For any distribution $Q$ on $\mathcal{H}$,
\[
\KL(Q \,\|\, \pi) + \lambda\,\E_{h\sim Q}[\emprisk_S(h)]
= \sum_h Q(h)\Bigl[\ln \tfrac{Q(h)}{\pi(h)} + \lambda\,\emprisk_S(h)\Bigr]
= \sum_h Q(h)\,\ln \frac{Q(h)}{\pi(h)\,e^{-\lambda \emprisk_S(h)}}.
\]
Substituting $\pi(h)\,e^{-\lambda \emprisk_S(h)} = Z_\lambda\,Q_\lambda(h)$ and splitting the logarithm,
\[
= \sum_h Q(h)\,\ln \frac{Q(h)}{Z_\lambda\,Q_\lambda(h)}
= \sum_h Q(h)\,\ln \frac{Q(h)}{Q_\lambda(h)} - \ln Z_\lambda
= \KL(Q \,\|\, Q_\lambda) - \ln Z_\lambda .
\]
By Gibbs' inequality $\KL(Q \,\|\, Q_\lambda) \geq 0$, with equality if and only if $Q = Q_\lambda$. Hence $F_\lambda(Q) \geq -\ln Z_\lambda$, with equality exactly at $Q = Q_\lambda$.
\end{proof}

The PAC-Bayes bound's right-hand side is, up to the monotone square root and constants, an instance of $F_\lambda$: minimizing the complexity-penalized empirical risk over posteriors yields $Q_\lambda$. The bound thus names the exact posterior to use, which is the precise sense in which the Gibbs posterior converts a Bayesian learner into a PAC learner.\formal{FLT_Proofs.Theorem.PACBayes.gibbs_bound_of_pointwise}

\begin{remark}[Gibbs $\neq$ MAP]
\label{rmk:gibbs-vs-map}
The Gibbs posterior draws a \emph{random} hypothesis weighted by empirical performance; MAP takes the \emph{mode}. In high-dimensional settings the Gibbs posterior gives tighter guarantees because PAC-Bayes controls expected risk under the full posterior, while a MAP bound requires uniform convergence over a point estimate. As $\lambda \to \infty$ the Gibbs posterior concentrates on the empirical risk minimizers and recovers MAP in the limit; for finite $\lambda$ it strictly hedges, which is the source of the improved bound.
\end{remark}

The full development of PAC-Bayes bounds, including the McAllester and Catoni variants, appears in \Cref{ch:generalization}.

% ================================================================
% SECTION 4: Teams and Meta-Learning
% ================================================================

\section{Teams and Meta-Learning}
\label{sec:teams-meta}

\subsection{Team learning}

The team model asks what happens when several learners collaborate, where success requires only that \emph{at least one} member converges to a correct hypothesis.

\begin{defn}[Team Learner]
\label{def:team-learner}
A \emph{team} of size $k$ is a tuple of learners $(L_1, \ldots, L_k)$. The team \emph{identifies} $c \in \mathcal{C}$ if at least one $L_i$ converges to a correct hypothesis; $\mathcal{C}$ is \emph{team-identifiable} if some finite team identifies every $c \in \mathcal{C}$.\formal{FLT_Proofs.Learner.Properties.TeamLearner}
\end{defn}

\begin{thm}[Teams $>$ Individuals {\cite{Smith1982}}]
\label{thm:team-power}
There exist concept classes identifiable by a team of size~2 but by no single learner.
\end{thm}

The team members receive the \emph{same} data; they neither partition the instance space nor communicate, so the power comes from the disjunctive success criterion alone: ``at least one succeeds'' is a weaker demand than ``the one learner succeeds,'' and the weaker demand enlarges the learnable classes. The separation lives inside Gold's standard framework, with ordinary recursively enumerable languages as the witnesses. The proof builds a class on which any single learner is forced into infinitely many mind changes by diagonalization (as in the text/informant separation of \Cref{ch:data}), while two learners with complementary strategies (one guessing large languages, one guessing small) guarantee that one converges. The argument appears in \Cref{ch:separations}.

\begin{remark}[Teams and probabilistic learners]
\label{rmk:team-prob}
A team of deterministic learners is strictly stronger than a single deterministic learner (\cref{thm:team-power}). A probabilistic learner succeeding with probability $> 1/2$ can simulate a $k$-team by running $k$ internal copies; but in the Gold setting, where success means convergence with certainty, this simulation is unavailable, and the team hierarchy is strict.
\end{remark}

\subsection{Meta-learning}

\begin{defn}[Meta-Learner]
\label{def:meta-learner}
A \emph{meta-learner} operates over a family of tasks $\{\mathcal{C}_j\}$ drawn from an environment $\mathcal{E}$: from experience on $\mathcal{C}_1, \ldots, \mathcal{C}_{n-1}$ it produces a bias (a hypothesis space or prior) that performs well on a fresh task $\mathcal{C}_n$ from the same environment.\formal{FLT_Proofs.Theorem.Extended.MetaLearnerPAC}
\end{defn}

Baxter~\cite{Baxter2000} established the first PAC-style bounds for meta-learning: if $\mathcal{E}$ has bounded complexity (covering numbers over a family of hypothesis spaces) and the meta-learner sees $n$ tasks of $m$ examples each, the expected excess risk on a new task is
\[
O\!\left(\sqrt{\frac{\mathrm{complexity}(\mathcal{E})}{n}} + \sqrt{\frac{\VC(\mathcal{H})}{m}}\right),
\]
a \emph{task-level} term (estimating the environment from $n$ tasks) plus a \emph{within-task} term (learning one task from $m$ examples). This two-level decomposition is the defining feature of meta-learning bounds, and whether a single combinatorial dimension characterizes meta-learnability is open (\cref{op:meta-char}).

% ================================================================
% SECTION 5: Teachers and the Teaching Game
% ================================================================

\section{Teachers and the Teaching Game}
\label{sec:teachers}

In the models above, data arrives passively or through the learner's queries. The \emph{teacher} reverses the direction: an agent chooses examples to help (or hinder) the learner.

\begin{table}[ht]
\centering
\small
\caption{Teacher types. Each row gives the teacher's strategy, the induced data model, and the chapter where the concept plays its primary role.}
\label{tab:teacher-types}
\begin{tabularx}{\textwidth}{>{\bfseries\small}l X l}
\toprule
\textbf{Name} & \textbf{Strategy and Key Property} & \textbf{Primary Role} \\
\midrule
Random &
  The \emph{random teacher}\index{random teacher} draws examples i.i.d.\ from a distribution $D$. The standard PAC assumption; the teacher has no intention, and ``teaching'' is an anthropomorphic misnomer for sampling. &
  Ch.~\ref{ch:pac} \\[4pt]

Adversarial &
  Chooses worst-case examples adaptively. The online model treats the data stream as an adversarial teacher; performance is worst-case mistake count. &
  Ch.~\ref{ch:online} \\[4pt]

Optimal &
  Selects the smallest set of examples from which the target is uniquely identified. The size of this set is the \emph{teaching dimension} of the concept within the class. &
  Ch.~13 \\[4pt]

Minimally adequate &
  Angluin's model: answers membership and equivalence queries. ``Minimally adequate'' because MQ + EQ is the minimal oracle configuration that makes DFAs efficiently learnable. &
  Ch.~\ref{ch:exact} \\[4pt]
\bottomrule
\end{tabularx}
\end{table}

\subsection{The optimal teacher and teaching dimension}

\begin{defn}[Teaching Dimension]
\label{def:teaching-dim}
The \emph{teaching dimension} of $c \in \mathcal{C}$ is the least number of labeled examples that uniquely identify $c$ within $\mathcal{C}$:
\[
\mathrm{TD}(c, \mathcal{C}) = \min\{|S| : S \subseteq X \times Y,\; c \text{ is the only } c' \in \mathcal{C} \text{ consistent with } S\}.
\]
The teaching dimension of the class is $\mathrm{TD}(\mathcal{C}) = \max_{c \in \mathcal{C}} \mathrm{TD}(c, \mathcal{C})$, and the \emph{optimal teacher} for $c$ achieves $\mathrm{TD}(c, \mathcal{C})$.\formal{FLT_Proofs.Complexity.Structures.TeachingDim}
\end{defn}

The teaching dimension measures the cooperative communication complexity of learning: how many examples a helpful teacher must provide so that a learner who knows $\mathcal{C}$ can identify the target. It is developed in full in \Cref{ch:ordinals}, where it meets the recursive teaching dimension and the teaching--learning duality.

\subsection{The minimally adequate teacher}

\begin{defn}[Minimally Adequate Teacher {\cite{Angluin1988}}]
\label{def:mat-agents}
A \emph{minimally adequate teacher} (MAT) for $\mathcal{C}$ provides a membership oracle $\mathrm{MQ}$ and an equivalence oracle $\mathrm{EQ}$ (\Cref{ch:data}). A class is MAT-learnable if some algorithm exactly identifies any $c \in \mathcal{C}$ using polynomially many queries.
\end{defn}

Angluin's $L^\ast$ algorithm (\cref{thm:lstar}) makes DFAs MAT-learnable. The teacher is ``minimal'' precisely: dropping either oracle makes the class unlearnable in polynomial time (\cref{prop:mq-alone,prop:eq-alone}). The full treatment of $L^\ast$, including the observation table and the Myhill--Nerode connection, is in \Cref{ch:exact}.

% ================================================================
% SECTION 6: Deferred Agents
% ================================================================

\section{Deferred Agent Types}
\label{sec:deferred}

Three agent concepts in the knowledge graph, \nde{synthesizer} (a program-synthesis agent), \nde{verifier} (a correctness checker for synthesizer output), and \nde{llm\_critic} (a large language model used as a critic in a learning loop), belong to the application layer. They are built \emph{from} the learner and teacher primitives above, and their treatment is deferred to \Cref{ch:frontiers}, in the context of program synthesis, neurosymbolic learning, and LLM-in-the-loop verification.

% ================================================================
% SECTION 7: Summary
% ================================================================

\section{What This Chapter Established}
\label{sec:ch4-summary}

The taxonomy rests on three substantive points.

\begin{enumerate}[leftmargin=*,itemsep=3pt]
\item \textbf{Most learner types are constraints on one abstraction.} The eight variants in Table~\ref{tab:learner-types} are orthogonal constraints (data access, state management, hypothesis selection) on the base type $L \colon (X \times Y)^\ast \to \mathcal{H}$, each a typed predicate in the development. The measurability condition (\cref{def:measurable-learner}) is the silent hypothesis of the statistical chapters, and a finite-head attention learner satisfies it (\cref{prop:attention-mbl}), placing transformers inside the abstraction on which the PAC apparatus operates.

\item \textbf{The Bayesian--PAC bridge is a variational fact.} The Bayesian learner outputs a posterior $Q \in \Delta(\mathcal{H})$; PAC learning wants a single $h \in \mathcal{H}$. The Gibbs posterior resolves this as the unique minimizer of the complexity-penalized empirical risk (\cref{prop:gibbs-variational}), which is why it is the right posterior to feed the PAC-Bayes bound.

\item \textbf{Teams are strictly more powerful than individuals.} In the Gold setting, two learners identify classes no single learner can (\cref{thm:team-power}); the mechanism is disjunctive success alone, the same data reaching both members, who never split the instance space or exchange a message.
\end{enumerate}

Teachers range from passive (random) through adversarial to cooperative (optimal, minimally adequate), connecting this chapter to the teaching dimension of \Cref{ch:ordinals} and the exact learning algorithms of \Cref{ch:exact}.

\section{Exercises and Open Problems}
\label{sec:ch4-open}

Each problem is anchored to a concept above and tagged: a \emph{known unknown} (\textsc{ku}) is an articulated open question; an \emph{unknown known} (\textsc{uk}) is a result in the literature or the verified development not yet absorbed into a clean general statement. Verified handles are named where they exist.

\begin{openprob}[A combinatorial dimension for meta-learning, \textsc{ku}]
\label{op:meta-char}
No analogue of the fundamental theorem of PAC learning exists for meta-learning: Baxter's bounds are stated through covering numbers of the environment, and no single dimension is known to stand in their place. Find a single complexity measure of the environment $\mathcal{E}$ that characterizes meta-learnability the way $\VC$ characterizes PAC learnability, building on the typed \leanname{FLT_Proofs.Theorem.Extended.MetaLearnerPAC}, or prove no such measure exists.
\end{openprob}

\begin{openprob}[A combinatorial characterization of MAT-learnability, \textsc{ku}]
\label{op:mat-char}
Which classes are learnable from a minimally adequate teacher (\cref{def:mat-agents}) in polynomial time? Beyond DFAs the picture is partial: decision trees, certain Boolean formulas, and residual finite-state automata are known positive. Give a combinatorial characterization analogous to the VC dimension for PAC learning, or show the query model admits none.
\end{openprob}

\begin{openprob}[Capacity of attention learners, \textsc{uk}]
\label{op:attention-capacity}
\Cref{prop:attention-mbl} shows a finite-head attention learner is a measurable batch learner, so the PAC theory applies. Determine its sample complexity: bound the VC or Rademacher complexity of the $k$-head attention hypothesis class in terms of $k$, the parameter dimension, and the score temperature, using the chaining and covering-number machinery already proved for attention (\leanname{TLT_Proofs.Capacity.Chaining.RademacherSymmetrization}), and decide whether the dependence on $k$ is logarithmic or linear.
\end{openprob}

\begin{openprob}[Optimal temperature for the Gibbs posterior, \textsc{ku}]
\label{op:gibbs-temperature}
\Cref{prop:gibbs-variational} fixes the optimal posterior for each $\lambda$. Determine the $\lambda$ that minimizes the resulting PAC-Bayes risk bound as a function of sample size $m$ and prior $\pi$, and characterize when the optimal $\lambda$ is finite (strict hedging) versus $\infty$ (MAP), sharpening \cref{rmk:gibbs-vs-map}.
\end{openprob}

\begin{openprob}[The team hierarchy, \textsc{ku}]
\label{op:team-hierarchy}
\Cref{thm:team-power} separates teams of size~2 from single learners. Is the team hierarchy strict at every level: does a team of size $k+1$ identify classes that no team of size $k$ can, and is there a single parameter of $\mathcal{C}$ that determines the least team size required? Formalize the hierarchy on top of \leanname{FLT_Proofs.Learner.Properties.TeamLearner}.
\end{openprob}

\begin{openprob}[Teaching dimension versus VC dimension, \textsc{ku}]
\label{op:td-vc}
The teaching dimension (\cref{def:teaching-dim}) and the VC dimension both measure a class, but neither bounds the other in general. Characterize the classes for which $\mathrm{TD}(\mathcal{C}) = O(\VC(\mathcal{C}))$, and determine the extremal gap between the two over classes of fixed VC dimension, connecting to the recursive teaching dimension of \Cref{ch:ordinals}.
\end{openprob}

\begin{openprob}[Measurable selection for Bayesian learners, \textsc{uk}]
\label{op:bayes-measurable}
\Cref{def:measurable-learner} is the regularity the PAC theory needs, and \cref{def:bayesian-learner} outputs a posterior. Determine conditions on the prior and likelihood under which the MAP read-off (\leanname{FLT_Proofs.Bridge.bayesianToBatch}) is a \emph{measurable} batch learner, so that Bayesian learners inherit the PAC apparatus without passing through the Gibbs construction.
\end{openprob}

\begin{openprob}[Order-sensitivity outside i.i.d., \textsc{ku}]
\label{op:order-sensitivity}
In the PAC setting the set-driven/iterative distinction collapses because example order is irrelevant. Quantify the cost of order-sensitivity in the online and Gold settings: give a class on which a set-driven learner needs exponentially less memory than any iterative learner, and characterize the memory-versus-order trade-off.
\end{openprob}

\begin{openprob}[A measurability dimension for learners, \textsc{uk}]
\label{op:measurability-dimension}
\Cref{def:measurable-learner} is binary: a learner is measurable or not. The attention router is measurable (\cref{prop:attention-mbl}) but lies one quantifier away from a non-Borel construction (\Cref{ch:separations}). Define a graded measurability complexity for learner families that places attention precisely, and relate it to the analytic-but-not-Borel boundary the development witnesses for attention routing.
\end{openprob}

\begin{openprob}[Teams of attention learners, \textsc{uk}]
\label{op:team-attention}
Combine \cref{thm:team-power,prop:attention-mbl}: does a team of finite-head attention learners (a mixture-of-experts ensemble selected by the disjunctive criterion) identify classes that no single attention learner can, and is the resulting team itself a measurable batch learner? This connects the team hierarchy to the routed-expert architectures of modern practice.
\end{openprob}

\subsection*{Counterfactual exercises}

Each exercise imposes a constraint from Figure~\ref{fig:learner-taxonomy} on the base learner and asks how its effect differs between the PAC and Gold paradigms.

\begin{enumerate}[leftmargin=*,itemsep=3pt]
\item Impose consistency (\leanname{FLT_Proofs.Learner.Properties.IsConsistent}) on a PAC learner. Show it costs nothing: every PAC-learnable class is learnable by a consistent, proper empirical-risk minimizer. Trace the result through \leanname{erm_consistent_realizable}, \leanname{empError_zero_iff_consistent}, and \leanname{fundamental_theorem}.

\item Read conservativeness (\leanname{FLT_Proofs.Learner.Properties.IsConservative}) in the Gold setting. Some classes a conservative learner identifies from an informant it cannot identify from text (\cite{Angluin1980,Gold1967,CaseSmith1983}). Reconstruct the separation.

\item Take consistency alone in both paradigms. It is free in PAC (\leanname{erm_consistent_realizable}) yet in Gold can force a learner into infinitely many mind changes. Exhibit a class realizing the Gold side, and explain why the i.i.d.\ assumption removes the cost in PAC.

\item Compare set-drivenness (\leanname{FLT_Proofs.Learner.Properties.IsSetDriven}) with iterativeness (\leanname{FLT_Proofs.Learner.Properties.IsIterative}) in Gold and online learning. A set-driven class can force unbounded memory on any iterative learner (\cite{KotzingPalenta2016,Gold1967}); this separation is open in the kernel, recorded as \cref{op:order-sensitivity}. Move the same pair into PAC and show the distinction collapses because example order is irrelevant.

\item Read properness under cryptographic hardness, where the spent resource is computation time, not information. Properness preserves the finite-VC PAC class yet can cost strictly more computation (\cite{PittValiant1988}). Identify which axis the separation lives on.

\item Add internal randomness, the probabilistic variant (\leanname{FLT_Proofs.Learner.Properties.ProbabilisticLearner}). In PAC randomness derandomizes at no asymptotic cost; in the Gold limit probabilistic identification is a distinct success notion. Contrast the two.

\item Replace a single learner by a team (\leanname{FLT_Proofs.Learner.Properties.TeamLearner}). Unlike consistency, the team relaxation keeps its bite in Gold, where two learners identify classes no single learner can (\cref{thm:team-power}), and survives the move to PAC uncollapsed. Connect this to the team hierarchy of \cref{op:team-hierarchy}.

\item Ask whether ``impose consistency and properness'' is a closure operator under which PAC learnability is invariant and idempotent, so that a constraint is free exactly when it lies below the closure's fixed points. Such a closure would classify, by one Galois-style law, which constraints collapse in which paradigm. No such adjunction is defined; record what a definition would have to supply.
\end{enumerate}

\input{chapters/ch05_pac}
% Chapter 6: Online Learning and the Littlestone Dimension
% mistake_bounded = Profile B-light, regret_bounded = Profile C (Revelation),
% littlestone_characterization = Profile B (Load-Bearing).
% Centerpiece: the SOA algorithm and mistake tree diagram.
% Style: combinatorial and algorithmic, in contrast to Ch5's probabilistic flavor.

\chapter{Online Learning and the Littlestone Dimension}
\label{ch:online}

PAC learning can afford to be patient: nature draws training examples from a fixed distribution, the learner observes and generalizes, and the worst the world can do is an unlucky sample. That patience is a luxury. In online learning the luxury is revoked. There is no distribution to average over, no training phase to hide behind. At each round an adversary presents an instance, the learner commits to a label, and the adversary reveals the truth; every disagreement is a recorded mistake. The question is not ``what does the learner achieve in expectation?'' but ``how many mistakes can a clever adversary force, no matter what the learner does?''

This shift from distribution to adversary, from batch to sequential, from probabilistic to combinatorial, transforms the mathematics from the ground up. PAC learning is characterized by a number, the VC dimension, extracted through probabilistic concentration arguments. Online learning is characterized by a \emph{tree}, the mistake tree, whose depth the adversary can trace root-to-leaf to force exactly that many mistakes. The central result of this chapter, $\mathrm{Opt}(\mathcal{H}) = \Ldim(\mathcal{H})$, is proved by constructing both sides: an algorithm that limits mistakes to $\Ldim(\mathcal{H})$, and an adversary that forces at least $\Ldim(\mathcal{H})$. No tail probability appears anywhere; the entire argument is an explicit strategy meeting an explicit counter-strategy.

% ================================================================
% SECTION 1: THE ONLINE LEARNING GAME
% Profile G (Framework Bridge). Presented as a GAME, not a definition.
% ================================================================

\section{The Online Learning Game}
\label{sec:online-game}

The setup is a protocol between two players, a \emph{learner} and an \emph{adversary}, competing over a domain $X$, a label set $Y = \{0,1\}$, and a hypothesis class $\mathcal{H} \subseteq \{0,1\}^X$. One constraint ties the adversary's hands: \emph{realizability}, the requirement that some $h^* \in \mathcal{H}$ must be consistent with every label the adversary has assigned. Apart from that, the adversary is free to adapt, watching the learner's prior predictions before choosing the next instance $x_t$.

\begin{defn}[Online Learning Protocol]
\label{def:online-protocol}
The online learning game proceeds in rounds $t = 1, 2, 3, \ldots$\,:
\begin{enumerate}[leftmargin=*, itemsep=2pt]
\item The adversary selects an instance $x_t \in X$.
\item The learner observes $x_t$ and predicts a label $\hat{y}_t \in \{0,1\}$.
\item The adversary reveals the true label $y_t \in \{0,1\}$.
\item If $\hat{y}_t \neq y_t$, the learner incurs a \emph{mistake}.
\end{enumerate}
The game continues for as many rounds as the adversary chooses. The adversary must ensure that there exists $h^* \in \mathcal{H}$ with $h^*(x_t) = y_t$ for all $t$ (realizability).
\end{defn}

\begin{figure}[ht]
\centering
\begin{tikzpicture}[
    >={Stealth[length=2mm]},
    advbox/.style={draw, thick, rounded corners, fill=thmred!12, font=\small\bfseries,
      minimum width=2.5cm, minimum height=0.65cm},
    lrnbox/.style={draw, thick, rounded corners, fill=defblue!12, font=\small\bfseries,
      minimum width=2.5cm, minimum height=0.65cm},
    life/.style={semithick, densely dashed, gray!55},
    rstep/.style={-{Stealth[length=2mm]}, semithick},
    msg/.style={font=\scriptsize, fill=white, inner sep=1.5pt},
    val/.style={draw, semithick, rounded corners, align=center, font=\scriptsize,
      inner sep=3pt, minimum height=0.85cm, minimum width=1.9cm},
    vedge/.style={-{Stealth[length=2mm]}, sepgreen!55!black, semithick},
    sedge/.style={-{Stealth[length=2mm]}, edgeorange!80!black, semithick, densely dashed},
    hdl/.style={font=\tiny, text=black!70, align=center}
]
% --- protocol: two lifelines, time flows down ---
\node[advbox] (adv) at (-3.7,2.4) {Adversary};
\node[lrnbox] (lrn) at (3.7,2.4) {Learner};
\draw[life] (adv) -- (-3.7,-0.55);
\draw[life] (lrn) -- (3.7,-0.55);
\draw[rstep,thmred!70!black] (-3.7,1.45) -- node[msg]{1.~presents $x_t$} (3.7,1.45);
\draw[rstep,defblue!70!black] (3.7,0.7) -- node[msg]{2.~predicts $\hat y_t$} (-3.7,0.7);
\draw[rstep,thmred!70!black] (-3.7,-0.05) -- node[msg]{3.~reveals $y_t$; mistake if $\hat y_t\neq y_t$} (3.7,-0.05);
\node[font=\scriptsize\itshape, text=black!65] at (0,-0.95) {realizability: $\exists\,h^{\ast}\!\in\mathcal{H},\ h^{\ast}(x_t)=y_t$ for all $t$};
\node[hdl, text=gray!60, anchor=west] at (-6.4,0.7) {round\\$t=1,2,\dots$};
% --- link: why the value is combinatorial ---
\draw[gray!55, ->, densely dashed] (0,-1.3) -- node[hdl, right, pos=0.5]
  {optimal adversary $=$ shattered mistake tree (\cref{fig:mistake-tree})} (0,-2.05);
% --- value chain (what the game is worth) ---
\node[val, fill=sepgreen!10] (opt)  at (-5.2,-2.7) {game value\\$\mathrm{Opt}(\mathcal{H})$};
\node[val, fill=sepgreen!10] (ldim) at (-1.7,-2.7) {$\Ldim(\mathcal{H})$};
\node[val, fill=defblue!8]  (onl)  at (1.8,-2.7) {online\\learnable};
\node[val, fill=white, densely dashed] (pac) at (5.3,-2.7) {PAC\\learnable};
\draw[vedge] (opt) -- node[hdl, above]{$=$} (ldim);
\draw[vedge] (ldim) -- node[hdl, above]{finite $\Leftrightarrow$} (onl);
\draw[sedge] (onl) -- node[hdl, above, text=edgeorange!80!black]{$\Rightarrow$ strict} (pac);
\node[hdl] at (-3.45,-3.45) {\leanname{optimal_mistake_bound_eq_ldim}};
\node[hdl] at (0.05,-3.45) {\leanname{littlestone_characterization}};
\node[hdl, text=edgeorange!75!black] at (3.55,-3.45) {\leanname{online_strictly_stronger_pac}};
\end{tikzpicture}
\caption[The online learning game, whose value is the Littlestone dimension]{The online learning game: round by round the adversary presents an instance, the learner predicts, and the adversary reveals the truth, recording a mistake on every disagreement, with realizability as the only constraint on the adversary. The lower chain shows what the game is worth: the optimal mistake bound equals the Littlestone dimension (\leanname{optimal_mistake_bound_eq_ldim}), the class is online learnable if and only if that dimension is finite, and the one-way implication to PAC learnability is strict, with thresholds on $\R$ (VC dimension one, Littlestone dimension infinite) as the separating witness.}
\label{fig:online-protocol}
\end{figure}

Three features separate this game from the PAC framework of \Cref{ch:pac}, and each one tightens the screws:

\begin{enumerate}[leftmargin=*, itemsep=3pt]
\item \textbf{No distribution.} PAC learning grades the learner against a fixed, unknown distribution $D$; averaging over randomness is baked in. Here there is no distribution. The adversary's choices need not come from any probability measure.

\item \textbf{The adversary adapts.} The adversary watches the learner's predictions and targets the next instance accordingly. This is strictly harder than random sampling: an adaptive adversary can simulate any distribution and then do worse, exploiting whatever the learner just revealed.

\item \textbf{Worst-case over every sequence.} The mistake bound must hold for every sequence the adversary might produce, subject only to realizability. Averaging is off the table.
\end{enumerate}

\begin{defn}[Mistake Bound]
\label{def:mistake-bound}
A learner $A$ is \emph{mistake-bounded} with bound $M$ on $\mathcal{H}$ if, for every adversary strategy consistent with some $h^* \in \mathcal{H}$, the total number of mistakes made by $A$ is at most $M$. The \emph{optimal mistake bound} of $\mathcal{H}$ is $\mathrm{Opt}(\mathcal{H}) = \min_A \max_{\mathrm{adversary}} \#\{\text{mistakes of } A\}$.
\end{defn}

The central question is: what pins $\mathrm{Opt}(\mathcal{H})$? The answer is a combinatorial object, the Littlestone dimension, and it is a tree, not a set.

% ================================================================
% SECTION 2: MISTAKE TREES AND THE LITTLESTONE DIMENSION
% The key combinatorial object. TikZ diagram essential.
% ================================================================

\section{Mistake Trees and the Littlestone Dimension}
\label{sec:mistake-trees}

The VC dimension asks: how large a \emph{set} can $\mathcal{H}$ shatter? The Littlestone dimension asks a harder question: how deep a \emph{tree} can $\mathcal{H}$ shatter? The difference is adaptivity. A shattered set fixes its points before asking whether the class can realize all labelings. A shattered tree lets the points at each depth depend on the labels above, which is exactly the power an adaptive adversary exercises. Replacing sets with trees is the move that captures the right object.

\begin{defn}[Mistake Tree]
\label{def:mistake-tree}
A \emph{mistake tree} for $\mathcal{H}$ over $X$ is a complete binary tree $T$ in which:
\begin{itemize}[leftmargin=*, itemsep=2pt]
\item Each internal node $v$ is labeled with an instance $x_v \in X$.
\item The left child of $v$ corresponds to the label $y = 0$ and the right child to $y = 1$.
\item For every root-to-leaf path $(v_1, y_1), (v_2, y_2), \ldots, (v_d, y_d)$, there exists a hypothesis $h \in \mathcal{H}$ consistent with all labels along the path: $h(x_{v_i}) = y_i$ for all $i$.
\end{itemize}
The \emph{depth} of the tree is the number of internal nodes on any root-to-leaf path (all paths have the same length since the tree is complete).\formal{FLT_Proofs.Complexity.GameInfra.LTree.isShattered}
\end{defn}

A mistake tree is precisely an adversary strategy written down as a tree. Starting at the root, the adversary presents $x_{v_1}$. Whatever the learner predicts, the adversary labels $x_{v_1}$ to contradict that prediction and descends to the corresponding child. Because every root-to-leaf path is consistent with some $h \in \mathcal{H}$, realizability holds throughout. A tree shattered at depth $d$ therefore witnesses an adversary strategy that forces \emph{any} learner to make at least $d$ mistakes.

\begin{defn}[Littlestone Dimension {\cite{Littlestone1988}}]
\label{def:ldim}
The \emph{Littlestone dimension} of a hypothesis class $\mathcal{H}$, denoted $\Ldim(\mathcal{H})$, is the maximum depth of a complete mistake tree for $\mathcal{H}$. If arbitrarily deep mistake trees exist, $\Ldim(\mathcal{H}) = \infty$.\formal{FLT_Proofs.Complexity.GameInfra.LittlestoneDim}
\end{defn}

\begin{figure}[ht]
\centering
\begin{tikzpicture}[
    level distance=1.6cm,
    sibling distance=1.2cm,
    every node/.style={font=\small},
    treenode/.style={draw, circle, thick, minimum size=0.7cm, fill=onlineblue!10, font=\small\bfseries},
    leafnode/.style={draw, rectangle, rounded corners, thick, minimum width=0.6cm, minimum height=0.45cm, fill=sepgreen!15, font=\scriptsize},
    edgelabel/.style={font=\scriptsize, midway},
    advedge/.style={draw=thmred, line width=1.3pt},
    level 1/.style={sibling distance=5.6cm},
    level 2/.style={sibling distance=2.8cm},
    level 3/.style={sibling distance=1.4cm},
]
\node[treenode] (root) {$4$}
    child { node[treenode] (n2) {$2$}
        child { node[treenode] (n1) {$1$}
            child { node[leafnode] (h0) {$h_{\leq 0}$} edge from parent node[edgelabel, left] {$0$} }
            child { node[leafnode] {$h_{\leq 1}$} edge from parent node[edgelabel, right] {$1$} }
            edge from parent node[edgelabel, left] {$0$}
        }
        child { node[treenode] {$3$}
            child { node[leafnode] {$h_{\leq 2}$} edge from parent node[edgelabel, left] {$0$} }
            child { node[leafnode] {$h_{\leq 3}$} edge from parent node[edgelabel, right] {$1$} }
            edge from parent node[edgelabel, right] {$1$}
        }
        edge from parent node[edgelabel, left] {$0$}
    }
    child { node[treenode] {$6$}
        child { node[treenode] {$5$}
            child { node[leafnode] {$h_{\leq 4}$} edge from parent node[edgelabel, left] {$0$} }
            child { node[leafnode] {$h_{\leq 5}$} edge from parent node[edgelabel, right] {$1$} }
            edge from parent node[edgelabel, left] {$0$}
        }
        child { node[treenode] {$7$}
            child { node[leafnode] {$h_{\leq 6}$} edge from parent node[edgelabel, left] {$0$} }
            child { node[leafnode] {$h_{\leq 7}$} edge from parent node[edgelabel, right] {$1$} }
            edge from parent node[edgelabel, right] {$1$}
        }
        edge from parent node[edgelabel, right] {$1$}
    };
% --- highlighted adversary run: root -> 2 -> 1 -> h_{<=0}, forces depth-many (3) mistakes ---
\draw[advedge] (root) -- (n2);
\draw[advedge] (n2) -- (n1);
\draw[advedge] (n1) -- (h0);
\node[font=\scriptsize\itshape, text=thmred, anchor=north, align=center] at ($(h0.south)+(0,-0.12)$)
  {adversary run:\\$3$ mistakes forced};
% --- the general meaning + verified handles (A4: verified >= vs borrowed =) ---
\node[font=\scriptsize, align=left, anchor=north west] at (3.7,1.05) {%
  \textbf{A shattered mistake tree.}\\[2pt]
  \textit{$X=\{1,\dots,8\}$; thresholds $h_{\leq\theta}(x)=\ind[x\leq\theta]$.}\\[1pt]
  \textit{internal node: an instance; leaf: a consistent $h$.}\\[1pt]
  \textit{every path realizable $\Rightarrow$ tree shattered.}\\[3pt]
  depth $=3=\Ldim(\mathcal{H})$\quad\leanname{LittlestoneDim}\\[1pt]
  a depth-$d$ shattered tree forces $d$ mistakes:\\
  \quad\leanname{adversary_core}\\[2pt]
  here $\Ldim\geq 3$\quad\leanname{thresholdTree_shattered}\\[1pt]
  \textcolor{notegray}{$=\lceil\log_2 8\rceil$\quad(\cite{Littlestone1988})}
};
\end{tikzpicture}
\caption[A shattered mistake tree of depth three witnessing $\Ldim=3$ for thresholds on eight points]{A complete mistake tree of depth $3$ for thresholds on $\{1,\ldots,8\}$, witnessing $\Ldim(\mathcal{H}) \geq 3$ (\leanname{thresholdTree_shattered}). Each internal node is an instance; each edge is a label; every root-to-leaf path is realized by the threshold named at its leaf. The highlighted run shows the adversary executing binary search: at each node it assigns the label opposite to the learner's prediction and descends, forcing $3$ mistakes in $3$ rounds while maintaining realizability throughout (\leanname{adversary_core}). The Littlestone dimension is the maximum depth of any such shattered tree (\leanname{LittlestoneDim}); the fact that this tree has depth $\lceil\log_2 8\rceil = 3$ confirms the threshold bound.}
\label{fig:mistake-tree}
\end{figure}

\begin{example}[Littlestone dimension of fundamental classes]
\label{ex:ldim-examples}
\begin{enumerate}[label=(\alph*), leftmargin=*, itemsep=4pt]
\item \textbf{Finite classes.} If $|\mathcal{H}| < \infty$, then $\Ldim(\mathcal{H}) \leq \lfloor \log_2 |\mathcal{H}| \rfloor$, since a complete binary tree of depth $d$ has $2^d$ leaves and each leaf requires a distinct consistent hypothesis. This bound is tight for classes that are ``richly structured enough'' (e.g., all functions on a domain of size $\log_2 |\mathcal{H}|$).

\item \textbf{Thresholds on $\{1, \ldots, n\}$.} The class $\{h_{\leq \theta} : \theta \in \{0, 1, \ldots, n\}\}$ has $\Ldim = \lceil \log_2(n+1) \rceil$. The mistake tree in \Cref{fig:mistake-tree} illustrates the case $n = 8$. The adversary performs binary search on the threshold.

\item \textbf{Thresholds on $\R$.} Now $\Ldim = \infty$. The key point: the domain is dense, so the adversary can always find a point between any two previous thresholds. At each round, the adversary presents a point that bisects the current interval of uncertainty, forcing a mistake no matter what the learner predicts. This produces mistake trees of unbounded depth.

\item \textbf{Intervals on $\R$.} The class $\{x \mapsto \ind[a \leq x \leq b] : a, b \in \R\}$ also has $\Ldim = \infty$. The adversary can adaptively narrow down both endpoints.

\item \textbf{Linear classifiers in $\R^d$.} The class of halfspaces has $\Ldim = d$, matching the VC dimension. This is one of the rare cases where the two dimensions coincide.
\end{enumerate}
\end{example}

\begin{remark}[Sets vs.\ trees]
\label{rmk:sets-vs-trees}
Shattering a \emph{set} $\{x_1, \ldots, x_d\}$ means $\mathcal{H}$ can produce all $2^d$ labelings of these \emph{fixed} points. Shattering a \emph{tree} of depth $d$ means $\mathcal{H}$ can produce a consistent labeling along every root-to-leaf path, but the points themselves may \emph{depend on previous labels}. The tree structure captures adaptivity: the adversary's choice at depth $k$ depends on the learner's responses at depths $1, \ldots, k-1$. A shattered set is a special case of a mistake tree (one in which every node at the same depth is labeled with the same instance), so $\VC(\mathcal{H}) \leq \Ldim(\mathcal{H})$ always holds. The converse fails: thresholds on $\R$ have $\VC = 1$ but $\Ldim = \infty$.
\end{remark}

% ================================================================
% SECTION 3: THE STANDARD OPTIMAL ALGORITHM (SOA)
% The constructive heart of online learning. Read like describing an algorithm.
% ================================================================

\section{The Standard Optimal Algorithm}
\label{sec:soa}

The simplest reasonable strategy is Halving: maintain the version space (all hypotheses consistent with everything seen so far), predict the majority vote. Every mistake eliminates at least half the surviving hypotheses, so Halving makes at most $\lfloor \log_2 |\mathcal{H}| \rfloor$ mistakes when $\mathcal{H}$ is finite.

For infinite classes, cardinality collapses and Halving offers nothing. The Standard Optimal Algorithm (SOA), due to Littlestone~\cite{Littlestone1988}, replaces ``majority vote'' with a sharper criterion: predict the label whose sub-version-space carries the larger \emph{Littlestone dimension}. The switch from counting hypotheses to measuring mistake-tree depth is what makes the algorithm optimal for every hypothesis class, finite or infinite.

\begin{defn}[Standard Optimal Algorithm (\SOA)]
\label{def:soa}
Given a hypothesis class $\mathcal{H}$, the $\SOA$ maintains a version space $V_t \subseteq \mathcal{H}$, initially $V_1 = \mathcal{H}$. At round $t$, upon receiving instance $x_t$:
\begin{enumerate}[leftmargin=*, itemsep=2pt]
\item Partition $V_t$ into $V_t^0 = \{h \in V_t : h(x_t) = 0\}$ and $V_t^1 = \{h \in V_t : h(x_t) = 1\}$.
\item Predict $\hat{y}_t = \arg\max_{b \in \{0,1\}} \Ldim(V_t^b)$.
\item After observing the true label $y_t$, update $V_{t+1} = V_t^{y_t}$.
\end{enumerate}
Ties in step 2 are broken arbitrarily.\formal{FLT_Proofs.Complexity.GameInfra.SOA}
\end{defn}

The SOA reasons about future adversarial damage. At each round it asks: if this prediction is wrong, which sub-version-space does the adversary inherit, and how deeply can they search from there? By predicting the label whose sub-version-space has the larger Littlestone dimension, the algorithm guarantees that every mistake lands it in the smaller sub-version-space, reducing the Littlestone dimension by at least one. The game ends when the dimension reaches zero.

\begin{thm}[Upper Bound: $\SOA$ achieves $\Ldim(\mathcal{H})$ mistakes {\cite{Littlestone1988}}]
\label{thm:soa-upper}
For any hypothesis class $\mathcal{H}$ with $\Ldim(\mathcal{H}) = d < \infty$, the $\SOA$ makes at most $d$ mistakes against any adversary.\formal{FLT_Proofs.Theorem.Online.soa_mistakes_bounded}
\end{thm}

\begin{proof}
We show that the Littlestone dimension of the version space drops by at least one at each mistake. Define $d_t = \Ldim(V_t)$. We claim:
\begin{equation}
\label{eq:ldim-drop}
\text{If the $\SOA$ makes a mistake at round } t, \text{ then } d_{t+1} \leq d_t - 1.
\end{equation}%
\formal{FLT_Proofs.Theorem.Online.ldim_strict_decrease_on_mistake}

\emph{Proof of the claim.} Suppose the $\SOA$ makes a mistake at round $t$. This means $\hat{y}_t \neq y_t$. By the algorithm's rule, $\hat{y}_t$ was chosen so that $\Ldim(V_t^{\hat{y}_t}) \geq \Ldim(V_t^{y_t})$, i.e., the $\SOA$ predicted the side with the \emph{larger} (or equal) Littlestone dimension. After the mistake, the version space updates to $V_{t+1} = V_t^{y_t}$, the side with the \emph{smaller} (or equal) Littlestone dimension.

Now we must show that $\Ldim(V_t^{y_t}) \leq d_t - 1$. Suppose for contradiction that $\Ldim(V_t^0) \geq d_t$ and $\Ldim(V_t^1) \geq d_t$. Then there exists a complete mistake tree $T_0$ of depth $d_t$ for $V_t^0$ and a complete mistake tree $T_1$ of depth $d_t$ for $V_t^1$. We can construct a complete mistake tree of depth $d_t + 1$ for $V_t$: the root is labeled $x_t$, its left subtree (corresponding to label $0$) is $T_0$, and its right subtree (corresponding to label $1$) is $T_1$. Every root-to-leaf path through the left subtree is consistent with some $h \in V_t^0 \subseteq V_t$ (and $h(x_t) = 0$), and similarly for the right subtree. This yields $\Ldim(V_t) \geq d_t + 1$, contradicting $\Ldim(V_t) = d_t$.

Therefore $\min\{\Ldim(V_t^0), \Ldim(V_t^1)\} \leq d_t - 1$. Since $\hat{y}_t$ selects the side with the larger dimension, $V_{t+1} = V_t^{y_t}$ is the side with the smaller dimension, so $d_{t+1} = \Ldim(V_{t+1}) \leq d_t - 1$.

\emph{Completing the proof.} We have $d_1 = \Ldim(\mathcal{H}) = d$. Each mistake decreases $d_t$ by at least $1$. Since the Littlestone dimension is a non-negative integer (the version space always contains the target $h^*$, so $V_t \neq \emptyset$ and $\Ldim(V_t) \geq 0$), the total number of mistakes is at most $d$.
\end{proof}

\begin{remark}[SOA vs.\ Halving]
\label{rmk:soa-halving}
For finite $\mathcal{H}$, the Halving algorithm predicts by majority vote: $\hat{y}_t = \arg\max_b |V_t^b|$. Each mistake halves the version space, giving at most $\lfloor \log_2 |\mathcal{H}| \rfloor$ mistakes. The $\SOA$ replaces cardinality $|V_t^b|$ with $\Ldim(V_t^b)$. For finite classes, this distinction is often unimportant (both give $O(\log |\mathcal{H}|)$ mistakes), but for infinite classes, cardinality is meaningless while the Littlestone dimension is finite and well-behaved.
\end{remark}

\begin{remark}[Computability of the $\SOA$]
\label{rmk:soa-computability}
The $\SOA$ is an \emph{information-theoretically} optimal algorithm, but it is not necessarily \emph{computationally} efficient. Computing $\Ldim(V_t^b)$ at each round may be undecidable for some hypothesis classes. The $\SOA$ is thus an existence proof: it shows that $d$ mistakes suffice, even if finding the optimal prediction at each step is computationally hard. Efficient online algorithms for specific classes (Perceptron for halfspaces, Winnow for disjunctions) achieve near-optimal mistake bounds through class-specific structure.
\end{remark}

% ================================================================
% SECTION 4: THE LOWER BOUND - ADVERSARY STRATEGY
% The adversary plays the mistake tree from root to leaf.
% ================================================================

\section{The Lower Bound: The Adversary Strategy}
\label{sec:lower-bound}

The SOA limits the learner's mistakes to $\Ldim(\mathcal{H})$. That is the ceiling. The lower bound closes the gap: no learner can stay below that ceiling, because for any learner an adversary exists that forces at least $\Ldim(\mathcal{H})$ mistakes. Together the two bounds pin $\mathrm{Opt}(\mathcal{H})$ exactly.

\begin{thm}[Lower Bound: Any learner makes at least $\Ldim(\mathcal{H})$ mistakes {\cite{Littlestone1988}}]
\label{thm:lower-bound}
For any hypothesis class $\mathcal{H}$ with $\Ldim(\mathcal{H}) = d$ and any (possibly randomized) learner $A$, there exists an adversary strategy that forces $A$ to make at least $d$ mistakes.\formal{FLT_Proofs.Theorem.Online.adversary_lower_bound}
\end{thm}

\begin{proof}
Let $T$ be a complete mistake tree for $\mathcal{H}$ of depth $d$. The adversary plays $T$ as follows.

The adversary maintains a pointer to a node of $T$, starting at the root. At round $t$, the adversary presents the instance $x_{v_t}$ labeling the current node $v_t$. The learner predicts $\hat{y}_t$. The adversary then sets $y_t = 1 - \hat{y}_t$ (the opposite of the learner's prediction) and descends to the child of $v_t$ corresponding to label $y_t$.

This strategy has two properties:
\begin{enumerate}[leftmargin=*, itemsep=2pt]
\item \textbf{Every round is a mistake.} By construction, $y_t = 1 - \hat{y}_t \neq \hat{y}_t$.
\item \textbf{Realizability is maintained.} After $d$ rounds, the adversary has descended from the root to a leaf of $T$, tracing a path $(v_1, y_1), (v_2, y_2), \ldots, (v_d, y_d)$. By the definition of a mistake tree, there exists $h^* \in \mathcal{H}$ consistent with all labels along this path: $h^*(x_{v_i}) = y_i$ for all $i \in \{1, \ldots, d\}$. For any subsequent rounds $t > d$, the adversary can label consistently with this $h^*$.
\end{enumerate}

The adversary forces exactly $d$ mistakes in the first $d$ rounds, so $A$ makes at least $d$ mistakes.

For randomized learners, the argument extends by the minimax theorem (or directly): for any distribution over predictions $\hat{y}_t$, the adversary's strategy of playing the opposite forces an expected mistake at every round. Alternatively, one can apply Yao's minimax principle: the worst-case deterministic adversary against the best randomized learner equals the best deterministic learner against the worst-case adversary, and we have shown the latter is at least $d$.
\end{proof}

Combining \Cref{thm:soa-upper,thm:lower-bound}, we obtain the fundamental characterization.

\begin{thm}[Littlestone's Characterization {\cite{Littlestone1988}}]
\label{thm:littlestone-characterization}
For any hypothesis class $\mathcal{H}$:
\[
\mathrm{Opt}(\mathcal{H}) = \Ldim(\mathcal{H}).
\]
That is, the optimal mistake bound of $\mathcal{H}$ in the online learning game is exactly the Littlestone dimension of $\mathcal{H}$. In particular, $\mathcal{H}$ admits a finite mistake bound if and only if $\Ldim(\mathcal{H}) < \infty$.\formal{FLT_Proofs.Theorem.Online.optimal_mistake_bound_eq_ldim}
\end{thm}

\begin{historical}
Littlestone proved this characterization in 1988~\cite{Littlestone1988}, just four years after Valiant's PAC model. The result established online learning as a mathematically independent paradigm: the combinatorial quantity governing online learnability is genuinely different from the one governing PAC learnability. The $\SOA$ is sometimes called the ``Standard Optimal Algorithm'' precisely because it achieves the information-theoretic optimum; the name is due to the online learning community's convention.
\end{historical}

% ================================================================
% SECTION 5: REGRET BOUNDS AND MULTIPLICATIVE WEIGHTS
% Profile C (Revelation). Regret is surprising because it is relative.
% ================================================================

\section{Regret Bounds and the Multiplicative Weights Framework}
\label{sec:regret}

The mistake-bound framework rests on realizability: some $h^* \in \mathcal{H}$ must be consistent with every label. That is a substantial assumption. When it fails, the optimal mistake count is infinite and the framework offers no useful bound.

The \emph{regret} framework drops realizability and changes the benchmark. Rather than asking for an absolute mistake count, it grades the learner against the \emph{best fixed hypothesis in hindsight}: whoever single $h \in \mathcal{H}$ would have made the fewest errors on the realized sequence, even if that $h$ errs constantly. The learner's job is to stay close to that benchmark, not to be right.

\begin{defn}[Regret]
\label{def:regret}
Given a sequence $(x_1, y_1), \ldots, (x_T, y_T)$ and the learner's predictions $\hat{y}_1, \ldots, \hat{y}_T$, the \emph{regret} of the learner relative to $\mathcal{H}$ is
\[
\mathrm{Regret}_T = \sum_{t=1}^T \ind[\hat{y}_t \neq y_t] - \min_{h \in \mathcal{H}} \sum_{t=1}^T \ind[h(x_t) \neq y_t].
\]
\end{defn}

The definition holds a conceptual surprise. Regret can be small even when the learner makes many mistakes, because the benchmark also makes many mistakes. If the adversary is unrealizable and every hypothesis in $\mathcal{H}$ errs frequently, then keeping pace with the best of them is a genuine achievement, even with a large absolute error count. What the regret framework measures is the \emph{cost of not knowing which hypothesis is best}, not the absolute cost of prediction.

\begin{remark}[Mistake bound vs.\ regret bound]
\label{rmk:mistake-vs-regret}
In the realizable case, the best hypothesis $h^*$ makes zero mistakes, so $\mathrm{Regret}_T$ equals the total mistake count. The regret framework is therefore a strict generalization of the mistake-bound framework. The regret formulation becomes essential when $\mathcal{H}$ is infinite or when realizability fails: precisely the settings where the Littlestone dimension may be infinite and the mistake-bound framework provides no guarantee.
\end{remark}

The Weighted Majority algorithm of Littlestone and Warmuth~\cite{Littlestone1988} achieves the following regret bound for finite $\mathcal{H}$.

\begin{thm}[Weighted Majority / Hedge]
\label{thm:weighted-majority}
For a finite hypothesis class $\mathcal{H}$ with $|\mathcal{H}| = N$, the Weighted Majority algorithm achieves, for any sequence of $T$ rounds:
\[
\mathrm{Regret}_T \leq 2\sqrt{T \ln N}.
\]
In particular, the per-round regret $\mathrm{Regret}_T / T \to 0$ as $T \to \infty$.
\end{thm}

The algorithm maintains a weight $w_t(h)$ for each $h \in \mathcal{H}$, initially $w_1(h) = 1$. At each round: predict the weighted majority vote; after observing $y_t$, multiply the weight of each hypothesis that erred by a factor $(1 - \eta)$ for a learning rate $\eta \in (0, 1)$. Setting $\eta = \sqrt{\ln N / T}$ (or using the doubling trick when $T$ is unknown) yields the bound above.

This approach generalizes to the \emph{multiplicative weights} framework (also called Hedge), which extends beyond binary prediction to decision-making over finitely many ``experts.'' The framework is one of the most widely applicable algorithmic paradigms in theoretical computer science, with applications to zero-sum games, boosting, and linear programming~\cite{ShalevShwartzBenDavid2014}.

\begin{remark}[Connection to Littlestone dimension]
\label{rmk:regret-ldim}
Ben-David, P\'al, and Shalev-Shwartz~\cite{BenDavidPalShalevShwartz2009} showed that for infinite hypothesis classes, the Littlestone dimension also governs the optimal regret rate. Specifically, the minimax regret over $T$ rounds is $\Theta(\sqrt{\Ldim(\mathcal{H}) \cdot T})$ in the agnostic (non-realizable) setting. This connects the combinatorial tree structure of the Littlestone dimension to the sequential prediction framework even beyond realizability.
\end{remark}

% ================================================================
% SECTION 6: THE PAC-ONLINE GAP
% VC ≤ Ldim always; thresholds show the gap can be infinite.
% ================================================================

\subsection{Attention as exponential weights}
\label{sec:attention-online}

The multiplicative-weights update of \cref{thm:weighted-majority} is the computational core of the attention mechanism.  Fix scores $s(x) \in \R^k$ assigned to $k$ keys by a parametric router on input $x$.  The soft-attention weight on key $i$ is the softmax $w_i(x) = e^{s_i(x)} / \sum_j e^{s_j(x)}$, which is exactly the exponential-weights (Hedge) distribution with the scores in the role of cumulative gains.

\begin{prop}[Softmax attention is an exponential-weights step]
\label{prop:attention-hedge}
The soft-attention weights $w_i(x) \propto \exp(s_i(x))$ are the Hedge distribution of \cref{thm:weighted-majority} over the $k$ keys.\formal{TLT_Proofs.Routing.MeasurableScoreRouting.softmaxWeight} The weight order matches the score order ($w_i(x) \le w_j(x)$ if and only if $s_i(x) \le s_j(x)$), so the hard, argmax route attends to the maximally-weighted key, the Halving-style greedy specialization of the same update.\formal{TLT_Proofs.Routing.MeasurableScoreRouting.softmaxWeight_le_iff_score}
\end{prop}

\begin{proof}
Both claims unfold the softmax.  The normalization $\sum_j e^{s_j(x)}$ is positive, so $w_i \propto e^{s_i}$ is a probability vector, and it coincides termwise with the exponential-weights distribution that assigns mass $\propto e^{\eta G_i}$ to expert $i$ of cumulative gain $G_i$ at learning rate $\eta$ (the inverse temperature, absorbed into the score scale).  Since $t \mapsto e^t$ is strictly increasing, $w_i \le w_j \Leftrightarrow e^{s_i} \le e^{s_j} \Leftrightarrow s_i \le s_j$, and the $\arg\max$ of the weights is the $\arg\max$ of the scores.
\end{proof}

The correspondence makes the regret framework of \cref{sec:regret} the right setting for a routed transformer: a layer attending over $k$ keys runs one exponential-weights step, and the regret bound $2\sqrt{T \ln k}$ bounds how much worse the routed predictions are than the single best key in hindsight.

\begin{remark}[Route stability]
\label{rmk:route-stability}
A tempered attention router's hard route is stable: away from a score tie, a small perturbation of the scores, or of the parameters and temperature that produce them, leaves the routed key unchanged.\formal{TLT_Proofs.TemperedDesignLaw.Stability.TD7_route_stable_proof} The inputs at which the route flips form the score-tie boundary, of the same measure-zero character that the mistake-tree analysis of \cref{sec:soa} exploits when it splits the version space by a single instance's label.  Stability is what lets the discrete, argmax-routed learner inherit the regret guarantee of its smooth softmax surrogate.
\end{remark}

\section{The PAC--Online Gap}
\label{sec:pac-online-gap}

Both the VC dimension and the Littlestone dimension measure the complexity of a hypothesis class, but each captures a different aspect. The VC dimension asks what the class can do to a fixed set of points; the Littlestone dimension asks what an adaptive adversary can do when the points are chosen sequentially in response to the learner. The two measures agree on some classes and diverge wildly on others.

\begin{prop}[$\VC \leq \Ldim$]
\label{prop:vc-leq-ldim}
For any hypothesis class $\mathcal{H}$, $\VC(\mathcal{H}) \leq \Ldim(\mathcal{H})$.\formal{FLT_Proofs.Complexity.Littlestone.BranchWiseLittlestoneDim_ge_VCDim}
\end{prop}

\begin{proof}
Let $S = \{x_1, \ldots, x_d\}$ be a set shattered by $\mathcal{H}$, so $\VC(\mathcal{H}) \geq d$. We construct a complete mistake tree of depth $d$. The root is labeled $x_1$. Every node at depth $k$ is labeled $x_{k+1}$ (the same instance at every node of a given depth). For any root-to-leaf path with labels $(y_1, \ldots, y_d)$, shattering guarantees that some $h \in \mathcal{H}$ satisfies $h(x_i) = y_i$ for all $i$. Thus this is a valid mistake tree of depth $d$, so $\Ldim(\mathcal{H}) \geq d$.
\end{proof}

\begin{figure}[ht]
\centering
\begin{tikzpicture}[
    level distance=1.4cm,
    every node/.style={font=\small},
    treenode/.style={draw, circle, thick, minimum size=0.65cm, fill=defblue!10, font=\small},
    leafnode/.style={draw, rectangle, rounded corners, thick, minimum width=0.5cm, minimum height=0.4cm, fill=sepgreen!15, font=\scriptsize},
    edgelabel/.style={font=\scriptsize, midway},
    setbox/.style={draw, thick, rounded corners, fill=defblue!6, align=center, font=\scriptsize, inner sep=4pt},
    embed/.style={-{Stealth[length=2.5mm]}, thick, defblue!75!black},
    gap/.style={-{Stealth[length=2.5mm]}, thick, densely dashed, thmred!80!black},
    hdl/.style={font=\tiny, text=black!70, align=center},
    level 1/.style={sibling distance=3.6cm},
    level 2/.style={sibling distance=1.8cm},
]
% --- the non-adaptive tree: same instance x_2 at every depth-1 node ---
\node[treenode] (root) {$x_1$}
    child { node[treenode] {$x_2$}
        child { node[leafnode] {$(0,0)$} edge from parent node[edgelabel, left] {$0$} }
        child { node[leafnode] {$(0,1)$} edge from parent node[edgelabel, right] {$1$} }
        edge from parent node[edgelabel, left] {$0$}
    }
    child { node[treenode] {$x_2$}
        child { node[leafnode] {$(1,0)$} edge from parent node[edgelabel, left] {$0$} }
        child { node[leafnode] {$(1,1)$} edge from parent node[edgelabel, right] {$1$} }
        edge from parent node[edgelabel, right] {$1$}
    };
% --- (A) shattered set source + solid embedding arrow (set -> non-adaptive tree) ---
\node[setbox] (set) at (-4.8,-0.35) {Shattered set\\$S=\{x_1,x_2\}$\\\textsc{vc}: all $2^2$\\labelings realized};
\node[font=\scriptsize\itshape, text=notegray] at (-1.95,0.95) {embed (same instance per depth)};
\node[font=\small] at (-1.95,0.55) {$\VC(\mathcal{H})\le\Ldim(\mathcal{H})$};
\node[hdl] at (-1.95,0.18) {\leanname{BranchWiseLittlestoneDim_ge_VCDim}};
\draw[embed] (set.east) to[out=0,in=168] (root.west);
\node[hdl, text=notegray, align=center, anchor=north] at (-4.8,-1.3) {(branch-wise $=$ path-\\consistent on this\\non-adaptive tree)};
% --- (C) strict-gap dashed arrow + note: adaptivity escapes the image ---
\node[draw, thick, densely dashed, thmred!75!black, rounded corners, fill=thmred!4,
      align=left, font=\scriptsize, inner sep=4pt] (gapnote) at (5.0,-2.3)
  {\textbf{One-way only.}\\an adaptive tree need not\\come from a set: thresholds\\on $\R$ give $\VC=1$, $\Ldim=\infty$.\\[2pt]
   {\tiny\leanname{ldim_threshold_top}}\\{\tiny\leanname{online_strictly_stronger_pac}}};
\draw[gap] (2.0,-1.7) to[out=-12,in=165] (gapnote.west);
% --- (D) one-liner under the tree ---
\node[font=\scriptsize\itshape, text=notegray, align=center, anchor=north, text width=10.5cm] at (-0.3,-3.5)
  {Non-adaptive: the same instance at each depth, so the tree \emph{is} the set; dropping that constraint is adaptivity, which is what lets $\Ldim$ exceed $\VC$.};
\end{tikzpicture}
\caption[Shattered set embeds as a non-adaptive tree, proving $\VC\le\Ldim$; thresholds on $\R$ show the gap is infinite]{The argument for $\VC(\mathcal{H})\le\Ldim(\mathcal{H})$ made visible: a shattered set $\{x_1,x_2\}$ becomes a non-adaptive mistake tree of the same depth by placing the same instance at every node of each level (\leanname{BranchWiseLittlestoneDim_ge_VCDim}). Every root-to-leaf path is realized because shattering supplies a consistent hypothesis for each labeling. The embedding is one-way: an adaptive tree lets the instances vary with the labels, escaping the image of any shattered set. Thresholds on $\R$ demonstrate that the gap is unbounded, with $\VC=1$ and $\Ldim=\infty$ (\leanname{ldim_threshold_top}, \leanname{online_strictly_stronger_pac}). The slack $\Ldim - \VC$ is the precise measure of what adaptivity adds.}
\label{fig:vc-leq-ldim}
\end{figure}

The inequality $\VC \leq \Ldim$ is tight for some classes: for halfspaces in $\R^d$, both dimensions equal $d$. But the gap can be infinite, and a single class witnesses it.

\begin{prop}[The PAC--Online separation: thresholds on $\R$]
\label{prop:pac-online-sep}
Let $\mathcal{H}_{\mathrm{thr}} = \{x \mapsto \ind[x \leq \theta] : \theta \in \R\}$. Then $\VC(\mathcal{H}_{\mathrm{thr}}) = 1$ and $\Ldim(\mathcal{H}_{\mathrm{thr}}) = \infty$.\formal{FLT_Proofs.Theorem.Separation.pac_not_implies_online}
\end{prop}

\begin{proof}
$\VC = 1$: Any single point $x$ is shattered (take thresholds $\theta < x$ and $\theta > x$). No two points $x_1 < x_2$ are shattered, because $h(x_1) = 0, h(x_2) = 1$ requires $x_2 \leq \theta < x_1$, which contradicts $x_1 < x_2$.

$\Ldim = \infty$: We construct mistake trees of arbitrary depth. For any $d$, consider the domain restricted to $\{1, 2, \ldots, 2^d\}$. The adversary can perform binary search: at depth $k$, present the midpoint of the current interval. Whatever the learner predicts, the adversary labels to force a mistake (and narrows the interval by half). After $d$ rounds, the adversary has traced a path from root to leaf, consistent with the threshold at the boundary of the final interval. Since $d$ was arbitrary, $\Ldim = \infty$.
\end{proof}

This separation is the most important non-implication in the concept graph. It witnesses the edge $\nde{pac\_learning} \xrightarrow{\rel{does\_not\_imply}} \nde{online\_learning}$: a class can be PAC learnable (finite VC dimension) yet not online learnable (infinite Littlestone dimension). The converse implication does hold: $\Ldim < \infty$ implies $\VC < \infty$ (by $\VC \leq \Ldim$), so online learnability implies PAC learnability.\formal{FLT_Proofs.Theorem.Separation.online_imp_pac}

\begin{separation}
\textbf{PAC $\not\Rightarrow$ Online.} Witness: $\mathcal{H}_{\mathrm{thr}}$ on $\R$. In the online game, the adversary uses the order structure of $\R$ to perform binary search on the threshold, forcing a mistake at every round. A random sample from a fixed distribution cannot exploit this structure; with high probability, the sample locates the threshold to within $\varepsilon$. The gap between PAC and online learnability reflects the difference between a world that is fixed and a world that watches you.
\end{separation}

\begin{remark}[Why the gap arises]
\label{rmk:gap-intuition}
In PAC learning the distribution $D$ is fixed before the game begins; the learner faces a static environment. In online learning the adversary adapts to the learner's behavior, and adaptivity is the source of the gap. A threshold on $\R$ is easy to learn from random data: one sample locates it approximately. Against an adaptive adversary the same class is impossible to learn: the adversary always presents a point that bisects the remaining uncertainty, and the bisection never ends. The Littlestone dimension measures the depth of that adaptive search, and it can exceed the VC dimension by an unbounded amount. A full treatment of this and related separations appears in \Cref{ch:separations}.
\end{remark}

\section{What This Chapter Established}
\label{sec:ch6-summary}

The sequential game between learner and adversary pins its value to a single combinatorial quantity:

\begin{enumerate}[leftmargin=*, itemsep=3pt]
\item \textbf{The Littlestone characterization.} The optimal mistake bound for online learning equals the Littlestone dimension $\Ldim(\mathcal{H})$, the maximum depth of a complete mistake tree. This is a \emph{combinatorial} characterization, in contrast to the VC dimension's \emph{set-combinatorial} characterization of PAC learning.

\item \textbf{The SOA algorithm.} The upper bound is constructive: the Standard Optimal Algorithm achieves $\Ldim(\mathcal{H})$ mistakes by predicting the label whose version space has the larger Littlestone dimension. Each mistake provably reduces the Littlestone dimension of the version space by at least one.

\item \textbf{The adversary strategy.} The lower bound is also constructive: the adversary traverses a complete mistake tree from root to leaf, labeling each instance to contradict the learner's prediction.

\item \textbf{The PAC--Online gap.} $\VC(\mathcal{H}) \leq \Ldim(\mathcal{H})$ always, but the gap can be infinite (thresholds: $\VC = 1$, $\Ldim = \infty$). Online learnability implies PAC learnability, but not conversely.

\item \textbf{The regret framework.} Regret bounds, measuring performance relative to the best fixed hypothesis, extend online learning beyond the realizable setting. The Littlestone dimension governs optimal regret rates even in the agnostic case.
\end{enumerate}

The combinatorial and algorithmic character of these results is structural. Online learning theory runs on trees and explicit strategies; PAC learning theory runs on sets and tail probabilities. Both measure the complexity of the same object, a hypothesis class $\mathcal{H}$, but the instruments are genuinely different: the adversary adapts to the learner's behavior round by round, while a fixed distribution cannot.

\section{Open Problems}
\label{sec:ch6-open}

Each problem is anchored to a result above and tagged: a \emph{known unknown} (\textsc{ku}) is an articulated open question; an \emph{unknown known} (\textsc{uk}) is a result in the literature not yet absorbed into a clean statement.

\begin{openprob}[The Littlestone dimension of attention, \textsc{uk}]
\label{op:ldim-attention}
\Cref{prop:attention-hedge} identifies soft attention with exponential weights.  Determine the Littlestone dimension of a finite-head attention class as a function of the architecture and the inverse temperature: is it finite (so the class is online learnable, \cref{thm:littlestone-characterization}), and does it diverge as the temperature grows and the router approaches a hard threshold, the regime in which thresholds on $\R$ already give $\Ldim = \infty$ (\cref{prop:pac-online-sep})?
\end{openprob}

\begin{openprob}[Regret of online attention routing, \textsc{uk}]
\label{op:attention-regret}
Via \cref{prop:attention-hedge}, a layer routing over $k$ keys runs one Hedge step.  Prove the $O(\sqrt{T \ln k})$ regret bound for online attention routing with the inverse temperature as the learning rate, and determine the optimal temperature schedule, connecting it to the route-stability margin (\cref{rmk:route-stability}, \leanname{TLT_Proofs.TemperedDesignLaw.Stability.TD7_route_stable_proof}).
\end{openprob}

\begin{openprob}[Computable optimal online learners, \textsc{ku}]
\label{op:soa-computable}
\Cref{rmk:soa-computability} notes the $\SOA$ may be uncomputable because $\Ldim(V_t^b)$ can be undecidable.  Characterize the classes for which the $\SOA$ is computable in polynomial time per round, and decide whether every class of finite Littlestone dimension admits \emph{some} efficient optimal online learner, or whether an information--computation gap separates online learnability from efficient online learnability as it does in the PAC setting (\Cref{ch:pac}).
\end{openprob}

\begin{openprob}[The optimal regret constant, \textsc{ku}]
\label{op:regret-constant}
Ben-David--P\'al--Shalev-Shwartz~\cite{BenDavidPalShalevShwartz2009} give minimax regret $\Theta(\sqrt{\Ldim(\mathcal{H}) \cdot T})$ in the agnostic setting (\cref{rmk:regret-ldim}).  Determine the exact leading constant, and the finite-$T$ correction, matching upper and lower bounds in terms of the Littlestone dimension alone.
\end{openprob}

\begin{openprob}[The multiclass online characterization, \textsc{uk}]
\label{op:online-multiclass}
The mistake tree and $\SOA$ are defined for binary labels.  Determine the multiclass online characterization: which tree dimension (a multiclass Littlestone dimension over $|Y| > 2$) governs the optimal mistake bound, and whether the clean $\mathrm{Opt} = \Ldim$ identity of \cref{thm:littlestone-characterization} survives, paralleling the multiclass PAC question of \Cref{ch:extensions}.
\end{openprob}

\begin{openprob}[Online learnability and privacy, \textsc{ku}]
\label{op:private-online}
Alon--Livni--Malliaris--Moran established that a class is online learnable if and only if it is differentially-privately PAC learnable, with finite Littlestone dimension on both sides.  Formalize this equivalence on top of the verified Littlestone characterization (\leanname{FLT_Proofs.Theorem.Online.optimal_mistake_bound_eq_ldim}), and determine the quantitative dependence of the private sample complexity on $\Ldim(\mathcal{H})$.
\end{openprob}

\begin{openprob}[Quantitative route stability, \textsc{uk}]
\label{op:route-stability-quant}
\Cref{rmk:route-stability} asserts stability away from score ties.  Quantify it: bound the measure of the score-tie boundary set, and the perturbation margin that preserves the route, as functions of the inverse temperature, exhibiting the trade-off between routing sharpness (which improves the Hedge approximation) and stability (which preserves the inherited regret guarantee).
\end{openprob}

\begin{openprob}[Bandit attention, \textsc{ku}]
\label{op:bandit-attention}
In the partial-information (bandit) setting only the loss of the routed key is observed, not the full label.  Determine the regret of online attention routing under bandit feedback (the $\sqrt{kT}$-versus-$\sqrt{T\ln k}$ gap between bandit and full-information exponential weights) and whether route stability mitigates the exploration cost.
\end{openprob}

\begin{openprob}[Halving versus the $\SOA$, \textsc{ku}]
\label{op:halving-soa}
\Cref{rmk:soa-halving} contrasts cardinality with Littlestone dimension.  Characterize the classes on which Halving is sub-optimal by an unbounded factor, where $\log_2|\mathcal{H}|$ (or its effective analogue) exceeds $\Ldim(\mathcal{H})$ arbitrarily, and give the structural property of $\mathcal{H}$ that the $\SOA$ exploits and Halving does not.
\end{openprob}

\begin{openprob}[Online-to-batch for attention, \textsc{uk}]
\label{op:online-to-batch-attention}
The implication online $\Rightarrow$ PAC is verified (\leanname{FLT_Proofs.Theorem.Separation.online_imp_pac}).  Make it quantitative for attention: convert the online mistake bound of a finite-Littlestone attention class into a PAC sample-complexity bound via online-to-batch, and decide whether the conversion is tight or loses a factor, sharpening the relation between \cref{op:ldim-attention} here and the VC question \cref{op:attention-vc} of \Cref{ch:pac}.
\end{openprob}


\subsection*{Counterfactual exercises}

The characterization $\mathrm{Opt}(\mathcal{H})=\Ldim(\mathcal{H})$ rests on several interlocking assumptions. Each exercise below removes or modifies one feature of the game and asks what survives. Each exercise asks you to locate which piece of the argument a given assumption is carrying.

\begin{enumerate}[leftmargin=*,itemsep=3pt]
\item \textbf{The value is two-sided.} The equality $\mathrm{Opt}=\Ldim$ is a conjunction of two theorems: the Standard Optimal Algorithm makes at most $\Ldim$ mistakes (\leanname{soa_mistakes_bounded}) and the adversary forces at least $\Ldim$ (\leanname{adversary_lower_bound}). Show that either half alone gives only a one-sided bound, and that the matching of the two is what makes the value exact.

\item \textbf{Drop realizability.} Remove the constraint that some hypothesis fits every label. The mistake bound becomes vacuous and the right benchmark is regret against the best fixed hypothesis. Show why the agnostic value calls for the regret framework rather than the Littlestone characterization (\cref{op:regret-constant}).

\item \textbf{The gap to PAC.} Online learnability implies PAC learnability but the converse fails. Exhibit the witness: thresholds on $\R$ have $\VC=1$ and $\Ldim=\infty$ (\leanname{online_strictly_stronger_pac}, \leanname{pac_not_implies_online}). Conclude that an adaptive adversary is strictly stronger than any random distribution.

\item \textbf{The conservative leg.} A finite Littlestone bound forces a finite VC bound, so online learnability gives PAC learnability and adds no statistical power (\leanname{online_imp_pac}). Prove this one-way half of the separation directly from $\VC\le\Ldim$.

\item \textbf{An efficient optimal learner.} The $\SOA$ achieves the value information-theoretically, but computing the version-space Littlestone dimensions it needs may be intractable. Investigate which classes admit a polynomial-time-per-round optimal online learner (\cref{op:soa-computable}).

\item \textbf{Depth exceeds set-size.} On thresholds on $\R$ the adversary bisects the surviving interval forever, so the tree dimension is infinite while the shattered-set dimension is one. Show that $\Ldim(\mathcal{H})$ can exceed $\VC(\mathcal{H})$ without bound, and that the slack measures adaptivity (\leanname{ldim_threshold_top}).

\item \textbf{The embedding $\VC\le\Ldim$.} A shattered set of size $d$ becomes a non-adaptive mistake tree of depth $d$ by presenting the same instance at every node of a given depth. Verify that this map preserves size as depth and validity as shattering, so $\VC(\mathcal{H})\le\Ldim(\mathcal{H})$ always (\leanname{BranchWiseLittlestoneDim_ge_VCDim}).

\item \textbf{Multiclass labels.} Replace binary labels by a finite label set, giving each internal node one child per label. Conjecture a multiclass tree-dimension pinning the optimal mistake bound, and determine whether the clean $\mathrm{Opt}=\Ldim$ identity survives (\cref{op:online-multiclass}).
\end{enumerate}

% Chapter 7: Identification in the Limit
% ex_learning = Profile B (load-bearing), bc_learning = Profile C (revelation),
% mind_change_characterization = Profile B, online_pac_gold_separation = Profile D
%
% Texture: recursion-theoretic. Proofs use diagonalization, not concentration
% or combinatorial trees. Infinite horizon, no sample complexity, ordinal-valued
% mind-change bounds. This chapter should feel fundamentally different from
% Ch5 (PAC) and Ch6 (Online).

\chapter{Identification in the Limit}
\label{ch:gold}

In 1967, seventeen years before Valiant introduced PAC learning, E.\ Mark Gold posed a question that remains the oldest open research programme in computational learning theory: which classes of recursive functions can a machine identify from examples?

Gold's framework differs from PAC and online learning not merely in its definitions but in its \emph{proof technique}. PAC theory rests on concentration inequalities: the probabilistic convergence of empirical risk to true risk. Online learning rests on combinatorial game trees: the Littlestone dimension measures the depth of a binary tree that the adversary can force. Gold-style identification rests on \emph{diagonalization}: the recursion-theoretic technique of defeating every candidate learner by constructing an adversarial input that forces failure. The proof method shapes the entire chapter.

Three features distinguish this paradigm from everything that precedes it in this book:

\begin{enumerate}[leftmargin=*,itemsep=3pt]
\item \textbf{Infinite horizon.} The learner receives an infinite stream of data and must eventually converge. There is no sample complexity bound, no $\varepsilon$, no $\delta$. The question is whether the learner converges, not how fast.

\item \textbf{Ordinal-valued complexity.} When we ask ``how many times does the learner change its mind?'' the answer is not an integer but a \emph{transfinite ordinal}. This is the first appearance of ordinals in learning theory.

\item \textbf{A lattice of success criteria.} PAC learning has one definition (modulo agnostic, realizable, improper variants). Gold-style identification has a rich hierarchy: $\FIN < \Ex < \BC$, with anomalous, monotonic, and vacillatory learning branching off. The hierarchy itself is mathematical content.
\end{enumerate}

\begin{historical}
Gold's 1967 paper~\cite{Gold1967} was preceded by his 1965 paper on limiting recursion~\cite{Gold1965}, which introduced the idea that a computation could ``converge in the limit'' even if it never halts in the classical sense. The 1967 paper applied this idea to language learning. It was the first mathematical theory of learning from examples, predating Valiant's PAC model by 17 years, Vapnik and Chervonenkis's foundational work on uniform convergence by 4 years, and Littlestone's online model by 21 years. The diagonalization technique Gold used to prove his impossibility theorem later became standard in computational complexity theory, but Gold's use predates most of those applications.
\end{historical}

% ================================================================
% SECTION 1: Gold's Question
% ================================================================

\section{Gold's Question}
\label{sec:gold-question}

Fix a countable domain. In Gold's original setting, this is the set of all strings over a finite alphabet $\Sigma$, and the objects to be learned are formal languages $L \subseteq \Sigma^*$. We work in this setting when it clarifies the recursion-theoretic content, and in the general setting of concept classes $\mathcal{C} \subseteq 2^\N$ when it does not matter.

The learner receives data sequentially, one datum at a time, forever. Two data presentation modes are standard:

\begin{defn}[Text and Informant]
\label{def:text-informant}
Let $L \subseteq \Sigma^*$ be a language.
\begin{itemize}[leftmargin=*,itemsep=2pt]
\item A \emph{text} for $L$ is an infinite sequence $t = (t_0, t_1, t_2, \ldots)$ with $\{t_i : i \in \N\} = L$. Every element of $L$ appears at least once; no element of $\Sigma^* \setminus L$ ever appears. The text may contain repetitions and need not follow any fixed ordering.
\item An \emph{informant} for $L$ is an infinite sequence of labelled pairs $((s_0, b_0), (s_1, b_1), \ldots)$ where $\{s_i : i \in \N\} = \Sigma^*$ and $b_i = 1$ if and only if $s_i \in L$. Every string is eventually presented together with its membership status.
\end{itemize}
\end{defn}

Text presents only positive examples; informant presents both positive and negative examples. This distinction matters enormously: Gold showed that informant is strictly more powerful than text. We focus on text, which is the harder and more interesting case.

After seeing the first $n$ data points $(t_0, \ldots, t_{n-1})$, the learner outputs a hypothesis $h_n$. The hypothesis is an \emph{index}, a natural number naming a program that computes a language. The learner has no deadline, no accuracy target, and no probability distribution on targets. It simply sees more and more of $L$ and must eventually guess correctly.

% ================================================================
% SECTION 2: Ex-Learning and Gold's Impossibility Theorem
% ================================================================

\section{Ex-Learning and Gold's Impossibility Theorem}
\label{sec:ex-gold}

\begin{defn}[Explanatory Learning ($\Ex$)]
\label{def:ex-learning}
A learner $M$ \emph{$\Ex$-identifies} a language $L$ from text if, for every text $t$ for $L$, there exists an index $e$ and a time $n_0$ such that $M(t_0, \ldots, t_n) = e$ for all $n \geq n_0$, and $W_e = L$ (where $W_e$ is the language computed by program $e$).

A class $\mathcal{L}$ of languages is \emph{$\Ex$-identifiable} if there exists a single learner $M$ that $\Ex$-identifies every $L \in \mathcal{L}$ from text.\formal{FLT_Proofs.Criterion.Gold.EXLearnable}
\end{defn}

The definition requires \emph{syntactic convergence}: the learner must eventually output the same index forever. It is not enough to output a sequence of indices that all happen to compute the correct language; the index itself must stabilize. (Relaxing this requirement gives BC-learning; see \Cref{sec:bc-learning}.)

\begin{example}[$\Ex$-identification of finite languages]
\label{ex:finite-languages}
Let $\mathcal{L}_{\mathrm{fin}}$ be the class of all finite subsets of $\N$. Here is an $\Ex$-learner for $\mathcal{L}_{\mathrm{fin}}$: at time $n$, output an index for $\{t_0, \ldots, t_n\}$. After all elements of $L$ have appeared (which happens at some finite time $n_0$, since $L$ is finite), the set $\{t_0, \ldots, t_n\}$ equals $L$ for all $n \geq n_0$. The hypothesis stabilizes.
\end{example}

This simple algorithm illustrates an asymmetry: the learner does not know \emph{when} it has converged. It has no way to announce ``I am done.'' It merely stabilizes, silently. This lack of a convergence signal is what makes Gold's impossibility theorem possible.

\begin{thm}[Gold's Impossibility Theorem {\cite{Gold1967}}]
\label{thm:gold}
Let $\mathcal{L}$ be any class of languages that contains all finite languages and at least one infinite language. Then $\mathcal{L}$ is not $\Ex$-identifiable from text.\formal{FLT_Proofs.Theorem.Gold.gold_theorem}
\end{thm}

This is the centerpiece of the chapter. The proof is a diagonalization argument: we defeat \emph{every} candidate learner by constructing a text that forces it to fail. The construction is adversarial, not probabilistic: there is no distribution, no $\varepsilon$-net, no covering number. The adversary simply watches the learner and manipulates the future of the stream to prevent convergence.

\begin{proof}
Let $M$ be any learner. We construct a text $t$ for some language $L \in \mathcal{L}$ such that $M$ fails to $\Ex$-identify $L$ on $t$.

We build $t$ in stages. Let $L_\infty \in \mathcal{L}$ be the infinite language that $\mathcal{L}$ contains by assumption. Enumerate $L_\infty = \{w_0, w_1, w_2, \ldots\}$.

\medskip
\noindent\textbf{Stage 0.} Present $w_0$ to $M$. The learner outputs some hypothesis $h_0 = M(w_0)$.

\medskip
\noindent\textbf{Stage $k \geq 1$.} At this point we have presented the finite sequence $(w_0, \ldots, w_{n_{k-1}})$ for some $n_{k-1}$, and $M$ has output hypothesis $h_{n_{k-1}}$. Let $F_k = \{w_0, \ldots, w_{n_{k-1}}\}$ be the set of strings presented so far.

Consider two cases:
\begin{itemize}[leftmargin=*]
\item \textbf{Case A:} $W_{h_{n_{k-1}}} = F_k$ (the current hypothesis computes exactly the finite set seen so far). Then present $w_{n_{k-1}+1}$ (the next element of $L_\infty$), enlarging $F_k$. Let $n_k = n_{k-1} + 1$. Proceed to stage $k+1$.

\item \textbf{Case B:} $W_{h_{n_{k-1}}} \neq F_k$. Then the current hypothesis is already wrong for the finite language $F_k$. We may define $L = F_k$: since $F_k$ is finite and $\mathcal{L}$ contains all finite languages, $F_k \in \mathcal{L}$. We extend the text by repeating elements of $F_k$ forever. The learner $M$ has output a hypothesis $h_{n_{k-1}}$ with $W_{h_{n_{k-1}}} \neq F_k = L$, and the text for $L$ can be arranged so that $M$ never converges to a correct index (we continue to the next stage rather than stopping, as shown below).
\end{itemize}

The key observation is the \emph{diagonalization}: we never let the learner rest.

If Case~A occurs at every stage, then every element of $L_\infty$ is eventually presented, so $t$ is a text for $L_\infty$. At each stage, $M$'s current hypothesis $h_{n_k}$ is checked: does $W_{h_{n_k}} = F_k$? If yes, we expand. But $F_k \subsetneq L_\infty$ for all $k$ (since $L_\infty$ is infinite), so $W_{h_{n_k}} = F_k \neq L_\infty$. Therefore $M$ outputs a hypothesis that is wrong at every stage; it cannot converge to a correct index for $L_\infty$.

If Case~B occurs at some stage $k$, then $M$ has already failed for $F_k$. We commit to $L = F_k$ and repeat elements of $F_k$ forever. But we can do more: before committing, we wait to see if $M$ changes its mind. If $M$ later outputs a hypothesis $h'$ with $W_{h'} = F_k$, we switch to Case~A and add the next element of $L_\infty$. This forces $M$ to fail for $L_\infty$ instead.

In either case, $M$ fails. Since $M$ was an arbitrary learner, no learner $\Ex$-identifies $\mathcal{L}$.
\end{proof}

The proof has a recursive structure that repays careful study. The adversary maintains a ``threat'' at every stage: either the current language is the finite set seen so far, or it is the infinite language $L_\infty$. Whenever the learner commits to one possibility, the adversary switches to the other. The learner is forced to change its mind infinitely often, and $\Ex$-identification requires that it eventually stop changing its mind.

\begin{remark}[Diagonalization vs.\ concentration]
\label{rmk:diag-vs-conc}
Compare this proof to the lower bound for PAC learning (\Cref{ch:pac}). The PAC lower bound constructs a \emph{distribution} that fools the learner with positive probability; it is probabilistic. Gold's proof constructs a \emph{text} that fools the learner with certainty; it is adversarial. The PAC proof uses counting (how many hypotheses can be distinguished by $m$ samples); Gold's proof uses the halting problem implicitly (the adversary must check whether $W_e = F$ for arbitrary $e$). These are entirely different mathematical worlds.
\end{remark}

\begin{historical}
The proof technique of Theorem~\ref{thm:gold} is a \emph{priority argument} in the sense of recursion theory, though a simple one. More sophisticated priority arguments appear in the study of learning with resource bounds (Case and Smith~\cite{CaseSmith1983}). The connection between Gold-style learning and recursion theory runs deep: Barzdin\v{s}~\cite{Barzdinsh1974} showed that the class of languages $\Ex$-identifiable from informant is exactly the class of $\Sigma^0_2$-definable families, establishing a precise link to the arithmetical hierarchy.
\end{historical}

% ================================================================
% SECTION 3: The Hierarchy: BC > Ex > FIN
% ================================================================

\section{The Identification Hierarchy}
\label{sec:hierarchy}

Gold's impossibility theorem says that $\Ex$-identification is limited. Two natural responses: strengthen the success criterion (demand faster convergence) or weaken it (demand less of the final hypothesis). Both directions are fruitful.

\subsection{Finite Identification}
\label{sec:fin-learning}

\begin{defn}[Finite Identification ($\FIN$)]
\label{def:fin-learning}
A learner $M$ \emph{$\FIN$-identifies} a language $L$ from text if, for every text $t$ for $L$, there exists a time $n_0$ such that $M$ outputs exactly one hypothesis: $M$ outputs ``?'' (no guess) for $n < n_0$ and outputs a single index $e$ at time $n_0$ with $W_e = L$, never changing its mind. A class $\mathcal{L}$ is \emph{$\FIN$-identifiable} if a single learner $\FIN$-identifies every $L \in \mathcal{L}$.\formal{FLT_Proofs.Criterion.Gold.FiniteLearnable}
\end{defn}

$\FIN$ is the strongest identification criterion: the learner must produce the correct answer in finitely many steps with no subsequent revision. The class of $\FIN$-identifiable families is a strict subset of $\Ex$-identifiable families; for instance, the class of all finite languages is $\Ex$-identifiable (\Cref{ex:finite-languages}) but not $\FIN$-identifiable, since the learner cannot know when the last element has been seen.

\subsection{Behaviorally Correct Learning}
\label{sec:bc-learning}

\begin{defn}[Behaviorally Correct Learning ($\BC$)]
\label{def:bc-learning}
A learner $M$ \emph{$\BC$-identifies} a language $L$ from text if, for every text $t$ for $L$, there exists a time $n_0$ such that $W_{M(t_0, \ldots, t_n)} = L$ for all $n \geq n_0$. That is, from time $n_0$ onward, every hypothesis the learner outputs is \emph{extensionally correct} (computes the right language), but the indices may keep changing.\formal{FLT_Proofs.Criterion.Gold.BCLearnable}
\end{defn}

The distinction between $\Ex$ and $\BC$ is subtle but consequential. $\Ex$ requires \emph{syntactic} convergence: the index stabilizes. $\BC$ requires only \emph{semantic} convergence: the computed language stabilizes, even if the learner keeps switching between different programs that all compute the same language. At first glance, this seems like a minor relaxation. It is not.

\begin{thm}[Case--Smith {\cite{CaseSmith1983}}]
\label{thm:bc-strictly-stronger}
$\BC$ is strictly more powerful than $\Ex$: there exists a class of languages that is $\BC$-identifiable but not $\Ex$-identifiable from text.
\end{thm}

\begin{proof}
We construct a class $\mathcal{L}$ that separates $\BC$ from $\Ex$.

For each total recursive function $f : \N \to \N$, define the language $L_f = \{\langle n, f(n)\rangle : n \in \N\}$, where $\langle \cdot, \cdot \rangle$ is a computable pairing function. Let $\mathcal{L} = \{L_f : f \text{ is total recursive}\}$.

\medskip\noindent\textbf{$\mathcal{L}$ is $\BC$-identifiable.} Given a text $t$ for some $L_f$, at time $n$ the learner has seen finitely many pairs $\langle n_i, m_i\rangle$. Define the partial function $g_n$ by $g_n(n_i) = m_i$ for all pairs seen so far, and $g_n(k) = 0$ for all other $k$. Output an index for the language $L_{g_n} = \{\langle k, g_n(k)\rangle : k \in \N\}$. Once all pairs $\langle k, f(k)\rangle$ for $k \leq K$ have appeared (for any $K$), the hypothesis computes $L_f$ correctly on $\{0, \ldots, K\}$. As $n \to \infty$, $g_n$ converges pointwise to $f$. Since $f$ is total, for every $k$ there exists a time after which $g_n(k) = f(k)$. The language computed by the hypothesis is eventually $L_f$, though the \emph{index} keeps changing as new pairs are absorbed. This is $\BC$-identification.

\medskip\noindent\textbf{$\mathcal{L}$ is not $\Ex$-identifiable.} Suppose for contradiction that a learner $M$ $\Ex$-identifies $\mathcal{L}$. We diagonalize against $M$. Define a total recursive function $f$ as follows.

Begin presenting pairs $\langle 0, 0\rangle, \langle 1, 0\rangle, \langle 2, 0\rangle, \ldots$ (corresponding to the zero function). Wait until $M$ converges to some index $e$. Since $M$ $\Ex$-identifies $\mathcal{L}$ and the zero function is total recursive, $M$ must converge. Now $W_e = L_{\mathbf{0}}$ (the language of the zero function).

Modify $f$: set $f(k) = 1$ for some large $k$ not yet presented. This changes the target to a different total recursive function, but the text presented so far is consistent with both targets. The learner $M$ has already committed to index $e$, which computes $L_{\mathbf{0}} \neq L_f$. By a careful effectivization of this argument (choosing $k$ computably based on $M$'s behavior), we construct $f \in \mathcal{L}$ on which $M$ fails. This contradicts $M$ $\Ex$-identifying all of $\mathcal{L}$.

The essential point is that $\BC$-identification does not require the \emph{index} to stabilize, only the \emph{extension}. The class $\mathcal{L}$ exploits this: the learner must continually update its program as new pairs arrive, and the programs keep changing, but they all eventually compute the same language. $\Ex$-identification cannot tolerate this perpetual updating of indices.
\end{proof}

\begin{figure}[ht]
\centering
\begin{tikzpicture}[
    criterion/.style={draw, thick, rounded corners, minimum width=2.2cm, minimum height=0.85cm,
                      font=\small\bfseries, align=center},
    branch/.style={draw, rounded corners, minimum width=2cm, minimum height=0.7cm,
                   font=\small, align=center},
    strict/.style={-{Stealth[length=2mm]}, thick, sepgreen!70!black},
    weak/.style={-{Stealth[length=2mm]}, thin, dashed, notegray},
    witlabel/.style={font=\tiny, text=edgeorange, fill=white, inner sep=1.5pt, align=center},
    hdl/.style={font=\tiny, fill=white, inner sep=1pt, align=center}
]
% --- spine (layered, x=0; bottom-to-top = increasing power) ---
\node[criterion, fill=defblue!15] (fin) at (0,0)   {$\FIN$};
\node[criterion, fill=defblue!15] (ex)  at (0,2.7) {$\Ex$};
\node[criterion, fill=defblue!15] (bc)  at (0,5.4) {$\BC$};
% --- branches: restricts LEFT, extends RIGHT ---
\node[branch, fill=edgeorange!8] (mono)  at (-3.5,2.7) {Monotonic\\$\mathrm{Mon}\Ex$};
\node[branch, fill=edgeorange!8] (anom)  at ( 3.6,1.9) {Anomalous\\$\Ex^*$};
\node[branch, fill=goldcolor!10] (vac)   at ( 3.6,4.0) {Vacillatory\\$\mathrm{Vex}$};
\node[branch, fill=goldcolor!10] (trial) at ( 3.6,6.0) {Trial \&\\Error};
% --- spine edges: FIN<Ex SOLID (verified), Ex<BC DASHED (literature) ---
\draw[strict] (fin) -- node[hdl, left, xshift=-2pt]
   {all finite langs\\are $\Ex$, not $\FIN$\\{\tiny\leanname{ex_strictly_stronger_finite}}} (ex);
\draw[weak] (ex) -- node[hdl, left, xshift=-2pt]
   {Case--Smith 1983\\(\cref{thm:bc-strictly-stronger})\\\itshape literature, not kernel} (bc);
% --- relaxation edges: all DASHED, oriented by constraint-toggle direction ---
\draw[weak] (anom) -- node[witlabel, above, pos=0.45]{\rel{extends}\\$-$\,zero-error} (ex);
\draw[weak] (mono) -- node[witlabel, above]{\rel{restricts}\\$+$\,no-retraction} (ex);
\draw[weak] (vac) -- node[witlabel, below right, xshift=-3pt]{\rel{extends}} (ex);
\draw[weak] (vac) -- node[witlabel, above right, xshift=-3pt]{\rel{restricts}} (bc);
\draw[weak] (bc) -- node[witlabel, above right, xshift=-3pt]{$\Delta^0_2$\\(Putnam--Gold)} (trial);
\end{tikzpicture}
\caption[The identification hierarchy as a chain of separations]{The identification hierarchy, read as a
chain of strict separations refined by a lattice of single-constraint relaxations. Each spine node
($\FIN$, $\Ex$, $\BC$) is a success criterion; each spine edge is a separation theorem, a class one
criterion learns and the other cannot, not a more liberal definition. The bottom rung
$\FIN\subsetneq\Ex$ is verified in the kernel (\leanname{ex_strictly_stronger_finite}), drawn solid;
the witness is the class of all finite languages, which is $\Ex$-identifiable but not
$\FIN$-identifiable. The upper rung $\Ex\subsetneq\BC$ is the Case--Smith separation, where $\Ex$
demands the index converge (syntactic) and $\BC$ demands only the language converge (semantic); it is
established in the literature and not yet in the kernel, so it is drawn dashed and cited. The side nodes
are single-constraint relaxations of $\Ex$: anomalous $\Ex^*$ removes the zero-error constraint and so
extends $\Ex$; monotonic $\mathrm{Mon}\Ex$ adds a no-retraction constraint and so restricts $\Ex$;
vacillatory $\mathrm{Vex}$ allows finitely many final hypotheses and sits strictly between $\Ex$ and
$\BC$; trial-and-error learning pins the limit-computable predicates to the $\Delta^0_2$ level. Every
branch and the $\Delta^0_2$ link are literature, not kernel theorems, hence dashed.}
\label{fig:identification-hierarchy}
\end{figure}

\begin{traversal}
\textbf{Path:} $\nde{fin\_learning} \xrightarrow{\rel{strictly\_stronger}} \nde{ex\_learning} \xrightarrow{\rel{strictly\_stronger}} \nde{bc\_learning}$.\\[3pt]
Each arrow is witnessed by an explicit separation. The graph encodes these as \rel{strictly\_stronger} edges with the separation witness as metadata. The hierarchy is not a sequence of increasingly liberal definitions; it is a chain of \emph{theorems}, each proved by constructing a class that one criterion can learn and the other cannot.
\end{traversal}

% ================================================================
% SECTION 4: Relaxations and Their Lattice
% ================================================================

\section{Relaxations of Identification}
\label{sec:relaxations}

The $\FIN$--$\Ex$--$\BC$ chain is the spine of the hierarchy. Several natural relaxations branch off from $\Ex$, each modifying a different aspect of the convergence requirement.

\subsection{Anomalous Learning}
\label{sec:anomalous-gold}

\begin{defn}[Anomalous Learning ($\Ex^*$)]
\label{def:anomalous}
A learner \emph{$\Ex^*$-identifies} a language $L$ if it $\Ex$-converges to an index $e$ such that $W_e$ differs from $L$ on at most finitely many strings. That is, the final hypothesis may contain finitely many errors: finitely many strings incorrectly included or excluded.\formal{FLT_Proofs.Criterion.Gold.AnomalousLearnable}
\end{defn}

Anomalous learning relaxes $\Ex$ by removing the requirement of exact correctness, replacing it with correctness up to a finite set. In graph terms, this is a \rel{restricts} edge from $\nde{anomalous\_learning}$ to $\nde{ex\_learning}$: the constraint (zero anomalies) is removed, and no new grammatical structure is introduced. The class of $\Ex^*$-identifiable families strictly contains the class of $\Ex$-identifiable families~\cite{CaseSmith1983}.

\subsection{Monotonic Learning}
\label{sec:monotonic}

\begin{defn}[Monotonic Learning ($\mathrm{Mon}\Ex$)]
\label{def:monotonic}
A learner \emph{monotonically} $\Ex$-identifies $L$ if it $\Ex$-identifies $L$ and, whenever it changes its hypothesis from $e$ to $e'$, the new hypothesis is at least as inclusive on the data seen so far: $W_e \cap \{t_0, \ldots, t_n\} \subseteq W_{e'} \cap \{t_0, \ldots, t_n\}$. Once a datum is correctly classified, that classification is never retracted.\formal{FLT_Proofs.Criterion.Gold.MonotonicLearnable}
\end{defn}

Monotonic learning \emph{strengthens} $\Ex$ by adding a constraint. The class of $\mathrm{Mon}\Ex$-identifiable families is a strict subset of $\Ex$-identifiable families. In the graph, $\nde{monotonic\_learning} \xrightarrow{\rel{restricts}} \nde{ex\_learning}$.

\subsection{Vacillatory Learning}
\label{sec:vacillatory}

\begin{defn}[Vacillatory Learning ($\mathrm{Vex}$)]
\label{def:vacillatory}
A learner \emph{vacillatorily} identifies $L$ if it eventually oscillates among finitely many indices $e_1, \ldots, e_k$, all satisfying $W_{e_i} = L$. The learner never settles on a single index, but all its eventual outputs are extensionally correct and drawn from a finite set.\formal{FLT_Proofs.Criterion.Gold.VacillatoryLearnable}
\end{defn}

Vacillatory learning sits between $\Ex$ (which requires convergence to a single index) and $\BC$ (which allows infinitely many correct indices). It is strictly more powerful than $\Ex$ and strictly less powerful than $\BC$.

\subsection{Trial and Error}
\label{sec:trial-error}

\begin{defn}[Trial-and-Error Learning]
\label{def:trial-error}
A \emph{trial-and-error} predicate for a set $A \subseteq \N$ is a computable function $f : \N \to \{0, 1\}$ such that $\lim_{s \to \infty} f_s(n)$ exists for all $n$ and equals the characteristic function of $A$. The learner may ``retract'' previous outputs: it learns by making mistakes and correcting them.\formal{FLT_Proofs.Criterion.Gold.TrialAndErrorLearnable}
\end{defn}

Trial-and-error learning connects identification in the limit to the arithmetical hierarchy. A set $A$ is trial-and-error learnable if and only if $A \in \Delta^0_2$, the class of sets that are both $\Sigma^0_2$ and $\Pi^0_2$. This is the precise recursion-theoretic characterization, due to Putnam~\cite{Gold1965} and Gold: the limit-computable functions are exactly the $\Delta^0_2$ functions.\formal{FLT_Proofs.Computation.Delta02Class} This connection places Gold-style learning firmly within the landscape of classical recursion theory.

% ================================================================
% SECTION 5: Mind-Change Complexity and Ordinals
% ================================================================

\section{Mind-Change Complexity}
\label{sec:mind-change}

Gold's impossibility theorem tells us that certain classes cannot be identified at all. For classes that \emph{can} be identified, a natural quantitative question arises: how many times must the learner change its mind before converging?

\begin{defn}[Mind Change]
\label{def:mind-change}
A \emph{mind change} occurs at time $n$ if the learner's output satisfies $M(t_0, \ldots, t_n) \neq M(t_0, \ldots, t_{n-1})$. The \emph{mind-change count} of $M$ on text $t$ is the number of times $M$ changes its hypothesis.\formal{FLT_Proofs.Complexity.MindChange.MindChangeCount}
\end{defn}

For finite mind-change bounds, the theory is straightforward: we say that $M$ identifies $\mathcal{L}$ with at most $k$ mind changes if, on every text for every $L \in \mathcal{L}$, the mind-change count is at most $k$. The class of all finite languages is identifiable with 0 mind changes after the first hypothesis (the learner in \Cref{ex:finite-languages} changes its mind every time a new element appears, but a more careful learner can be designed with bounded mind changes for restricted subclasses).

The surprise (genuinely unexpected for readers coming from PAC theory) is that integer-valued bounds are \emph{not sufficient} to characterize the full landscape.

\begin{thm}[Freivalds--Smith {\cite{FreivaldsSmith1993}}]
\label{thm:mind-change-ordinal}
The mind-change complexity of $\Ex$-identification is naturally measured by \emph{countable ordinals}. Specifically:
\begin{enumerate}[label=(\roman*),leftmargin=*]
\item For every countable ordinal $\alpha$, one can define what it means for a learner to identify a class with mind-change bound $\alpha$.
\item There exist classes identifiable with mind-change bound $\omega$ (the first infinite ordinal) that cannot be identified with any finite mind-change bound.
\item More generally, for every ordinal $\alpha < \omega_1$, there exist classes identifiable with mind-change bound $\alpha$ but not with any bound $\beta < \alpha$.
\end{enumerate}\formal{FLT_Proofs.Complexity.MindChange.MindChangeOrdinal}\formal{FLT_Proofs.Theorem.Gold.mind_change_characterization}
\end{thm}

The idea behind ordinal mind-change bounds is as follows. A learner with mind-change bound $\omega$ begins with an ordinal counter set to $\omega$. At each mind change, the counter must decrease, but it may decrease to any smaller ordinal. Since there is no infinite descending sequence of ordinals (ordinals are well-ordered), the learner must eventually stop changing its mind. The counter $\omega$ allows finitely many mind changes, but the \emph{number} of mind changes need not be bounded in advance by any fixed integer; it may depend on the input.

This is fundamentally different from an integer bound. With a bound of $k \in \N$, the learner may change its mind at most $k$ times on \emph{every} input. With a bound of $\omega$, the learner may change its mind $k$ times for input-dependent $k$, with no uniform finite upper bound. The ordinal hierarchy continues: $\omega \cdot 2$ allows the learner to ``reset'' its finite counter once; $\omega^2$ allows nested levels of resetting; and so on up through the constructive ordinals.

\begin{remark}[Ordinals in learning theory]
\label{rmk:ordinals-surprise}
The appearance of transfinite ordinals in learning theory is a genuine structural surprise. PAC learning theory uses real-valued parameters ($\varepsilon, \delta, m(\varepsilon, \delta)$). Online learning uses integer-valued dimensions ($\Ldim$). Gold-style learning requires ordinal-valued complexity measures. Each proof technique brings its own number system. The full development of ordinal mind-change complexity is deferred to \Cref{ch:ordinals}, where it interacts with the constructive ordinal notation systems of Kleene.
\end{remark}

\begin{traversal}
\textbf{Path:} $\nde{mind\_change\_characterization} \xrightarrow{\rel{measures}} \nde{ex\_learning}$.\\[3pt]
The mind-change ordinal is a complexity measure on $\Ex$-identifiable classes, analogous to the VC dimension for PAC-learnable classes. But where VC dimension is a single integer, mind-change complexity is an ordinal, potentially transfinite. This is a \rel{measures} edge of a qualitatively different kind than $\nde{vc\_dimension} \xrightarrow{\rel{measures}} \nde{concept\_class}$.
\end{traversal}

% ================================================================
% SECTION 6: The Three-Paradigm Separation
% ================================================================

\subsection{In-context learning as identification in the limit}
\label{sec:icl-gold}

Gold's framework (infinite horizon, a stream of data, a sequence of conjectures) describes a learning mode that modern transformers exhibit directly: in-context learning.  A transformer fed a growing prompt of labelled examples produces, at each prefix length, a prediction rule: its conjecture given what it has seen so far.  The conjecture map ``finite data so far $\mapsto$ predictor'' is exactly the type of a Gold learner.

\begin{prop}[A transformer in context is a Gold learner]
\label{prop:transformer-gold}
A parametric attention learner $T$, read as the map from a finite prefix $((x_1,y_1),\dots,(x_n,y_n))$ to the predictor it computes in context, is a Gold learner: a function from finite labelled sequences to concepts.\formal{FLT_Proofs.Learner.Core.GoldLearner} Its behaviour on a target and data stream is therefore classified by the criteria of this chapter ($\Ex$, $\BC$, vacillatory), and its mind-change complexity is the ordinal-valued measure of \cref{thm:mind-change-ordinal} applied to $T$.\formal{FLT_Proofs.Complexity.MindChange.MindChangeOrdinal}
\end{prop}

\begin{proof}
The Gold-learner type is a conjecture map from $\mathrm{List}(X \times Y)$ to concepts, and the in-context predictor of $T$ at prefix $s$ is such a map: feed $s$ as context and return the function $x \mapsto T(s)(x)$.  The mind-change count and ordinal of \cref{def:mind-change,thm:mind-change-ordinal} are defined for every Gold learner, so they apply to $T$ verbatim.
\end{proof}

The proposition turns a practical question (does in-context learning ``converge''?) into the precise question of which identification criterion $T$ meets.  Syntactic convergence, where the in-context predictor stops changing, is $\Ex$; extensional convergence, where the predictor keeps changing but computes a fixed function, is $\BC$; oscillation among finitely many correct predictors is vacillatory.  Gold's impossibility theorem then bites: no transformer $\Ex$-identifies, from positive data alone, a class containing all finite concepts and one infinite concept (\cref{thm:gold}).  Whether realistic in-context learning $\Ex$-converges, $\BC$-converges, or vacillates, and with what mind-change ordinal, is open (\cref{op:icl-identification}).

\section{Three Paradigms: Online Inside PAC, Gold Apart}
\label{sec:three-paradigm}

We now have three learning paradigms: PAC (\Cref{ch:pac}), online (\Cref{ch:online}), and Gold. Each defines ``learnable'' differently. The question is: what is the relationship?

The answer is a partial order, not an antichain. Online learnability is \emph{strictly contained} in PAC learnability: every online-learnable class is PAC-learnable, and the converse fails, so the two are comparable. Gold-style identification, by contrast, is incomparable to both: it identifies classes of unbounded VC and Littlestone dimension that no statistical learner can handle, while it fails on classes with no convergence signal that PAC and online learners handle easily. The deep structural fact is that statistical capacity (PAC and online) and convergence in the limit (Gold) are orthogonal axes, and that adaptivity, the only thing separating online from PAC, buys nothing here: a finite Littlestone dimension forces a finite VC dimension, so online sits properly inside PAC.

\begin{separation}
\textbf{Theorem} (PAC, online, and Gold).\formal{FLT_Proofs.Theorem.Separation.online_pac_gold_separation} \textit{Online learnability is strictly contained in PAC learnability, so the two are comparable; Gold-style identification is incomparable to both. The machine-checked statement establishes the threshold separation (PAC but not online) and separation (a) below (the finite languages, $\Ex$ but not PAC); the remaining separations (b)--(d) are classical:}

\begin{enumerate}[label=(\alph*),leftmargin=*,itemsep=4pt]
\item \textbf{$\Ex$-learnable but not PAC-learnable.} The class $\mathcal{L}_{\mathrm{fin}}$ of all finite subsets of $\N$ is $\Ex$-identifiable from text (\Cref{ex:finite-languages}). But the associated concept class has infinite VC dimension (any finite set of points can be shattered), so it is not PAC-learnable.

The witness exploits the fundamental difference in how the two paradigms treat ``all finite subsets.'' $\Ex$-identification succeeds because on any \emph{fixed} finite set, the learner eventually sees all elements. PAC fails because the learner must handle \emph{all} finite sets simultaneously with bounded sample complexity, and the VC dimension is infinite.

\item \textbf{PAC-learnable but not $\Ex$-identifiable.} The class of threshold functions $\mathcal{C} = \{x \mapsto \ind[x \geq \theta] : \theta \in \R\}$ over $\R$ is PAC-learnable ($\VC = 1$). But it is not $\Ex$-identifiable from text, because a text for a threshold language $\{\theta, \theta+1, \theta+2, \ldots\}$ reveals the threshold only in the limit, and the class of all such languages together with all finite languages triggers Gold's impossibility theorem.

\item \textbf{Online-learnable but not $\Ex$-identifiable.} The class of singletons $\mathcal{C} = \{\{x\} : x \in X\}$ has Littlestone dimension~1 (online-learnable with at most 1 mistake). But the class of all singletons together with the empty set is not $\Ex$-identifiable from text for reasons analogous to Gold's theorem: the learner cannot distinguish ``no more positive examples will come'' from ``the next positive example has not yet arrived.''

\item \textbf{$\Ex$-identifiable but not online-learnable.} Consider any class $\mathcal{L}$ of recursive languages that is $\Ex$-identifiable and has infinite Littlestone dimension. Such classes exist: the class of all pattern languages (Angluin, 1980) is $\Ex$-identifiable from text but has infinite Littlestone dimension over suitable domains.
\end{enumerate}
\end{separation}

\begin{figure}[ht]
\centering
\begin{tikzpicture}[
    >={Stealth[length=2mm]},
    territory/.style={draw, thick, rounded corners=8pt, align=center, font=\small},
    wit/.style={font=\scriptsize, text=edgeorange, fill=white, inner sep=1.5pt, align=center},
    hdl/.style={font=\tiny, text=black!65, align=center},
    imp/.style={-{Stealth[length=2.2mm]}, very thick, defblue!60!black},
    sepv/.style={-{Stealth[length=2mm]}, thick, thmred},
    sepl/.style={-{Stealth[length=2mm]}, thick, thmred!85, densely dashed}
]
% --- PAC outer region; Online nested inside it; Gold a separate region ---
\node[territory, fill=defblue!8, minimum width=5.2cm, minimum height=4.0cm] (pac) at (0,0) {};
\node[font=\small\bfseries, defblue!75!black, anchor=north west, align=left]
   at ([xshift=5pt,yshift=-4pt]pac.north west) {PAC\\[-1pt]{\scriptsize\mdseries $\varepsilon,\delta$; VC dim}};
\node[territory, fill=onlineblue!14, minimum width=3.0cm, minimum height=1.7cm] (online) at (0.45,-0.7)
   {Online\\[-1pt]{\scriptsize\mdseries Littlestone dim}};
\node[territory, fill=goldcolor!12, minimum width=3.2cm, minimum height=4.0cm] (gold) at (7.6,0)
   {Gold / $\Ex$\\[-1pt]{\scriptsize\mdseries identification}\\{\scriptsize\mdseries in the limit}\\[3pt]{\scriptsize\mdseries ordinal /}\\{\scriptsize\mdseries mind-change}};
% --- E1: Online => PAC (verified containment), strict ---
\draw[imp] (0.45,0.15) -- (0.45,1.5);
\node[hdl, anchor=south, text=defblue!65!black] at (0.45,1.52) {\leanname{online_imp_pac}};
\node[hdl, anchor=north] at (0.45,-1.7)
   {strict: thresholds $\VC{=}1$, $\Ldim{=}\infty$\\{\tiny\leanname{pac_not_implies_online}}};
% --- E2: Gold =/=> PAC (verified separation, solid) ---
\draw[sepv] (6.0,1.3) -- (2.6,1.3);
\node[wit, fill=white] at (4.3,1.3) {finite sets: $\Ex$ yes, PAC no\\{\tiny\leanname{ex_not_implies_pac}}};
% --- E3,E4,E5: literature separations (dashed) ---
\draw[sepl] (2.6,0.45) -- (6.0,0.45);
\node[wit, fill=white] at (4.3,0.45) {thresholds: PAC yes, $\Ex$ no\\{\tiny\cite{Gold1967}}};
\draw[sepl] (1.95,-0.55) -- (6.0,-0.55);
\node[wit, fill=white] at (4.3,-0.55) {singletons$+\emptyset$: Online yes, $\Ex$ no\\{\tiny\cite{Gold1967}}};
\draw[sepl] (6.0,-1.4) -- (1.95,-1.4);
\node[wit, fill=white] at (4.3,-1.4) {patterns: $\Ex$ yes, Online no\\{\tiny\cite{Angluin1980}}};
\end{tikzpicture}
\caption[The three paradigms: online strictly inside PAC, Gold incomparable to both]{The relationship
among the three paradigms, corrected to the verified structure: it is a partial order, not an antichain.
Online learnability is strictly contained in PAC learnability: every online-learnable class is PAC-learnable
(\leanname{online_imp_pac}, solid), and the containment is strict because the threshold class over $\R$ is
PAC-learnable but has infinite Littlestone dimension, hence is not online-learnable
(\leanname{pac_not_implies_online}); so the Online region sits inside the PAC region. There is no class that
is online-learnable but not PAC-learnable, and none with finite Littlestone but infinite VC dimension, since
$\VC\le\Ldim$. Gold-style identification is incomparable to both: it identifies the finite languages, which
have infinite VC dimension and so escape every statistical learner (\leanname{ex_not_implies_pac}, solid),
while it cannot identify classes with no convergence signal that PAC and online learners handle (thresholds,
singletons with the empty set; dashed, \cite{Gold1967}), and pattern languages give the reverse (dashed,
\cite{Angluin1980}). Solid arrows are kernel theorems; dashed arrows are literature.}
\label{fig:three-paradigm}
\end{figure}

\begin{remark}[What the separations mean]
\label{rmk:separation-meaning}
The three-paradigm separation is not a deficiency of the definitions. It reflects a genuine mathematical fact: different notions of ``learning from data'' capture different aspects of the learning problem, and no single notion is universal. PAC learning measures statistical generalization under distributional assumptions. Online learning measures worst-case sequential prediction. Gold-style identification measures eventual convergence without any accuracy or efficiency guarantee. These are three different mathematical questions, and they have three different answers.
\end{remark}

% ================================================================
% SECTION 7: Chapter Summary
% ================================================================

\section{What This Chapter Established}
\label{sec:ch7-summary}

This chapter introduced Gold-style identification in the limit, the oldest paradigm in formal learning theory, and established five structural facts:

\begin{enumerate}[leftmargin=*,itemsep=3pt]
\item \textbf{Gold's impossibility theorem} (\Cref{thm:gold}): any class containing all finite languages and at least one infinite language is not $\Ex$-identifiable from text. The proof is by diagonalization, the first impossibility result in learning theory and one whose proof technique (adversarial stream construction) has no analogue in PAC or online theory.

\item \textbf{The identification hierarchy} (\Cref{sec:hierarchy}): the chain $\FIN \subsetneq \Ex \subsetneq \BC$ is strict, with the Case--Smith separation (\Cref{thm:bc-strictly-stronger}) providing the witness for $\BC > \Ex$. The hierarchy is not a sequence of definitions but a chain of theorems.

\item \textbf{Relaxations form a lattice} (\Cref{sec:relaxations}): anomalous, monotonic, and vacillatory learning branch off from $\Ex$, each modifying a different aspect of the convergence requirement. Trial-and-error learning connects to $\Delta^0_2$ in the arithmetical hierarchy.

\item \textbf{Mind-change complexity is ordinal-valued} (\Cref{sec:mind-change}): the Freivalds--Smith characterization shows that the right complexity measure for $\Ex$-identification uses transfinite ordinals, not integers. The full treatment is in \Cref{ch:ordinals}.

\item \textbf{The three paradigms form a partial order, not an antichain} (\Cref{sec:three-paradigm}): online learnability is strictly contained in PAC learnability (every online-learnable class is PAC-learnable, and the converse fails), while Gold-style identification is incomparable to both. Statistical capacity and convergence in the limit are orthogonal axes.
\end{enumerate}

The recursion-theoretic character of this chapter (diagonalization proofs, connections to the arithmetical hierarchy, ordinal-valued measures) contrasts sharply with the probabilistic character of \Cref{ch:pac} and the combinatorial character of \Cref{ch:online}. Each paradigm brings its own mathematical world. The separations of \Cref{sec:three-paradigm} are consequences of this fact: the paradigms use different proof techniques because they formalize genuinely different questions.


% ================================================================
% EXERCISES
% ================================================================

\subsection*{Exercises}

\begin{enumerate}[leftmargin=*,itemsep=6pt]

\item \textbf{The power of informant over text.}
  Gold showed that the class $\mathcal{L}$ of all recursive languages
  is $\Ex$-identifiable from \emph{informant} (both positive and
  negative examples) but not from text (positive examples only).
  \begin{enumerate}[label=(\alph*)]
  \item Prove the informant direction: construct an $\Ex$-learner $M$
    that identifies any recursive language $L$ from informant.
    \emph{Hint:}  Dovetail over all programs $\varphi_0, \varphi_1,
    \ldots$, and at time $t$ hypothesize the smallest index~$e$
    consistent with all labeled data seen so far.  Use the fact that
    the informant eventually presents every string with its correct
    label, so incorrect indices are eventually refuted.
  \item Explain precisely why this learner fails from text.  Identify
    the step in the argument that relies on negative examples, and
    show that no text-based substitute exists.  (\emph{Hint:} The
    learner can refute ``$\varphi_e$ includes $x$'' from an
    informant presentation of $(x, 0)$, but a text for $L$ never
    explicitly says ``$x \notin L$''; it merely fails to present
    $x$, which is indistinguishable from delay.)
  \end{enumerate}

\item \textbf{$\Ex$-identification of co-finite languages.}
  Let $\mathcal{L}_{\mathrm{cof}}$ be the class of all co-finite
  languages over $\N$: $L \in \mathcal{L}_{\mathrm{cof}}$ iff
  $\N \setminus L$ is finite.
  \begin{enumerate}[label=(\alph*)]
  \item Prove that $\mathcal{L}_{\mathrm{cof}}$ is
    $\Ex$-identifiable from text.  (\emph{Hint:} Every co-finite
    language $L$ contains all but finitely many natural numbers.
    A text for $L$ eventually presents every element of $L$, so after
    time $n_0$, every natural number $\leq n_0$ has either appeared
    in the text (and is in $L$) or has not appeared (and may or
    may not be in $L$).  Design a learner that waits long enough
    to conclude that unseen small numbers are \emph{not} in $L$.)
  \item Prove that $\mathcal{L}_{\mathrm{fin}} \cup
    \mathcal{L}_{\mathrm{cof}}$ (all finite and all co-finite
    languages together) is \emph{not} $\Ex$-identifiable from text.
    (\emph{Hint:} Apply Gold's impossibility theorem.  Verify
    that this class contains all finite languages and at least one
    infinite language.)
  \item The result of~(b) is sharper than Gold's theorem applied
    to ``all finite + one infinite'': the single infinite language
    is itself very structured (co-finite).  Explain why the structure
    of the infinite language does not help: what makes the
    diagonalization work is not the complexity of $L_\infty$ but
    the learner's inability to distinguish ``finite set, done
    growing'' from ``co-finite set, still growing.''
  \end{enumerate}

\end{enumerate}

\section{Open Problems}
\label{sec:ch7-open}

Each problem is anchored to a result above and tagged: a \emph{known unknown} (\textsc{ku}) is an articulated open question; an \emph{unknown known} (\textsc{uk}) is a result in the literature or the verified development not yet absorbed into a clean statement.  Verified handles are named where they exist.

\begin{openprob}[The identification criterion of in-context learning, \textsc{uk}]
\label{op:icl-identification}
\Cref{prop:transformer-gold} casts a transformer's in-context predictor as a Gold learner.  Determine which criterion it meets on natural concept classes: does a fixed transformer $\Ex$-converge (the in-context rule stabilizes), $\BC$-converge (it stabilizes only extensionally), or vacillate, and what is its mind-change ordinal (\leanname{FLT_Proofs.Complexity.MindChange.MindChangeOrdinal}) as a function of context length and model size?
\end{openprob}

\begin{openprob}[Formalizing the $\BC > \Ex$ separation, \textsc{uk}]
\label{op:bc-ex-formal}
Gold's impossibility theorem is verified (\leanname{FLT_Proofs.Theorem.Gold.gold_theorem}) and the criterion hierarchy is typed, but the Case--Smith strict separation (\cref{thm:bc-strictly-stronger}) is not.  Formalize the construction of a class that is $\BC$-identifiable but not $\Ex$-identifiable, on top of \leanname{FLT_Proofs.Criterion.Gold.BCLearnable} and \leanname{FLT_Proofs.Criterion.Gold.EXLearnable}.
\end{openprob}

\begin{openprob}[The Freivalds--Smith ordinal hierarchy, \textsc{uk}]
\label{op:fs-ordinal}
\Cref{thm:mind-change-ordinal} is stated but its strictness is cited.  Formalize the transfinite mind-change hierarchy on top of the typed ordinal measure (\leanname{FLT_Proofs.Complexity.MindChange.MindChangeOrdinal}): prove that for each constructive ordinal $\alpha$ the $\alpha$-bounded classes strictly contain the $\beta$-bounded classes for $\beta < \alpha$, connecting to the constructive ordinal notations of \Cref{ch:ordinals}.
\end{openprob}

\begin{openprob}[Angluin's telltale characterization, \textsc{ku}]
\label{op:telltale}
Gold's theorem is an impossibility result; the matching characterization is Angluin's: a class is $\Ex$-identifiable from text if and only if every language in it has a finite ``telltale'' subset that no proper sublanguage in the class contains~\cite{Angluin1980}.  Formalize this characterization on top of \leanname{FLT_Proofs.Criterion.Gold.EXLearnable}, giving the exact condition that separates the identifiable classes from those Gold's theorem rules out.
\end{openprob}

\begin{openprob}[The vacillatory hierarchy, \textsc{ku}]
\label{op:vex-hierarchy}
Vacillatory identification (\cref{def:vacillatory}) allows oscillation among at most $k$ correct indices.  Determine whether the $\mathrm{Vex}_k$ hierarchy is strict in $k$ and where $\mathrm{Vex} = \bigcup_k \mathrm{Vex}_k$ sits strictly between $\Ex$ and $\BC$, formalizing on top of \leanname{FLT_Proofs.Criterion.Gold.VacillatoryLearnable}.
\end{openprob}

\begin{openprob}[A meta-measure across the three paradigms, \textsc{ku}]
\label{op:meta-measure}
The three-paradigm structure (\cref{sec:three-paradigm}, \leanname{FLT_Proofs.Theorem.Separation.online_pac_gold_separation}) places online learnability strictly inside PAC and Gold-style identification incomparable to both.  Determine whether a single parameterized complexity measure specializes to the VC dimension, the Littlestone dimension, and the mind-change ordinal under three settings of a data-model parameter, or prove that the three invariants are provably non-unifiable.
\end{openprob}

\begin{openprob}[Anomalous in-context learning, \textsc{uk}]
\label{op:anomalous-icl}
Anomalous identification (\cref{def:anomalous}, \leanname{FLT_Proofs.Criterion.Gold.AnomalousLearnable}) permits finitely many errors in the final hypothesis.  Determine whether relaxing exactness to bounded anomaly enlarges what a transformer can identify in context, and relate the anomaly count to the approximation error of the attention learner on the tail of the data stream.
\end{openprob}

\begin{openprob}[Identification from informant versus text for attention, \textsc{ku}]
\label{op:informant-attention}
Exercise~1 shows informant is strictly more powerful than text.  Quantify the gap for a transformer: how much does access to negative examples (an informant-style prompt) reduce the mind-change ordinal of in-context identification relative to a positive-only (text-style) prompt, and is the separation as sharp for bounded-context transformers as it is in the recursive setting?
\end{openprob}

\begin{openprob}[Monotonic in-context learning, \textsc{uk}]
\label{op:monotonic-icl}
Monotonic identification (\cref{def:monotonic}, \leanname{FLT_Proofs.Criterion.Gold.MonotonicLearnable}) forbids retracting a correct classification.  Determine whether attention learners can be constrained to be monotonic without loss of identification power on natural classes, and whether monotonicity corresponds to a stability property of the routing under prefix extension.
\end{openprob}

\begin{openprob}[Effective three-paradigm witnesses, \textsc{uk}]
\label{op:effective-three-paradigm}
The three-paradigm separation produces its witnesses existentially.  Give a single parameterized family of concept classes that realizes each of the separations of \cref{sec:three-paradigm} (online strictly inside PAC, Gold incomparable to both) with explicit witnesses and explicit obstructions, turning the separation lattice into one constructive theorem rather than several independent constructions.
\end{openprob}


\subsection*{Counterfactual exercises: generalizing Figure~\ref{fig:identification-hierarchy}}

Each exercise perturbs the hierarchy. Show, in the figure's grammar (a strict-power partial order on identification criteria whose every edge is a separation theorem, refined by a relaxation lattice of single-constraint edits to the success notion, solid where a kernel theorem certifies the separation and dashed where it is borrowed from the literature or open), that the generalized picture subsumes the concrete one as an \emph{instance} (an edge whose separation is determined), a \emph{boundary} (where one constraint is toggled and the class strictly grows or shrinks), or a \emph{frontier} (a dashed, open or literature-only edge).

\begin{enumerate}[leftmargin=*,itemsep=3pt]
\item \textbf{The one verified rung.} Read the bottom edge $\FIN \to \Ex$ not as a containment of definitions but as a theorem with a witness class. Show this is the figure's only solid instance: the kernel proves $\FIN$-learnability implies $\Ex$-learnability by dropping the no-mind-change conjunct (\leanname{ex_strictly_stronger_finite}), and the class of all finite languages is the separating witness (\Cref{ex:finite-languages}), so this rung alone is certified rather than cited.

\item \textbf{Syntactic versus semantic convergence is the whole gap.} Replace the labels $\Ex$ and $\BC$ by their quantifier skeletons. Show this is the boundary that names the upper rung: $\Ex$ requires the index to satisfy $h = c$ (\leanname{EXLearnable}) while $\BC$ requires only $\forall x,\ h\,x = c\,x$ (\leanname{BCLearnable}), one clause toggled from equality of programs to equality of extensions, and that single toggle is exactly what the Case--Smith class exploits.

\item \textbf{Why the upper rung is dashed.} Ask the kernel for the $\Ex \subsetneq \BC$ separation. Show this is the figure's principal frontier: the criteria are typed (\leanname{EXLearnable}, \leanname{BCLearnable}) but no theorem separates them, so the edge is drawn dashed and cited (\cite{CaseSmith1983}, \cref{op:bc-ex-formal}); contrast this with the solid rung of Exercise~1 to see that line style is evidence grade, not decoration.

\item \textbf{Anomalous learning extends, not restricts.} Toggle the exactness clause of $\Ex$ from $h = c$ to ``$h$ errs on at most $a$ strings''. Show this is a boundary on the larger side: removing a constraint can only enlarge the learnable class, so $\Ex^*$ \rel{extends} $\Ex$ (\leanname{AnomalousLearnable}), the separation is strict but literature-only (\cite{CaseSmith1983}, \cref{op:anomalous-icl}), hence the $\Ex^* \to \Ex$ edge is dashed and oriented as an extension.

\item \textbf{Monotonic learning restricts.} Conjoin to $\Ex$ the no-retraction clause $h\,\mathrm{data}\,x = \mathrm{true} \Rightarrow h\,(\mathrm{data}\,{+}{+}\,[xy])\,x = \mathrm{true}$. Show this is the boundary on the smaller side: adding a constraint can only shrink the class, so $\mathrm{Mon}\Ex$ \rel{restricts} $\Ex$ (\leanname{MonotonicLearnable}); the strictness is literature (\cite{OshersonStobWeinstein1986}, \cref{op:monotonic-icl}), so the edge is dashed and oriented as a restriction, the mirror image of Exercise~4.

\item \textbf{Vacillation is two toggles at once.} Replace the single final index by membership in a finite set of correct indices. Show this is the instance that sits the branch between the spine nodes: relaxing $\Ex$'s one index to finitely many extends $\Ex$, while restricting $\BC$'s arbitrarily many indices to finitely many is the dual restriction of $\BC$, so $\mathrm{Vex}$ (\leanname{VacillatoryLearnable}) lies strictly between them, with both legs dashed and cited (\cite{CaseSmith1983}, \cref{op:vex-hierarchy}).

\item \textbf{The arithmetical anchor.} Read the $\Delta^0_2$ tag on the trial-and-error node as a characterization, not a containment. Show this is a frontier that re-types the value: trial-and-error learnability is pointwise convergence (\leanname{TrialAndErrorLearnable}) and equals membership in $\Delta^0_2$ (\leanname{Delta02Class}) by the Putnam--Gold theorem, but the kernel carries the class as a definition and not the iff, so the characterization edge is dashed and cited (\cite{Putnam1965}).

\item \textbf{Gold's theorem is the pressure behind the chain.} Ask why the hierarchy branches at $\Ex$ at all. Show this is the instance that motivates every edge: Gold's impossibility (\leanname{gold_theorem}) proves $\neg\,\Ex$-identifiability for any class with all finite languages and one infinite language, so strengthening to $\FIN$ and weakening to $\BC$ are the two escape routes the figure draws, and the chain exists because $\Ex$ is provably bounded.

\item \textbf{Mind changes turn the spine into an ordinal.} Refine the $\FIN$-to-$\Ex$ gap by counting retractions rather than forbidding them. Show this is a frontier that grades the bottom rung: $\FIN$ is zero mind changes and $\Ex$ is finitely many, and the kernel types the intermediate measure (\leanname{MindChangeOrdinal}), but the strictness of the transfinite hierarchy between them is cited, not proved (\cite{FreivaldsSmith1993}, \cref{op:fs-ordinal}), so the graded edge is dashed.

\item \textbf{Is the lattice one graded criterion?} Ask whether all seven nodes are evaluations of a single success criterion parameterized by four toggles: error budget, retraction flag, index cardinality, and convergence mode (syntactic or semantic). Show this is the figure's most general open frontier: it would present the hierarchy not as a fixed chain plus branches but as one criterion space with the four constraints as a basis, yet the kernel formalizes the seven criteria separately with no graded predicate unifying them, so the meta-edge is dashed over an otherwise typed lattice.
\end{enumerate}

\subsection*{Counterfactual exercises: generalizing Figure~\ref{fig:three-paradigm}}

Each exercise perturbs the corrected picture. Show, in the figure's grammar (a strict-power partial order on learning paradigms whose every edge is either an implication with a strict witness or a separation with a witness, solid where a kernel theorem certifies it and dashed where it is borrowed from the literature), that the generalized picture subsumes the concrete one as an \emph{instance} (an edge whose direction and evidence grade are determined), a \emph{boundary} (where one constraint is toggled and the learnable class strictly grows or shrinks), or a \emph{frontier} (a dashed, open or literature-only edge).

\begin{enumerate}[leftmargin=*,itemsep=3pt]
\item \textbf{The picture is an order, not an antichain.} Read the three paradigms not as three incomparable territories but as nodes of a partial order under ``strictly more powerful than''. Show this is the figure's organizing instance: online learnability implies PAC learnability (\leanname{online_imp_pac}) and the converse fails (\leanname{pac_not_implies_online}), so the Online and PAC nodes are \emph{comparable} and the whole picture is a containment plus an incomparable Gold, never a symmetric triangle of six arrows.

\item \textbf{The impossible witness.} Try to draw the edge ``some class is online-learnable but not PAC-learnable'', certified by ``finite Littlestone dimension, infinite VC dimension''. Show this is the frontier that does not exist: $\VC \le \Ldim$ (\leanname{BranchWiseLittlestoneDim_ge_VCDim}) makes the certificate ``$\Ldim<\infty$ and $\VC=\infty$'' an empty predicate, and \leanname{online_imp_pac} makes ``online but not PAC'' an empty class, so this edge is removed rather than dashed. A separation drawn with an impossible witness is not a weak edge; it is no edge.

\item \textbf{Strictness is the threshold class.} Read the strictness of $\text{Online}\subsetneq\mathrm{PAC}$ as a theorem with a witness, not a gap. Show this is the boundary on the larger side: the threshold class over $\N$ is PAC-learnable yet has $\Ldim=\infty$, hence is not online-learnable (\leanname{pac_not_implies_online}), so toggling from ``i.i.d.\ sample'' to ``adversarial order'' strictly shrinks the class, and this single witness is what makes the Online region sit \emph{properly} inside PAC.

\item \textbf{One solid arrow certifies two separations.} Ask the kernel only for $\Ex\not\Rightarrow\mathrm{PAC}$. Show this is the instance that does double duty: the finite-languages witness is $\Ex$-identifiable but has $\VC=\infty$ (\leanname{ex_not_implies_pac}), and because Online lies inside PAC, the same class is a fortiori not online-learnable, so the lone solid arrow $\Ex\not\Rightarrow\mathrm{PAC}$ simultaneously certifies $\Ex\not\Rightarrow\text{Online}$ and the cluster $\{\mathrm{PAC},\text{Online}\}$ needs no separate Gold-escapes-Online theorem.

\item \textbf{What the three-way theorem actually says.} Unfold \leanname{online_pac_gold_separation}. Show this is the frontier that names the figure's solid content exactly: it is the conjunction $(\exists\,\mathrm{PAC}\wedge\neg\,\text{Online})\wedge(\exists\,\Ex\wedge\neg\,\mathrm{PAC})$, that is \leanname{pac_not_implies_online} and \leanname{ex_not_implies_pac} and \emph{nothing else}, so any sixth ``mutual'' edge the eye expects from a triangle is not in the theorem and must be drawn dashed or not at all.

\item \textbf{Gold's incomparability is a different axis.} Ask why Gold escapes PAC ``upward'' (unbounded VC) yet is escaped by PAC ``downward'' (thresholds). Show this is the boundary that separates two axes: $\Ex\not\Rightarrow\mathrm{PAC}$ is a statistical-capacity fact certified in the kernel (\leanname{ex_not_implies_pac}), while $\mathrm{PAC}\not\Rightarrow\Ex$ is a convergence-signal fact that is literature only (\cite{Gold1967}), so the two legs at the Gold node measure capacity and convergence respectively, not one symmetric ``incomparability''.

\item \textbf{Why three of the legs are dashed.} Ask the kernel for $\mathrm{PAC}\not\Rightarrow\Ex$, $\text{Online}\not\Rightarrow\Ex$, and $\Ex\not\Rightarrow\text{Online}$. Show this is the figure's principal frontier: the paradigms are typed (\leanname{PACLearnable}, \leanname{OnlineLearnable}, \leanname{EXLearnable}) but no theorem in the kernel separates them in these directions, so all three edges are drawn dashed and cited (\cite{Gold1967} for the first two, \cite{Angluin1980} for the third); contrast the solid legs of Exercises~3 and~4 to read line style as evidence grade.

\item \textbf{Pattern languages are the reverse witness.} Toggle the Gold-versus-Online comparison to the other direction. Show this is the boundary that fills the $\Ex\not\Rightarrow\text{Online}$ leg: the class of pattern languages is $\Ex$-identifiable from positive data yet has infinite Littlestone dimension (\cite{Angluin1980}), so it is a class Gold learns that no online learner can, the mirror of the singletons-with-empty-set class that online learns and Gold cannot.

\item \textbf{Singletons sit inside PAC but outside Gold.} Locate the singletons-with-empty-set class on the corrected picture. Show this is the instance that tests the geometry: the class has $\Ldim=1$ so it is online-learnable, hence (by \leanname{online_imp_pac}) it lies in both the Online and PAC regions, yet it is not $\Ex$-identifiable from text (\cite{Gold1967}), so it is a single class witnessing membership in the nested $\{\text{Online}\subset\mathrm{PAC}\}$ region together with exclusion from the Gold region.

\item \textbf{Is the order one graded measure?} Ask whether the VC dimension, the Littlestone dimension, and the mind-change ordinal are three readings of a single parameterized complexity measure under one data-model parameter. Show this is the figure's most general open frontier: the kernel proves the comparabilities it can ($\VC\le\Ldim$ via \leanname{BranchWiseLittlestoneDim_ge_VCDim}, $\text{Online}\Rightarrow\mathrm{PAC}$ via \leanname{online_imp_pac}) but carries the three invariants as separate definitions with no graded predicate unifying them (\cref{op:meta-measure}, \cref{op:effective-three-paradigm}), so the meta-edge collapsing the order into one measure is dashed over an otherwise typed lattice.
\end{enumerate}

% Chapter 8: Exact Learning and Query Models
% Centerpiece: Angluin's L* algorithm and the passive-active gap.
% DFAs are NOT PAC-learnable (Kearns-Valiant 1994) but ARE exactly
% learnable with MQ+EQ (Angluin 1987). This separation is the content.

\chapter{Exact Learning and Query Models}
\label{ch:exact}

The previous chapters study learners that receive data passively: drawn from a distribution (PAC), or presented by an adversary (online), or enumerated by nature (Gold). This chapter studies what happens when the learner can \emph{ask questions}.

The exact learning model, introduced by Angluin~\cite{Angluin1987}, equips the learner with two oracles: a membership query oracle and an equivalence query oracle. The learner's goal is not approximate generalization but \emph{exact identification}: the game ends when the learner proposes a hypothesis that the equivalence oracle accepts. The central result of this chapter is that DFAs, computationally intractable under passive data~\cite{KearnsValiant1994}, become efficiently learnable with query access. This gap between passive and active learning is one of the sharpest structural results in the field.

% ================================================================
% SECTION 1: THE EXACT LEARNING FRAMEWORK
% ================================================================

\section{The Exact Learning Framework}
\label{sec:exact-framework}

The query model was introduced in \Cref{sec:queries}, where membership and equivalence oracles were defined (\Cref{def:mq,def:eq}). We recall the framework here with the precision needed for algorithmic analysis.

\begin{defn}[Minimally Adequate Teacher]
\label{def:mat}
A \emph{minimally adequate teacher} (MAT) for a target concept $c$ over domain $X$ provides two oracles:
\begin{itemize}[leftmargin=*, itemsep=2pt]
\item $\mathrm{MQ}(x)$: returns $c(x)$ for any $x \in X$ chosen by the learner.\formal{FLT_Proofs.Data.MembershipOracle}
\item $\mathrm{EQ}(h)$: if $h = c$, returns \textsc{yes}. Otherwise, returns a counterexample $z \in X$ with $h(z) \neq c(z)$.\formal{FLT_Proofs.Data.EquivalenceOracle}
\end{itemize}
The teacher is ``minimally adequate'' in the sense that it provides exactly these two capabilities and nothing more: no distributional information, no structural hints, no preference among counterexamples.\formal{FLT_Proofs.Learner.Active.MinimallyAdequateTeacher}
\end{defn}

\begin{defn}[Exact Learning]
\label{def:exact-learning-formal}
A concept class $\mathcal{C}$ is \emph{exactly learnable in polynomial time} if there exists an algorithm $A$ such that for every target $c \in \mathcal{C}$, algorithm $A$ interacts with a MAT for $c$ and eventually receives \textsc{yes} from $\mathrm{EQ}$, using a number of queries and computation steps polynomial in:
\begin{enumerate}[leftmargin=*, nosep]
\item the representation size $n$ of the target $c$, and
\item the length $m$ of the longest counterexample returned by $\mathrm{EQ}$.
\end{enumerate}\formal{FLT_Proofs.Criterion.PAC.ExactLearnable}
\end{defn}

\begin{remark}[The role of counterexample length]
\label{rmk:counterexample-length}
The dependence on $m$ is necessary because the equivalence oracle's counterexamples are adversarial: the oracle may return arbitrarily long counterexamples. The algorithm must process them, so the complexity must account for their length.
\end{remark}

% ================================================================
% SECTION 2: THE L* ALGORITHM
% ================================================================

\section{Angluin's $L^*$ Algorithm}
\label{sec:lstar}

The $L^*$ algorithm learns the minimal DFA for an unknown regular language $L \subseteq \Sigma^*$ using membership and equivalence queries. It is the canonical example of exact learning and the centerpiece of this chapter.

\subsection{Observation Tables}

The algorithm maintains an \emph{observation table} $(S, E, T)$, where:
\begin{itemize}[leftmargin=*, itemsep=2pt]
\item $S \subseteq \Sigma^*$ is a finite, prefix-closed set of \emph{access strings} (candidate state representatives).
\item $E \subseteq \Sigma^*$ is a finite, suffix-closed set of \emph{distinguishing extensions}.
\item $T : (S \cup S \cdot \Sigma) \times E \to \{0,1\}$ is the table, where $T(s, e) = \mathrm{MQ}(s \cdot e)$.
\end{itemize}
For each string $s \in S \cup S \cdot \Sigma$, define the \emph{row} $\mathrm{row}(s) = (T(s, e_1), T(s, e_2), \ldots, T(s, e_{|E|}))$.

\begin{defn}[Closed and Consistent Table]
\label{def:closed-consistent}
The observation table $(S, E, T)$ is:
\begin{itemize}[leftmargin=*, itemsep=2pt]
\item \emph{Closed} if for every $s \in S$ and $a \in \Sigma$, there exists $s' \in S$ with $\mathrm{row}(s \cdot a) = \mathrm{row}(s')$.
\item \emph{Consistent} if for all $s_1, s_2 \in S$ with $\mathrm{row}(s_1) = \mathrm{row}(s_2)$, we have $\mathrm{row}(s_1 \cdot a) = \mathrm{row}(s_2 \cdot a)$ for all $a \in \Sigma$.
\end{itemize}
\end{defn}

Closedness ensures that every one-step extension of a state representative is equivalent to some existing representative, so the transition function of the conjectured DFA is well-defined. Consistency ensures that equivalent representatives behave identically under extension, so the transition function is well-defined as a function of equivalence classes, not of particular strings.

\subsection{The Algorithm}

\begin{figure}[ht]
\centering
\begin{tikzpicture}[
    >={Stealth[length=2.2mm]},
    pole/.style={draw, very thick, rounded corners, align=center, font=\small\bfseries,
                 minimum width=2.3cm, minimum height=1.3cm},
    lstep/.style={draw, semithick, rounded corners, align=center, font=\scriptsize,
                 fill=sepgreen!10, draw=sepgreen!55!black, minimum width=3.0cm,
                 minimum height=0.9cm, inner sep=3pt},
    cyc/.style={-{Stealth[length=2mm]}, semithick, defblue!60!black},
    q/.style={-{Stealth[length=2mm]}, semithick, thmred!70!black},
    lit/.style={draw, densely dashed, thmred!70!black, rounded corners, fill=thmred!4,
                align=left, font=\scriptsize, text=black!72, inner sep=4pt}
]
% --- agents (poles) ---
\node[pole, fill=defblue!10, draw=defblue!60!black] (learner) at (-0.3,0)
   {Learner\\[-1pt]{\scriptsize\mdseries table $(S,E,T)$}};
\node[pole, fill=thmred!10, draw=thmred!60!black, minimum height=1.6cm] (mat) at (11,0)
   {MAT\\[-1pt]{\scriptsize\mdseries knows $A^{\ast}$}\\[-1pt]{\tiny\mdseries\leanname{MinimallyAdequateTeacher}}};
% --- the L* state-discovery loop (verified data-type nodes) ---
\node[lstep] (n1) at (3.3,2.1) {1.~Fill table via $\mathrm{MQ}$\\{\tiny\leanname{MembershipOracle}}};
\node[lstep] (n2) at (3.3,0) {2.~Close \& make consistent\\{\tiny\cref{def:closed-consistent}}};
\node[lstep] (n3) at (3.3,-2.1) {3.~Conjecture DFA $M$\\{\tiny\leanname{DFA'}}};
\node[lstep] (n4) at (7.7,-2.1) {4.~$\mathrm{EQ}(M)$\\{\tiny\leanname{EquivalenceOracle}}};
\node[lstep] (n5) at (7.7,2.1) {5.~counterexample $z$\\{\tiny\leanname{Counterexample}}};
% --- cycle edges ---
\draw[cyc] (n1) -- (n2);  \draw[cyc] (n2) -- (n3);  \draw[cyc] (n3) -- (n4);
\draw[cyc] (n4) -- (n5);  \draw[cyc] (n5) -- (n1);
\draw[cyc, densely dashed] (n2.west) to[out=158,in=202]
   node[font=\tiny, left, align=center]{not closed/\\consistent} (n1.west);
% --- query interaction with the MAT ---
\draw[q] (learner.north) to[out=65,in=115] node[font=\scriptsize, above]{membership queries $\mathrm{MQ}$} (mat.north);
\draw[q] (n4.east) to[out=0,in=205] node[font=\tiny, above right, pos=0.6]{$\mathrm{EQ}$} (mat.south west);
\draw[q, densely dashed] (mat.west) to[out=155,in=0] node[font=\tiny, above, pos=0.45]{$z$ / \textsc{yes}} (n5.east);
% --- the criterion (Will) ---
\node[font=\tiny, text=defblue!55!black, align=center, anchor=north] at (3.3,-2.85)
   {halt iff $\mathrm{EQ}=$\textsc{yes} (\leanname{ExactLearnable})};
% --- literature: correctness + termination (Angluin 1987, not kernel) ---
\node[lit, text width=3.5cm, anchor=north] at (7.9,-2.9)
   {\textbf{Correctness \& the $n-1$ bound:} converges to the minimal DFA in $\le n-1$ equivalence queries ($n=$ Myhill--Nerode index). Angluin 1987 (\cref{thm:lstar-complexity}); open in the kernel (\cref{op:lstar-formal}).};
\end{tikzpicture}
\caption[The $L^*$ exact-learning loop and its Myhill--Nerode reading]{The $L^*$ protocol as the
exact-learning-via-a-MAT loop. A learner and a minimally adequate teacher cooperate to identify the
minimal DFA for an unknown regular language. The observation table $(S,E,T)$ is a finite approximation of
the Myhill--Nerode right-congruence of the target: its distinct rows are the equivalence classes
discovered so far. Membership queries fill the table cell by cell; an equivalence query returns a
counterexample that forces a new state class, splitting the current approximation. Closedness and
consistency are exactly the conditions under which the table induces a well-defined DFA, and termination
in at most $n-1$ equivalence queries is the index of the congruence ($n$ states in the minimal DFA). The
data-types are verified scaffolding: the membership and equivalence oracles
(\leanname{MembershipOracle}, \leanname{EquivalenceOracle}), the teacher
(\leanname{MinimallyAdequateTeacher}), the algorithm object (\leanname{LStar}), the conjecture type
(\leanname{DFA'}), and the success criterion (\leanname{ExactLearnable}). The correctness of the loop and
the $n-1$ bound are Angluin's theorem (\cite{Angluin1987}), cited and not yet kernel-proved
(\cref{op:lstar-formal}); they are drawn dashed.}
\label{fig:lstar-protocol}
\end{figure}

The $L^*$ algorithm proceeds as follows.

\begin{enumerate}[leftmargin=*, itemsep=4pt]
\item \textbf{Initialize.} Set $S = E = \{\varepsilon\}$ (the empty string). Fill the table $T$ using membership queries.

\item \textbf{Close and make consistent.} Repeat until the table is both closed and consistent:
\begin{itemize}[nosep]
\item If the table is not closed (there exists $s \in S$, $a \in \Sigma$ with $\mathrm{row}(s \cdot a) \neq \mathrm{row}(s')$ for any $s' \in S$), add $s \cdot a$ to $S$ and fill new table entries via $\mathrm{MQ}$.
\item If the table is not consistent (there exist $s_1, s_2 \in S$ with $\mathrm{row}(s_1) = \mathrm{row}(s_2)$ but $\mathrm{row}(s_1 \cdot a) \neq \mathrm{row}(s_2 \cdot a)$ for some $a \in \Sigma$), then there exists $e \in E$ witnessing the difference. Add $a \cdot e$ to $E$ and fill new entries via $\mathrm{MQ}$.
\end{itemize}

\item \textbf{Conjecture.} When the table is closed and consistent, construct a DFA $M$:
\begin{itemize}[nosep]
\item States: the distinct rows $\mathrm{row}(s)$ for $s \in S$.
\item Initial state: $\mathrm{row}(\varepsilon)$.
\item Transition: $\delta(\mathrm{row}(s), a) = \mathrm{row}(s \cdot a)$ (well-defined by closedness and consistency).
\item Accepting states: those $\mathrm{row}(s)$ with $T(s, \varepsilon) = 1$.
\end{itemize}

\item \textbf{Query.} Submit $\mathrm{EQ}(M)$.
\begin{itemize}[nosep]
\item If the teacher returns \textsc{yes}, halt and output $M$.
\item If the teacher returns a counterexample $z$, add $z$ and all its prefixes to $S$, fill the table, and return to step~2.
\end{itemize}
\end{enumerate}

\begin{thm}[$L^*$ correctness and complexity {\cite{Angluin1987}}]
\label{thm:lstar-complexity}
Let $A^*$ be the minimal DFA for the target language $L$, with $n$ states over alphabet $\Sigma$. The $L^*$ algorithm terminates and outputs the minimal DFA for $L$, using:
\begin{itemize}[nosep]
\item at most $n - 1$ equivalence queries, and
\item at most $O(n^2 |\Sigma| \cdot m)$ membership queries,
\end{itemize}
where $m$ is the length of the longest counterexample received.
\end{thm}

\begin{proof}[Proof sketch]
Each counterexample introduces at least one new distinct row in the observation table, because the counterexample witnesses a string that the current conjecture misclassifies. Since the minimal DFA has $n$ states and the table's distinct rows correspond to Myhill--Nerode equivalence classes, the number of distinct rows is bounded by $n$. The table starts with one row and gains at least one per counterexample, so at most $n - 1$ equivalence queries occur before the table has $n$ distinct rows, at which point the conjecture DFA is isomorphic to $A^*$.

The membership query bound follows from the table dimensions: $|S| \leq n$, $|S \cdot \Sigma| \leq n|\Sigma|$, and $|E| \leq m \cdot n$ (each counterexample adds at most $m$ suffixes). The total number of cells is $O(n|\Sigma|) \times O(nm) = O(n^2 |\Sigma| m)$, each requiring one membership query.

The full correctness proof (that a closed, consistent table yields a DFA consistent with all membership queries, and that the Myhill--Nerode theorem guarantees uniqueness) is given in Angluin~\cite{Angluin1987}.
\end{proof}

\begin{remark}[Why the table works]
\label{rmk:table-intuition}
The observation table is a finite approximation to the Myhill--Nerode equivalence relation of the target language. Two strings $s_1, s_2$ are Myhill--Nerode equivalent if $s_1 \cdot w \in L \iff s_2 \cdot w \in L$ for all $w \in \Sigma^*$. The table approximates this by testing only the extensions in $E$. Closedness ensures the approximation is fine enough to define transitions; consistency ensures it is coherent. Each counterexample reveals that the current approximation conflates two genuinely distinct equivalence classes, forcing a refinement.
\end{remark}

% ================================================================
% SECTION 3: THE PASSIVE-ACTIVE GAP
% ================================================================

\section{The Passive--Active Gap}
\label{sec:passive-active-gap}

The structural significance of $L^*$ lies not in the algorithm itself but in what it implies when combined with a hardness result.

\begin{thm}[DFAs are not PAC-learnable {\cite{KearnsValiant1994}}]
\label{thm:dfa-pac-hard}
Under standard cryptographic assumptions (specifically, the hardness of inverting RSA or factoring Blum integers), the class of DFAs with $n$ states is not efficiently PAC-learnable from random examples alone. No polynomial-time algorithm can PAC-learn DFAs, even improperly.
\end{thm}

The proof reduces the problem of breaking a cryptographic primitive to the problem of PAC-learning DFAs: if a polynomial-time PAC learner existed, it could be used to invert the cryptographic function, contradicting the hardness assumption. The reduction is non-trivial and we defer to Kearns and Valiant~\cite{KearnsValiant1994} for the full argument.

\begin{separation}
\textbf{Passive $\not\Leftrightarrow$ Active.} Witness: the class of DFAs.
\begin{itemize}[leftmargin=*, nosep]
\item \textbf{Active learning succeeds.} $L^*$ exactly identifies any $n$-state DFA using $O(n^2 |\Sigma| m)$ membership queries and $n - 1$ equivalence queries (\Cref{thm:lstar-complexity}).
\item \textbf{Passive learning fails.} Under cryptographic assumptions, no polynomial-time algorithm PAC-learns DFAs from random examples (\Cref{thm:dfa-pac-hard}).
\end{itemize}
The gap is not quantitative (more data needed) but \emph{qualitative} (no amount of passive data suffices for efficient learning, while a polynomial number of queries does). The mechanism is clear: membership queries let the learner probe the target's transition structure at chosen points, and equivalence queries provide directed feedback at the learner's current frontier of error. Random examples provide neither capability.
\end{separation}

This separation witnesses a fundamental edge in the concept graph: $\nde{pac\_learning} \not\Rightarrow \nde{exact\_learning}$ and $\nde{exact\_learning} \not\Rightarrow \nde{pac\_learning}$. Neither paradigm subsumes the other. The reverse direction (classes that are PAC-learnable but not efficiently exactly learnable) also holds, because simulating an equivalence oracle in the PAC setting requires enumerating hypotheses and checking consistency, which may be computationally intractable for some classes.

\begin{remark}[The role of the cryptographic assumption]
\label{rmk:crypto-assumption}
The Kearns--Valiant result is \emph{conditional}: it assumes that certain cryptographic primitives are hard to break. Without this assumption, the PAC-hardness of DFAs is open. This is typical of computational hardness results in learning theory: unconditional lower bounds for PAC learning remain rare, and most known hardness results reduce from cryptography or from $\mathsf{NP}$-hardness assumptions. The exact learning model sidesteps computational hardness entirely; query access changes the information-theoretic landscape so fundamentally that the cryptographic barrier does not apply.
\end{remark}

\begin{figure}[ht]
\centering
\begin{tikzpicture}[
    box/.style={draw, thick, rounded corners, minimum width=2.6cm, minimum height=0.9cm, font=\small\bfseries, align=center},
    arr/.style={-{Stealth[length=2mm]}, thick},
    noarr/.style={thick, thmred, dashed},
    note/.style={font=\scriptsize, align=center, text=notegray}
]
\node[box, fill=defblue!10] (exact) {Exact Learning\\{\scriptsize (MQ + EQ)}};
\node[box, fill=edgeorange!10, right=4cm of exact] (pac) {PAC Learning\\{\scriptsize (random examples)}};
\node[box, fill=sepgreen!10, below=2.5cm of $(exact)!0.5!(pac)$] (dfa) {DFAs};

% Arrows
\draw[noarr] (pac) -- node[above, font=\scriptsize, text=thmred] {$\not\Rightarrow$} (exact);
\draw[noarr] (exact) -- node[above, font=\scriptsize, text=thmred] {$\not\Rightarrow$} (pac);

\draw[arr, sepgreen!60!black] (dfa) -- node[left, note, text width=2.8cm] {Learnable:\\$L^*$ in poly queries\\{\tiny (Angluin 1987)}} (exact);
\draw[arr, thmred!60!black, dashed] (dfa) -- node[right, note, text width=2.8cm] {Not learnable:\\crypto hardness\\{\tiny (Kearns--Valiant 1994)}} (pac);
\end{tikzpicture}
\caption{The passive--active separation. DFAs are exactly learnable with query access but not PAC-learnable under cryptographic assumptions. Neither paradigm implies the other.}
\label{fig:passive-active-gap}
\end{figure}

% ================================================================
% SECTION 4: QUERY COMPLEXITY AND FORWARD REFERENCES
% ================================================================

\subsection{Extracting an automaton from a transformer}
\label{sec:dfa-extraction}

The passive--active gap has a direct modern use: it is the principle behind \emph{automaton extraction} from trained neural sequence models.  A trained attention model computes a binary string function (accept/reject, or a thresholded next-token decision), and we may treat that function as the unknown target of an exact-learning game, with the model itself playing the minimally adequate teacher.

\begin{prop}[A trained transformer is an $L^*$ target]
\label{prop:transformer-extraction}
Let $f$ be the binary string function computed by a trained attention model.  Answer membership queries by a forward pass, $\mathrm{MQ}(s) = f(s)$, and equivalence queries by sampling-based testing of a conjectured DFA against $f$.  Running $L^*$ against this teacher (\cref{thm:lstar-complexity}) yields a DFA consistent with $f$ on the queried strings, and the minimal DFA for $f$ if the equivalence oracle is exact.\formal{FLT_Proofs.Computation.DFA'} The teacher is a minimally adequate teacher in the sense of \cref{def:mat}, with the model's forward pass as the membership oracle.\formal{FLT_Proofs.Learner.Active.MinimallyAdequateTeacher}
\end{prop}

\begin{proof}
The membership oracle $\mathrm{MQ}(s) = f(s)$ is total and truthful: a forward pass computes $f(s)$ exactly up to the bounded finite-precision forward error of the attention layer.\formal{TLT.Fp32AttnLit.attnLiteralForwardError} The equivalence oracle is approximate: it certifies $L(M) = f$ on a sampled test set rather than on all strings, so the conjecture is exact on the queries made and probably-approximately correct off them; with an exact equivalence oracle $L^*$ returns the minimal DFA for $f$ by \cref{thm:lstar-complexity}.  Either way the model is a legitimate $L^*$ teacher.
\end{proof}

This is how Weiss, Goldberg, and Yahav~\cite{WeissGoldbergYahav2018} extract interpretable automata from recurrent networks, and the construction applies to attention models verbatim.  It turns the passive--active gap into a methodology: a model opaque under passive inspection becomes a queryable teacher, and $L^*$ distills its regular core into a DFA one can read.  The caveat is the one above: the equivalence oracle is only approximate, so extraction is sound on the queries and heuristic beyond them, and whether a transformer's string function even \emph{has} a finite regular core is itself open (\cref{op:transformer-regular}).

\section{Query Complexity}
\label{sec:query-complexity}

The $L^*$ algorithm's complexity is measured in \emph{query complexity}: the number of membership and equivalence queries as a function of the target's representation size. This measure is the exact learning analogue of sample complexity in PAC learning.

\begin{defn}[Query Complexity]
\label{def:query-complexity}
The \emph{membership query complexity} $\mathrm{MQ}(\mathcal{C}, n)$ and \emph{equivalence query complexity} $\mathrm{EQ}(\mathcal{C}, n)$ of a class $\mathcal{C}$ are the minimum numbers of membership and equivalence queries, respectively, that any exact learning algorithm must make in the worst case to identify any target $c \in \mathcal{C}$ of representation size $n$.
\end{defn}

For DFAs, $L^*$ achieves $\mathrm{EQ} \leq n - 1$ and $\mathrm{MQ} = O(n^2 |\Sigma| m)$. The equivalence query count has a direct connection to the combinatorial structure of the hypothesis class.

\begin{prop}[Equivalence queries alone are expensive {\cite{Angluin1988}}]
\label{prop:eq-lower}
With equivalence queries alone (no membership queries), exactly learning a finite class of $N$ concepts requires at most $N - 1$ queries always, and $N - 1$ in the worst case. Membership queries break this bound: $L^*$ uses only $n - 1$ equivalence queries for an $n$-state DFA, although $|\mathcal{C}_n|$ is exponential in $n$.
\end{prop}

\begin{proof}
Maintain the version space $V$ of candidates consistent with every counterexample received, initially the whole class.  Each query $\mathrm{EQ}(h)$ either returns \textsc{yes} (and the learner halts) or a counterexample that is inconsistent with $h$, so it removes at least $h$ from $V$; hence $|V|$ drops by at least one per query and $N - 1$ queries suffice to reach $|V| = 1$.  For the matching lower bound, take a class in which each concept differs from every other at a point unique to that pair, so that for any conjecture $h$ the adversary can return a counterexample consistent with all surviving candidates except $h$; then exactly one candidate is eliminated per query and $N - 1$ are forced.  The contrast with $L^*$ is the content of the chapter: membership queries supply the structural information (which Myhill--Nerode class a string belongs to) that equivalence queries alone cannot, collapsing an exponential query count to a linear one (\cref{prop:eq-alone} of \Cref{ch:data}).
\end{proof}

\begin{remark}[Certificate complexity]
\label{rmk:certificate}
The exact learning framework connects to \emph{certificate complexity} in Boolean function theory: the minimum number of input bits that must be queried to certify a function's value. Teaching dimension (\Cref{ch:dimensions}), which measures the minimum number of examples a teacher must provide to uniquely identify a concept, provides a lower bound on the number of equivalence queries. The connections among query complexity, teaching dimension, and certificate complexity are explored in \Cref{ch:dimensions,ch:computational}.
\end{remark}

\section{What This Chapter Established}
\label{sec:ch8-summary}

\begin{enumerate}[leftmargin=*, itemsep=3pt]
\item \textbf{The exact learning framework.} A minimally adequate teacher provides membership and equivalence queries. Exact learning requires identifying the target concept precisely, not approximately.

\item \textbf{The $L^*$ algorithm.} Angluin's algorithm learns the minimal DFA for any regular language using polynomially many queries. The observation table approximates the Myhill--Nerode equivalence relation; each counterexample refines it by introducing a new state distinction.

\item \textbf{The passive--active gap.} DFAs are exactly learnable with MQ~+~EQ (Angluin, 1987) but not PAC-learnable under cryptographic assumptions (Kearns--Valiant, 1994). This separation is qualitative, not quantitative: query access changes the complexity landscape fundamentally.

\item \textbf{Query complexity.} The number of queries required to exactly learn a class is the natural complexity measure, analogous to sample complexity in PAC learning.
\end{enumerate}

The gap between passive and active learning demonstrated by DFAs is not an isolated curiosity. It reflects a structural principle: the information geometry of a learning problem can change qualitatively when the learner moves from passive observation to active interrogation. This principle recurs in statistical experimental design, active learning in the PAC setting, and reinforcement learning: all settings where the learner's ability to choose which data to collect transforms the computational landscape.

\section{Open Problems}
\label{sec:ch8-open}

Each problem is anchored to a result above and tagged: a \emph{known unknown} (\textsc{ku}) is an articulated open question; an \emph{unknown known} (\textsc{uk}) is a result in the literature or the verified development not yet absorbed into a clean statement.  Verified handles are named where they exist.

\begin{openprob}[The regular core of a transformer, \textsc{uk}]
\label{op:transformer-regular}
\Cref{prop:transformer-extraction} extracts a DFA from a trained attention model, exact only on the queried strings.  Determine when a transformer's binary string function has a finite Myhill--Nerode index (a genuine regular core) so that $L^*$ extraction is exact rather than heuristic, and bound the number of states as a function of the architecture (depth, heads, context length).
\end{openprob}

\begin{openprob}[An exact equivalence oracle for neural models, \textsc{ku}]
\label{op:exact-eq-oracle}
Extraction is approximate because the equivalence oracle samples rather than verifies.  Determine for which model classes an \emph{exact} equivalence query (``does this DFA agree with the model on all strings?'') is decidable, turning approximate extraction into exact $L^*$ learning, and relate the cost to formal verification of the attention layer.
\end{openprob}

\begin{openprob}[Formalizing $L^*$, \textsc{uk}]
\label{op:lstar-formal}
The $L^*$ correctness and complexity theorem (\cref{thm:lstar-complexity}) is cited, not verified.  Formalize $L^*$ (the observation table, closedness, consistency, and the Myhill--Nerode termination argument) on top of the typed minimally adequate teacher (\leanname{FLT_Proofs.Learner.Active.MinimallyAdequateTeacher}) and \leanname{FLT_Proofs.Computation.DFA'}, with the $n-1$ equivalence-query bound proved.
\end{openprob}

\begin{openprob}[A combinatorial dimension for MAT-learnability, \textsc{ku}]
\label{op:mat-dimension}
Which classes are exactly learnable from a minimally adequate teacher in polynomial time?  Beyond DFAs (decision trees, certain Boolean formulas, residual automata are known), give a combinatorial dimension that characterizes polynomial-query MAT-learnability the way $\VC$ characterizes PAC learnability, or prove the query model admits none.
\end{openprob}

\begin{openprob}[The reverse separation, \textsc{ku}]
\label{op:reverse-separation}
The passive--active gap (\cref{sec:passive-active-gap}) is verified in one direction (DFAs: exact yes, PAC no).  Make the reverse precise: exhibit a class that is efficiently PAC-learnable but not efficiently exactly learnable, pinning down why simulating an equivalence oracle in the PAC setting is computationally intractable for it.
\end{openprob}

\begin{openprob}[Sharp equivalence-query complexity, \textsc{ku}]
\label{op:eq-sharp}
\Cref{prop:eq-lower} gives the worst-case $N-1$ bound for equivalence queries alone.  Characterize the exact equivalence-query complexity of a class by a combinatorial parameter (a ``certificate'' or ``extended teaching'' dimension), interpolating between the $N-1$ of unstructured classes and the $n-1$ of DFAs under added membership queries.
\end{openprob}

\begin{openprob}[Query complexity versus teaching dimension, \textsc{uk}]
\label{op:query-teaching}
\Cref{rmk:certificate} connects equivalence queries to teaching dimension and certificate complexity.  Make the connection quantitative: prove two-sided bounds relating $\mathrm{EQ}(\mathcal{C}, n)$ to the teaching dimension (\leanname{FLT_Proofs.Learner.Active.IsOptimalTeacher}) and the certificate complexity of $\mathcal{C}$, unifying the active-learning and teaching pictures of \Cref{ch:dimensions}.
\end{openprob}

\begin{openprob}[Test complexity of attention extraction, \textsc{uk}]
\label{op:extraction-samples}
The approximate equivalence oracle of \cref{prop:transformer-extraction} samples strings to test a conjecture.  Determine the number of test samples needed to certify $L(M) = f$ with confidence $1 - \delta$ as a function of the DFA size and the model's margin, giving a PAC guarantee on the extracted automaton.
\end{openprob}

\begin{openprob}[Counterexample length, \textsc{ku}]
\label{op:counterexample-length}
The $L^*$ membership-query count carries a factor $m$, the longest counterexample length (\cref{rmk:counterexample-length}).  Determine whether a learner or a cooperative teacher can avoid the linear dependence on $m$, for instance by binary-searching counterexamples (Rivest--Schapire), and what the optimal dependence is.
\end{openprob}

\begin{openprob}[An LLM as a minimally adequate teacher, \textsc{uk}]
\label{op:llm-teacher}
The development types an LLM critic as a teacher (\leanname{FLT_Proofs.Learner.Active.LLMCritic}).  Determine the soundness conditions under which a large language model can serve as the membership and equivalence oracle for a downstream exact learner (when its answers are reliable enough that $L^*$ converges to the intended concept), bounding the error introduced by an unreliable teacher.
\end{openprob}


\subsection*{Counterfactual exercises: generalizing Figure~\ref{fig:lstar-protocol}}

Each exercise perturbs the protocol. Show, in the figure's grammar (a cyclic state-discovery loop in which a learner refines a finite approximation of the target's Myhill--Nerode right-congruence by membership queries and is corrected by counterexamples from an equivalence oracle, with the flowchart's nodes drawn solid where a verified data-type certifies the object and the correctness-and-termination claims drawn dashed where they are Angluin's theorem, cited not proved), that the generalized picture subsumes the concrete one as an \emph{instance} (the loop run on a determined target), a \emph{boundary} (where one capability of the teacher or one condition on the table is toggled and the learnable class strictly grows or shrinks), or a \emph{frontier} (a dashed, open or literature-only edge).

\begin{enumerate}[leftmargin=*,itemsep=3pt]
\item \textbf{The table is the right congruence.} Read the observation table not as scratch space but as the object being learned. Show this is the figure's defining instance: two access strings share a row exactly when they are indistinguishable by the suffixes in $E$, a finite slice of the Myhill--Nerode relation $s_1 \sim s_2 \iff (\forall w,\ s_1 w \in L \Leftrightarrow s_2 w \in L)$ (\cref{rmk:table-intuition}), so the distinct rows \emph{are} the equivalence classes discovered so far and the conjecture DFA is their quotient; the row map is typed by the truthful membership oracle (\leanname{FLT_Proofs.Data.MembershipOracle}).

\item \textbf{Closedness and consistency are the quotient's well-definedness.} Toggle off closedness: let some $\mathrm{row}(s\cdot a)$ match no representative in $S$. Show this is the boundary that makes the conjecture illegal: without closedness the transition $\delta(\mathrm{row}(s),a)=\mathrm{row}(s\cdot a)$ has no target state, and without consistency it depends on the chosen representative (\cref{def:closed-consistent}); the two conditions are precisely what licenses the quotient automaton (\leanname{FLT_Proofs.Computation.DFA'}), so they are well-definedness constraints, not optimizations.

\item \textbf{A counterexample is a new-class witness.} Replace the teacher's \textsc{yes}/counterexample reply by a bare bit ``correct or not''. Show this is the boundary that disables refinement: the equivalence oracle's counterexample carries the contract that $z$ is a point where the conjecture and target disagree (\leanname{FLT_Proofs.Data.EquivalenceOracle}; the witness is typed by \leanname{FLT_Proofs.Data.Counterexample}), and adding $z$ with its prefixes forces at least one new distinct row, i.e.\ splits a Myhill--Nerode class; a bare bit removes the witness and the loop can no longer discover the missing state.

\item \textbf{The $n-1$ bound is the index of the congruence.} Read ``at most $n-1$ equivalence queries'' not as a runtime figure but as a structural identity. Show this is the frontier that grades termination: the minimal DFA has $n$ states, hence the congruence has index $n$ (\cref{rmk:table-intuition}, uniqueness by \cite{RabinScott1959}), the table starts at one class and gains at least one per counterexample, so termination after $\le n-1$ rounds \emph{is} the finiteness of the index (\cref{thm:lstar-complexity}); the bound is a literature theorem (\cite{Angluin1987}), not yet kernel-proved (\cref{op:lstar-formal}), so the termination annotation is drawn dashed.

\item \textbf{The whole loop is cited, not verified.} Ask the kernel for the statement ``$L^*$ converges to the minimal DFA''. Show this is the figure's principal frontier: the nodes are verified types ($L^*$ itself is the structure \leanname{LStar} with fields \texttt{teacher}, \texttt{numStates}, \texttt{result}, all sorry-free, but with no method and no correctness lemma), so the objects are solid while the convergence-and-termination edges are dashed and carry \cite{Angluin1987}; closing this gap on top of the typed teacher (\leanname{FLT_Proofs.Learner.Active.MinimallyAdequateTeacher}) and \leanname{FLT_Proofs.Computation.DFA'} is exactly \cref{op:lstar-formal}.

\item \textbf{The teacher is the only source of state-splits.} Read the minimally adequate teacher as the entire interface between learner and target. Show this is the instance that names the model: a MAT is a membership oracle and an equivalence oracle answering about one shared target (\leanname{FLT_Proofs.Learner.Active.MinimallyAdequateTeacher}, field \texttt{consistent} forcing $\texttt{mq.target}=\texttt{eq.target}$; \cref{def:mat}), and nothing else, so every state the learner discovers is forced by a query, never by distributional luck; this is what makes exact identification (\leanname{FLT_Proofs.Criterion.PAC.ExactLearnable}, \cref{def:exact-learning-formal}) attainable here.

\item \textbf{Drop membership queries and the loop blows up.} Run the same cycle with equivalence queries alone. Show this is the boundary that collapses the discovery rate: without membership queries the learner cannot place a string in its Myhill--Nerode class, so each equivalence query removes only one candidate from the version space and an $N$-concept class costs $N-1$ queries in the worst case (\cref{prop:eq-lower}, \cite{Angluin1988}), whereas membership queries supply the structural bit that collapses an exponential count to the linear $n-1$ of \cref{thm:lstar-complexity}.

\item \textbf{Passive data cannot drive the loop at all.} Replace the active teacher with random labelled examples drawn from a distribution. Show this is the frontier that separates the Wills: the same target class (DFAs) is exactly learnable by this loop (\cref{thm:lstar-complexity}) yet not PAC-learnable under cryptographic assumptions (\cref{thm:dfa-pac-hard}, \cite{KearnsValiant1994}), so the passive--active edge is qualitative, and the reverse direction (PAC-easy yet exact-hard) is itself open (\cref{op:reverse-separation}).

\item \textbf{The counterexample length is a hidden cost on the arrows.} Refine the $z$-arrow by asking how long $z$ may be. Show this is a frontier on the membership budget: the equivalence oracle may return an adversarially long counterexample (\cref{rmk:counterexample-length}), so the membership-query count carries a factor $m$ (\cref{thm:lstar-complexity}); whether a learner can binary-search the counterexample to remove the linear dependence on $m$ is open (\cref{op:counterexample-length}, \cite{RivestSchapire1993}).

\item \textbf{Any queryable function is an $L^*$ target.} Replace the unknown regular language by the binary string function of a trained attention model, answering membership by a forward pass and equivalence by sampling. Show this is the figure's most general instance: the model is a legitimate minimally adequate teacher (\cref{prop:transformer-extraction}, \cite{WeissGoldbergYahav2018}), so the same state-discovery loop extracts a DFA exact on the queried strings; whether the model's function even has a finite Myhill--Nerode index (a genuine regular core) so that extraction is exact rather than heuristic is open (\cref{op:transformer-regular}), and so is the soundness of an unreliable oracle such as an LLM critic (\cref{op:llm-teacher}).
\end{enumerate}

% Chapter 9: Universal Learning and the Trichotomy
% universal_learning = Profile C: what is the optimal learning RATE over all distributions?
% universal_trichotomy = Profile C: the modern capstone -- three regimes,
%   classified by Littlestone trees (from online learning!) applied to the statistical setting.
% Centerpiece: the SURPRISE that the Littlestone dimension controls statistical rates.

\chapter{Universal Learning and the Trichotomy}
\label{ch:universal}

The Fundamental Theorem of \Cref{ch:pac} answers a binary question: \emph{is}
$\mathcal{H}$ learnable?  The answer is yes if and only if
$\VC(\mathcal{H}) < \infty$.  But the binary answer hides a quantitative
question.  Among all learnable classes, some are learned faster than others.
How fast?

PAC learning theory phrases this quantitatively through the sample complexity
$m(\varepsilon, \delta)$: how many samples suffice for accuracy $\varepsilon$
with confidence $1 - \delta$?  But the answer, $\Theta(d/\varepsilon)$ in
the realizable case, is a function of the \emph{accuracy parameter}.  It
describes how accuracy improves as data grows, with $d$ as a constant.  It
does not describe the \emph{rate} at which the risk itself decreases as a
function of the sample size $n$.

This chapter asks the rate question directly: given $n$ samples from an
unknown distribution $D$, how fast can the risk $\risk_D(A(S))$ be driven to
zero, \emph{uniformly over all distributions $D$ for which learning is
possible}?  The answer is the trichotomy theorem, which classifies every
hypothesis class into exactly one of three rate regimes.  The classification
involves a combinatorial object from an unexpected source.


% ================================================================
% SECTION 1: THE QUESTION
% ================================================================

\section{Beyond PAC Sample Complexity}
\label{sec:universal-question}

Fix a hypothesis class $\mathcal{H}$ with $\VC(\mathcal{H}) < \infty$, so
that $\mathcal{H}$ is PAC learnable.  The PAC framework guarantees a learner
$A$ and a sample complexity $m(\varepsilon, \delta)$ such that, for any
distribution $D$ and any $\varepsilon, \delta > 0$, drawing
$m \geq m(\varepsilon, \delta)$ samples ensures
$\Pr[\risk_D(A(S)) > \varepsilon] \leq \delta$.

Now invert the perspective.  Instead of asking ``how many samples for
accuracy $\varepsilon$?'', ask: ``given $n$ samples, what is the best
accuracy achievable?''  More precisely, define the \emph{learning rate} of a
class $\mathcal{H}$ as follows.

\begin{defn}[Universal Learning Rate]
\label{def:universal-rate}
A function $R \colon \N \to [0,1]$ is a \emph{universal learning rate} for
$\mathcal{H}$ if there exists a learner $A$ such that, for every realizable
distribution $D$ (i.e., $\risk_D(h^*) = 0$ for some $h^* \in \mathcal{H}$)
and every $n \in \N$:
\[
  \E_{S \sim D^n}\!\bigl[\risk_D(A(S))\bigr] \leq R(n).
\]
The \emph{optimal universal rate} is the fastest-decaying $R$ achievable by
any learner.\formal{FLT_Proofs.Criterion.Extended.UniversalLearnable}
\end{defn}

Three features distinguish this from the PAC formulation:

\begin{enumerate}[leftmargin=*,itemsep=3pt]
\item \textbf{No fixed $\varepsilon$.}  The rate $R(n)$ describes how the
  expected risk shrinks with $n$, rather than asking for a threshold sample
  size at a given $\varepsilon$.

\item \textbf{Uniformity over distributions.}  The same learner $A$ and the
  same rate $R(n)$ must work for \emph{every} realizable distribution $D$.
  The learner cannot be tuned to a specific distribution.

\item \textbf{Expected risk, not high-probability.}  We use
  $\E[\risk_D(A(S))]$ rather than a $(\varepsilon,\delta)$ guarantee.  This
  is a convenience that simplifies the rate classification without changing
  the qualitative picture.
\end{enumerate}

For a class with $\VC(\mathcal{H}) = d < \infty$, the PAC upper bound gives
$R(n) = O(d/n)$: a $\Theta(1/n)$ rate.  But is $1/n$ always achievable?
Can some classes be learned faster, exponentially fast?  And are there
learnable classes where no uniform $1/n$ rate is possible, so that the rate
degrades depending on the distribution?

All three phenomena occur.  The trichotomy theorem says these are the
\emph{only} three possibilities.


% ================================================================
% SECTION 2: THE TRICHOTOMY THEOREM
% ================================================================

\section{The Trichotomy Theorem}
\label{sec:trichotomy}

The classification involves a combinatorial object that generalizes the
Littlestone tree of \Cref{ch:online}.  Recall that a mistake tree
(\Cref{def:mistake-tree}) is a complete binary tree whose internal nodes are
labeled with instances $x \in X$, such that every root-to-leaf path is
realized by some $h \in \mathcal{H}$.  We now define a variant with a
weaker consistency requirement.

\begin{defn}[VCL Tree]
\label{def:vcl-tree}
A \emph{VCL tree} (Vapnik--Chervonenkis--Littlestone tree) for
$\mathcal{H}$ over $X$ is a complete binary tree $T$ in which:
\begin{itemize}[leftmargin=*,itemsep=2pt]
\item Each internal node $v$ is labeled with an instance $x_v \in X$.
\item The left child corresponds to label $0$, the right to label $1$.
\item For every root-to-leaf path $\pi = (v_1, y_1), \ldots, (v_d, y_d)$,
  there exists a \emph{set} $\{x_{v_1}, \ldots, x_{v_d}\}$ that is
  \emph{shattered} by $\mathcal{H}$, not merely a single hypothesis
  consistent with $\pi$, but the full shattering condition: for every
  labeling $b \in \{0,1\}^d$ of these $d$ points, some $h \in \mathcal{H}$
  realizes $b$.
\end{itemize}
The tree is \emph{infinite} if it has unbounded depth.
\end{defn}

\begin{remark}[VCL trees vs.\ Littlestone trees]
\label{rmk:vcl-vs-littlestone}
Every Littlestone tree is a VCL tree (shattering is a stronger condition
than path-consistency, but the shattering requirement in the VCL definition
applies to each path individually, and a Littlestone tree already provides a
consistent hypothesis for every path, since the tree is complete, all
$2^d$ paths are realized, which implies shattering of the instances along
any single path).  The converse fails: a VCL tree requires shattering along
each path, which is a weaker global condition than the Littlestone tree's
requirement that \emph{every} root-to-leaf path has a \emph{single}
consistent hypothesis.  The distinction matters: infinite VCL trees can
exist even when $\Ldim(\mathcal{H}) < \infty$.  The combined VC--Littlestone measure underlying these trees is typed in the development as a pair of ordinal-valued dimensions.\formal{FLT_Proofs.Complexity.Ordinal.OrdinalLittlestoneDim}
\end{remark}

\begin{thm}[The Trichotomy Theorem {\cite{BousquetHannekeMoranVanHandelYehudayoff2021}}]
\label{thm:trichotomy}
Let $\mathcal{H}$ be a hypothesis class with $|\mathcal{H}| \geq 3$.
Exactly one of the following three cases holds:
\begin{enumerate}[label=\textnormal{(T\arabic*)}]
\item \label{T:exp} \textbf{Exponential rate.}  $\mathcal{H}$ has no
  infinite Littlestone tree (equivalently, $\Ldim(\mathcal{H}) < \infty$).
  Then the optimal universal learning rate is \emph{exponential}:
  \[
    R(n) = \Theta\!\left(e^{-n}\right).
  \]

\item \label{T:linear} \textbf{Linear rate.}  $\mathcal{H}$ has an infinite
  Littlestone tree but no infinite VCL tree.  Then the optimal universal
  learning rate is \emph{linear}:
  \[
    R(n) = \Theta\!\left(\frac{1}{n}\right).
  \]

\item \label{T:arb} \textbf{Arbitrarily slow rates.}  $\mathcal{H}$ has an
  infinite VCL tree.  Then no uniform rate exists: for every function
  $R(n) \to 0$ (no matter how slowly), every learner $A$ fails to achieve
  rate $R(n)$ on some realizable distribution $D$.
\end{enumerate}
\end{thm}

\begin{remark}[Formalization status]
\label{rmk:trichotomy-formal}
The trichotomy's three \emph{regimes} are typed in the companion development (the universal-rate criterion (\cref{def:universal-rate}) and the ordinal VC and Littlestone dimensions), and the boundary case ``$\Ldim < \infty$ iff online learnable'' is verified (\cref{thm:littlestone-characterization}).  The trichotomy theorem itself is \emph{not} yet machine-checked: the middle branch (infinite Littlestone, finite VC, hence linear rate) rests on the Bousquet--Hanneke--Moran--Ramanan doubling construction, which the development marks as pending, and the universal-to-PAC direction routes through a boosting lemma that is currently incomplete.  Closing these is \cref{op:trichotomy-formal}; we present the proof as a sketch throughout, following~\cite{BousquetHannekeMoranVanHandelYehudayoff2021}.
\end{remark}

The three cases are exhaustive and mutually exclusive.  Every learnable class
(i.e., every class with $\VC(\mathcal{H}) < \infty$ and
$|\mathcal{H}| \geq 3$) falls into exactly one regime.  The boundaries are
sharp: there is no class whose optimal rate is, say, $1/\sqrt{n}$ or
$1/n^2$.

\begin{figure}[ht]
\centering
\begin{tikzpicture}[
    decision/.style={draw, thick, diamond, aspect=2.2, inner sep=2pt, font=\small,
      fill=defblue!8, draw=defblue!70!black, align=center, minimum width=3cm},
    outcome/.style={draw, thick, densely dashed, rounded corners, minimum width=2.7cm,
      minimum height=1.25cm, font=\small, align=center},
    sedge/.style={-{Stealth[length=2.5mm]}, thick},
    dedge/.style={-{Stealth[length=2.5mm]}, thick, densely dashed},
    yeslabel/.style={font=\scriptsize\bfseries, text=sepgreen!60!black},
    nolabel/.style={font=\scriptsize\bfseries, text=thmred},
    hdl/.style={font=\tiny, text=defblue!55!black, align=center},
]
% --- decisions: verified dimensional thresholds (solid) ---
\node[decision] (d1) at (0,0) {Infinite\\Littlestone tree?};
\node[hdl, below=1pt of d1] {\leanname{OrdinalLittlestoneDim}};
\node[decision] (d2) at (3.6,-3.1) {Infinite\\VCL tree?};
\node[hdl, below=1pt of d2] {\cref{def:vcl-tree}, \leanname{VCLTree}};
% --- rate outcomes: typed regimes, but the rate THEOREMS are pending -> dashed boxes ---
\node[outcome, fill=sepgreen!12, draw=sepgreen!60!black] (exp) at (-3.5,-6.1) {%
  \textbf{Exponential}\\$R(n)=\Theta(e^{-n})$\\{\scriptsize $\Ldim<\infty$}};
\node[outcome, fill=defblue!12, draw=defblue!60!black] (lin) at (1.5,-6.1) {%
  \textbf{Linear}\\$R(n)=\Theta(1/n)$\\{\scriptsize $\Ldim=\infty$, no $\infty$-VCL}};
\node[outcome, fill=thmred!12, draw=thmred!60!black] (slow) at (6.3,-6.1) {%
  \textbf{Arbitrarily slow}\\$\forall\,R(n)\to0$: fails\\{\scriptsize $\infty$-VCL exists}};
% --- edges: D1 solid (verified boundary), D2 dashed (rate consequence is literature) ---
\draw[sedge] (d1) -- node[nolabel, above left, pos=0.35]{No}
   node[hdl, below, pos=0.62]{\leanname{littlestone_characterization}} (exp);
\draw[sedge] (d1) -- node[yeslabel, above right, pos=0.4]{Yes} (d2);
\draw[dedge] (d2) -- node[nolabel, above left, pos=0.4]{No} (lin);
\draw[dedge] (d2) -- node[yeslabel, above right, pos=0.4]{Yes} (slow);
% --- the inclusion ladder + the formalization-status annotation ---
\node[hdl, text=black!55, anchor=west] at (4.55,-1.15)
   {$\{\Ldim{<}\infty\}\subsetneq\{$no $\infty$-VCL$\}\subsetneq\{$all$\}$};
\node[font=\scriptsize\itshape, text=thmred!75!black, align=center, anchor=north, text width=11cm] at (1.4,-7.5)
   {Solid: the decisions are verified thresholds, and \leanname{littlestone_characterization} is the one verified rate boundary. Dashed: the $\Theta$-rates and the trichotomy are \cite{BousquetHannekeMoranVanHandelYehudayoff2021}, typed but not yet machine-checked (\cref{op:trichotomy-formal}).};
\end{tikzpicture}
\caption[The universal-rate trichotomy: two tree-conditions, three regimes]{The trichotomy decision tree.
Two combinatorial questions partition every hypothesis class with $|\mathcal{H}|\geq 3$ into exactly three
universal learning-rate regimes. The first question, does $\mathcal{H}$ admit an infinite Littlestone
tree, is a verified dimensional threshold: by the Littlestone characterization
(\cref{thm:littlestone-characterization}, \leanname{littlestone_characterization}) a finite Littlestone
dimension is exactly online learnability, and this is the boundary of the exponential regime
$R(n)=\Theta(e^{-n})$. The second question, does $\mathcal{H}$ admit an infinite VCL tree, uses the VCL
tree of \cref{def:vcl-tree}, which generalizes the Littlestone tree by requiring that the instances along
each root-to-leaf path be shattered rather than merely consistent with one hypothesis; its finiteness
separates the linear regime $R(n)=\Theta(1/n)$ from the arbitrarily-slow regime. The two decision nodes
are drawn solid because the dimensions and the VCL tree are typed in the companion development; the three
rate values and the ``exactly one of three'' claim are drawn dashed because they are the result of
Bousquet--Hanneke--Moran--van Handel--Yehudayoff~\cite{BousquetHannekeMoranVanHandelYehudayoff2021}, cited
here and not yet machine-checked (\cref{op:trichotomy-formal}). There is no fourth regime and no
intermediate rate.}
\label{fig:trichotomy-tree}
\end{figure}

\subsection{An Online Concept in a Statistical Theorem}
\label{sec:surprise}

The trichotomy theorem is the structural capstone of this book.  To
appreciate what it says, consider the expectations each tradition brought
to the table.

The PAC tradition (\Cref{ch:pac}) characterized learnability through the VC
dimension and established the sample complexity
$\Theta(d/\varepsilon)$.  The natural expectation: the VC dimension should
also determine the learning rate.  It does not.
Two classes with the same VC dimension can have different optimal
rates: one exponential, one $1/n$.  So what draws the boundary?

The online learning tradition (\Cref{ch:online}) introduced the Littlestone
dimension to characterize mistake-bounded learning against an adversary.
A worst-case combinatorial quantity, built for a worst-case adversarial
game.  No connection to statistical convergence rates.  None expected.

And yet.  The Littlestone dimension, this adversarial, combinatorial
object, is exactly what governs the boundary between exponential and
polynomial convergence in the i.i.d.\ setting.

Neither tradition predicted this.  The trichotomy does not merely classify
rates; it reveals that the deepest structural boundary in statistical
learning theory is drawn by a quantity from a different paradigm entirely.

\subsection{Examples}
\label{sec:trichotomy-examples}

\begin{example}[The three regimes, witnessed]
\label{ex:three-regimes}
\begin{enumerate}[label=(\alph*),leftmargin=*,itemsep=4pt]
\item \textbf{Exponential: finite classes.}  Any finite class
  $\mathcal{H}$ with $|\mathcal{H}| = k$ has
  $\Ldim(\mathcal{H}) \leq \lfloor\log_2 k\rfloor < \infty$.  No infinite
  Littlestone tree exists, so $\mathcal{H}$ falls into regime~\ref{T:exp}.
  Concretely, the Halving algorithm achieves risk $\leq k \cdot 2^{-n}$
  after $n$ rounds.

\item \textbf{Exponential: halfspaces.}  The class of halfspaces in
  $\R^d$ has $\Ldim = d < \infty$ (\Cref{ex:ldim-examples}).  This class
  is learned at an exponential rate, despite having infinite cardinality.

\item \textbf{Linear: thresholds on $\R$.}  The class
  $\mathcal{H}_{\mathrm{thr}} = \{x \mapsto \ind[x \geq \theta] :
  \theta \in \R\}$ has $\VC = 1$ but $\Ldim = \infty$
  (\Cref{prop:pac-online-sep}).  It has no infinite VCL tree (the
  one-dimensional ordering prevents shattering along arbitrarily long
  adaptive paths).  The optimal rate is $\Theta(1/n)$.

\item \textbf{Arbitrarily slow: unions of intervals.}  For each $k$, let
  $\mathcal{H}_k$ be the class of unions of at most $k$ intervals on $\R$.
  The class $\mathcal{H}_\infty = \bigcup_k \mathcal{H}_k$ (unions of
  finitely many intervals, with no bound on the number) has
  $\VC < \infty$ but admits infinite VCL trees.  No uniform rate $R(n) \to 0$
  is achievable.
\end{enumerate}
\end{example}


% ================================================================
% SECTION 3: PROOF SKETCH
% ================================================================

\section{Proof Architecture}
\label{sec:proof-sketch}

The full proof of \Cref{thm:trichotomy} occupies much of
\cite{BousquetHannekeMoranVanHandelYehudayoff2021}.  We sketch the key ideas
for each regime, emphasizing why the boundaries fall where they do.


\subsection{The Exponential Case: Finite Littlestone Dimension}
\label{sec:exp-case}

Suppose $\Ldim(\mathcal{H}) = d < \infty$.  The Standard Optimal Algorithm
($\SOA$) from \Cref{ch:online} was designed for adversarial online learning,
but its structure can be adapted to the statistical setting.

The key insight is that a learner with access to i.i.d.\ samples can
simulate a version-space strategy.  Consider the following approach: given a
sample $S$ of size $n$, maintain the version space
$V = \{h \in \mathcal{H} : h \text{ consistent with } S\}$.  Return a
hypothesis from $V$ using an $\SOA$-derived selection rule.

When $\Ldim(\mathcal{H}) = d$, the version space after seeing $n$ i.i.d.\
samples has the following property: the probability (over $D$) that there
exists a hypothesis $h \in V$ with $\risk_D(h) > \varepsilon$ decreases
exponentially in $n$.  This is because each sample point, with probability
at least $\varepsilon$ under $D$, eliminates all hypotheses in $V$ that
disagree with the true label at that point.  After $n$ samples, the survival
probability of any ``bad'' hypothesis (one with risk $> \varepsilon$) is at
most $(1 - \varepsilon)^n \leq e^{-\varepsilon n}$.  The finite Littlestone
dimension controls the effective complexity of $V$ through the mistake tree
structure, ensuring that a union bound over the relevant version space
subsets remains finite.

The result:
\[
  \E[\risk_D(A(S))] \leq C \cdot 2^d \cdot e^{-n/C'}
\]
for constants $C, C'$ depending on $d$.  The rate is exponential in $n$.

The matching lower bound shows that exponential convergence is impossible
when $\Ldim(\mathcal{H}) = \infty$: an infinite Littlestone tree provides an
adversary-like construction (within the distributional framework) that
forces the learner's expected risk to decay no faster than $\Omega(1/n)$.


\subsection{The Linear Case: Infinite Littlestone, Finite VCL}
\label{sec:linear-case}

Suppose $\Ldim(\mathcal{H}) = \infty$ but $\mathcal{H}$ has no infinite VCL
tree.  The exponential strategy fails because the $\SOA$-based approach
requires finite Littlestone dimension.  But a different algorithm, the
\emph{one-inclusion predictor}, achieves the $1/n$ rate.

\begin{defn}[One-Inclusion Predictor (Informal)]
\label{def:oip}
Given a sample $S = \{(x_1, y_1), \ldots, (x_n, y_n)\}$ and a new point
$x$, form the multiset $\{x_1, \ldots, x_n, x\}$.  Consider the one-inclusion
graph $G$: the graph whose vertices are all labelings of this multiset
consistent with some $h \in \mathcal{H}$, with edges connecting labelings
that differ on exactly one point.  Predict the label $\hat{y}$ that
minimizes the maximum degree of the vertex in $G$ corresponding to the
completed labeling.
\end{defn}

Haussler, Littlestone, and Warmuth~\cite{Haussler1988} introduced the
one-inclusion graph technique.  Bousquet et
al.~\cite{BousquetHannekeMoranVanHandelYehudayoff2021} showed that when
$\mathcal{H}$ has no infinite VCL tree, the one-inclusion predictor achieves
\[
  \E[\risk_D(A(S))] \leq \frac{C}{n}
\]
for a constant $C$ depending on $\mathcal{H}$.  The absence of infinite VCL
trees is precisely the combinatorial condition needed for the one-inclusion
graph to have bounded degree growth, which in turn controls the expected
risk.

The matching lower bound is the PAC lower bound: any class with
$\VC(\mathcal{H}) \geq 1$ requires $\Omega(1/n)$ expected risk, since
distinguishing hypotheses that differ on a $1/n$-fraction of the domain
requires $\Omega(n)$ samples.


\subsection{The Slow Case: Infinite VCL Trees}
\label{sec:slow-case}

Suppose $\mathcal{H}$ has an infinite VCL tree.  The theorem claims that no
uniform rate $R(n) \to 0$ is achievable: for every learner $A$ and every
function $R(n) \to 0$, there exists a realizable distribution $D$ such that
$\E[\risk_D(A(S_n))] \geq R(n)$ for infinitely many $n$.

The proof constructs, for any candidate learner $A$ and any target rate
$R(n)$, a distribution that defeats $A$ at the desired rate.  The
construction uses the infinite VCL tree as a source of combinatorial
richness.

Descend the VCL tree adaptively, choosing at each node the branch that is
harder for $A$.  Because the tree is infinite and each path corresponds to a
shattered set, the adversary can build a distribution at any depth $d$ whose
effective VC dimension on the relevant support is $d$.  By choosing $d$
large enough (as a function of $n$), the adversary ensures that the PAC
lower bound $\Omega(d/n)$ exceeds $R(n)$.  The infinite VCL tree provides
arbitrarily large effective VC dimension along adaptive paths, preventing
any uniform rate from holding.

This is the sense in which infinite VCL trees are an obstruction: they
provide an inexhaustible reservoir of complexity that any learner eventually
confronts.


% ================================================================
% SECTION 4: THE CROSS-PARADIGM MAP
% ================================================================

\section{The Cross-Paradigm Map}
\label{sec:cross-paradigm}

The trichotomy theorem connects three chapters of this book through a single
combinatorial object.

\begin{figure}[ht]
\centering
\begin{tikzpicture}[
    regime/.style={draw, thick, ellipse, minimum width=3.6cm,
      minimum height=2.2cm, font=\small, align=center},
    example/.style={font=\scriptsize, align=center, text=notegray},
    label/.style={font=\small\bfseries, align=center},
    connection/.style={thick, dashed, notegray},
]
% Three regimes as overlapping regions
\node[regime, fill=sepgreen!10] (exp) at (-3.5, 0) {};
\node[label] at (-3.5, 0.4) {$e^{-n}$};
\node[example] at (-3.5, -0.3) {finite classes\\halfspaces\\conjunctions};

\node[regime, fill=defblue!10] (lin) at (0, 0) {};
\node[label] at (0, 0.4) {$1/n$};
\node[example] at (0, -0.3) {thresholds on $\R$\\intervals on $\R$};

\node[regime, fill=thmred!10] (slow) at (3.5, 0) {};
\node[label] at (3.5, 0.4) {arb.\ slow};
\node[example] at (3.5, -0.3) {$\bigcup_k \mathcal{H}_k$\\unbounded\\unions};

% Boundaries
\draw[very thick, sepgreen!70!black, {Stealth[length=2mm]}-{Stealth[length=2mm]}]
  (-1.6, -1.8) -- (-1.6, 1.8)
  node[above, font=\scriptsize, align=center, text=sepgreen!70!black]
  {$\Ldim < \infty$\\vs.\ $\Ldim = \infty$};

\draw[very thick, thmred!70!black, {Stealth[length=2mm]}-{Stealth[length=2mm]}]
  (1.7, -1.8) -- (1.7, 1.8)
  node[above, font=\scriptsize, align=center, text=thmred!70!black]
  {no $\infty$-VCL\\vs.\ $\infty$-VCL};

% Chapter connections
\node[font=\scriptsize\itshape, text=defblue, align=center] at (-3.5, -2.6)
  {Ch.~6: Online\\(Littlestone dim)};
\node[font=\scriptsize\itshape, text=defblue, align=center] at (0, -2.6)
  {Ch.~5: PAC\\(VC dim)};
\node[font=\scriptsize\itshape, text=defblue, align=center] at (3.5, -2.6)
  {Ch.~9: Universal\\(VCL trees)};

\draw[connection, -{Stealth[length=1.5mm]}] (-3.5, -2.2) -- (-1.8, -1.6);
\draw[connection, -{Stealth[length=1.5mm]}] (0, -2.2) -- (0, -1.2);
\draw[connection, -{Stealth[length=1.5mm]}] (3.5, -2.2) -- (1.9, -1.6);
\end{tikzpicture}
\caption{The three rate regimes and their governing combinatorial conditions.
The left boundary is drawn by the Littlestone dimension (from online
learning, \Cref{ch:online}); the right boundary by VCL trees (a hybrid of
VC shattering and Littlestone trees).  The VC dimension (\Cref{ch:pac})
determines \emph{whether} learning is possible; the Littlestone dimension
and VCL trees determine \emph{how fast}.}
\label{fig:rate-regimes}
\end{figure}

The structural lesson is a hierarchy of three combinatorial conditions,
each finer than the last:

\begin{center}\small
\renewcommand{\arraystretch}{1.3}
\begin{tabular}{llll}
\toprule
\textbf{Condition} & \textbf{Governs} & \textbf{Chapter} & \textbf{Paradigm} \\
\midrule
$\VC(\mathcal{H}) < \infty$ & Learnability (yes/no) & \Cref{ch:pac} & PAC \\
$\Ldim(\mathcal{H}) < \infty$ & Rate: $e^{-n}$ vs.\ $1/n$ & \Cref{ch:online} & Online \\
No infinite VCL tree & Rate: $1/n$ vs.\ arb.\ slow & This chapter & Universal \\
\bottomrule
\end{tabular}
\end{center}

Each row refines the previous one.  The VC dimension draws the coarsest
line (learnable vs.\ not).  The Littlestone dimension draws a finer line
within the learnable classes (fast vs.\ slow).  The VCL tree draws the
finest line (uniformly slow vs.\ hopelessly slow).

This hierarchy was not predicted by any single tradition.  It emerged from
the trichotomy theorem's proof, which required techniques from both the PAC
and online learning literatures.  The result is a unified picture in which
the three paradigms (PAC, online, and universal) are not separate theories
with accidental parallels, but different projections of a single
combinatorial landscape.


\subsection{Where attention sits in the trichotomy}
\label{sec:attention-trichotomy}

The trichotomy classifies every fixed-architecture attention class, and the classification turns on a single question carried over from \Cref{ch:online}: is the Littlestone dimension of the attention class finite?

\begin{prop}[A fixed attention class is not arbitrarily slow]
\label{prop:attention-trichotomy}
A finite-head attention class on $\R^d$ has finite VC dimension (\cref{prop:attention-pac}), so it has no infinite VCL tree and is not in the arbitrarily-slow regime~\ref{T:arb}: it is universally learnable at a uniform rate.\formal{FLT_Proofs.Criterion.Extended.UniversalLearnable} Its regime is exponential~\ref{T:exp} when its Littlestone dimension is finite and linear~\ref{T:linear} when its Littlestone dimension is infinite.
\end{prop}

\begin{proof}
Finite VC dimension rules out infinite VCL trees (an infinite VCL tree shatters sets of every size, forcing infinite VC dimension), so by \cref{thm:trichotomy} the class lies in regime~\ref{T:exp} or~\ref{T:linear}.  Which of the two is decided by the Littlestone dimension: finite gives the exponential rate via the version-space strategy of \cref{sec:exp-case}, infinite gives the linear rate via the one-inclusion predictor of \cref{sec:linear-case}.
\end{proof}

The rate at which a transformer head learns its target (exponential or merely linear) is therefore governed by exactly the invariant whose finiteness for attention is open (\cref{op:ldim-attention} of \Cref{ch:online}).  A fixed attention class behaves like halfspaces (exponential) if its score-induced order structure has bounded Littlestone depth, and like thresholds on $\R$ (linear) if the temperature lets the router resolve an unbounded ordering; deciding which is the content of \cref{op:attention-regime}.

\section{What This Chapter Established}
\label{sec:ch9-summary}

Two results, one structural lesson.

\begin{enumerate}[leftmargin=*,itemsep=3pt]
\item \textbf{Universal learning rates.}  The optimal learning rate asks a
  harder question than PAC sample complexity: not how many samples for a
  fixed accuracy, but how fast the risk decreases as a function of $n$,
  uniformly over all realizable distributions.

\item \textbf{The trichotomy.}  Every hypothesis class with
  $|\mathcal{H}| \geq 3$ falls into exactly one of three regimes:
  exponential ($e^{-n}$), linear ($1/n$), or arbitrarily slow.  The
  boundaries are determined by the Littlestone dimension and VCL trees.
  There are no intermediate rates.

\item \textbf{The cross-paradigm bridge.}  The Littlestone dimension, born
  in the adversarial online setting of \Cref{ch:online}, controls the
  boundary between exponential and polynomial convergence in the i.i.d.\
  statistical setting.  This was the central open question resolved by
  Bousquet et al.~\cite{BousquetHannekeMoranVanHandelYehudayoff2021}: not
  just what the rate regimes are, but what combinatorial objects govern
  them; and the answer came from an unexpected paradigm.
\end{enumerate}

\section{Open Problems}
\label{sec:ch9-open}

Each problem is anchored to a result above and tagged: a \emph{known unknown} (\textsc{ku}) is an articulated open question; an \emph{unknown known} (\textsc{uk}) is a result in the literature or the verified development not yet absorbed into a clean statement.  Verified handles are named where they exist.

\begin{openprob}[Formalizing the trichotomy, \textsc{uk}]
\label{op:trichotomy-formal}
The regimes are typed (\leanname{FLT_Proofs.Criterion.Extended.UniversalLearnable}, the ordinal dimensions) but the trichotomy theorem is not machine-checked (\cref{rmk:trichotomy-formal}).  Formalize the Bousquet--Hanneke--Moran--Ramanan middle branch (infinite Littlestone, finite VC, hence linear rate) via the doubling/aggregation construction, and complete the boosting lemma in the universal-to-PAC direction (\leanname{FLT_Proofs.Theorem.Separation.universal_imp_pac}) that currently rests on an unfinished step.
\end{openprob}

\begin{openprob}[The trichotomy regime of attention, \textsc{uk}]
\label{op:attention-regime}
\Cref{prop:attention-trichotomy} places a fixed-head attention class in the exponential or linear regime according to its Littlestone dimension.  Decide which: compute the Littlestone dimension of the $k$-head attention class as a function of architecture and temperature, settling whether a transformer head learns its target exponentially or only linearly (tied to \cref{op:ldim-attention} of \Cref{ch:online}).
\end{openprob}

\begin{openprob}[A single dimension for the linear/slow boundary, \textsc{ku}]
\label{op:vcl-characterization}
The boundary between linear and arbitrarily-slow rates is ``no infinite VCL tree''.  Give a clean combinatorial dimension (an ordinal-valued measure on the class) whose finiteness is equivalent to the absence of an infinite VCL tree, the way $\Ldim < \infty$ captures the exponential/linear boundary, building on the typed ordinal dimensions (\leanname{FLT_Proofs.Complexity.Ordinal.OrdinalLittlestoneDim}).
\end{openprob}

\begin{openprob}[Intermediate rates beyond the realizable case, \textsc{ku}]
\label{op:intermediate-rates}
The trichotomy is sharp in the realizable setting: only $e^{-n}$, $1/n$, and arbitrarily slow occur.  Determine the rate landscape in the agnostic and bounded-noise settings: do intermediate rates such as $1/\sqrt{n}$ appear, and is there a corresponding finer trichotomy or a genuine continuum?
\end{openprob}

\begin{openprob}[The rate constants, \textsc{ku}]
\label{op:rate-constants}
\Cref{sec:exp-case,sec:linear-case} give $\Theta(e^{-n})$ and $\Theta(1/n)$ with unspecified constants.  Determine the leading constant in each regime as a function of the Littlestone dimension (exponential) and of the one-inclusion graph degree (linear), matching upper and lower bounds.
\end{openprob}

\begin{openprob}[The multiclass trichotomy, \textsc{uk}]
\label{op:universal-multiclass}
The trichotomy is stated for binary classification.  Determine whether it survives for multiclass and real-valued prediction, with the DS or fat-shattering tree in place of the VCL tree, and whether the three-regime structure persists or splits further (\Cref{ch:extensions}).
\end{openprob}

\begin{openprob}[Distribution-dependent universal rates, \textsc{ku}]
\label{op:adaptive-rates}
The universal rate is uniform over all realizable distributions.  Characterize the finer, distribution-dependent rates (the fastest rate achievable on each fixed $D$) and determine how the per-distribution rate relates to the worst-case regime of \cref{thm:trichotomy}.
\end{openprob}

\begin{openprob}[Temperature and the rate transition, \textsc{uk}]
\label{op:attention-temperature-rate}
\Cref{prop:attention-trichotomy} ties an attention class's regime to its Littlestone dimension.  Determine how the score temperature drives the transition: as the temperature grows and the router approaches a hard argmax, does the Littlestone dimension diverge (pushing the class from exponential to linear), and is there a sharp temperature threshold?
\end{openprob}

\begin{openprob}[The one-inclusion predictor for attention, \textsc{uk}]
\label{op:oneinclusion-attention}
The linear regime is achieved by the one-inclusion predictor (\cref{def:oip}).  Give an efficient implementation of the one-inclusion predictor for a finite-head attention class and bound the constant $C$ in its $C/n$ rate in terms of the architecture, since the abstract predictor is not obviously tractable on attention hypothesis classes.
\end{openprob}

\begin{openprob}[Why online governs statistical rates, \textsc{ku}]
\label{op:universal-online-bridge}
The trichotomy's central surprise is that the Littlestone dimension, an adversarial, online quantity, draws the exponential/linear boundary in the i.i.d.\ setting (\cref{sec:surprise}).  Find a single structural principle, beyond the trichotomy proof, that explains the cross-paradigm bridge: a reason the same combinatorial object should control both worst-case mistakes and average-case rates.
\end{openprob}


\subsection*{Counterfactual exercises: generalizing Figure~\ref{fig:trichotomy-tree}}

Each exercise perturbs the decision tree.  Argue, in the figure's grammar (a two-question ladder of combinatorial tree-conditions of increasing permissiveness, finite Littlestone dimension inside no-infinite-VCL-tree inside all classes, whose two decision nodes are drawn solid where a verified ordinal dimension or the VCL tree definition certifies the threshold, and whose three rate leaves are drawn dashed where the rate value and the partition are Bousquet--Hanneke--Moran--van Handel--Yehudayoff's theorem, cited not proved), that the perturbed picture subsumes the concrete one as an \emph{instance} (a class routed to its regime), a \emph{boundary} (one tree-condition toggled so the regime jumps), or a \emph{frontier} (a dashed, open or literature-only edge).

\begin{enumerate}[leftmargin=*,itemsep=3pt]
\item \textbf{The first diamond is the online-learnable boundary.} Read decision one not as a yes/no on trees but as a verified dimensional threshold. Show this is the figure's defining instance on the exponential branch: ``no infinite Littlestone tree'' is exactly $\Ldim(\mathcal{H}) < \infty$, which is exactly online learnability by the Littlestone characterization (\cref{thm:littlestone-characterization}, \leanname{littlestone_characterization}), so the \textsc{no} edge into the $\Theta(e^{-n})$ leaf is the only rate-regime boundary the kernel certifies; the ordinal refinement used for the universal setting is \leanname{FLT_Proofs.Complexity.Ordinal.OrdinalLittlestoneDim}.

\item \textbf{The VCL tree generalizes the Littlestone tree by shattering.} Replace the second diamond's ``per-path consistency'' by ``per-path shattering''. Show this is the instance that names the second object: a VCL tree (\cref{def:vcl-tree}) demands that the instances along each root-to-leaf path be shattered, not merely consistent with a single hypothesis, so every Littlestone tree is a VCL tree and the converse fails (\cref{rmk:vcl-vs-littlestone}); this is why an infinite VCL tree can exist with $\Ldim < \infty$, making decision two strictly weaker than decision one and the ladder $\{\Ldim<\infty\} \subsetneq \{\text{no }\infty\text{-VCL}\} \subsetneq \{\text{all}\}$ strict at both steps.

\item \textbf{The leaves are typed regimes, the rates are not theorems yet.} Ask the kernel for the statement ``this class learns at rate $\Theta(e^{-n})$''. Show this is the figure's principal frontier: the three regimes are typed (the universal-rate criterion \cref{def:universal-rate} is \leanname{FLT_Proofs.Criterion.Extended.UniversalLearnable}, the ordinal dimensions exist), so the leaf \emph{boxes} are well-typed objects, but the rate \emph{values} $\Theta(e^{-n})$, $\Theta(1/n)$, and arbitrarily-slow, together with the ``exactly one of three'' claim, are not machine-checked; closing them, by formalizing the middle-branch doubling construction and completing the boosting lemma in the universal-to-PAC direction, is exactly \cref{op:trichotomy-formal} (\cite{BousquetHannekeMoranVanHandelYehudayoff2021}), so those edges are dashed.

\item \textbf{Collapse the second diamond to a single dimension.} Suppose someone exhibits an ordinal-valued measure whose finiteness equals ``no infinite VCL tree''. Show this would redraw decision two as a clean threshold matching decision one: the linear/arbitrarily-slow boundary would read ``$\dim < \infty$'' the way the exponential/linear boundary reads ``$\Ldim < \infty$'', unifying both diamonds into one ordinal invariant with three value-ranges; this single-invariant question is open (\cref{op:vcl-characterization}, building on \leanname{FLT_Proofs.Complexity.Ordinal.OrdinalLittlestoneDim}).

\item \textbf{Shrink the class to a single threshold family.} Instantiate $\mathcal{H}$ as thresholds on $\mathbb{N}$. Show this is the boundary that lands on the middle leaf: thresholds have $\VC = 1 < \infty$ (\cref{def:vc-dim}) yet $\Ldim = \infty$, so decision one answers \textsc{yes} and decision two answers \textsc{no} (no infinite VCL tree), placing the class in the linear regime $\Theta(1/n)$; the witness is the threshold class of \Cref{ch:online}, whose $\Ldim=\infty$ is constructive while its $\Theta(1/n)$ rate is the dashed literature claim (\cref{op:trichotomy-formal}).

\item \textbf{Force the bottom leaf with an infinite VC dimension.} Take a class that shatters arbitrarily large finite sets. Show this is the boundary into the arbitrarily-slow regime: infinite VC dimension forces an infinite VCL tree (a tree whose every path shatters its instances), so decision two answers \textsc{yes} and no learner achieves any uniform rate; the leaf is dashed because the lower bound ``for every $R(n)\to 0$, some realizable $D$ defeats every learner'' is the theorem, not kernel-proved (\cref{op:trichotomy-formal}, \cite{BousquetHannekeMoranVanHandelYehudayoff2021}).

\item \textbf{A fixed attention head is never arbitrarily slow.} Replace $\mathcal{H}$ by a finite-head attention class on $\mathbb{R}^d$. Show this is the instance that prunes the bottom branch: finite VC dimension rules out an infinite VCL tree (\cref{prop:attention-trichotomy}, \leanname{FLT_Proofs.Criterion.Extended.UniversalLearnable}), so decision two answers \textsc{no} and the head is universally learnable; whether it sits in the exponential or the linear leaf is decided by its Littlestone dimension, whose finiteness for attention is itself open (\cref{op:attention-regime}, \cref{op:ldim-attention}).

\item \textbf{The Littlestone dimension is an online quantity drawing a statistical line.} Read decision one across paradigms. Show this is the frontier that makes the tree surprising: the same invariant that bounds worst-case mistakes in the adversarial online setting of \Cref{ch:online} draws the exponential/linear boundary in the i.i.d.\ statistical setting (\cref{sec:surprise}), and a single structural reason for this cross-paradigm coincidence, beyond the trichotomy proof, is open (\cref{op:universal-online-bridge}).

\item \textbf{Leave the realizable case and the partition may refine.} Relabel the leaves for agnostic or bounded-noise data. Show this is the frontier that may add regimes: the trichotomy is sharp only in the realizable setting, and whether intermediate rates such as $1/\sqrt{n}$ appear, splitting a leaf or introducing a fourth, is open (\cref{op:intermediate-rates}); the figure's ``no fourth regime'' annotation is therefore a realizable-case claim, not a universal one.

\item \textbf{Swap the VCL tree for a multiclass tree.} Replace binary labels by a finite or real-valued label set and the VCL tree by the DS or fat-shattering tree. Show this is the frontier that tests the ladder's stability: whether the three-regime structure persists with the new combinatorial object in the second diamond, or splits further, is open (\cref{op:universal-multiclass}, \Cref{ch:extensions}); decision two's object is paradigm-relative, while decision one's $\Ldim$ threshold is the fixed verified rung.
\end{enumerate}

% Chapter 10: Combinatorial Dimensions
% 3 revelations (ds_dimension, natarajan_dimension, star_number),
% 10 definitions with brief treatment.
% Technical core of Part III. Full Sauer--Shelah proof. Multiclass detective story.

\chapter{Combinatorial Dimensions}
\label{ch:dimensions}

Part~II characterized learnability in four paradigms, and in each case the
answer was a number: the VC dimension for PAC learning, the Littlestone
dimension for online learning, the mind-change ordinal for Gold-style
identification, the teaching dimension for exact learning.  This chapter and
the three that follow develop these numbers (and the two dozen others in the
concept graph) as objects of study in their own right.

The present chapter treats \emph{combinatorial dimensions}: quantities defined
by asking how richly a hypothesis class can label finite configurations of
points.  The simplest such quantity, the VC dimension, asks how many points
can be shattered.  The others refine, extend, or replace this question for
settings beyond binary classification.  We prove the two foundational
combinatorial results in full (the Sauer--Shelah lemma and the
VC dimension characterization already established in \Cref{ch:pac}), recall the
Littlestone dimension from \Cref{ch:online}, and then turn to the chapter's
narrative centerpiece: the thirty-year search for the correct combinatorial
dimension characterizing multiclass learnability, which ended only in 2022.

\medskip

The chapter is organized as follows.

\begin{enumerate}[leftmargin=*,itemsep=2pt]
\item \textbf{VC dimension and shattering} (\Cref{sec:vc-dim}): the
  definitions, the growth function, and the full proof of the Sauer--Shelah
  lemma.
\item \textbf{The Littlestone dimension} (\Cref{sec:ldim-recall}): a brief
  recall from \Cref{ch:online} and the relationship $\VC \leq \Ldim$.
\item \textbf{Beyond binary: pseudodimension and fat-shattering}
  (\Cref{sec:beyond-binary}): scale-sensitive dimensions for real-valued
  function classes.
\item \textbf{The multiclass story: from Natarajan to DS}
  (\Cref{sec:multiclass}): the detective story of multiclass PAC
  characterization.
\item \textbf{Other dimensions} (\Cref{sec:other-dims}): a concise catalog of
  further combinatorial dimensions, with forward references.
\end{enumerate}


% ================================================================
% SECTION 1: VC DIMENSION AND SHATTERING
% ================================================================

\section{VC Dimension and Shattering}
\label{sec:vc-dim}

The definitions in this section were introduced informally in \Cref{ch:objects}
and used without proof in \Cref{ch:pac}.  We now state them precisely and prove
the combinatorial backbone that supports the Fundamental Theorem.

\begin{defn}[Restriction and Shattering]
\label{def:shattering}
Let $\mathcal{H} \subseteq \{0,1\}^X$ be a hypothesis class and let
$S = \{x_1, \ldots, x_n\} \subseteq X$ be a finite set.  The
\emph{restriction} of $\mathcal{H}$ to $S$ is
\[
  \mathcal{H}|_S \;=\; \bigl\{\bigl(h(x_1), \ldots, h(x_n)\bigr) : h \in \mathcal{H}\bigr\}
  \;\subseteq\; \{0,1\}^n.
\]
We say $\mathcal{H}$ \emph{shatters} $S$ if $\mathcal{H}|_S = \{0,1\}^n$,
i.e., every labeling of $S$ is realized by some hypothesis in $\mathcal{H}$.\formal{FLT_Proofs.Complexity.VCDimension.Shatters}
\end{defn}

\begin{defn}[VC Dimension {\cite{VapnikChervonenkis1971}}]
\label{def:vc-dim}
The \emph{Vapnik--Chervonenkis dimension} of $\mathcal{H}$, denoted
$\VC(\mathcal{H})$, is the largest cardinality of a set shattered by
$\mathcal{H}$:
\[
  \VC(\mathcal{H}) \;=\; \max\bigl\{|S| : S \subseteq X,\; \mathcal{H}
  \text{ shatters } S\bigr\}.
\]
If $\mathcal{H}$ shatters arbitrarily large finite sets,
$\VC(\mathcal{H}) = \infty$.\formal{FLT_Proofs.Complexity.VCDimension.VCDim}
\end{defn}

\begin{example}[Halfplanes in $\R^2$]
\label{ex:halfplanes-shatter}
Let $\mathcal{H}$ be the class of halfplanes in $\R^2$:
$\mathcal{H} = \bigl\{x \mapsto \ind[\langle w, x \rangle \geq b] : w \in
\R^2,\, b \in \R\bigr\}$.  Any three non-collinear points are shattered
(\Cref{fig:shatter-halfplane}), but no four points can be: for any four
points, at least one is inside the convex hull of the other three, and the
labeling that assigns it the opposite label from the other three cannot be
realized by a halfplane.  Hence $\VC(\mathcal{H}) = 3$.
\end{example}

\begin{figure}[ht]
\centering
\begin{tikzpicture}[scale=0.7,
    pospt/.style={circle, fill=thmred, inner sep=1.8pt},
    negpt/.style={circle, fill=defblue, inner sep=1.8pt},
    label0/.style={font=\scriptsize, text=defblue},
    label1/.style={font=\scriptsize, text=thmred},
    sep/.style={thick, black!65},
    poshalf/.style={fill=thmred!9},
]
  % Header note: the eight witnesses give VC >= 3
  \node[font=\footnotesize, text=black!78, align=center] at (5.5,2.95)
    {All $2^3=8$ labelings realized by a halfplane: the three points are shattered, so $\VC\ge 3$.};
  % 8 labelings, each WITH its separating halfplane, in 2 rows of 4.
  % fields: col / la / lb / lc / lt(1=horiz,2=vert) / lp(line position) / pd(1=upper|right shaded,0=lower|left)
  \foreach \col/\la/\lb/\lc/\lt/\lp/\pd in {
    0/0/0/0/1/2.0/1,
    1/0/0/1/1/0.85/1,
    2/0/1/0/2/1.5/1,
    3/0/1/1/2/0.5/1,
    4/1/0/0/2/0.5/0,
    5/1/0/1/2/1.5/0,
    6/1/1/0/1/0.85/0,
    7/1/1/1/1/-0.6/1%
  } {
    \pgfmathtruncatemacro{\cx}{mod(\col,4)*3}
    \pgfmathtruncatemacro{\cyr}{int(\col/4)}
    \begin{scope}[shift={(\cx,{-\cyr*3.4})}]
      % separating line + lightly shaded label-1 (positive) halfplane, drawn first (behind points)
      \ifnum\lt=1
        \ifnum\pd=1
          \fill[poshalf] (-0.45,\lp) rectangle (2.45,2.2);
        \else
          \fill[poshalf] (-0.45,-0.9) rectangle (2.45,\lp);
        \fi
        \draw[sep] (-0.45,\lp) -- (2.45,\lp);
      \else
        \ifnum\pd=1
          \fill[poshalf] (\lp,-0.9) rectangle (2.45,2.2);
        \else
          \fill[poshalf] (-0.45,-0.9) rectangle (\lp,2.2);
        \fi
        \draw[sep] (\lp,-0.9) -- (\lp,2.2);
      \fi
      % the three points, on top of the shade
      \ifnum\la=1 \node[pospt] (a\col) at (0,0) {}; \node[label1, below=1pt] at (a\col) {$1$};
      \else        \node[negpt] (a\col) at (0,0) {}; \node[label0, below=1pt] at (a\col) {$0$};\fi
      \ifnum\lb=1 \node[pospt] (b\col) at (2,0) {}; \node[label1, below=1pt] at (b\col) {$1$};
      \else        \node[negpt] (b\col) at (2,0) {}; \node[label0, below=1pt] at (b\col) {$0$};\fi
      \ifnum\lc=1 \node[pospt] (c\col) at (1,1.5) {}; \node[label1, above=1pt] at (c\col) {$1$};
      \else        \node[negpt] (c\col) at (1,1.5) {}; \node[label0, above=1pt] at (c\col) {$0$};\fi
    \end{scope}
  }
  % The d=4 obstruction (Radon): VC < 4, hence VC = 3.
  \begin{scope}[shift={(4.0,-7.9)}]
    \coordinate (p1) at (0,0);
    \coordinate (p2) at (3,0);
    \coordinate (p3) at (1.5,2.2);
    \coordinate (pin) at (1.5,0.75);
    \fill[black!6] (p1) -- (p2) -- (p3) -- cycle;
    \draw[black!40, dashed] (p1) -- (p2) -- (p3) -- cycle;
    \node[negpt] at (p1) {}; \node[label0, below=1pt] at (p1) {$0$};
    \node[negpt] at (p2) {}; \node[label0, below=1pt] at (p2) {$0$};
    \node[negpt] at (p3) {}; \node[label0, above=1pt] at (p3) {$0$};
    \node[pospt] at (pin) {}; \node[label1, right=1pt] at (pin) {$1$};
    \node[font=\footnotesize, text=black!78, align=center] at (1.5,-1.2)
      {No four points are shattered: the inner point lies in the hull of the others, so the\\labeling isolating it is unrealizable by a halfplane (Radon). Hence $\VC<4$, so $\VC=3$.};
  \end{scope}
\end{tikzpicture}
\caption[Shattering three points: each of the eight labelings realized by a separating halfplane, and no four points shattered, so the VC dimension is three.]{The eight labelings of three non-collinear points, each drawn with a separating halfplane that realizes it: the boundary splits the label-one points (red) from the label-zero points (blue), and the label-one side is lightly shaded.  A witness line in every cell means the restriction of the halfplane class to these three points is the full set of eight labelings, so the points are shattered (\leanname{Shatters}) and the VC dimension is at least three.  The lower panel shows why no four points are shattered: one of any four lies in the convex hull of the others, and the labeling isolating it is unrealizable by a halfplane (Radon), so the VC dimension (\leanname{VCDim}) equals three, which is $d+1$ for halfspaces in $d$ dimensions.  The shattering and VC-dimension concepts are the verified definitions; the value three and the four-point obstruction are classical.}
\label{fig:shatter-halfplane}
\end{figure}


\subsection{The Growth Function}

The growth function measures how many distinct labelings $\mathcal{H}$ can
produce on sets of a given size.

\begin{defn}[Growth Function]
\label{def:growth-function}
The \emph{growth function} (or \emph{shatter coefficient}) of $\mathcal{H}$ is
\[
  \growth_\mathcal{H}(n) \;=\; \max_{S \subseteq X,\; |S|=n}
  \bigl|\mathcal{H}|_S\bigr|.
\]
Equivalently, $\growth_\mathcal{H}(n)$ is the maximum number of distinct
behaviors $\mathcal{H}$ can exhibit on any $n$ points.\formal{FLT_Proofs.Complexity.VCDimension.GrowthFunction}
\end{defn}

The growth function satisfies $\growth_\mathcal{H}(n) \leq 2^n$ always, with
equality when $\mathcal{H}$ shatters some set of size~$n$.  The remarkable
content of the Sauer--Shelah lemma is that once shattering fails (at
$n = \VC(\mathcal{H}) + 1$), the growth function drops from exponential to
polynomial \emph{immediately and permanently}.


\subsection{The Sauer--Shelah Lemma}

\begin{lemma}[Sauer--Shelah {\cite{Sauer1972,Shelah1972}}]
\label{lem:sauer-shelah}
Let $\mathcal{H} \subseteq \{0,1\}^X$ with $\VC(\mathcal{H}) = d$.  Then for
all $n \geq 1$,
\[
  \growth_\mathcal{H}(n) \;\leq\; \sum_{i=0}^{d} \binom{n}{i}.
\]
In particular, $\growth_\mathcal{H}(n) \leq \bigl(\frac{en}{d}\bigr)^d$ for
$n \geq d \geq 1$, so the growth function is polynomial in $n$ when
$d < \infty$.\formal{FLT_Proofs.Theorem.PAC.sauer_shelah_quantitative}
\end{lemma}

The proof is by double induction and is one of the most elegant arguments in
combinatorics.  The key idea is to split the hypothesis class along one
coordinate and apply the inductive hypothesis to the two resulting classes,
whose VC dimensions are carefully controlled.

\begin{proof}
\label{lem:sauer-shelah-proof}
We prove a stronger statement: for any finite set $S$ with $|S| = n$ and
any class $\mathcal{H}' \subseteq \{0,1\}^S$ (not necessarily arising from
a restriction),
\[
  |\mathcal{H}'| \;\leq\; \bigl|\{A \subseteq S : \mathcal{H}' \text{ shatters } A\}\bigr|.
  \tag{$\star$}
\]
This implies the lemma: if $\VC(\mathcal{H}') \leq d$, then the right-hand
side counts only subsets $A$ with $|A| \leq d$, of which there are at most
$\sum_{i=0}^{d}\binom{n}{i}$.

We prove $(\star)$ by induction on $|S|$.

\smallskip\noindent
\textbf{Base case} ($|S| = 1$, say $S = \{x\}$).  If
$\mathcal{H}' = \{0\}$ or $\mathcal{H}' = \{1\}$, then $|\mathcal{H}'| = 1$
and $\mathcal{H}'$ shatters $\emptyset$, giving at least one shattered set.
If $\mathcal{H}' = \{0,1\}$, then $|\mathcal{H}'| = 2$ and $\mathcal{H}'$
shatters both $\emptyset$ and $\{x\}$, giving at least two shattered sets.
In both cases $(\star)$ holds.

\smallskip\noindent
\textbf{Inductive step.}  Let $|S| = n > 1$.  Pick any element $x \in S$
and set $S' = S \setminus \{x\}$.  Partition $\mathcal{H}'$ according to
behavior on $x$:
\begin{align*}
  \mathcal{H}_0 &= \{h|_{S'} : h \in \mathcal{H}',\; h(x) = 0\}, \\
  \mathcal{H}_1 &= \{h|_{S'} : h \in \mathcal{H}',\; h(x) = 1\}, \\
  \mathcal{H}_\cap &= \mathcal{H}_0 \cap \mathcal{H}_1.
\end{align*}
Here $\mathcal{H}_\cap$ consists of the functions on $S'$ that appear with
\emph{both} labels on $x$.

Count: each $h \in \mathcal{H}'$ contributes to exactly one of $\mathcal{H}_0$
or $\mathcal{H}_1$ (as a function on $S'$), but a function in
$\mathcal{H}_\cap$ is contributed by two elements of $\mathcal{H}'$ (one
with $h(x)=0$, one with $h(x)=1$).  Therefore,
\[
  |\mathcal{H}'| \;=\; |\mathcal{H}_0| + |\mathcal{H}_1| - |\mathcal{H}_\cap|
  + |\mathcal{H}_\cap|
  \;=\; |\mathcal{H}_0 \cup \mathcal{H}_1| + |\mathcal{H}_\cap|.
\]
More precisely, $|\mathcal{H}'| = |\mathcal{H}_0| + |\mathcal{H}_1|$, since
functions in $\mathcal{H}_0 \setminus \mathcal{H}_\cap$ appear only from
$h(x) = 0$, functions in $\mathcal{H}_1 \setminus \mathcal{H}_\cap$ appear
only from $h(x) = 1$, and functions in $\mathcal{H}_\cap$ appear twice.  So
\[
  |\mathcal{H}'| = |\mathcal{H}_0 \cup \mathcal{H}_1| + |\mathcal{H}_\cap|.
\]

Apply the inductive hypothesis (on $S'$, which has $n-1$ elements) to both
$\mathcal{H}_0 \cup \mathcal{H}_1$ and $\mathcal{H}_\cap$:
\begin{align*}
  |\mathcal{H}_0 \cup \mathcal{H}_1|
    &\leq \bigl|\{A \subseteq S' : (\mathcal{H}_0 \cup \mathcal{H}_1) \text{ shatters } A\}\bigr|, \\
  |\mathcal{H}_\cap|
    &\leq \bigl|\{A \subseteq S' : \mathcal{H}_\cap \text{ shatters } A\}\bigr|.
\end{align*}

Now we connect the shattered sets of the restricted classes back to those
of $\mathcal{H}'$:

\begin{enumerate}[leftmargin=*,itemsep=3pt]
\item If $A \subseteq S'$ is shattered by $\mathcal{H}_0 \cup \mathcal{H}_1$
  but \emph{not} by $\mathcal{H}_\cap$, then $A$ is shattered by $\mathcal{H}'$
  as well (since every labeling of $A$ is realized by some $h|_{S'} \in
  \mathcal{H}_0 \cup \mathcal{H}_1$, and at least one such $h$ exists in
  $\mathcal{H}'$).

\item If $A \subseteq S'$ is shattered by $\mathcal{H}_\cap$, then every
  labeling of $A$ is realized by functions appearing with both labels on $x$.
  Hence $\mathcal{H}'$ shatters $A \cup \{x\}$: for any labeling of
  $A \cup \{x\}$, pick the label $b$ assigned to $x$, then pick $h|_{S'} \in
  \mathcal{H}_\cap$ realizing the labeling of $A$; the corresponding
  $h \in \mathcal{H}'$ with $h(x) = b$ realizes the full labeling.
\end{enumerate}

The sets in (1) do not contain $x$, and the sets in (2) do
contain $x$, so they are disjoint families.  Therefore,
\[
  |\mathcal{H}'|
  \;\leq\; \bigl|\{A \subseteq S : \mathcal{H}' \text{ shatters } A,\;
  x \notin A\}\bigr|
  + \bigl|\{A \subseteq S : \mathcal{H}' \text{ shatters } A,\;
  x \in A\}\bigr|
  \;=\; \bigl|\{A \subseteq S : \mathcal{H}' \text{ shatters } A\}\bigr|.
\]
This completes the induction.
\end{proof}

\begin{remark}
The lemma was proved independently by Sauer~\cite{Sauer1972} and
Shelah~\cite{Shelah1972}, with a related result by Perles and Shelah
appearing earlier.  Vapnik and Chervonenkis~\cite{VapnikChervonenkis1971}
proved a weaker bound in their 1971 paper.  The proof above follows the
elegant ``number of shattered sets'' formulation due to Alon and
Frankl.
\end{remark}

\begin{figure}[ht]
\centering
\begin{tikzpicture}
\begin{axis}[
    width=12cm, height=7cm,
    xlabel={$n$ (number of points)},
    ylabel={Number of labelings $\growth_\mathcal{H}(n)$},
    ymode=log, log basis y=10,
    xmin=0, xmax=16, ymin=1, ymax=1e5,
    xtick={0,2,4,6,8,10,12,14,16},
    ytick={1,10,100,1000,10000,100000},
    legend pos=south east,
    legend cell align={left},
    legend style={font=\small, fill=white, fill opacity=0.9, draw=gray!40, text opacity=1},
    grid=major,
    grid style={gray!20},
    thick,
    clip mode=individual,
]
% 2^n reference ray: the ceiling, attained while the class still shatters (dashed)
\addplot[domain=0:16, samples=120, thick, thmred, dashed] {2^x};
\addlegendentry{$2^n$ (shatters forever)}

% Actual growth of a maximum VC=3 class = the Sauer--Shelah bound:
% equals 2^n for n<=3 (binomials i>n vanish), then peels permanently below (solid, integer marks)
\addplot[domain=0:16,
         samples at={0,1,2,3,4,5,6,7,8,9,10,11,12,13,14,15,16},
         thick, defblue, mark=*, mark size=1pt]
    {1 + x + x*(x-1)/2 + x*(x-1)*(x-2)/6};
\addlegendentry{$\growth_\mathcal{H}(n)=\sum_{i=0}^{3}\binom{n}{i}$ \;($\VC=3$)}

% Marker: n = d = 3 (last shattered size)
\draw[thick, dashed, notegray] (axis cs:3,1) -- (axis cs:3,1e5);
\node[font=\scriptsize, text=notegray, anchor=south, rotate=90]
    at (axis cs:3,5e2) {$n=d=3$ (last shattered)};

% Marker: n = d+1 = 4 (Sauer--Shelah bound bites; permanent drop)
\draw[thick, dashed, notegray] (axis cs:4,1) -- (axis cs:4,1e5);
\node[font=\scriptsize, text=notegray, anchor=south, rotate=90]
    at (axis cs:4,5e2) {$n=d{+}1$: bound bites};

% The widening exponential gap for n > d
\draw[<->, thin, gray!70] (axis cs:11,2048) -- (axis cs:11,232);
\node[font=\scriptsize, text=gray!75, anchor=west]
    at (axis cs:11.3,720) {exponential saving};

\end{axis}
\end{tikzpicture}
\caption[The growth function: exponential while the class shatters, then polynomial of degree~$\VC$.]{The
growth function $\growth_\mathcal{H}(n)$ (\leanname{GrowthFunction}) counts the most
distinct labelings a class can produce on any $n$ points.  It never exceeds $2^n$ (the
dashed ray) and equals $2^n$ exactly while the class still shatters a set of that size,
that is for $n \leq d = \VC(\mathcal{H})$.  The solid curve is the largest a class of VC
dimension three can grow, the Sauer--Shelah bound $\sum_{i=0}^{3}\binom{n}{i}$
(\leanname{sauer_shelah_quantitative}): it rides the $2^n$ ray through $n=3$, then at
$n = d+1 = 4$ peels permanently away into a polynomial arc whose degree is the VC
dimension.  The drop is immediate and there is no intermediate growth rate.  On the
logarithmic vertical axis the exponential $2^n$ is a straight line and the polynomial
bends below it, so the widening gap is the exponential saving a finite VC dimension buys.
This collapse is the combinatorial step the Fundamental Theorem turns into PAC
learnability, since the growth function is the quantity a uniform-convergence bound
controls.  The definition and the bound are both verified; the asymptotic form
$(en/d)^d$ and the no-middle-regime reading are classical.}
\label{fig:growth-function}
\end{figure}

The Sauer--Shelah lemma is the combinatorial engine behind the Fundamental
Theorem of Statistical Learning (\Cref{thm:fundamental}): it converts
a finite VC dimension into a polynomial bound on the growth function, which in
turn yields uniform convergence of empirical risks via a symmetrization
argument.


% ================================================================
% SECTION 2: THE LITTLESTONE DIMENSION (CROSS-REFERENCE)
% ================================================================

\section{The Littlestone Dimension}
\label{sec:ldim-recall}

The VC dimension measures shatterability of \emph{sets}: it asks for the
largest set whose every labeling is realized.  The Littlestone dimension,
introduced in \Cref{ch:online} and fully characterized there, measures
shatterability of \emph{trees}: it asks for the deepest complete binary tree
whose every root-to-leaf path is realized.  We recall the definition for
reference.

\begin{defn}[Mistake Tree, Recall]
\label{def:mistake-tree-recall}
A \emph{mistake tree} for $\mathcal{H} \subseteq \{0,1\}^X$ is a complete
binary tree in which each internal node is labeled with an instance $x \in X$,
left edges correspond to label~$0$, right edges to label~$1$, and every
root-to-leaf path is consistent with some $h \in \mathcal{H}$
(see \Cref{def:mistake-tree}).
\end{defn}

\begin{defn}[Littlestone Dimension, Recall]
\label{def:ldim-recall}
The \emph{Littlestone dimension} $\Ldim(\mathcal{H})$ is the maximum depth
of a mistake tree for $\mathcal{H}$ (\Cref{def:ldim}).  The
Littlestone characterization (\Cref{thm:littlestone-characterization}) establishes that
$\Ldim(\mathcal{H})$ equals the optimal worst-case mistake bound.
\end{defn}

The relationship between the two fundamental dimensions is strict:

\begin{prop}
\label{prop:vc-leq-ldim-dim}
For any hypothesis class $\mathcal{H}$,
$\VC(\mathcal{H}) \leq \Ldim(\mathcal{H})$.\formal{FLT_Proofs.Complexity.Littlestone.BranchWiseLittlestoneDim_ge_VCDim}
\end{prop}

\begin{proof}
If $\mathcal{H}$ shatters $\{x_1, \ldots, x_d\}$, one can construct a
mistake tree of depth $d$ by placing $x_i$ at every node of depth $i$:
since every labeling is realized, every root-to-leaf path is consistent with
some $h \in \mathcal{H}$.  Hence any shattered set yields a mistake tree of
the same depth.
\end{proof}

The converse fails: threshold classifiers on $\R$ have $\VC = 1$ but
$\Ldim = \infty$ (\Cref{ch:separations}).  Intuitively, the Littlestone
dimension is larger because it accounts for adaptive adversaries: the
adversary in the online game chooses instances sequentially, seeing the
learner's responses, while the VC dimension considers all labelings of a
fixed set simultaneously.


% ================================================================
% SECTION 3: BEYOND BINARY
% ================================================================

\section{Beyond Binary: Pseudodimension and Fat-Shattering}
\label{sec:beyond-binary}

When the hypothesis class consists of real-valued functions
$\mathcal{F} \subseteq \R^X$ rather than binary classifiers, shattering must
be redefined.  Two progressively refined notions capture the relevant
complexity.

\begin{defn}[Pseudodimension]
\label{def:pseudodim}
Let $\mathcal{F} \subseteq \R^X$.  A set $S = \{x_1, \ldots, x_n\} \subseteq X$
is \emph{pseudo-shattered} (or \emph{P-shattered}) by $\mathcal{F}$ if there
exist thresholds $t_1, \ldots, t_n \in \R$ such that, for every labeling
$b \in \{0,1\}^n$, there exists $f \in \mathcal{F}$ with
\[
  \mathrm{sgn}\bigl(f(x_i) - t_i\bigr) = (-1)^{1 - b_i}
  \quad \text{for all } i = 1, \ldots, n.
\]
The \emph{pseudodimension} $\Pdim(\mathcal{F})$ is the largest cardinality of
a P-shattered set.\formal{FLT_Proofs.Complexity.Structures.Pseudodimension}
\end{defn}

The thresholds $t_i$ act as ``witnesses'': they reduce a real-valued
shattering problem to a binary one.  When $\mathcal{F}$ is already
$\{0,1\}$-valued, the thresholds $t_i = 1/2$ recover the standard VC
dimension, so $\Pdim$ is a proper generalization.

\begin{defn}[Fat-Shattering Dimension {\cite{AlonBenDavidCesaBianchiHaussler1997}}]
\label{def:fat-shattering}
Let $\mathcal{F} \subseteq \R^X$ and $\gamma > 0$.  A set
$S = \{x_1, \ldots, x_n\}$ is \emph{$\gamma$-fat-shattered} by $\mathcal{F}$
if there exist thresholds $t_1, \ldots, t_n \in \R$ such that, for every
labeling $b \in \{0,1\}^n$, there exists $f \in \mathcal{F}$ with
\[
  f(x_i) \;\geq\; t_i + \gamma \;\;\text{when } b_i = 1,
  \qquad
  f(x_i) \;\leq\; t_i - \gamma \;\;\text{when } b_i = 0,
\]
for all $i$.  The \emph{fat-shattering dimension at scale $\gamma$}, denoted
$\fatshat_\gamma(\mathcal{F})$, is the largest cardinality of a
$\gamma$-fat-shattered set.\formal{FLT_Proofs.Complexity.Structures.FatShatteringDim}
\end{defn}

The margin parameter $\gamma$ makes fat-shattering \emph{scale-sensitive}:
large margins are harder to achieve, so
$\fatshat_\gamma(\mathcal{F}) \leq \fatshat_{\gamma'}(\mathcal{F})$ whenever
$\gamma \geq \gamma'$.  As $\gamma \to 0$, the fat-shattering dimension
approaches the pseudodimension.

\begin{thm}[Characterization of Real-Valued Learnability
  {\cite{AlonBenDavidCesaBianchiHaussler1997}}]
\label{thm:fat-char}
A real-valued function class $\mathcal{F} \subseteq [0,1]^X$ is
\emph{agnostically PAC learnable} (with respect to the absolute loss) if and
only if $\fatshat_\gamma(\mathcal{F}) < \infty$ for every $\gamma > 0$.
\end{thm}

This is the real-valued analogue of the Fundamental Theorem: just as finite VC
dimension characterizes binary PAC learnability, finite fat-shattering
dimension at every scale characterizes real-valued learnability.  The
pseudodimension alone is not sufficient for characterization in the agnostic
setting, as it ignores the scale structure.

\begin{remark}[Dimension hierarchy]
For a class $\mathcal{F} \subseteq [0,1]^X$:
\[
  \fatshat_\gamma(\mathcal{F})
  \;\leq\; \Pdim(\mathcal{F})
  \;\leq\; \VC\!\bigl(\mathrm{subgraph}(\mathcal{F})\bigr)
\]
where $\mathrm{subgraph}(\mathcal{F}) = \{(x,t) \mapsto \ind[f(x) \geq t] :
f \in \mathcal{F}\}$.  When $\mathcal{F}$ is $\{0,1\}$-valued, all three
quantities coincide with $\VC(\mathcal{F})$.
\end{remark}


% ================================================================
% SECTION 4: THE MULTICLASS STORY
% ================================================================

\section{The Multiclass Story: From Natarajan to DS}
\label{sec:multiclass}

Binary classification has a clean answer: a class $\mathcal{H} \subseteq
\{0,1\}^X$ is PAC learnable if and only if $\VC(\mathcal{H}) < \infty$.
What is the corresponding answer when the label set $Y$ has $k \geq 3$
elements?  This question, first posed in the 1980s, was open for thirty
years.  Its resolution required not merely a new proof technique but a
fundamentally new combinatorial dimension, because the obvious candidate turned out to be wrong.


\subsection{The Natarajan Dimension: The Obvious Candidate}

\begin{defn}[Natarajan Dimension]
\label{def:natarajan}
Let $\mathcal{H} \subseteq Y^X$ with $|Y| \geq 2$.  The \emph{Natarajan
dimension} $\Ndim(\mathcal{H})$ is the largest $d$ such that there exist
\begin{itemize}[leftmargin=*,itemsep=1pt]
\item a set $S = \{x_1, \ldots, x_d\} \subseteq X$, and
\item two functions $f_0, f_1 \colon S \to Y$ with $f_0(x_i) \neq f_1(x_i)$
  for all $i$,
\end{itemize}
such that for every labeling $b \in \{0,1\}^d$, there exists $h \in
\mathcal{H}$ with $h(x_i) = f_{b_i}(x_i)$ for all $i$.\formal{FLT_Proofs.Complexity.Structures.NatarajanDim}
\end{defn}

The Natarajan dimension generalizes the VC dimension in a natural way: instead
of requiring all $2^d$ binary labelings, it requires all $2^d$ labelings
drawn from two specified ``opposing'' functions.  When $|Y| = 2$, the
Natarajan dimension equals the VC dimension exactly.

In the 1980s and 1990s, the working conjecture was:
\begin{quote}
\itshape A multiclass hypothesis class $\mathcal{H} \subseteq Y^X$ is PAC
learnable if and only if $\Ndim(\mathcal{H}) < \infty$.
\end{quote}
This conjecture was supported by partial results.  Ben-David, Cesa-Bianchi,
Haussler, and Long proved that finite Natarajan dimension is
\emph{sufficient} for multiclass PAC learnability when combined with a bound
on $|Y|$.  The ``if'' direction seemed solid.  The conjecture survived for
decades without a counterexample.

It was wrong.


\subsection{The Counterexample: When the Obvious Fails}

\begin{thm}[Failure of Natarajan Dimension
  {\cite{DanielyShalevShwartz2014}}]
\label{thm:natarajan-failure}
There exists a hypothesis class $\mathcal{H} \subseteq Y^X$ with
$\Ndim(\mathcal{H}) = 1$ that is \emph{not} PAC learnable.
\end{thm}

This is a dramatic failure: the Natarajan dimension can be as small as
possible (just one) and still the class can be unlearnable.  The problem lies
in the label set: the construction uses $|Y| = \infty$, and the Natarajan
dimension, which examines only pairs of opposing functions, fails to capture
the combinatorial explosion that arises when the label set is large.

The witness is a class constructed via a combinatorial argument.  The key
insight is that the Natarajan dimension constrains the number of ``binary
slices'' through the labeling space but does not constrain the total
complexity when the label set grows without bound.

\begin{separation}
\textbf{Natarajan vs.\ PAC learnability.}  The Natarajan dimension does not
characterize multiclass PAC learnability:
\[
  \Ndim(\mathcal{H}) < \infty
  \quad\not\!\!\!\implies\quad
  \mathcal{H} \text{ is PAC learnable}.
\]
The gap is witnessed by a class with $\Ndim = 1$, $|Y| = \infty$, that
requires unbounded sample complexity.  The correct characterization requires
the DS dimension (\Cref{def:ds-dim} below).
\end{separation}


\subsection{The DS Dimension: The Correct Answer}

\begin{defn}[DS Dimension {\cite{BrukhimCarmonDinurMoranYehudayoff2022}}]
\label{def:ds-dim}
Let $\mathcal{H} \subseteq Y^X$.  The \emph{DS dimension} (named for
Daniely and Shalev-Shwartz) $\DSdim(\mathcal{H})$ is defined via
pseudo-cubes in the one-inclusion hypergraph of $\mathcal{H}$.

Construct the \emph{one-inclusion hypergraph} $G_\mathcal{H}$ on a finite set
$S = \{x_1, \ldots, x_n\} \subseteq X$ as follows: the vertices are the
hypotheses $\mathcal{H}|_S$ (restrictions of $\mathcal{H}$ to $S$), and for
each $i \in \{1, \ldots, n\}$, add a hyperedge connecting all hypotheses
that agree on $S \setminus \{x_i\}$ (i.e., they differ only at $x_i$).

A \emph{pseudo-cube of dimension $d$} in $G_\mathcal{H}$ is a subhypergraph
isomorphic to the complete $d$-dimensional hypercube graph $\{0,1\}^d$.

The DS dimension $\DSdim(\mathcal{H})$ is the supremum over all finite
$S \subseteq X$ of the maximum dimension of a pseudo-cube in $G_\mathcal{H}$.\footnote{The genuine DS dimension is not machine-checked. The kernel types a related but distinct notion: the \emph{all-labelings} (strong shattering) dimension \leanname{FLT_Proofs.Complexity.Structures.StrongShatteringDim}, the largest $S$ on which every $f \colon S \to Y$ is realizable, whose value is at most $\Ndim$, the reverse of the genuine DS relationship $\Ndim \le \DSdim$. Formalizing the pseudo-cube DS dimension is \cref{op:ds-formal}.}
\end{defn}

\begin{thm}[Multiclass Characterization
  {\cite{BrukhimCarmonDinurMoranYehudayoff2022}}]
\label{thm:ds-char}
A hypothesis class $\mathcal{H} \subseteq Y^X$ is PAC learnable if and only
if $\DSdim(\mathcal{H}) < \infty$.
\end{thm}

This theorem, proved by Brukhim, Chepoi, Daniel, Moran, and Yehudayoff in
2022, resolved the multiclass characterization problem.  It is the
multiclass analogue of the Fundamental Theorem: the DS dimension plays
exactly the role that the VC dimension plays for binary classification.

\begin{historical}
\textbf{Why the Natarajan dimension fails and the DS dimension succeeds.}
The Natarajan dimension examines the labeling structure of $\mathcal{H}$ by
projecting onto pairs of opposing functions, a fundamentally binary lens
applied to a multiclass problem.  The DS dimension, by contrast, examines
the full combinatorial structure of the one-inclusion hypergraph, which
encodes all the ways hypotheses can disagree on single points.

The one-inclusion hypergraph is not new: Haussler, Littlestone, and Warmuth
introduced it in 1994 for binary classification, where it yields an optimal
PAC learner (the one-inclusion algorithm).  What Brukhim et al.\ discovered
is that the right way to measure the complexity of a multiclass hypothesis
class is to look for high-dimensional pseudo-cubes inside this hypergraph.
Pseudo-cubes are the multiclass analogue of shattered sets: they represent
configurations where the labeling structure is rich enough to prevent
learning.
\end{historical}

The dimension hierarchy for multiclass problems is:

\begin{figure}[ht]
\centering
\begin{tikzpicture}[
    dimbox/.style={draw, thick, rounded corners=3pt, minimum width=2.6cm,
      minimum height=0.95cm, font=\small\bfseries, align=center},
    dimboxdash/.style={dimbox, dashed},
    ladsolid/.style={-{Stealth[length=2.5mm]}, thick},
    laddash/.style={-{Stealth[length=2.5mm]}, thick, dashed, draw=notegray},
    branch/.style={thick, defblue},
    solidcall/.style={thick, draw=defblue},
    dashcall/.style={thick, dashed, draw=notegray},
    failcall/.style={thick, dashed, draw=thmred},
    implabel/.style={font=\scriptsize, text=notegray, align=center},
    gaplabel/.style={font=\scriptsize\itshape, text=thmred, align=center},
]
  % --- the magnitude ladder spine: verified floor -> literature ceiling ---
  \node[dimbox, fill=teal!14] (sst) at (0,0) {Strong\\shattering\\{\scriptsize dim.}};
  \node[dimbox, fill=edgeorange!14] (nat) at (4,0) {$\Ndim$\\{\scriptsize Natarajan}};
  \node[dimboxdash, fill=notegray!8] (ds) at (8,0) {$\DSdim$\\{\scriptsize DS / pseudo-cube}};
  % --- VC: the binary face of Ndim, above the spine ---
  \node[dimbox, fill=defblue!14, minimum width=2.2cm] (vc) at (4,2.0) {$\VC$\\{\scriptsize binary}};

  % --- ladder edges: solid = kernel-verified, dashed = literature ---
  \draw[ladsolid] (sst) -- node[implabel, above] {$\le$} (nat);
  \draw[laddash] (nat) -- node[implabel, above] {$\le$} (ds);
  \draw[branch] (vc) -- node[implabel, right, text=defblue] {\scriptsize $|Y|{=}2$:\ $\Ndim{=}\VC$} (nat);

  % --- characterization call-outs (each verdict hangs off its node) ---
  \node[implabel, text=defblue, above=0.55cm of vc] (vcnote)
    {$\VC<\infty \Leftrightarrow$ learnable\\{\scriptsize binary, verified}};
  \draw[solidcall] (vc) -- (vcnote);

  \node[gaplabel, below=0.6cm of nat] (gap)
    {$\Ndim<\infty \not\Rightarrow$ learnable\\{\scriptsize witness $\Ndim=1$, $|Y|=\infty$}\\{\scriptsize (Daniely--Shalev-Shwartz 2014)}};
  \draw[failcall] (nat) -- (gap);

  \node[implabel, below=0.6cm of ds] (dscorrect)
    {$\DSdim<\infty \Leftrightarrow$ learnable\\{\scriptsize (Brukhim et al.\ 2022,}\\{\scriptsize not formalized)}};
  \draw[dashcall] (ds) -- (dscorrect);

  % --- line-style legend: the evidence-grade key ---
  \begin{scope}[shift={(-0.3,-3.75)}]
    \draw[ladsolid] (0,0) -- (0.7,0);
    \node[font=\scriptsize, anchor=west] at (0.85,0) {kernel-verified (sorry-free)};
    \draw[laddash] (5.1,0) -- (5.8,0);
    \node[font=\scriptsize, anchor=west] at (5.95,0) {literature / not yet formalized};
  \end{scope}
\end{tikzpicture}
\caption[The multiclass dimension ladder: verified rungs versus literature.]{The
multiclass dimension ladder, drawn at its evidence grade.  For binary classes
($|Y|=2$) the VC dimension characterizes PAC learnability completely: a class is
learnable if and only if its VC dimension is finite, and this equivalence is
machine-checked (solid).  The obvious generalization to $k \ge 3$ labels is the
Natarajan dimension (\leanname{NatarajanDim}), which equals the VC dimension when
$|Y|=2$; but finite Natarajan dimension does \emph{not} imply learnability
(Daniely--Shalev-Shwartz 2014, with a witness of Natarajan dimension one over an
infinite label set), so the obvious dimension fails to characterize the multiclass
regime.  The genuine characterization is the DS dimension, defined through
pseudo-cubes in the one-inclusion hypergraph (Brukhim et al.\ 2022); it is correct
but not yet formalized, so it and its equivalence are drawn dashed.  In its place the
kernel verifies the all-labelings strong-shattering dimension
(\leanname{StrongShatteringDim}) and the inequality
\leanname{strongShatteringDim_le_NatarajanDim}, which is the \emph{reverse} of the
genuine relationship $\Ndim \le \DSdim$.  Hence the full ladder
$\mathrm{StrongShattering} \le \Ndim \le \DSdim$ has only its lower rung
machine-checked: solid edges are kernel-verified, dashed edges are literature or not
yet formalized.}
\label{fig:dimension-hierarchy}
\end{figure}


% ================================================================
% SECTION 5: OTHER DIMENSIONS
% ================================================================

\section{Other Dimensions}
\label{sec:other-dims}

Several further combinatorial dimensions appear in the concept graph.  We
collect them here in catalog form; each is developed more fully in the chapter
where it plays a characterizing role.

\begin{table}[ht]
\centering
\small
\caption{Catalog of combinatorial dimensions beyond VC, Littlestone,
pseudodimension, fat-shattering, Natarajan, and DS.}
\label{tab:other-dims}
\begin{tabularx}{\textwidth}{@{}l l X l@{}}
\toprule
\textbf{Dimension} & \textbf{Notation} & \textbf{What it measures} & \textbf{Reference} \\
\midrule
Star number
  & $\mathfrak{s}(\mathcal{H})$
  & The largest set on which each point can be isolated by some concept.
    It satisfies $\mathfrak{s}(\mathcal{H}) \geq \VC(\mathcal{H})$, so finite
    star number implies finite VC, but \emph{not} conversely (singletons:
    $\mathfrak{s} = \infty$, $\VC = 1$).  It governs active-learning label
    complexity.
  & \Cref{ch:objects} \\
\addlinespace
Eluder dimension
  & $\dim_E(\mathcal{F})$
  & The length of the longest sequence of points such that each point is
    ``independent'' of the previous ones, given the function class.
    Measures exploration complexity in sequential decision problems.
  & \Cref{ch:extensions} \\
\addlinespace
SQ dimension
  & $\SQdim(\mathcal{C})$
  & The largest set of nearly uncorrelated concepts.  Determines the number
    of statistical queries needed to learn $\mathcal{C}$ in the statistical
    query model, where the learner receives expectations
    $\E[\phi(x,y)]$ rather than individual examples.
  & \Cref{ch:extensions} \\
\addlinespace
VCL tree
  & (none)
  & A tree structure used in the trichotomy theorem (\Cref{ch:universal})
    to separate exponential from arbitrarily slow learners.  Combines VC
    and Littlestone structure.
  & \Cref{ch:universal} \\
\addlinespace
Ordinal VC dim
  & $\VC^\alpha$
  & A transfinite refinement of VC dimension, using ordinal-indexed
    Cantor--Bendixson ranks on the space of consistent hypotheses.
  & \Cref{ch:ordinals} \\
\addlinespace
Ordinal Littlestone dim
  & $\Ldim^\alpha$
  & A transfinite refinement of Littlestone dimension, paralleling
    ordinal VC dimension for the online setting.
  & \Cref{ch:ordinals} \\
\bottomrule
\end{tabularx}
\end{table}

\begin{remark}[Star number versus VC dimension]
\label{rmk:star-number}
It is tempting to believe that finite star number and finite VC dimension
coincide.  They do not.  The star number dominates the VC dimension,
$\mathfrak{s}(\mathcal{H}) \geq \VC(\mathcal{H})$ (\Cref{prop:star-ge-vc}),
because a shattered set is in particular a set on which each point can be
isolated; but the reverse fails, and fails infinitely.  The class of
singletons on an infinite domain has $\VC = 1$ yet $\mathfrak{s} = \infty$
(\Cref{rmk:star-not-equiv}).\formal{FLT_Proofs.Complexity.VCDimension.StarNumber} Finiteness of the star number is therefore
\emph{strictly stronger} than PAC learnability, not equivalent to it; its
proper role is in active learning, where it controls the label
complexity~\cite{HannekeYang2015}.
\end{remark}


% ================================================================
% CLOSING REMARKS
% ================================================================

\bigskip
\begin{center}
\rule{0.5\textwidth}{0.4pt}
\end{center}
\bigskip

The combinatorial dimensions of this chapter are the coordinates of a map.
The VC dimension locates a class in the PAC landscape; the Littlestone
dimension locates it in the online landscape; the DS dimension locates it in
the multiclass landscape; the fat-shattering dimension locates it in the
real-valued landscape.  No single dimension suffices for all purposes, because
the landscapes are genuinely different, as the separation results of
\Cref{ch:separations} will make precise.

The next chapter turns from combinatorial dimensions to a different family of
complexity measures: those based on \emph{compression} and \emph{sample
complexity}.

\section{The Dimensions of an Attention Class}
\label{sec:attention-dimensions}

The combinatorial dimensions of this chapter take concrete values on the attention classes of a transformer, and those values control its sample complexity, its online learnability, and its learning rate.  Fix a finite-head attention class $\mathcal{A}_{d,n_K}$ on $\R^d$ with $n_K$ keys.

\begin{prop}[The VC dimension of attention is finite]
\label{prop:attention-vc-dim}
A fixed-architecture soft-attention class has finite VC dimension, so the Sauer--Shelah bound (\cref{lem:sauer-shelah}) applies: its growth function is polynomial.\formal{FLT_Proofs.Complexity.VCDimension.VCDim}
\end{prop}

\begin{proof}
The class is a well-behaved VC-measurable target (\cref{prop:attention-pac}), and its concepts are cut out by a bounded number of smooth inequalities in finitely many parameters; such parametric families shatter no arbitrarily large set, so $\VC(\mathcal{A}_{d,n_K}) < \infty$.  Sauer--Shelah then bounds $\growth_{\mathcal{A}}(n) \leq (en/\VC)^{\VC}$.
\end{proof}

The other dimensions of this chapter then read off the rest of the attention class's profile.  Its \emph{Littlestone} dimension (finite or infinite) decides whether the class is online learnable and whether its universal rate is exponential or linear (\Cref{ch:universal}); finiteness is open (\cref{op:attention-dim-profile}).  Its \emph{star number} governs how cheaply an active learner can extract the head's behaviour by querying, the methodology of \Cref{ch:exact}.  And when attention routes over a label set of size $n_K \geq 3$ (a genuine multiclass head), its \emph{DS} dimension, not its Natarajan dimension, is what decides learnability, by the same thirty-year correction told above.  The single architecture is thus a point in every one of this chapter's dimension spaces at once, and the spaces disagree about it in exactly the ways the separation results of \Cref{ch:separations} predict.

\section{Open Problems}
\label{sec:ch10-open}

Each problem is anchored to a result above and tagged: a \emph{known unknown} (\textsc{ku}) is an articulated open question; an \emph{unknown known} (\textsc{uk}) is a result in the literature or the verified development not yet absorbed into a clean statement.  Verified handles are named where they exist.

\begin{openprob}[The dimension profile of attention, \textsc{uk}]
\label{op:attention-dim-profile}
\Cref{prop:attention-vc-dim} pins the VC dimension of a finite-head attention class as finite.  Compute the full dimension profile (VC, Littlestone, star number, and, for multiclass heads, DS) as functions of the architecture and score temperature, settling the regime questions of \Cref{ch:online,ch:universal} for attention in one stroke.
\end{openprob}

\begin{openprob}[The optimal compression of VC classes, \textsc{ku}]
\label{op:od-compression-dim}
The Sauer--Shelah lemma (\cref{lem:sauer-shelah}) converts finite VC dimension into a polynomial growth function.  Determine whether every class of VC dimension $d$ admits a sample compression scheme of size $O(d)$ (Littlestone--Warmuth), the open conjecture isolated in \leanname{FLT_Proofs.Complexity.Structures.Open_NoInfoCompressionStrengthening} and recurring from \Cref{ch:objects,ch:compression}.
\end{openprob}

\begin{openprob}[A combinatorial proof that DS characterizes multiclass learning, \textsc{uk}]
\label{op:ds-formal}
The DS characterization (\cref{thm:ds-char}) is cited, not verified.  The development types the Natarajan dimension and an \emph{all-labelings} shattering dimension (\leanname{StrongShatteringDim}, whose value is at most $\Ndim$) but \emph{not} the genuine Daniely--Shalev-Shwartz DS dimension, whose pseudo-cube definition gives the reverse bound $\Ndim \le \DSdim$.  Formalize the true DS dimension and the Brukhim et al.\ proof that $\DSdim(\mathcal{H}) < \infty$ iff $\mathcal{H}$ is multiclass PAC learnable.
\end{openprob}

\begin{openprob}[Closing the Natarajan--DS gap quantitatively, \textsc{ku}]
\label{op:natarajan-ds-gap}
\Cref{thm:natarajan-failure} shows $\Ndim$ can be $1$ while the class is unlearnable ($\DSdim = \infty$).  Characterize exactly how large the DS dimension can be as a function of the Natarajan dimension and the label-set size $|Y|$, and identify the structural feature of the label set that makes the gap appear.
\end{openprob}

\begin{openprob}[Fat-shattering rates for attention, \textsc{uk}]
\label{op:fat-attention}
Soft attention outputs real-valued scores before thresholding.  Compute the fat-shattering dimension $\fatshat_\gamma$ (\cref{def:fat-shattering}) of the real-valued attention-score class as a function of the margin $\gamma$ and the temperature, giving margin-based generalization bounds for attention via \cref{thm:fat-char}.
\end{openprob}

\begin{openprob}[A unifying dimension functor, \textsc{ku}]
\label{op:dimension-functor}
The chapter catalogs a dozen dimensions, each governing one paradigm.  Determine whether they are values of a single functor from concept classes (with label-preserving maps) to ordinals, recovering VC, Littlestone, DS, and fat-shattering by specializing a data-model and a scale parameter, or prove the dimensions are provably non-unifiable.
\end{openprob}

\begin{openprob}[The exact star-number--VC gap, \textsc{ku}]
\label{op:star-vc-gap}
\Cref{rmk:star-number} corrects the folklore: $\mathfrak{s} \geq \VC$ with an infinite gap for singletons.  Characterize the classes for which $\mathfrak{s}(\mathcal{H}) = O(\VC(\mathcal{H}))$, and determine the active-learning label-complexity speed-up that finiteness of the star number buys over the worst case, on top of \leanname{FLT_Proofs.Complexity.VCDimension.StarNumber}.
\end{openprob}

\begin{openprob}[Pseudodimension versus fat-shattering for neural classes, \textsc{uk}]
\label{op:pdim-fat}
\Cref{thm:fat-char} shows the pseudodimension alone does not characterize agnostic real-valued learnability; the scale-sensitive fat-shattering dimension does.  Determine, for the score class of a transformer, how far the pseudodimension over-estimates the learnable complexity, and whether the gap shrinks with depth.
\end{openprob}

\begin{openprob}[Ordinal dimensions and the slow regime, \textsc{uk}]
\label{op:ordinal-dim-slow}
\Cref{tab:other-dims} lists the ordinal VC and Littlestone dimensions (\leanname{FLT_Proofs.Complexity.Ordinal.OrdinalLittlestoneDim}).  Determine whether a single ordinal-valued dimension characterizes the boundary between the linear and arbitrarily-slow universal rates of \Cref{ch:universal}, completing the dimension story of the trichotomy.
\end{openprob}

\begin{openprob}[The dual VC dimension and attention routing, \textsc{uk}]
\label{op:dual-vc-attention}
The development types the dual class and proves a dual-to-primal shattering transfer (\leanname{FLT_Proofs.Complexity.DualVC.dual_shatters_imp_original_shatters}).  Interpret the dual VC dimension of an attention class (shattering by the \emph{keys} rather than the inputs) and determine what it measures about a transformer's routing capacity.
\end{openprob}

\subsection*{Counterfactual exercises: generalizing Figure~\ref{fig:shatter-halfplane}}

Each exercise perturbs the shattering picture.  Argue, in the figure's grammar (a $2^d$-cell grid in which each cell carries one labeling of a fixed $d$-point set together with the separating halfplane that realizes it, so that filling every cell with a witness exhibits the restriction $\mathcal{H}|_S$ as the full cube $\{0,1\}^d$, drawn solid where the kernel certifies the concept (the definitions \leanname{Shatters} and \leanname{VCDim} and the deep role \leanname{vc_characterization}) and dashed where the \emph{value} $\VC=3$, its generalization $\VC=d+1$, and the four-point obstruction are classical, cited not proved), that the perturbed picture subsumes the concrete one as an \emph{instance} (a class whose largest fillable grid is exhibited), a \emph{boundary} (one labeling or one point toggled so a cell becomes unrealizable), or a \emph{frontier} (a dashed, open or literature-only edge).

\begin{enumerate}[leftmargin=*,itemsep=3pt]
\item \textbf{The grid is the definition of shattering, not a fact about the plane.} Read each filled cell not as a colored triangle but as one element of the restriction. Show this is the figure's defining instance of the verified concept: a witness in every cell says exactly $\mathcal{H}|_S = \{0,1\}^{|S|}$, which is the definition of shattering (\cref{def:shattering}, \leanname{Shatters}), so the eight halfplanes are not decoration but the constructive content of ``$S$ is shattered''; the picture certifies \leanname{Shatters} cell by cell rather than asserting it.

\item \textbf{The value three is the largest fillable grid, and it is literature.} Read the whole figure as a single number. Show this is the instance that names the second object: $\VC(\mathcal{H})$ is the largest $|S|$ for which the grid fills (\cref{def:vc-dim}, \leanname{VCDim}), so the figure exhibits $\VC(\mathcal{H}) \ge 3$ by construction, but the matching upper bound $\VC = 3$ is the classical value (\cite{VapnikChervonenkis1971}, \cite{BlumerEhrenfeuchtHausslerWarmuth1989}), drawn dashed because the kernel types \leanname{VCDim} without proving the halfplane value; the concept is verified, the number is cited.

\item \textbf{The fourth point empties a cell.} Add a fourth point inside the convex hull of the other three and ask for the labeling that flips it against the rest. Show this is the figure's principal boundary: that single cell of the $2^4$ grid can never carry a separating line, because one point in the hull of the others forbids the opposite-label halfplane (Radon, \Cref{ex:halfplanes-shatter}), so the grid for $d=4$ has a permanently empty cell and $\VC = 3$; the obstruction is convex-geometric and classical, with no kernel theorem, so it is drawn dashed.

\item \textbf{Why filling the grid is what makes the class learnable.} Ask why the largest-fillable-grid number governs anything. Show this is the frontier that gives the figure its weight: the very finiteness the grid measures is what the Fundamental Theorem turns into PAC learnability, since $\mathcal{H}$ is PAC learnable if and only if $\VC(\mathcal{H}) < \infty$ (\cref{thm:fundamental}, \leanname{vc_characterization}, \leanname{fundamental_theorem}); the halfplane class, with its grid stalling at $d=4$, is therefore learnable, and this if-and-only-if is the one solidly verified edge attached to the picture.

\item \textbf{Shrink to intervals on the line.} Replace halfplanes in $\R^2$ by halflines, or by intervals, on $\R$. Show this is the instance one dimension down: a single threshold shatters one point but not two of the right pattern, so the largest fillable grid is the $2^1$ grid and $\VC = 1$ for halflines, while intervals fill the $2^2$ grid and $\VC = 2$; both values instantiate the same \leanname{VCDim} definition and are the classical geometric values (\cite{BlumerEhrenfeuchtHausslerWarmuth1989}), so the picture's $d$ is a free parameter, not fixed at three.

\item \textbf{Grow to axis-aligned rectangles.} Replace halfplanes by axis-aligned rectangles in $\R^2$. Show this is the instance that raises the number without changing the grammar: four points placed as the extreme north, south, east, and west witnesses fill the $2^4$ grid, while no five points do, so $\VC = 4$ for rectangles (\cite{BlumerEhrenfeuchtHausslerWarmuth1989}); the figure's $2^d$ grid is unchanged, only the witness shape (a rectangle in each cell instead of a line) and the largest fillable $d$ differ, both still instances of \leanname{VCDim}.

\item \textbf{Climb to halfspaces in $d$ dimensions.} Replace the plane by $\R^d$ and halfplanes by halfspaces. Show this is the instance that states the figure's general value: $d+1$ points in general position fill the $2^{d+1}$ grid and no $d+2$ do, so $\VC = d+1$ for halfspaces (\cite{BlumerEhrenfeuchtHausslerWarmuth1989}), with the $d=2$ picture as the case $\VC = 3$; the value $d+1$ is classical and not in the kernel (there is no halfspace dimension theorem), so it is dashed, and formalizing it together with the Radon upper bound is open.

\item \textbf{The grid never reaches the Littlestone grid.} Ask whether filling the static $2^d$ grid is the same as winning the adaptive online game. Show this is the boundary between two dimensions: shattering a set is non-adaptive (one fixed $S$, all labelings at once), so $\VC(\mathcal{H}) \le \Ldim(\mathcal{H})$ always (\cref{prop:vc-leq-ldim-dim}, \leanname{BranchWiseLittlestoneDim_ge_VCDim}), and the gap is infinite for thresholds, whose static grid fills only to $d=1$ yet whose adaptive mistake tree is unbounded; the figure draws the static witness, and the adaptive refinement is a strictly larger object certified by the same kernel.

\item \textbf{Swap the points and the halfplanes.} Read the figure with roles reversed: let the eight labelings shatter the family rather than the three points shatter the cube. Show this is the frontier of the dual view: the one-inclusion and dual-shattering reading swaps inputs for hypotheses, and the kernel types the dual class with a dual-to-primal shattering transfer (\cref{op:dual-vc-attention}, \leanname{FLT_Proofs.Complexity.DualVC.dual_shatters_imp_original_shatters}), but the self-dual value $\VC = d+1$ under this swap is not drawn and not formalized, so the dual reading of the grid is a dashed, literature-and-open edge.

\item \textbf{Is the grid the rank of a functor?} Read ``the largest finite set on which $\mathcal{H}$ realizes every labeling'' as a universal property. Show this is the figure's deepest frontier: $\VC$ would be the rank of a representable functor on finite sets, the largest free object (full cube) the class can hit, and the chapter's other dimensions (Littlestone, Natarajan \leanname{NatarajanDim}, fat-shattering) would be the same functor specialized by a data-model and a scale; whether a single such dimension functor exists, recovering all of them, or whether the dimensions are provably non-unifiable, is open (\cref{op:dimension-functor}), so this reading of the grid is drawn dashed.
\end{enumerate}

\subsection*{Counterfactual exercises: generalizing Figure~\ref{fig:growth-function}}

Each exercise perturbs the growth-function picture.  Argue, in the figure's grammar (a single curve $\growth_\mathcal{H}(n)$ that rides the $2^n$ reference ray exactly while the class still shatters, $n \le d = \VC(\mathcal{H})$, then peels permanently into the degree-$d$ polynomial $\sum_{i=0}^{d}\binom{n}{i}$ once $n$ passes $d$, drawn solid where the kernel certifies the object the curve traces (the definition \leanname{GrowthFunction} and the bound \leanname{sauer_shelah_quantitative}) and dashed where the asymptotic form $(en/d)^d$ and the immediate-and-permanent reading are classical, cited not proved), that the perturbed picture subsumes the concrete $d=3$ one as an \emph{instance} (a class whose peeling curve bends at its own $\VC$), a \emph{boundary} (one regime, point, or label toggled so the transition moves or vanishes), or a \emph{frontier} (a dashed, open or literature-only edge).

\begin{enumerate}[leftmargin=*,itemsep=3pt]
\item \textbf{The peeling curve is the growth function.} Read the solid curve not as a bound drawn around the class but as the quantity itself. Show this is the figure's defining instance of the verified object: the curve plots $\growth_\mathcal{H}(n) = \max_{|S|=n}|\mathcal{H}|_S|$ for a maximum class of VC dimension three (\cref{def:growth-function}, \leanname{GrowthFunction}), so its coincidence with $2^n$ up to $n=d$ and its descent afterward are the behavior of a single defined function, and the figure certifies \leanname{GrowthFunction} value by value rather than asserting two disconnected ceilings.

\item \textbf{The bend is the bound, and the bound is verified.} Read the moment the curve leaves the $2^n$ ray. Show this is the instance that names the second object: for $n > d$ the curve follows the Sauer--Shelah bound $\sum_{i=0}^{d}\binom{n}{i}$ (\cref{lem:sauer-shelah}, \leanname{sauer_shelah_quantitative}), a genuine kernel theorem and not the legacy existential bound, so this peeling edge of the figure is solidly verified, while the closed form $(en/d)^d$ named beside it is the classical simplification (\cite{Sauer1972}, \cite{Shelah1972}), drawn as annotation because no kernel declaration proves that form.

\item \textbf{Thresholds bend at $d=1$: a line.} Replace the $d=3$ class by halflines on $\R$. Show this is the instance one degree down: a single threshold shatters one point but not two, so $\VC = 1$, the curve rides $2^n$ only at $n \le 1$ and then follows the degree-one polynomial $\sum_{i=0}^{1}\binom{n}{i} = 1+n$, a straight line, peeling from $2^n$ at $n=2$; the same \leanname{GrowthFunction} and \leanname{sauer_shelah_quantitative} apply with $d=1$, so the figure's bend-degree is a free parameter, not fixed at three (\cite{BlumerEhrenfeuchtHausslerWarmuth1989}).

\item \textbf{Rectangles bend at $d=4$.} Replace the class by axis-aligned rectangles in $\R^2$. Show this is the instance that raises the degree without changing the grammar: four extreme witnesses fill the $2^4$ grid and no five do, so $\VC = 4$, the curve coincides with $2^n$ through $n=4$ and peels into a quartic $\sum_{i=0}^{4}\binom{n}{i}$ for $n>4$; only the transition point ($n=d+1=5$) and the terminal slope (degree four) move, both still instances of \leanname{sauer_shelah_quantitative} (\cite{BlumerEhrenfeuchtHausslerWarmuth1989}).

\item \textbf{A finite class flattens to a constant.} Replace the class by a finite $\mathcal{H}$ with $|\mathcal{H}| = N$. Show this is the degenerate instance: $\growth_\mathcal{H}(n) \le \min(2^n, N)$, so the curve rides $2^n$ until $2^n$ exceeds $N$ and is then pinned flat at $N$, the polynomial of degree $\VC \le \log_2 N$ collapsing to a constant ceiling; the same \leanname{GrowthFunction} bounds it, and the bend is here a hard horizontal cap rather than a polynomial arc, the smallest nontrivial case of the dichotomy.

\item \textbf{Infinite VC: the ray that never bends.} Take a class that shatters arbitrarily large sets ($\VC = \infty$), for instance all finite and cofinite subsets of $\N$. Show this is the principal boundary: $\growth_\mathcal{H}(n) = 2^n$ for every $n$, so the curve is the straight $2^n$ ray with no peel, the Sauer--Shelah hypothesis $\VC = d$ never holds for any finite $d$, and \cref{lem:sauer-shelah} simply does not apply; this is the case the Fundamental Theorem excludes, since $\mathcal{H}$ is not PAC learnable when $\VC = \infty$ (\cref{thm:fundamental}, \leanname{vc_characterization}).

\item \textbf{The transition sits exactly at $n=d+1$.} Ask where the bound first constrains the curve below $2^n$. Show this is the boundary that locates the phase change: at $n=d$ a set can still be shattered so $\growth$ may equal $2^d$, but at $n=d+1$ the side condition $d \le m$ of \leanname{sauer_shelah_quantitative} fires and $\sum_{i=0}^{d}\binom{n}{i} < 2^n$ for the first time, forcing the peel; the drop is immediate, with no $n$ of intermediate growth, and permanent, which is the classical content (\cite{Sauer1972}, \cite{Shelah1972}) the marker makes visible.

\item \textbf{$\growth = 2^n$ is the definition of shattering, read on the curve.} Ask why the curve equals $2^n$ exactly on $n \le d$ and not by approximation. Show this is the definitional edge: $\growth_\mathcal{H}(n) = 2^n$ holds iff $\mathcal{H}$ shatters some size-$n$ set (\cref{def:shattering}, \leanname{Shatters}), so the on-the-ray portion of the curve is not a bound being nearly tight but the literal statement that $\mathcal{H}|_S = \{0,1\}^n$ for some $S$; the coincidence with $2^n$ is a definition unfolding, not a theorem, and the figure shows shattering and its failure as one continuous picture.

\item \textbf{Is the bound tight from below?} Ask whether every class of VC dimension three grows as fast as the solid curve, or can stay strictly below it. Show this is the figure's tightness frontier: the kernel certifies only the upper bound (\leanname{sauer_shelah_quantitative}) and the solid curve plots a maximum class that attains it, but no declaration proves a matching lower bound, so whether a given class's growth reaches $\Theta(n^d)$ or drops further is classical for the extremal witness classes (\cite{Sauer1972}) and unformalized in general; this tightness is a dashed, open edge, and the degree it pins is what an optimal compression scheme of size $O(d)$ would also measure (\cref{op:od-compression-dim}).

\item \textbf{Is the bend-degree the rank of a graded object?} Read ``the polynomial that bounds $\growth$ has degree $\VC$'' as a statement about a graded structure. Show this is the figure's deepest frontier: the growth function would be the Hilbert or Poincar\'e series of a graded object attached to $\mathcal{H}$ whose dimension is $\VC$, so ``polynomial of degree $d$'' becomes ``the representable functor has rank $d$'', recovering the chapter's other dimensions (Littlestone, Natarajan \leanname{NatarajanDim}, fat-shattering) as the same series specialized by a data-model and a scale; whether one such dimension functor exists, with the bend-slope as its rank, or the dimensions are provably non-unifiable, is open (\cref{op:dimension-functor}), so this reading of the curve is drawn dashed.
\end{enumerate}

\subsection*{Counterfactual exercises: generalizing Figure~\ref{fig:dimension-hierarchy}}

Each exercise perturbs the dimension ladder.  Argue, in the figure's grammar (a magnitude ladder of combinatorial dimensions $\mathrm{StrongShattering} \le \Ndim \le \DSdim$, with the VC dimension as the binary specialization of $\Ndim$, drawn solid where the kernel certifies the object (the definitions \leanname{VCDim}, \leanname{NatarajanDim}, \leanname{StrongShatteringDim}, the verified rung \leanname{strongShatteringDim_le_NatarajanDim}, and the binary characterization \leanname{vc_characterization}) and dashed where the genuine DS dimension, the upper rung $\Ndim \le \DSdim$, the multiclass characterization, and the Natarajan failure are literature, cited not proved), that the perturbed picture subsumes the concrete ladder as an \emph{instance} (a label-set regime or class on which a rung collapses or is attained), a \emph{boundary} (one label-set size or one inequality toggled so a rung reverses or a characterization breaks), or a \emph{frontier} (a dashed, open or literature-only edge).

\begin{enumerate}[leftmargin=*,itemsep=3pt]
\item \textbf{Two labels collapse the ladder to the Fundamental Theorem.} Set $|Y| = 2$ and read the whole ladder. Show this is the figure's defining instance of the binary regime: when the label set has two elements the Natarajan dimension equals the VC dimension exactly (\cref{def:natarajan}, \leanname{NatarajanDim}), the strong-shattering dimension equals both, and the three rungs degenerate to the single binary characterization $\VC(\mathcal{H}) < \infty \Leftrightarrow \mathcal{H}$ is PAC learnable (\cref{thm:fundamental}, \leanname{vc_characterization}, \leanname{fundamental_theorem}); the ladder is the multiclass refinement of one verified equivalence, and the $|Y|=2$ slice is the only place the entire figure is solid.

\item \textbf{The lower rung is the figure's only solid multiclass edge.} Read the bottom inequality alone. Show this is the instance that names the verified handle: from any set on which every labeling is realizable one builds a two-colouring Natarajan witness, so $\mathrm{StrongShatteringDim}(\mathcal{H}) \le \Ndim(\mathcal{H})$ for $|Y| \ge 2$ (\leanname{strongShatteringDim_le_NatarajanDim}), a genuine kernel theorem proved by exhibiting the witness, not asserted; this is the single machine-checked rung of the multiclass ladder, and the figure draws it solid while everything above it is cited.

\item \textbf{The Natarajan dimension is the obvious lift, and it equals VC at the base.} Read the $\Ndim$ node as the natural generalization. Show this is the instance that states the figure's central object: the Natarajan dimension replaces ``all $2^d$ binary labelings'' by ``all $2^d$ labelings drawn from two opposing functions'' (\cref{def:natarajan}, \leanname{NatarajanDim}), which is exactly the VC dimension when $|Y|=2$ and a strict generalization otherwise; the figure draws it solid as a verified definition, leaving open only whether finiteness of it suffices for learning, which the next exercise denies.

\item \textbf{Finite Natarajan dimension does not imply learnability.} Toggle the question from ``is $\Ndim$ defined'' to ``does $\Ndim < \infty$ characterize learning''. Show this is the figure's principal boundary, the red dashed failure: there is a class with $\Ndim(\mathcal{H}) = 1$ that is not PAC learnable (\cref{thm:natarajan-failure}, \cite{DanielyShalevShwartz2014}), so the obvious dimension can be as small as possible and still miss unlearnability; the implication $\Ndim < \infty \Rightarrow$ learnable is drawn dashed red because it is false, the one edge the figure refutes rather than asserts.

\item \textbf{An infinite label set is where the obvious dimension fails.} Ask where the witness of the previous exercise lives. Show this is the boundary in the label-set sweep: the $\Ndim = 1$ unlearnable class requires $|Y| = \infty$ (\cref{thm:natarajan-failure}), because the Natarajan dimension reads only pairs of opposing functions and cannot see the combinatorial explosion of an unbounded label set; the gap is not informational but label-set-driven, and quantifying how large the DS dimension grows as a function of $|Y|$ when $\Ndim$ stays finite is open (\cref{op:natarajan-ds-gap}).

\item \textbf{The DS dimension is the correct characterization, and it is literature.} Replace the failed rung by the true one. Show this is the instance that names the third object: a class is PAC learnable if and only if its DS dimension is finite (\cref{thm:ds-char}, \cite{BrukhimCarmonDinurMoranYehudayoff2022}), the multiclass analogue of the Fundamental Theorem, with the DS dimension defined via pseudo-cubes in the one-inclusion hypergraph (\cref{def:ds-dim}); both the dimension and its equivalence are drawn dashed because the genuine pseudo-cube notion is not machine-checked, so the figure cites the characterization rather than verifying it.

\item \textbf{The kernel's handle is the reverse inequality, not the DS dimension.} Read the bottom rung against the top one. Show this is the boundary that separates the verified object from the literature object: the kernel proves $\mathrm{StrongShatteringDim} \le \Ndim$ (\leanname{strongShatteringDim_le_NatarajanDim}), whereas the genuine DS relationship is $\Ndim \le \DSdim$, the \emph{opposite} direction, and the DS dimension can be infinite while the Natarajan dimension is finite so no $\DSdim \le \Ndim$ bound can hold; the all-labelings strong-shattering dimension is therefore \emph{not} the DS dimension, and the figure must show the kernel's actual handle rather than silently substitute it for $\DSdim$.

\item \textbf{The full ladder has a verified floor and a literature ceiling.} Read all three rungs as one chain. Show this is the boundary between evidence grades: the complete ladder is $\mathrm{StrongShattering} \le \Ndim \le \DSdim$, of which only the left inequality is kernel-verified (\leanname{strongShatteringDim_le_NatarajanDim}) while the right inequality $\Ndim \le \DSdim$ is literature (\cite{BrukhimCarmonDinurMoranYehudayoff2022}); the line-style is the statement that multiclass learnability is currently half machine-checked, the binary characterization and the lower rung solid, the multiclass characterization and the upper rung dashed.

\item \textbf{Formalize the genuine DS dimension and its characterization.} Ask what it would take to make the dashed ceiling solid. Show this is the figure's headline frontier: formalizing the pseudo-cube DS dimension and the Brukhim et al.\ proof that $\DSdim(\mathcal{H}) < \infty$ iff $\mathcal{H}$ is multiclass PAC learnable would turn the dashed top rung and the dashed characterization into verified edges (\cref{op:ds-formal}), supplying the multiclass analogue of \leanname{vc_characterization}; until then the one-inclusion hypergraph and its pseudo-cubes (\cref{def:ds-dim}) are a typed-but-unproved target, and this is the chapter's principal open formalization.

\item \textbf{Is there a single dimension functor under the whole ladder?} Read the three dimensions as one object specialized three ways. Show this is the figure's deepest frontier: the VC, Natarajan, and DS dimensions would be values of a single functor from concept classes (with label-preserving maps) to ordinals, recovered by specializing a data-model and a label-set, with the one-inclusion hypergraph as the representable structure whose pseudo-cube rank is the learnability dimension; whether one such dimension functor exists, recovering all rungs of the ladder, or the dimensions are provably non-unifiable, is open (\cref{op:dimension-functor}), so this reading of the ladder is drawn dashed.
\end{enumerate}

% Chapter 11: Sample Complexity, Compression, and Occam's Razor
% sample_complexity = Profile B-light, unlabeled_compression = Profile D (Negative),
% occam_algorithm = Profile B-light, description_length/kolmogorov/kl = Profile A,
% covering_number = Profile A-plus.
%
% Texture: question-driven. The open conjecture (compression ≈ VC dimension)
% is the gravitational center. Everything else orbits it: sample complexity
% provides the motivation, Occam's razor provides the algorithmic bridge,
% information-theoretic notions provide the vocabulary. The chapter should
% feel like a detective story where the case remains unsolved.

\chapter{Sample Complexity, Compression, and Occam's Razor}
\label{ch:compression}

A learner that processes a thousand training examples may ultimately lean on
three.  The Fundamental Theorem of \Cref{ch:pac} tells us \emph{which} classes
are PAC learnable, but it says nothing about what the learner actually needs to
retain.  The sharper question is: \emph{how little of the data must survive
for learning to succeed?}

The tight sample complexity bounds of \Cref{sec:sample-bounds} establish that
the VC dimension $d$ determines how many examples are needed, up to constant
factors.  Compression schemes (\Cref{sec:compression}) then reframe the same
question from the learner's side: a class is learnable if and only if any
training sample can be reduced to a bounded subsample that fully determines the
hypothesis.  The size of that subsample is what the chapter is really about.
The natural conjecture, that $O(d)$ points always suffice, was stated by
Littlestone and Warmuth in 1986.  It has not been resolved in the forty years
since.  Every known learnable class compresses to $O(d)$ points; no one can
prove that all of them do.

Two landmarks bracket the frontier.  Moran and Yehudayoff proved in 2016 that
compression to $2^{O(d)}$ points always exists: a fact unknown before their
work.  For $d = 20$, that bound is roughly $10^6$ points; the conjecture asks
for $20$.  That distance is the gap the conjecture has not closed in forty years.

Occam's razor (\Cref{sec:occam}) provides the bridge between compression and
generalization: any algorithm that finds a description of the data short enough
relative to the sample size generalizes.  The information-theoretic foundations
of \Cref{sec:info-foundations} (description length, Kolmogorov complexity,
covering numbers) give precise meaning to ``short enough.''

\medskip

The chapter is organized around a single question:

\begin{center}
\itshape
Does every concept class of VC dimension $d$ admit a compression scheme of
size $O(d)$?
\end{center}

Everything we prove either approaches this question or illuminates why it
resists resolution.


% ================================================================
% SECTION 1: TIGHT SAMPLE COMPLEXITY BOUNDS
% Profile B-light. State the bounds, give proof sketch for upper bound.
% ================================================================

\section{Tight Sample Complexity Bounds}
\label{sec:sample-bounds}

The sample complexity of PAC learning was bounded in \Cref{ch:pac} through the
VC characterization.  Here we state the tight bounds precisely.

\begin{thm}[Realizable PAC sample complexity {\cite{Hanneke2016,SimonZilles2015}}]
\label{thm:sample-realizable}
Let $\mathcal{H}$ be a hypothesis class with $\VC(\mathcal{H}) = d < \infty$.
In the realizable setting, the optimal sample complexity of PAC learning
$\mathcal{H}$ satisfies
\[
  m(\varepsilon, \delta) = \Theta\!\left(\frac{d}{\varepsilon}
    + \frac{\log(1/\delta)}{\varepsilon}\right).
\]
The upper bound is achieved by any consistent learner (including $\ERM$); the
lower bound holds for every learner.
\end{thm}

\begin{remark}[History of the tight bound]
\label{rmk:sample-history}
The classical upper bound of Blumer et al.~\cite{BlumerEhrenfeuchtHausslerWarmuth1989}
gave $O\!\bigl(\frac{d}{\varepsilon}\log\frac{1}{\varepsilon} +
\frac{1}{\varepsilon}\log\frac{1}{\delta}\bigr)$, with a logarithmic overhead
in the $d/\varepsilon$ term.  Removing this overhead (proving that the correct
dependence on $d$ is linear, not $d \log(1/\varepsilon)$) required substantial
effort.  The matching lower bound was established by Ehrenfeucht et
al.~\cite{EhrenfeuchtHausslerKearnsValiant1989}.  The tight upper bound without
the logarithmic factor was completed by Simon and Zilles~\cite{SimonZilles2015}
and, in a sharper form addressing the constant, by
Hanneke~\cite{Hanneke2016}.  Hanneke's result shows that \emph{any} consistent
learner achieves the optimal bound: no algorithmic cleverness beyond
consistency is needed.
\end{remark}

\begin{proof}[Proof sketch of the upper bound]
The argument proceeds in two stages.  First, one establishes the \emph{double
sampling} bound: the probability that $\ERM$ returns a hypothesis with error
$> \varepsilon$ is at most
\[
  \Pr[\exists h \in \mathcal{H} : \emprisk_S(h) = 0 \text{ and }
    \risk_D(h) > \varepsilon]
  \;\leq\; 2\,\growth_\mathcal{H}(2m) \cdot 2^{-\varepsilon m / 2},
\]
where $\growth_\mathcal{H}(n) \leq (en/d)^d$ by the Sauer--Shelah lemma
(\Cref{lem:sauer-shelah}).  Setting the right-hand side $\leq \delta$ and
solving for $m$ gives $m = O\!\bigl(\frac{d}{\varepsilon}
\log\frac{1}{\varepsilon} + \frac{1}{\varepsilon}\log\frac{1}{\delta}\bigr)$.

The removal of the $\log(1/\varepsilon)$ factor uses a more delicate argument.
One decomposes the hypothesis class into subsets of bounded ``disagreement mass''
and applies uniform convergence separately to each piece.  The key insight
(Hanneke~\cite{Hanneke2016}) is that the number of effectively distinct
hypotheses at accuracy scale $\varepsilon$ is controlled by the \emph{star
number} of $\mathcal{H}$, which is $O(d)$ rather than $O(d \log(1/\varepsilon))$.
\end{proof}

In the agnostic setting, the sample complexity worsens:

\begin{thm}[Agnostic PAC sample complexity]
\label{thm:sample-agnostic}
In the agnostic setting, the optimal sample complexity satisfies
\[
  m(\varepsilon, \delta) = \Theta\!\left(\frac{d}{\varepsilon^2}
    + \frac{\log(1/\delta)}{\varepsilon^2}\right).
\]
\end{thm}

The gap between $d/\varepsilon$ (realizable) and $d/\varepsilon^2$ (agnostic)
is the ``$\varepsilon^2$ price'' discussed in \Cref{sec:agnostic} of
\Cref{ch:pac}.  It is not an artifact of loose analysis: the lower bound
matches.


% ================================================================
% SECTION 2: COMPRESSION SCHEMES
% Profile F (Open Frontier). CENTERPIECE of the chapter.
% ================================================================

\section{Compression Schemes}
\label{sec:compression}

Sample complexity tells us how much data learning requires.  Compression asks
what happens once the data has arrived: \emph{how little of it does the learner
actually need to keep?}

\begin{defn}[Sample Compression Scheme {\cite{LittlestoneWarmuth1986,FloydWarmuth1995}}]
\label{def:compression-scheme}
A \emph{sample compression scheme} of size $k$ for a hypothesis class
$\mathcal{H}$ over domain $X$ consists of two maps:
\begin{enumerate}[leftmargin=*, itemsep=2pt]
\item A \emph{compression function} $\kappa$ that, given any sample $S$
  consistent with some $h \in \mathcal{H}$, selects a subsample
  $\kappa(S) \subseteq S$ with $|\kappa(S)| \leq k$.
\item A \emph{reconstruction function} $\rho$ that, given $\kappa(S)$, produces
  a hypothesis $\rho(\kappa(S))$ consistent with $S$.
\end{enumerate}
The reconstruction function $\rho$ sees \emph{only} the compressed subsample
$\kappa(S)$; it has no access to the remaining points.\formal{FLT_Proofs.Complexity.Structures.CompressionScheme}
\end{defn}

\begin{example}[Support vectors as compression]
\label{ex:svm-compression}
A support vector machine in $\R^d$ with a hard margin computes a
hyperplane determined by at most $d+1$ support vectors.  These support vectors
form a compression scheme of size $d+1$: the subsample $\kappa(S)$ consists of
the support vectors, and the reconstruction function $\rho$ returns the
maximum-margin hyperplane through them.
\end{example}

The fundamental connection between compression and learnability was established
by Littlestone and Warmuth:

\begin{thm}[Compression implies learnability {\cite{LittlestoneWarmuth1986,FloydWarmuth1995}}]
\label{thm:compression-pac}
If $\mathcal{H}$ admits a sample compression scheme of size $k$, then
$\mathcal{H}$ is PAC learnable with sample complexity
\[
  m(\varepsilon, \delta) = O\!\left(\frac{k}{\varepsilon}
    \log\frac{k}{\varepsilon} + \frac{1}{\varepsilon}\log\frac{1}{\delta}\right).
\]
In particular a finite compression scheme forces finite VC dimension, $\VC(\mathcal{H}) \leq 2(k+1)^2 - 1$, so compressibility implies learnability.\formal{FLT_Proofs.Complexity.Generalization.compression_imp_vcdim_finite}
\end{thm}

\begin{proof}[Proof sketch]
Fix a sample $S$ of size $m$ drawn i.i.d.\ from $D$.  The compression function
selects a subsample $\kappa(S)$ of size $\leq k$.  The number of possible
subsamples is $\binom{m}{k} \leq m^k$.  For each fixed subsample, the
reconstruction $\rho(\kappa(S))$ is a fixed hypothesis, and by a Hoeffding
bound, a single fixed hypothesis with empirical risk zero on the remaining
$m - k$ points has true risk $> \varepsilon$ with probability at most
$e^{-\varepsilon(m-k)}$.  Taking a union bound over all $m^k$ possible
subsamples:
\[
  \Pr[\risk_D(\rho(\kappa(S))) > \varepsilon] \leq m^k \cdot e^{-\varepsilon(m-k)}.
\]
Setting this $\leq \delta$ and solving for $m$ gives the stated bound.
\end{proof}

The converse direction (does learnability imply compression?) is where the
story becomes a forty-year open problem.

\subsection{The Compression Conjecture}

\begin{openprob}[Sample Compression Conjecture {\cite{LittlestoneWarmuth1986,FloydWarmuth1995}}]
\label{op:compression}
Does every concept class $\mathcal{H}$ with $\VC(\mathcal{H}) = d$ admit a
sample compression scheme of size $O(d)$?\formal{FLT_Proofs.Complexity.Structures.Open_NoInfoCompressionStrengthening}
\end{openprob}

This conjecture, stated by Littlestone and Warmuth in 1986, remains open.  It
asks for the converse of \Cref{thm:compression-pac}: if a class is learnable
(which, by the Fundamental Theorem, is equivalent to finite VC dimension $d$),
can the learner always compress its evidence to $O(d)$ points?

Four decades on, the statement is clean, the motivation immediate, and every
natural proof strategy fails.  Topological arguments reach $2^{O(d)}$ but not
$O(d)$.  Algebraic approaches founder on classes with no linear structure.  The
gap between what compression \emph{exists} and what the conjecture \emph{demands}
has not narrowed since the question was posed.

\begin{thm}[Exponential compression exists {\cite{MoranYehudayoff2016}}]
\label{thm:moran-yehudayoff}
Every concept class $\mathcal{H}$ with $\VC(\mathcal{H}) = d$ admits a sample
compression scheme of size $2^{O(d)}$.\formal{FLT_Proofs.Theorem.PAC.fundamental_vc_compression}
\end{thm}

This result of Moran and Yehudayoff~\cite{MoranYehudayoff2016}
established that compression schemes \emph{exist} for every learnable class, a
fact unknown before 2016.  The exponential bound $2^{O(d)}$ is vastly larger
than the conjectured $O(d)$.  The scheme it produces carries side information;
whether the side bits can be removed, and whether the size can be brought to
$O(d)$, are the open strengthenings (\cref{op:compression}).

\begin{proof}[Proof sketch]
The proof uses a topological argument.  By a result of
Radon, any set of $d+1$ points in a $d$-dimensional space admits a Radon
partition.  Moran and Yehudayoff generalize this to VC classes via a
``staircase'' construction: they build a sequence of hypothesis classes of
increasing VC dimension and show that each admits a compression scheme, using
the Sauer--Shelah lemma (\Cref{lem:sauer-shelah}) to control the growth of
the compression size.  The exponential blowup arises from iterating this
construction $d$ times.
\end{proof}

\begin{remark}[The gap]
\label{rmk:compression-gap}
The gap between the conjecture ($O(d)$) and the best known upper bound
($2^{O(d)}$) is enormous.  For $d = 20$, the conjecture predicts compression to
$\sim 20$ points; the Moran--Yehudayoff bound gives $\sim 10^6$.  No
polynomial-in-$d$ bound is known for general classes.
\end{remark}

\subsection{The One-Inclusion Graph}
\label{sec:one-inclusion}

The Moran--Yehudayoff proof rests on a combinatorial object that deserves
explicit treatment, both because it illuminates why maximum classes are
special and because it localizes the difficulty of the conjecture at a
precise structural point.

\begin{defn}[One-Inclusion Graph {\cite{HausslerLittlestonWarmuth1994}}]
\label{def:one-inclusion-graph}
Let $\mathcal{H}$ be a concept class over domain $X$, and let
$S = (x_1, \ldots, x_m)$ be a sequence of points from $X$.  The
\emph{one-inclusion graph} $G(\mathcal{H}, S)$ is the graph whose
vertex set is
$\{(h(x_1), \ldots, h(x_m)) : h \in \mathcal{H}\}$
(the distinct restrictions of $\mathcal{H}$ to $S$)
and whose edge set connects pairs of restrictions that differ on
exactly one coordinate.
\end{defn}

The one-inclusion graph mediates between the VC dimension and
compression.  The key property is that its density (specifically, the
minimum over all orientations of the maximum in-degree) controls
whether a compression scheme of size $d$ exists.

\begin{thm}[One-Inclusion Graph Orientation {\cite{RubinsteinKroese2007,KuzminWarmuth2007}}]
\label{thm:one-inclusion-orientation}
Let $\mathcal{H}$ be a \emph{maximum class} of VC dimension $d$
(i.e., $\growth_\mathcal{H}(m) = \sum_{i=0}^d \binom{m}{i}$ for all
$m$).  Then for any sample $S$ of size $m \geq d$, the one-inclusion
graph $G(\mathcal{H}, S)$ can be oriented so that every vertex has
in-degree at most $d$.
\end{thm}

This orientation directly yields a compression scheme of size $d$.
Given a sample $S$ consistent with concept $h^*$, the vertex
$v = h^*|_S$ in the oriented one-inclusion graph has at most $d$
incoming edges.  Each incoming edge corresponds to a sample point
$x_i$ where some neighboring concept $h'$ disagrees with $h^*$.
The compressed subsample $\kappa(S)$ consists of these $\leq d$
``witness points.''  The reconstruction function $\rho$ recovers
$h^*$ from the witness points using the acyclicity of the oriented
graph: the witness points uniquely determine $h^*$ among all
concepts in $\mathcal{H}$ consistent with $S$.

\begin{example}[One-inclusion graph for intervals]
\label{ex:oig-intervals}
Let $\mathcal{H}$ be the class of intervals $[a, b] \subseteq [0, 1]$
(labeling $x$ as $1$ iff $x \in [a,b]$), with $\VC(\mathcal{H}) = 2$.
This is a maximum class.  Given a sorted sample
$S = (x_1 < x_2 < \cdots < x_m)$, each concept restricts to a binary
string of the form $0^i 1^j 0^k$ (a contiguous block of 1s flanked by
0s).  The one-inclusion graph connects patterns that differ in a single
bit, which can only happen at the boundaries of the block: either the
leftmost $1$ becomes $0$ or the rightmost $1$ becomes $0$ (or the
reverse).  The graph can be oriented so that each vertex points toward
its two boundary points.  The compressed subsample consists of the two
boundary points of the interval: exactly the information needed to
reconstruct $[a, b]$ on the sample.

The mechanism is subtle: what matters is the two \emph{sample points}
at the boundary of the positive region.  The \emph{endpoints} of the
interval in $[0,1]$ play no role (and recovering them would require
real-valued parameters, which the scheme never stores).  The
one-inclusion graph identifies the two boundary sample points
combinatorially, without reference to the underlying geometry.
\end{example}

For general (non-maximum) classes, the one-inclusion graph does
\emph{not} admit an orientation with in-degree $\leq d$.  This is
the precise technical point where the conjecture resists.

\begin{remark}[The maximum-to-general gap]
\label{rmk:max-general-gap}
Moran and Yehudayoff's strategy is to reduce the general case to the
maximum case.  Any class $\mathcal{H}$ of VC dimension $d$ can be
``fattened'' to a maximum class $\mathcal{M}$ with
$\VC(\mathcal{M}) = d$ by adding concepts until the Sauer--Shelah
bound is achieved with equality.  The maximum class has compression of
size $d$, and the compression scheme can in principle be pulled back to
$\mathcal{H}$.  However, the fattening introduces auxiliary
concepts; the reconstruction function, given a compressed subsample,
must decide which of many $\mathcal{M}$-consistent concepts actually
belongs to $\mathcal{H}$.  The number of such choices is exponential
in $d$, and encoding this choice inflates the compression size to
$2^{O(d)}$.

Closing the gap to $O(d)$ requires either working directly with
non-maximum classes (for which no orientation theorem is known) or
proving that the reconstruction ambiguity can be resolved with
$O(d)$ additional bits rather than $2^{O(d)}$.
\end{remark}


The conjecture has been verified for specific class families:

\begin{enumerate}[leftmargin=*, itemsep=3pt]
\item \textbf{Maximum classes.}  A class $\mathcal{H}$ with
  $\growth_\mathcal{H}(n) = \sum_{i=0}^d \binom{n}{i}$ for all $n$ (achieving
  the Sauer--Shelah bound with equality) admits a compression scheme of size
  $d$~\cite{FloydWarmuth1995,RubinsteinKroese2007}.

\item \textbf{Intersection-closed classes.}  If $\mathcal{H}$ is closed under
  intersections, compression of size $d$ exists.

\item \textbf{Classes of VC dimension 1.}  Trivially: a single point suffices.

\item \textbf{Halfspaces in $\R^d$.}  The support vector compression of
  \Cref{ex:svm-compression} gives size $d + 1$.

\item \textbf{Balls, rectangles, and other geometric classes.}  Known
  case-by-case.
\end{enumerate}

The intersection-closed case is worth examining in detail, because its
compression mechanism is fundamentally different from the SVM case and
reveals what structure makes compression ``easy.''

\begin{example}[Compression for intersection-closed classes]
\label{ex:intersection-closed-compression}
Let $\mathcal{H}$ be a concept class closed under intersection: if
$h_1, h_2 \in \mathcal{H}$, then $h_1 \cap h_2 \in \mathcal{H}$
(where $(h_1 \cap h_2)(x) = h_1(x) \wedge h_2(x)$).

Given a sample $S$ consistent with target $h^* \in \mathcal{H}$,
define
$\hat{h} = \bigcap\{h \in \mathcal{H} : h \text{ consistent with } S\}$.
By intersection-closure, $\hat{h} \in \mathcal{H}$, and $\hat{h}$ is
the \emph{unique minimal} concept consistent with~$S$.  The
compression scheme exploits this uniqueness.

Let $P^+ = \{x_i \in S : h^*(x_i) = 1\}$ be the positive examples.
The minimal concept $\hat{h}$ is determined by the constraint that it
must label all points in $P^+$ as positive while being minimal in
$\mathcal{H}$.  A \emph{spanning set} for $P^+$ is a subset
$T \subseteq P^+$ such that
$\bigcap\{h \in \mathcal{H} : h(x) = 1\;\forall x \in T\}
 = \bigcap\{h \in \mathcal{H} : h(x) = 1\;\forall x \in P^+\}$%
that is, $T$ imposes the same constraints as the full positive set.
A minimal spanning set has size $\leq d = \VC(\mathcal{H})$, because
the VC dimension bounds the number of independent constraints.  The
compression function selects a minimal spanning set; the
reconstruction function computes
$\hat{h} = \bigcap\{h \in \mathcal{H} : h(x) = 1\;\forall (x, 1) \in
\kappa(S)\}$.

The mechanism is structurally different from SVMs.  Support vector
machines use \emph{geometric} structure (margin maximization) to
select anchor points.  Intersection-closed classes use \emph{lattice}
structure (existence of meets) to guarantee a unique minimum.  Both
achieve compression of size $d$, but via incomparable mathematical
facts.  This incomparability is a symptom of the conjecture's
difficulty: there is no single mechanism that explains compression
across all class families.
\end{example}


\subsection{Compression for Attention Classes}
\label{sec:attention-compression}

The compression question is not a relic of classical concept classes; it
reappears, still unsolved, for the hypothesis classes computed by modern
attention layers.  Fix an embedding dimension $d$ and a context of $n_K$
key--value pairs.  A single softmax-attention head maps a query $x \in \R^d$ to
the weighted readout
$\sum_{i=1}^{n_K} \mathrm{softmax}(\langle x, k_i\rangle)_i\, v_i$, and
thresholding this readout at zero yields a concept
$c_{K,V}\colon \R^d \to \{0,1\}$.  Letting the keys $K = (k_1, \ldots, k_{n_K})$
and values $V = (v_1, \ldots, v_{n_K})$ range over all configurations gives the
\emph{attention concept class} $\mathcal{A}_{d,n_K} = \{c_{K,V}\}$.

\begin{prop}[Attention classes are compressible]
\label{prop:attention-compression}
For every $d$ and $n_K$, the attention concept class $\mathcal{A}_{d,n_K}$ has
finite VC dimension and therefore admits a sample compression scheme of size
$2^{O(\VC(\mathcal{A}_{d,n_K}))}$.
\end{prop}

\begin{proof}
The readout map
$(K, V, x) \mapsto \sum_i \mathrm{softmax}(\langle x,k_i\rangle)_i\, v_i$ is
jointly continuous, hence Borel measurable, so the bad-event sets the
uniform-convergence argument quantifies over are measurable.  This is the
well-behaved measurable VC target
condition,\formal{TLT.Bridge.softAttention_wellBehaved} so the empirical process
indexed by $\mathcal{A}_{d,n_K}$ converges uniformly and the class has finite VC
dimension.  Finite VC dimension is the hypothesis of the existence
theorem,\formal{FLT_Proofs.Theorem.PAC.fundamental_vc_compression} which supplies
a compression scheme of size exponential in the VC dimension.
\end{proof}

The exponential bound is, as for general classes, far from tight, and the argmax
(hard-attention) router suggests why.  When the softmax is sharpened to a top-$1$
selection, the head's output on a query $x$ is determined by the single key whose
score is largest, and the routing rule that selects this key is again
well-behaved,\formal{TLT_Proofs.Routing.MeasurableScoreRouting.multiHeadArgmax_wellBehaved}
partitioning the query space into finitely many Borel cells.  A query's label
depends only on which cell it lands in, and each cell is pinned down by a bounded
number of support keys (the attention analogue of the support vectors of
\Cref{ex:svm-compression}).  This predicts a compression of size $O(n_K)$, linear
in the context length.  Turning the heuristic into a theorem runs into the
obstacle of \Cref{op:compression}: the support-key set is selected with knowledge
of the whole sample, and pulling it back to a labeled subsample of size $O(\VC)$
is the open problem.

\begin{remark}[The conjecture, modernized]
\label{rmk:attention-compression}
\Cref{prop:attention-compression} places attention classes inside the landscape
of \Cref{fig:compression-landscape} at the compression-size $2^{O(d)}$ node, with
the descent to size $O(d)$ open.  The Littlestone--Warmuth conjecture is thus
load-bearing for contemporary architectures.  A positive resolution would bound
how few training examples an attention head must retain to reproduce its
behavior; a negative one would exhibit a learnable attention class that is
incompressible.
\end{remark}


\subsection{Why the Conjecture Resists Resolution}
\label{sec:why-hard}

The compression conjecture has resisted attack for four decades.  Techniques
are not scarce: the difficulty is structural.  Every known approach either
works for specific, well-structured class families but cannot reach arbitrary
VC classes, or works for all VC classes but pays an exponential price in scheme
size.  The two regimes do not overlap, and no method yet bridges them.

\begin{obstbox}
\textbf{Three approaches and their failure modes.}

\medskip
\emph{Approach 1: Direct construction.}  For each verified class
family (halfspaces, intervals, intersection-closed, maximum classes),
one can construct a compression scheme tailored to the family's
structure.  The constructions exploit distinct properties: geometric
anchoring for SVMs, lattice minimality for intersection-closed classes,
one-inclusion graph orientation for maximum classes.  But these
properties are incomparable: no single structural assumption covers all
families, let alone arbitrary VC classes.  Every direct construction
exploits something that arbitrary classes lack.

\medskip
\emph{Approach 2: One-inclusion graph orientation.}  For maximum
classes, \Cref{thm:one-inclusion-orientation} provides an orientation
with in-degree $d$, directly yielding compression of size $d$.  For
non-maximum classes, no such orientation theorem is known.
Non-maximum one-inclusion graphs can have vertices with degree much
larger than $d$, and no known orientation strategy achieves in-degree
$O(d)$ in general.  The maximum class case is tight, but the
combinatorial extremality that makes it work (achieving the
Sauer--Shelah bound with equality) is exactly what arbitrary classes
lack.

\medskip
\emph{Approach 3: Topological embedding.}  Moran and Yehudayoff's
proof uses a topological argument related to Radon partitions in
convex geometry.  The topological machinery is powerful enough to
handle all VC classes but pays an exponential price: the Radon-type
argument iterates a doubling construction $d$ times, producing
compression of size $2^{O(d)}$.  Improving this requires either a
fundamentally different topological tool or a way to avoid the
iterative blowup; the iterative structure appears intrinsic to
the proof method itself, the doubling having to run $d$ times before a
Radon partition is available, so a tighter analysis of the same
construction would not remove it.

\medskip
The conjecture sits at the intersection of three mathematical
territories: combinatorics (VC dimension, growth functions, maximum
classes), topology (Radon partitions, convex position), and
information theory (description length, reconstruction
complexity); and each territory contributes techniques that
illuminate a piece of the puzzle while leaving the central question
open.  A resolution may require connecting these territories in a way
that current methods do not.
\end{obstbox}


\subsection{Labeled vs.\ Unlabeled Compression}
\label{sec:unlabeled-compression}

In \Cref{def:compression-scheme}, the compressed subsample $\kappa(S)$ retains
both the points and their labels.  Labels carry information, and a natural
question is whether the labels are doing real work: can the reconstruction
function recover from the point positions alone?

\begin{defn}[Unlabeled Compression Scheme]
\label{def:unlabeled-compression}
An \emph{unlabeled compression scheme} of size $k$ for $\mathcal{H}$ is a pair
$(\kappa, \rho)$ where $\kappa$ selects a set of $\leq k$ \emph{points}
(without labels) from the sample, and $\rho$ reconstructs a hypothesis
consistent with $S$ from these points alone.
\end{defn}

The compression conjecture has an unlabeled analogue, in which the stored
subsample carries no labels and the reconstruction recovers them from the point
positions alone.

\begin{thm}[Bounded unlabeled compression {\cite{MoranYehudayoff2016}}]
\label{thm:unlabeled-bounded}
Every concept class $\mathcal{H}$ with $\VC(\mathcal{H}) = d$ admits an
unlabeled compression scheme of size bounded by a function of $d$ alone,
independent of $|\mathcal{H}|$.
\end{thm}

This matches the labeled case qualitatively: labeled compression of size
$2^{O(d)}$ always exists (\Cref{thm:moran-yehudayoff}), and the unlabeled scheme
stays bounded by a function of $d$ as well.  Whether the unlabeled size can be
pushed to $O(d)$, matching the labeled conjecture, is the Kuzmin--Warmuth
unlabeled compression conjecture~\cite{KuzminWarmuth2007}, open in general and
proved for maximum classes.

\begin{remark}
The unlabeled conjecture sharpens the labeled one (\Cref{op:compression}):
recovering labels from the point positions alone is the extra demand, and
whether it costs more than a constant factor in scheme size is unknown.
\end{remark}


% ================================================================
% SECTION 3: OCCAM'S RAZOR
% Profile B-light. Occam → PAC and partial converse.
% ================================================================

\section{Occam's Razor}
\label{sec:occam}

Compression schemes embody one form of Occam's razor: prefer the hypothesis
determined by the fewest data points.  A description-length version of the same
principle, formalized by Blumer, Ehrenfeucht, Haussler, and Warmuth, shifts the
criterion from subsample size to the length of the binary encoding: the output
hypothesis should be shorter than the sample itself.  Both forms connect
simplicity to generalization, and the formal connection runs precisely.

\begin{defn}[Occam Algorithm {\cite{BlumerEhrenfeuchtHausslerWarmuth1987}}]
\label{def:occam}
Let $\mathcal{H}$ be a hypothesis class and let $\mathcal{R}$ be a
\emph{representation class} (a set of hypotheses with associated description
lengths).  An algorithm $A$ is an \emph{Occam algorithm} for $\mathcal{H}$ with
respect to $\mathcal{R}$ if, given a sample $S$ of size $m$ consistent with some
$h \in \mathcal{H}$ of description length $n$, the algorithm outputs a
hypothesis $h' \in \mathcal{R}$ that:
\begin{enumerate}[leftmargin=*, nosep]
\item is consistent with $S$, and
\item has description length $|h'| \leq n^\alpha m^\beta$ for some constants
  $\alpha \geq 1$ and $0 \leq \beta < 1$.
\end{enumerate}
The sub-linear dependence on $m$ (exponent $\beta < 1$) is the key requirement:
the output hypothesis must be ``simpler'' than the data.
\end{defn}

\begin{thm}[Occam's Razor Theorem {\cite{BlumerEhrenfeuchtHausslerWarmuth1987}}]
\label{thm:occam}
If $A$ is an Occam algorithm for $\mathcal{H}$ with parameters $\alpha, \beta$,
then $A$ PAC-learns $\mathcal{H}$ with sample complexity
\[
  m(\varepsilon, \delta) = O\!\left(\frac{1}{\varepsilon}\left(
    n^\alpha \log\frac{1}{\varepsilon}
    + \log\frac{1}{\delta}\right)\right)^{1/(1-\beta)}.
\]
\end{thm}

\begin{proof}[Proof sketch]
The argument is a counting and union bound.  The number of hypotheses in
$\mathcal{R}$ with description length $\leq L$ is at most $2^L$.  An Occam
algorithm outputs a hypothesis of length $\leq n^\alpha m^\beta$.  Fixing a
sample of size $m$, the number of possible outputs is at most
$2^{n^\alpha m^\beta}$.  For each fixed hypothesis with empirical risk zero, the
probability that its true risk exceeds $\varepsilon$ is at most
$(1-\varepsilon)^m \leq e^{-\varepsilon m}$.  A union bound gives
\[
  \Pr[\risk_D(A(S)) > \varepsilon] \leq 2^{n^\alpha m^\beta} \cdot e^{-\varepsilon m}.
\]
Setting the right-hand side $\leq \delta$ and solving for $m$ (using $\beta < 1$
to ensure that the exponential decay in $m$ dominates the polynomial growth of
$2^{m^\beta}$) gives the stated bound.
\end{proof}

\begin{remark}[Occam and compression]
\label{rmk:occam-compression}
Every compression scheme of size $k$ induces an Occam algorithm: the compressed
subsample has description length $O(k \log m)$, which is sub-linear in $m$.
Conversely, an Occam algorithm is not necessarily a compression scheme, because
the output hypothesis need not be determined by a subsample of the input.
Occam's razor is thus a relaxation of compression.
\end{remark}

The converse of the Occam theorem (does every PAC-learnable class admit an
Occam algorithm?) was partially answered by Board and Pitt:

\begin{thm}[Partial converse {\cite{BoardPitt1992}}]
\label{thm:board-pitt}
If the class of all polynomial-size circuits is not PAC learnable by
polynomial-size circuits (a widely believed complexity-theoretic assumption),
then there exist PAC-learnable classes that do not admit polynomial-time Occam
algorithms.
\end{thm}

Thus the Occam property is \emph{strictly stronger} than PAC learnability
(under standard complexity assumptions): there exist learnable classes where
no efficient algorithm can find short descriptions of the data.  The gap is
purely computational: a short consistent description still exists, but the
work of producing it outruns any polynomial-time algorithm.


% ================================================================
% SECTION 4: INFORMATION-THEORETIC FOUNDATIONS
% Profile A and A-plus. Brief definitions.
% ================================================================

\section{Information-Theoretic Foundations}
\label{sec:info-foundations}

``Short description'' and ``efficiently coverable function class'' are the
informal language Occam's razor and compression schemes both rely on.  This
section gives each notion a precise, computable (or at least formal) definition.

\subsection{Description Length}

\begin{defn}[Description Length]
\label{def:description-length}
The \emph{description length} of a hypothesis $h$ in a representation class
$\mathcal{R}$ is the length $|h|$ of the shortest binary string encoding $h$ in
$\mathcal{R}$.  A \emph{prefix-free} representation ensures that
$\sum_{h \in \mathcal{R}} 2^{-|h|} \leq 1$ (the Kraft inequality).
\end{defn}

Description length connects to PAC learning through the observation that a
class of hypotheses with bounded description lengths has bounded effective size:
$|\{h \in \mathcal{R} : |h| \leq L\}| \leq 2^L$.  This makes union bounds
over the class tractable.

\subsection{Kolmogorov Complexity}

\begin{defn}[Kolmogorov Complexity]
\label{def:kolmogorov}
The \emph{Kolmogorov complexity} $K(x)$ of a string $x$ is the length of the
shortest program (on a fixed universal Turing machine) that outputs $x$ and
halts.  The \emph{conditional Kolmogorov complexity} $K(x \mid y)$ is the length
of the shortest program that outputs $x$ given $y$ as input.
\end{defn}

Kolmogorov complexity is uncomputable but provides the theoretical gold standard
for description length: $K(x)$ is the ultimate compression of $x$.

\subsection{KL Complexity}

\begin{defn}[KL Complexity]
\label{def:kl-complexity}
The \emph{Kullback--Leibler divergence} (KL divergence) between distributions
$P$ and $Q$ on a measurable space is
\[
  \KL(P \| Q) = \E_P\!\left[\log\frac{dP}{dQ}\right],
\]
when $P \ll Q$, and $+\infty$ otherwise.  In the PAC-Bayes framework
(\Cref{ch:generalization}), $\KL(\rho \| \pi)$ measures the ``complexity'' of a
posterior $\rho$ relative to a prior $\pi$: it is the number of additional
nats needed to encode a hypothesis drawn from $\rho$ using a code designed for
$\pi$.
\end{defn}

\subsection{Covering Numbers}

\begin{defn}[Covering Number]
\label{def:covering-number}
Let $(\mathcal{F}, d)$ be a (pseudo)metric space.  The \emph{$\varepsilon$-covering
number} $\mathcal{N}(\varepsilon, \mathcal{F}, d)$ is the minimum number of
balls of radius $\varepsilon$ (under $d$) needed to cover $\mathcal{F}$.
\end{defn}

Covering numbers measure the ``effective size'' of a function class at
resolution $\varepsilon$.  They connect to sample complexity through the
classical chaining argument:

\begin{thm}[Dudley's entropy integral {\cite{Dudley1967}}]
\label{thm:dudley}
Let $\mathcal{F}$ be a class of functions $f \colon X \to [0,1]$, and let
$d_n(f,g) = \bigl(\frac{1}{n}\sum_{i=1}^n (f(x_i) - g(x_i))^2\bigr)^{1/2}$
be the empirical $L^2$ metric on a sample $(x_1, \ldots, x_n)$.  Then
\[
  \E\!\left[\sup_{f \in \mathcal{F}} \left|\frac{1}{n}\sum_{i=1}^n f(x_i)
    - \E[f]\right|\right]
  \leq \frac{C}{\sqrt{n}} \int_0^1
    \sqrt{\log \mathcal{N}(\varepsilon, \mathcal{F}, d_n)}\; d\varepsilon,
\]
where $C$ is a universal constant.
\end{thm}

The integral $\int_0^1 \sqrt{\log \mathcal{N}(\varepsilon, \mathcal{F}, d_n)}\;
d\varepsilon$ is the \emph{Dudley entropy integral}.  It bounds the expected
supremum of the empirical process indexed by $\mathcal{F}$, the quantity that
controls uniform convergence and hence generalization.  When
$\mathcal{N}(\varepsilon, \mathcal{F}, d_n) \leq (C/\varepsilon)^d$ (as for
VC classes), the integral evaluates to $O(\sqrt{d})$, recovering the standard
sample complexity bounds.


% ================================================================
% SECTION 5: CHAPTER SUMMARY AND OPEN QUESTIONS
% ================================================================

\section{The State of the Art}
\label{sec:compression-summary}

The chapter's results cluster around a single qualitative fact and a single open
question.  The qualitative fact is settled and kernel-verified: PAC learnability,
finite VC dimension, and the existence of a finite compression scheme with side
information are equivalent.  The open question is how short the scheme can be
made.  \Cref{fig:compression-landscape} shows where each result sits and which
descent steps remain unproven.

\begin{figure}[ht]
\centering
\begin{tikzpicture}[
    box/.style={draw, thick, rounded corners, minimum width=2.5cm,
      minimum height=0.85cm, font=\small, align=center},
    verified/.style={box, fill=defblue!12},
    literature/.style={box, fill=defblue!4, draw=gray!55},
    openq/.style={box, dashed, fill=edgeorange!10, text=edgeorange!85!black},
    refuted/.style={box, fill=thmred!10, draw=thmred!70},
    veq/.style={{Stealth[length=2.5mm]}-{Stealth[length=2.5mm]}, thick, defblue!75},
    vedge/.style={-{Stealth[length=2.5mm]}, thick, defblue!75},
    litedge/.style={-{Stealth[length=2.5mm]}, semithick, gray!60},
    openedge/.style={-{Stealth[length=2.5mm]}, thick, dashed, edgeorange},
    rededge/.style={-{Stealth[length=2.5mm]}, thick, dashed, thmred},
    lbl/.style={font=\scriptsize, midway, fill=white, inner sep=1pt},
]
% --- verified equivalence core: THAT a short summary exists (kernel sc=0) ---
\node[verified] (pac) at (0,0) {PAC learnable};
\node[verified] (vc) at (3.5,0) {$\VC<\infty$};
\node[verified] (comp) at (7.0,0) {Finite\\compression\\{\scriptsize with side info}};
\draw[veq] (pac) -- (vc);
\draw[veq] (vc) -- (comp);

% --- Occam: a literature alternative route to the same qualitative conclusion ---
\node[literature] (occam) at (3.5,-2.0) {Occam algorithm};
\draw[litedge] (comp) -- node[lbl, above right] {implies} (occam);
\draw[litedge] (occam) -- node[lbl, below left] {\cref{thm:occam}} (pac);

% --- the size ladder: HOW short? descent = a strictly stronger claim ---
\node[literature] (sexp) at (7.0,-2.0) {Size $2^{O(d)}$};
\node[openq] (spoly) at (7.0,-3.4) {Size $\mathrm{poly}(d)$?};
\node[openq] (slin) at (7.0,-4.8) {Size $O(d)$?};
\draw[litedge] (comp) -- node[lbl, right] {\cite{MoranYehudayoff2016}} (sexp);
\draw[openedge] (sexp) -- node[lbl, right] {\cref{op:poly-compression}} (spoly);
\draw[openedge] (spoly) -- node[lbl, right] {\cref{op:compression}} (slin);

% --- side-info axis: an open strengthening of the verified core ---
\node[openq] (noinfo) at (10.1,-2.0) {Drop side info?};
\draw[openedge] (comp) -- node[lbl, pos=0.72] {\cref{op:side-info}} (noinfo);

% --- the unlabeled variant: bounded size known, O(d) open ---
\node[openq] (unlab) at (7.0,-6.2) {Unlabeled $O(d)$?};
\draw[openedge] (slin) -- node[lbl, right] {\cite{KuzminWarmuth2007}} (unlab);

% --- line-style legend: the evidence grade ---
\begin{scope}[shift={(0,-3.7)}]
  \draw[vedge] (0,0) -- (0.6,0);
  \node[font=\scriptsize, anchor=west] at (0.72,0) {kernel-verified ($sc{=}0$)};
  \draw[litedge] (0,-0.5) -- (0.6,-0.5);
  \node[font=\scriptsize, anchor=west] at (0.72,-0.5) {literature (cited)};
  \draw[openedge] (0,-1.0) -- (0.6,-1.0);
  \node[font=\scriptsize, anchor=west] at (0.72,-1.0) {open};
\end{scope}
\end{tikzpicture}
\caption[Compression and learnability.]{The top row is settled and kernel-verified
  (\leanname{fundamental_vc_compression}): PAC learnability, finite VC dimension, and
  the existence of a finite compression scheme with side information are all equivalent.
  This three-way equivalence is the qualitative core; it certifies that a short summary
  exists but says nothing about its length. The vertical descent tracks how short. A
  compression of size $2^{O(d)}$ is known, a result that was open until Moran and
  Yehudayoff established it in 2016; whether the size can reach $\mathrm{poly}(d)$ or
  the conjectured $O(d)$ is unresolved. The right branch asks whether the side information
  can be dropped at the same size, also open. Unlabeled compression of bounded size exists
  for every VC class; whether it reaches $O(d)$ is the Kuzmin--Warmuth question, open in
  general and settled only for maximum classes.}
\label{fig:compression-landscape}
\end{figure}

The following table summarizes the best known compression bounds:

\begin{table}[ht]
\centering\small
\caption{Known compression bounds by class family.}
\label{tab:compression-bounds}
\begin{tabular}{lcc}
\toprule
\textbf{Class family} & \textbf{VC dim.} & \textbf{Compression size} \\
\midrule
Maximum classes & $d$ & $d$ \\
Intersection-closed & $d$ & $d$ \\
Halfspaces in $\R^d$ & $d+1$ & $d+1$ \\
VC dim.\ 1 & 1 & 1 \\
Arbitrary VC class & $d$ & $2^{O(d)}$ \\
Unlabeled (arbitrary) & $d$ & No $f(d)$ bound \\
\bottomrule
\end{tabular}
\end{table}

\begin{openprob}[Polynomial compression]
\label{op:poly-compression}
Does every concept class of VC dimension $d$ admit a compression scheme of
size $\mathrm{poly}(d)$?  Even this weaker form of the compression conjecture
remains open.
\end{openprob}

Tight sample complexity, compression, and Occam's razor all express a single
structural fact: \emph{learnability is equivalent to the existence of a short
summary of the evidence}.  The VC dimension certifies that such a summary
exists.  The compression conjecture asks how short it can be.  After forty
years, that question stands unanswered.


% ================================================================
% OPEN PROBLEMS
% ================================================================

\section{Open Problems}
\label{sec:compression-open}

The $O(d)$ conjecture (\Cref{op:compression}) is the gravitational center, but
the frontier it anchors is wider.  The following ten problems branch off it,
some sharply posed with the obstruction articulated, others still missing the
right invariant or vocabulary.

\begin{openprob}[Compression without side information]
\label{op:side-info}
The existence theorem produces a compression scheme \emph{with} side
information: the reconstruction is handed $O(d)$ extra bits beyond the subsample.
Does every class of VC dimension $d$ admit a scheme of size $2^{O(d)}$ with
\emph{no} side bits at all (\leanname{Open_NoInfoCompressionStrengthening})?  A
positive answer is the first step toward
\Cref{op:compression}~\cite{MoranYehudayoff2016}.
\end{openprob}

\begin{openprob}[Orientations of non-maximum graphs]
\label{op:oig-orientation}
\Cref{thm:one-inclusion-orientation} orients the one-inclusion graph of a
\emph{maximum} class with in-degree $\leq d$.  Does every (non-maximum) class of
VC dimension $d$ admit an orientation of in-degree $\mathrm{poly}(d)$?  An
affirmative answer would yield polynomial compression.  No orientation theorem
is known off the maximum case: the combinatorial object is in hand, the extremal
lever that tames it is not~\cite{KuzminWarmuth2007}.
\end{openprob}

\begin{openprob}[Agnostic compression]
\label{op:agnostic-compression}
The schemes of this chapter are realizable: they reconstruct a hypothesis
\emph{consistent} with the sample.  Does every VC-$d$ class admit an
\emph{agnostic} compression scheme of size $O(d)$, one whose reconstruction has
error at most $\mathrm{opt} + \varepsilon$ when no consistent hypothesis
exists?  Sharpened by David, Moran, and
Yehudayoff~\cite{DavidMoranYehudayoff2016}.
\end{openprob}

\begin{openprob}[Compression for attention classes]
\label{op:attention-compression}
\Cref{prop:attention-compression} certifies that the softmax-attention class
$\mathcal{A}_{d,n_K}$ is compressible at size $2^{O(d)}$
(\leanname{softAttention_wellBehaved}), and the support-key heuristic predicts
$O(n_K)$.  Does $\mathcal{A}_{d,n_K}$ admit a compression scheme of size
$O(n_K)$, or even $O(\VC)$?  This is \Cref{op:compression} instantiated for a
contemporary architecture, with no compression invariant specific to attention
yet known~\cite{MoranYehudayoff2016}.
\end{openprob}

\begin{openprob}[Compression and recursive teaching]
\label{op:rtd}
The recursive teaching dimension $\mathrm{RTD}(\mathcal{H})$ is sandwiched with
the VC dimension but the two can differ.  Is the optimal compression size
$\Theta(\mathrm{RTD})$ rather than $\Theta(\VC)$?  This would re-route the
conjecture through teaching rather than shattering~\cite{SimonZilles2015}.
\end{openprob}

\begin{openprob}[Stable and replicable compression]
\label{op:stable-compression}
Does every VC class admit a compression scheme that is \emph{stable} under
resampling, such that changing one training example changes $O(1)$ of the selected
indices?  Stability would link compression to differential privacy and
replicability, where an equivalence with online learnability is already
known~\cite{BunLivniMoran2020}.  The stability invariant for compression has no
accepted definition yet.
\end{openprob}

\begin{openprob}[List compression and the DS dimension]
\label{op:list-compression}
In multiclass learning, finite Daniely--Shalev-Shwartz (DS) dimension
characterizes learnability.  Is finite DS dimension equivalent to a bounded
\emph{list} compression scheme, where reconstruction outputs a short list of
hypotheses containing a consistent one?  The multiclass compression vocabulary
is still being built~\cite{BrukhimCarmonDinurMoranYehudayoff2022}.
\end{openprob}

\begin{openprob}[Fat-shattering compression]
\label{op:fat-compression}
For $[0,1]$-valued function classes, learnability is governed by the
fat-shattering dimension $\mathrm{fat}_\gamma$.  Does finite $\mathrm{fat}_\gamma$
imply a real-valued compression scheme of size $\mathrm{poly}(\mathrm{fat}_\gamma)$?
The Boolean machinery does not transfer verbatim, and the scale parameter
$\gamma$ has no settled role in the compressed object (Alon, Ben-David,
Cesa-Bianchi, and Haussler).
\end{openprob}

\begin{openprob}[Super-linear lower bounds]
\label{op:compression-lower}
Every known VC-$d$ class compresses to $O(d)$ points; the obstruction to the
conjecture is constructing schemes, not ruling them out.  Is there a VC-$d$
class provably \emph{requiring} compression size $\omega(d)$?  No super-linear
lower bound is known, and a resolution either way would be decisive.
\end{openprob}

\begin{openprob}[Boosting weak compression]
\label{op:weak-compression}
Call a scheme \emph{weak} if its reconstruction is correct on a
$\tfrac{1}{2}+\gamma$ fraction of relabelings of the sample.  Does weak
compression boost to full compression with only $O(\log(1/\gamma))$ overhead in
size, mirroring the boosting of weak PAC learners?  A boosting operator on
compression schemes has not been formulated~\cite{MoranYehudayoff2016}.
\end{openprob}


% ================================================================
% EXERCISES
% ================================================================

\subsection*{Exercises}

\begin{enumerate}[leftmargin=*,itemsep=6pt]

\item \textbf{Compression lower bound.}
  Let $\mathcal{H}$ be a concept class with $\VC(\mathcal{H}) = d$, and
  let $(\kappa, \rho)$ be a sample compression scheme of size $k$ for
  $\mathcal{H}$.  Prove that $k \geq d$.

  \emph{Hint:} Let $C = \{x_1, \ldots, x_d\}$ be a shattered set.
  For each of the $2^d$ labelings $b \in \{0,1\}^d$, the compression
  function selects a subsequence $\kappa(S_b) \subseteq S_b$ of size
  $\leq k$.  Show that if $k < d$, a pigeonhole argument forces two
  distinct labelings $b \neq b'$ to compress to the same subsequence
  $\kappa(S_b) = \kappa(S_{b'})$ with the same labels, contradicting
  the requirement that reconstruction produces hypotheses consistent
  with both $S_b$ and $S_{b'}$.

\item \textbf{Compression for unions of intervals.}
  Let $\mathcal{H}_k$ be the class of unions of at most $k$ disjoint
  intervals on $[0,1]$: $h \in \mathcal{H}_k$ iff $h^{-1}(1)$ is a
  union of at most $k$ intervals.
  \begin{enumerate}[label=(\alph*)]
  \item Prove that $\VC(\mathcal{H}_k) = 2k$.
  \item Construct an explicit compression scheme of size $2k$ for
    $\mathcal{H}_k$.  (\emph{Hint:} Generalize
    \Cref{ex:oig-intervals}; the compressed subsample consists of
    the $2k$ boundary points of the positive regions.)
  \item Use the result of Exercise~1 to conclude that no compression
    scheme of size $< 2k$ exists.  This verifies the compression
    conjecture for $\mathcal{H}_k$ with equality.
  \end{enumerate}

\item \textbf{Compression and teaching dimension.}
  The \emph{teaching dimension} $\mathrm{TD}(\mathcal{H})$ is the
  minimum $k$ such that every concept $h \in \mathcal{H}$ has a
  \emph{teaching set} of size $k$: a set $T \subseteq X$ with
  $|T| = k$ such that $h$ is the unique concept in $\mathcal{H}$
  consistent with $h|_T$.
  \begin{enumerate}[label=(\alph*)]
  \item Show that if $\mathcal{H}$ has teaching dimension $k$, then
    $\mathcal{H}$ admits a compression scheme of size $k$.  (The
    teaching set is the compressed subsample.)
  \item Show that a compression scheme of size $k$ does \emph{not}
    imply teaching dimension $\leq k$: construct a class $\mathcal{H}$
    with a compression scheme of size $1$ but teaching dimension
    $> 1$.  (\emph{Hint:} The compression function may depend on the
    full sample, not just on the target concept; the teaching set must
    work for a single concept without knowing the other data.)
  \item Prove that $\mathrm{TD}(\mathcal{H}) \geq \VC(\mathcal{H})$
    for all classes $\mathcal{H}$.  Combined with the compression
    conjecture, this would imply
    $\VC \leq \text{compression size} \leq \mathrm{TD}$: compression
    sits between VC dimension and teaching dimension.  But whether
    the compression conjecture implies
    $\text{compression size} \leq O(\mathrm{TD})$ or
    $\text{compression size} \leq O(\VC)$ remains open; these are
    distinct (and incomparable) questions.
  \end{enumerate}

\end{enumerate}

\subsection*{Counterfactual exercises}

\Cref{fig:compression-landscape} encodes an argument: each node carries a claim
at a specific evidence grade (kernel-verified, literature-cited, or open).  The exercises below stress-test that structure: take a known result,
vary one of its hypotheses or conclusions, and track exactly where the shifted
claim lands.  Each variation reveals why a rung on the size ladder is a real
boundary, why the side-information qualifier matters, and why the distance from
$2^{O(d)}$ to $O(d)$ has resisted forty years of effort.

\begin{enumerate}[leftmargin=*,itemsep=3pt]
\item \textbf{The qualitative core.} Show that the top row is a single three-way equivalence: PAC learnability, finite VC dimension, and the existence of \emph{some} finite compression scheme with side information (\cref{thm:fundamental}, \leanname{fundamental_vc_compression}). Explain why this asserts \emph{that} a short summary exists but says nothing about its length.

\item \textbf{The easy half.} Prove that any finite compression scheme forces finite VC dimension, with the explicit bound $\VC(\mathcal{H}) \le 2(k+1)^2-1$ (\cref{thm:compression-pac}, \leanname{compression_imp_vcdim_finite}). Contrast this settled direction with the converse, finite VC implying \emph{small} compression, which splinters into the literature bound and the open conjectures.

\item \textbf{Halfspaces hit the target.} For halfspaces in $\R^d$, the maximum-margin hyperplane is determined by at most $d+1$ support vectors (\cref{ex:svm-compression}), so the family compresses to size $d+1$ at $\VC=d+1$, attaining the $O(d)$ target. Argue that this support-vector mechanism has no known analogue for an arbitrary VC class.

\item \textbf{Maximum classes collapse the ladder.} For a maximum class of VC dimension $d$ (growth function meeting Sauer--Shelah with equality), show that compression to exactly $d$ points exists, so the rungs $2^{O(d)}$, $\mathrm{poly}(d)$, and $O(d)$ coincide. The vertical separation between rungs is a statement about \emph{arbitrary} classes; explain why it vanishes here.

\item \textbf{The side-information axis.} The existence theorem produces a scheme \emph{with} side information (\leanname{fundamental_vc_compression}); whether the side bits can be removed at the same size is open (\cref{op:side-info}, \leanname{Open_NoInfoCompressionStrengthening}). Show that the two endpoints differ by a structural qualifier alone, the size held fixed.

\item \textbf{The exponential gap.} The best known bound is $2^{O(d)}$ while the conjecture demands $O(d)$ (\cref{rmk:compression-gap}): at $d=20$ this is a gap from $\sim 10^6$ to $\sim 20$. Explain why the $2^{O(d)}$ bound (\cref{thm:moran-yehudayoff}, \cite{MoranYehudayoff2016}), unknown before 2016, is a literature result rather than a refinement of the conjecture.

\item \textbf{The unlabeled variant.} Bounded-size unlabeled compression exists for every VC class (\cref{def:unlabeled-compression}, \cref{thm:unlabeled-bounded}); whether it reaches the labeled $O(d)$ bound is the Kuzmin--Warmuth question, open in general and settled for maximum classes. Explain what dropping the stored labels costs, and why it leaves the size question open rather than closing it.

\item \textbf{Compression as a dimension functor.} Read ``the smallest summary that reconstructs the data'' as a universal property: the optimal compression size as the minimal generating section of an evidence-functor on finite samples whose rank is the VC dimension. Discuss whether such a single dimension functor exists, unifying the size ladder, the side-information qualifier, and the alternative dimensions (recursive teaching, DS), or whether the notions are provably non-unifiable (\cref{op:dimension-functor}, \cref{op:rtd}).
\end{enumerate}

% Chapter 12: Generalization Bounds
%   uniform_convergence: Profile B (load-bearing theorem: symmetrization
%     + concentration proof IS the content)
%   Five bound traditions: Profile G (framework bridge: parallax)
%   rademacher_complexity: Profile A-to-B bridge (definition matters, but
%     comparison to VC is the content)
%   pac_bayes_bound: Profile E (obstruction: looks Bayesian but prior
%     has a fundamentally different role)
%
% Centerpiece: the symmetrization argument (Sec 2) and the five-framework
% parallax (Sec 8).

\chapter{Generalization Bounds}
\label{ch:generalization}

\Cref{ch:pac} established \emph{that} finite VC dimension characterizes PAC
learnability.  \Cref{ch:compression} refined the quantitative dependence
through compression.  This chapter asks a different question: given a
specific learning algorithm~$A$ and a specific training sample~$S$, what
can we say about the generalization error of the hypothesis~$A(S)$?

The answer is not unique.  Five distinct research traditions have developed
bounds on the same quantity (the gap between true and empirical
risk), each using different information about the learner:

\begin{enumerate}[leftmargin=*, itemsep=3pt]
\item \textbf{VC/uniform convergence} measures the combinatorial richness
  of the hypothesis class through the growth function.
\item \textbf{Rademacher complexity} refines this to a data-dependent
  measure of how well the class correlates with random noise.
\item \textbf{PAC-Bayes bounds} measure the ``distance'' of a learned
  posterior from a data-independent prior.
\item \textbf{Algorithmic stability} measures how much the output changes
  when a single training example is perturbed.
\item \textbf{Margin theory} measures the ``confidence'' of predictions
  and connects to scale-sensitive dimensions.
\end{enumerate}

\noindent
We also discuss information-theoretic bounds, which measure the mutual
information between the training sample and the learned hypothesis.

Each framework produces a bound on the same quantity but requires
different structural assumptions and succeeds in different regimes.  The
centerpiece of this chapter is twofold: (i)~the \emph{symmetrization
argument} that powers uniform convergence (\Cref{sec:uniform-convergence}),
proved in full as the engine behind classical generalization theory; and
(ii)~the \emph{comparison} among the five frameworks
(\Cref{sec:comparison}), which reveals what each sees that the others miss.

\medskip

Throughout this chapter, $\mathcal{H}$ denotes a hypothesis class,
$D$~a distribution on $X \times Y$,
$S = \{(x_i, y_i)\}_{i=1}^m$ an i.i.d.\ sample from~$D$,
$A$~a learning algorithm, and
$\ell \colon \mathcal{H} \times (X \times Y) \to [0,1]$ a bounded loss
function.


% ================================================================
% SECTION 1: GENERALIZATION ERROR
% Profile A.  Definition.
% ================================================================

\section{Generalization Error}
\label{sec:gen-error}

\begin{defn}[Generalization Error]
\label{def:gen-error}
The \emph{generalization error} (or \emph{generalization gap}) of a
hypothesis $h$ with respect to distribution~$D$ and sample~$S$ is
\[
  \mathrm{gen}(h, S) \;=\; \risk_D(h) - \emprisk_S(h)
    \;=\; \E_{(x,y) \sim D}[\ell(h, (x,y))]
    - \frac{1}{m}\sum_{i=1}^m \ell(h, (x_i, y_i)).
\]
A \emph{generalization bound} is any high-probability or in-expectation
upper bound on $|\mathrm{gen}(h, S)|$ or on $\mathrm{gen}(A(S), S)$.
\end{defn}

The generalization error is a random variable (through~$S$), and when
$h = A(S)$ depends on~$S$, controlling it requires accounting for the
dependence between the hypothesis and the data that produced it.  This is
the core technical challenge that each framework addresses differently.


% ================================================================
% SECTION 2: UNIFORM CONVERGENCE
% Profile B (load-bearing theorem).  FULL proof.
% ================================================================

\section{Uniform Convergence: The Engine}
\label{sec:uniform-convergence}

The classical approach to generalization does not look at the algorithm at
all.  Instead, it asks: does the empirical risk converge to the true risk
\emph{uniformly over all hypotheses}?  If so, any hypothesis with small
empirical risk automatically has small true risk, including the one the
algorithm selects.

\begin{defn}[Uniform Convergence]
\label{def:uniform-convergence}
A hypothesis class~$\mathcal{H}$ has the \emph{uniform convergence
property} with respect to loss~$\ell$ if, for every $\varepsilon > 0$,
\[
  \sup_{h \in \mathcal{H}} |\risk_D(h) - \emprisk_S(h)|
  \xrightarrow{\;m \to \infty\;} 0
  \quad \text{in probability, uniformly over all distributions } D.
\]
\end{defn}

The content of uniform convergence is the proof technique: the
\emph{symmetrization argument} introduced by Vapnik and Chervonenkis in
1971.  The argument proceeds in three stages, each with a distinct role.

\begin{thm}[VC Uniform Convergence Bound]
\label{thm:vc-uniform}
Let $\mathcal{H} \subseteq \{0,1\}^X$ with $\VC(\mathcal{H}) = d < \infty$.
Then for any distribution~$D$ on $X \times \{0,1\}$, any $\varepsilon > 0$,
and any $\delta > 0$, with probability at least $1 - \delta$ over
$S \sim D^m$:
\[
  \sup_{h \in \mathcal{H}} |\risk_D(h) - \emprisk_S(h)|
  \;\leq\; \sqrt{\frac{8\,d \log(2em/d) + 8\log(4/\delta)}{m}}.
\]
\end{thm}

\begin{proof}
The proof has three stages: symmetrization, reduction to a finite
problem, and concentration.

\textbf{Stage 1: Symmetrization (the ghost sample).}
Let $S = (z_1, \ldots, z_m)$ be the training sample and let
$S' = (z_1', \ldots, z_m')$ be an independent ``ghost'' sample drawn
from the same distribution.  Define
$\Phi(S) = \sup_{h \in \mathcal{H}} (\risk_D(h) - \emprisk_S(h))$.

For any fixed~$h$, the true risk equals
$\risk_D(h) = \E_{S'}[\emprisk_{S'}(h)]$, so
\[
  \Phi(S) = \sup_{h \in \mathcal{H}} \E_{S'}\!\left[
    \emprisk_{S'}(h) - \emprisk_S(h)\right].
\]
Moving the supremum inside the expectation (Jensen's inequality applied
to the convex functional $\sup$) gives:
\begin{equation}
\label{eq:symm-step}
  \E_S[\Phi(S)]
  \;\leq\; \E_{S,S'}\!\left[\sup_{h \in \mathcal{H}}
    \bigl(\emprisk_{S'}(h) - \emprisk_S(h)\bigr)\right].
\end{equation}

\emph{Why this helps.}  The left side involves the population
quantity~$\risk_D(h)$, which we cannot compute.  The right side involves
only \emph{two finite samples}; no population quantities remain.  This is
the symmetrization trick: the ghost sample replaces the unknown
distribution.

\textbf{Stage 1b: Rademacher sign-flip.}
The difference $\emprisk_{S'}(h) - \emprisk_S(h) =
\frac{1}{m}\sum_{i=1}^m (\ell(h,z_i') - \ell(h,z_i))$ is a sum of
symmetric random variables (since $z_i$ and $z_i'$ are identically
distributed, swapping them does not change the joint distribution).
Introducing i.i.d.\ Rademacher variables $\sigma_1,\ldots,\sigma_m$
(uniform on $\{-1,+1\}$):
\begin{align}
  \E_{S,S'}\!\left[\sup_{h} \frac{1}{m}\sum_{i=1}^m
    \bigl(\ell(h,z_i') - \ell(h,z_i)\bigr)\right]
  &= \E_{S,S',\sigma}\!\left[\sup_{h} \frac{1}{m}\sum_{i=1}^m
    \sigma_i\bigl(\ell(h,z_i') - \ell(h,z_i)\bigr)\right]
    \nonumber \\
  &\leq 2\,\E_{S,\sigma}\!\left[\sup_{h} \frac{1}{m}\sum_{i=1}^m
    \sigma_i \ell(h,z_i)\right],
    \label{eq:rademacher-step}
\end{align}
where the last step uses the triangle inequality for suprema and the
identical distribution of $S$ and~$S'$.

\textbf{Stage 2: Reduction to a finite class.}
The supremum in \eqref{eq:rademacher-step} ranges over all
$h \in \mathcal{H}$, but on the fixed sample~$S$, only finitely many
distinct behavior patterns occur.  Specifically, the loss vectors
$(\ell(h,z_1), \ldots, \ell(h,z_m))$ take at most $\growth_\mathcal{H}(m)$
distinct values, where $\growth_\mathcal{H}$ is the growth function.

For binary classification with 0--1 loss, each behavior pattern is a
binary vector in $\{0,1\}^m$, and the Sauer--Shelah lemma
(\Cref{ch:dimensions}) gives
\begin{equation}
\label{eq:sauer-shelah-app}
  \growth_\mathcal{H}(m) \;\leq\; \left(\frac{em}{d}\right)^d.
\end{equation}

Applying Massart's finite lemma to the at most $\growth_\mathcal{H}(m)$
distinct vectors, each with Euclidean norm at most
$\sqrt{m}$,\formal{FLT_Proofs.Complexity.Rademacher.finite_massart_lemma}
\begin{equation}
\label{eq:massart}
  \E_\sigma\!\left[\sup_{h} \frac{1}{m}\sum_{i=1}^m \sigma_i \ell(h,z_i)
  \right]
  \;\leq\; \sqrt{\frac{2\log \growth_\mathcal{H}(m)}{m}}
  \;\leq\; \sqrt{\frac{2d\log(em/d)}{m}}.
\end{equation}

\textbf{Stage 3: Concentration.}
The bound so far controls $\E[\Phi(S)]$.  To obtain a high-probability
bound, observe that $\Phi(S) = \sup_h (\risk_D(h) - \emprisk_S(h))$
satisfies the bounded-differences condition: replacing any single~$z_i$
changes $\Phi(S)$ by at most $1/m$ (since each loss value lies in
$[0,1]$).  By McDiarmid's inequality,
\[
  \Pr\!\left[\Phi(S) > \E[\Phi(S)] + t\right]
  \leq \exp(-2mt^2).
\]
Setting $t = \sqrt{\log(2/\delta)/(2m)}$ and combining with
\eqref{eq:rademacher-step}--\eqref{eq:massart}, the one-sided bound
follows.  Applying the same argument to
$\sup_h (\emprisk_S(h) - \risk_D(h))$ and taking a union bound yields
the two-sided result.
\end{proof}

\begin{historical}
The symmetrization argument was introduced by Vapnik and Chervonenkis
(1971) in the original paper that defined what is now called the VC
dimension.  The three-stage structure (ghost sample, sign-flip,
finite reduction) was not initially presented in this clean form.  The
modern presentation, which passes through Rademacher complexity as an
intermediate quantity, crystallized in the work of Koltchinskii (2001)
and Bartlett--Mendelson (2002).  The Rademacher step (Stage~1b) was
implicit in the original argument but was not isolated as a self-standing
concept until much later.
\end{historical}

\begin{traversal}
The uniform convergence theorem ties three threads together: the growth
function characterizes uniform convergence, the Sauer--Shelah lemma is the
step that turns this into the VC generalization bound, and the VC dimension
upper-bounds the Rademacher complexity.  The proof above traces exactly this
path: the growth function enters at Stage~2, Sauer--Shelah provides the
polynomial bound, and the Rademacher intermediate appears at Stage~1b.
\end{traversal}


\begin{example}[Halfspaces in $\R^d$]
\label{ex:halfspace-uc}
Let $\mathcal{H} = \{x \mapsto \ind[\langle w, x \rangle \geq b]
: w \in \R^d, b \in \R\}$ be the class of halfspaces.
Then $\VC(\mathcal{H}) = d+1$ (\Cref{ch:dimensions}), and the uniform
convergence bound gives: with probability $\geq 1 - \delta$,
\[
  \sup_{h \in \mathcal{H}} |\risk_D(h) - \emprisk_S(h)|
  \leq \sqrt{\frac{8(d+1)\log(2em/(d+1)) + 8\log(4/\delta)}{m}}.
\]
For $d = 10$ and $m = 10{,}000$, this gives a bound of approximately
$0.12$, a nontrivial and meaningful result.  For a two-layer neural network
with $p = 10^6$ parameters (where naive VC bounds give
$\VC \approx O(p \log p)$), the same formula gives a bound exceeding~$1$:
the uniform convergence bound is \emph{vacuous}.
\end{example}


% ================================================================
% SECTION 3: RADEMACHER COMPLEXITY
% Profile A-to-B bridge.  Definition + comparison to VC.
% ================================================================

\section{Rademacher Complexity}
\label{sec:rademacher}

The Rademacher intermediate that appeared in Stage~1b of the
symmetrization proof is itself a fundamental complexity measure.  It
refines VC dimension in two ways: it is \emph{data-dependent} (computed on
the actual sample, not worst-case) and it applies to \emph{real-valued}
function classes (not just binary hypotheses).

\begin{defn}[Empirical Rademacher Complexity]
\label{def:rademacher}
Let $\mathcal{F}$ be a class of functions $f \colon Z \to \R$ and let
$S = (z_1, \ldots, z_m)$ be a fixed sample.  The \emph{empirical
Rademacher complexity} of $\mathcal{F}$ on~$S$ is
\[
  \Rad_S(\mathcal{F}) = \E_{\sigma}\!\left[\sup_{f \in \mathcal{F}}
    \frac{1}{m}\sum_{i=1}^m \sigma_i f(z_i)\right],
\]
where $\sigma_1, \ldots, \sigma_m$ are independent Rademacher random
variables (uniform on $\{-1, +1\}$).  The \emph{Rademacher complexity} of
$\mathcal{F}$ at sample size~$m$ is
$\mathcal{R}_m(\mathcal{F}) = \E_S[\Rad_S(\mathcal{F})]$.\formal{FLT_Proofs.Complexity.Rademacher.EmpiricalRademacherComplexity}
\end{defn}

The Rademacher complexity measures how well $\mathcal{F}$ can
\emph{correlate with pure noise}.  A class that achieves high correlation
with random $\pm 1$ labels is rich enough to fit arbitrary patterns,
including the noise in the training data.

\begin{thm}[Rademacher Generalization Bound
{\cite{BartlettMendelson2002}}]
\label{thm:rademacher-bound}
Let $\mathcal{F}$ be a class of functions $f \colon Z \to [0,1]$, let
$D$ be a distribution on~$Z$, and let $S = (z_1,\ldots,z_m) \sim D^m$.
Then for any $\delta > 0$, with probability at least $1-\delta$ over~$S$,
for all $f \in \mathcal{F}$ simultaneously:
\[
  \E_D[f] \leq \frac{1}{m}\sum_{i=1}^m f(z_i)
    + 2\,\Rad_S(\mathcal{F})
    + 3\sqrt{\frac{\log(2/\delta)}{2m}}.
\]
\end{thm}

\begin{proof}
The proof follows the same three-stage architecture as the uniform
convergence theorem: symmetrization, Rademacher sign-flip, and
concentration.

Define $\Phi(S) = \sup_{f \in \mathcal{F}} (\E_D[f] -
\frac{1}{m}\sum_i f(z_i))$.

\textbf{Step 1: Concentration of $\Phi$.}
Replacing any single $z_i$ in~$S$ changes $\Phi(S)$ by at most $1/m$
(since $f$ is bounded in $[0,1]$).  By McDiarmid's inequality:
\begin{equation}
\label{eq:rad-conc}
  \Pr[\Phi(S) > \E[\Phi(S)] + t] \leq \exp(-2mt^2).
\end{equation}
Setting $t = \sqrt{\log(2/\delta)/(2m)}$ gives
$\Phi(S) \leq \E[\Phi(S)] + \sqrt{\log(2/\delta)/(2m)}$
with probability $\geq 1 - \delta/2$.

\textbf{Step 2: Symmetrization.}
Let $S' = (z_1', \ldots, z_m')$ be an independent ghost sample.  Then
\begin{align*}
  \E_S[\Phi(S)]
  &= \E_S\!\left[\sup_{f \in \mathcal{F}} \left(\E_{S'}
    \!\left[\frac{1}{m}\sum_{i=1}^m f(z_i')\right]
    - \frac{1}{m}\sum_{i=1}^m f(z_i)\right)\right] \\
  &\leq \E_{S, S'}\!\left[\sup_{f \in \mathcal{F}}
    \frac{1}{m}\sum_{i=1}^m \bigl(f(z_i') - f(z_i)\bigr)\right].
\end{align*}
Since $z_i$ and $z_i'$ are i.i.d., the differences are symmetric
random variables, so introducing Rademacher signs does not change the
distribution:
\begin{align*}
  \E_{S, S'}\!\left[\sup_{f}
    \frac{1}{m}\sum_i \bigl(f(z_i') - f(z_i)\bigr)\right]
  &= \E_{S, S', \sigma}\!\left[\sup_{f}
    \frac{1}{m}\sum_i \sigma_i\bigl(f(z_i') - f(z_i)\bigr)\right] \\
  &\leq 2\,\E_{S, \sigma}\!\left[\sup_{f}
    \frac{1}{m}\sum_i \sigma_i f(z_i)\right]
  = 2\,\mathcal{R}_m(\mathcal{F}).
\end{align*}

\textbf{Step 3: Data-dependent form.}
Apply McDiarmid again: $\Rad_S(\mathcal{F})$ has bounded differences
$1/m$ (changing one sample point changes the Rademacher average by at
most $1/m$), so
$\mathcal{R}_m(\mathcal{F}) \leq \Rad_S(\mathcal{F})
  + \sqrt{\log(2/\delta)/(2m)}$
with probability $\geq 1 - \delta/2$.

A union bound over the two $\delta/2$ events yields the stated bound.
\end{proof}


\subsection{Rademacher Complexity versus VC Dimension}
\label{sec:rad-vs-vc}

The following proposition makes the relationship precise.

\begin{prop}[VC dimension bounds Rademacher complexity]
\label{prop:rad-vc}
For any binary hypothesis class $\mathcal{H} \subseteq \{0,1\}^X$ with
$\VC(\mathcal{H}) = d$:
\[
  \mathcal{R}_m(\mathcal{H})
  \;\leq\; \sqrt{\frac{2d \log(em/d)}{m}}.
\]
This quantitative bound is machine-checked.\formal{FLT_Proofs.Complexity.Rademacher.vcdim_bounds_rademacher_quantitative}
\end{prop}

\begin{proof}
On any fixed sample of size~$m$, the restriction $\mathcal{H}|_S$ has at
most $\growth_\mathcal{H}(m) \leq (em/d)^d$ elements (Sauer--Shelah).
Each element is a binary vector of Euclidean norm at most $\sqrt{m}$.
By Massart's finite lemma, the Rademacher complexity of a finite set~$V$
of vectors in $\R^m$ satisfies
$\E_\sigma[\sup_{v \in V} \frac{1}{m}\langle \sigma, v\rangle]
\leq \frac{\max_v \|v\|_2}{m}\sqrt{2\log|V|}$.
Substituting $|V| \leq (em/d)^d$ and $\|v\|_2 \leq \sqrt{m}$ gives
the result.
\end{proof}

\begin{remark}[When Rademacher is tighter than VC]
\label{rmk:rad-tighter}
The VC bound is worst-case: it uses the maximum of
$|\mathcal{H}|_S|$ over all samples of size~$m$.  The empirical
Rademacher complexity $\Rad_S(\mathcal{H})$ uses the \emph{actual}
restriction $\mathcal{H}|_S$, which may be much smaller.  Consider the
class of all threshold functions on $\R$:
$\mathcal{H} = \{x \mapsto \ind[x \geq \theta] : \theta \in \R\}$.
This class has $\VC(\mathcal{H}) = 1$ and
$\growth_\mathcal{H}(m) = m + 1$, so the VC bound gives
$\mathcal{R}_m \leq \sqrt{2\log(em)/m}$.  But on a sample where all
points lie in $\{0,1\}$, the restriction has at most $3$ elements, and the
empirical Rademacher complexity is $O(1/\sqrt{m})$ without the
logarithmic factor.  The gain is modest here but becomes significant for
classes with distribution-dependent structure.
\end{remark}

\begin{computational}
\textbf{Rademacher complexity of linear classifiers.}
Let $\mathcal{F} = \{x \mapsto \langle w, x \rangle : \|w\|_2 \leq W\}$
and suppose $\|x_i\|_2 \leq B$ for all $i$.  Then
\[
  \Rad_S(\mathcal{F})
  = \E_\sigma\!\left[\sup_{\|w\| \leq W}
    \frac{1}{m}\left\langle w, \sum_i \sigma_i x_i\right\rangle\right]
  = \frac{W}{m}\,\E_\sigma\!\left[\left\|\sum_i \sigma_i x_i\right\|_2
    \right]
  \leq \frac{WB}{\sqrt{m}},
\]
where the last step uses
$\E[\|\sum_i \sigma_i x_i\|_2^2] = \sum_i \|x_i\|_2^2 \leq mB^2$ and
Jensen's inequality.  Note: no dependence on ambient dimension.  A
linear classifier in $\R^{10^6}$ with bounded norm has the same
Rademacher complexity as one in $\R^{10}$; only the norm and the data
radius matter.
\end{computational}


% ================================================================
% SECTION 4: THE GROWTH FUNCTION AND SAUER--SHELAH
% Brief recall (proved in Ch 10).
% ================================================================

\section{The Growth Function}
\label{sec:growth-function}

The growth function $\growth_\mathcal{H}(m) = \max_{|S|=m}
|\mathcal{H}|_S|$ mediates between VC dimension and generalization bounds.
It is the growth function, not the VC dimension directly, that enters
the uniform convergence argument (Stage~2 of the proof of
\Cref{thm:vc-uniform}).  The VC dimension is a single number; the growth
function is a whole sequence.

The Sauer--Shelah lemma (\Cref{ch:dimensions}) provides the critical
link: if $\VC(\mathcal{H}) = d$, then
$\growth_\mathcal{H}(m) \leq \sum_{i=0}^d \binom{m}{i}
\leq (em/d)^d$ for $m \geq d$.\formal{FLT_Proofs.Complexity.Rademacher.sauer_shelah_exp_bound}
The dichotomy is sharp: either $\growth_\mathcal{H}(m) = 2^m$ for all~$m$
(when $\VC = \infty$) or $\growth_\mathcal{H}(m) = O(m^d)$
(when $\VC = d < \infty$).  There is no intermediate growth rate.


% ================================================================
% SECTION 5: PAC-BAYES BOUNDS
% Profile E (obstruction).  The prior is NOT a Bayesian prior.
% ================================================================

\section{PAC-Bayes Bounds}
\label{sec:pac-bayes}

The PAC-Bayes framework replaces the worst-case uniform convergence of
Rademacher bounds with a \emph{distribution over hypotheses}.  Instead of
bounding the risk of every $h \in \mathcal{H}$ simultaneously, it bounds
the expected risk of a hypothesis drawn from a learned posterior~$\rho$,
measured relative to a data-independent prior~$\pi$.

\begin{defn}[PAC-Bayes Setting]
\label{def:pac-bayes-setting}
Fix a hypothesis space~$\mathcal{H}$, a loss
$\ell \colon \mathcal{H} \times (X \times Y) \to [0,1]$,
a \emph{prior} distribution~$\pi$ over~$\mathcal{H}$ chosen
\emph{before} seeing data, and a training sample $S \sim D^m$.
A \emph{posterior}~$\rho$ is any distribution over~$\mathcal{H}$ that
may depend on~$S$.  Define the Gibbs (mixture)
risks\formal{FLT_Proofs.Theorem.PACBayes.gibbsError}
\[
  \risk_D(\rho) = \E_{h \sim \rho}[\risk_D(h)], \qquad
  \emprisk_S(\rho) = \E_{h \sim \rho}[\emprisk_S(h)].
\]
\end{defn}

\begin{thm}[McAllester's PAC-Bayes Bound
{\cite{McAllester1999}}]
\label{thm:pac-bayes}
For any prior~$\pi$ over~$\mathcal{H}$ (chosen independently of~$S$),
any $\delta > 0$, with probability at least $1 - \delta$ over
$S \sim D^m$: for all posteriors~$\rho$ over~$\mathcal{H}$
simultaneously,
\[
  \risk_D(\rho) \leq \emprisk_S(\rho)
    + \sqrt{\frac{\KL(\rho \| \pi) + \log(2\sqrt{m}/\delta)}{2m}}.
\]
For finite~$\mathcal{H}$, the bound (with a cross-entropy complexity term in
place of $\KL(\rho\|\pi)$) is machine-checked.\formal{FLT_Proofs.Theorem.PACBayes.pac_bayes_finite}
\end{thm}

\begin{proof}
The proof has four steps: change-of-measure, moment generating function
bound, prior expectation, and optimization.

\textbf{Step 1: Change-of-measure inequality.}
For any measurable function $\phi \colon \mathcal{H} \to \R$ and
distributions $\rho$,~$\pi$, the Donsker--Varadhan variational formula
gives
\[
  \E_{h \sim \rho}[\phi(h)] \leq \KL(\rho \| \pi)
    + \log \E_{h \sim \pi}[e^{\phi(h)}].
\]
This holds for every~$\rho$; it converts a bound on the moment generating
function under~$\pi$ into a bound under any~$\rho$, at the cost of
$\KL(\rho \| \pi)$.

\textbf{Step 2: Moment generating function bound.}
Set $\phi(h) = \lambda(\risk_D(h) - \emprisk_S(h))$ for a parameter
$\lambda > 0$ to be optimized.  For a fixed~$h$ (independent of~$S$),
$\emprisk_S(h) = \frac{1}{m}\sum_{i=1}^m \ell(h, z_i)$ is an average of
i.i.d.\ $[0,1]$-valued random variables with mean~$\risk_D(h)$.  By
Hoeffding's lemma,
\[
  \E_S\!\left[e^{\lambda(\risk_D(h) - \emprisk_S(h))}\right]
  \leq e^{\lambda^2/(8m)}.
\]

\textbf{Step 3: Prior expectation.}
Since $\pi$ is independent of~$S$, we can exchange the expectations:
\[
  \E_S\!\left[\E_{h \sim \pi}\!\left[
    e^{\lambda(\risk_D(h) - \emprisk_S(h))}\right]\right]
  = \E_{h \sim \pi}\!\left[\E_S\!\left[
    e^{\lambda(\risk_D(h) - \emprisk_S(h))}\right]\right]
  \leq e^{\lambda^2/(8m)}.
\]
By Markov's inequality applied to the non-negative random variable
$\E_{h \sim \pi}[e^{\lambda(\risk_D(h) - \emprisk_S(h))}]$: with
probability $\geq 1 - \delta$,
\[
  \log \E_{h \sim \pi}\!\left[
    e^{\lambda(\risk_D(h) - \emprisk_S(h))}\right]
  \leq \frac{\lambda^2}{8m} + \log\frac{1}{\delta}.
\]

\textbf{Step 4: Combine and optimize.}
Applying Step~1 to the event from Step~3: with probability
$\geq 1 - \delta$, for all~$\rho$,
\[
  \lambda\bigl(\risk_D(\rho) - \emprisk_S(\rho)\bigr)
  \leq \KL(\rho \| \pi) + \frac{\lambda^2}{8m} + \log\frac{1}{\delta}.
\]
Dividing by~$\lambda$ and optimizing
$\lambda = \sqrt{8m(\KL(\rho\|\pi) + \log(1/\delta))}$ yields
\[
  \risk_D(\rho) - \emprisk_S(\rho)
  \leq \sqrt{\frac{2(\KL(\rho \| \pi) + \log(1/\delta))}{m}}.
\]
The displayed change-of-measure argument yields the constant
$\sqrt{2(\KL(\rho\|\pi) + \log(1/\delta))/m}$.  The sharper form stated in
\Cref{thm:pac-bayes} (with $1/(2m)$ rather than $2/m$ under the root) requires
either Catoni's thermodynamic refinement~\cite{Catoni2007} or, for
finite~$\mathcal{H}$, the direct union bound over hypotheses that the
machine-checked companion follows; we do not reproduce either sharpening here.
\end{proof}


\begin{obstbox}
\textbf{The ``Bayesian'' in PAC-Bayes is an obstruction, not a
feature.}

The PAC-Bayes bound looks Bayesian: it involves a ``prior''~$\pi$
and a ``posterior''~$\rho$, and the complexity term $\KL(\rho\|\pi)$
penalizes departure from the prior.  But the prior plays a
fundamentally different role from a Bayesian prior.

In Bayesian inference, the prior $\pi$ represents \emph{beliefs}
about which hypothesis is true.  The posterior is obtained by Bayes'
rule: $\rho(h) \propto \pi(h) \cdot \prod_i p(z_i \mid h)$.  The
prior must be honest: it should reflect actual uncertainty.

In PAC-Bayes, the prior is a \emph{reference measure} chosen for
technical convenience.  It need not represent beliefs.  The ``posterior''
is \emph{any} distribution over~$\mathcal{H}$, not necessarily the
Bayesian posterior.  The bound holds for all~$\rho$ simultaneously,
so one can optimize~$\rho$ to minimize the bound.  The optimal
$\rho$ is the \emph{Gibbs posterior}
$\rho_\lambda \propto e^{-\lambda \emprisk_S(h)} \pi(h)$, which
resembles a Bayesian posterior but with a free temperature
parameter~$\lambda$ that a true Bayesian would set to~$m$.

\emph{Type mismatch:} Bayesian prior lives in the space of beliefs;
PAC-Bayes prior lives in the space of reference measures.  They have
the same mathematical type (distributions over~$\mathcal{H}$) but
different \emph{epistemic types}.  This is why PAC-Bayes bounds hold
for \emph{any} prior, including bad ones; the bound simply becomes
large.  A Bayesian analysis with a bad prior can give arbitrarily
wrong posteriors; a PAC-Bayes analysis with a bad prior gives a
correct but loose bound.
\end{obstbox}


% ================================================================
% SECTION 6: ALGORITHMIC STABILITY
% ================================================================

\section{Algorithmic Stability}
\label{sec:stability}

Stability bounds abandon the hypothesis class entirely and focus on the
\emph{algorithm}: if replacing a single training example changes the
output hypothesis only slightly, the algorithm generalizes.

\begin{defn}[Uniform Stability {\cite{BousquetElisseeff2002}}]
\label{def:uniform-stability}
A (possibly randomized) algorithm~$A$ is \emph{$\beta$-uniformly stable}
with respect to loss~$\ell$ if, for all samples
$S = (z_1, \ldots, z_m)$ and all samples~$S^{(i)}$ obtained by
replacing~$z_i$ with an independent copy~$z_i'$:
\[
  \sup_{z \in X \times Y}
    \bigl|\ell(A(S), z) - \ell(A(S^{(i)}), z)\bigr|
  \leq \beta.
\]
\end{defn}

\begin{thm}[Stability Generalization Bound
{\cite{BousquetElisseeff2002}}]
\label{thm:stability-bound}
If $A$ is $\beta$-uniformly stable and $\ell$ takes values in $[0,1]$,
then for any $\delta > 0$, with probability at least $1 - \delta$:
\[
  |\risk_D(A(S)) - \emprisk_S(A(S))|
  \leq 2\beta + (4m\beta + 1)\sqrt{\frac{\log(1/\delta)}{2m}}.
\]
In particular, if $\beta = O(1/m)$, then the generalization gap is
$O(1/\sqrt{m})$.
\end{thm}

\begin{proof}
Define $\Delta(S) = \risk_D(A(S)) - \emprisk_S(A(S))$.  The proof
proceeds in two stages.

\emph{Stage 1: Bound $\E[\Delta(S)]$.}
Write $\risk_D(A(S)) = \E_{z'}[\ell(A(S), z')]$ and
$\emprisk_S(A(S)) = \frac{1}{m}\sum_i \ell(A(S), z_i)$.  By symmetry of
the i.i.d.\ draw,
$\E[\ell(A(S), z_i)] = \E[\ell(A(S^{(i)}), z_i')]$
up to an error of~$\beta$ (by uniform stability applied to~$S$
vs.~$S^{(i)}$).  Summing over~$i$ and using
$\E[\ell(A(S^{(i)}), z_i')] = \E[\risk_D(A(S^{(i)}))] =
\E[\risk_D(A(S))]$ (by identical distribution), one obtains
$|\E[\Delta(S)]| \leq \beta$.

\emph{Stage 2: Concentration.}
Replacing~$z_i$ changes $\emprisk_S(A(S))$ by at most $1/m$ (direct
computation) and changes $\risk_D(A(S))$ by at most~$\beta$ (via
stability).  So $\Delta(S)$ has bounded differences
$c_i \leq 2\beta + 1/m$.  McDiarmid's inequality gives
\[
  \Pr[|\Delta(S) - \E[\Delta(S)]| > t]
  \leq 2\exp\!\left(\frac{-2t^2}{\sum_i c_i^2}\right)
  \leq 2\exp\!\left(\frac{-2t^2}{m(2\beta+1/m)^2}\right).
\]
Combining with $|\E[\Delta(S)]| \leq \beta$ gives the stated bound.
\end{proof}

\begin{example}[Regularized ERM is stable]
\label{ex:regularized-erm-stable}
Consider regularized empirical risk minimization:
$A(S) = \arg\min_{h \in \mathcal{H}} \emprisk_S(h) +
\lambda \|h\|^2$.  For $\lambda$-strongly convex regularizers and
$L$-Lipschitz losses, $\beta \leq L^2/(2\lambda m)$, giving
$\beta = O(1/m)$ and hence $O(1/\sqrt{m})$
generalization~\cite{BousquetElisseeff2002}.
\end{example}

\begin{remark}[Stability beyond binary classification]
\label{rmk:stability-beyond}
A central result in stability theory is due to Shalev-Shwartz, Shamir,
Srebro, and Sridharan~\cite{ShalevShwartzShamir2010}: beyond binary
classification with 0--1 loss, \emph{learnability is equivalent to the
existence of a stable empirical risk minimizer}.  In this more general
setting (for instance, certain convex stochastic optimization
problems); uniform convergence provably fails to characterize learnability,
yet stability succeeds.  We state this without proof; the construction
separating learnability from uniform convergence is the technical core of
the cited paper.
\end{remark}


% ================================================================
% SECTION 7: INFORMATION-THEORETIC BOUNDS
% ================================================================

\section{Information-Theoretic Bounds}
\label{sec:info-bound}

The most recent entry in the generalization-bound landscape connects
generalization to the mutual information between the training sample and
the learned hypothesis.

\begin{thm}[Information-Theoretic Generalization Bound
{\cite{XuRaginsky2017}}]
\label{thm:xu-raginsky}
Let $A$ be a (possibly randomized) learning algorithm, $\ell$ a loss
taking values in $[0,1]$, and $W = A(S)$ the output hypothesis.  Let
$I(S; W)$ denote the mutual information between~$S$ and~$W$.  Then
\[
  \bigl|\E[\risk_D(W) - \emprisk_S(W)]\bigr|
  \leq \sqrt{\frac{I(S; W)}{2m}}.
\]
\end{thm}

\begin{proof}[Proof sketch]
The argument uses the transportation lemma, a consequence of Pinsker's
inequality and the Donsker--Varadhan representation of KL divergence.

For each index~$i$, define
$g_i(w) = \E_Z[\ell(w, Z)] - \ell(w, Z_i)$, where the first expectation
is over a fresh test point $Z \sim D$.  Then
$\E[\mathrm{gen}(W,S)] = \frac{1}{m}\sum_{i=1}^m \E[g_i(W)]$.

For each~$i$, the function $g_i$ satisfies $|g_i| \leq 1$, so $g_i(W)$
is $1$-sub-Gaussian under the marginal distribution of~$W$ given all
other samples.  The transportation lemma gives
$|\E[g_i(W)]| \leq \sqrt{2 I(W; Z_i)}$.
Averaging over~$i$ and using $\sum_i I(W; Z_i) \leq I(W; S)$ yields the
result.  This last inequality is \emph{not} the data-processing inequality:
it follows from the chain rule
$I(W; S) = \sum_i I(W; Z_i \mid Z_{<i})$ together with
$I(W; Z_i \mid Z_{<i}) \geq I(W; Z_i)$, the latter holding because the $Z_i$
are independent (equivalently, $I(Z_i; Z_{<i} \mid W) \geq 0$).
\end{proof}

\begin{remark}[Conditional mutual information]
\label{rmk:cmi}
The Xu--Raginsky bound can be vacuous for deterministic algorithms,
since $I(S; W) = H(W)$ when $W = A(S)$ is a deterministic function
of~$S$.  Steinke and Zakynthinou (2020) introduced the \emph{conditional
mutual information} (CMI) framework, which replaces $I(S; W)$ with
$I(A(\tilde{S}_U); U \mid \tilde{S})$, where $\tilde{S}$ is a
``supersample'' of $2m$ points and $U \in \{0,1\}^m$ selects which half
to train on.  The CMI is always bounded by $m \log 2$ and gives
nontrivial bounds even for deterministic algorithms.  This framework
unifies PAC-Bayes, VC, and stability bounds within a single
information-theoretic language.
\end{remark}


% ================================================================
% SECTION 8: MARGIN THEORY
% Profile F (Open Frontier).
% ================================================================

\section{Margin Theory}
\label{sec:margin}

Margin-based bounds were the original motivation for support vector
machines and remain the primary tool for analyzing neural network
generalization.

\begin{defn}[Margin]
\label{def:margin}
For a classifier $f \colon X \to \R$ and a labeled example $(x, y)$
with $y \in \{-1, +1\}$, the \emph{margin} is
$\gamma(f, x, y) = y \cdot f(x)$.  The margin is positive when the
prediction is correct, and its magnitude measures the ``confidence'' of
the prediction.
\end{defn}

\begin{thm}[Margin-Based Generalization Bound]
\label{thm:margin-bound}
Let $\mathcal{F}$ be a class of functions $f \colon X \to \R$ with
$\|f\|_\infty \leq B$, and let $\gamma > 0$.  Then for any $\delta > 0$,
with probability at least $1 - \delta$, for all $f \in \mathcal{F}$:
\[
  \Pr_{(x,y) \sim D}[y \cdot f(x) \leq 0]
  \leq \frac{1}{m}\sum_{i=1}^m \ind[y_i f(x_i) \leq \gamma]
  + \frac{4}{\gamma}\,\Rad_S(\mathcal{F})
  + \sqrt{\frac{\log(2/\delta)}{2m}}.
\]
\end{thm}

\begin{proof}[Proof sketch]
Define the ramp loss $\ell_\gamma(t) = \min(1, \max(0, 1 - t/\gamma))$,
which is $1/\gamma$-Lipschitz and satisfies
$\ind[t \leq 0] \leq \ell_\gamma(t) \leq \ind[t \leq \gamma]$.  Then
\begin{align*}
  \Pr_D[y f(x) \leq 0]
  &\leq \E_D[\ell_\gamma(y f(x))] \\
  &\leq \frac{1}{m}\sum_i \ell_\gamma(y_i f(x_i))
    + 2\Rad_S(\ell_\gamma \circ \mathcal{F})
    + \sqrt{\frac{\log(2/\delta)}{2m}} \\
  &\leq \frac{1}{m}\sum_i \ind[y_i f(x_i) \leq \gamma]
    + \frac{2}{\gamma}\cdot 2\Rad_S(\mathcal{F})
    + \sqrt{\frac{\log(2/\delta)}{2m}},
\end{align*}
using the Rademacher contraction lemma
($\Rad_S(\phi \circ \mathcal{F}) \leq L \cdot \Rad_S(\mathcal{F})$
for $L$-Lipschitz~$\phi$) in the last step.
\end{proof}

\begin{example}[Halfspaces with margin]
\label{ex:halfspace-margin}
Let $\mathcal{F} = \{x \mapsto \langle w, x \rangle : \|w\|_2 \leq W\}$
with data satisfying $\|x_i\| \leq B$.  From the computational
illustration in \Cref{sec:rademacher},
$\Rad_S(\mathcal{F}) \leq WB/\sqrt{m}$.  The margin bound gives
\[
  \Pr_D[y \langle w, x \rangle \leq 0]
  \leq \hat{L}_\gamma(w) + \frac{4WB}{\gamma\sqrt{m}}
    + \sqrt{\frac{\log(2/\delta)}{2m}},
\]
where $\hat{L}_\gamma(w)$ is the fraction of training points with margin
$\leq \gamma$.  Observe: the bound depends on $WB/\gamma$ (the ratio of
norm to margin), not on the ambient dimension.  This is why SVMs with
large-margin solutions generalize well even in very high dimensions.
\end{example}

\begin{example}[Neural networks: the open frontier]
\label{ex:neural-margin}
For a depth-$L$ network $f_W(x) = W_L \sigma(W_{L-1} \sigma(\cdots
W_1 x \cdots))$ with ReLU activation~$\sigma$, Bartlett, Foster, and
Telgarsky (2017) showed that the Rademacher complexity satisfies
\[
  \Rad_S(\mathcal{F})
  \leq \frac{B_x \cdot (\sqrt{2\log(2d_{\max})})^L
    \cdot \prod_{i=1}^L \|W_i\|_\sigma}{\sqrt{m}},
\]
where $\|W_i\|_\sigma$ is the spectral norm of the $i$-th layer and
$d_{\max}$ is the maximum width.  Combined with the margin bound, this
gives a generalization bound scaling as the product of spectral norms
divided by the margin.

The problem: for practical networks, this product is large.  A ResNet-18
on CIFAR-10 has $\prod_i \|W_i\|_\sigma \approx 10^3$, making the bound
vacuous (larger than~1).  Yet the network generalizes perfectly well.  The
resolution of this tension (finding bounds that are both non-vacuous and
predictive) remains one of the central open problems in learning theory.
\end{example}

\begin{openprob}[Tight deep network bounds]
\label{op:deep-margins}
For deep neural networks, all known generalization bounds (margin-based,
PAC-Bayes, stability, and information-theoretic) are \emph{loose} in
the following sense: the bounds scale with quantities that grow with
network size, yet empirically, larger networks generalize \emph{better}.

The specific open problem: find a generalization bound for deep networks
trained by stochastic gradient descent that (a)~is non-vacuous (smaller
than~1) on standard benchmarks, (b)~correctly predicts that increasing
width improves generalization, and (c)~does not require post-hoc
compression or quantization.  As of 2025, no known framework achieves all
three simultaneously.
\end{openprob}


% ================================================================
% SECTION 9: META-LEARNING BOUNDS
% Profile A-plus.  Brief.
% ================================================================

\section{Meta-Learning Bounds}
\label{sec:meta-pac}

\begin{thm}[Meta-PAC Bound {\cite{Baxter2000}}]
\label{thm:meta-pac}
In the meta-learning setting, the learner observes $T$ tasks, each
providing $m$ samples from a task-specific distribution~$D_t$.  The tasks
are drawn from a task distribution~$\mathcal{T}$.  Baxter's meta-PAC
bound states that if the bias class~$\mathcal{B}$ has finite complexity,
then with $T$ tasks and $m$ samples per task:
\[
  \risk_\mathcal{T}(\hat{b}) \leq \emprisk_{T,m}(\hat{b})
  + O\!\left(\sqrt{\frac{\log \mathcal{N}(\varepsilon,\mathcal{B},d)}{T}}
    + \sqrt{\frac{d_{\mathrm{task}}}{m}}\right),
\]
where $d_{\mathrm{task}}$ is the VC dimension of the task-level
hypothesis class induced by the bias.  The first term controls
meta-generalization; the second controls within-task generalization.  We
state the bound without proof; the argument combines a covering-number bound
at the meta level with a VC bound within each task~\cite{Baxter2000}.
\end{thm}


% ================================================================
% SECTION 10: THE COMPARISON
% CENTERPIECE: Profile G.  Five frameworks in parallax.
% ================================================================

\section{Five Frameworks Compared}
\label{sec:comparison}

The five frameworks bound the same quantity but use different information
and succeed in different regimes.  \Cref{tab:five-frameworks} compares
them along five axes.

\begin{table}[ht]
\centering\small
\caption{Five generalization-bound frameworks compared.  Each row
describes what the framework measures, what it bounds, where it is
tightest, and where it fails.}
\label{tab:five-frameworks}
\renewcommand{\arraystretch}{1.4}
\begin{tabularx}{\textwidth}{>{\bfseries\small}l X X X X}
\toprule
\textbf{Framework} & \textbf{Key quantity} & \textbf{What it bounds}
  & \textbf{Tight for} & \textbf{Fails on} \\
\midrule
Rademacher
  & $\Rad_S(\mathcal{F})$: noise correlation
  & Uniform over all $f \in \mathcal{F}$
  & Finite VC classes; linear models
  & Overparameterized models (vacuous) \\
\addlinespace
PAC-Bayes
  & $\KL(\rho \| \pi)$: posterior--prior divergence
  & Expected risk under posterior~$\rho$
  & Bayesian and ensemble methods
  & Requires meaningful prior \\
\addlinespace
Stability
  & $\beta$: one-sample sensitivity
  & $|\risk - \emprisk|$ for algorithm~$A$
  & Regularized ERM; SGD
  & Non-regularized ERM (unstable) \\
\addlinespace
Information
  & $I(S; W)$: mutual information
  & $|\E[\risk - \emprisk]|$
  & Noisy algorithms (DP-SGD)
  & Deterministic algorithms (vacuous) \\
\addlinespace
Margin
  & $\gamma$: prediction confidence
  & $\Pr[\text{error}]$ via margin loss
  & SVMs; kernel methods
  & Deep networks (norms too large) \\
\bottomrule
\end{tabularx}
\end{table}


\subsection{The Parallax Diagram}

\Cref{fig:five-frameworks} displays the implication and obstruction
relationships among the five frameworks.  The diagram should be read as a
\emph{parallax}: five views of the same phenomenon (generalization),
each projecting the three-way interaction between data, class, and
algorithm onto a different axis.

\begin{figure}[ht]
\centering
\begin{tikzpicture}[
    vfbox/.style={draw=defblue, line width=1pt, rounded corners=4pt,
      minimum width=2.5cm, minimum height=0.9cm, font=\small\bfseries,
      align=center, fill=defblue!10},
    litbox/.style={draw=black!45, thin, dashed, rounded corners=4pt,
      minimum width=2.5cm, minimum height=0.9cm, font=\small, align=center,
      fill=black!3},
    vedge/.style={-{Stealth[length=2.5mm]}, line width=1pt, defblue},
    ledge/.style={-{Stealth[length=2.5mm]}, semithick, black!55, densely dashed},
    cedge/.style={{Stealth[length=2.2mm]}-{Stealth[length=2.2mm]}, thick,
      black!45, dotted},
    elbl/.style={font=\scriptsize, midway, fill=white, inner sep=1.5pt},
]
% --- frameworks: verified-bound (thick blue border) vs literature-only (thin dashed) ---
\node[vfbox] (vc) at (-4.2,2.3) {VC / Uniform\\Convergence};
\node[vfbox] (rad) at (0,2.3) {Rademacher};
\node[litbox] (margin) at (4.2,2.3) {Margin};
\node[litbox] (info) at (0,0) {Information-\\Theoretic};
\node[litbox] (stab) at (-2.7,-2.5) {Stability};
\node[vfbox] (pb) at (2.7,-2.5) {PAC-Bayes};

% --- the one kernel-verified bridge: Rademacher controlled by VC via Sauer-Shelah ---
\draw[vedge] (vc) -- node[elbl] {Sauer--Shelah} (rad);

% --- literature implications / analogies ---
\draw[ledge] (margin) -- node[elbl] {specializes} (rad);
\draw[ledge] (stab) -- node[elbl] {implies} (info);
\draw[ledge, {Stealth[length=2.5mm]}-{Stealth[length=2.5mm]}]
  (info) -- node[elbl] {dual} (pb);

% --- conceptual orthogonalities (argued, not proved) ---
\draw[cedge] (rad) -- node[elbl] {orthogonal} (info);
\draw[cedge] (stab) -- node[elbl] {type mismatch} (pb);

% --- evidence-grade legend ---
\begin{scope}[shift={(-4.4,-4.0)}]
  \draw[vedge] (0,0) -- (0.7,0);
  \node[font=\scriptsize, anchor=west] at (0.82,0) {kernel-verified bound / bridge};
  \draw[ledge] (5.0,0) -- (5.7,0);
  \node[font=\scriptsize, anchor=west] at (5.82,0) {literature (cited)};
  \draw[cedge] (0,-0.52) -- (0.7,-0.52);
  \node[font=\scriptsize, anchor=west] at (0.82,-0.52) {conceptual orthogonality (argued)};
  \node[vfbox, minimum width=0.55cm, minimum height=0.3cm] at (5.27,-0.52) {};
  \node[font=\scriptsize, anchor=west] at (5.62,-0.52) {verified-bound};
  \node[litbox, minimum width=0.55cm, minimum height=0.3cm] at (8.5,-0.52) {};
  \node[font=\scriptsize, anchor=west] at (8.85,-0.52) {literature-only};
\end{scope}
\end{tikzpicture}
\caption[Six generalization frameworks graded by evidence: one risk gap, many certificates.]{The
six frameworks all bound the \emph{same} quantity, the gap between true and empirical risk, but
each reads a different slice of the data--class--algorithm interaction.  The kernel machine-checks
the bounds for the class-based frameworks: VC uniform convergence (\leanname{finite_massart_lemma}),
Rademacher complexity (\leanname{EmpiricalRademacherComplexity}), and PAC-Bayes
(\leanname{pac_bayes_finite}); these three nodes are drawn solid.  The one machine-checked
cross-framework bridge is that Rademacher complexity is controlled by VC dimension through
Sauer--Shelah (\leanname{vcdim_bounds_rademacher_quantitative}), the single solid edge.  The
remaining links are literature, proved on paper but not verified (margin specializes Rademacher,
stability implies an information bound, PAC-Bayes and the information bound are dual views), or
conceptual, argued rather than proved (information and Rademacher are orthogonal because one reads
the algorithm and the other the class; stability and PAC-Bayes are a type mismatch between
algorithm space and distribution space).  Solid marks a kernel-verified bound or bridge, a lighter
dashed style a cited literature result, and a dotted style a conceptual orthogonality that is
argued rather than proved.}
\label{fig:five-frameworks}
\end{figure}


\subsection{Cross-Framework Relationships}

The five frameworks are not independent.  Several implication, analogy,
and obstruction relationships connect them:

\begin{enumerate}[leftmargin=*, itemsep=3pt]
\item \textbf{Rademacher $\Rightarrow$ VC.}  The Rademacher complexity of
  a binary class satisfies $\mathcal{R}_m(\mathcal{H}) \leq
  \sqrt{2\VC(\mathcal{H}) \log(em/\VC(\mathcal{H}))/m}$
  (\Cref{prop:rad-vc}), so Rademacher bounds subsume and refine VC
  bounds.  The refinement is strict: Rademacher complexity adapts to the
  sample, whereas VC dimension does not.

\item \textbf{Margin $\Rightarrow$ Rademacher.}  The margin bound
  (\Cref{thm:margin-bound}) is derived \emph{from} Rademacher complexity
  applied to the ramp loss, so margin bounds are a specialization.

\item \textbf{Stability $\Rightarrow$ Information.}  A $\beta$-uniformly
  stable randomized algorithm satisfies $I(S; W) = O(m \beta^2)$.  The
  operative step is a stability-to-divergence bound (replacing one example
  perturbs the output distribution by $O(\beta)$ in KL divergence), not the
  data-processing inequality; summing the per-coordinate contributions gives
  the $m\beta^2$ scaling.  Thus, for noisy algorithms, stability bounds imply
  information-theoretic bounds.

\item \textbf{PAC-Bayes $\leftrightarrow$ Information.}  The PAC-Bayes
  bound with a data-independent prior can be interpreted as an
  information-theoretic bound: the KL divergence
  $\KL(\rho \| \pi)$ upper-bounds a specific mutual information
  term.  Conversely, the information-theoretic framework recovers
  PAC-Bayes via a change-of-measure argument.  The two are essentially
  dual views~\cite{HellstromDurisi2020}.

\item \textbf{Obstruction: Information $\not\Rightarrow$ Rademacher.}
  Information-theoretic bounds depend on the \emph{algorithm}; Rademacher
  bounds depend on the \emph{class}.  A deterministic ERM algorithm can have
  $I(S; W)$ as large as $H(W)$ (the full entropy of its output) while
  operating on a class with small Rademacher complexity.  The two frameworks
  capture orthogonal aspects of generalization.

\item \textbf{Obstruction: Stability $\not\Rightarrow$ PAC-Bayes.}
  Stability measures perturbation sensitivity of the \emph{algorithm}.
  PAC-Bayes measures divergence of the \emph{posterior} from the
  \emph{prior}.  An algorithm defines a posterior (connecting them), but
  the primary objects are different types: stability lives in the space of
  algorithms, PAC-Bayes in the space of distributions.
\end{enumerate}


\subsection{What Each Framework Sees}

\begin{figure}[ht]
\centering
\begin{tikzpicture}[
    vfbox/.style={draw=defblue, line width=1pt, rounded corners=3pt,
      minimum width=1.7cm, minimum height=0.58cm, font=\small\bfseries,
      align=center, fill=defblue!10, inner sep=2pt},
    litbox/.style={draw=black!45, thin, dashed, rounded corners=3pt,
      minimum width=1.7cm, minimum height=0.58cm, font=\small, align=center,
      fill=black!3, inner sep=2pt},
    cdot/.style={circle, fill, inner sep=2.4pt},
    clbl/.style={font=\footnotesize\bfseries},
    note/.style={font=\scriptsize, text=black!55},
]
% --- the three information sources (corners of the triangle) ---
\coordinate (D) at (0,0);
\coordinate (H) at (7,0);
\coordinate (A) at (3.5,5);
\draw[thick, fill=layergray] (D) -- (H) -- (A) -- cycle;
\node[cdot, defblue] at (D) {};
\node[cdot, thmred] at (H) {};
\node[cdot, sepgreen] at (A) {};
\node[clbl, defblue, below left=0pt and -3pt] at (D) {Data $D$};
\node[clbl, thmred, below right=0pt and -3pt] at (H) {Class $\mathcal{H}$};
\node[clbl, sepgreen, above=2pt] at (A) {Algorithm $A$};
% --- each framework placed at the source it reads; solid = verified bound, dashed = literature ---
% interior channel: Information reads the data-to-output channel I(S;W)
\node[litbox] (info) at (3.5,1.65) {Information};
\node[note, below=1pt of info] {$I(S;W)$ channel};
% class corner (right): VC and Rademacher (verified), Margin (literature, class geometry)
\node[vfbox] (vc) at (7.8,0.55) {VC};
\node[vfbox] (rad) at (7.8,1.4) {Rademacher};
\node[litbox] (margin) at (7.8,2.25) {Margin};
% algorithm-class edge (upper right): PAC-Bayes (verified, posterior over the class)
\node[vfbox] (pb) at (5.5,3.4) {PAC-Bayes};
% algorithm-data edge (upper left): Stability (literature)
\node[litbox] (stab) at (0.5,3.3) {Stability};
% --- evidence-grade legend (same convention as Fig 12.1) ---
\begin{scope}[shift={(-0.3,-1.35)}]
  \node[vfbox, minimum width=0.85cm, minimum height=0.3cm] at (0,0) {};
  \node[font=\scriptsize, anchor=west] at (0.55,0) {verified-bound};
  \node[litbox, minimum width=0.85cm, minimum height=0.3cm] at (3.6,0) {};
  \node[font=\scriptsize, anchor=west] at (4.15,0) {literature-only};
\end{scope}
\end{tikzpicture}
\caption[Generalization is read from three sources: data, class, or algorithm.]{Every
generalization bound reads one of three sources of information, the data, the hypothesis
class, or the algorithm, drawn here as the three corners of a triangle, so the choice of
framework is the choice of which source one can measure.  The class-based bounds sit at the
class corner and are machine-checked: VC uniform convergence (\leanname{finite_massart_lemma})
and Rademacher complexity (\leanname{EmpiricalRademacherComplexity}), drawn solid.  PAC-Bayes
(\leanname{pac_bayes_finite}) reads the posterior over the class induced by the algorithm and
is also machine-checked, so it sits on the algorithm--class edge and is drawn solid as well.
Stability reads the algorithm--data edge, information reads the data-to-output channel in the
interior, and margin reads the class geometry; these three are proved in the literature but
not verified, and are drawn lighter and dashed.  The obstructions of the companion
\Cref{fig:five-frameworks} are exactly these projections landing on different corners:
information and Rademacher read opposite sources, and stability and PAC-Bayes read different
edges, so the two figures are spatial complements of one another.}
\label{fig:triangle}
\end{figure}


\subsection{Which Framework When?}

No single framework dominates.  The choice depends on what is known:

\begin{itemize}[leftmargin=*, itemsep=3pt]
\item \textbf{If you know the class but not the algorithm}: use
  Rademacher or margin bounds.
\item \textbf{If you know the algorithm but not the class}: use
  stability or information-theoretic bounds.
\item \textbf{If you can choose a prior}: use PAC-Bayes.
\item \textbf{If you want algorithm-specific bounds for SGD}: stability
  and information-theoretic bounds apply; Rademacher bounds do not.
\item \textbf{If you want non-vacuous bounds for deep networks}:
  PAC-Bayes with learned priors is currently the most successful
  approach, but the bounds remain far from tight
  (\Cref{op:deep-margins}).
\end{itemize}


\begin{example}[Three bounds on the same problem]
\label{ex:three-bounds}
Consider linear classification with $\|w\| \leq 1$ on data with
$\|x\| \leq 1$ in $\R^d$, with $m$ training points and 0--1 loss.

\textbf{VC bound.}  $\VC = d+1$, giving generalization gap
$O(\sqrt{d \log m / m})$.  This scales linearly with dimension.

\textbf{Rademacher bound.}  $\Rad_S \leq 1/\sqrt{m}$
(from the computation in \Cref{sec:rademacher}), giving generalization
gap $O(1/\sqrt{m})$.  No dimension dependence.

\textbf{Margin bound.}  If the margin is $\gamma$, the bound is
$\hat{L}_\gamma + O(1/(\gamma\sqrt{m}))$.  If the data is
well-separated ($\gamma = 0.1$, $\hat{L}_\gamma = 0$), this gives
$O(10/\sqrt{m})$, a worse constant, but the zero margin-loss term
compensates.

The three bounds are complementary.  In dimension $d = 1000$ with
$m = 5000$, the VC bound gives approximately $0.6$ (nearly vacuous), the
Rademacher bound gives approximately $0.04$ (useful), and the margin
bound with $\gamma = 0.1$ gives approximately $0.14$ plus the margin
loss.
\end{example}


% ================================================================
% SECTION 11: CONCRETE COMPARISON ON NEURAL NETWORKS
% ================================================================

\section{A Case Study: Neural Networks}
\label{sec:nn-comparison}

The five frameworks yield strikingly different results on neural
networks, making them the ideal testbed for understanding what each
framework captures.

\begin{computational}
\textbf{Two-layer ReLU network on $\R^{100}$.}
Consider $f_W(x) = v^\top \sigma(Wx)$ with $W \in \R^{k \times 100}$,
$v \in \R^k$, width $k = 1000$, and ReLU activation $\sigma$.

\smallskip
\begin{tabular}{@{}lll@{}}
\toprule
\textbf{Framework} & \textbf{Bound expression} & \textbf{With $m=10^4$} \\
\midrule
VC & $O(\sqrt{k \cdot 100 \cdot \log(m) / m})$ & $\approx 3.4$ (vacuous) \\
Rademacher & $O(\|v\|_1 \|W\|_\sigma B_x / \sqrt{m})$
  & $\approx 0.3$ (if norms bounded) \\
PAC-Bayes & $O(\sqrt{\KL(\rho\|\pi)/m})$
  & $\approx 0.05$ (with good prior) \\
Stability & $O(L^2/(\lambda m))$ for reg.\ ERM
  & $\approx 0.1$ (if regularized) \\
Margin & $O(\|v\|_1 \|W\|_\sigma / (\gamma\sqrt{m}))$
  & depends on $\gamma$ \\
\bottomrule
\end{tabular}

\smallskip
The VC bound is vacuous because it scales with the parameter count
($\sim 10^5$, via the proxy $\VC = \tilde{O}(\#\text{params})$).
Rademacher and margin bounds depend on norms, which can be controlled by
regularization.  PAC-Bayes gives the tightest bound when a meaningful
prior is available (e.g., the initialization distribution of SGD).
Stability bounds require regularization or early stopping to
ensure~$\beta = O(1/m)$.  The information-theoretic bound depends on the
noise level of SGD: more noise means less $I(S;W)$ and a tighter bound,
but also worse training performance.
\end{computational}


\subsection{A Certified Bound for Attention}
\label{sec:certified-attention}

The five frameworks are theorems about hypothesis classes and algorithms in
the abstract.  For one concrete modern architecture, a softmax-attention
head, the pipeline of \Cref{sec:uniform-convergence,sec:rademacher} has been
carried out inside a proof assistant, producing a generalization guarantee
that is machine-checked down to the floating-point execution of the head.

\begin{prop}[Certified generalization for an attention head]
\label{prop:certified-attention}
Fix a context length~$n$, embedding dimension~$d$, and a bounded,
$L_\ell$-Lipschitz loss~$\ell$ with $|\ell| \leq b$.  Consider a single
attention head $f_\theta(x) = \mathrm{attn}\bigl(W_{\mathrm{dec}}(\theta);
\mathrm{prep}(x)\bigr)$ with a continuous, $L_W$-Lipschitz parameter map
$W_{\mathrm{dec}}$, over parameters $\theta$ in a Euclidean ball of
radius~$R$ in $\R^p$.  Set $\kappa = 4R\,L_\ell\,(dB)\,L_W$.  Then for every
$\varepsilon \geq 0$ and sample size~$m$, with probability at least
$1 - \exp\!\bigl(-\varepsilon^2 m / (2b^2)\bigr)$ over $S \sim P^m$,
\[
  \E_{x \sim P}[\ell(f_\theta(x))]
  \;\leq\; \frac{1}{m}\sum_{i=1}^m \ell(f_\theta(S_i))
    \;+\; \frac{24\sqrt{2}}{\sqrt{m}}
      \int_0^{2b}\!\Bigl(\sqrt{\log 2}
        + \sqrt{p\,\kappa}\;\eta^{-1/2}\Bigr)\mathrm{d}\eta
    \;+\; 2\varepsilon
    \;+\; 2 L_\ell\,\mathrm{env}(f_\theta),
\]
where $\mathrm{env}(f_\theta)$ is the floating-point execution envelope of
the head.
\end{prop}

\begin{proof}[Proof idea]
The bound instantiates the chapter's machinery on the attention head viewed
as a real-valued function class indexed by~$\theta$.  The middle term is the
\emph{Dudley entropy integral} (\Cref{ch:compression}) of the
parameter-perturbation cover of the head: the Lipschitz constant $\kappa$
controls the covering numbers, and the integral $\int_0^{2b}\sqrt{\log
\mathcal{N}(\eta)}\,\mathrm{d}\eta$ takes the displayed closed form because a
radius-$R$ ball in $\R^p$ has metric entropy $O(p\log(1/\eta))$.  The
high-probability lift is the bounded-differences (McDiarmid) step of
\Cref{thm:rademacher-bound}, with per-example increment $2b/m$.  The one new
ingredient is the rounding term: the verified statement bounds the
\emph{executed}, finite-precision head, so it carries an explicit envelope
$\mathrm{env}(f_\theta)$ separating the idealized map from its floating-point
realization.  The full derivation is
formalized.\formal{TLT.attnHead_certified_generalization}
\end{proof}

The construction lifts from one head to a depth-graded multi-head encoder
stack: the budget becomes the Dudley capacity of the depth-replicated block,
growing through the per-layer perturbation envelope with depth and head
count, and the resulting predicate (true risk exceeds empirical risk plus
budget plus rounding only on a $\delta$-event) is discharged for the stack
indexed by its architectural
configuration.\formal{TLT.config_capacity_law}  That predicate is itself the
machine-checked definition of a generalization
certificate.\formal{TLT.CertifiedGeneralization}

\begin{remark}[What the certificate settles, and what it does not]
\label{rmk:certified-attention}
\Cref{prop:certified-attention} is an end-to-end verified generalization
bound for a transformer component: the abstract theory of this chapter,
instantiated and checked for attention.  It does \emph{not} resolve
\Cref{op:deep-margins}.  The Dudley budget still grows with the parameter
radius~$R$, the depth, and the head count, so on a large pretrained model the
certificate is correct but loose, exactly as the classical bounds are.
What it adds is a different axis of progress: a \emph{trustworthy} bound,
whose constants and finite-precision corrections have been audited by a
kernel rather than estimated by hand.
\end{remark}


% ================================================================
% OPEN PROBLEMS
% ================================================================

\section{Open Problems}
\label{sec:generalization-open}

The deep-network gap of \Cref{op:deep-margins} is the headline question, but
the chapter's five frameworks and their machine-checked fragments open a
wider frontier.  Each problem below is tagged as a \emph{known unknown} (KU),
sharply posed but unresolved, or an \emph{unknown known} (UK), where the
right invariant or formal object is still missing; a design-lab or literature
anchor is named where one exists.

\begin{openprob}[Formalizing the data-dependent Rademacher bound]
\label{op:formalize-rademacher}
The empirical Rademacher complexity, the Massart finite lemma, and the VC
bound $\mathcal{R}_m \leq \sqrt{2d\log(em/d)/m}$ are machine-checked
(\Cref{def:rademacher,prop:rad-vc}), but the assembled high-probability bound
of \Cref{thm:rademacher-bound} (symmetrization plus the two McDiarmid
applications) is so far only a placeholder in the companion development.  Can
the full data-dependent bound be formalized end to end?  A known unknown:
classical on paper~\cite{BartlettMendelson2002}, open as a verified artifact.
\end{openprob}

\begin{openprob}[A non-vacuous certified bound for attention]
\label{op:tighten-certified-attention}
The certificate of \Cref{prop:certified-attention} is correct but its Dudley
budget grows with the parameter radius, depth, and head count
(\leanname{config_capacity_law}).  Can one machine-check a bound for a
realistic attention stack that is also \emph{non-vacuous} on a benchmark: a
verified analogue of the PAC-Bayes results for deep networks?  A known
unknown, and the verified counterpart of \Cref{op:deep-margins}.
\end{openprob}

\begin{openprob}[Discharging the expressivity leg]
\label{op:expressivity-leg}
The transformer design law assembles a certified capacity bound but states an
\emph{expressivity} obligation as an open hypothesis: that the stack can
realize a target risk at all.  Can the expressivity leg be discharged with
explicit width/depth, turning the design law into an unconditional
generalization-and-approximation guarantee?  An unknown known: the
approximation invariant for the certified stack is not yet
formalized~\cite{XuRaginsky2017}.
\end{openprob}

\begin{openprob}[Making the CMI unification precise]
\label{op:cmi-unification}
The conditional-mutual-information framework (\Cref{rmk:cmi}) is said to
unify PAC-Bayes, VC, and stability bounds.  Is there a single CMI inequality
from which all three are recovered as corollaries, with explicit constants?
An unknown known: the unification is asserted but no master theorem with the
three recoveries written out exists~\cite{SteinkeZakynthinou2020}.
\end{openprob}

\begin{openprob}[When is uniform convergence necessary?]
\label{op:uniform-necessity}
Shalev-Shwartz et al.\ show that beyond binary classification, learnability
can hold without uniform convergence (\Cref{rmk:stability-beyond}).  Exactly
which loss/hypothesis structures force uniform convergence to be necessary,
not merely sufficient?  A known unknown: the boundary between the two regimes
is uncharted~\cite{ShalevShwartzShamir2010}.
\end{openprob}

\begin{openprob}[Generalization of interpolators]
\label{op:benign-overfitting}
Modern networks interpolate the training data ($\emprisk_S = 0$) yet
generalize, contradicting the intuition behind every bound in this chapter.
Is there a complexity measure that is small precisely for the interpolators
that generalize (benign overfitting) and large for those that do not (the
double-descent peak)?  A known unknown driving current theory.
\end{openprob}

\begin{openprob}[Optimal data-dependent priors]
\label{op:pac-bayes-prior}
PAC-Bayes is tightest with a prior close to the learned posterior, but the
prior must be data-independent.  Is there a principled, optimal way to build
a data-dependent prior (via data splitting or differential privacy) that
provably minimizes the bound?  A known unknown opened by the non-vacuous
deep-network results~\cite{DziugaiteRoy2017}.
\end{openprob}

\begin{openprob}[A formal margin bound for attention]
\label{op:margin-attention}
The spectral-norm margin bound for deep networks (\Cref{ex:neural-margin})
has no verified counterpart for attention.  Can one machine-check a
margin-based bound for a softmax-attention class, with the product of
query/key/value operator norms in place of the layer spectral
norms~\cite{BartlettFosterTelgarsky2017}?  A known unknown on the transformer
side, complementary to \Cref{prop:certified-attention}.
\end{openprob}

\begin{openprob}[Rademacher lower bounds]
\label{op:rademacher-lower}
Every bound here is an upper bound on the generalization gap.  For which
classes is the Rademacher complexity also a \emph{lower} bound on the
worst-case gap, so that $\mathcal{R}_m$ exactly captures learnability rather
than merely bounding it?  An unknown known: a matching lower-bound invariant
is missing for general classes.
\end{openprob}

\begin{openprob}[Non-vacuous information bounds for deterministic learners]
\label{op:info-deterministic}
The Xu--Raginsky bound is vacuous when $W = A(S)$ is deterministic, since
$I(S;W) = H(W)$ (\Cref{rmk:cmi}).  Apart from the CMI construction, is there a
generalization bound for deterministic algorithms phrased in genuinely
information-theoretic quantities that stays finite?  A known unknown at the
frontier of the information-theoretic tradition.
\end{openprob}


% ================================================================
% SECTION 12: HISTORICAL CONTEXT
% ================================================================

\begin{historical}
The study of generalization bounds spans five decades and five research
communities.

The VC theory (Vapnik--Chervonenkis, 1971) established the first uniform
convergence bounds.  The connection to PAC learning was made explicit by
Blumer et al.\ (1989).  Rademacher complexity was developed independently
by Koltchinskii (2001) and Bartlett--Mendelson (2002) as a data-dependent
refinement.

PAC-Bayes bounds originated with Shawe-Taylor and Williamson (1997) and
were formalized by McAllester (1998, 1999).  Catoni (2007) developed the
optimal ``thermodynamic'' form.  The framework gained renewed interest
around 2017 when Dziugaite and Roy showed that PAC-Bayes bounds can be
made non-vacuous for deep networks.

Algorithmic stability was introduced by Devroye and Wagner (1979) in a
simpler form, and the modern uniform stability framework was established
by Bousquet and Elisseeff (2002).  The connection to learnability beyond
binary classification was established by Shalev-Shwartz et al.\ (2010).

Information-theoretic bounds are the most recent tradition, beginning
with Russo and Zou (2016) and Xu and Raginsky (2017).  The conditional
mutual information framework of Steinke and Zakynthinou (2020) provided
the unifying perspective.

Margin theory has the longest applied history (Vapnik, 1995; Bartlett,
1998) and was reinvigorated by Bartlett, Foster, and Telgarsky (2017)
for deep networks.

The unification of these traditions (showing they are all projections
of a single underlying phenomenon) remains an active research
direction.
\end{historical}


% ================================================================
% NOTES
% ================================================================

\bigskip

\begin{remark}[The chapter in context]
This chapter completes Part~III's treatment of quantitative measures of
learning complexity.  \Cref{ch:dimensions} developed the combinatorial
dimensions (VC, Littlestone, DS).  \Cref{ch:compression} connected
learnability to sample compression.  The present chapter adds the
\emph{analytic} layer: bounds that read the algorithm, the data, and the
geometry of the learned hypothesis on top of the class itself.
The five frameworks exemplify a recurring theme of this book: the same
mathematical phenomenon (in this case, the gap between training and
test performance) looks different depending on which structural
information one chooses to measure.  No single framework captures the
full picture, and the obstructions between frameworks (\Cref{fig:five-frameworks})
are as informative as the implications.
\end{remark}

\subsection*{Counterfactual exercises: generalizing Figure~\ref{fig:five-frameworks}}

The parallax of \Cref{fig:five-frameworks} is unusually lopsided as figures in this book go: most of its content is a frontier.  Of six framework nodes only three are drawn solid (VC uniform convergence \leanname{finite_massart_lemma}, Rademacher complexity \leanname{EmpiricalRademacherComplexity}, and PAC-Bayes \leanname{pac_bayes_finite}, all class-based); the other three (stability, information, margin) are literature.  Of six cross-framework links exactly one is solid, the Sauer--Shelah bridge by which VC dimension controls Rademacher complexity \leanname{vcdim_bounds_rademacher_quantitative}; the rest are either literature implications (margin specializing Rademacher, stability implying an information bound, PAC-Bayes and information being dual) or merely-argued conceptual edges (the algorithm-versus-class orthogonality and the stability-versus-PAC-Bayes type mismatch).  So when each exercise below perturbs the picture, ask which of three things the perturbation produces: an \emph{instance} that the one verified core or a single graded edge already exhibits, a \emph{boundary} where toggling one node sort, edge grade, or information axis flips the evidence grade or the direction, or a \emph{frontier} where a literature or conceptual edge still waits to be promoted to verified.  Most of these land on the frontier, which is the honest shape of the subject.

\begin{enumerate}[leftmargin=*,itemsep=3pt]
\item \textbf{The three solid nodes are exactly the class-based bounds the kernel checks.} The figure's three solid borders are a claim about evidence, and this is the defining verified instance of the whole diagram: VC uniform convergence, Rademacher complexity, and PAC-Bayes are the frameworks whose bounds are machine-checked (\cref{thm:vc-uniform}, \leanname{finite_massart_lemma}; \cref{def:rademacher}, \leanname{EmpiricalRademacherComplexity}; \cref{thm:pac-bayes}, \leanname{pac_bayes_finite}), so a solid border asserts exactly \emph{that the bound is a kernel theorem}. The three lighter-bordered nodes, stability, information, and margin, are drawn weaker precisely because their bounds (\cref{thm:stability-bound}, \cref{thm:xu-raginsky}, \cref{thm:margin-bound}) are literature with no verified handle.

\item \textbf{The one solid edge is the only machine-checked cross-framework bridge.} Take the single solid arrow between Rademacher and VC on its own. It is the lone settled link in the diagram: Rademacher complexity is bounded by VC dimension through Sauer--Shelah, $\mathcal{R}_m(\mathcal{H}) \le \sqrt{2\VC(\mathcal{H})\log(em/\VC(\mathcal{H}))/m}$ (\cref{prop:rad-vc}, \leanname{vcdim_bounds_rademacher_quantitative}), the one cross-framework relationship the kernel certifies. The old figure's separate ``refines'' arrow and dotted ``Sauer--Shelah'' arrow named the \emph{same} fact, which is why the re-graded figure collapses them into one solid edge and leaves nothing else solid between frameworks. This is the figure's instance.

\item \textbf{The parallax is one risk-gap quantity seen from six sides.} Read the whole figure as six projections of one object rather than six theories, and it becomes the organizing instance: every framework bounds the single quantity $|\risk - \emprisk|$ (or its expectation), and the rows of \cref{tab:five-frameworks} differ only in the \emph{key quantity} each reads (noise correlation, posterior--prior divergence, one-sample sensitivity, mutual information, margin). The diagram is the abstract witness that generalization is one quantity admitting many certificates, three of them, the class-based ones, machine-checked, and the rest borrowed from the literature.

\item \textbf{PAC-Bayes is the verified bound that lives in distribution space.} The PAC-Bayes node anchors the posterior axis: the McAllester finite bound on the Gibbs risk is machine-checked (\cref{thm:pac-bayes}, \leanname{pac_bayes_finite}, with \leanname{gibbsError} the verified mean error of the posterior-sampled predictor), so PAC-Bayes joins VC and Rademacher among the solid nodes even though it reads the \emph{posterior} rather than the class. That same kernel handle is what keeps the conceptual stability-to-PAC-Bayes edge a type relation and not a theorem, since stability reads the algorithm and PAC-Bayes the distribution it induces. This is the instance that fixes the posterior reading.

\item \textbf{The two obstructions are orthogonal information sources, not separations.} Toggle the figure's reading of the conceptual obstructions from ``proved gap'' to ``different information axis,'' and the figure's principal boundary appears: information-theoretic bounds depend on the \emph{algorithm} and Rademacher bounds on the \emph{class}, so a deterministic ERM can carry $I(S;W)=H(W)$ on a class of small Rademacher complexity (\cref{sec:comparison}), and stability lives in algorithm space while PAC-Bayes lives in distribution space. Neither obstruction is a separation theorem nor even a formal literature separation, so both are drawn in the conceptual style, argued not proved, and reversing either direction does not turn it into a verified edge.

\item \textbf{Margin is a Rademacher specialization, and that specialization is literature.} The margin-to-Rademacher edge is the boundary that fixes the edge's grade: the margin bound is \emph{derived from} Rademacher complexity applied to the ramp loss (\cref{thm:margin-bound}), so the arrow points from margin into Rademacher as a specialization, yet because thm:margin-bound is not machine-checked the edge is drawn lighter than the solid rad--vc bridge. Promoting it would require a verified ramp-loss contraction, exactly the kind of step still missing for the assembled bounds.

\item \textbf{The stability-to-information edge is one direction only, and unverified.} Follow the stability-to-information arrow and ask whether the reverse is drawn; it is not, and that asymmetry is the boundary between a literature implication and an absent one. A $\beta$-uniformly stable randomized algorithm satisfies $I(S;W)=O(m\beta^2)$ (\cref{sec:comparison}, \cref{thm:stability-bound}), a one-way implication with no kernel handle, so the figure draws stability into information in the literature style and leaves the converse undrawn. Both endpoints are lighter-bordered nodes, which confirms that an edge between two unverified frameworks cannot be solid even when the implication is classical.

\item \textbf{The Rademacher generalization bound is the figure's nearest open frontier.} Ask what it would take to draw a \emph{second} solid edge out of the Rademacher node, the full generalization bound rather than the VC bridge: this is the figure's nearest verification frontier. The empirical Rademacher complexity, the Massart lemma, and the VC bound are machine-checked, but the assembled high-probability bound of \cref{thm:rademacher-bound} (symmetrization plus the two McDiarmid steps) is still a placeholder (\cref{op:formalize-rademacher}, \cite{BartlettMendelson2002}), so the Rademacher node is solid for its \emph{complexity measure} while its generalization \emph{guarantee} stays literature. Closing this would promote Rademacher's own bound to solid, alongside its already-solid bridge to VC.

\item \textbf{Formalizing stability and information would flip two lighter nodes to solid.} Which perturbation promotes the algorithm-based frameworks? Here is the frontier that the conceptual obstructions hide: the stability bound (\cref{thm:stability-bound}) and the Xu--Raginsky information bound (\cref{thm:xu-raginsky}, \cite{XuRaginsky2017}) have no verified handle, and a single conditional-mutual-information master inequality is conjectured to recover VC, PAC-Bayes, and stability at once (\cref{op:cmi-unification}, \cite{SteinkeZakynthinou2020}). Proving and formalizing it would recolor the stability and information nodes solid and replace the stability-to-information literature edge by a verified one, leaving only the genuinely orthogonal algorithm-versus-class edge conceptual.

\item \textbf{Is every framework a section of one risk-gap functor over its own fibration?} Read ``one quantity, many certificates'' as a universal property, and you reach the figure's deepest frontier. Each framework would be a section of a single risk-gap functor over a different fibration of the learning problem, the class for VC and Rademacher, the algorithm for stability and information, the posterior for PAC-Bayes, the operator norms for margin, so the verified core would read ``the functor admits a section over the class fibration'' (the three solid nodes) and the conceptual obstructions would read ``the class and algorithm fibrations do not refine one another,'' explaining the orthogonality as a fact about base categories rather than a separation. Whether such a single generalization-certificate object exists, of which the certified attention bound (\cref{prop:certified-attention}, \leanname{CertifiedGeneralization}) is a sixth, algorithm-specific section, or the fibrations are provably non-unifiable, is open (\cref{op:cmi-unification}, \cref{op:tighten-certified-attention}), so this reading is drawn dashed.
\end{enumerate}

\subsection*{Counterfactual exercises: generalizing Figure~\ref{fig:triangle}}

Where \Cref{fig:five-frameworks} graded the links, \Cref{fig:triangle} grades the geometry.  Its grammar is the placement: generalization information has three sources, the data, the hypothesis class, and the algorithm, set at the corners of a triangle, and every bound is a projection that reads one corner or one edge.  Three framework nodes are solid because their bounds are machine-checked (VC and Rademacher at the class corner, \leanname{finite_massart_lemma} and \leanname{EmpiricalRademacherComplexity}; PAC-Bayes on the algorithm--class edge, \leanname{pac_bayes_finite}); three are lighter and dashed because their bounds are literature (stability on the algorithm--data edge, information in the interior channel, margin on the class geometry).  No edge here is a theorem; the figure is pure projection, and the only verified content it carries is the grade of each node.  The obstructions of \Cref{fig:five-frameworks} reappear here as projections that simply land on different corners.  So as each exercise below perturbs the projection, classify what it yields: an \emph{instance} whose corner placement the verified core or the literature already exhibits, a \emph{boundary} where toggling one node grade or information source flips the evidence grade or the source being read, or a \emph{frontier} where a literature reading still waits to be promoted to a verified source.

\begin{enumerate}[leftmargin=*,itemsep=3pt]
\item \textbf{The class corner carries both verified bounds.} Take the two solid nodes at the class corner together and you have the figure's defining verified instance: VC uniform convergence and Rademacher complexity both read the hypothesis class and nothing else, and both have machine-checked bounds (\cref{thm:vc-uniform}, \leanname{finite_massart_lemma} and \leanname{sauer_shelah_exp_bound}; \cref{def:rademacher}, \leanname{EmpiricalRademacherComplexity} and \leanname{vcdim_bounds_rademacher_quantitative}). The class corner is therefore the most fully certified source in the diagram, and the solid border at that corner asserts exactly that the class-source reading is a kernel theorem.

\item \textbf{PAC-Bayes is the verified reading on the algorithm--class edge.} Why is the PAC-Bayes node solid yet off the class corner? This is the instance that anchors the posterior reading: the McAllester finite bound on the Gibbs risk is machine-checked (\cref{thm:pac-bayes}, \leanname{pac_bayes_finite}, with \leanname{gibbsError} the verified mean error of the posterior-sampled predictor), and PAC-Bayes reads the posterior over the class \emph{induced by the algorithm}, so its node sits on the edge joining the algorithm and class corners, solid like VC and Rademacher but reading two sources at once.

\item \textbf{One problem read from three corners is the three-bounds example.} \Cref{ex:three-bounds} is a tour of the triangle, and so the figure's organizing instance: linear classification with bounded norm admits a VC bound (class corner, $\VC=d+1$), a Rademacher bound (class corner, $\Rad_S\le 1/\sqrt{m}$), and a margin bound (class geometry), all on the \emph{same} problem. The diagram is the abstract witness that a single learning problem is bounded by reading whichever corner one can measure, and the three numbers in the example differ only because they project different sources.

\item \textbf{Which framework when is a corner-selection rule.} The ``Which Framework When?'' guide is a map back onto the triangle, which makes it the instance that names the figure's use: knowing the class but not the algorithm selects the class corner (Rademacher or margin), knowing the algorithm but not the class selects the algorithm readings (stability or information), and being able to choose a prior selects the algorithm--class edge (PAC-Bayes). The choice of framework is literally the choice of which corner of \cref{fig:triangle} one can measure.

\item \textbf{The obstructions are different-corner projections, not separations.} Toggle the companion figure's two conceptual edges onto this triangle and the figure's principal boundary appears: information reads the data-to-output channel in the interior while Rademacher reads the class corner, so the \cref{fig:five-frameworks} ``orthogonal'' edge is just these two nodes sitting at different places (\cref{sec:comparison}), and stability reads the algorithm--data edge while PAC-Bayes reads the algorithm--class edge, so the ``type mismatch'' edge is two nodes on different edges sharing only the algorithm corner. Neither is a separation theorem, which is why both figures draw them without a verified arrow.

\item \textbf{Stability and information read the algorithm, not the class.} Why are the two interior or algorithm-edge nodes dashed? This is the boundary that fixes their grade: stability measures the algorithm's sensitivity to one data point (\cref{thm:stability-bound}) and information measures the mutual information $I(S;W)$ between the sample and the algorithm's output (\cref{thm:xu-raginsky}, \cite{XuRaginsky2017}), so both read the algorithm rather than the class and neither has a verified handle, hence both nodes are drawn lighter. A deterministic ERM can carry $I(S;W)=H(W)$ on a class of small Rademacher complexity, which confirms the interior reading is genuinely a different source from the class corner.

\item \textbf{Margin reads class geometry, and that reading is literature.} The margin node near the class--algorithm side marks the boundary between a class-adjacent reading and a verified one: the margin bound is derived from Rademacher complexity applied to the ramp loss (\cref{thm:margin-bound}), so it reads the class geometry together with the predictor's margin, yet because thm:margin-bound is not machine-checked the node is drawn dashed even though it sits beside the solid class corner. Promoting it would require a verified ramp-loss contraction, the same step missing for the assembled Rademacher bound.

\item \textbf{Formalizing the algorithm readings would recolor the interior.} Which perturbation turns the dashed interior and algorithm-edge nodes solid? Here is the frontier the dashed nodes mark: the stability bound (\cref{thm:stability-bound}) and the Xu--Raginsky information bound (\cref{thm:xu-raginsky}, \cite{XuRaginsky2017}) have no verified handle, and the Rademacher node is solid for its complexity measure while its assembled high-probability \emph{guarantee} (\cref{thm:rademacher-bound}) is still a placeholder (\cref{op:formalize-rademacher}, \cite{BartlettMendelson2002}). Formalizing these would recolor the algorithm-reading and interior nodes solid and certify sources beyond the class corner.

\item \textbf{Is there a single bound that reads all three corners at once?} Could one certificate project the whole triangle rather than one corner or edge? That is the figure's unifying frontier: a single conditional-mutual-information master inequality is conjectured to recover VC, PAC-Bayes, and stability simultaneously (\cref{op:cmi-unification}, \cite{SteinkeZakynthinou2020}), which would be one bound reading the class, the algorithm, and the posterior together, collapsing the three solid-and-dashed readings into a single source-spanning certificate. Until it is proved and formalized the corners remain separate projections.

\item \textbf{Is the triangle the base of a fibration with each framework a section?} Read ``three sources, many projections'' as a universal property and you reach the figure's deepest frontier. Each framework would be a section of one risk-gap functor over a different face of the learning problem, the class face for VC and Rademacher, the algorithm face for stability and information, the posterior for PAC-Bayes, the operator norms for margin, so the verified core would read ``a section exists over the class face'' (the solid corner) and the obstructions would read ``the class and algorithm faces do not refine one another,'' explaining the orthogonality as a fact about base faces rather than a separation. Whether such a single generalization-certificate object exists, of which the certified attention bound (\cref{prop:certified-attention}, \leanname{CertifiedGeneralization}) is an algorithm-face section, or the faces are provably non-unifiable, is open (\cref{op:cmi-unification}, \cref{op:tighten-certified-attention}), so this reading is drawn dashed.
\end{enumerate}

% Chapter 13: Mind-Change Ordinals and Transfinite Hierarchies
%   mind_change_ordinal = Profile C (revelation): narrate the Freivalds-Smith discovery
%   constructive_ordinal = Profile A (vocabulary): brief setup of notation
%   bc_learning = Profile D (negative result): BC \ EX separation witness is the content
%   ordinal hierarchy = Profile B (load-bearing): full proof of strict separation
%   mind_change_bound = operational, brief
%   anomalous_learning = definition + hierarchy, brief
%
% Texture: recursion-theoretic with ordinal arithmetic. This chapter extends
% the diagonalization methods of Ch7 into the transfinite. The surprise is
% that ordinals are not merely convenient bookkeeping but the *right* measure.

\chapter{Mind-Change Ordinals and Transfinite Hierarchies}
\label{ch:ordinals}

How many times must a learner change its mind before converging to a correct
hypothesis?  In \Cref{ch:gold} we saw that $\Ex$-learners are permitted
finitely many mind changes, but the number may be unbounded: no single
integer $n$ suffices to bound the mind changes across all target functions in
the class.  This observation raises a natural question: can one \emph{measure}
the mind-change complexity of a class, and does that measure yield a strict
hierarchy of learning power?

The answer, due to Freivalds and Smith~\cite{FreivaldsSmith1993}, is that
the correct measure is not an integer but a \emph{constructive ordinal}.
The mind-change ordinal of a class assigns to each learnable class a
position in a transfinite hierarchy that refines the gap between finite
learning ($\FIN$) and identification in the limit ($\Ex$).  The hierarchy
is strict at every level, and its union does not exhaust $\Ex$.

This chapter develops the ordinal mind-change hierarchy in full.  We begin
with the minimal ordinal vocabulary needed (\Cref{sec:constructive-ordinals}),
then tell the story of how ordinals entered learning theory
(\Cref{sec:mco-revelation}), prove that the hierarchy is strict
(\Cref{sec:ordinal-hierarchy}), and close with the separation between
behaviorally correct and explanatory learning (\Cref{sec:bc-separation}).

\medskip

\begin{enumerate}[leftmargin=*,itemsep=2pt]
\item \textbf{Constructive ordinals} (\Cref{sec:constructive-ordinals}):
  Kleene's system~$\mathcal{O}$, notation for computable ordinals.
\item \textbf{The mind-change ordinal} (\Cref{sec:mco-revelation}):
  how ordinals became the right measure of learning complexity.
\item \textbf{The strict hierarchy} (\Cref{sec:ordinal-hierarchy}):
  $\forall \alpha < \beta$, $\Ex_\alpha \subsetneq \Ex_\beta$, with full proof.
\item \textbf{Anomalous learning} (\Cref{sec:anomalous}): the $\Ex_a$
  hierarchy and its interaction with mind changes.
\item \textbf{Behaviorally correct learning} (\Cref{sec:bc-separation}):
  $\BC \setminus \Ex$ via diagonal witness.
\item \textbf{The full landscape} (\Cref{sec:ordinal-landscape}):
  a diagram and summary of all inclusions.
\end{enumerate}


% ================================================================
% SECTION 1: CONSTRUCTIVE ORDINALS (Profile A -- vocabulary)
% ================================================================

\section{Constructive Ordinals}
\label{sec:constructive-ordinals}

We need ordinals that a Turing machine can manipulate.  The full class of
countable ordinals is too large: many have no computable description.  Kleene's
system~$\mathcal{O}$ provides the computable fragment.

\begin{defn}[Kleene's System $\mathcal{O}$]
\label{def:system-O}
The \emph{constructive ordinals} are a subset $\mathcal{O} \subseteq \N$
together with a partial ordering $<_{\mathcal{O}}$ and a map
$|\cdot| : \mathcal{O} \to \omega_1^{\mathrm{CK}}$ (the Church--Kleene
ordinal) defined inductively:
\begin{enumerate}[leftmargin=*,itemsep=2pt]
\item $1 \in \mathcal{O}$ with $|1| = 0$.
\item If $a \in \mathcal{O}$, then $2^a \in \mathcal{O}$ with
  $|2^a| = |a| + 1$, and $a <_{\mathcal{O}} 2^a$.
\item If $\varphi_e$ is a total computable function with
  $\varphi_e(0) <_{\mathcal{O}} \varphi_e(1) <_{\mathcal{O}} \cdots$,
  then $3 \cdot 5^e \in \mathcal{O}$ with
  $|3 \cdot 5^e| = \sup_n |\varphi_e(n)|$, and
  $\varphi_e(n) <_{\mathcal{O}} 3 \cdot 5^e$ for all $n$.
\end{enumerate}
\end{defn}

\begin{notation}
\label{not:ordinal-shorthand}
We write ordinals in standard notation ($0, 1, 2, \ldots, \omega,
\omega + 1, \ldots, \omega \cdot 2, \ldots, \omega^2, \ldots,
\omega^\omega, \ldots$) rather than their codes in $\mathcal{O}$.
When we say ``ordinal $\alpha$'' in the context of mind-change bounds,
we always mean a constructive ordinal: one that has a notation in
$\mathcal{O}$.
\end{notation}

\begin{remark}
\label{rem:system-O-noncomputable}
Membership in $\mathcal{O}$ is $\Pi_1^1$-complete and hence far from
decidable.  This does not obstruct its use as a mind-change counter:
the learner is given a specific notation $a \in \mathcal{O}$ and needs
only to compute the predecessor relation, which \emph{is} computable
given the notation.  The learner never needs to decide whether an
arbitrary number belongs to $\mathcal{O}$.
\end{remark}


% ================================================================
% SECTION 2: THE MIND-CHANGE ORDINAL (Profile C -- revelation)
% ================================================================

\section{The Mind-Change Ordinal: How Ordinals Entered Learning Theory}
\label{sec:mco-revelation}

The story begins with a natural frustration.

\subsection*{The problem that was stuck}

After Gold's 1967 theorem established the basic framework and Case and
Smith's 1983 paper~\cite{CaseSmith1983} introduced the hierarchy of
success criteria ($\FIN$, $\Ex$, $\BC$, and their variants), a question
remained open throughout the 1980s: \emph{how should one measure the
mind-change complexity of a class?}

The naive approach is to count mind changes with natural numbers.  Define
$\Ex_n$ as the collection of classes learnable by a machine that changes
its mind at most $n$ times on any input.  This gives a hierarchy:
\[
  \FIN = \Ex_0 \subsetneq \Ex_1 \subsetneq \Ex_2 \subsetneq \cdots
  \subsetneq \Ex.
\]
Velauthapillai~\cite{Velauthapillai1989} proved these inclusions are
strict.  But the hierarchy has a defect: the union
$\bigcup_{n \in \N} \Ex_n$ does \emph{not} equal $\Ex$.

\begin{defn}[Mind-Change Count]
\label{def:mc-count}
Let $M$ be a learner and $f$ a target function.  The \emph{mind-change
count} of $M$ on $f$, denoted $\mathrm{mc}(M, f)$, is the number of
times $M$ changes its hypothesis while processing a text for $f$:
\[
  \mathrm{mc}(M, f) \;=\; \bigl|\{t \in \N : h_t \neq h_{t+1}\}\bigr|.
\]
The \emph{mind-change complexity} of a class $\mathcal{C}$ with respect
to $M$ is $\mathrm{mc}(M, \mathcal{C}) = \sup_{f \in \mathcal{C}}
\mathrm{mc}(M, f)$.
\end{defn}

The problem is that $\mathrm{mc}(M, \mathcal{C})$ may be finite for every
$f \in \mathcal{C}$ yet unbounded across the class.  The supremum is
$\omega$, but no single integer bounds all instances.  There exist
$\Ex$-learnable classes in $\Ex \setminus \bigcup_n \Ex_n$: the learner
converges on every input but there is no uniform bound on the number of
mind changes.

\begin{computational}
Consider the class $\mathcal{C} = \{f_n : n \in \N\}$ where $f_n(x) = 0$
for $x < n$ and $f_n(x) = 1$ for $x \geq n$.  A learner that guesses
``$f_n$ for the largest $n$ seen so far'' is an $\Ex$-learner for
$\mathcal{C}$.  On input $f_n$, it changes its mind exactly $n$ times
(from $f_0$ to $f_1$ to $\ldots$ to $f_n$).  The mind-change count is
$n$, which is finite for each $f_n$ but unbounded across $\mathcal{C}$.
No $\Ex_k$ contains $\mathcal{C}$.
\end{computational}

\subsection*{The obvious approach that fails}

One might try to rescue the integer hierarchy by allowing $\Ex_\omega$
as a ``catch-all'' for classes learnable with finitely many but unbounded
mind changes.  But this merely relabels the gap; it does not give a
\emph{refined} hierarchy.  The gap between $\bigcup_n \Ex_n$ and $\Ex$
may itself contain structure: some classes might require the learner
to perform mind changes that depend on mind changes, a kind of
second-order revision.  Is there a hierarchy that captures this structure?

\subsection*{The Freivalds--Smith construction}

In 1993, Freivalds and Smith~\cite{FreivaldsSmith1993} answered this
question with a construction that remains surprising.  They observed
that the \emph{ordinal arithmetic of mind changes} is the correct
framework.

\begin{defn}[Ordinal Mind-Change Bound]
\label{def:ordinal-mc}
Let $\alpha$ be a constructive ordinal.  A learner $M$
\emph{$\alpha$-bounds its mind changes} if $M$ carries a counter
initialized to some notation for $\alpha$, and at each mind change,
$M$ replaces the counter with a strictly smaller ordinal
(in $<_{\mathcal{O}}$).  Since there is no infinite descending
sequence in the ordinals, $M$ must eventually stop changing its mind.
\end{defn}

\begin{defn}[$\Ex_\alpha$]
\label{def:ex-alpha}
For a constructive ordinal $\alpha$, the class $\Ex_\alpha$ consists
of all concept classes $\mathcal{C}$ such that some learner
$\alpha$-bounds its mind changes and $\Ex$-identifies $\mathcal{C}$.
\end{defn}

\begin{remark}[The hierarchy is notation-relative]
\label{rem:notation-dependence}
A subtlety underlies every statement that follows.  The class
$\Ex_\alpha$ is defined through a \emph{notation} $a \in \mathcal{O}$
for~$\alpha$, not through the ordinal $|a|$ alone, because the learner
decrements along the computable predecessor sequence that the notation
supplies.  Two notations $a, a'$ with $|a| = |a'|$ can give
$\Ex_a \neq \Ex_{a'}$ (Exercise~6).  When the theorems below write
``$\Ex_\alpha$'' and ``for all constructive~$\alpha$,'' they mean: fix a
notation for each ordinal; the strictness statements hold for these
fixed notations.  Which ordinals (as opposed to notations) admit a
notation-invariant $\Ex_\alpha$ is the subject
of~\cite{MartinSharmaStephan2006} and is only partially understood.
\end{remark}

The key insight is operational: a counter initialized to $\omega$ can
be decremented to any finite value $n$, but a counter initialized to
$n$ can only be decremented $n$ times.  So $\Ex_\omega$ is strictly
stronger than every $\Ex_n$: a learner with an $\omega$-counter can
``spend'' its first mind change moving from $\omega$ to some finite $n$,
after which it has $n$ mind changes remaining.  The value $n$ can
depend on the data seen so far.

\begin{historical}
Freivalds and Smith's construction was motivated by an analogy with
proof theory, where ordinal analysis assigns ordinal ``proof-theoretic
strengths'' to formal systems.  Just as the consistency strength of
Peano arithmetic is measured by the ordinal $\varepsilon_0$, the
mind-change complexity of a class is measured by a constructive ordinal.
The analogy is not perfect (the ordinals in learning theory stay well
below $\varepsilon_0$), but the structural parallel guided the discovery.
\end{historical}

\begin{example}
\label{ex:omega-learner}
Recall the class $\mathcal{C} = \{f_n : n \in \N\}$ from the
computational illustration above.  A learner with an $\omega$-counter
operates as follows:
\begin{enumerate}[leftmargin=*,itemsep=2pt]
\item Initialize the counter to $\omega$.
\item Upon seeing the first datum $f(0)$, make an initial guess and
  set the counter to some finite value (say, the current time step).
\item Each subsequent mind change decrements the (now finite) counter.
\end{enumerate}
After the first mind change, the learner has a finite counter, so it
changes its mind at most finitely many times.  The key point is
that the size of the finite counter is chosen \emph{adaptively}: it
depends on the data.  This is exactly what an $\omega$-counter provides
and a fixed integer counter does not.
\end{example}

\begin{remark}
\label{rem:ordinal-vs-integer}
The passage from $\bigcup_n \Ex_n$ to $\Ex_\omega$ is
not a mere notational trick.  It reflects a genuine increase in
computational power: the learner can defer its commitment to a finite
bound until it has seen enough data to determine what the bound should
be.  This deferred-commitment mechanism is the operational content of
transfinite mind changes.
\end{remark}


% ================================================================
% SECTION 3: THE STRICT HIERARCHY (Profile B -- load-bearing)
% ================================================================

\section{The Ordinal Mind-Change Hierarchy}
\label{sec:ordinal-hierarchy}

The central theorem of this chapter asserts that the ordinal-indexed
hierarchy is strict at every level.

\begin{thm}[Freivalds--Smith 1993; Ambainis--Jain--Sharma 1999]
\label{thm:ordinal-hierarchy}
For all constructive ordinals $\alpha < \beta$,
\[
  \Ex_\alpha \subsetneq \Ex_\beta \subsetneq \Ex.
\]
Moreover, $\bigcup_{\alpha \in \mathcal{O}} \Ex_\alpha \subsetneq \Ex$.
\end{thm}

The proof has two components: the inclusion $\Ex_\alpha \subseteq
\Ex_\beta$ (easy), and the strictness $\Ex_\alpha \neq \Ex_\beta$
(which requires constructing a witness class in $\Ex_\beta \setminus
\Ex_\alpha$).  We develop each in turn.

\subsection{Inclusions}

\begin{lemma}
\label{lem:inclusion}
If $\alpha <_{\mathcal{O}} \beta$, then $\Ex_\alpha \subseteq \Ex_\beta$.
\end{lemma}

\begin{proof}
Let $\mathcal{C} \in \Ex_\alpha$, witnessed by learner $M$ with an
$\alpha$-counter.  Since $\alpha < \beta$, we can initialize a
$\beta$-counter and immediately decrement it to $\alpha$.  (Formally:
replace the counter $\beta$ with $\alpha$ at the first step, which is
a valid decrement since $\alpha <_{\mathcal{O}} \beta$.)  Then run $M$
with the remaining $\alpha$-counter.  This witnesses
$\mathcal{C} \in \Ex_\beta$.
\end{proof}

\subsection{Witness classes for strictness}

The hard direction requires, for each pair $\alpha < \beta$, a class
$\mathcal{C}_{\alpha,\beta} \in \Ex_\beta \setminus \Ex_\alpha$.  We
construct this via a \emph{self-referential coding} argument.

\begin{lemma}[Witness Construction]
\label{lem:witness}
For every constructive ordinal $\alpha$, there exists a class
$\mathcal{C}_\alpha \in \Ex_{\alpha+1} \setminus \Ex_\alpha$.
\end{lemma}

\begin{proof}
We construct $\mathcal{C}_\alpha$ so that:
\begin{enumerate}[label=(\roman*)]
\item an $(\alpha+1)$-bounded learner can identify every function
  in $\mathcal{C}_\alpha$, but
\item no $\alpha$-bounded learner can identify all functions in
  $\mathcal{C}_\alpha$.
\end{enumerate}

\textbf{Construction.}
Let $M_0, M_1, M_2, \ldots$ be an effective enumeration of all
$\alpha$-bounded learners (i.e., all Turing machines equipped with
an $\alpha$-counter).  We define, for each $e \in \N$, a function
$f_e \in \mathcal{C}_\alpha$ designed to defeat learner $M_e$.

The function $f_e$ is defined by stages.  At each stage~$s$, $f_e$ has
been defined on arguments $0, 1, \ldots, s-1$.  We simulate $M_e$ on
the text $(f_e(0), f_e(1), \ldots, f_e(s-1))$ and observe $M_e$'s
hypothesis and counter value.  There are two cases:

\begin{itemize}[leftmargin=*]
\item \textbf{Case 1:} $M_e$ has exhausted its counter (the counter has
  reached~$0$).  Then $M_e$ can no longer change its mind.  It is now
  committed to some hypothesis $h$.  We define $f_e$ on the remaining
  arguments to disagree with $\varphi_h$: choose $f_e(s)$ so that
  $f_e(s) \neq \varphi_h(s)$ (this is possible since $\varphi_h$ is a
  fixed computable function and we are free to define $f_e$).  Then
  $M_e$ fails to $\Ex$-identify $f_e$.

\item \textbf{Case 2:} $M_e$ still has counter $> 0$.  We extend $f_e$
  by one more value, choosing $f_e(s)$ to force $M_e$ to change its mind
  (if possible) or simply setting $f_e(s) = 0$ if $M_e$ does not change
  its mind on any extension.
\end{itemize}

\textbf{Defeating $M_e$.}
By iterating Case~2, we either:
\begin{itemize}[leftmargin=*]
\item force $M_e$ to exhaust its $\alpha$-counter (each forced mind
  change decrements the counter, and there is no infinite descending
  chain in the ordinals), after which Case~1 applies; or
\item reach a point where $M_e$ refuses to change its mind no matter
  how we extend $f_e$, in which case $M_e$ has committed to a
  hypothesis $h$ and we can diagonalize as in Case~1.
\end{itemize}

In both sub-cases, $M_e$ fails on $f_e$.  Hence no $\alpha$-bounded
learner identifies all of $\mathcal{C}_\alpha = \{f_e : e \in \N\}$.

\medskip
\textbf{An $(\alpha+1)$-bounded learner for $\mathcal{C}_\alpha$.}
An $(\alpha+1)$-bounded learner $M^*$ for $\mathcal{C}_\alpha$ works
as follows.  Given text for some $f_e$, the learner $M^*$ searches
for the index $e$ by dovetailing: it simulates the construction of
$f_0, f_1, f_2, \ldots$ in parallel with the incoming data, and
hypothesizes the smallest $e$ consistent with the data seen so far.
Each time the current hypothesis $e$ is refuted by new data, $M^*$
moves to the next candidate, decrementing its counter.

The counter is initialized to $\alpha + 1$.  The first mind change
(from the initial hypothesis to the first genuine candidate)
decrements the counter from $\alpha + 1$ to $\alpha$.  That the
subsequent mind changes are bounded by~$\alpha$ is the crux of the
construction, and it is exactly here that this proof is a sketch: the
$\alpha$-bound on $M^*$'s residual revisions follows by reconstructing
the diagonal construction from the data via the recursion theorem, a
priority argument we do not carry out.  See Freivalds and
Smith~\cite{FreivaldsSmith1993} for the complete construction.
\end{proof}

\begin{proof}[Proof of \Cref{thm:ordinal-hierarchy}]
Inclusions hold by \Cref{lem:inclusion}.  For strictness: given
$\alpha < \beta$, the class $\mathcal{C}_\alpha$ from
\Cref{lem:witness} lies in $\Ex_{\alpha+1} \subseteq \Ex_\beta$ but
not in $\Ex_\alpha$.

For the final strict inclusion $\bigcup_\alpha \Ex_\alpha \subsetneq
\Ex$: this requires a class $\mathcal{C}$ that is $\Ex$-learnable but
not $\Ex_\alpha$-learnable for any constructive $\alpha$.  The
construction uses a priority argument: define $\mathcal{C}$ by
simultaneously diagonalizing against all $\alpha$-bounded learners
for all $\alpha$.  The learner for $\mathcal{C}$ is an ordinary
$\Ex$-learner with no ordinal counter; its mind-change count is finite
on each input but grows faster than any constructive ordinal as a
function of the input.  We refer to Ambainis, Jain, and
Sharma~\cite{AmbainisJainSharma1999} for the full priority construction,
which uses techniques from $\alpha$-recursion theory beyond the scope
of this chapter.
\end{proof}


\subsection{The hierarchy at limit ordinals}

The successor case $\Ex_\alpha \subsetneq \Ex_{\alpha+1}$ is the
engine of the hierarchy.  But limit ordinals also mark strict jumps.

\begin{prop}
\label{prop:limit-ordinals}
For every limit ordinal $\lambda$ with a notation in $\mathcal{O}$,
\[
  \bigcup_{\alpha < \lambda} \Ex_\alpha \subsetneq \Ex_\lambda.
\]
\end{prop}

\begin{proof}
The inclusion $\bigcup_{\alpha < \lambda} \Ex_\alpha \subseteq
\Ex_\lambda$ follows from \Cref{lem:inclusion}.  For strictness,
let $\lambda = \sup_n \alpha_n$ where $(\alpha_n)$ is the computable
sequence witnessing $\lambda \in \mathcal{O}$.

Define a class $\mathcal{C}_\lambda = \{f_n : n \in \N\}$ where
$f_n$ requires exactly $\alpha_n$ mind changes to identify.  (Each
$f_n$ is drawn from the witness class $\mathcal{C}_{\alpha_n}$ of
\Cref{lem:witness}.)  Then $\mathcal{C}_\lambda$ is not in
$\Ex_{\alpha_k}$ for any $k$ (since $f_{k+1}$ requires more than
$\alpha_k$ mind changes), but an $\Ex_\lambda$-learner can handle
it: upon seeing enough data to determine which $f_n$ is the target,
the learner sets its counter to $\alpha_n$ (which is
$<_{\mathcal{O}} \lambda$) and proceeds.
\end{proof}

\begin{remark}
\label{rem:omega-to-n}
The case $\lambda = \omega$ is particularly instructive.
$\Ex_\omega \setminus \bigcup_n \Ex_n$ consists of classes where the
learner needs finitely many mind changes on each input, but the
number of mind changes must be chosen adaptively.  The $\omega$-counter
captures exactly this deferred-commitment pattern.
\end{remark}


% ================================================================
% SECTION 4: ANOMALOUS LEARNING (brief)
% ================================================================

\section{Anomalous Learning}
\label{sec:anomalous}

The mind-change hierarchy refines $\Ex$ by bounding the number of
hypothesis revisions.  A different axis of relaxation tolerates
\emph{errors} in the final hypothesis.

\begin{defn}[$\Ex_a^*$ -- Anomalous Learning {\cite{CaseSmith1983}}]
\label{def:anomalous-ord}
For $a \in \N$, the class $\Ex_a^*$ consists of all concept classes
$\mathcal{C}$ for which there exists a learner that converges to a
hypothesis correct on all but at most $a$ inputs.  $\Ex^*$ allows
finitely many anomalies without a fixed bound:
$\Ex^* = \bigcup_{a \in \N} \Ex_a^*$.
\end{defn}

\begin{thm}[Anomaly Hierarchy {\cite{CaseSmith1983}}]
\label{thm:anomaly-hierarchy}
\[
  \Ex = \Ex_0^* \subsetneq \Ex_1^* \subsetneq \Ex_2^* \subsetneq \cdots
  \subsetneq \Ex^* \subsetneq \BC_0 \subsetneq \BC.
\]
Each inclusion is strict.
\end{thm}

\begin{remark}
\label{rem:anomaly-ordinal}
The anomaly hierarchy and the mind-change hierarchy are \emph{orthogonal}
axes of relaxation.  One can combine them: $\Ex_\alpha^a$ denotes
learning with an $\alpha$-bounded mind-change counter and tolerance for
$a$ anomalies.  The resulting two-dimensional hierarchy was studied by
Sharma, Stephan, and Ventsov~\cite{SharmaStephanVentsov2004}, who
showed that it yields a strict partial order under inclusion.
\end{remark}


% ================================================================
% SECTION 5: BC LEARNING (Profile D -- negative result)
% ================================================================

\section{Behaviorally Correct Learning and the Separation $\BC \setminus \Ex$}
\label{sec:bc-separation}

In $\Ex$-learning, convergence is \emph{syntactic}: the learner must
eventually output the same hypothesis index forever.  $\BC$-learning
relaxes this to \emph{semantic} convergence: the learner must eventually
output hypotheses that are all extensionally correct, but the index
may keep changing.

\begin{defn}[$\BC$-Learning {\cite{CaseSmith1983}}]
\label{def:bc-learning-ord}
A learner $M$ \emph{$\BC$-identifies} a function $f$ if there exists
$t_0$ such that for all $t \geq t_0$, $\varphi_{h_t} = f$
(extensional equality).  The hypothesis index $h_t$ may continue
to change.  The class $\BC$ consists of all concept classes
$\BC$-identifiable by some learner.
\end{defn}

The separation $\BC \setminus \Ex \neq \emptyset$ is a \emph{negative
result about $\Ex$}: there are classes where semantic convergence is
achievable but syntactic convergence is not.  The witness construction
is the content.

\begin{thm}[Case--Smith 1983]
\label{thm:bc-separation}
$\BC \setminus \Ex \neq \emptyset$.  There exists a class of total
recursive functions that is $\BC$-identifiable but not
$\Ex$-identifiable.
\end{thm}

\begin{proof}
\textbf{The witness class.}
For each total recursive function $g$ and each index $e$ with
$\varphi_e = g$, define the function $f_{g,e} : \N \to \N$ by
\[
  f_{g,e}(x) \;=\;
  \begin{cases}
    e & \text{if } x = 0, \\
    g(x) & \text{if } x \geq 1.
  \end{cases}
\]
The class is
\[
  \mathcal{C} \;=\; \bigl\{f_{g,e} : g \text{ total recursive},\;
  \varphi_e = g\bigr\}.
\]
Each total recursive function $g$ has infinitely many indices, so
$\mathcal{C}$ contains infinitely many functions that agree on
$\{1, 2, 3, \ldots\}$ but differ at argument~$0$.

\medskip
\textbf{$\mathcal{C} \in \BC$.}
A $\BC$-learner for $\mathcal{C}$ proceeds as follows.  Given text
for some $f_{g,e} \in \mathcal{C}$, the learner eventually sees
$f_{g,e}(0) = e$.  At time $t$, having observed
$f_{g,e}(0), f_{g,e}(1), \ldots, f_{g,e}(t)$, the learner outputs
an index $h_t$ for the function that agrees with the observed data
on $\{0, \ldots, t\}$ and follows $\varphi_e$ on $\{t+1, t+2, \ldots\}$.

Since $\varphi_e = g$ and $f_{g,e}(x) = g(x)$ for $x \geq 1$, each
hypothesis $h_t$ (for $t$ large enough that $f_{g,e}(0)$ has been
observed) computes $f_{g,e}$.  The indices $h_t$ may differ (the
learner may construct different programs at each step), but they all
compute the same function.  This is $\BC$-identification.

\medskip
\textbf{$\mathcal{C} \notin \Ex$.}
Suppose for contradiction that learner $M$ $\Ex$-identifies
$\mathcal{C}$.  Then for every $f_{g,e} \in \mathcal{C}$, $M$ must
converge to a single index $h^*$ with $\varphi_{h^*} = f_{g,e}$.

Fix any total recursive $g$.  For each index $e$ with $\varphi_e = g$,
the function $f_{g,e}$ belongs to $\mathcal{C}$.  The learner $M$,
on input $f_{g,e}$, sees $(e, g(1), g(2), g(3), \ldots)$ and must
converge to some $h_e^*$ with $\varphi_{h_e^*} = f_{g,e}$.

Consider two distinct indices $e \neq e'$ with $\varphi_e = \varphi_{e'} = g$.
The functions $f_{g,e}$ and $f_{g,e'}$ differ only at argument~$0$:
$f_{g,e}(0) = e \neq e' = f_{g,e'}(0)$.  So $M$ must converge to
distinct hypotheses $h_e^*$ and $h_{e'}^*$ (since
$\varphi_{h_e^*}(0) = e \neq e' = \varphi_{h_{e'}^*}(0)$).

But after time~$0$, both data streams are identical: $g(1), g(2), \ldots$.
The learner's behavior from time~$1$ onward depends only on its
hypothesis after seeing $f(0)$ and the shared tail.  Since $g$ has
infinitely many indices $e_0, e_1, e_2, \ldots$, the learner $M$ must
produce infinitely many distinct convergent hypotheses $h_{e_0}^*,
h_{e_1}^*, h_{e_2}^*, \ldots$, all determined by the single initial
datum $f(0)$ and the same infinite tail.

The contradiction is now information-theoretic, sharpened by the
recursion theorem.  Fix any total recursive~$g$ and let
$e_0, e_1, e_2, \ldots$ be infinitely many indices with
$\varphi_{e_i} = g$.  The inputs $f_{g,e_0}, f_{g,e_1}, \ldots$ share
the identical tail $g(1), g(2), \ldots$ and differ only in the initial
label.  After processing that label, $M$'s subsequent behavior is driven
by the \emph{same} tail, so it follows a single hypothesis trajectory;
yet $\Ex$-identification demands a \emph{distinct} final index
$h_{e_i}^*$ for each~$i$, since $\varphi_{h_{e_i}^*}(0) = e_i$.  A single
convergent trajectory cannot end at infinitely many distinct indices.
To turn this tension into a concrete failure, the operator recursion
theorem~\cite{CaseSmith1983} supplies a total recursive~$g$ together with
two of its indices $e \neq e'$ on which $M$ produces the \emph{same}
limiting hypothesis~$h^*$; then $\varphi_{h^*}(0)$ cannot equal both
$e$ and~$e'$, so $M$ misidentifies at least one of
$f_{g,e}, f_{g,e'}$.  We do not reproduce the priority bookkeeping of the
operator recursion theorem here.
\end{proof}

\begin{separation}
The class $\mathcal{C}$ witnesses $\BC \setminus \Ex \neq \emptyset$.
Here is why the $\Ex$-learner cannot escape.

A $\BC$-learner has room to maneuver: it may output different indices
for the same function, exploiting the many-to-one nature of G\"odel
numbering.  An $\Ex$-learner has no such room.  It must converge to a
\emph{single} index.

Now fill $\mathcal{C}$ with infinitely many functions that differ only
in their self-referential label at argument~$0$.  The $\BC$-learner
reads the label, builds a correct program, moves on.  The $\Ex$-learner
reads the label too, but it must commit.  Commit to which index?
Every candidate is correct extensionally for the current function but
wrong syntactically for all the others.  The learner converges; the
class shifts; the index is wrong.  The trap is shut.

The obstruction is the gap between extensional and intensional identity:
$\BC$ asks only that the program compute the right function; $\Ex$
demands the right program.
\end{separation}


% ================================================================
% SECTION 6: THE FULL LANDSCAPE (diagram)
% ================================================================

\section{The Full Landscape}
\label{sec:ordinal-landscape}

\Cref{fig:ordinal-hierarchy} displays the complete hierarchy of
mind-change-bounded and anomaly-bounded learning criteria.

\begin{figure}[htbp]
\centering
\begin{tikzpicture}[
  node distance=0.6cm and 1.2cm,
  every node/.style={font=\small},
  basecrit/.style={draw, very thick, rounded corners, minimum width=1.9cm,
    minimum height=0.58cm, align=center},
  litlevel/.style={draw, thin, dashed, rounded corners, minimum width=1.8cm,
    minimum height=0.54cm, align=center},
  litincl/.style={-{Stealth[length=5pt]}, thick, dashed, color=thmred},
  citetag/.style={font=\scriptsize\itshape, color=notegray},
]
% --- mind-change column (left, bottom to top): SOLID base FIN/Ex, DASHED ordinal levels ---
\node[basecrit, fill=defblue!10] (fin) {$\FIN = \Ex_0$};
\node[litlevel, fill=defblue!8, above=of fin] (ex1) {$\Ex_1$};
\node[litlevel, fill=defblue!8, above=of ex1] (ex2) {$\Ex_2$};
\node[above=0.26cm of ex2] (dots1) {$\vdots$};
\node[litlevel, fill=defblue!8, above=0.26cm of dots1] (exn) {$\Ex_n$};
\node[above=0.26cm of exn] (dots2) {$\vdots$};
\node[litlevel, fill=edgeorange!14, above=0.26cm of dots2] (exomega) {$\Ex_\omega$};
\node[litlevel, fill=edgeorange!14, above=of exomega] (exomega2) {$\Ex_{\omega \cdot 2}$};
\node[above=0.26cm of exomega2] (dots3) {$\vdots$};
\node[litlevel, fill=edgeorange!14, above=0.26cm of dots3] (exomega2a) {$\Ex_{\omega^2}$};
\node[above=0.26cm of exomega2a] (dots4) {$\vdots$};
\node[litlevel, fill=thmred!9, above=0.26cm of dots4] (exall) {$\bigcup_\alpha \Ex_\alpha$};
\node[basecrit, fill=thmred!14, above=of exall] (ex) {$\Ex$};
% --- anomaly arm (up-right from Ex): DASHED levels, SOLID BC ---
\node[litlevel, fill=sepgreen!10, right=2.4cm of ex] (exa1) {$\Ex_1^*$};
\node[litlevel, fill=sepgreen!10, above=of exa1] (exa2) {$\Ex_2^*$};
\node[above=0.26cm of exa2] (dots5) {$\vdots$};
\node[litlevel, fill=sepgreen!10, above=0.26cm of dots5] (exastar) {$\Ex^*$};
\node[litlevel, fill=analogypurple!10, above=of exastar] (bc0) {$\BC_0$};
\node[basecrit, fill=analogypurple!14, above=of bc0] (bc) {$\BC$};
% --- inclusion arrows: ALL dashed-literature (every rung is classical, none kernel-verified) ---
\draw[litincl] (fin) -- (ex1);
\draw[litincl] (ex1) -- (ex2);
\draw[litincl] (exn.north) -- (dots2.south);
\draw[litincl] (dots2.north) -- (exomega.south);
\draw[litincl] (exomega) -- (exomega2);
\draw[litincl] (exomega2a.north) -- (dots4.south);
\draw[litincl] (dots4.north) -- (exall.south);
\draw[litincl] (exall) -- (ex);
\draw[litincl] (ex) -- (exa1);
\draw[litincl] (exa1) -- (exa2);
\draw[litincl] (exastar) -- (bc0);
\draw[litincl] (bc0) -- (bc);
% --- one cite tag per author-result, at the regime boundary ---
\node[citetag, right=2pt of ex1] {Velauthapillai};
\node[citetag, left=2pt of exomega] {Freivalds--Smith};
\node[citetag, left=2pt of exall] {Ambainis--Jain--Sharma};
\node[citetag, right=2pt of exa2] {Case--Smith};
\node[font=\scriptsize, color=edgeorange, right=2pt of exomega2] {$\supsetneq \bigcup_n \Ex_n$};
% --- braces ---
\draw[decorate, decoration={brace, amplitude=6pt, mirror}, color=defblue!70]
  ([xshift=-1.7cm]fin.south west) -- ([xshift=-1.7cm]ex.north west)
  node[midway, left=8pt, align=center, font=\scriptsize, color=defblue!70] {mind-change\\hierarchy};
\draw[decorate, decoration={brace, amplitude=6pt}, color=sepgreen!70]
  ([xshift=1.6cm]exa1.south east) -- ([xshift=1.6cm]bc.north east)
  node[midway, right=8pt, align=center, font=\scriptsize, color=sepgreen!70] {anomaly\\hierarchy};
% --- VERIFIED ordinal-rank inset: the kernel's actual transfinite contribution, a SEPARATE measure ---
\node[draw, very thick, rounded corners, fill=defblue!5, align=left, font=\scriptsize,
  text width=6.1cm, below=1.4cm of fin.south west, anchor=north west] (rankbox) {%
  \textbf{Verified ordinal rank} (dimension side, a separate measure)\\
  \leanname{withTopNatToOrdinal} embeds $\N\cup\{\infty\}$ into the ordinals, sending $\infty$ to $\omega$\\
  \leanname{vcdim_to_ordinal_vcdim}: finite VC gives a finite rank, $\VC{=}\infty$ gives $\omega$\\
  finite-capacity classes (softmax attention, halfspaces) sit at the bottom, below $\omega$};
% --- evidence-grade legend ---
\coordinate (leg) at ([yshift=-0.42cm]rankbox.south west);
\draw[very thick] (leg) -- ++(0.55,0);
\node[font=\scriptsize, anchor=west] at ([xshift=0.62cm]leg) {kernel-verified ($sc{=}0$)};
\draw[thick, dashed, color=thmred] ([xshift=4.0cm]leg) -- ++(0.55,0);
\node[font=\scriptsize, anchor=west] at ([xshift=4.62cm]leg) {classical literature (cited)};
\draw[thick, dashed, color=edgeorange] ([yshift=-0.42cm]leg) -- ++(0.55,0);
\node[font=\scriptsize, anchor=west] at ([xshift=0.62cm, yshift=-0.42cm]leg)
  {open: not yet machine-checked (\cref{op:formalize-ordinal-hierarchy}, \cref{op:formalize-bc})};
\end{tikzpicture}
\caption[The ordinal complexity spine of limit identification: a classical tower with a verified rank at its base.]{The
number of mind changes a limit learner needs (left column) or the number of anomalies it
tolerates (right column) is a constructive ordinal, and that ordinal stratifies the
identification criteria into a strict tower: from finite identification $\FIN = \Ex_0$ up
through the transfinite levels $\Ex_\alpha$ into explanatory learning $\Ex$, and on through the
anomaly levels $\Ex_a^*$ into behaviorally correct learning $\BC$.  Every strict inclusion is a
classical theorem (the finite levels by Velauthapillai, the first limit jump
$\bigcup_n \Ex_n \subsetneq \Ex_\omega$ by Freivalds--Smith, the transfinite levels and
$\bigcup_\alpha \Ex_\alpha \subsetneq \Ex$ by Ambainis--Jain--Sharma, the anomaly axis by
Case--Smith), and each inclusion is drawn dashed because none is machine-checked and formalizing
the tower is still open.  The three base criteria $\FIN$, $\Ex$, $\BC$ are formalized, so their
nodes are drawn solid; the ordinal and anomaly levels are literature definitions and drawn
lighter.  The kernel certifies a separate ordinal rank on the dimension side, in the lower box:
an order-preserving embedding sends an infinite VC dimension to the first transfinite rank
$\omega$ and fixes every finite dimension below it, so every finite-capacity class of this book,
including softmax attention and the geometric classes, sits at a finite rank at the very bottom.
The transfinite tower thus becomes necessary only for the exact identification of programs.}
\label{fig:ordinal-hierarchy}
\end{figure}

We summarize the key structural facts.

\begin{thm}[Summary of Inclusions]
\label{thm:summary}
The following inclusions are all strict:
\begin{enumerate}[label=(\roman*)]
\item $\FIN = \Ex_0 \subsetneq \Ex_1 \subsetneq \cdots \subsetneq
  \Ex_n \subsetneq \cdots$ \quad (finite levels,
  Velauthapillai~\cite{Velauthapillai1989}).
\item $\bigcup_{n < \omega} \Ex_n \subsetneq \Ex_\omega$ \quad
  (first limit ordinal jump, Freivalds--Smith~\cite{FreivaldsSmith1993}).
\item $\forall \alpha < \beta$ constructive:
  $\Ex_\alpha \subsetneq \Ex_\beta$ \quad
  (\Cref{thm:ordinal-hierarchy}).
\item $\bigcup_{\alpha \in \mathcal{O}} \Ex_\alpha \subsetneq \Ex$
  \quad (Ambainis--Jain--Sharma~\cite{AmbainisJainSharma1999}).
\item $\Ex = \Ex_0^* \subsetneq \Ex_1^* \subsetneq \cdots \subsetneq
  \Ex^* \subsetneq \BC_0 \subsetneq \BC$ \quad
  (Case--Smith~\cite{CaseSmith1983}).
\end{enumerate}
\end{thm}


\subsection*{What the ordinal measures}

The mind-change ordinal of a class $\mathcal{C}$ is the least
constructive ordinal $\alpha$ such that $\mathcal{C} \in \Ex_\alpha$,
if such an ordinal exists.  If $\mathcal{C} \in \Ex \setminus
\bigcup_\alpha \Ex_\alpha$, then $\mathcal{C}$ has no ordinal
mind-change bound; its mind-change complexity exceeds all constructive
ordinals.

\begin{defn}[Mind-Change Ordinal of a Class]
\label{def:mco-class}
The \emph{mind-change ordinal} of a class $\mathcal{C}$ is
\[
  \mathrm{mco}(\mathcal{C}) \;=\;
  \min\bigl\{\alpha \in \mathcal{O} : \mathcal{C} \in \Ex_\alpha\bigr\}
\]
if the minimum exists, and $\infty$ otherwise.
\end{defn}

\begin{example}
\label{ex:pattern-languages}
Unions of $n$ pattern languages have mind-change complexity exactly
$\omega^n$ (Ambainis--Jain--Sharma~\cite{AmbainisJainSharma1999}, who
establish both the $\omega^n$ upper bound and a matching lower bound).
This provides a natural example of classes at each level $\omega^n$
in the hierarchy.
\end{example}

\begin{graphnode}[title={\nde{mind\_change\_ordinal}}]
\textbf{Graph node.} Category: \texttt{complexity\_measure}, Layer~5.\\
\textbf{Edges:}
$\nde{mind\_change\_ordinal}
  \xrightarrow{\rel{extends\_grammar}} \nde{mind\_change\_count}$
(new primitives required: constructive ordinals).\\
$\nde{mind\_change\_characterization}
  \xrightarrow{\rel{characterizes}} \nde{mind\_change\_ordinal}$
(Ambainis--Jain--Sharma 1999).\\
\textbf{Provenance:} Freivalds--Smith 1993.
\end{graphnode}

\begin{graphnode}[title={\nde{bc\_learning}}]
\textbf{Graph node.} Category: \texttt{success\_criterion}, Layer~4.\\
\textbf{Edges:}
$\nde{bc\_learning} \xrightarrow{\rel{restricts}} \nde{ex\_learning}$
(constraint removal: drops syntactic convergence requirement).\\
\textbf{Provenance:} Case--Smith 1983.
\end{graphnode}


\subsection{A Machine-Checked Ordinal Rank}
\label{sec:ordinal-rank-formal}

The mind-change ordinal is one ordinal-valued complexity measure, native
to recursion-theoretic identification.  The combinatorial dimensions of
\Cref{ch:dimensions} carry a second one, and that one has been formalized.

The Littlestone and VC dimensions take values in $\N \cup \{\infty\}$.
There is a machine-checked, order-preserving embedding of this value type
into the ordinals that fixes each finite dimension and sends the infinite
value to~$\omega$,\formal{FLT_Proofs.Bridge.withTopNatToOrdinal} and under
it the VC dimension is dominated by its ordinal lift,
$\mathrm{withTopNatToOrdinal}(\VC(\mathcal{C})) \leq
\mathrm{OrdinalVCDim}(\mathcal{C})$.\formal{FLT_Proofs.Bridge.vcdim_to_ordinal_vcdim}
This is the verified core of the ``ordinal VC dimension'' of Martin,
Sharma, and Stephan~\cite{MartinSharmaStephan2006}: a class of infinite VC
dimension sits at the first transfinite rank~$\omega$, finite-dimension
classes strictly below it.

\begin{prop}[Online mistakes descend an ordinal rank]
\label{prop:online-ordinal-descent}
In the online (Littlestone) model, the optimal mistake bound of a class
equals its Littlestone dimension.\formal{FLT_Proofs.Theorem.Online.optimal_mistake_bound_eq_ldim}
For a class of \emph{finite} Littlestone dimension, the Standard Optimal
Algorithm realizes this bound by treating the dimension as a rank: each
time it errs, the Littlestone dimension of the surviving version space
strictly decreases.  Lifted through the embedding above, this rank is a
finite ordinal, and the strict decrease on each mistake is precisely the
well-founded descent that \Cref{def:ordinal-mc} carries into the
transfinite.
\end{prop}

\begin{proof}[Proof idea]
The mistake-bound equality is an adversary argument matched against the
algorithm's potential.  Its engine is the machine-checked descent step:
a mistake on a version space of finite Littlestone dimension strictly
lowers that dimension, so iterating bounds the mistakes by the initial
rank.  The online counter never passes~$\omega$; it starts at a finite
rank (or, for an infinite-dimensional class, at~$\omega$) and counts down
by ordinary induction.  The transfinite levels $\Ex_{\omega^2},
\Ex_{\omega^\omega}, \ldots$ of this chapter are the price of
\emph{identifying programs} rather than \emph{predicting labels}; the
prediction game stays at the bottom of the ordinal tower.
\end{proof}

\begin{remark}[Attention sits at a finite rank]
\label{rem:attention-ordinal}
The softmax-attention concept class of \Cref{ch:generalization} is a
machine-checked well-behaved class,\formal{TLT.Bridge.softAttention_wellBehaved}
and \Cref{prop:certified-attention} gives it a finite Dudley capacity once
its parameters are norm-bounded.  Its ordinal rank is therefore finite:
attention lives at the very bottom of the ordinal picture, alongside
halfspaces and the other geometric classes.  The transfinite mind-change
hierarchy is invisible to such classes.  It becomes necessary only when
the objects learned are programs and the success criterion is exact
identification (the recursion-theoretic regime of this chapter), which no
finite-capacity concept class can reach.  Placing attention in the ordinal
landscape thus also marks where that landscape first becomes non-trivial.
\end{remark}


% ================================================================
% OPEN PROBLEMS
% ================================================================

\section{Open Problems}
\label{sec:ordinals-open}

The mind-change hierarchy is classically complete in its main lines, but
its formal status, its interaction with the verified online rank, and its
two-dimensional refinement remain open.  Each problem is tagged a
\emph{known unknown} (KU), sharply posed but unresolved, or an
\emph{unknown known} (UK), where the right object or invariant is still
missing; a design-lab or literature anchor is named where one exists.

\begin{openprob}[The mind-change/anomaly lattice]
\label{op:two-dim-lattice}
Does every point of the two-dimensional mind-change$\,\times\,$anomaly
lattice admit a natural characterization, as unions of $n$ pattern
languages characterize $\omega^n$?  A known unknown left open by Sharma,
Stephan, and Ventsov~\cite{SharmaStephanVentsov2004}.
\end{openprob}

\begin{openprob}[Formalizing the transfinite hierarchy]
\label{op:formalize-ordinal-hierarchy}
The dimension-to-ordinal bridge is machine-checked
(\leanname{withTopNatToOrdinal}), but the Freivalds--Smith hierarchy
itself (constructive ordinals, $\alpha$-bounded learners, the strict
inclusions of \Cref{thm:ordinal-hierarchy}) has no verified counterpart.
Can the ordinal mind-change hierarchy be formalized in a proof assistant,
including the priority constructions?  An unknown known: the recursion-
theoretic machinery has no formal home yet.
\end{openprob}

\begin{openprob}[Notation-invariant levels]
\label{op:notation-invariance}
\Cref{rem:notation-dependence} records that $\Ex_a$ can depend on the
notation~$a$, not just the ordinal~$|a|$.  Characterize exactly which
ordinals admit a notation-invariant $\Ex_\alpha$, and whether the
strictness theorem survives the quotient by notation.  A known
unknown~\cite{MartinSharmaStephan2006}.
\end{openprob}

\begin{openprob}[The ordinal rank of attention]
\label{op:attention-ordinal-rank}
\Cref{rem:attention-ordinal} places the attention class at a finite VC
rank, but its \emph{Littlestone} dimension (the online mind-change
rank) may still be infinite (thresholds already are).  Is the Littlestone
dimension of a softmax-attention class finite, and if so what is its exact
ordinal rank?  A known unknown on the transformer side
(\leanname{softAttention_wellBehaved}, \leanname{optimal_mistake_bound_eq_ldim}).
\end{openprob}

\begin{openprob}[A verified online-to-transfinite bridge]
\label{op:online-transfinite}
The online descent (each mistake strictly lowers the Littlestone
dimension) is machine-checked (\leanname{ldim_strict_decrease_on_mistake})
and stays below~$\omega$.  Is there a single formal framework in which this
verified finite descent and the transfinite mind-change counter are
instances of one well-founded-recursion principle?  An unknown known: the
unifying ordinal-recursion abstraction is not yet written down.
\end{openprob}

\begin{openprob}[Formalizing $\BC \setminus \Ex$]
\label{op:formalize-bc}
The separation of \Cref{thm:bc-separation} rests on the operator recursion
theorem, which we sketched rather than proved.  Can the Case--Smith
separation be machine-checked, operator recursion theorem
included~\cite{CaseSmith1983}?  A known unknown of formalization.
\end{openprob}

\begin{openprob}[Tightness of the Hessenberg bound]
\label{op:hessenberg-tight}
Exercise~3 gives $\mathcal{C}_1 \cup \mathcal{C}_2 \in
\Ex_{\alpha \oplus \beta}$ for the natural (Hessenberg) sum.  Is this bound
tight: are there classes whose union's mind-change ordinal is exactly
$\alpha \oplus \beta$ and no smaller?  A known unknown about resource
composition.
\end{openprob}

\begin{openprob}[Natural classes at each ordinal]
\label{op:natural-witnesses}
Pattern-language unions populate the levels $\omega^n$.  Is there a natural
family of concept classes populating the levels up to $\varepsilon_0$, or
does naturalness break down past $\omega^\omega$?  A known
unknown~\cite{AmbainisJainSharma1999}.
\end{openprob}

\begin{openprob}[An ordinal anomaly hierarchy]
\label{op:ordinal-anomaly}
The anomaly count~$a$ in $\Ex_a^*$ is a natural number.  Is there a
well-founded transfinite anomaly hierarchy (an $\Ex_\beta^*$ indexed by
ordinals), and does it interleave with the mind-change ordinals?  An
unknown known: the transfinite anomaly object has not been defined.
\end{openprob}

\begin{openprob}[Mind-change ordinals and ordinal VC dimension]
\label{op:mco-vs-ordinalvc}
The chapter's mind-change ordinal and the verified ordinal VC dimension
(\leanname{vcdim_to_ordinal_vcdim}) are two ordinal-valued measures on
overlapping domains.  For classes where both are defined, how are they
related: is one always a bound on the other?  A known
unknown~\cite{MartinSharmaStephan2006}.
\end{openprob}


% ================================================================
% EXERCISES
% ================================================================

\subsection*{Exercises}

\begin{enumerate}[leftmargin=*,itemsep=4pt]

\item \textbf{Finite mind-change hierarchy.}
  Prove directly (without ordinals) that $\Ex_n \subsetneq \Ex_{n+1}$
  for every $n \in \N$.  \emph{Hint:} Diagonalize against all learners
  with at most $n$ mind changes.

\item \textbf{$\omega$-counter mechanics.}
  Let $\mathcal{C} = \{f_n : n \in \N\}$ where $f_n$ is the
  characteristic function of $\{0, 1, \ldots, n\}$.  Show that
  $\mathcal{C} \in \Ex_\omega \setminus \bigcup_n \Ex_n$.

\item \textbf{Ordinal arithmetic of mind changes.}
  Show that if $\mathcal{C}_1 \in \Ex_\alpha$ and
  $\mathcal{C}_2 \in \Ex_\beta$, then
  $\mathcal{C}_1 \cup \mathcal{C}_2 \in \Ex_{\alpha \oplus \beta}$
  where $\oplus$ denotes natural (Hessenberg) sum.

\item \textbf{BC does not collapse.}
  Show that $\BC_0 \subsetneq \BC_1$: construct a class that is
  $\BC$-identifiable with at most one anomaly but not with zero
  anomalies.

\item \textbf{Anomaly vs.\ mind change.}
  Prove that $\Ex_1^* \not\subseteq \Ex_\omega$: tolerance for one
  anomaly cannot be simulated by any ordinal mind-change bound (without
  anomaly tolerance).  \emph{Hint:} Use a class where the anomalous
  input cannot be detected in finite time.

\item \textbf{Counter initialization matters.}
  Show that the learning power of $\Ex_\alpha$ depends on the
  \emph{notation} for $\alpha$, not just the ordinal.  That is,
  two different notations $a, a' \in \mathcal{O}$ with $|a| = |a'|$
  may yield different classes $\Ex_a \neq \Ex_{a'}$.
  \emph{Hint:} Different notations for $\omega$ give different
  computable sequences for decrementing.

\end{enumerate}

\subsection*{Notes and Further Reading}

The mind-change hierarchy was introduced by Freivalds and
Smith~\cite{FreivaldsSmith1993}, building on the finite mind-change
bounds of Velauthapillai~\cite{Velauthapillai1989} and the systematic
study of mind changes by Fulk, Jain, and
Osherson~\cite{FulkJainOsherson1995}.  The full characterization using
constructive ordinals, including the strictness of the hierarchy and
its non-exhaustion of $\Ex$, was completed by Ambainis, Jain, and
Sharma~\cite{AmbainisJainSharma1999}.

The four types of mind-change counters (constant, ordinal, linearly
ordered, partially ordered) were introduced by Sharma, Stephan, and
Ventsov~\cite{SharmaStephanVentsov2004}, who showed that these form a
strict hierarchy of learning power.  The connection between
mind-change ordinals and ordinal VC dimension was explored by Martin,
Sharma, and Stephan~\cite{MartinSharmaStephan2006}.

$\BC$-learning was introduced by Case and Smith~\cite{CaseSmith1983}.
The comprehensive treatment of learning criteria hierarchies is in
Jain, Osherson, Royer, and Sharma~\cite{JainOshersonRoyerSharma1999},
which remains the definitive reference for Gold-style learning theory.

The interaction between mind changes and anomalies produces a
two-dimensional lattice that is still not fully mapped.  Whether every
point in this lattice has a natural characterization (analogous to
pattern language unions for $\omega^n$) remains an open question.


\subsection*{Counterfactual exercises: generalizing Figure~\ref{fig:ordinal-hierarchy}}

Each exercise perturbs the tower.  Argue, in the figure's grammar (the number of mind changes or anomalies a limit learner needs is a constructive ordinal that stratifies the criteria into a strict tower from $\FIN = \Ex_0$ up through $\Ex$ and on into $\BC$; every strict inclusion is a classical theorem with no machine-checked counterpart, so every inclusion arrow is drawn dashed and cited to Velauthapillai, Freivalds--Smith, Ambainis--Jain--Sharma, or Case--Smith; the three base criteria $\FIN$, $\Ex$, $\BC$ are formalized and drawn solid, the ordinal and anomaly levels are literature definitions and drawn lighter; and the kernel's only verified transfinite content is a \emph{separate} ordinal rank on the dimension side, the lower box, that sends an infinite VC dimension to $\omega$ and places every finite-capacity class at a finite rank below it), that the perturbed picture subsumes the concrete one as an \emph{instance} (a criterion or rank placement the verified core or the literature exhibits), a \emph{boundary} (one inclusion or one node grade toggled so the evidence grade or the axis of relaxation changes), or a \emph{frontier} (a literature reading whose promotion to a verified rank is open).

\begin{enumerate}[leftmargin=*,itemsep=3pt]
\item \textbf{The base criteria are the figure's solid nodes.} Read the three boxes $\FIN$, $\Ex$, $\BC$ apart from the arrows between them. Show this is the figure's defining verified instance at the node level: these three are kernel definitions (\leanname{BCLearnable} and its $\FIN$ and $\Ex$ companions in the criterion library, the same base criteria formalized in \Cref{ch:gold}), so they are drawn with solid borders, whereas the levels $\Ex_\alpha$ (\cref{def:ex-alpha}) and $\Ex_a^*$ (\cref{def:anomalous-ord}) are literature definitions drawn lighter; the node grade is independent of the edge grade, and at the nodes alone the endpoints of the tower are certified while its rungs are not.

\item \textbf{The finite base levels are the integer instance of the ordinal spine.} Restrict the left column to $\FIN = \Ex_0 \subsetneq \Ex_1 \subsetneq \cdots \subsetneq \Ex_n$. Show this is the figure's most concrete instance of ``mind changes are an ordinal'': a learner bounded to $n$ revisions sits at the finite ordinal $n$, the strict inclusions hold (\cref{thm:summary}~(i), \cite{Velauthapillai1989}), and the arrows are dashed because even this integer fragment, the easiest in the tower, has no machine-checked counterpart; the whole transfinite picture is the closure of this fragment under limits.

\item \textbf{Attention sits at the bottom finite rank of the verified inset.} Read the boxed dimension-side rank, not the Gold tower. Show this is the instance the inset exists to display: the softmax-attention class is machine-checked well-behaved (\cref{rem:attention-ordinal}, \leanname{softAttention_wellBehaved}) and has finite Dudley capacity once norm-bounded (\cref{prop:certified-attention}), so its ordinal rank is finite and it lives at the very bottom of the picture beside halfspaces, while the embedding \leanname{withTopNatToOrdinal} reserves the first transfinite rank $\omega$ for an infinite VC dimension; the transfinite tower is invisible to every finite-capacity class.

\item \textbf{The online mistake count is the verified finite-rank shadow.} Read \cref{prop:online-ordinal-descent} as the inset's mechanism. Show this is the instance that certifies a descending rank: the optimal online mistake bound equals the Littlestone dimension (\leanname{optimal_mistake_bound_eq_ldim}), and each mistake strictly lowers the surviving version space's Littlestone dimension, a well-founded descent that stays below $\omega$; lifted through \leanname{withTopNatToOrdinal} this is a finite ordinal counted down by ordinary induction, the machine-checked finite shadow of the transfinite mind-change counter the tower draws above $\omega$.

\item \textbf{The verified rank and the mind-change ordinal are two measures, drawn apart on purpose.} Ask why the inset is fenced off from the left column rather than placed on it. Show this is the figure's principal boundary: the left column is the recursion-theoretic mind-change ordinal (\cref{def:mco-class}), a literature object, while the inset is the dimension-side ordinal VC rank (\leanname{vcdim_to_ordinal_vcdim}, \leanname{OrdinalVCDim}), the verified core of Martin--Sharma--Stephan \cite{MartinSharmaStephan2006}; they are two ordinal-valued measures on overlapping domains and whether one bounds the other is open (\cref{op:mco-vs-ordinalvc}), so collapsing them onto a single axis would assert a comparison the kernel does not certify.

\item \textbf{The first limit jump is the boundary between integers and ordinals.} Toggle the rung $\bigcup_{n<\omega}\Ex_n \subsetneq \Ex_\omega$. Show this is the boundary where the spine becomes genuinely transfinite: a learner may need finitely many mind changes on each input with no uniform integer bound, so the supremum is $\omega$ and no $\Ex_n$ suffices (\cref{rem:omega-to-n}, \cref{thm:summary}~(ii), \cite{FreivaldsSmith1993}); the arrow is dashed like every other, and it is exactly this jump that the integer hierarchy of Exercise~2 cannot reach, which is why the ordinal index is forced.

\item \textbf{The anomaly axis is an orthogonal relaxation, not a higher rung.} Turn from the mind-change column to the right column $\Ex \subsetneq \Ex_1^* \subsetneq \cdots \subsetneq \Ex^* \subsetneq \BC_0 \subsetneq \BC$. Show this is the boundary between two axes of relaxation: the mind-change column bounds the number of hypothesis revisions, while the anomaly column tolerates errors in the final hypothesis and then drops syntactic for semantic convergence (\cref{thm:anomaly-hierarchy}, \cref{thm:bc-separation}, \cite{CaseSmith1983}); both columns are dashed literature, and the two axes are independent, so the figure draws them as separate columns rather than one chain.

\item \textbf{The tower is classically complete but entirely unformalized.} Ask which arrow in the figure is solid. Show this is the dashed-versus-solid boundary the figure is built to make legible: none of the inclusion arrows is solid, because the Freivalds--Smith hierarchy, including its priority constructions, has no verified counterpart (\cref{op:formalize-ordinal-hierarchy}) and the Case--Smith separation rests on an operator recursion theorem that is sketched rather than machine-checked (\cref{op:formalize-bc}, \cite{CaseSmith1983}); the tower is a complete body of classical mathematics whose formalization is open, which is the precise content of drawing every rung dashed.

\item \textbf{Formalizing the tower would recolor the whole left column.} Ask which single advance turns the dashed rungs solid. Show this is the frontier the dashed arrows mark: a proof-assistant development of constructive ordinal notations, $\alpha$-bounded learners, and the strict inclusions of \cref{thm:ordinal-hierarchy} (\cref{op:formalize-ordinal-hierarchy}), together with the operator recursion theorem behind $\BC \setminus \Ex$ (\cref{op:formalize-bc}), would promote the literature tower to a verified one and make its arrows solid; until then the contrast in the figure is exactly the contrast between the solid dimension-side inset and the dashed recursion-theoretic spine.

\item \textbf{Is the mind-change ordinal the rank of a well-founded version-space descent functor whose finite shadow is the verified online descent?} Read the left column and the inset as two instances of one principle. Show this is the figure's deepest frontier: the verified online descent lowers a finite Littlestone rank on each mistake (the mechanism behind \leanname{optimal_mistake_bound_eq_ldim}, \cref{prop:online-ordinal-descent}) and stays below $\omega$, while the mind-change counter descends a constructive ordinal that may exceed every $\omega^n$; whether both are instances of a single well-founded-recursion principle, a version-space descent functor whose ordinal rank is the mind-change ordinal and whose finite-rank restriction is the certified online descent, is open (\cref{op:online-transfinite}), as is whether every point of the two-dimensional mind-change-by-anomaly lattice admits such a rank (\cref{op:two-dim-lattice}, \cite{SharmaStephanVentsov2004}); proving and formalizing it would place the dashed tower and the solid inset under one verified abstraction, which is why this reading is drawn dashed.
\end{enumerate}

% Chapter 14: What Does Not Imply What
% This chapter is witness-centered: each separation lives or dies by
% the construction that demonstrates it.
% Content: ALL 9 does_not_imply edges, ALL 4 strictly_stronger edges.

\chapter{What Does Not Imply What}
\label{ch:separations}

Most textbooks in learning theory treat separation results as scattered
remarks: brief asides after a characterization theorem, noting that
some plausible implication fails. This chapter reverses that convention.
Here, the separations are first-class citizens: each one receives a
formal statement, a witness construction, and an analysis of what
structural feature the witness exploits.

A separation result has two components. The \emph{statement} asserts
that some implication $A \Rightarrow B$ does not hold. The
\emph{witness} is a concrete mathematical object (a concept class,
a dimension pair, a computational reduction) that demonstrates the
failure. The witness is the mathematics; the statement is merely its
summary. Throughout this chapter, we privilege the construction over the
claim.

We organize the chapter's 13 edges into two groups: 9
\rel{does\_not\_imply} edges, where the non-implication is the content,
and 4 \rel{strictly\_stronger} edges, where an implication does hold but
is provably non-reversible. Together they form the \emph{separation
lattice} of formal learning theory.

% ================================================================
% SECTION 1: THE SEPARATION LATTICE
% ================================================================

\section{The Separation Lattice}
\label{sec:sep-lattice}

\Cref{fig:separation-lattice} displays all 13 edges as a single
diagram. Dashed red arrows denote \rel{does\_not\_imply} edges: the
source does \emph{not} entail the target, and the label names the
witness. Solid blue arrows denote \rel{strictly\_stronger} edges: the
source strictly contains the target as a special case, with the witness
demonstrating the gap.

The first observation is structural: the lattice is \emph{sparse}.
Thirteen edges connect concepts drawn from six paradigms and a dozen
complexity measures. Most paradigm pairs are simply incomparable: they
neither imply nor contradict each other, because they operate on
different mathematical objects. The sparsity is itself informative:
learning theory is not a linear hierarchy from weak to strong, but a
partially ordered collection of largely independent formalisms.

\begin{figure}[ht]
\centering
\begin{tikzpicture}[
  node distance=1.8cm and 2.4cm,
  concept/.style={
    rectangle, rounded corners=3pt, draw, thick,
    fill=layergray, font=\small\sffamily,
    minimum width=2.2cm, minimum height=0.7cm,
    align=center
  },
  dni/.style={-{Stealth[length=5pt]}, dashed, thmred, thick},
  ss/.style={-{Stealth[length=5pt]}, defblue, thick},
  witness/.style={font=\tiny\itshape, text=notegray, align=center}
]

% --- Nodes ---
\node[concept] (ul) {Universal\\Learning};
\node[concept, below left=2.0cm and 1.2cm of ul] (ol) {Online\\Learning};
\node[concept, below right=2.0cm and 1.2cm of ul] (pac) {PAC\\Learning};
\node[concept, right=2.8cm of pac] (exact) {Exact\\Learning};
\node[concept, below=1.8cm of ol] (mb) {Mistake\\Bounded};
\node[concept, below=1.8cm of pac] (ex) {Ex-\\Learning};
\node[concept, below=1.8cm of ex] (fin) {FIN-\\Learning};

\node[concept, left=2.5cm of ol] (vc) {VC\\Dimension};
\node[concept, below=1.8cm of vc] (sq) {SQ\\Dimension};
\node[concept, below left=1.8cm and 0.0cm of vc] (comp) {Comput.\\Hardness};

\node[concept, right=2.8cm of exact] (nd) {Natarajan\\Dimension};
\node[concept, above=1.8cm of nd] (ds) {DS\\Dimension};

\node[concept, below=1.8cm of mb] (ft) {Fundamental\\Theorem};
\node[concept, right=2.8cm of ft] (stab) {Algorithmic\\Stability};

\node[concept, below left=1.4cm and 0.5cm of pac] (pi) {Proper/Improper\\Separation};
\node[concept, below right=1.8cm and 0.0cm of exact] (uc) {Unlabeled\\Compression};
\node[concept, below=1.5cm of uc] (cs) {Compression\\Scheme};

% --- does_not_imply edges (dashed red) ---
\draw[dni] (pac) -- node[witness, above, sloped] {thresholds} (mb);
\draw[dni] (nd) -- node[witness, above, sloped] {$d_N{=}1$, not learnable} (pac);
\draw[dni] (pac) -- node[witness, above, sloped] {approx.\ vs.\ exact} (exact);
\draw[dni] (vc) -- node[witness, left, pos=0.4] {crypto} (comp);
\draw[dni] (vc) -- node[witness, left] {parities} (sq);
\draw[dni] (ex) -- node[witness, above, sloped] {$\VC = \infty$} (pac);
\draw[dni] (pi) -- node[witness, below, sloped] {3-DNF} (pac);
\draw[dni] (uc) -- node[witness, right, pos=0.4] {P\'alv\"olgyi--Tardos} (cs);
\draw[dni] (ft) -- node[witness, above, sloped] {beyond binary} (stab);

% --- strictly_stronger edges (solid blue) ---
\draw[ss] (ul) -- node[witness, above left, sloped] {removes uniformity} (pac);
\draw[ss] (ol) -- node[witness, above, sloped] {$\Ldim \Rightarrow \VC$} (pac);
\draw[ss] (ex) -- node[witness, right] {unbounded mind changes} (fin);
\draw[ss] (ds) -- node[witness, right] {pseudo-manifold} (nd);

\end{tikzpicture}

\caption{The separation lattice. Dashed red:
\rel{does\_not\_imply} (9~edges). Solid blue:
\rel{strictly\_stronger} (4~edges). Each label names the witness.}
\label{fig:separation-lattice}
\end{figure}


% ================================================================
% SECTION 2: SEPARATIONS BETWEEN PARADIGMS
% ================================================================

\section{Separations Between Paradigms}
\label{sec:separations}

We now present the 9 \rel{does\_not\_imply} edges in the graph, each
with its witness construction. The order moves from the most elementary
witness (a one-parameter family on $\R$) to the most conceptually
involved (the breakdown of the fundamental theorem beyond binary
classification).

% --- Separation 1 ---

\begin{separation}
\textbf{PAC learning $\not\Rightarrow$ mistake-bounded learning.}\formal{FLT_Proofs.Theorem.Separation.pac_not_implies_online}

\medskip\noindent
\emph{Witness.}
Let $\mathcal{H} = \{h_\theta : \theta \in \R\}$ be the class of
thresholds on $\R$, where $h_\theta(x) = \ind[x \geq \theta]$.

This class has $\VC(\mathcal{H}) = 1$: a single point $x$ is shattered
(choose $\theta < x$ or $\theta > x$), but no two points can be
simultaneously shattered.\formal{FLT_Proofs.Theorem.Separation.vcdim_threshold_finite} By the fundamental theorem of statistical
learning, $\mathcal{H}$ is PAC learnable with sample complexity
$O(1/\varepsilon)$.

However, $\Ldim(\mathcal{H}) = \infty$.\formal{FLT_Proofs.Theorem.Separation.ldim_threshold_top} An adaptive adversary can force
arbitrarily many mistakes by binary search: maintain an interval
$(a, b)$ of uncertainty, present the midpoint, and whichever label the
learner predicts, place the true threshold on the opposite side. Each
round forces a mistake and halves the interval, but the reals admit
no finite termination.

\medskip\noindent
\emph{What the witness exploits.}
The gap between VC and Littlestone dimension: batch learnability
(distribution-free, i.i.d.\ samples) does not imply sequential
learnability (adversarial, adaptive instances). The reals are
``too dense'' for any online strategy to pin down the threshold,
but a random sample reveals it with high probability.

\hfill\cite{Littlestone1988}
\end{separation}

% --- Separation 2 ---

\begin{separation}
\textbf{Ex-learning $\not\Rightarrow$ PAC learning.}\formal{FLT_Proofs.Theorem.Separation.ex_not_implies_pac}

\medskip\noindent
\emph{Witness.}
Let $\mathcal{C} = \{ S \subseteq \N : S \text{ is finite} \}$, the
class of all finite subsets of $\N$. A Gold-style learner can identify
any finite $S$ in the limit from positive data: enumerate elements as
they appear, and at each stage conjecture the set of elements seen so
far. After the last new element arrives, the conjecture stabilizes on
$S$.

However, $\VC(\mathcal{C}) = \infty$: for any $n$ points
$\{x_1, \ldots, x_n\} \subset \N$, every subset is itself a finite set
in $\mathcal{C}$, so the class shatters every finite set. The
fundamental theorem then implies $\mathcal{C}$ is not PAC learnable.

\medskip\noindent
\emph{What the witness exploits.}
The data models are incompatible. Gold-style identification receives an
infinite enumeration of the target's elements and succeeds in the limit;
PAC learning receives a finite i.i.d.\ sample and must generalize
immediately. A class can be identifiable from exhaustive enumeration yet
have no finite VC dimension.

\hfill\cite{Gold1967}
\end{separation}

% --- Separation 3 ---

\begin{separation}
\textbf{PAC learning $\not\Rightarrow$ exact learning.}

\medskip\noindent
\emph{Witness.}
Exact learning in the sense of Angluin~\cite{Angluin1988} identifies the
target \emph{exactly}, with zero error, from membership and equivalence
queries.  It departs from PAC learning along \emph{two} axes at once: the
success criterion (zero error versus $\varepsilon$-approximate) and the
data model (active queries versus a passive i.i.d.\ sample).  A
PAC-learnable class need not be exactly learnable on either count.

The criterion axis alone already bites computationally. Consider a class
$\mathcal{C}$ that is PAC learnable (finite VC dimension) but for which
finding a hypothesis \emph{exactly} consistent in the proper representation
is NP-hard (intersections of halfspaces in $\R^d$ are the standard
example~\cite{PittValiant1988}). An improper PAC learner attains low error
through a surrogate class, yet exact identification of the target is
intractable.

\medskip\noindent
\emph{What the witness exploits.}
Two mismatches stacked: approximate-versus-exact in the criterion, and
i.i.d.-sample-versus-query in the data model. PAC tolerates
$\varepsilon$~error and reads a passive sample; exact learning tolerates
none and may interrogate the target. Neither model subsumes the other.

\hfill\cite{Angluin1988, PittValiant1988}
\end{separation}

% --- Separation 4 ---

\begin{separation}
\textbf{Finite VC dimension $\not\Rightarrow$ computational
tractability.}

\medskip\noindent
\emph{Witness.}
Kearns and Valiant~\cite{KearnsValiant1994} constructed concept classes
based on polynomial-size Boolean circuits whose VC dimension is
polynomial in the circuit size, yet for which PAC learning is
computationally intractable under the assumption that the
\emph{decisional Diffie--Hellman} (or related cryptographic) assumption
holds.

The construction proceeds as follows. Fix a one-way function family.
Define $\mathcal{C}$ to be the class of functions computed by circuits
of size $s$. This class has $\VC(\mathcal{C}) \leq O(s \log s)$, so it
is information-theoretically PAC learnable with $\mathrm{poly}(s)$
samples. But any polynomial-time PAC learner for $\mathcal{C}$ could be
used to invert the one-way function, contradicting the cryptographic
assumption.

\medskip\noindent
\emph{What the witness exploits.}
The fundamental theorem is \emph{information-theoretic}: it
characterizes learnability in terms of sample complexity, not
computational complexity. Cryptographic hardness lives in a different
layer entirely. Finite VC dimension guarantees that enough data suffices;
it says nothing about whether a polynomial-time algorithm can find a
good hypothesis from that data.

\hfill\cite{KearnsValiant1994}
\end{separation}

% --- Separation 5 ---

\begin{separation}
\textbf{Finite VC dimension $\not\Rightarrow$ low SQ dimension.}

\medskip\noindent
\emph{Witness.}
The class of parities over $\{0,1\}^n$. Each parity function $\chi_S$,
for $S \subseteq [n]$, maps $x \mapsto \bigoplus_{i \in S} x_i$. The
parity class has $\VC = n$ (it shatters the standard basis). But any
statistical query algorithm requires queries of tolerance
$\tau = O(2^{-n/2})$ or uses $2^{\Omega(n)}$ queries of polynomial
tolerance, because distinct parities are pairwise uncorrelated under the
uniform distribution.

Formally, $\SQdim(\text{Parities}_n) = 2^n$, which is exponential in
$\VC = n$.

\medskip\noindent
\emph{What the witness exploits.}
VC dimension measures combinatorial shattering capacity
(distribution-free). SQ dimension measures pairwise correlation
structure under a specific distribution. Parities are maximally
uncorrelated under the uniform distribution, forcing any SQ learner to
make exponentially many queries, even though the class is small in the
VC sense.

\hfill\cite{BlumFurstJacksonKearnsMansourRudich1994}
\end{separation}

% --- Separation 6 ---

\begin{separation}
\textbf{Natarajan dimension $\not\Rightarrow$ multiclass PAC
learnability.}

\medskip\noindent
\emph{Witness.}
Brukhim et al.~\cite{BrukhimCarmonDinurMoranYehudayoff2022}
constructed a concept class $\mathcal{C}$ with label space
$Y$ of infinite cardinality, Natarajan dimension $\Ndim(\mathcal{C})=1$,
yet $\mathcal{C}$ is not PAC learnable.

The Natarajan dimension, which
generalizes VC dimension to multiclass settings by requiring two
distinct labelings of shattered points, is too coarse when $|Y|$ is
infinite. The construction builds a class in which every pair of points
can be 2-colored in only one way (so $\Ndim = 1$), but the global
combinatorial structure is rich enough to prevent uniform convergence.

\medskip\noindent
\emph{What the witness exploits.}
The Natarajan dimension captures pairwise shattering. When the label
space is infinite, pairwise structure does not determine global
learnability. The DS dimension, which accounts for higher-order
combinatorial structure via oriented hypergraphs, is the correct
characterization.

\hfill\cite{BrukhimCarmonDinurMoranYehudayoff2022}
\end{separation}

% --- Separation 7 ---

\begin{separation}
\textbf{Proper learnability $\not\Rightarrow$ efficient proper PAC
learning.}

\medskip\noindent
\emph{Witness.}
The class of 3-term DNF formulas over $\{0,1\}^n$. Pitt and
Valiant~\cite{PittValiant1988} showed that \emph{proper} PAC
learning\formal{FLT_Proofs.Basic.IsProper} (where the hypothesis must
itself be a 3-term DNF) is NP-hard, assuming
$\mathrm{RP} \neq \mathrm{NP}$. Yet the same class
is efficiently PAC learnable \emph{improperly}: every 3-term DNF can be
represented as a 3-CNF, and 3-CNFs are efficiently PAC learnable.

\medskip\noindent
\emph{What the witness exploits.}
The gap between proper and improper learning is purely
computational: the information-theoretic sample complexity is
identical, but the representation constraint of proper learning
introduces NP-hardness. The concept class is simple enough to learn
with the right representation, but finding a hypothesis in the
\emph{same} representation class is intractable.

\hfill\cite{PittValiant1988}
\end{separation}

% --- Separation 8 ---

\begin{separation}
\textbf{Labeled compression $\not\Rightarrow$ unlabeled compression.}

\medskip\noindent
\emph{Witness.}
P\'alv\"olgyi and Tardos~\cite{PalvolgyiTardos2020} exhibited a
concept class with $\VC = 2$ that admits a labeled sample compression
scheme of size~2 but does \emph{not} admit an unlabeled compression
scheme of size~2.

In an unlabeled compression scheme, the compression function stores only
the \emph{points} (not their labels) from the training sample; the
reconstruction function must infer both the hypothesis and the labels
from the point set alone. The witness class is constructed so that the
label information is essential for reconstruction: two subsets of the
same size can correspond to different concepts, and only the labels
disambiguate them.

\medskip\noindent
\emph{What the witness exploits.}
The label information carries entropy that cannot always be recovered
from the geometry of the point set alone. Unlabeled compression demands
that the point configuration determines the concept, which is a strictly
stronger requirement than labeled compression.

\hfill\cite{PalvolgyiTardos2020}
\end{separation}

% --- Separation 9 ---

\begin{separation}
\textbf{The fundamental theorem $\not\Rightarrow$ stability
characterization.}

\medskip\noindent
\emph{Witness.}
The Fundamental Theorem of Statistical Learning (\Cref{ch:pac}) ties
together VC dimension, uniform convergence, and PAC learnability into a
single equivalence.  Nine conditions, all the same condition.  That
theorem is the crown jewel of binary classification.

It does not survive contact with real-valued loss.

Shalev-Shwartz et al.~\cite{ShalevShwartzShamir2010} construct a
learning problem (arbitrary loss, multi-valued output) that is
learnable, yet for which uniform convergence fails outright.  The
nine-way equivalence collapses.  In its place they prove a different
characterization: in the general learning setting, a problem is learnable
if and only if it admits a \emph{stable}, asymptotically
empirical-risk-minimizing learner.  Uniform convergence is demoted from
characterization to a sufficient condition.

\medskip\noindent
\emph{What the witness exploits.}
The symmetrization argument at the heart of the Fundamental Theorem's
proof requires the loss to take finitely many values.  Specifically:
two.  That assumption is invisible in the binary setting: it looks like
a harmless feature of the problem.  Make the loss real-valued and
symmetrization breaks.  Uniform convergence stops characterizing
learnability; the stable-AERM condition takes over as the
characterization, indifferent to the cardinality of the loss.

The Fundamental Theorem is not wrong.  It is \emph{local}: an artifact
of the 0-1 loss that does not announce itself as such.

\hfill\cite{ShalevShwartzShamir2010}
\end{separation}


% ================================================================
% SECTION 3: THE STRICT STRENGTH HIERARCHY
% ================================================================

\section{Strict Strength Hierarchy}
\label{sec:strict-hierarchy}

The 4 \rel{strictly\_stronger} edges assert that one concept
\emph{does} imply another, but the reverse fails: the stronger concept
is a proper generalization. Each edge requires two proofs: one for the
forward implication, and one (the witness) for the strictness of the
gap.

% --- Strict 1 ---

\begin{separation}
\textbf{Ex-learning $\supsetneq$ FIN-learning.}\formal{FLT_Proofs.Theorem.Separation.ex_strictly_stronger_finite}

\medskip\noindent
\emph{Forward implication.} Every FIN-learnable class is trivially
Ex-learnable: a learner that outputs its final hypothesis after finitely
many data points and never changes it again is, a fortiori, a learner
that converges in the limit.

\medskip\noindent
\emph{Witness for strictness.}
FIN-learning requires zero mind changes: the learner must output its
final, correct hypothesis at some point and never retract it. Ex-learning
allows unbounded mind changes before convergence. The gap is witnessed
by any class requiring unbounded mind changes for identification.

Consider the class of pattern languages with growing structural
complexity. A Gold-style learner can identify any member in the limit
(Ex-learn it) by successively refining its conjecture as new data
arrives, but no learner can commit to a correct hypothesis after seeing
only finitely many initial elements without ever revising. The revision
process is essential, and FIN-learning forbids it.

In the Gold identification hierarchy, this separation is the first step:
$\FIN \subsetneq \Ex$, and the gap is measured by the mind-change
ordinal. FIN corresponds to $0$ mind changes; Ex allows $< \omega$ mind
changes (any finite number, not bounded in advance).

\hfill\cite{Gold1967}
\end{separation}

% --- Strict 2 ---

\begin{separation}
\textbf{Online learning $\supsetneq$ PAC learning.}\formal{FLT_Proofs.Theorem.Separation.online_strictly_stronger_pac}

\medskip\noindent
\emph{Forward implication.}
If $\mathcal{H}$ has finite Littlestone dimension $\Ldim(\mathcal{H})
= d$, then in particular $\VC(\mathcal{H}) \leq d$.\formal{FLT_Proofs.Theorem.Separation.vcdim_le_of_mistake_bounded} (Every shattered
set in the VC sense extends to a path in a Littlestone tree, so a
mistake-bounded learner with bound $d$ forces $\VC \leq d$.) By the
fundamental theorem, $\mathcal{H}$ is PAC learnable.

\medskip\noindent
\emph{Witness for strictness.}
Thresholds on finite domains provide the simplest witness. Fix
$X = \{1, 2, \ldots, n\}$ and let $\mathcal{H}_n = \{h_\theta :
\theta \in \{0, 1, \ldots, n\}\}$ where $h_\theta(x) = \ind[x >
\theta]$. Then $\VC(\mathcal{H}_n) = 1$ for every $n$, but
$\Ldim(\mathcal{H}_n) = \lfloor \log_2(n+1) \rfloor$, which grows
without bound as $n \to \infty$.

More dramatically, thresholds on $\R$ (the witness from Separation~1)
have $\VC = 1$ but $\Ldim = \infty$: PAC learnable but not online
learnable with any finite mistake bound. The containment is strict at
every level of the hierarchy.

\hfill\cite{Littlestone1988}
\end{separation}

% --- Strict 3 ---

\begin{separation}
\textbf{Universal learning $\supsetneq$ PAC learning.}

\medskip\noindent
\emph{Forward implication.}
Universal learnability is the \emph{weaker} requirement: it asks only for a
learner whose error decays at some rate under each distribution
\emph{separately}, whereas PAC learning demands a single sample-complexity
bound uniform over all distributions. Removing the uniformity, every
PAC-learnable class (finite VC dimension) is therefore universally
learnable, indeed at an exponential rate. The strictness below shows the
converse fails.

\medskip\noindent
\emph{Witness for strictness.}
Bousquet et al.~\cite{BousquetHannekeMoranVanHandelYehudayoff2021}
established a trichotomy for universal learning rates: every concept
class has either an exponential rate, an arbitrarily slow rate, or is
not universally learnable at all. The trichotomy theorem shows that
classes with $\VC(\mathcal{H}) = \infty$ but containing no infinite
Littlestone tree achieve exponential universal learning rates; they
are universally learnable but \emph{not} PAC learnable (since PAC
requires finite VC dimension).

The critical structural difference is the quantifier order.
PAC learning demands: $\exists$ learner $\forall$ distributions
$\forall \varepsilon, \delta$, the learner succeeds. Universal learning
demands: $\forall$ distributions $\exists$ rate at which \emph{some}
learner succeeds. By removing the uniformity requirement over
distributions, universal learning captures a strictly larger class
of learning problems.

\hfill\cite{BousquetHannekeMoranVanHandelYehudayoff2021}
\end{separation}

% --- Strict 4 ---

\begin{separation}
\textbf{DS dimension $\supsetneq$ Natarajan dimension.}

\medskip\noindent
\emph{Forward implication.}
The DS dimension refines the Natarajan dimension: $\Ndim(\mathcal{H})
\leq \DSdim(\mathcal{H})$ for every hypothesis class $\mathcal{H}$.
Both measure the combinatorial complexity of multiclass concept classes,
but the DS dimension accounts for the full oriented hypergraph structure,
not just pairwise shattering.

\medskip\noindent
\emph{Witness for strictness.}
Brukhim et al.~\cite{BrukhimCarmonDinurMoranYehudayoff2022} constructed
a concept class using a hyperbolic pseudo-manifold with $\Ndim = 1$
but $\DSdim = \infty$. The construction builds a concept class over an
infinite label space where every pair of points admits only one
non-trivial two-coloring (giving $\Ndim = 1$), but the global
combinatorial structure (encoded in the oriented faces of the
pseudo-manifold) is rich enough to drive the DS dimension to infinity.

This is the witness that simultaneously demonstrates the separation
$\Ndim \not\Rightarrow \text{PAC learnable}$ from \Cref{sec:separations}:
the same class has $\Ndim = 1$ but is not learnable, because learnability
is characterized by the DS dimension, not the Natarajan dimension.

\hfill\cite{BrukhimCarmonDinurMoranYehudayoff2022}
\end{separation}


% ================================================================
% SECTION 4: WHAT THE NEGATIVE LAYER REVEALS
% ================================================================

\section{What the Negative Layer Reveals}
\label{sec:neg-layer-summary}

The separation lattice of \Cref{fig:separation-lattice} carries a
meta-theorem about the structure of formal learning theory, which we
can now state explicitly.

\begin{remark}[The separation lattice is sparse]
\label{thm:sparsity}
This chapter establishes thirteen edges with explicit witnesses among the
major learning paradigms and complexity measures of
\Cref{fig:separation-lattice}, far fewer than the $\binom{k}{2}$ relations
that $k$ such concepts could \emph{a priori} exhibit.  We do not claim a
complete census of the lattice; the point is that the witnessed edges are
scarce, and that many remaining pairs appear unrelated: they operate on
different mathematical types, or are connected only by non-implicational
relations (analogy, measurement, assumption) rather than by an implication
that could be witnessed or refuted.
\end{remark}

This scarcity is not merely an artifact of incomplete knowledge.  Many of
the major paradigms operate on different mathematical types: PAC learning
and Gold-style identification take different data models (i.i.d.\ samples
versus infinite enumerations); online and exact learning use different
interaction protocols (adversarial instances versus query access); VC and
SQ dimension measure different properties (combinatorial shattering versus
correlation structure).  Between such concepts a clean implication is often
not even the right question.  Where we \emph{can} ask it, this chapter
witnesses one direction (Gold identification does not imply PAC learning,
say), while establishing the reverse non-implication is, in several cases, a
separate open problem rather than a settled fact.

The witnesses in this chapter make the incomparability concrete.
Thresholds on $\R$ separate PAC from online learning. All finite subsets
of $\N$ separate Gold identification from PAC learning. Parities
separate VC dimension from SQ dimension. Each witness exploits a specific
structural mismatch in the data model, the success criterion, the
adversarial model, or the computational model; and these mismatches are
not removable by clever proof techniques. They are features of the
mathematical landscape.

The four strict strength edges, taken together, form two chains:
\[
\text{Universal} \supsetneq \text{Online} \supsetneq \text{PAC}
\qquad \text{and} \qquad
\DSdim \supsetneq \Ndim.
\]
The first chain orders three paradigms by the stringency of their
learnability requirements. The second orders two dimensions by their
sensitivity to multiclass structure. In both cases, the strict
containment is proved by a single, explicit construction.

These are the theorems that textbooks usually omit. They deserve
their chapter.


\subsection{Where Attention Sits}
\label{sec:attention-lattice}

A separation lattice is a map, and it is natural to ask where a modern
hypothesis class lands on it. The softmax-attention concept class of
\Cref{ch:generalization} is machine-checked to be
well-behaved,\formal{TLT.Bridge.softAttention_wellBehaved} and with
norm-bounded parameters it has finite capacity
(\Cref{prop:certified-attention}); it therefore sits on the PAC node: it
is PAC learnable.

Its position relative to the \emph{online} node is open.  The
witness for Separation~1 was a class of finite VC dimension and infinite
Littlestone dimension (thresholds on~$\R$).  Attention compares query--key
scores, and threshold behavior on a score is exactly the structure that
sends the Littlestone dimension to infinity while the VC dimension stays
finite.  It is therefore possible that the softmax-attention class is its
\emph{own} witness for $\text{PAC} \not\Rightarrow \text{online}$, joining
thresholds on that dashed red edge.  Whether it is so is
\Cref{op:attention-littlestone}; the machine-checked
separation~\leanname{pac_not_implies_online} already supplies the template
a proof would instantiate.


% ================================================================
% OPEN PROBLEMS
% ================================================================

\section{Open Problems}
\label{sec:separations-open}

A separation chapter generates two kinds of open problem: the missing
\emph{reverse} of a one-sided non-implication, and the missing
\emph{formal proof} of a separation so far only cited.  Each problem below
is tagged a \emph{known unknown} (KU), sharply posed but unresolved, or an
\emph{unknown known} (UK), where the right object or invariant is still
missing; a design-lab or literature anchor is named where one exists.

\begin{openprob}[The Littlestone dimension of attention]
\label{op:attention-littlestone}
Is the Littlestone dimension of the softmax-attention concept class finite
or infinite?  If infinite, attention is its own witness for
$\text{PAC} \not\Rightarrow \text{online}$
(\leanname{pac_not_implies_online}); if finite, attention is online
learnable and the score comparisons are tamer than thresholds.  A known
unknown on the transformer side (\leanname{softAttention_wellBehaved}).
\end{openprob}

\begin{openprob}[Does PAC imply Ex?]
\label{op:reverse-pac-ex}
The chapter proves Gold identification does not imply PAC learning
(\leanname{ex_not_implies_pac}), but never the reverse.  Is there a
PAC-learnable concept class that is not Ex-learnable as a class of
functions, so that PAC and Gold are genuinely incomparable rather than
one-sidedly separated?  A known unknown.
\end{openprob}

\begin{openprob}[Formalizing the DS/Natarajan gap]
\label{op:formalize-ds-natarajan}
The pseudo-manifold witness with $\Ndim = 1$ and $\DSdim = \infty$ is cited
from Brukhim et al., not derived.  Can the construction, and the
non-learnability it implies, be machine-checked?  An unknown known: the
oriented-hypergraph combinatorics has no formal
home~\cite{BrukhimCarmonDinurMoranYehudayoff2022}.
\end{openprob}

\begin{openprob}[Formalizing the universal-learning trichotomy]
\label{op:formalize-trichotomy}
The forward edge universal$\,\Rightarrow\,$PAC is machine-checked, but the
trichotomy that witnesses Strict~3 (exponential, arbitrarily slow, or not
at all) is not.  Can the Bousquet et al.\ trichotomy be formalized,
closing the strictness direction?  An unknown
known~\cite{BousquetHannekeMoranVanHandelYehudayoff2021}.
\end{openprob}

\begin{openprob}[VC versus SQ, quantitatively]
\label{op:sq-vc-quantitative}
Parities give $\VC = n$, $\SQdim = 2^n$ under the uniform distribution.
Across all distributions, what is the tightest possible relationship
between VC and SQ dimension: is the exponential gap always achievable, or
distribution-dependent?  A known
unknown~\cite{BlumFurstJacksonKearnsMansourRudich1994}.
\end{openprob}

\begin{openprob}[Formalizing labeled versus unlabeled compression]
\label{op:formalize-compression-sep}
The P\'alv\"olgyi--Tardos $\VC = 2$ class separating labeled from unlabeled
size-2 compression is cited, not constructed here.  Can the separation be
machine-checked, and does it persist at larger VC dimension and scheme
size?  An unknown known~\cite{PalvolgyiTardos2020}.
\end{openprob}

\begin{openprob}[Is the lattice complete?]
\label{op:lattice-completeness}
For each pair of paradigms with no witnessed edge, is the pair provably
incomparable, or merely unstudied?  A census distinguishing genuine
incomparability from undiscovered implications would turn the informal
sparsity of \Cref{thm:sparsity} into a theorem.  A known unknown.
\end{openprob}

\begin{openprob}[Characterizing the proper/improper gap]
\label{op:proper-improper-sep}
3-DNF has an efficient improper learner but no efficient proper one
(\leanname{IsProper}, under $\mathrm{RP}\neq\mathrm{NP}$).  Which structural
property of a class forces a proper/improper efficiency gap, and which
classes have none?  A known unknown~\cite{PittValiant1988}.
\end{openprob}

\begin{openprob}[The scope of the stability characterization]
\label{op:stability-scope}
Stable AERM characterizes learnability in the general learning setting
(\Cref{sec:separations}, Separation~9).  Exactly which loss families and
hypothesis structures fall inside that characterization, and where does
even stability cease to characterize?  A known
unknown~\cite{ShalevShwartzShamir2010}.
\end{openprob}

\begin{openprob}[The SQ dimension of attention]
\label{op:attention-sq}
Parities resist statistical-query learning despite small VC.  Does the
softmax-attention class behave like a VC class under SQ, or can it encode
parity-like correlations that make it SQ-hard?  An unknown known on the
transformer side (\leanname{softAttention_wellBehaved}).
\end{openprob}

% Chapter 15: Analogies and Their Obstructions
% Organized by OBSTRUCTION TYPE, not by analogy source.
% Content: ALL 32 analogy edges from the graph.

\chapter{Analogies and Their Obstructions}
\label{ch:analogies}

When you read that ``compression is like VC dimension'' or that ``stability
is related to PAC-Bayes,'' you are being handed an analogy without its
expiry date. The remark could mean that both quantities measure the same
underlying thing; it could mean that one bounds the other in one direction
only; or it could mean that the objects involved belong to different parts
of the type system and no formal connection will ever hold. Without knowing
which, you cannot tell whether to try proving the equivalence, whether to
expect a counterexample, or whether to search for a new bridging concept
entirely. The analogy, stated informally, is a map with the legend missing.

This chapter supplies the legend. The concept graph contains 32 edges of
type \rel{analogy}, and each is treated here as a formal object with three
components:

\begin{enumerate}[leftmargin=*, nosep]
\item A \emph{source} and \emph{target}: two concepts that share a
structural parallel.
\item A \emph{shared structure}: the precise mathematical property that
makes the analogy plausible.
\item An \emph{obstruction}: the precise reason the analogy fails to be
a theorem, classified by \emph{type}.
\end{enumerate}

The obstruction type carries more information than the analogy itself. Two
analogies can look identical at the surface level (``X is like Y because
both measure complexity'') while failing for entirely different reasons:
one because the objects live in different type systems, another because the
conjectured equivalence is plausible but still unproved. The type identifies the reason the analogy fails, and that
determines whether the gap is a research target, a structural impossibility,
or something in between.

The 32 analogy edges distribute across six obstruction types:

\begin{center}\small
\begin{tabular}{llr}
\toprule
\textbf{Obstruction type} & \textbf{Meaning} & \textbf{Count} \\
\midrule
Type mismatch & Objects live in different formal categories & 12 \\
Missing equivalence witness & Plausible but unproved & 9 \\
One-way theorem only & One direction proved, reverse open/false & 4 \\
Data model mismatch & Incompatible data-generating assumptions & 3 \\
Proof method mismatch & Same structure, non-transferable proofs & 3 \\
Success criterion mismatch & Different optimization objectives & 1 \\
\bottomrule
\end{tabular}
\end{center}

The plurality are type mismatches: the objects being compared belong to
different formal categories, and making the analogy a theorem would require
new bridging concepts that do not yet exist. The second largest group, 9
missing equivalence witnesses, is the frontier: analogies where someone
proving the right conjecture would convert an informal remark into a
theorem. The sample compression conjecture sits here. So does the
MML--MDL relationship. These are not structurally blocked; they are open.


% ================================================================
% SECTION 1: TYPE MISMATCH (12 edges)
% ================================================================

\section{Type Mismatch}
\label{sec:type-mismatch}

A type mismatch arises when the source and target of an analogy belong
to different formal categories in the concept graph. The structural
parallel is real: both concepts do something similar. But they operate on
different mathematical objects, and no bridging theorem connects the two
type systems.

These are the most common obstruction, accounting for 12 of the 32 edges,
and the least likely to be resolved. The gap is not waiting for a clever
proof; it is a mismatch in the formalism itself. The analogy is not a
conjecture. It is a description of two things that happen to resemble each
other across a category boundary.

\begin{table}[ht]
\centering\small
\caption{Twelve analogies that fail because the source and target live in
different formal categories. The ``type gap'' column names the specific
category boundary that blocks a formal theorem. For each pair, the shared
structure is real; the gap is why that shared structure cannot be stated
as an equivalence without a bridging concept that does not yet exist.}
\label{tab:type-mismatch}
\begin{tabularx}{\textwidth}{@{}p{4.2cm}p{3.2cm}X@{}}
\toprule
\textbf{Analogy} & \textbf{Shared structure} & \textbf{Type gap} \\
\midrule

\nde{algorithmic\_stability} $\leftrightarrow$ \nde{pac\_bayes\_bound}
& Both bound generalization gap
& Stability is algorithm-dependent (perturbation of $A$). PAC-Bayes is
posterior-dependent (divergence $\KL(Q \| P)$). An algorithm defines a
posterior, but the primary objects differ. \\[6pt]

\nde{information\_theoretic\_bound} $\leftrightarrow$
\nde{pac\_bayes\_bound}
& Both use information-divergence measures
& Info-theoretic bounds use $I(W; S)$ (mutual information between
algorithm output and training set). PAC-Bayes uses $\KL(Q \| P)$
(posterior vs.\ prior). Different information measures on different
objects. Hellstr\"om et al.\ (2023) partially unify via change-of-measure,
but the objects remain distinct. \\[6pt]

\nde{fundamental\_theorem} $\leftrightarrow$
\nde{algorithmic\_stability}
& Both characterize learnability
& The fundamental theorem is a characterization (equivalence result).
Stability is a complexity measure. The structural parallel exists but the
formal types are incompatible without a bridging theorem. \\[6pt]

\nde{sq\_dimension} $\leftrightarrow$ \nde{vc\_dimension}
& Both measure concept class complexity
& SQ dimension measures pairwise correlation structure under a
distribution. VC dimension measures distribution-free shattering.
Exponential gap possible (parities: $\VC = n$, $\SQdim = 2^n$). \\[6pt]

\nde{grammar\_induction} $\leftrightarrow$ \nde{gold\_theorem}
& Both concern language learning
& Grammar induction is a process; Gold's theorem is an impossibility
result. Different concept categories. \\[6pt]

\nde{concept\_drift} $\leftrightarrow$ \nde{ex\_under\_drift}
& Both concern non-stationary learning
& Concept drift is a process; Ex-under-drift is a success criterion.
Different concept categories. \\[6pt]

\nde{lifelong\_learning} $\leftrightarrow$ \nde{meta\_learner}
& Both concern learning across tasks
& Lifelong learning is a process; meta-learner is a learner type.
Different concept categories. \\[6pt]

\nde{background\_knowledge} $\leftrightarrow$ \nde{advice}
& Both provide side information to a learner
& Background knowledge is a process; advice is a data presentation.
Different concept categories. \\[6pt]

\nde{mdl} $\leftrightarrow$ \nde{kolmogorov\_complexity}
& Both measure description length
& MDL uses finite computable codes. Kolmogorov complexity uses the
shortest program on a universal Turing machine (generally
incomputable). The parallel is ``short description = good model,''
but MDL's codes are constructive while Kolmogorov's are not. \\[6pt]

\nde{srm} $\leftrightarrow$ \nde{inductive\_bias}
& Both constrain hypothesis selection
& SRM is a model selection principle; inductive bias is a base type.
The structural parallel exists but the categories are incompatible. \\[6pt]

\nde{bayesian\_inference} $\leftrightarrow$ \nde{inductive\_bias}
& Both encode prior assumptions
& Bayesian inference is a model selection method (probability over
hypotheses). Inductive bias is a base type (constraint on learner
behavior). The prior \emph{is} a specific form of bias, but the type
systems differ. \\[6pt]

\nde{occam\_algorithm} $\leftrightarrow$ \nde{compression\_scheme}
& Both embody ``short representation $\Rightarrow$ generalization''
& Occam uses description length of hypotheses (syntactic: bits to
encode $h$). Compression uses a subset of training data (extensional:
examples to reconstruct $h$). The notion of shortness differs. \\

\bottomrule
\end{tabularx}
\end{table}

In each of these 12 cases, the analogy identifies something real: a
structural parallel that is not an accident. But the source and target live
in different parts of the type system, and formalizing the parallel requires
a bridging theorem relating the two types. No impossibility result rules out
such a bridge for any individual pair. One pair, information-theoretic
bounds versus PAC-Bayes, already has a partial bridge via change-of-measure
\cite{HellstromDurisi2020}. The table records where bridges are currently
absent, not where they are provably out of reach.


% ================================================================
% SECTION 2: MISSING EQUIVALENCE WITNESS (9 edges)
% ================================================================

\section{Missing Equivalence Witness}
\label{sec:missing-equiv}

A missing equivalence witness marks an analogy where the formal equivalence
looks plausible but has no proof. These are the active frontier of the
field: the places where a theorem is conceivably available, and where
finding one would matter.

Not all nine edges here are genuinely open. Three are settled and appear in
this section only because the analogy is commonly stated as if it were an
equivalence. The mistake-bound/Littlestone pair \emph{is} an equivalence
(Littlestone's theorem, proved below). The Rademacher/VC and
label-complexity/sample-complexity pairs are established
\emph{non}-equivalences: the separations are known, not conjectural.
The genuinely open conjectures are compression/VC and the two MML edges.
The distinction between ``open conjecture'' and ``known separation carrying
the same color'' is load-bearing; the exercises at the end of the chapter
press directly on it.

\begin{obstbox}
\textbf{Compression $\leftrightarrow$ VC dimension.}

\medskip\noindent
The sample compression conjecture asserts that every concept class with
VC dimension $d$ admits a sample compression scheme of size $O(d)$.
Moran and Yehudayoff~\cite{MoranYehudayoff2016} proved that finite VC
dimension implies \emph{some} finite compression scheme (of size
exponential in $d$), establishing one
direction.\formal{FLT_Proofs.Theorem.PAC.fundamental_vc_compression} The conjecture that
compression size is $O(d)$ remains open. If proved, this would upgrade
the analogy to an equivalence; if disproved, it would become a
separation. It is the most important open problem connected to the
analogy edges in the graph.
\end{obstbox}

\begin{obstbox}
\textbf{Rademacher complexity $\leftrightarrow$ VC dimension.}

\medskip\noindent
The two are asymptotically related: $\Rad_n(\mathcal{H}) \leq
O\!\left(\sqrt{\VC(\mathcal{H})/n}\right)$. Both measure the complexity
of a hypothesis class, and for binary classification both control
generalization. But Rademacher complexity is data-dependent
(computed from a specific sample and distribution), while VC dimension is
distribution-free. No equivalence holds: Rademacher complexity can be
much smaller than the VC bound for benign distributions, and the two
quantities can disagree on the ranking of hypothesis classes.
\end{obstbox}

\begin{obstbox}
\textbf{VC characterization $\leftrightarrow$ Littlestone
characterization.}

\medskip\noindent
Both theorems have the same logical form: ``a class is learnable in
paradigm $P$ if and only if combinatorial dimension $d$ is finite.''
The parallel is exact at the level of logical structure. But the
dimensions measure different adversarial strengths (batch (VC) vs.\
adaptive (Littlestone)), and $\VC \leq \Ldim$ with no reverse
bound.\formal{FLT_Proofs.Theorem.Separation.vcdim_le_of_mistake_bounded}
The gap can be infinite (thresholds on $\R$). Alon et
al.~\cite{AlonLivniMalliarisMoran2019} connected the two via
differential privacy, but no equivalence between the dimensions exists.
\end{obstbox}

\begin{obstbox}
\textbf{MML $\leftrightarrow$ MDL.}

\medskip\noindent
Minimum Message Length (Wallace--Freeman) and Minimum Description Length
(Rissanen) both select hypotheses by minimizing a two-part code: the
model plus the data given the model. They are often treated as variants
of the same idea. However, MML uses a Bayesian prior explicitly and
integrates over the parameter space, while MDL uses a minimax-optimal
universal code. No formal theorem establishes their equivalence or
separation; the two traditions developed independently and use different
mathematical frameworks.
\end{obstbox}

\begin{obstbox}
\textbf{MML $\leftrightarrow$ Bayesian inference.}

\medskip\noindent
MML selects the hypothesis that minimizes total message length; Bayesian
inference selects (or averages over) hypotheses according to posterior
probability. For certain prior--likelihood pairs, MML and MAP
(maximum a posteriori) estimates coincide asymptotically. But MML is a
coding-theoretic criterion and Bayesian inference is a probabilistic
criterion; the equivalence is approximate and model-dependent, holding for
particular prior--likelihood pairs rather than as a general theorem.
\end{obstbox}

\begin{obstbox}
\textbf{PAC learning $\leftrightarrow$ posterior consistency.}

\medskip\noindent
Both are success criteria requiring convergence to the truth. PAC
demands uniform convergence in probability over all distributions;
Bayesian consistency demands that the posterior concentrates on the true
parameter. The structural parallel is ``learning = convergence,'' but
PAC convergence is frequentist and distribution-free, while posterior
consistency is Bayesian and prior-dependent. No general theorem
connects the two: a class can be PAC learnable but Bayesian-inconsistent
under a misspecified prior, and vice versa.
\end{obstbox}

\begin{obstbox}
\textbf{Mistake bound $\leftrightarrow$ Littlestone dimension.}

\medskip\noindent
The optimal mistake bound for a concept class equals its Littlestone
dimension (this is Littlestone's
theorem).\formal{FLT_Proofs.Theorem.Online.optimal_mistake_bound_eq_ldim}
However, the \emph{analogy} edge in the graph refers to the structural
parallel between the two as complexity measures in their own right. The mistake bound is defined
operationally (worst-case mistakes of the best algorithm); the
Littlestone dimension is defined combinatorially (depth of the largest
shattered tree). Their equivalence is a theorem, but as stand-alone
objects they measure different things, and extending either to new
settings (e.g., partial monitoring, bandit feedback) produces divergent
generalizations.
\end{obstbox}

\begin{obstbox}
\textbf{Label complexity $\leftrightarrow$ sample complexity.}

\medskip\noindent
Both count the amount of data needed for learning, but label complexity
(from active learning) counts the number of \emph{labels queried},
while sample complexity counts the number of \emph{labeled examples
drawn}. In active learning, the learner sees many unlabeled points but
requests labels selectively, so label complexity can be exponentially
smaller than sample complexity. No equivalence exists; the gap is the
content of active learning theory.
\end{obstbox}

\begin{obstbox}
\textbf{Lifelong learning $\leftrightarrow$ concept drift.}

\medskip\noindent
Both concern learning in non-stationary environments. Lifelong learning
emphasizes knowledge transfer across a sequence of tasks; concept drift
emphasizes adaptation to a changing target within a single task. The
structural parallel is ``the learning problem changes over time,'' but
the formalization differs: lifelong learning uses a task distribution,
concept drift uses a time-varying concept. No formal theorem connects
the two frameworks.
\end{obstbox}


% ================================================================
% SECTION 3: ONE-WAY THEOREM ONLY (4 edges)
% ================================================================

\section{One-Way Theorem Only}
\label{sec:one-way}

In these four analogies, one direction has been proved and the other has
not. The structural parallel is real, but it runs one way. What is at stake
here is whether the learner conflates a proved implication with a
biconditional: assuming symmetry that the mathematics does not support, and
then treating the converse as obvious when it is either open or false.

\begin{obstbox}
\textbf{Sample complexity $\leftrightarrow$ VC dimension}
(a theorem that an identity-reading would erase).

\medskip\noindent
The fundamental theorem establishes the biconditional
$\VC(\mathcal{H}) < \infty \Leftrightarrow$ finite sample
complexity.\formal{FLT_Proofs.Theorem.PAC.fundamental_theorem} Unlike the
other entries in this section, this relationship runs \emph{both} ways; it
appears here as a cautionary case.  The analogy is usually stated as an
identity (``sample complexity just \emph{is} VC dimension''), which
collapses a two-directional theorem (an explicit upper bound one way, a
matching lower bound the other, proved by different arguments) into a
definition, erasing exactly the content the fundamental theorem supplies.
\end{obstbox}

\begin{obstbox}
\textbf{Compression scheme $\rightarrow$ sample complexity} (but not
$\leftarrow$).

\medskip\noindent
A compression scheme of size $k$ implies PAC sample complexity
$O(k \cdot \log m / \varepsilon)$ (Littlestone--Warmuth, Moran--Yehudayoff).
But no theorem says that optimal sample complexity determines optimal
compression size. A class might be learnable with few samples yet require
a large compression scheme. The compression direction is proved;
the reverse is open.
\end{obstbox}

\begin{obstbox}
\textbf{Covering number $\rightarrow$ VC dimension} (but not
$\leftarrow$).

\medskip\noindent
The Sauer--Shelah lemma gives: $\log N(\varepsilon, \mathcal{H},
L_1(P)) \leq d \cdot \log(2e/\varepsilon)$, bounding covering numbers
from VC dimension.\formal{FLT_Proofs.Complexity.Rademacher.sauer_shelah_exp_bound} But covering numbers are defined for arbitrary
metric spaces, while VC dimension applies only to binary hypothesis
classes. Finite VC dimension suffices for covering number
bounds; it is not required for them, since the covering number framework is
strictly more general and bounds spaces that have no VC dimension at all.
\end{obstbox}

\begin{obstbox}
\textbf{Covering number $\rightarrow$ Rademacher complexity} (but not
$\leftarrow$).

\medskip\noindent
Dudley's entropy integral bounds Rademacher complexity from above using
covering numbers: $\Rad_n(\mathcal{H}) \leq \inf_{\alpha > 0}
\left[ 4\alpha + \frac{12}{\sqrt{n}} \int_\alpha^1
\sqrt{\log N(\varepsilon, \mathcal{H}, L_2)} \, d\varepsilon \right]$.
This bound can be loose; tighter chaining methods (Talagrand's generic
chaining) and direct Rademacher estimates sometimes give better results.
The Dudley bound is one-directional: covering numbers control
Rademacher, but Rademacher complexity does not determine covering
numbers.
\end{obstbox}


% ================================================================
% SECTION 4: DATA MODEL MISMATCH (3 edges)
% ================================================================

\section{Data Model Mismatch}
\label{sec:data-model}

These three analogies fail at the level of assumptions about how data
arrives. The structural parallel is real: the results look similar and the
intuitions transfer. But the data-generating processes are incompatible.
A result proved under i.i.d.\ sampling does not migrate to adversarial
sequences, and a result that assumes stationarity cannot be directly applied
where non-stationarity is the explicit object of study. The analogy works
as intuition; it breaks as a formal argument the moment the data model is
stated precisely.

\begin{obstbox}
\textbf{Concept drift $\leftrightarrow$ online learning.}

\medskip\noindent
Both are sequential learning settings, but concept drift uses i.i.d.\
draws per time step from a slowly changing distribution, while online
learning uses adversarially chosen instances from a fixed concept class.
The shared structure is ``sequential prediction under non-stationarity,''
but the non-stationarity is stochastic in one case and adversarial in
the other. Results from one setting do not transfer: online regret
bounds assume a fixed comparator class, while drift bounds assume a
rate of distributional change.
\end{obstbox}

\begin{obstbox}
\textbf{Advice reduction $\leftrightarrow$ sample complexity.}

\medskip\noindent
Both quantify ``how much side information shrinks the learner's search
space.'' But advice operates in the Gold model (infinite stream,
computable learner), while sample complexity operates in the PAC model
(finite i.i.d.\ sample, distribution-free guarantees). No formal theorem
bridges the two frameworks. The structural parallel (``side information
helps'') is too vague to formalize without choosing a specific data
model, and the two models are incompatible.

\hfill\cite{KaufmannStephan2001}
\end{obstbox}

\begin{obstbox}
\textbf{Granger causality $\leftrightarrow$ concept drift.}

\medskip\noindent
Both operate on time-series data, but Granger causality assumes
stationarity within estimation windows (testing whether past values
of $X$ improve prediction of $Y$), while concept drift explicitly
models non-stationarity. The structural parallel is ``temporal
dependence in sequential data,'' but the assumptions about the data
process are contradictory: Granger requires stationarity; drift requires
its absence.
\end{obstbox}


% ================================================================
% SECTION 5: PROOF METHOD MISMATCH (3 edges)
% ================================================================

\section{Proof Method Mismatch}
\label{sec:proof-method}

These three analogies have identical theorem structure. The result shapes
are parallel and the logical form is the same. But the proofs inhabit
different mathematical worlds: concentration inequalities in one case,
change-of-measure arguments in another, Ramsey theory in a third. The
analogy holds at the level of what is proved; the obstruction is that
proving one tells you nothing about how to prove the other. Knowing that
stability gives a generalization bound does not, by itself, give you a
proof technique for information-theoretic bounds, even though the conclusion
has the same form.

\begin{obstbox}
\textbf{Information-theoretic bound $\leftrightarrow$ algorithmic
stability.}

\medskip\noindent
Both bound the generalization gap via properties of the learning
algorithm. Stability uses sup-norm perturbation sensitivity (how much
does the output change when one training point is replaced?).
Information-theoretic bounds use mutual information $I(W; S)$ (how much
does the algorithm's output depend on the training set?). These are
connected via the chain: low stability $\Rightarrow$ low $I(W; S)$
$\Rightarrow$ low generalization gap. But the proof methods are
fundamentally different: stability arguments use coupling and
McDiarmid-type concentration; information-theoretic arguments use
change-of-measure and data processing inequalities. Neither proof
technique subsumes the other.
\end{obstbox}

\begin{obstbox}
\textbf{Ordinal VC dimension $\leftrightarrow$ mind-change ordinal.}

\medskip\noindent
Both use ordinal-valued measures to characterize learnability in
transfinite settings. Ordinal VC dimension (Martin--Sharma--Stephan)
characterizes predictive complexity; mind-change ordinals characterize
identification complexity in the Gold model. Both use ordinals, but they
measure different things. The connection is indirect, routed through predictive
complexity: ordinal VC dimension governs prediction error
bounds, while mind-change ordinals govern convergence to a correct
hypothesis. The proof techniques (Ramsey-theoretic for ordinal VC,
topological for mind-change) do not transfer.
\end{obstbox}

\begin{obstbox}
\textbf{Occam's razor $\leftrightarrow$ MDL.}

\medskip\noindent
Both assert that simpler models generalize better. Occam's razor (in the
PAC sense of Blumer et al.) is a worst-case generalization theorem:
if the hypothesis is short, its error is low with high probability.
MDL is an information-theoretic model selection principle: minimize the
total description length of model plus data-given-model. The theorem
\emph{structure} is parallel (``short $\Rightarrow$ good''), but Occam
gives frequentist, distribution-free guarantees while MDL gives
Bayesian-flavored, coding-theoretic guarantees. The proof methods
inhabit different mathematical worlds.

\hfill\cite{BlumerEhrenfeuchtHausslerWarmuth1987}
\end{obstbox}


% ================================================================
% SECTION 6: SUCCESS CRITERION MISMATCH (1 edge)
% ================================================================

\section{Success Criterion Mismatch}
\label{sec:criterion-mismatch}

A single analogy edge in the graph carries this obstruction. The type does not
reduce to the others: the two concepts share a setting and a procedure, but
they answer different questions about that setting.

\begin{obstbox}
\textbf{Granger causality $\leftrightarrow$ online learning.}

\medskip\noindent
Granger causality tests causal structure: does the past of time series
$X$ improve prediction of time series $Y$ beyond what $Y$'s own past
provides? Online learning minimizes cumulative prediction error against
an adversary. Both involve sequential prediction, but they optimize
different objectives. Granger seeks a \emph{structural} conclusion
(causal influence exists or does not); online learning seeks a
\emph{performance} guarantee (low regret). The success criteria are
fundamentally different: statistical significance of a causal test vs.\
sub-linear regret growth. No formal theorem connects them because they
answer different questions about the same sequential data.
\end{obstbox}


% ================================================================
% SECTION 7: THE OBSTRUCTION MAP
% ================================================================

\section{The Obstruction Map}
\label{sec:obstruction-map}

\Cref{fig:obstruction-map} puts all 32 analogy edges in one place, colored
by obstruction type, with a filled dot marking the five edges where
formalization has already produced a verified bridge. Reading the color
distribution is reading the map of what formal learning theory does not yet
know how to unify, and the five dots mark exactly where that effort has
succeeded so far.

\begin{figure}[ht]
\centering
\begin{tikzpicture}[
  x=1.5cm, y=1.2cm,
  concept/.style={
    rectangle, rounded corners=2pt, draw, thick,
    fill=layergray, font=\scriptsize\sffamily,
    minimum width=1.3cm, minimum height=0.5cm,
    align=center, inner sep=2pt
  },
  tm/.style={-, edgeorange, thick},        % type mismatch
  mew/.style={-, defblue, thick, dashed},   % missing equiv witness
  owt/.style={-, sepgreen, thick, dotted},  % one-way theorem
  dm/.style={-, thmred, thick, dashdotted}, % data model mismatch
  pm/.style={-, analogypurple, thick, densely dashed}, % proof method
  sc/.style={-, boundbrown, thick, densely dotted},    % success criterion
  ev/.style={circle, fill=black, draw=white, line width=0.25pt,
             inner sep=0pt, minimum size=3.6pt}, % verified-handle dot
]

% Layout by connected component: the analogy graph is disconnected, so each
% component sits in its own region (zero inter-region crossings) and is drawn
% planar. Color/style = obstruction type; a filled dot = a verified sorry-free
% kernel handle backs that edge.

% --- Complexity core (VC-hub component) ---
\node[concept] (vc)      at (1.5,3.5)  {VC dim};
\node[concept] (sq)      at (0.4,3.5)  {SQ dim};
\node[concept] (rad)     at (1.5,4.5)  {Rademacher};
\node[concept] (cov)     at (2.7,4.5)  {Covering\\number};
\node[concept] (comp)    at (1.5,2.5)  {Compression};
\node[concept] (sc_node) at (2.85,3.0) {Sample\\complexity};
\node[concept] (lc)      at (4.1,3.4)  {Label\\complexity};
\node[concept] (ar)      at (4.1,2.4)  {Advice\\reduction};

% --- Model-selection tail (hangs off compression via Occam) ---
\node[concept] (occam)   at (1.5,1.5)  {Occam};
\node[concept] (mdl)     at (2.7,1.5)  {MDL};
\node[concept] (kolm)    at (2.7,0.6)  {Kolmogorov};
\node[concept] (mml)     at (3.9,1.5)  {MML};
\node[concept] (bayes)   at (3.9,0.6)  {Bayesian};
\node[concept] (ib)      at (5.1,1.05) {Inductive\\bias};
\node[concept] (srm)     at (5.1,2.0)  {SRM};

% --- Verified dimension-equivalence pairs (isolated edges, both dotted) ---
\node[concept] (ldim)    at (3.9,4.6)  {Littlestone\\dim};
\node[concept] (mb)      at (5.1,4.6)  {Mistake\\bound};
\node[concept] (vcc)     at (3.5,5.5)  {VC charac.};
\node[concept] (lc_char) at (4.9,5.5)  {Ldim charac.};

% --- Sequential component (drift-hub) ---
\node[concept] (drift)   at (7.0,3.4)  {Concept\\drift};
\node[concept] (ol)      at (6.3,2.6)  {Online};
\node[concept] (gc)      at (7.7,2.6)  {Granger};
\node[concept] (exd)     at (7.7,4.2)  {Ex under\\drift};
\node[concept] (ll)      at (6.3,4.2)  {Lifelong};
\node[concept] (ml)      at (6.3,5.1)  {Meta-\\learner};

% --- Paradigm component (stability-hub) ---
\node[concept] (stab)    at (3.1,-0.4) {Stability};
\node[concept] (pb)      at (4.1,0.1)  {PAC-Bayes};
\node[concept] (itb)     at (4.1,-0.9) {Info-theoretic};
\node[concept] (ft)      at (2.0,-0.4) {Fund.\\Theorem};

% --- Remaining isolated edges ---
\node[concept] (pac)     at (0.4,1.5)  {PAC};
\node[concept] (post)    at (0.4,0.6)  {Posterior\\consistency};
\node[concept] (gi)      at (1.1,-1.7) {Grammar\\induction};
\node[concept] (gt)      at (2.3,-1.7) {Gold's\\theorem};
\node[concept] (bk)      at (3.7,-1.7) {Background\\knowledge};
\node[concept] (adv)     at (4.9,-1.7) {Advice};

% --- Type mismatch edges (orange, 12) ---
\draw[tm] (stab) -- (pb);
\draw[tm] (itb) -- (pb);
\draw[tm] (ft) -- (stab);
\draw[tm] (sq) -- (vc);
\draw[tm] (gi) -- (gt);
\draw[tm] (drift) -- (exd);
\draw[tm] (ll) -- (ml);
\draw[tm] (bk) -- (adv);
\draw[tm] (mdl) -- (kolm);
\draw[tm] (srm) -- (ib);
\draw[tm] (bayes) -- (ib);
\draw[tm] (occam) -- (comp);

% --- Missing equivalence witness edges (blue dashed, 9); 3 verified-backed ---
\draw[mew] (comp) -- (vc)      node[ev,pos=0.5]{}; % fundamental_vc_compression
\draw[mew] (rad) -- (vc);
\draw[mew] (vcc) -- (lc_char)  node[ev,pos=0.5]{}; % vcdim_le_of_mistake_bounded
\draw[mew] (mml) -- (mdl);
\draw[mew] (mml) -- (bayes);
\draw[mew] (pac) -- (post);
\draw[mew] (mb) -- (ldim)      node[ev,pos=0.5]{}; % optimal_mistake_bound_eq_ldim
\draw[mew] (lc) -- (sc_node);
\draw[mew] (ll) -- (drift);

% --- One-way theorem edges (green dotted, 4); 2 verified-backed ---
\draw[owt] (sc_node) -- (vc)  node[ev,pos=0.5]{}; % fundamental_theorem
\draw[owt] (comp) -- (sc_node);
\draw[owt] (cov) -- (vc)      node[ev,pos=0.5]{}; % sauer_shelah_exp_bound
\draw[owt] (cov) -- (rad);                        % Dudley: no kernel handle, undotted

% --- Data model mismatch edges (red dashdotted, 3) ---
\draw[dm] (drift) -- (ol);
\draw[dm] (ar) -- (sc_node);
\draw[dm] (gc) -- (drift);

% --- Proof method mismatch edges (purple, 3) ---
\draw[pm] (itb) -- (stab);
\draw[pm] (occam) -- (mdl);
% ordinal_vc_dim -> mind_change_ordinal not shown (nodes would clutter)

% --- Success criterion mismatch (brown, 1) ---
\draw[sc] (gc) -- (ol);

\end{tikzpicture}

\caption[Obstruction map: thirty-two analogies, six failure types, five
verified bridges]{Thirty-two cross-paradigm analogies, none of them a
theorem, fail for six classifiable reasons. Color and line style name the
obstruction type; a filled dot marks the five edges where formalization has
already produced a verified one-way bridge or equivalence
(\leanname{fundamental_theorem}). The argument the figure makes is that
analogy-failure is not case-by-case: it is taxonomic, and the six types
exhaust what the field currently has names for. The five dots show the
frontier where the taxonomy has been broken by a proof. Key:
\textcolor{edgeorange}{\rule[0.4ex]{1.1em}{0.9pt}} type mismatch (12);
\textcolor{defblue}{\textbf{-\,-\,-}} missing equivalence witness (9);
\textcolor{sepgreen}{\textbf{$\cdots$}} one-way theorem (4);
\textcolor{thmred}{\textbf{-$\cdot$-}} data-model mismatch (3);
\textcolor{analogypurple}{\textbf{-\,-}} proof-method mismatch (3);
\textcolor{boundbrown}{\textbf{$\cdots$}} success-criterion mismatch (1);
$\bullet$~verified bridge handle (5).
One proof-method edge (ordinal VC dimension $\leftrightarrow$ mind-change
ordinal) is omitted for clarity.}
\label{fig:obstruction-map}
\end{figure}

\subsection*{What the distribution tells us}

The obstruction type distribution reveals three structural facts about
formal learning theory.

\emph{First}, 12 of the 32 failures are type mismatches: the objects being
compared are not the same kind of thing. This is not a failure of proof
technique. The proofs have not simply not been found; the objects live in
different categories. Future unification will require new bridging concepts
that mediate between the existing type systems.

\emph{Second}, 9 missing equivalence witnesses mark the open frontier.
The sample compression conjecture (compression $\leftrightarrow$ VC
dimension) is the most prominent: Moran and Yehudayoff proved the
exponential-size direction in 2016, and the $O(d)$ conjecture remains open.
The Rademacher--VC relationship and the characterization-theorem parallel
(VC $\leftrightarrow$ Littlestone) are equally fundamental. A proof on any
of these would restructure the analogy map.

\emph{Third}, the 4 one-way theorems show where the asymmetry is built into
the objects themselves. Covering numbers bound Rademacher complexity via
Dudley's entropy integral, and covering numbers are bounded by VC dimension
via Sauer--Shelah. Both directions flow one way only. The asymmetry is not
an artifact of proof technique; it reflects a generality hierarchy: covering
numbers apply to metric spaces, Rademacher complexity to function classes,
VC dimension to binary classes. Each step downward tightens the bound and
loses applicability.

The 32 edges and their obstructions together form a map of what formal
learning theory does not yet know how to unify. Understanding why an analogy
fails is, in this field, as informative as understanding why a theorem
holds.


% ================================================================
% SECTION 8: REALIZED ANALOGIES (the verified counterpoint)
% ================================================================

\section{When the Analogy Is a Theorem}
\label{sec:realized-analogies}

Some analogies carry no obstruction at all. They began as loose parallels
and someone proved them. The cleanest is now machine-checked.

The flagship is the correspondence \emph{online learning
$\leftrightarrow$ zero-sum games}. A learner facing an adversary and a row
player in a matrix game look, at first glance, merely analogous: both react
to a worst-case opponent. The analogy is discharged by the
multiplicative-weights (Hedge) algorithm. Run on the game matrix, it
produces a distribution over rows whose guaranteed payoff is within
$\varepsilon$ of the game's value. No-regret online play is not
\emph{like} approximate game solving; it \emph{is} approximate game solving.

\begin{thm}[Multiplicative weights solve the game]
\label{thm:mwu-minimax}
Let $M \colon R \times C \to \{0,1\}$ be a finite Boolean game whose value
is at least~$v$ (every column distribution admits a row response of payoff
$\geq v$). For every $\varepsilon > 0$, the multiplicative-weights
algorithm returns a row distribution~$p$ with payoff $\geq v - \varepsilon$
against every column~$c$.\formal{ProbabilityTheory.mwu_approx_minimax} In
particular the minimax value is attained in the limit, recovering the
finite minimax theorem.\formal{ProbabilityTheory.finite_approx_minimax}
\end{thm}

\begin{proof}[Proof idea]
The algorithm keeps a weight per row, multiplied by a factor $e^{\eta}$
whenever that row would have scored against the observed column, then
renormalized to a distribution. A potential argument bounds the gap
between the algorithm's cumulative payoff and the best fixed row by
$O(\sqrt{\log|R| / T})$ after $T$ rounds; taking $T$ large enough drives
the gap below~$\varepsilon$. Those same weights, read as a mixed strategy,
certify the minimax value: the boosting-as-a-game and regret-as-minimax
correspondence, formalized end to end.
\end{proof}

This realized analogy has a transformer face. The multiplicative-weights
distribution is a softmax over accumulated scores, exactly the operation
a softmax-attention head performs over its key scores. Attention is, in
this precise sense, one step of the Hedge update, and the head's induced
concept class is
well-behaved.\formal{TLT.Bridge.softAttention_wellBehaved} The
online$\,\leftrightarrow\,$game analogy thus reaches into the architecture
of \Cref{ch:generalization}: the potential argument that solves a matrix
game is the same normalization that sits inside every attention layer.

\begin{remark}[Why the realized case matters]
\label{rem:realized-matters}
A taxonomy of obstructions would be vacuous if every analogy were blocked.
\Cref{thm:mwu-minimax} shows the frontier of \Cref{sec:missing-equiv} is
worth pushing: an analogy that looked like a metaphor (``learning against an
adversary is like playing a game'') is a theorem, with a softmax at its
center. The obstruction types classify only the analogies not yet discharged.
\end{remark}


% ================================================================
% OPEN PROBLEMS
% ================================================================

\section{Open Problems}
\label{sec:analogies-open}

The open problems here come in two varieties. The first is to discharge a
parallel into a theorem: find the proof that converts a blue dashed edge
into a verified bridge. The second is harder: prove that no such theorem
can exist, that a specific type mismatch is permanently blocked and not
merely waiting for the right idea. Each problem below is tagged a
\emph{known unknown} (KU), sharply posed but unresolved, or an
\emph{unknown known} (UK), where the right bridging object has not yet
been named.

\begin{openprob}[The sample compression conjecture]
\label{op:compression-vc-analogy}
Does every class of VC dimension~$d$ admit a compression scheme of
size~$O(d)$, discharging the compression$\,\leftrightarrow\,$VC analogy
into an equivalence?  The exponential-size direction is machine-checked
(\leanname{fundamental_vc_compression}); the $O(d)$ conjecture is the most
prominent missing witness.  A known unknown~\cite{MoranYehudayoff2016}.
\end{openprob}

\begin{openprob}[MML versus MDL]
\label{op:mml-mdl}
Is there a formal theorem relating Minimum Message Length and Minimum
Description Length: an equivalence on a restricted class of priors, or a
separation?  Both coding objects are formalizable (\leanname{MDLPrinciple},
\leanname{KolmogorovComplexity}); neither direction has a witness.  A known
unknown.
\end{openprob}

\begin{openprob}[Completing the information--PAC-Bayes bridge]
\label{op:itb-pacbayes-bridge}
Hellstr\"om et al.\ partially bridge information-theoretic and PAC-Bayes
bounds via change-of-measure.  Can the partial bridge be completed to a
theorem that recovers each bound from the other with explicit constants,
turning a ``type mismatch'' into a realized analogy?  An unknown
known~\cite{HellstromDurisi2020}.
\end{openprob}

\begin{openprob}[Formalizing the obstruction taxonomy]
\label{op:formalize-taxonomy}
Can ``type mismatch,'' ``one-way theorem,'' and the other obstruction
types be made precise (categorically, say), so that the classification of
the 32 edges becomes a theorem rather than an expert judgment?  An unknown
known: the meta-level vocabulary has no formal definition yet.
\end{openprob}

\begin{openprob}[Attention as Hedge, at scale]
\label{op:attention-hedge}
\Cref{thm:mwu-minimax} identifies a single softmax with one Hedge step.
Does the no-regret guarantee survive composition? Does a stacked
attention model inherit a minimax/regret certificate, or does depth break
the correspondence?  A known unknown on the transformer side
(\leanname{mwu_approx_minimax}, \leanname{softAttention_wellBehaved}).
\end{openprob}

\begin{openprob}[Provable permanence of a type mismatch]
\label{op:type-mismatch-impossibility}
For a specific pair (algorithmic stability versus PAC-Bayes, say), is
there a proof that \emph{no} bridging theorem exists, or is the
``permanence'' of \Cref{sec:type-mismatch} merely the current absence of a
bridge?  A known unknown: the chapter conjectures permanence but proves
none.
\end{openprob}

\begin{openprob}[Where Rademacher beats VC]
\label{op:rademacher-vc-dist}
Characterize the distributions under which the empirical Rademacher
complexity is exponentially smaller than the VC bound.  A sharp
distributional characterization would convert the
Rademacher$\,\leftrightarrow\,$VC ``missing witness'' into a precise
separation theorem.  A known unknown.
\end{openprob}

\begin{openprob}[Ordinal VC versus mind-change ordinals]
\label{op:ordinal-vc-mindchange}
The two transfinite measures are linked here only by a proof-method
mismatch.  Is there a direct theorem relating the ordinal VC dimension
(\leanname{vcdim_to_ordinal_vcdim}) to the mind-change ordinal of
\Cref{ch:ordinals}?  An unknown known~\cite{MartinSharmaStephan2006}.
\end{openprob}

\begin{openprob}[The active--passive gap]
\label{op:active-passive}
Characterize exactly the classes for which label complexity is
exponentially smaller than sample complexity.  A structural
characterization would turn the label$\,\leftrightarrow\,$sample
``missing witness'' into a theorem about which classes benefit from active
querying.  A known unknown.
\end{openprob}

\begin{openprob}[Minimax for continuous attention]
\label{op:minimax-attention}
The verified minimax theorem (\leanname{finite_approx_minimax}) is finite.
Does it lift to the continuous game a softmax-attention head plays over an
unbounded key set, giving a verified value for attention as an adversarial
prediction game?  An unknown known on the transformer side
(\leanname{softAttention_wellBehaved}).
\end{openprob}

\subsection*{Counterfactual exercises}

The figure's central claim is that analogy-failure is taxonomic: thirty-two
independent failures collapse to six types. The exercises that follow probe
whether you can read that claim off the figure itself, using only the two
codes each edge carries. The color and line style name the obstruction type
(why the analogy is not a theorem). The filled dot marks the five edges
where a verified bridge already exists. The exercises work these codes
against each other, especially the cases where both codes apply to the same
edge and appear to conflict.

\begin{enumerate}[leftmargin=*,itemsep=3pt]
\item Count colors and ignore node identities. Thirty-two analogies fail across
only six obstruction types, with type mismatch the plurality (twelve), so ``why
does this analogy fail?'' has a finite-codomain answer rather than a
case-by-case story. Argue that this taxonomic structure is the figure's central
claim.

\item Read only the edges carrying a filled dot. Exactly five of the
thirty-two edges carry a verified bridge, and they cluster on the
complexity-measure side of the graph. Explain why this small dotted subgraph is
the machine-checked fragment while the rest is conceptual analysis or open
conjecture.

\item The pair \nde{mistake\_bound}--\nde{littlestone\_dimension} \emph{is} a
verified equivalence (\leanname{optimal_mistake_bound_eq_ldim}, hence the dot),
yet the edge is blue. Show that the ``missing witness'' names a distinct fact:
as complexity measures in their own right the two diverge when extended to
bandit or partial-monitoring feedback (\Cref{sec:missing-equiv}). Why does the
diagram need both codes here?

\item On \nde{compression\_scheme}--\nde{vc\_dimension}, the dot marks
\leanname{fundamental_vc_compression} (finite VC dimension implies \emph{some}
finite compression scheme, exponential in $d$), but the edge stays blue dashed
because the $O(d)$ sample compression conjecture is open
(\Cref{op:compression-vc-analogy}). Explain how one edge can carry a verified
direction and an open equivalence at once.

\item The edges \nde{rademacher\_complexity}--\nde{vc\_dimension} and
\nde{label\_complexity}--\nde{sample\_complexity} are blue but carry no dot.
These are \emph{established} non-equivalences (Rademacher complexity can fall
far below the VC bound; label complexity can be exponentially smaller than
sample complexity, \Cref{op:active-passive}), not open conjectures. Explain how
the absence of a dot plus the prose distinguishes a known separation from a
live conjecture inside the same color class.

\item Both \nde{covering\_number}--\nde{vc\_dimension} and
\nde{covering\_number}--\nde{rademacher\_complexity} are green ``one-way
theorem'' edges, but only the first carries a dot
(\leanname{sauer_shelah_exp_bound}); the second is Dudley's entropy integral,
stated classically (\Cref{sec:one-way}). Show how the evidence code splits one
obstruction type into a verified bridge and a literature bridge.

\item On \nde{sample\_complexity}--\nde{vc\_dimension} the dot marks
\leanname{fundamental_theorem}, the equivalence between finite VC dimension and
finite sample complexity. The analogy is usually stated as an identity
(``sample complexity just \emph{is} VC dimension''), which erases the two
directions proved by different arguments (\Cref{sec:one-way}). Explain why the
dot guards a two-directional theorem against being read as a definition.

\item Ask whether the six obstruction types are exhaustive: must every
cross-paradigm analogy fail for one of exactly these reasons? A completeness
theorem would require defining the types formally
(\Cref{op:formalize-taxonomy}) and proving the permanence of a single mismatch
(\Cref{op:type-mismatch-impossibility}). Until both are settled, the counts six
and five are a census, not an invariant of the field.
\end{enumerate}

% Chapter 16: Computational vs. Information-Theoretic Learnability
% Centerpiece: the Kearns-Valiant 1994 gap — classes with finite VC dim
% that are NOT efficiently learnable (under cryptographic assumptions).
% Concepts: computational_hardness, proper_improper_separation, requires_assumption edges.
% SHORT chapter (~3-4 pages).

\chapter{Computational vs.\ Information-Theoretic Learnability}
\label{ch:computational}

The Fundamental Theorem of \Cref{ch:pac} characterizes learnability in
information-theoretic terms: $\mathcal{H}$ is PAC learnable if and only if
$\VC(\mathcal{H}) < \infty$.  The characterization is silent about
\emph{computation}.  It guarantees the \emph{existence} of a sample size
$m(\varepsilon, \delta)$ and a learner $A$ that achieves accuracy $\varepsilon$
with confidence $1 - \delta$, but it says nothing about how long $A$ takes to
run.

This silence conceals a gap.  There exist hypothesis classes with finite VC
dimension, hence information-theoretically learnable, that no polynomial-time
algorithm can PAC learn, unless widely believed cryptographic assumptions are
false.  The gap between what is \emph{learnable} and what is \emph{efficiently
learnable} is the subject of this chapter.  It is a negative result, but a
structurally important one: it shows that information-theoretic
characterizations, however elegant, do not tell the whole story.

\medskip

The chapter has three sections.

\begin{enumerate}[leftmargin=*,itemsep=2pt]
\item \textbf{The information--computation gap} (\Cref{sec:gap}): the
  Kearns--Valiant result and the cryptographic assumptions it requires.
\item \textbf{Proper vs.\ improper learning} (\Cref{sec:proper-improper-comput}): the
  separation between proper and improper learnability, and why the
  representation matters.
\item \textbf{The \rel{requires\_assumption} edges}
  (\Cref{sec:requires-assumption}): what this edge type means structurally in
  the concept graph.
\end{enumerate}


% ================================================================
% SECTION 1: THE INFORMATION--COMPUTATION GAP
% ================================================================

\section{The Information--Computation Gap}
\label{sec:gap}

\begin{graphnode}[title={Graph Node: \nde{computational\_hardness}}]
Layer: \texttt{meta}.  Incoming edge:
$\edge{computational\_hardness}{requires\_assumption}{cryptographic\_hardness}$.
This node represents the structural fact that efficient learnability requires
assumptions beyond finite VC dimension, specifically, cryptographic hardness.
\end{graphnode}

The central result is due to Kearns and Valiant~\cite{KearnsValiant1994}.

\begin{thm}[Kearns--Valiant 1994]
\label{thm:kearns-valiant}
Under a cryptographic hardness assumption (the original construction uses the
hardness of factoring; later strengthenings use decisional composite
residuosity or learning with errors, detailed below), there exist concept
classes $\mathcal{C}$ with:
\begin{enumerate}[leftmargin=*,nosep]
\item $\VC(\mathcal{C}) < \infty$, so that $\mathcal{C}$ is PAC learnable
  information-theoretically,\formal{FLT_Proofs.Theorem.PAC.fundamental_theorem} but
\item no polynomial-time algorithm PAC learns $\mathcal{C}$.
\end{enumerate}
\end{thm}

The original construction uses Boolean formulae: the class of polynomial-size
Boolean formulae has VC dimension polynomial in the number of variables, yet
PAC learning this class is at least as hard as breaking certain public-key
cryptosystems.  The reduction goes through the construction of
\emph{cryptographic pseudorandom function families}: if a learner could
efficiently distinguish the target concept from random, it could break the
cryptographic assumption.

\begin{remark}[The structure of the reduction]
\label{rmk:kv-reduction}
The Kearns--Valiant reduction is from a cryptographic primitive to a learning
problem, not from a worst-case problem.  The learner is given random examples
$(x, c(x))$ where $x$ is drawn from a distribution and $c$ is the target
concept.  If the learner succeeds in learning, it implicitly solves the
cryptographic problem on random instances.  This is why the hardness is
\emph{average-case}, not worst-case.
\end{remark}

\subsection{Why P $\neq$ NP Does Not Suffice}

One might hope that the information--computation gap follows from
$\mathrm{P} \neq \mathrm{NP}$.  It does not.  Applebaum, Barak, and
Xiao~\cite{ApplebaumBarakXiao2008} showed the following.

\begin{thm}[Applebaum--Barak--Xiao 2008]
\label{thm:abx}
There is no \emph{black-box} reduction from $\mathrm{NP}$-hardness to the
hardness of PAC learning.  Specifically, if PAC learning a class $\mathcal{C}$
is hard, this hardness cannot be established by a reduction that treats the
learner as an oracle, unless $\mathrm{NP} \subseteq \mathrm{coAM}$.
\end{thm}

The reason is structural.  PAC learning is an \emph{average-case} problem: the
learner receives random examples from a distribution and must generalize.
$\mathrm{NP}$-hardness is a \emph{worst-case} notion.  Reductions from
worst-case to average-case problems are notoriously difficult and, in many
settings, provably impossible via relativizing techniques.

\subsection{The Necessary Assumptions}

The assumptions used in the literature fall into three families, in
roughly increasing order of generality:

\begin{enumerate}[leftmargin=*,itemsep=3pt]
\item \textbf{Factoring and RSA.}  The original Kearns--Valiant result
  assumes that factoring $n$-bit integers requires super-polynomial time.
  This is the most classical assumption but also the most fragile: it is
  broken by quantum computers (Shor's algorithm).

\item \textbf{Decisional Composite Residuosity (DCRA).}  Used in subsequent
  refinements, DCRA assumes that it is hard to distinguish $n$th residues
  modulo $n^2$.  Like factoring, this is vulnerable to quantum attacks.

\item \textbf{Learning With Errors (LWE).}  Introduced by
  Regev~\cite{Regev2005}, the LWE assumption states that it is hard to recover
  a secret vector $\mathbf{s}$ from noisy inner products
  $\langle \mathbf{a}_i, \mathbf{s} \rangle + e_i \pmod{q}$.  LWE is
  believed to resist quantum computers and has become the standard assumption
  in modern hardness-of-learning results~\cite{Daniely2021}.
\end{enumerate}

\begin{obstbox}
The information--computation gap is an \emph{obstruction} in the concept
graph: the edge
$\edge{vc\_dimension}{characterizes}{pac\_learnability}$ is
information-theoretically valid, but no analogous computational edge is known.
Whether \emph{any} combinatorial dimension characterizes efficient PAC
learnability is open; no such dimension has been found, and no proof rules
one out.  What \emph{is} settled, and machine-checked, is a sharper local
fact: finite VC dimension does not control the statistical-query dimension.
Singleton indicators on $\N$ have $\VC = 1$ yet infinite SQ
dimension,\formal{FLT_Proofs.Theorem.Extended.vcdim_not_implies_hardness} so
VC dimension already fails to capture hardness in the SQ model, one concrete
instance of the gap, even though its full extent is unknown.
\end{obstbox}


% ================================================================
% SECTION 2: PROPER VS. IMPROPER LEARNING
% ================================================================

\section{Proper vs.\ Improper Learning}
\label{sec:proper-improper-comput}

The information--computation gap has a subtler cousin: the gap between
\emph{proper} and \emph{improper} learning.

\begin{defn}[Proper and Improper Learning]
\label{def:proper-improper}
A PAC learner for $\mathcal{H}$ is \emph{proper} if it always outputs a
hypothesis $h \in \mathcal{H}$.  It is \emph{improper} if it may output any
efficiently evaluable function $h$, not necessarily in $\mathcal{H}$.
\end{defn}

The distinction matters because the \emph{representation} of hypotheses affects
computational complexity.

\begin{thm}[Proper--Improper Separation]
\label{thm:proper-improper}
There exist concept classes $\mathcal{C}$ such that:
\begin{enumerate}[leftmargin=*,nosep]
\item proper PAC learning of $\mathcal{C}$ is $\mathrm{NP}$-hard (so, unless
  $\mathrm{RP} = \mathrm{NP}$, no efficient proper learner exists), but
\item there exists a polynomial-time improper learner for $\mathcal{C}$.
\end{enumerate}
\end{thm}

The canonical example is the class of $3$-term DNF formulae.  Finding a
consistent $3$-term DNF is $\mathrm{NP}$-hard~\cite{PittValiant1988}, but the
class of $3$-term DNFs is contained in the class of $3$-CNF formulae, which
\emph{is} efficiently learnable.  An improper learner simply outputs a $3$-CNF
formula that is consistent with the data.  The hypothesis is not a DNF, but it
classifies correctly.

Information-theoretically, a proper learner always exists: any class of finite
VC dimension is learned by proper empirical risk minimization, which returns a
consistent member of the class whenever one
exists.\formal{FLT_Proofs.Complexity.Compression.vcdim_finite_imp_proper_finite_support_learner}
The separation is therefore purely computational: the proper learner exists
but cannot be run efficiently, whereas the improper learner can.

\begin{remark}[Representation vs.\ class]
\label{rmk:representation}
The proper--improper separation illustrates a principle that runs throughout
computational learning theory: what matters for computational complexity is not
the \emph{set of functions} $\mathcal{H} \subseteq \{0,1\}^X$ but the
\emph{representation} of those functions.  Two representations of the same set
of Boolean functions can have dramatically different computational properties.
This is why the VC dimension, which depends only on the set
$\mathcal{H}$, cannot capture computational learnability.
\end{remark}

The structural consequence is that efficient learnability is a property of
\emph{represented} hypothesis classes, not of hypothesis classes as sets.  The
concept graph captures this through the node \nde{proper\_improper\_separation},
which has type \texttt{separation} and connects to
\nde{computational\_hardness} via a \rel{motivates} edge.


\subsection{The Transformer Face of the Gap}
\label{sec:attention-gap}

The information--computation gap is not a museum piece; it is the daily
condition of deep learning.  A softmax-attention head induces a concept class
that is machine-checked to be well-behaved, hence of finite VC dimension and
information-theoretically PAC
learnable.\formal{TLT.Bridge.softAttention_wellBehaved} A proper learner for
it therefore exists in principle.  Whether that learner runs in polynomial
time is the open, and widely doubted, question: hardness-of-learning results
for neural networks under LWE~\cite{Daniely2021} place architectures of this
kind on the conjecturally efficient-unlearnable side of
\Cref{thm:kearns-valiant}.

The proper/improper distinction is what makes practice possible. Recovering
the exact parameters of a target attention head, which is proper learning, is the
regime where hardness bites. Gradient descent does something improper
instead: it returns \emph{a} network that predicts well, not the target
network, exactly as the $3$-CNF learner returns a CNF rather than the target
DNF. The gap of this chapter is, in this sense, the precise reason that
training large models is an improper-learning enterprise.


% ================================================================
% SECTION 3: THE REQUIRES_ASSUMPTION EDGES
% ================================================================

\section{The \rel{requires\_assumption} Edges}
\label{sec:requires-assumption}

The concept graph contains several edges of type \rel{requires\_assumption}.
These edges encode a specific logical structure: the source node's validity
depends on an assumption that is not provable from the axioms of the theory
alone.

\begin{defn}[The \rel{requires\_assumption} edge type]
\label{def:requires-assumption}
An edge $\edge{A}{requires\_assumption}{B}$ means: the result or property
associated with node $A$ holds \emph{conditional on} the assumption encoded
by node $B$.  If $B$ is falsified, the status of $A$ becomes indeterminate.
\end{defn}

The principal instances in the graph are:

\begin{enumerate}[leftmargin=*,itemsep=3pt]
\item $\edge{computational\_hardness}{requires\_assumption}{cryptographic\_hardness}$.
  The information--computation gap of \Cref{thm:kearns-valiant} holds only
  under cryptographic assumptions (hardness of factoring, or LWE).  Without
  them it remains conceivable, though it would be a startling
  collapse, that every class of finite VC dimension is efficiently learnable.
  The assumption is genuinely cryptographic: the hardness of a specific
  computational problem, not a constraint on the learner's hypothesis space.

\item $\edge{posterior\_consistency}{requires\_assumption}{bayesian\_inference}$.
  Posterior consistency (\Cref{ch:agents}, \Cref{sec:bayesian-pac}) holds under
  conditions on the prior and the model class.  When these conditions
  fail, for example when the prior assigns zero mass to neighborhoods of
  the true parameter, posterior consistency fails
  (Diaconis--Freedman~\cite{DiaconisFreedman1986}).  The edge records this
  conditional dependence.
\end{enumerate}

\begin{remark}[Structural meaning]
\label{rmk:requires-assumption-structure}
The \rel{requires\_assumption} edge type is the graph's mechanism for
representing \emph{conditional} theorems.  Unlike \rel{implies} edges, which
are unconditional, a \rel{requires\_assumption} edge signals that the source
node's content is only as secure as the target node's assumption.  This is
the graph-theoretic encoding of the phrase ``under the assumption that\ldots''
that pervades computational learning theory.
\end{remark}


% ================================================================
% OPEN PROBLEMS
% ================================================================

\section{Open Problems}
\label{sec:computational-open}

The computational layer is where formal learning theory is least complete:
its central hardness results are conditional, and its central
characterization question is open.  Each problem below is tagged a
\emph{known unknown} (KU), sharply posed but unresolved, or an \emph{unknown
known} (UK), where the right object is still missing; a design-lab or
literature anchor is named where one exists.

\begin{openprob}[A dimension for efficient learnability]
\label{op:efficient-dimension}
Is there \emph{any} combinatorial parameter whose finiteness characterizes
efficient (polynomial-time) PAC learnability, as VC dimension characterizes
the information-theoretic case?  None is known, and none is ruled out.  A
known unknown, the central question of the computational layer.
\end{openprob}

\begin{openprob}[Formalizing the cryptographic gap]
\label{op:formalize-kv}
The information-theoretic side of \Cref{thm:kearns-valiant} is machine-checked
(\leanname{vcdim_finite_imp_proper_finite_support_learner}); the cryptographic
hardness side is not.  Can a Kearns--Valiant-style hardness result, a finite-VC
class with no efficient learner under LWE, be formalized in a proof
assistant?  An unknown known~\cite{KearnsValiant1994,Daniely2021}.
\end{openprob}

\begin{openprob}[Distribution-free SQ hardness, verified]
\label{op:sq-distribution-free}
The machine-checked separation finite-VC$\,\not\Rightarrow\,$finite-SQ uses a
fixed witness distribution (\leanname{vcdim_not_implies_hardness}).  Can one
verify a finite-VC class whose SQ dimension is super-polynomial under
\emph{every} distribution, matching the parities lower
bound~\cite{BlumFurstJacksonKearnsMansourRudich1994}?  A known unknown.
\end{openprob}

\begin{openprob}[Is attention efficiently learnable?]
\label{op:attention-efficient}
A softmax-attention class is information-theoretically learnable
(\leanname{softAttention_wellBehaved}).  Is it efficiently PAC learnable, or
does learning attention reduce to a cryptographic problem?  A known unknown on
the transformer side, and the modern instance of \Cref{thm:kearns-valiant}.
\end{openprob}

\begin{openprob}[Past the black-box barrier]
\label{op:worst-to-average}
Applebaum--Barak--Xiao rule out \emph{black-box} reductions from
$\mathrm{NP}$-hardness to learning hardness.  Is there a non-black-box route
that bases hardness of learning on a worst-case assumption such as
$\mathrm{P} \neq \mathrm{NP}$?  A known unknown~\cite{ApplebaumBarakXiao2008}.
\end{openprob}

\begin{openprob}[Characterizing improper learnability]
\label{op:improper-characterization}
Efficient \emph{improper} learning can be strictly easier than proper learning
(\Cref{thm:proper-improper}).  Is there a structural characterization of the
classes that are efficiently improperly learnable?  An unknown known: no
combinatorial handle on improper efficiency is known.
\end{openprob}

\begin{openprob}[Minimal assumptions for hardness]
\label{op:minimal-assumption}
Hardness-of-learning currently rests on strong primitives (factoring, LWE).
Is some cryptographic assumption \emph{necessary} for the
information--computation gap: does an efficient learner for every finite-VC
class imply the collapse of one-way functions?  A known
unknown~\cite{Daniely2021}.
\end{openprob}

\begin{openprob}[Learning implies cryptography?]
\label{op:learning-to-crypto}
The Kearns--Valiant reduction builds hardness \emph{from} cryptography. Does
the converse hold: does the hardness of learning a natural concept class
\emph{yield} a cryptographic primitive?  An unknown known: a generic
learning-to-cryptography transformation is not established.
\end{openprob}

\begin{openprob}[Size of the proper/improper gap]
\label{op:proper-improper-size}
For \Cref{thm:proper-improper}, how large can the separation be: how much
harder is proper than improper learning across natural classes, measured in
running time?  A known unknown~\cite{PittValiant1988}.
\end{openprob}

\begin{openprob}[A proper/improper gap for attention]
\label{op:attention-proper-improper}
Does an attention class exhibit a proper/improper gap: is there an efficient
improper learner (a surrogate model) even where recovering the target head is
intractable, formalizing the intuition that gradient descent learns
improperly?  A known unknown on the transformer side
(\leanname{softAttention_wellBehaved},
\leanname{vcdim_finite_imp_proper_finite_support_learner}).
\end{openprob}

\bigskip

This chapter has established the structural gap between
information-theoretic and computational learnability.  The gap is real: it
produces classes that are provably learnable but (conditionally) not
efficiently learnable.  And no combinatorial dimension is yet known to
characterize efficient PAC learnability; whether one can exist is itself an
open problem.  The next chapter examines a different kind of gap: the distance
between binary classification and the richer learning settings that require
new mathematical primitives.

% Chapter 17: Extensions Beyond Binary Classification
%   ds_dimension — Profile C (Revelation): The 30-year gap from Natarajan to DS
%   multiclass_learning — Profile G (Framework Bridge): Multiple paradigms tried
%   agnostic_learning — Profile B-light: The agnostic fundamental theorem
%   proper_improper_separation — Profile D (Witness Construction)
% Theme: GRAMMAR GROWTH — each extension requires new primitives.
% The extends_grammar edge type is the structural backbone.

\chapter{Extensions Beyond Binary Classification}
\label{ch:extensions}

The PAC theory of \Cref{ch:pac} and its characterization through VC dimension
apply to binary classification: the label set is $\{0,1\}$, the hypothesis
class is $\mathcal{H} \subseteq \{0,1\}^X$, and the loss is $0$-$1$.  This is
the setting where the theory is cleanest and most complete.  But most learning
problems are not binary classification.  Labels may come from a finite set of
$k > 2$ classes; the target may be a real-valued function; the data may be
noisy; the realizability assumption may fail.

This chapter asks: how does the binary theory extend to these richer settings?
The answer is not ``straightforwardly.''  Each extension requires
\emph{new mathematical primitives}: new notions of shattering, new complexity
measures, new assumptions.  The binary theory generalizes by \emph{grammar
growth}: substituting richer labels into the old definitions is never enough,
because each setting forces a primitive the old language cannot name.  In the
concept graph,
this is captured by edges of type \rel{extends\_grammar}: the source dimension
extends the grammar of the VC dimension by introducing primitives that have no
analogue in the binary case.

\medskip

The chapter has six sections.

\begin{enumerate}[leftmargin=*,itemsep=2pt]
\item \textbf{Multiclass learning} (\Cref{sec:multiclass-ext}): the Natarajan
  dimension, the DS dimension, and the resolution of a thirty-year open problem.
\item \textbf{Real-valued functions} (\Cref{sec:real-valued}):
  pseudodimension and fat-shattering dimension.
\item \textbf{Agnostic learning} (\Cref{sec:agnostic-ext}): dropping realizability
  and the agnostic fundamental theorem.
\item \textbf{Noise-tolerant and partial concept learning}
  (\Cref{sec:noise-partial}): classification noise, malicious noise,
  and promise problems.
\item \textbf{Proper versus improper learning}
  (\Cref{sec:proper-improper}): the computational gap.
\item \textbf{The \rel{extends\_grammar} pattern}
  (\Cref{sec:extends-grammar}): what the extensions have in common.
\end{enumerate}


% ================================================================
%     CONCEPT MAP FIGURE
% ================================================================

\begin{figure}[ht]
\centering
\begin{tikzpicture}[
  >={Stealth[length=5pt]},
  concept/.style={
    rectangle, rounded corners=3pt, draw, thick,
    fill=layergray, font=\small\sffamily,
    minimum width=2.3cm, minimum height=0.7cm,
    align=center
  },
  % COLOR = relationship type ; LINE = evidence grade (the second axis)
  chr_solid/.style  ={->, defblue,    thick},            % characterizes, kernel-verified
  chr_lit/.style    ={->, defblue,    thick, dashed},    % characterizes, literature
  ext_concept/.style={->, edgeorange, thick, dotted},    % extends_grammar, conceptual
  ss_lit/.style     ={->, thmred,     thick, dashed},    % strictly_stronger, literature
  res_concept/.style={->, sepgreen,   thick, dotted},    % restricts, conceptual
  lb_lit/.style     ={->, boundbrown, thick, dashed},    % lower_bounds, literature
]

% Layout by structure: two hubs (vc, pac) on a horizontal spine; the multiclass
% cycle (nat, ds) fans UP, the real-valued branch (pdim, fat) fans DOWN, the
% settings (agn, pi, noise) sit on the right. Planar straight-line drawing =
% ZERO crossings. Two codes: color = relationship type, line = evidence grade.

% --- hubs ---
\node[concept] (vc)    at (0,0)      {VC\\Dimension};
\node[concept] (pac)   at (6.5,0)    {PAC\\Learning};
% --- multiclass branch (fans up) ---
\node[concept] (nat)   at (2.2,1.9)  {Natarajan\\Dimension};
\node[concept] (ds)    at (4.3,3.1)  {DS\\Dimension};
% --- real-valued branch (fans down-left) ---
\node[concept] (pdim)  at (1.8,-1.9) {Pseudo-\\dimension};
\node[concept] (fat)   at (2.8,-3.1) {Fat-Shattering\\Dimension};
% --- agnostic: small triangle just under the verified spine (clears fat->pac) ---
\node[concept] (agn)   at (3.5,-0.8) {Agnostic\\PAC};
% --- settings (right) ---
\node[concept] (noise) at (8.7,1.4)  {Noise-Tolerant\\Learning};
\node[concept] (pi)    at (8.7,-1.5) {Proper vs\\Improper};

% --- edges: color = relationship type, line = evidence grade ---
\draw[chr_solid]   (vc)    -- (pac);   % VERIFIED rung: fundamental_theorem (sc=0)
\draw[ext_concept] (nat)   -- (vc);
\draw[ext_concept] (pdim)  -- (vc);
\draw[ext_concept] (noise) -- (pac);   % NEW edge: connects the formerly floating noise node
\draw[ss_lit]      (ds)    -- (nat);
\draw[chr_lit]     (ds)    -- (pac);
\draw[chr_lit]     (fat)   -- (pac);
\draw[chr_lit]     (vc)    -- (agn);
\draw[res_concept] (agn)   -- (pac);
\draw[lb_lit]      (pi)    -- (pac);

\end{tikzpicture}

\caption[Concept map: binary PAC generalizes by grammar growth]{Concept map for
the chapter: color marks the relationship type and line style the evidence grade.
Binary PAC theory generalizes by grammar growth, each extension adding a
shattering primitive with no binary analogue, so substituting richer labels into
the binary definitions never suffices.  The one solid edge is VC dimension
characterizing binary PAC learning (\leanname{fundamental_theorem}); the
multiclass, real-valued, agnostic, and proper-versus-improper characterizations
are literature results.}
\label{fig:ch17-concept-map}
\end{figure}


% ================================================================
% SECTION 1: MULTICLASS LEARNING
% ================================================================

\section{Multiclass Learning: From Natarajan to DS}
\label{sec:multiclass-ext}

Let $Y = [k] = \{1, 2, \ldots, k\}$ for some $k \geq 2$, and let
$\mathcal{H} \subseteq Y^X$ be a multiclass hypothesis class.  The question
is: what characterizes PAC learnability of $\mathcal{H}$?

For $k = 2$, the answer is the VC dimension.  For $k > 2$, the answer was only
fully resolved in 2022, after a thirty-year journey through two competing
dimensions, each introducing a new shattering primitive.  That journey, and its
resolution, is the narrative backbone of this section.


\subsection{The Natarajan Dimension}

Natarajan~\cite{Natarajan1989} introduced the first multiclass complexity
measure by generalizing shattering to require two functions instead of one.

\begin{defn}[Natarajan Dimension {\cite{Natarajan1989}}]
\label{def:natarajan-dim}
The \emph{Natarajan dimension} $\Ndim(\mathcal{H})$ of a multiclass hypothesis
class $\mathcal{H} \subseteq Y^X$ is the largest size $d$ of a set
$\{x_1, \ldots, x_d\} \subseteq X$ for which there exist two functions
$f_0, f_1 \colon \{x_1, \ldots, x_d\} \to Y$ satisfying:
\begin{enumerate}[leftmargin=*,nosep]
\item $f_0(x_i) \neq f_1(x_i)$ for all $i \in [d]$, and
\item for every binary string $b \in \{0,1\}^d$, there exists $h \in
  \mathcal{H}$ with $h(x_i) = f_{b_i}(x_i)$ for all $i$.\formal{FLT_Proofs.Complexity.Structures.NatarajanDim}
\end{enumerate}
\end{defn}

The new primitive is the \emph{pair of functions} $(f_0, f_1)$.  In binary
classification, $f_0$ and $f_1$ are forced to be the constant functions $0$
and $1$, and the definition reduces to VC shattering.  In the multiclass case,
the choice of which two labels to use at each point is part of the
combinatorial structure.

\begin{example}[Natarajan shattering with $k=3$]
\label{ex:natarajan-k3}
Let $X = \{a, b\}$, $Y = \{1,2,3\}$.  Define $f_0(a)=1$, $f_0(b)=2$
and $f_1(a)=3$, $f_1(b)=1$.  To N-shatter $\{a,b\}$, we need four
hypotheses realizing every binary string:
\[
  h_{00}(a)=1,\, h_{00}(b)=2;\quad
  h_{01}(a)=1,\, h_{01}(b)=1;\quad
  h_{10}(a)=3,\, h_{10}(b)=2;\quad
  h_{11}(a)=3,\, h_{11}(b)=1.
\]
Any $\mathcal{H} \supseteq \{h_{00}, h_{01}, h_{10}, h_{11}\}$ has $\Ndim \geq 2$.
\end{example}

\begin{historical}
\textbf{The Natarajan gap (1989--2022).}
Natarajan showed that $\Ndim(\mathcal{H}) < \infty$ is \emph{sufficient}
for PAC learnability and provided sample complexity bounds.  The best
known bounds, due to Ben-David et al.~\cite{BenDavidCesaBianchiHausslerLong1995},
were:
\[
  C_1 \cdot \frac{\Ndim}{\varepsilon}
  \;\leq\; m_{\mathrm{PAC}}(\varepsilon,\delta)
  \;\leq\; C_2 \cdot \frac{\Ndim \cdot \ln k \cdot \ln(1/\varepsilon)}{\varepsilon}.
\]
Two gaps persisted.  First, the $\Theta(\log k)$ factor between upper
and lower bounds.  Second, and more fundamentally: is finiteness of $\Ndim$
also \emph{necessary}?  For finite label sets $|Y| = k < \infty$, the answer
is yes (the $\log k$ factor is at most a polynomial overhead).  But when
$|Y| = \infty$, a class can have $\Ndim = 1$ and yet fail to be PAC learnable.
This was the open problem that persisted for over thirty years.
\end{historical}


The \nde{natarajan\_dimension} extends the grammar of the \nde{vc\_dimension} by
the two-function shattering primitive $(f_0, f_1)$, and leaves open a tight
multiclass sample complexity without the $\log k$ factor.


\subsection{The DS Dimension: A Thirty-Year Resolution}
\label{subsec:ds-dim}

The correct characterization of multiclass PAC learnability was established
by Brukhim et al.~\cite{BrukhimCarmonDinurMoranYehudayoff2022} using the
\emph{DS dimension}, named after Daniely and
Shalev-Shwartz~\cite{DanielyShalevShwartz2014} who introduced the
underlying concept of pseudo-cubes in one-inclusion hypergraphs.

\begin{defn}[One-Inclusion Hypergraph]
\label{def:oig}
Let $\mathcal{H} \subseteq Y^X$ and let $S = \{x_1, \ldots, x_n\} \subseteq X$
be a finite set.  The \emph{one-inclusion hypergraph} $\mathrm{OIG}(\mathcal{H}, S)$
has vertex set $\mathcal{H}|_S = \{(h(x_1), \ldots, h(x_n)) : h \in \mathcal{H}\}$
and, for each coordinate $i \in [n]$, a hyperedge connecting all vertices that
agree on all coordinates except possibly coordinate $i$.
\end{defn}

\begin{defn}[DS Dimension {\cite{DanielyShalevShwartz2014,BrukhimCarmonDinurMoranYehudayoff2022}}]
\label{def:ds-dim-ext}
A set $\{x_1, \ldots, x_d\} \subseteq X$ is \emph{DS-shattered} by
$\mathcal{H} \subseteq Y^X$ if there exists a function
$f \colon \{x_1, \ldots, x_d\} \to Y$ such that for every $i \in [d]$,
there exists $h_i \in \mathcal{H}$ with:
\begin{enumerate}[leftmargin=*,nosep]
\item $h_i(x_i) \neq f(x_i)$, and
\item $h_i(x_j) = f(x_j)$ for all $j \neq i$.
\end{enumerate}
The \emph{DS dimension} $\DSdim(\mathcal{H})$ is the size of the largest
DS-shattered set.
\end{defn}

The structural content of DS-shattering is different from Natarajan shattering
in a revealing way.  Natarajan shattering asks for a \emph{global} pair of
labeling functions with \emph{all} $2^d$ interpolations realized.
DS-shattering asks for a single ``default'' labeling $f$ and, for each
point, a \emph{local} witness that deviates at exactly that point.  The
requirement is weaker per point (only one deviation, not arbitrary binary
combinations) but more structural: it encodes a property of the
one-inclusion hypergraph of $\mathcal{H}$.

\begin{figure}[ht]
\centering
\begin{tikzpicture}[
  scale=0.82,
  point/.style={circle, fill=defblue, inner sep=1.8pt},
  hyp/.style={font=\small},
  agree/.style={circle, fill=defblue!20, draw=defblue, inner sep=1.8pt},
  deviate/.style={circle, fill=thmred!20, draw=thmred, inner sep=1.8pt},
  brace/.style={decorate, decoration={brace, amplitude=5pt}},
  oigv/.style={circle, draw=defblue, fill=defblue!12, inner sep=2.4pt, font=\scriptsize},
  cedge/.style={defblue, thick}
]
% ===== LEFT PANEL: the DS witness table (the definition, d = 3) =====
\node[font=\bfseries\small] at (1.5, 3.15) {Witness table};

\node[hyp] at (-1.5, 2.0) {Default $f$:};
\foreach \i/\lbl in {1/a, 2/b, 3/c} {
  \node[point] (x\i) at (1.5*\i, 2.0) {};
  \node[above=3pt] at (x\i) {$x_{\i}$};
  \node[below=5pt, font=\small, text=defblue] at (x\i) {$\lbl$};
}

% Witness h1 : deviates at x1 only
\node[hyp] at (-1.5, 0.6) {$h_1$:};
\node[deviate] (h11) at (1.5, 0.6) {};  \node[below=5pt, font=\small, text=thmred] at (h11) {$\neq a$};
\node[agree]   (h12) at (3.0, 0.6) {};  \node[below=5pt, font=\small, text=defblue] at (h12) {$b$};
\node[agree]   (h13) at (4.5, 0.6) {};  \node[below=5pt, font=\small, text=defblue] at (h13) {$c$};

% Witness h2 : deviates at x2 only
\node[hyp] at (-1.5, -0.4) {$h_2$:};
\node[agree]   (h21) at (1.5, -0.4) {}; \node[below=5pt, font=\small, text=defblue] at (h21) {$a$};
\node[deviate] (h22) at (3.0, -0.4) {}; \node[below=5pt, font=\small, text=thmred] at (h22) {$\neq b$};
\node[agree]   (h23) at (4.5, -0.4) {}; \node[below=5pt, font=\small, text=defblue] at (h23) {$c$};

% Witness h3 : deviates at x3 only
\node[hyp] at (-1.5, -1.4) {$h_3$:};
\node[agree]   (h31) at (1.5, -1.4) {}; \node[below=5pt, font=\small, text=defblue] at (h31) {$a$};
\node[agree]   (h32) at (3.0, -1.4) {}; \node[below=5pt, font=\small, text=defblue] at (h32) {$b$};
\node[deviate] (h33) at (4.5, -1.4) {}; \node[below=5pt, font=\small, text=thmred] at (h33) {$\neq c$};

\draw[brace] (5.15, 0.9) -- (5.15, -1.7)
  node[midway, right=6pt, font=\scriptsize, align=left]
  {Each $h_i$ deviates\\from $f$ at $x_i$ only};

% ===== separator =====
\draw[gray!45, thin] (7.05, 3.3) -- (7.05, -2.5);

% ===== RIGHT PANEL: the one-inclusion hypergraph (star = the witness) =====
% OIG vertices = realized label-vectors; coordinate-i hyperedge joins vectors
% agreeing off coord i. Here f and h_i agree off coord i, so the edges are
% exactly {f,h_i}: a star. (h_i, h_j differ at TWO coords -> not adjacent.)
\node[font=\bfseries\small] at (10.0, 3.15) {One-inclusion hypergraph};

\node[oigv, fill=defblue!24] (Of)  at (10.0, 0.55) {$f$};
\node[oigv] (Oh1) at (8.55, 1.75) {$h_1$};
\node[oigv] (Oh2) at (11.45, 1.75) {$h_2$};
\node[oigv] (Oh3) at (10.0, -1.30) {$h_3$};

\draw[cedge] (Of) -- (Oh1) node[midway, above left=-2pt, font=\scriptsize, text=defblue] {$i{=}1$};
\draw[cedge] (Of) -- (Oh2) node[midway, above right=-2pt, font=\scriptsize, text=defblue] {$i{=}2$};
\draw[cedge] (Of) -- (Oh3) node[midway, right=1pt, font=\scriptsize, text=defblue] {$i{=}3$};

\node[align=center, font=\scriptsize] at (10.0, -2.15)
  {$f$ joins each $h_i$ along edge $i$; a minimum-outdegree\\
   proper orientation (the learner) runs on this structure};

\end{tikzpicture}
\caption[DS-shattering and its star in the one-inclusion hypergraph]{%
DS-shattering of a $d$-element set (here $d = 3$), drawn twice.  On the left a
default labeling $f$ assigns $a, b, c$ to $x_1, x_2, x_3$, and each witness $h_i$
agrees with $f$ except at coordinate $i$; on the right the same data is the star
joining $f$ to each $h_i$ along edge $i$ in the one-inclusion hypergraph
(\Cref{def:oig}).  Finiteness of the largest such set, the DS dimension,
characterizes multiclass PAC learnability (\Cref{thm:multiclass-char}, Brukhim et
al.\ 2022), where Natarajan shattering does not when the label set is infinite.}
\label{fig:ds-shattering}
\end{figure}


\begin{thm}[Multiclass Characterization {\cite{BrukhimCarmonDinurMoranYehudayoff2022}}]
\label{thm:multiclass-char}
A multiclass hypothesis class $\mathcal{H} \subseteq Y^X$ is PAC learnable if
and only if $\DSdim(\mathcal{H}) < \infty$.
\end{thm}

\begin{proof}[Proof sketch]
We outline both directions.

\medskip\noindent
\textbf{Sufficiency} ($\DSdim(\mathcal{H}) < \infty \Rightarrow$ learnable).
The proof proceeds through the one-inclusion hypergraph.  For a finite
restriction $\mathcal{H}|_S$ to a sample $S$ of size $m$, the one-inclusion
hypergraph $\mathrm{OIG}(\mathcal{H}, S)$ has vertices $\mathcal{H}|_S$ and
hyperedges indexed by coordinates.  The key insight of Daniely and
Shalev-Shwartz is that a \emph{proper orientation} of this hypergraph, an
assignment of each hyperedge to one of its vertices, yields a learning
algorithm.  Specifically:

\begin{enumerate}[leftmargin=*,nosep]
\item Given a sample $S = \{(x_1, y_1), \ldots, (x_m, y_m)\}$, restrict
  $\mathcal{H}$ to $\{x_1, \ldots, x_m\}$ and build $\mathrm{OIG}(\mathcal{H}, S)$.
\item Find a proper orientation with \emph{minimum outdegree}.
\item Return the hypothesis whose vertex has minimum outdegree in the
  orientation.
\end{enumerate}

If $\DSdim(\mathcal{H}) = d < \infty$, then the one-inclusion hypergraph
of any restriction has no $d$-dimensional pseudo-cube, and a combinatorial
argument (via the Helly property of CAT(0) cube complexes) shows that
a proper orientation with outdegree at most $O(d)$ exists.  The learner's
error is bounded by $O(d/m)$, yielding PAC learnability with sample
complexity $O(d/\varepsilon)$.

\medskip\noindent
\textbf{Necessity} ($\DSdim(\mathcal{H}) = \infty \Rightarrow$ not learnable).
If $\DSdim(\mathcal{H}) = \infty$, then for every $d$, there exists a
finite set that is DS-shattered.  The proof constructs, for each $d$, a
distribution $D_d$ over $X \times Y$ such that:
\begin{itemize}[leftmargin=*,nosep]
\item The Bayes-optimal classifier has zero error.
\item Any learner given $m < d$ samples has expected error
  $\Omega(d/m)$ on $D_d$.
\end{itemize}
The distribution is supported on the DS-shattered set and uses the default
labeling $f$ as the target.  The witness hypotheses $h_1, \ldots, h_d$ act
as ``confusers'': a learner that has not seen point $x_i$ cannot distinguish
$f$ from $h_i$, because they agree everywhere except at $x_i$.  Since this
works for arbitrarily large $d$, no finite sample suffices for
$\varepsilon$-accurate learning uniformly over $D_d$.
\end{proof}


\begin{obstbox}
\textbf{Natarajan $\not\Rightarrow$ PAC learnability when $|Y| = \infty$.}

\medskip\noindent
\emph{Witness (Brukhim et al.~2022).}
There exists a hypothesis class $\mathcal{H} \subseteq Y^X$ with
$\Ndim(\mathcal{H}) = 1$ but $\DSdim(\mathcal{H}) = \infty$.  The
construction uses a hyperbolic pseudo-manifold: the domain $X$ is the
vertex set of a high-dimensional simplicial complex with hyperbolic
geometry, and $\mathcal{H}$ encodes colorings of simplices.  The
hyperbolic geometry ensures that any two points can be N-shattered
(only two labels are needed locally), but the global structure admits
arbitrarily large DS-shattered sets because the one-inclusion hypergraph
contains arbitrarily large pseudo-cubes.

\medskip\noindent
\emph{Structural lesson.}
The Natarajan dimension is a \emph{local} measure: it examines pairs of
labels at each point.  The DS dimension is a \emph{global} measure: it
examines the connectivity structure of the entire one-inclusion hypergraph.
When $|Y| = \infty$, local two-label constraints do not capture global
learnability.

\medskip\noindent
\emph{Graph edge.}  Finiteness of the \nde{natarajan\_dimension} does not imply
\nde{pac\_learning}, witnessed by the Brukhim et al.\ (2022) non-learnable class
with $\Ndim = 1$.
\end{obstbox}


The relationship between the two dimensions is:

\begin{prop}[DS vs.\ Natarajan]
\label{prop:ds-vs-nat}
For any $\mathcal{H} \subseteq Y^X$:
\[
  \Ndim(\mathcal{H}) \;\leq\; \DSdim(\mathcal{H}).
\]
Moreover, when $|Y| = k < \infty$:
\[
  \DSdim(\mathcal{H}) \;\leq\; O\!\bigl(\Ndim(\mathcal{H}) \cdot \log k\bigr).
\]
\end{prop}

\begin{proof}
For the first inequality, we show that every N-shattered set is also
DS-shattered.  Suppose $S = \{x_1, \ldots, x_d\}$ is N-shattered by a pair
$(f_0, f_1)$.  Take the default labeling $f = f_0$.  For each $i$, the binary
string $e_i$ (with a $1$ in coordinate $i$ and $0$ elsewhere) is realized by
some $h_i \in \mathcal{H}$ with $h_i(x_i) = f_1(x_i) \neq f_0(x_i) = f(x_i)$
and $h_i(x_j) = f_0(x_j) = f(x_j)$ for $j \neq i$.  Thus $f$ together with
$h_1, \ldots, h_d$ witnesses that $S$ is DS-shattered.  Every N-shattered set
is therefore DS-shattered, so $\Ndim \leq \DSdim$.  DS-shattering is the
\emph{weaker} requirement (only $d$ single-coordinate deviations, not all
$2^d$ interpolations), which is exactly why $\DSdim$ can be infinite while
$\Ndim$ stays finite, as in the obstruction above.

For the second inequality (finite $k$), suppose $\DSdim = D$.  Any
DS-shattered set provides $D$ witness hypotheses deviating from a default.
Since each deviation chooses one of $k-1$ alternative labels, a pigeonhole
argument over the $k^D$ possible deviation patterns shows
$D \leq O(\Ndim \cdot \log k)$.
\end{proof}


The \nde{ds\_dimension} characterizes \nde{pac\_learning} in the multiclass
setting and is strictly stronger than the \nde{natarajan\_dimension}, the measure
that resolved a thirty-year open problem.


% ================================================================
% SECTION 2: REAL-VALUED FUNCTIONS
% ================================================================

\section{Real-Valued Functions}
\label{sec:real-valued}

Now let $\mathcal{F} \subseteq \R^X$ be a class of real-valued functions.
The learner observes samples $(x_i, y_i)$ with $y_i = f(x_i)$ (or
$y_i \approx f(x_i)$ in the noisy case) and aims to approximate $f$.
The $0$-$1$ loss is replaced by the squared loss or absolute loss, and the
question becomes: what is the analogue of VC dimension?

Two dimensions extend the binary theory.  Both introduce \emph{thresholds} as
a new primitive.


\subsection{Pseudodimension}

\begin{defn}[Pseudodimension {\cite{Pollard1984,Haussler1992}}]
\label{def:pseudodim-ext}
The \emph{pseudodimension} $\Pdim(\mathcal{F})$ of a function class
$\mathcal{F} \subseteq \R^X$ is the VC dimension of the class of subgraphs:
\[
  \Pdim(\mathcal{F}) \;=\; \VC\!\bigl(\bigl\{
    x \mapsto \ind[f(x) \geq t] : f \in \mathcal{F},\; t \in \R
  \bigr\}\bigr).
\]
Equivalently, $\Pdim(\mathcal{F}) = d$ if there exist $x_1, \ldots, x_d \in X$
and thresholds $t_1, \ldots, t_d \in \R$ such that for every
$b \in \{0,1\}^d$, there exists $f \in \mathcal{F}$ with
$f(x_i) \geq t_i \iff b_i = 1$.\formal{FLT_Proofs.Complexity.Structures.Pseudodimension}
\end{defn}

The new primitive is the \emph{threshold vector} $(t_1, \ldots, t_d)$.
In binary classification, thresholds are unnecessary: the outputs are
already $0$ or $1$.  For real-valued functions, thresholds convert the
continuous outputs into a binary shattering condition.

\begin{example}[Linear functions]
For the class of linear functions $\mathcal{F} = \{x \mapsto w^\top x : w \in \R^d\}$
on $\R^d$, we have $\Pdim(\mathcal{F}) = d$.  Any $d$ affinely independent
points $x_1, \ldots, x_d$ can be P-shattered with thresholds $t_i = 0$:
for each $b \in \{0,1\}^d$, one can find $w$ such that $\mathrm{sign}(w^\top x_i) = (-1)^{1-b_i}$.
\end{example}

\begin{prop}[Pseudodimension and covering numbers]
\label{prop:pdim-covering}
If $\Pdim(\mathcal{F}) = d < \infty$, then for every probability measure
$P$ on $X$ and every $\varepsilon > 0$:
\[
  \mathcal{N}(\varepsilon, \mathcal{F}, L_1(P))
  \;\leq\; \left(\frac{2e}{\varepsilon}\right)^d.
\]
\end{prop}

\begin{proof}[Proof sketch]
The subgraph class $\{x \mapsto \ind[f(x) \geq t] : f \in \mathcal{F},\,
t \in \R\}$ has VC dimension $\Pdim(\mathcal{F}) = d$ by definition.  Applying
the Sauer--Shelah bound and the standard covering-number estimate for VC
classes (\Cref{ch:generalization}) to this binary class, then transferring
back to $L_1(P)$ through the layer-cake identity
$f = \int_0^1 \ind[f \geq t]\,\mathrm{d}t$, yields the stated bound.
\end{proof}

As a complexity measure, the \nde{pseudodimension} extends the grammar of the
\nde{vc\_dimension}, and the extension is the threshold primitive.


\subsection{Fat-Shattering Dimension}

Pseudodimension controls learnability in the realizable case.  For the
agnostic setting with real-valued targets, a scale-sensitive
notion is needed.

\begin{defn}[Fat-Shattering Dimension {\cite{AlonBenDavidCesaBianchiHaussler1997}}]
\label{def:fat-shattering-ext}
For $\gamma > 0$, the \emph{$\gamma$-fat-shattering dimension}
$\fatshat_\gamma(\mathcal{F})$ is the largest $d$ such that there exist
$x_1, \ldots, x_d \in X$ and thresholds $t_1, \ldots, t_d \in \R$ satisfying:
for every $b \in \{0,1\}^d$, there exists $f \in \mathcal{F}$ with
\[
  f(x_i) \geq t_i + \gamma \;\text{if } b_i = 1,
  \qquad
  f(x_i) \leq t_i - \gamma \;\text{if } b_i = 0.
\]
\end{defn}

The new primitive beyond pseudodimension is the \emph{margin}
$\gamma$.\formal{FLT_Proofs.Complexity.Structures.FatShatteringDim}
As $\gamma \to 0$, the fat-shattering dimension approaches the
pseudodimension.

\begin{figure}[ht]
\centering
\begin{tikzpicture}[scale=0.9]
  % single margin value: forbidden band is [t_i - \g, t_i + \g], width 2\g, centred on t_i
  \def\g{0.5}

  % --- Axes ---
  \draw[->] (0,-0.5) -- (7,-0.5) node[right] {$x$};
  \draw[->] (0,-0.5) -- (0,4.2) node[above] {$f(x)$};

  % --- Per point: symmetric allowed zones, threshold, x-tick + label ---
  \foreach \i/\xpos/\tval in {1/1.5/1.5, 2/3.5/2.5, 3/5.5/1.8} {
    \fill[edgeorange!22] (\xpos-0.32, \tval+\g) rectangle (\xpos+0.32, \tval+\g+0.85);
    \fill[defblue!18]    (\xpos-0.32, \tval-\g-0.85) rectangle (\xpos+0.32, \tval-\g);
    \draw[dashed, notegray] (\xpos-0.5, \tval) -- (\xpos+0.5, \tval)
      node[right, font=\tiny, notegray] {$t_{\i}$};
    \draw (\xpos, -0.55) -- (\xpos, -0.45);
    \node[font=\tiny] at (\xpos, -0.75) {$x_{\i}$};
  }

  % --- zone labels + 2gamma brace, marked once at x_1 to keep it clean ---
  \node[font=\tiny, edgeorange!80!black] at (1.5, 1.5+\g+0.45) {$\geq t_i{+}\gamma$};
  \node[font=\tiny, defblue]             at (1.5, 1.5-\g-0.45) {$\leq t_i{-}\gamma$};
  \draw[<->, thick, sepgreen] (1.5+0.42, 1.5-\g) -- (1.5+0.42, 1.5+\g)
    node[midway, right, font=\tiny, text=sepgreen] {$2\gamma$};

  % ===== FUNCTION CURVES (the witnesses that exhibit the definition) =====
  % f_A realizes b=(1,0,1): >= t_1+g, <= t_2-g, >= t_3+g
  \draw[thmred, very thick]
    plot[smooth, tension=0.7] coordinates
    {(0.3,1.9) (1.5,2.4) (2.5,2.3) (3.5,1.7) (4.5,1.9) (5.5,2.6) (6.4,2.7)};
  \fill[thmred] (1.5,2.4) circle (2.2pt);
  \fill[thmred] (3.5,1.7) circle (2.2pt);
  \fill[thmred] (5.5,2.6) circle (2.2pt);
  \node[thmred, font=\scriptsize] at (2.5,3.2) {$f_A\colon b{=}101$};

  % f_B realizes b=(0,1,0): <= t_1-g, >= t_2+g, <= t_3-g
  \draw[analogypurple, very thick, dashed]
    plot[smooth, tension=0.7] coordinates
    {(0.3,1.1) (1.5,0.6) (2.5,1.6) (3.5,3.3) (4.5,1.9) (5.5,0.9) (6.4,0.7)};
  \fill[analogypurple] (1.5,0.6) circle (2.2pt);
  \fill[analogypurple] (3.5,3.3) circle (2.2pt);
  \fill[analogypurple] (5.5,0.9) circle (2.2pt);
  \node[analogypurple, font=\scriptsize] at (4.6,3.6) {$f_B\colon b{=}010$};
\end{tikzpicture}
\caption[$\gamma$-fat-shattering: curves realizing labelings against shared
thresholds]{$\gamma$-fat-shattering of $d=3$ points: the set is
$\gamma$-fat-shattered when thresholds $t_1,t_2,t_3$ admit, for every labeling, a
function clearing the margin (at least $t_i+\gamma$ where the label is $1$, at most
$t_i-\gamma$ where it is $0$).  The forbidden band of width $2\gamma$ makes the
notion scale-sensitive; the two curves realize $101$ and $010$ against a shared
threshold vector.  Finiteness for every $\gamma$ (\leanname{FatShatteringDim})
characterizes agnostic real-valued PAC learnability (Alon et al.\ 1997).}
\label{fig:fat-shattering}
\end{figure}

\begin{thm}[Real-Valued Characterization {\cite{AlonBenDavidCesaBianchiHaussler1997}}]
\label{thm:real-valued-char}
A class $\mathcal{F} \subseteq [0,1]^X$ of bounded real-valued functions is
agnostically PAC learnable with respect to the absolute loss if and only if
$\fatshat_\gamma(\mathcal{F}) < \infty$ for all $\gamma > 0$.
\end{thm}

\begin{proof}[Proof sketch]
\textbf{Sufficiency.}  If $\fatshat_\gamma(\mathcal{F}) = d_\gamma < \infty$
for all $\gamma > 0$, then the $L_2$ covering numbers satisfy
$\log \mathcal{N}_2(\varepsilon, \mathcal{F}, x_1^n) \leq
O(d_{\varepsilon/4} \cdot \log^2(n/d_{\varepsilon/4}))$.
Finite covering numbers imply uniform convergence of empirical risk to
true risk, which yields agnostic PAC learnability.

\textbf{Necessity.}  If $\fatshat_\gamma(\mathcal{F}) = \infty$ for some
$\gamma > 0$, then for every $n$, one can $\gamma$-fat-shatter $n$ points.
By a probabilistic argument, any empirical risk minimizer on a sample of
size $m$ has expected error $\Omega(\gamma)$ on a suitably chosen
distribution, regardless of $m$.  Hence $\mathcal{F}$ is not learnable
at accuracy $\gamma$.
\end{proof}

\begin{remark}[Scale families versus single numbers]
\label{rmk:scale-family}
The fat-shattering characterization is a \emph{family} of finiteness
requirements, one for each $\gamma > 0$, where a single threshold-free number
would have sufficed in the binary case.  This is
another instance of grammar growth: the real-valued setting requires a
scale parameter that has no analogue in the binary theory.  In practice,
$\fatshat_\gamma$ is often a decreasing function of $\gamma$; larger
margins are harder to maintain, so fewer points can be fat-shattered.
\end{remark}


% ================================================================
% SECTION 3: AGNOSTIC LEARNING
% ================================================================

\section{Agnostic Learning: Dropping Realizability}
\label{sec:agnostic-ext}

All results so far have assumed \emph{realizability}: there exists some
$h^* \in \mathcal{H}$ with $\risk_D(h^*) = 0$.  In practice, models are
always wrong; the question is how wrong.  The agnostic setting, introduced
by Haussler~\cite{Haussler1992}, removes the realizability assumption entirely.

\begin{defn}[Agnostic PAC Learning {\cite{Haussler1992}}]
\label{def:agnostic-pac-ext}
A hypothesis class $\mathcal{H}$ is \emph{agnostically PAC learnable}
if there exists a learner $A$ and a function $m \colon (0,1)^2 \to \N$ such
that for every distribution $D$ over $X \times \{0,1\}$ (with \emph{no}
assumption that the Bayes-optimal classifier lies in $\mathcal{H}$), for
every $\varepsilon, \delta > 0$: given $m(\varepsilon, \delta)$ i.i.d.\
samples from $D$,
\[
  \Pr_{S \sim D^m}\!\Bigl[
    \risk_D\!\bigl(A(S)\bigr)
    \;\leq\;
    \inf_{h \in \mathcal{H}} \risk_D(h) + \varepsilon
  \Bigr] \;\geq\; 1 - \delta.
\]
\end{defn}

The learner must compete with the \emph{best hypothesis in the class},
not with the Bayes-optimal classifier.  The gap
$\inf_{h \in \mathcal{H}} \risk_D(h)$ may be positive; the learner need
only close to within $\varepsilon$ of it.

\begin{thm}[Agnostic Fundamental Theorem]
\label{thm:agnostic-ft}
For binary classification with $0$-$1$ loss, a hypothesis class
$\mathcal{H} \subseteq \{0,1\}^X$ is agnostically PAC learnable if and only
if $\VC(\mathcal{H}) < \infty$.
\end{thm}

\begin{proof}[Proof sketch]
The characterizing condition is the same as in the realizable case, finite
VC dimension, but the proof mechanism and sample complexity differ.

\medskip\noindent
\textbf{Sufficiency.}  Suppose $\VC(\mathcal{H}) = d < \infty$.  The learner
runs empirical risk minimization: $A(S) = \arg\min_{h \in \mathcal{H}}
\emprisk_S(h)$.  By the uniform convergence theorem for VC classes, with
probability at least $1 - \delta$ over a sample of size
$m = O\!\bigl(\frac{d + \log(1/\delta)}{\varepsilon^2}\bigr)$,
\[
  \sup_{h \in \mathcal{H}} \bigl|\risk_D(h) - \emprisk_S(h)\bigr|
  \;\leq\; \varepsilon/2.
\]
Let $h^* = \arg\min_{h \in \mathcal{H}} \risk_D(h)$ and $\hat{h} = A(S)$.
Then:
\begin{align*}
  \risk_D(\hat{h})
  &\leq \emprisk_S(\hat{h}) + \varepsilon/2
    &\text{(uniform convergence)}\\
  &\leq \emprisk_S(h^*) + \varepsilon/2
    &\text{(ERM chose $\hat{h}$)}\\
  &\leq \risk_D(h^*) + \varepsilon/2 + \varepsilon/2
    &\text{(uniform convergence)}\\
  &= \risk_D(h^*) + \varepsilon.
\end{align*}

\medskip\noindent
\textbf{Necessity.}  If $\VC(\mathcal{H}) = \infty$, the No-Free-Lunch
argument of the realizable case applies \textit{a fortiori}: for any learner
and any sample size $m$, there exists a realizable distribution (which is
a special case of an arbitrary distribution) on which the learner fails.
\end{proof}

\begin{separation}
\textbf{Agnostic versus realizable sample complexity.}

\medskip\noindent
The agnostic and realizable settings are characterized by the \emph{same}
condition ($\VC < \infty$), but the sample complexities differ:
\[
  m_{\mathrm{realizable}}(\varepsilon, \delta)
  = \Theta\!\left(\frac{d}{\varepsilon} + \frac{\log(1/\delta)}{\varepsilon}\right),
  \qquad
  m_{\mathrm{agnostic}}(\varepsilon, \delta)
  = \Theta\!\left(\frac{d}{\varepsilon^2} + \frac{\log(1/\delta)}{\varepsilon^2}\right).
\]
The quadratic dependence on $1/\varepsilon$ in the agnostic case is inherent: a
matching lower bound forces it, so no sharper analysis can remove it.  The bound is achieved by
distributions where the optimal hypothesis has error $1/2 - \varepsilon$,
and distinguishing it from a hypothesis with error $1/2$ requires
$\Omega(1/\varepsilon^2)$ samples by a Chernoff-bound argument.

\medskip\noindent
\emph{Graph edge.}  \nde{agnostic\_pac\_learning} restricts \nde{pac\_learning}
in the graph: agnostic PAC learning \emph{generalizes} standard PAC learning by
removing the realizability constraint.
\end{separation}


% ================================================================
% SECTION 4: NOISE AND PARTIAL CONCEPTS
% ================================================================

\section{Noise-Tolerant and Partial Concept Learning}
\label{sec:noise-partial}

\subsection{Classification Noise}

The \emph{classification noise} model (Angluin and
Laird~\cite{AngluinLaird1988}) modifies PAC learning by flipping each
label independently with probability $\eta < 1/2$: the learner observes
$(x, y \oplus z)$ where $z \sim \mathrm{Bernoulli}(\eta)$ and $y$ is the true
label.  The question is whether finite VC dimension still suffices for
learning.

\begin{thm}[Noise-tolerant PAC learning]
\label{thm:noise-tolerant}
If $\VC(\mathcal{H}) = d < \infty$ and the noise rate $\eta < 1/2$ is known,
then $\mathcal{H}$ is PAC learnable under classification noise with sample
complexity
\[
  m = O\!\left(\frac{d}{(1-2\eta)^2 \varepsilon^2}
  + \frac{\log(1/\delta)}{(1-2\eta)^2 \varepsilon^2}\right).
\]
\end{thm}

The bound is an upper bound, achieved by a noise-robust empirical risk
minimizer~\cite{AngluinLaird1988}; the $(1-2\eta)^{-2}$ dependence is also
necessary, so the rate is tight.  The sample complexity degrades as
$\eta \to 1/2$, where the noise rate approaches the information-theoretic
limit at which labels carry no signal.

\begin{remark}[SQ connection]
\label{rmk:sq-noise}
Statistical query (SQ) learnable classes are automatically noise-tolerant:
an SQ algorithm estimates expectations $\E_D[\phi(x, y)]$ to tolerance
$\tau$, and these can be simulated from noisy samples with a $1/(1-2\eta)$
overhead in sample size.  However, not all noise-tolerant classes are
SQ learnable: parities are PAC learnable (via Gaussian elimination) and
hence noise-tolerant for known $\eta$, but not SQ learnable
($\SQdim = 2^n$).
\end{remark}


\subsection{Partial Concept Learning}

In \emph{partial concept learning} (also called \emph{promise problems}),
the hypothesis class $\mathcal{H}$ is defined only on a subset
$P \subseteq X$, the \emph{promise set}.  Points outside $P$ have no
defined label, and the distribution $D$ is supported on $P$.

\begin{defn}[Partial Concept Class]
\label{def:partial-concept}
A \emph{partial concept class} is a set $\mathcal{H} \subseteq \{0, 1, *\}^X$
where $*$ denotes ``undefined.''  A distribution $D$ is \emph{compatible}
with $h \in \mathcal{H}$ if $D$ is supported on $\{x : h(x) \neq *\}$.
\end{defn}

Partial concepts arise naturally in many settings: medical diagnosis where
only symptomatic patients have definite diagnoses, or feature spaces where
certain regions are ``don't care'' zones.

Alon, Hanneke, Holzman, and Moran~\cite{AlonHannekeHolzmanMoran2021}
developed the PAC theory of partial concept classes, and the picture is
genuinely new.  A partial class is PAC learnable if and only if its
(partial) VC dimension is finite, where shattering is restricted to the
promise region; so far, the binary analogue carries over.  What does
\emph{not} carry over is the reduction to total concepts: they exhibit
partial classes of VC dimension~$1$ that admit \emph{no} total extension of
finite VC dimension, and that cannot be learned by empirical risk
minimization at all.  Partial-concept learning is therefore not a special
case of the total theory; it is a setting where proper and ERM-based
learning provably fail while improper learning succeeds, and where the
computational landscape shifts: problems easy as total concepts can become
hard under a promise constraint.


% ================================================================
% SECTION 5: PROPER VS IMPROPER
% ================================================================

\section{Proper Versus Improper Learning}
\label{sec:proper-improper}

All PAC definitions permit the learner to output \emph{any} efficiently
evaluatable hypothesis, not just one from the target class $\mathcal{H}$.
When the learner is restricted to output a hypothesis from $\mathcal{H}$
itself, we call it a \emph{proper} learner.

\begin{defn}[Proper and Improper PAC Learning]
\label{def:proper-improper-ext}
A PAC learner $A$ for class $\mathcal{H}$ is \emph{proper} if
$A(S) \in \mathcal{H}$ for all samples $S$.  It is \emph{improper} if
$A(S)$ may lie in a larger class $\mathcal{H}' \supseteq \mathcal{H}$.
\end{defn}

\begin{remark}
In the information-theoretic (sample complexity) sense, there is no
separation between proper and improper learning for standard PAC: both are
characterized by finite VC dimension, and the sample complexities are
$\Theta(d/\varepsilon)$ in both cases.  The separation is \emph{computational}.
\end{remark}

\begin{thm}[Pitt--Valiant {\cite{PittValiant1988}}]
\label{thm:pitt-valiant}
For $k \geq 2$, $k$-term DNF is not properly PAC learnable unless $\mathrm{RP} = \mathrm{NP}$.
\end{thm}

\begin{proof}[Proof]
The proof reduces graph $k$-colorability to the problem of finding a
consistent $k$-term DNF formula.

Let $G = (V, E)$ be a graph with $V = \{v_1, \ldots, v_n\}$.  We construct
a learning instance over $X = \{0,1\}^n$ as follows.

\medskip\noindent
\textbf{Positive examples.}  For each vertex $v_i$, create the example
$e_i^+ \in \{0,1\}^n$ defined by $e_i^+(j) = 1$ if $j \neq i$ and
$e_i^+(i) = 0$.  Label: positive.

\medskip\noindent
\textbf{Negative examples.}  For each edge $(v_i, v_j) \in E$, create the
example $e_{ij}^- \in \{0,1\}^n$ defined by $e_{ij}^-(l) = 1$ if
$l \notin \{i, j\}$, and $e_{ij}^-(i) = e_{ij}^-(j) = 0$.  Label: negative.

\medskip\noindent
\textbf{Claim.}  A consistent $k$-term DNF exists for this labeled sample
if and only if $G$ is $k$-colorable.

\medskip\noindent
\emph{Forward direction.}  Given a proper $k$-coloring $c \colon V \to [k]$,
define the $k$-term DNF $\phi = T_1 \vee \cdots \vee T_k$ where
$T_j = \bigwedge_{i : c(v_i) = j} \overline{x_i}$.  Each positive example
$e_i^+$ satisfies $T_{c(v_i)}$ (since $e_i^+(i) = 0$ makes the literal
$\overline{x_i}$ true, and all other literals in $T_{c(v_i)}$ are true
because $e_i^+(l) = 1$ for $l \neq i$).  Each negative example $e_{ij}^-$
falsifies every term: for any term $T_j$, since $(v_i, v_j)$ is an edge
and $c$ is proper, $c(v_i) \neq c(v_j)$, so at most one of $v_i, v_j$
is in color class $j$, and the other has $e_{ij}^-(l) = 0$ which falsifies
a literal not in $T_j$.  (A careful case analysis confirms the argument.)

\medskip\noindent
\emph{Reverse direction.}  Given a consistent $k$-term DNF
$\phi = T_1 \vee \cdots \vee T_k$, assign color $c(v_i) = \min\{j : e_i^+
\text{ satisfies } T_j\}$.  Since $e_i^+$ is positive, at least one term is
satisfied.  For any edge $(v_i, v_j)$, the negative example $e_{ij}^-$
falsifies $\phi$, which forces $c(v_i) \neq c(v_j)$ (otherwise the term
satisfied by both would also satisfy $e_{ij}^-$).  Hence $c$ is a proper
$k$-coloring.

\medskip\noindent
\textbf{Conclusion.}  A proper PAC learner for $k$-term DNF, given the
above examples (which can be generated efficiently from $G$), would find
a consistent $k$-term DNF in polynomial time, thereby solving $k$-colorability
in polynomial time.  Since $k$-colorability is NP-complete for $k \geq 3$
(and NP-hard to decide in randomized polynomial time unless
$\mathrm{RP} = \mathrm{NP}$), proper PAC learning of $k$-term DNF is
computationally hard.
\end{proof}


\begin{separation}
\textbf{Proper $\not\Rightarrow$ efficient learning: the DNF witness.}

\medskip\noindent
The same class that is hard to learn properly is easy to learn improperly.
Observe that every $k$-term DNF over $n$ variables is a $k$-CNF formula
(by distributing), and $k$-CNF is learnable in polynomial time: simply
enumerate all $O(n^k)$ candidate clauses and filter by consistency with the
data.

\medskip\noindent
\emph{Summary.} For $k = 3$: 3-term DNF is NP-hard proper, polynomial-time
improper (via 3-CNF).  This is the cleanest known computational separation
in PAC learning.

\medskip\noindent
\emph{Graph edge.}  The \nde{proper\_improper\_separation} does not imply
\nde{pac\_learning}, witnessed by $3$-term DNF: NP-hard proper, polynomial-time
improper.
\end{separation}


\begin{figure}[ht]
\centering
\begin{tikzpicture}[
  >={Stealth[length=5pt]},
  % cells: fill = learner type (improper green / proper red);
  %        border = evidence (solid = machine-checked / dashed = literature)
  cell/.style={rectangle, rounded corners=3pt, draw, thick,
    minimum width=3.5cm, minimum height=1.3cm, font=\small, align=center, inner sep=3pt},
  vcell/.style={cell},                 % solid border = kernel-verified
  lcell/.style={cell, dashed},         % dashed border = literature
  impf/.style={fill=sepgreen!12},
  propf/.style={fill=thmred!12},
  hd/.style={font=\small\bfseries\sffamily, align=center},
  rl/.style={font=\small\bfseries\sffamily, align=center, text width=2.5cm},
  verdict/.style={font=\scriptsize\itshape, text=notegray, align=center},
  rel/.style={-{Stealth[length=5pt]}, dashed, notegray, thick}
]
% --- column headers ---
\node[hd, text=sepgreen] (hi) at (0,2.05)   {Improper\\[-1pt]{\scriptsize $h \in \mathcal{H}'$}};
\node[hd, text=thmred]   (hp) at (4.7,2.05)  {Proper\\[-1pt]{\scriptsize $h \in \mathcal{H}$}};

% --- row 1: STATISTICAL layer (SOLID border = machine-checked) ---
\node[vcell, impf]  (si) at (0,0.65)   {$\Theta(d/\varepsilon)$ samples};
\node[vcell, propf] (sp) at (4.7,0.65)  {$\Theta(d/\varepsilon)$ samples};

% --- row 2: COMPUTATIONAL layer (DASHED border = literature) ---
\node[lcell, impf]  (ci) at (0,-1.35)  {Polynomial time\\{\footnotesize (via $k$-CNF)}};
\node[lcell, propf] (cp) at (4.7,-1.35) {NP-hard\\{\footnotesize (via $k$-colouring)}};

% --- row labels (left) ---
\node[rl, text=defblue]    at (-2.95,0.65)  {Statistical\\[-1pt]{\scriptsize (machine-checked)}};
\node[rl, text=boundbrown] at (-2.95,-1.35) {Computational\\[-1pt]{\scriptsize (literature)}};

% --- verdicts (right) ---
\node[verdict] at (7.3,0.65)  {same:\\no gap};
\node[verdict] at (7.3,-1.35) {differs:\\the gap};

% --- restriction relation (row 1, between columns) ---
\draw[rel] (sp.west) -- node[above, font=\scriptsize\itshape, text=notegray] {$\subseteq\mathcal{H}$} (si.east);

% --- witness (bottom) ---
\node[font=\scriptsize, align=center, text width=10cm] at (2.35,-2.75)
  {\textbf{Witness.} $3$-term DNF is one class in all four cells: the same statistically,
   yet NP-hard to learn properly and polynomial-time improperly (as $3$-CNF).};
\end{tikzpicture}
\caption[Proper versus improper learning: no statistical gap, a computational
gap]{Proper versus improper PAC learning across two layers.  In the statistical
layer (top row) there is no separation: proper and improper learning are both
characterized by finite VC dimension and both need $\Theta(d/\varepsilon)$ samples
(\leanname{fundamental_theorem}).  The gap appears only in the computational layer
(bottom row), where $3$-term DNF is NP-hard properly (Pitt and Valiant 1988,
\Cref{thm:pitt-valiant}) yet polynomial-time improperly as $3$-CNF.}
\label{fig:proper-improper}
\end{figure}


As an impossibility result, the \nde{proper\_improper\_separation} both
lower-bounds \nde{pac\_learning} and fails to imply it.  Its provenance is
Pitt--Valiant 1988, and the proof is a direct reduction from graph coloring to
consistent hypothesis finding.


% ================================================================
% SECTION 6: THE EXTENDS_GRAMMAR PATTERN
% ================================================================

\section{The \rel{extends\_grammar} Pattern}
\label{sec:extends-grammar}

The extensions of the previous sections share a common structure.  Each
begins with the VC-dimension framework and introduces new primitives that
have no analogue in binary classification:

\begin{center}\small
\begin{tabular}{llp{5.5cm}}
\toprule
\textbf{Extension} & \textbf{New Primitive} & \textbf{Why Needed} \\
\midrule
Natarajan dim. & Two-function shattering $(f_0, f_1)$
  & With $k > 2$ labels, a single bit string does not
    describe a labeling. \\[4pt]
DS dimension & One-inclusion witness
  & Global hypergraph structure governs learnability
    when $|Y| = \infty$. \\[4pt]
Pseudodimension & Threshold vector $(t_1, \ldots, t_d)$
  & Real-valued outputs must be discretized before
    shattering can be defined. \\[4pt]
Fat-shattering dim. & Scale parameter $\gamma$
  & Agnostic real-valued learning requires margin;
    finiteness at all scales replaces a single
    finiteness condition. \\[4pt]
Agnostic PAC & Benchmark $\inf_{h \in \mathcal{H}} \risk(h)$
  & Without realizability, success is measured relative
    to the best in class. \\[4pt]
Noise model & Noise rate $\eta$
  & Labels are corrupted; sample complexity depends on
    the signal-to-noise ratio $(1-2\eta)^{-2}$. \\
\bottomrule
\end{tabular}
\end{center}

In each case, the extension is an \rel{extends\_grammar} edge in the concept
graph rather than a \rel{restricts} edge.  The distinction matters:

\begin{itemize}[leftmargin=*,itemsep=3pt]
\item A \rel{restricts} edge means the target is a special case of the source
  (no new vocabulary needed).
\item An \rel{extends\_grammar} edge means the source introduces
  \emph{new vocabulary}: concepts that cannot be expressed in the target's
  language.  The extension is irreducible.
\end{itemize}

\begin{remark}[The non-reduction principle]
\label{rmk:non-reduction}
A tempting strategy for multiclass or real-valued learning is to reduce to
binary classification (e.g., one-vs-all for multiclass, or
thresholding for real-valued).  Such reductions are computationally useful but
do not eliminate the need for new dimensions.  The VC dimension of the
reduced class does not generally equal the Natarajan, DS, pseudo-, or
fat-shattering dimension of the original class.  The reductions introduce
approximation gaps that the native dimensions avoid.  This is why the
\rel{extends\_grammar} edges exist: they mark the points where reduction to the
binary theory loses information.
\end{remark}


\begin{figure}[ht]
\centering
\begin{tikzpicture}[
  scale=0.85,
  level/.style={font=\small\sffamily, align=center},
  box/.style={
    rectangle, rounded corners=3pt, draw=defblue, thick,
    fill=defblue!5, minimum width=2.8cm, minimum height=0.8cm,
    font=\small, align=center
  },
  hubbox/.style={
    rectangle, rounded corners=3pt, draw=defblue, very thick,
    fill=defblue!12, minimum width=2.8cm, minimum height=0.8cm,
    font=\small, align=center
  },
  extarr/.style={-{Stealth[length=5pt]}, edgeorange, thick,
    decorate, decoration={snake, amplitude=1pt, segment length=6pt}}
]

% Binary core (the one machine-checked rung: fundamental_theorem, sc=0)
\node[hubbox] (bin) at (0,0) {Binary PAC\\$\VC(\mathcal{H}) < \infty$};
\node[font=\tiny\sffamily, text=defblue, anchor=north]
  at ([yshift=-2pt]bin.south) {\checkmark\,machine-checked};

% Extensions (literature characterizations; dimensions formalized as kernel defs)
\node[box] (mc) at (-4, 2.5) {Multiclass\\$\DSdim(\mathcal{H}) < \infty$};
\node[box] (rv) at (4, 2.5) {Real-valued\\$\fatshat_\gamma < \infty\;\forall\gamma$};
\node[box] (ag) at (-4, -2.5) {Agnostic\\$\VC(\mathcal{H}) < \infty$\\(same condition)};
\node[box] (no) at (4, -2.5) {Noise-tolerant\\$\VC(\mathcal{H}) < \infty$\\$+ \eta < 1/2$};

% Grammar growth arrows (each adds an irreducible primitive)
\draw[extarr] (bin) -- node[left, font=\tiny\itshape, text=edgeorange, align=center]
  {$+$ one-inclusion\\witness} (mc);
\draw[extarr] (bin) -- node[right, font=\tiny\itshape, text=edgeorange, align=center]
  {$+$ thresholds\\$+$ margins} (rv);
\draw[extarr] (bin) -- node[left, font=\tiny\itshape, text=edgeorange, align=center]
  {$+$ relative\\benchmark} (ag);
\draw[extarr] (bin) -- node[right, font=\tiny\itshape, text=edgeorange, align=center]
  {$+$ noise rate\\$\eta$} (no);

\end{tikzpicture}
\caption[Grammar growth from binary PAC learning]{Binary PAC learning sits at the
hub, characterized by finite VC dimension (\leanname{fundamental_theorem}).  Each
wavy arrow is an \rel{extends\_grammar} edge that grows the binary theory by
adding a primitive with no binary analogue: a one-inclusion witness for
multiclass learning, thresholds then a margin for real-valued learning, a relative
benchmark for agnostic learning, and a noise rate $\eta$.  The extension is
irreducible because reduction to the binary theory loses information
(\Cref{rmk:non-reduction}), proved in the multiclass case by the class with
$\Ndim=1$ and $\DSdim=\infty$ (\Cref{prop:ds-vs-nat}).}
\label{fig:grammar-growth}
\end{figure}


The grammar growth \emph{is} the extension.  In the multiclass case this
necessity is a theorem: the class with $\Ndim = 1$ and $\DSdim = \infty$ shows
that no two-function dimension can characterize learnability when $|Y| = \infty$.
For the real-valued, agnostic, and noisy settings no simpler characterization is
known, and the pattern is the same: each new primitive expresses what the existing
one cannot.


\subsection{Where Transformers Live}
\label{sec:transformer-grammar}

The extensions of this chapter are not abstractions a practitioner can avoid;
a transformer inhabits several at once.  Its next-token head is a
\emph{multiclass} hypothesis class over the vocabulary (tens of thousands of
labels, and unbounded in the limit of growing vocabularies), so it sits in
exactly the regime where the Natarajan dimension fails and the DS dimension is
the right measure.  Its logits are \emph{real-valued}, so its capacity is
governed by a scale-sensitive dimension; the certified generalization bound of
\Cref{ch:generalization} is, in this language, a fat-shattering and
covering-number control of the attention head.  And training is
\emph{agnostic}: no realizable $h^\star$ is assumed.

A softmax-attention head induces a concept class that is
well-behaved,\formal{TLT.Bridge.softAttention_wellBehaved} hence of finite VC
dimension in the binary setting.  Whether that control lifts to the multiclass
next-token head, that is, whether the DS dimension of an attention class over a
large label set is finite, is open (\Cref{open:attention-ds}).


% ================================================================
% OPEN PROBLEMS
% ================================================================

\section{Open Problems}
\label{sec:extensions-open}

Extending a characterization is where learning theory most often stumbles: a
dimension that works for binary classification need not, and the multiclass
history is the cautionary tale.  Each problem below is tagged a \emph{known
unknown} (KU), sharply posed but unresolved, or an \emph{unknown known} (UK),
where the right object or its formalization is still missing; a design-lab or
literature anchor is named where one exists.

\begin{openprob}[Tight multiclass sample complexity]
\label{open:multiclass-tight}
For finite $|Y| = k$, the best known bounds leave a $\Theta(\log k)$ gap
between the Natarajan-dimension-based upper and lower bounds on sample
complexity.  Does the DS dimension close this gap, yielding a sample
complexity of $\Theta(\DSdim / \varepsilon)$ without logarithmic factors?
\end{openprob}

\begin{openprob}[Multiclass online learning]
\label{open:multiclass-online}
The DS dimension characterizes PAC learnability in the multiclass setting.
What is the correct characterization of \emph{online} multiclass learnability?
Is there a multiclass analogue of the Littlestone dimension, and does it
relate to the DS dimension as the Littlestone dimension relates to VC?
\end{openprob}

\begin{openprob}[Formalizing the DS dimension]
\label{open:formalize-ds}
The Natarajan, pseudo-, and fat-shattering dimensions are formalized
(\leanname{NatarajanDim}, \leanname{Pseudodimension}, \leanname{FatShatteringDim}).
The DS dimension, defined by a default labeling perturbed by single-coordinate
deviations, satisfies $\Ndim \le \DSdim$ and characterizes multiclass
learnability: $\DSdim(\mathcal{H}) < \infty$ iff $\mathcal{H}$ is PAC learnable
(\cite{BrukhimCarmonDinurMoranYehudayoff2022}).  Formalize this dimension and
the characterization.
\end{openprob}

\begin{openprob}[The DS dimension of attention]
\label{open:attention-ds}
A softmax-attention head is verified well-behaved, hence finite-VC, in the
binary setting (\leanname{softAttention_wellBehaved}).  Is the DS dimension of
the multiclass next-token attention class (over a vocabulary of size $k$, and
as $k \to \infty$) finite?  A known unknown on the transformer side, and the
multiclass form of this chapter's central question.
\end{openprob}

\begin{openprob}[The fat-shattering dimension of attention]
\label{open:attention-fat}
The certified generalization bound for attention (\Cref{ch:generalization}) is
a covering-number control of real-valued logits.  Does it match the
fat-shattering dimension (\leanname{FatShatteringDim}) of the attention logit
class, and is that dimension finite at every scale $\gamma$?  A known unknown on
the transformer side.
\end{openprob}

\begin{openprob}[A realizable characterization via pseudodimension]
\label{open:pdim-realizable}
\Cref{prop:pdim-covering} bounds covering numbers, and
\Cref{thm:real-valued-char} characterizes the \emph{agnostic} case via
fat-shattering.  Is finiteness of the pseudodimension (\leanname{Pseudodimension})
exactly the characterization of \emph{realizable} real-valued learnability, and
can it be proved as a clean biconditional?  A known unknown.
\end{openprob}

\begin{openprob}[Formalizing partial-concept learnability]
\label{open:partial-formal}
The partial VC dimension characterizes partial-concept PAC learnability, with
the surprise that learnable partial classes need not extend to finite-VC total
classes~\cite{AlonHannekeHolzmanMoran2021}.  Can this be formalized, including
the no-finite-extension witness?  An unknown known: the partial-concept model
is absent from the kernel.
\end{openprob}

\begin{openprob}[List learning]
\label{open:list-learning}
A list learner outputs a short list of labels and succeeds if the true label
is among them.  What combinatorial dimension characterizes list PAC
learnability, and how does it relate to the DS
dimension~\cite{BrukhimCarmonDinurMoranYehudayoff2022}?  A known unknown.
\end{openprob}

\begin{openprob}[Agnostic proper versus improper]
\label{open:agnostic-proper}
In the realizable case, proper and improper PAC have the same sample
complexity (\Cref{sec:proper-improper}).  In the \emph{agnostic} case, is there
a class for which proper learning needs strictly more samples than improper
learning, an information-theoretic rather than merely computational separation?  A
known unknown.
\end{openprob}

\begin{openprob}[Noise-tolerant real-valued learning]
\label{open:noise-fat}
Classification noise tolerance is understood for binary VC classes
(\Cref{thm:noise-tolerant}).  What is the fat-shattering analogue: how does
label noise on real-valued targets degrade the scale-sensitive sample
complexity?  A known unknown.
\end{openprob}

\subsection*{Counterfactual exercises}

These exercises generalize the chapter's figures, each widening a concrete picture
into the chapter's argument that binary PAC theory extends only by grammar growth.

\begin{enumerate}[leftmargin=*,itemsep=3pt]
\item \textbf{The one verified rung.} The fundamental theorem, that finite VC
dimension characterizes realizable binary PAC learnability over $\mathcal{H}
\subseteq \{0,1\}^X$ (\leanname{fundamental_theorem}), is the single edge of the
concept map with a machine-checked proof.  Argue that every other characterization
in the map, multiclass, real-valued, agnostic, and proper-versus-improper, is a
literature result that grows outward from this one rung.

\item \textbf{Grammar growth as the thesis.} Read the \rel{extends\_grammar} edges
as one class.  Each marks a place where the binary language loses information and a
new primitive must be added: the pair $(f_0, f_1)$ for Natarajan
(\Cref{def:natarajan-dim}), the threshold vector then the margin $\gamma$ for the
real-valued branch (\Cref{def:pseudodim-ext}, \Cref{def:fat-shattering-ext}), and
the noise rate $\eta$ (\Cref{thm:noise-tolerant}).  Show that no reduction to the
binary theory (\Cref{rmk:non-reduction}) collapses these into VC dimension.

\item \textbf{Table and star are one object.} In Figure~\ref{fig:ds-shattering},
show that a witness table of $d$ rows, each deviating from the default $f$ at
exactly one coordinate, is the same data as a $d$-edge star out of $f$ in the
one-inclusion hypergraph $\mathrm{OIG}(\mathcal H, S)$ (\Cref{def:oig}).  Conclude
that generalizing to arbitrary $d$ means drawing a $d$-edge star.

\item \textbf{Weaker per point, larger dimension.} Using \Cref{prop:ds-vs-nat},
show every N-shattered set is DS-shattered (take $f = f_0$ and let $h_i$ realize
the basis string $e_i$), so $\Ndim(\mathcal H) \le \DSdim(\mathcal H)$.  Then
describe the Brukhim et al.\ (2022) class with $\Ndim = 1$ yet $\DSdim = \infty$,
and explain why local two-label structure fails to bound learnability when
$|Y| = \infty$.

\item \textbf{The margin is the scale primitive.} In
Figure~\ref{fig:fat-shattering}, show the symmetric band $(t_i-\gamma, t_i+\gamma)$
is the one object with no binary analogue.  Shrink it to zero width and show
$\gamma$-fat-shattering degenerates to the $P$-shattering condition of
\Cref{def:pseudodim-ext}, so $\fatshat_\gamma(\mathcal F) \to \Pdim(\mathcal F)$ as
$\gamma \to 0$.  Using \Cref{rmk:scale-family}, argue the characterizing object is
a scale family $\{\fatshat_\gamma\}_{\gamma > 0}$, one finiteness requirement per
scale.

\item \textbf{The proper-improper gap lives in the computational layer.} Trace
$3$-term DNF through Figure~\ref{fig:proper-improper}: learnable both ways
statistically with $\Theta(d/\varepsilon)$ samples, yet NP-hard properly and
polynomial-time improperly as $3$-CNF (\Cref{thm:pitt-valiant}).  Hold the columns
fixed and show that toggling only the layer turns a $\Theta(1)$ proper-to-improper
cost ratio into a superpolynomial one unless $\mathrm{RP} = \mathrm{NP}$.

\item \textbf{Is grammar growth a provable no-single-dimension invariant?} The
chapter proves grammar growth \emph{necessary} only for multiclass, via the
$\Ndim = 1$, $\DSdim = \infty$ class (\Cref{prop:ds-vs-nat}); for the real-valued,
agnostic, and noisy settings no simpler characterization is known.  Conjecture that
no single combinatorial dimension characterizes all extensions at once, a
no-free-lunch for shattering primitives, and identify what a proof would need
beyond the one multiclass instance.
\end{enumerate}

\bigskip

This chapter has traced the grammar growth required to extend binary PAC
learning theory to multiclass, real-valued, agnostic, noise-tolerant, and
computationally constrained settings.  Each extension introduces irreducible
new primitives and yields its own characterization theorem or separation
result.  The next chapter examines what lies beyond these
characterizations: the application-layer concepts, the open problems, and the
scope boundaries of the theory.

% Chapter 18: Frontiers and Open Problems
% Theme: WHAT WE DO NOT KNOW: question-driven structure.
% Each section leads with an open problem, shows what is known,
% identifies the gap, and asks what would close it.
% Concepts: compression_scheme, margin_theory, kolmogorov_complexity,
% algorithmic_probability, universal_learning, computational_hardness,
% differential_privacy_learning, active_learning, transfer/meta/continual.

\chapter{Frontiers and Open Problems}
\label{ch:frontiers}

Every preceding chapter has presented established results: definitions,
characterization theorems, separation witnesses, and tight bounds.  This
chapter presents none of these.  Its subject is the boundary of what we know,
and its organizing principle is the open question rather than the proven
theorem.

Each section leads with a question to which no complete answer is known,
presents the best partial results, identifies the precise gap between what is
known and what is conjectured, and asks what kind of argument might close it.

The problems chosen are those that
connect most tightly to the concept graph of this book: the compression
conjecture (Chapter~\ref{ch:compression}), deep network generalization
(Chapter~\ref{ch:generalization}), universal learning beyond countable classes
(Chapter~\ref{ch:universal}), computational--statistical tradeoffs
(Chapter~\ref{ch:computational}), and the role of Kolmogorov complexity in
learning.  Other important open problems (online multiclass learnability,
private agnostic learning, quantum sample complexity gaps) are noted but not
developed.

\medskip

The chapter has five sections.

\begin{enumerate}[leftmargin=*,itemsep=2pt]
\item \textbf{The Compression Conjecture} (\Cref{sec:compression-conjecture}):
  does $\VC(\mathcal{H}) = d$ imply a compression scheme of size $O(d)$?
\item \textbf{Deep Learning and Generalization}
  (\Cref{sec:deep-generalization}): why do overparameterized networks
  generalize?
\item \textbf{Universal Learning Beyond Countable Classes}
  (\Cref{sec:universal-frontiers}): does the trichotomy extend?
\item \textbf{Computational--Statistical Tradeoffs}
  (\Cref{sec:comp-stat-tradeoffs}): when does computational hardness prevent
  statistically optimal learning?
\item \textbf{Kolmogorov Complexity and the Ideal Learner}
  (\Cref{sec:kolmogorov-learning}): can the uncomputable yield computable
  insight?
\end{enumerate}


% ================================================================
% SECTION 1: THE COMPRESSION CONJECTURE
% ================================================================

\section{The Compression Conjecture}
\label{sec:compression-conjecture}

\begin{openprob}[Littlestone--Warmuth Compression Conjecture]
\label{op:compression-frontiers}
Let $\mathcal{H} \subseteq \{0,1\}^X$ with $\VC(\mathcal{H}) = d < \infty$.
Does there exist a sample compression scheme for $\mathcal{H}$ of size
$O(d)$?
\end{openprob}

This conjecture, stated by Littlestone and Warmuth in 1986, is among the
oldest major open problems in computational learning theory.  It asks whether
the sample complexity of PAC learning (which is $\Theta(d/\varepsilon)$) is
explained by compression: whether every learnable class can be learned by
storing only $O(d)$ examples from the training set and reconstructing a
consistent hypothesis from those examples alone.

\subsection{What is Known}

The positive direction is settled in a weak sense.

\begin{itemize}[leftmargin=*]
\item \textbf{Exponential compression exists.}
  Moran and Yehudayoff~\cite{MoranYehudayoff2016} proved that every class
  with $\VC(\mathcal{H}) = d$ admits a compression scheme of size
  $2^{O(d)}$.  The proof uses the Radon partition structure of finite VC
  classes, but the bound is exponentially larger than the conjecture predicts.

\item \textbf{Special cases are resolved.}  Maximum classes (those achieving
  the Sauer--Shelah bound $\growth_\mathcal{H}(n) = \binom{n}{\leq d}$)
  admit compression of size exactly $d$, as shown by Floyd and
  Warmuth~\cite{FloydWarmuth1995}.  Intersection-closed classes admit
  compression of size $d$~\cite{Helmbold2018}, and classes of small recursive
  teaching dimension compress to that dimension.

\item \textbf{Compression implies learning.}  The converse direction is
  established: a compression scheme of size $k$ yields PAC learning with
  sample complexity $O\bigl(\frac{k}{\varepsilon}\log\frac{k}{\varepsilon}
  + \frac{1}{\varepsilon}\log\frac{1}{\delta}\bigr)$.  Thus compression
  of size $O(d)$ would give an alternative proof of the fundamental theorem
  of PAC learning, with the added structural insight that the hypothesis
  depends on only $O(d)$ training points.
\end{itemize}

\begin{historical}
The compression conjecture predates the fundamental theorem of statistical
learning.  Littlestone and Warmuth formulated it in 1986, three years before
Blumer, Ehrenfeucht, Haussler, and Warmuth~\cite{BlumerEhrenfeuchtHausslerWarmuth1989}
completed the characterization of PAC learnability through VC dimension.
The conjecture was motivated by the observation that specific learnable
classes (halfspaces, rectangles, decision lists) all admit small compression
schemes, and the belief that this regularity was a property of finite-VC
classes as such.  Four decades later, the general case remains open.
\end{historical}

\subsection{Where the Gap Is}

The gap is between $O(d)$ and $2^{O(d)}$.  No class is known to require
compression size $\omega(d)$, but no proof technique is known to achieve
$o(2^d)$ in general.

The difficulty is structural.  The Moran--Yehudayoff proof proceeds through
a topological argument (the staircase lemma) that inherently produces
exponential-size compressions.  Improving the bound requires a fundamentally
different approach, one that exploits the combinatorial structure of VC
classes more directly.

Several partial strategies have been explored:

\begin{itemize}[leftmargin=*]
\item \textbf{Unlabeled compression.}  Allowing the compression to store
  unlabeled points (without their labels) does not help: the reconstruction
  map must still determine labels, and known reductions show that unlabeled
  compression of size $k$ implies labeled compression of size $O(k)$.

\item \textbf{Algebraic approaches.}  For classes definable by polynomial
  threshold functions, algebraic arguments yield compression of size $O(d)$.
  But this exploits algebraic structure absent from general VC classes.

\item \textbf{Recursive approaches.}  Decomposing a class into subclasses
  of smaller VC dimension and compressing each recursively is a natural
  strategy, but known decompositions (e.g., via Helly-type theorems) produce
  too many subclasses, losing the linear bound.
\end{itemize}

\subsection{What Would Close the Gap}

A resolution likely requires one of the following:

\begin{enumerate}[leftmargin=*]
\item A new structural theorem about VC classes that controls the geometry
  of their restrictions to finite sets more tightly than Sauer--Shelah.
\item A connection to a different branch of combinatorics (e.g., matroid
  theory or extremal set theory) that provides the right decomposition lemma.
\item A counterexample: a family of classes with VC dimension $d$ and
  no compression scheme of size $o(d^{1+c})$ for some $c > 0$.  Even a
  superlinear lower bound would be a breakthrough.
\end{enumerate}

\begin{remark}
The compression conjecture illustrates a recurring theme in this chapter:
the gap between \emph{existential} and \emph{constructive} knowledge.  We
know that every finite-VC-dimension class is PAC learnable (Chapter~\ref{ch:pac}).
We know that compression of size $2^{O(d)}$ exists
(Chapter~\ref{ch:compression}).  What we do not know is whether these two
facts are reflections of the same underlying structure at the right
quantitative scale.
\end{remark}


% ================================================================
% SECTION 2: DEEP LEARNING AND GENERALIZATION
% ================================================================

\section{Deep Learning and Generalization}
\label{sec:deep-generalization}

\begin{openprob}[Generalization in Overparameterized Networks]
\label{op:deep-gen}
Let $\mathcal{H}$ be the class of functions computed by a neural network
with $p \gg n$ parameters, trained by gradient descent on $n$ samples.
Why does the trained network generalize, i.e., why is $\risk(\hat{h})
\approx \emprisk(\hat{h})$, despite $\VC(\mathcal{H}) \geq p$?
\end{openprob}

This is among the most practically urgent open problems in learning theory.
Classical generalization bounds predict that a model with more parameters than
training examples should overfit catastrophically.  Deep networks routinely
violate this prediction.

\subsection{What is Known}

Several partial explanations exist, none fully satisfactory.

\begin{itemize}[leftmargin=*]
\item \textbf{Margin-based bounds.}  Bartlett, Foster, and
  Telgarsky~\cite{BartlettFosterTelgarsky2017} proved that for a
  depth-$L$ network with weight matrices $W_1, \ldots, W_L$, the
  generalization error is controlled by
  \[
    \frac{1}{n}\biggl(\frac{\prod_{i=1}^L \|W_i\|_\sigma}{\gamma}\biggr)^2
    \cdot \sum_{i=1}^L \frac{\|W_i\|_F^2}{\|W_i\|_\sigma^2}
  \]
  where $\gamma$ is the margin, $\|\cdot\|_\sigma$ is the spectral norm,
  and $\|\cdot\|_F$ is the Frobenius norm.  This bound is independent of the
  number of parameters $p$ and depends instead on the \emph{spectral
  complexity} of the learned weights.

\item \textbf{PAC-Bayes bounds.}  Treating the trained weights as a
  posterior distribution (e.g., by adding Gaussian noise), PAC-Bayes
  bounds (\Cref{ch:extensions}) can yield non-vacuous generalization
  certificates for specific trained networks.  Dziugaite and
  Roy~\cite{DziugaiteRoy2017} demonstrated this empirically.

\item \textbf{Algorithmic stability.}  SGD with bounded iterations satisfies
  uniform stability with $\beta = O(1/n)$ under convexity assumptions.
  For non-convex losses, Hardt, Recht, and Singer~\cite{HardtRechtSinger2016}
  showed that early stopping provides stability guarantees, but the bounds
  depend on the number of SGD steps and become vacuous for the training
  durations used in practice.

\item \textbf{Compression-based arguments.}  If the effective hypothesis
  after training can be compressed to $k \ll p$ bits, then compression
  bounds (Chapter~\ref{ch:compression}) apply.  Lottery ticket
  observations (that trained networks contain sparse subnetworks matching
  full performance) provide indirect evidence for compression, but no
  rigorous connection to sample complexity has been established.

\item \textbf{Machine-checked certificates.}  For the specific function class
  of a multi-head attention block, the generalization certificate of
  \Cref{ch:generalization} has been formalized: the executed
  (finite-precision) network's true risk exceeds its empirical risk by more
  than an explicit Dudley-capacity budget only on a sample event of
  probability at most~$\delta$.\formal{TLT.CertifiedGeneralization}  The budget
  is discharged from the layer norms and depth of the
  stack.\formal{TLT.config_capacity_law}  Like its hand-proved cousins, it
  grows with depth and norm and is not yet non-vacuous on large models.
\end{itemize}

\begin{computational}
\textbf{The double descent curve.}  Consider interpolation threshold: the
point where the number of parameters $p$ equals the number of training
examples $n$.  Classical theory predicts that test error is U-shaped
as a function of model complexity, with minimum at $p \approx n$.
Empirically, the test error \emph{decreases again} for $p \gg n$,
forming a double descent curve.  This was systematically documented
by Belkin et al.~\cite{BelkinHsuMaMandal2019}.

The double descent phenomenon is observed across several model families
(linear models, random features, deep networks)~\cite{BelkinHsuMaMandal2019}.
The classical complexity-based bounds (VC, Rademacher, and the parameter-count
bounds of \Cref{ch:generalization}) are monotone in model complexity and so
cannot predict the second descent; whether \emph{some} data-dependent bound can
is itself part of \Cref{op:deep-gen}.  The relevant notion of complexity for
overparameterized models is not a simple function of parameter count.
\end{computational}

\subsection{Where the Gap Is}

The gap has three aspects.

\begin{enumerate}[leftmargin=*]
\item \textbf{Quantitative vacuity.}  Most norm- and parameter-count-based
  bounds, when evaluated on practical networks, predict generalization error
  close to~$1$; that is, they are vacuous.  The PAC-Bayes certificates of
  Dziugaite and Roy~\cite{DziugaiteRoy2017} are a notable exception, giving
  non-vacuous numbers for specific small networks, but no bound is yet both
  non-vacuous and tight on large benchmark models.

\item \textbf{Architecture dependence.}  Existing bounds treat networks as
  generic function classes parameterized by weight norms.  They do not
  explain why certain architectures (e.g., ResNets, Transformers)
  generalize better than others with the same parameter count and weight
  norms.

\item \textbf{Optimization dependence.}  The generalization behavior of deep
  networks depends heavily on the optimization algorithm (SGD vs.\ Adam),
  learning rate schedule, batch size, and initialization.  No existing
  framework captures how the optimization trajectory selects hypotheses
  with low population risk from the vast set of interpolating solutions.
\end{enumerate}

\subsection{What Would Close the Gap}

The field needs a complexity measure that simultaneously accounts for the
architecture, the full optimization trajectory, and the data distribution.

Possible directions include:
\begin{itemize}[leftmargin=*]
\item \textbf{Implicit regularization theory:} proving that gradient descent
  on specific architectures converges to solutions with low complexity
  in a suitable measure (e.g., low rank, low spectral norm product, or
  low description length).
\item \textbf{Distribution-dependent bounds:} replacing uniform convergence
  (which bounds $\sup_D |\risk - \emprisk|$) with bounds that depend on the
  specific data distribution, potentially through conditional mutual
  information (Steinke--Zakynthinou, \Cref{ch:generalization}).
\item \textbf{Feature learning theory:} understanding how networks transform
  input representations during training, and why the learned features
  support generalization.
\end{itemize}


% ================================================================
% SECTION 3: UNIVERSAL LEARNING BEYOND COUNTABLE CLASSES
% ================================================================

\section{Universal Learning Beyond Countable Classes}
\label{sec:universal-frontiers}

\begin{openprob}[Universal Learning for Uncountable Classes]
\label{op:universal-uncountable}
The trichotomy theorem of Bousquet et
al.~\cite{BousquetHannekeMoranVanHandelYehudayoff2021} characterizes
universal learning rates for classes $\mathcal{H}$ with $|\mathcal{H}| \geq 3$.
Does an analogous characterization hold for uncountable classes of functions
$f \colon X \to [0,1]$ under general loss functions?
\end{openprob}

\subsection{What is Known}

The trichotomy theorem (\Cref{ch:universal}) establishes that for hypothesis
classes of size at least 3, the optimal universal learning rate is exactly one
of three types:
\begin{enumerate}[leftmargin=*,nosep]
\item Exponential ($e^{-n}$): if and only if the Littlestone dimension is
  finite.
\item Linear ($1/n$): if and only if the Littlestone dimension is infinite
  but no infinite VCL tree exists.
\item Arbitrarily slow: if and only if an infinite VCL tree exists.
\end{enumerate}
This characterization applies to binary classification with the $0$-$1$ loss.

For regression and general losses, partial results exist:
\begin{itemize}[leftmargin=*]
\item Hanneke, Moran, and others have extended the trichotomy to multiclass
  settings with finite label spaces, showing that the DS dimension
  (\Cref{ch:extensions}) plays a role analogous to the Littlestone dimension.
\item For real-valued regression under squared loss, the fat-shattering
  dimension at all scales characterizes PAC learnability, but the universal
  learning rate structure is not yet characterized.
\end{itemize}

\subsection{Where the Gap Is}

The gap has two components:

\begin{enumerate}[leftmargin=*]
\item \textbf{The loss function.}  The trichotomy proof relies on the
  discreteness of $0$-$1$ loss.  For continuous losses (squared loss,
  absolute loss, logistic loss), the proof technique does not directly
  apply.  The VCL tree, which is the combinatorial object governing the
  third regime, has no established analogue for continuous losses.

\item \textbf{The function class.}  For uncountable classes (e.g., all
  Lipschitz functions, all bounded-variation functions), measurability
  issues arise: the supremum of uncountably many random variables may not
  be measurable without additional regularity assumptions.  The universal
  learning framework requires that the learning rate hold for
  \emph{all} realizable distributions, and handling uncountable classes
  demands careful measure-theoretic foundations.
\end{enumerate}

\subsection{What Would Close the Gap}

Two ingredients appear necessary:
\begin{itemize}[leftmargin=*]
\item A combinatorial object for continuous losses that plays the role of
  the Littlestone tree and VCL tree in the binary case, a ``scale-sensitive
  game tree'' that captures the sequential complexity of the class at all
  accuracy levels.
\item A measurability framework (e.g., based on metric entropy or covering
  numbers) that allows the universal learning definition to extend cleanly
  to uncountable classes without requiring separability or other ad hoc
  assumptions.
\end{itemize}


% ================================================================
% SECTION 4: COMPUTATIONAL--STATISTICAL TRADEOFFS
% ================================================================

\section{Computational--Statistical Tradeoffs}
\label{sec:comp-stat-tradeoffs}

\begin{openprob}[Computational--Statistical Gap]
\label{op:comp-stat}
For which concept classes $\mathcal{C}$ does a gap exist between the
information-theoretic sample complexity of PAC learning $\mathcal{C}$ and
the sample complexity achievable by any polynomial-time algorithm?
Can such gaps be characterized combinatorially?
\end{openprob}

The results of Chapter~\ref{ch:computational} show that computational
hardness can prevent efficient PAC learning even when the VC dimension is
finite.  But those results are conditional: they assume the hardness of
specific cryptographic problems (factoring, LWE, discrete cube root).  The
deeper question is whether the computational--statistical gap is a structural
feature of learning or an artifact of our inability to prove $P \neq NP$.

\subsection{What is Known}

\begin{itemize}[leftmargin=*]
\item \textbf{Cryptographic hardness.}  Kearns and
  Valiant~\cite{KearnsValiant1994} showed that under the Discrete Cube Root
  Assumption, polynomial-size circuits are not PAC learnable in polynomial
  time, despite having polynomial VC dimension.  The gap is between
  $O(d/\varepsilon)$ samples (information-theoretically sufficient) and
  ``no polynomial-time algorithm suffices'' (computationally).

\item \textbf{Statistical query lower bounds.}  The SQ dimension provides
  \emph{unconditional} computational lower bounds within the SQ model
  (\Cref{ch:computational}).  Parities over $\{0,1\}^n$ have
  $\VC = n$ but $\SQdim = 2^n$: any SQ algorithm needs exponentially
  many queries.  Yet Gaussian elimination learns parities in polynomial
  time; it is simply not an SQ algorithm.  This shows that the SQ model
  captures one type of computational constraint but not all.

\item \textbf{Planted clique and related problems.}  A growing body of work
  studies problems where the information-theoretic threshold differs from the
  conjectured computational threshold: planted clique, sparse PCA, community
  detection.  For planted clique, a planted clique of size $\Theta(\sqrt{n})$
  is the conjectured threshold for polynomial-time detection, while one of
  size $\sim 2\log_2 n$ is already detectable information-theoretically.  None
  of these gaps is proven unconditionally; the evidence comes from restricted
  models (sum-of-squares, low-degree polynomials) rather than from
  unconditional separations.

\item \textbf{The barrier of Applebaum--Barak--Xiao.}  Applebaum, Barak, and
  Xiao~\cite{ApplebaumBarakXiao2008} showed that one cannot prove
  ``$\mathrm{P} \neq \mathrm{NP}$ implies hard PAC learning'' via standard
  (black-box) reductions.  Cryptographic or average-case assumptions are
  inherently necessary for hardness results in PAC learning.  This is a
  meta-theoretic constraint on the kind of hardness proofs that are possible.
\end{itemize}

\begin{historical}
The computational--statistical tradeoff emerged as a distinct research
direction in the 2010s, though its roots go back to Kearns and Valiant's
1994 cryptographic hardness results.  The key conceptual shift was the
recognition that sample complexity and computational complexity are not
independent axes of difficulty, but interact in subtle ways.  A problem
may have low sample complexity (few examples suffice) but high
computational complexity (no efficient algorithm can use those examples),
or vice versa.  The SQ model, introduced by Kearns~\cite{Kearns1998},
provided the first framework where computational lower bounds could be
proved unconditionally, but only within a restricted model of computation.
\end{historical}

\subsection{Where the Gap Is}

The fundamental gap is proof-theoretic: we lack techniques to prove
unconditional computational lower bounds for PAC learning.

\begin{enumerate}[leftmargin=*]
\item \textbf{No unconditional separations.}  Every known example of a
  computational--statistical gap relies on a complexity-theoretic assumption.
  Proving an unconditional gap would require proving $P \neq NP$ or something
  similarly difficult.

\item \textbf{No combinatorial characterization.}  For information-theoretic
  learnability, we have a complete combinatorial characterization: finite VC
  dimension.  For computationally efficient learnability, no analogous
  characterization exists.  We do not know what combinatorial property of a
  concept class determines whether it is efficiently learnable.

\item \textbf{Model dependence.}  Lower bounds proved in the SQ model, the
  low-degree polynomial model, or the sum-of-squares hierarchy apply only to
  those specific computational models.  An algorithm outside the model
  (like Gaussian elimination for parities) may circumvent the lower bound.
  No model of computation is known to be simultaneously natural, powerful
  enough to capture practical algorithms, and amenable to unconditional lower
  bounds.
\end{enumerate}

\subsection{What Would Close the Gap}

\begin{itemize}[leftmargin=*]
\item \textbf{A natural proof barrier for learning.}  Razborov and
  Rudich's natural proofs barrier explains why circuit lower bounds are hard.
  An analogous barrier for learning (explaining why computational lower
  bounds for PAC learning are hard) would clarify the structural limits of
  current proof techniques.
\item \textbf{A unifying computational model.}  A model of bounded
  computation that captures SQ algorithms, low-degree methods,
  sum-of-squares, and gradient descent, while admitting provable lower
  bounds, would provide a natural home for computational--statistical
  tradeoff results.
\item \textbf{Conditional characterizations.}  Even under assumptions
  like $P \neq NP$ or the hardness of LWE, a characterization of
  efficiently learnable classes (analogous to the fundamental theorem for
  information-theoretic learnability) would be a major advance.
\end{itemize}


% ================================================================
% SECTION 5: KOLMOGOROV COMPLEXITY AND THE IDEAL LEARNER
% ================================================================

\section{Kolmogorov Complexity and the Ideal Learner}
\label{sec:kolmogorov-learning}

\begin{openprob}[Computable Approximations to Solomonoff Induction]
\label{op:solomonoff-approx}
Solomonoff's universal prior $M(x) = \sum_{p:\, U(p)=x} 2^{-|p|}$ defines
an ideal but uncomputable learner.  To what extent can computable
approximations to $M$ achieve the learning-theoretic properties of the ideal
learner?
\end{openprob}

This section concerns the deepest stratum of the theory: the connection
between algorithmic information theory and learning.

\subsection{What is Known}

\begin{itemize}[leftmargin=*]
\item \textbf{Solomonoff's completeness theorem.}  For any computable
  measure $\mu$ on infinite binary sequences, the universal semimeasure $M$
  satisfies
  \[
    \sum_{n=1}^\infty \E_{x_{1:n} \sim \mu}
    \bigl[D_{\mathrm{KL}}(\mu(\cdot \mid x_{1:n}) \,\|\,
    M(\cdot \mid x_{1:n}))\bigr] < \infty.
  \]
  That is, the cumulative prediction error of $M$ relative to any
  computable environment is bounded.  This makes $M$ the ``ideal learner'':
  it converges to the true distribution at a rate dominated by the
  Kolmogorov complexity of the environment.

\item \textbf{MDL and MML as approximations.}  The minimum description
  length principle (Rissanen~\cite{Rissanen1984}) and minimum message length
  (Wallace and Boulton~\cite{WallaceBoulton1968}) are computable
  approximations to the Kolmogorov-optimal encoding.  Both embody a version
  of Occam's razor: prefer the hypothesis with the shortest total description
  (model + data given model).

\item \textbf{Occam bounds via Kolmogorov complexity.}  Li, Tromp, and
  Vit\'anyi~\cite{LiVitanyi2008} showed that replacing description length
  with Kolmogorov complexity in Occam-style PAC bounds yields tighter
  (but uncomputable) generalization guarantees.  The bound states that
  the generalization error of a hypothesis $h$ is controlled by
  $K(h)/n$, where $K(h)$ is the Kolmogorov complexity of $h$ and $n$
  is the sample size.

\item \textbf{Algorithmic probability and universal learning.}  Universal
  learning allows distribution-dependent rates; Solomonoff's prior achieves
  distribution-dependent convergence.  But the universal learning framework is
  information-theoretic (it asks only about sample complexity), while
  Solomonoff's prior is algorithmic (it asks about description complexity).
  The precise relationship between these two notions of universality is open.
\end{itemize}

\begin{historical}
Solomonoff's 1964 paper~\cite{Solomonoff1964} predates Gold's 1967
formalization of identification in the limit and Valiant's 1984
introduction of PAC learning.  It is among the first formal proposals
for a universal learning algorithm, though it was formulated in the language
of algorithmic information theory rather than learning theory.  The
intellectual genealogy runs:
Solomonoff~(1964) $\to$ Kolmogorov complexity~(1965) $\to$ MDL~(1978)
$\to$ Occam bounds in PAC~(1987) $\to$ compression schemes~(1986).
The compression conjecture of \Cref{sec:compression-conjecture} can be
viewed as asking whether the Occam principle (short description implies
good generalization) holds at the tightest possible quantitative level.
\end{historical}

\subsection{Where the Gap Is}

\begin{enumerate}[leftmargin=*]
\item \textbf{Uncomputability.}  Kolmogorov complexity and Solomonoff's
  prior are uncomputable.  Every computable approximation introduces a
  gap: MDL uses a fixed model class rather than all programs; bounded
  Solomonoff approximations (running programs for at most $t$ steps)
  introduce a time--complexity tradeoff that is poorly understood.

\item \textbf{No convergence rate for approximations.}  For the ideal
  learner $M$, the convergence rate is known: the expected KL divergence
  is $O(K(\mu)/n)$ where $K(\mu)$ is the complexity of the true
  environment.  For computable approximations, no analogous rate is
  established that holds uniformly over all computable environments.

\item \textbf{The gap between prediction and classification.}  Solomonoff's
  framework concerns \emph{sequence prediction}: predicting the next bit.
  PAC learning concerns \emph{concept identification}: finding a hypothesis
  close to the target.  The formal relationship between these two tasks
  is well understood in specific settings (e.g., the realizable case with
  i.i.d.\ data), but the general connection (especially in the agnostic
  or adversarial setting) remains underdeveloped.
\end{enumerate}

\subsection{What Would Close the Gap}

\begin{itemize}[leftmargin=*]
\item \textbf{Time-bounded Solomonoff induction.}  A theory of
  $M^t(x) = \sum_{p:\, U(p)=x,\; \mathrm{time}(p) \leq t} 2^{-|p|}$
  with provable convergence rates as a function of both $n$ (sample size)
  and $t$ (computation time) would bridge the gap between the ideal and
  the computable.  Such a theory would need to handle the fact that
  $M^t$ is not a semimeasure for finite $t$, and that the dominant programs
  may have very long running times.

\item \textbf{Kolmogorov complexity as a complexity measure for PAC
  learning.}  Formalizing $K(\mathcal{H})$ (the description complexity of
  a hypothesis class) as a complexity measure in the sense of
  Chapter~\ref{ch:dimensions} (with upper and lower bound relationships
  to VC dimension, Rademacher complexity, and compression size) would
  integrate algorithmic information theory into the concept graph of this
  book.

\item \textbf{A computable universal learner.}  An algorithm that, for
  every computable concept class $\mathcal{C}$ and every realizable
  distribution $D$, achieves PAC learning with sample complexity depending
  on $K(\mathcal{C})$, would be a computable analogue of the Solomonoff
  ideal.  Whether such an algorithm exists is itself an open question,
  closely related to the universal learning framework of
  Chapter~\ref{ch:universal}.
\end{itemize}


% ================================================================
% SECTION 6: FURTHER FRONTIERS
% ================================================================

\section{Further Frontiers}
\label{sec:further-frontiers}

The five problems above are those most tightly connected to the conceptual
framework of this book.  We note several other directions where the theory is
actively expanding, each involving mathematical structures that do not reduce
to the types established in Chapters~1--17.

\begin{openprob}[Private agnostic learning]
\label{op:private}
Differential privacy constrains the learner: the output must not reveal too
much about any individual example.  Kasiviswanathan et
al.~\cite{KasiviswanathanLeeNissimRaskhodnikovaSmith2011} showed that private
PAC learning under \emph{pure} differential privacy has sample complexity
$\Theta\bigl(d/\varepsilon + \log|\mathcal{X}|/(\varepsilon\alpha)\bigr)$, and
for \emph{approximate} privacy the Littlestone dimension characterizes private
learnability~\cite{AlonLivniMalliarisMoran2019,BunLivniMoran2020}, linking
online learning to privacy.  What is the tight sample complexity of private
\emph{agnostic} learning?
\end{openprob}

\begin{openprob}[Active learning under agnostic noise]
\label{op:active}
An active learner queries labels selectively, and the label complexity can be
exponentially smaller than passive sample complexity \emph{for favorable
classes and distributions} (for instance, halfspaces under a margin).
Hanneke~\cite{Hanneke2014} showed that the disagreement coefficient controls
label complexity for general classes.  What is the tight characterization of
active learning under \emph{agnostic} noise, and what is its computational
complexity?
\end{openprob}

\begin{openprob}[Transfer, meta-, and continual learning]
\label{op:meta}
Formal frameworks for learning across related tasks exist (Baxter's model of
learning to learn~\cite{Baxter2000}, and the PAC-Bayes approach to
meta-learning), but no characterization theorem analogous to the VC
characterization is known.  Can the combinatorial objects governing multi-task
learning be expressed through the single-task dimensions of this book, or are
genuinely new measures required?
\end{openprob}

\begin{openprob}[Online multiclass learnability]
\label{op:online-multiclass-frontiers}
The Littlestone dimension characterizes online binary learnability, but the
Littlestone dimension of the binary reductions does not always capture the
optimal multiclass mistake bound.  What combinatorial parameter characterizes
online multiclass learnability, and how does it relate to the DS dimension
(\Cref{ch:extensions})?
\end{openprob}

\begin{openprob}[Quantum sample complexity]
\label{op:quantum}
In quantum PAC learning the learner receives quantum superpositions of labeled
examples.  For binary classification the quantum and classical sample
complexities are equal up to constant factors, but for other tasks and for
\emph{time} complexity the relationship is unsettled.  Characterize, in terms
of classical complexity measures, exactly when quantum examples reduce the cost
of learning.
\end{openprob}

\begin{openprob}[A non-vacuous \emph{verified} bound for transformers]
\label{op:certified-transformer}
\Cref{sec:deep-generalization} notes that a machine-checked generalization
certificate already exists for an attention block
(\leanname{config_capacity_law}), but its Dudley budget grows with depth and
norm.  Can a generalization bound for a realistic transformer be \emph{both}
machine-checked and non-vacuous on a benchmark, thereby closing \Cref{op:deep-gen}?
\end{openprob}

\bigskip

\begin{computational}
\textbf{The shape of the frontier.}
The open problems of this chapter are not isolated.  They form a connected
subgraph in the concept graph, linked by shared dependencies.  The
compression conjecture involves the same VC dimension that governs PAC
learning, which connects to the computational hardness that creates
computational--statistical gaps, which connects to the SQ dimension, which
connects to the Kolmogorov complexity that underlies the ideal learner.
The deep learning generalization problem is linked to compression (through
lottery tickets), to margin theory (through spectral bounds), and to
algorithmic stability (through SGD analysis).  Universal learning connects
to the Littlestone dimension and therefore to private learning.

This interconnection is not coincidental.  A resolution of the compression
conjecture would yield insight into generalization bounds; a combinatorial
characterization of efficient learnability would require new structural
theorems about VC classes; a computable approximation to Solomonoff induction
would connect to both compression and universal learning.  Resolving any one
of these problems would reshape the others.
\end{computational}

\bigskip

This chapter has surveyed what we do not know.  The theory developed in
Chapters~1--17 provides a rigorous, interconnected framework for what we do
know: characterization theorems, tight bounds, separation results, and
structural analogies.  The open problems of this chapter mark the places
where that framework reaches its current limits, and where the next results
will be found.


% ============================================================
\appendix
% ============================================================

% ============================================================
% APPENDIX A: Complete Edge Inventory
% ============================================================
\chapter{Complete Edge Inventory}
\label{app:edge-inventory}

This appendix lists all 260 edges in the companion knowledge graph, organized by the 13 relation types.  Each table shows representative edges; the full inventory is machine-readable in \texttt{flt\_concept\_graph.json}.

\medskip
\noindent\textit{Column conventions.}
\nde{Source} and \nde{Target} give graph node identifiers.
\emph{Citation} gives the abbreviated bibliography key.
Relation-specific columns (\emph{Witness}, \emph{Obstruction type}, \emph{Generalization type}) appear where the schema requires them.

% ----------------------------------------------------------
\section*{A.1\quad \rel{defined\_using} (99 edges)}
\label{sec:edges-defined-using}

\begin{longtable}{lll}
\caption{\rel{defined\_using} edges (representative sample; 99 total).}\label{tab:defined-using}\\
\toprule
\textbf{Source} & \textbf{Target} & \textbf{Note} \\
\midrule
\endfirsthead
\toprule
\textbf{Source} & \textbf{Target} & \textbf{Note} \\
\midrule
\endhead
\midrule
\multicolumn{3}{r}{\footnotesize\itshape continued on next page} \\
\endfoot
\bottomrule
\endlastfoot
\nde{concept} & \nde{domain} & \\
\nde{concept} & \nde{label} & \\
\nde{concept\_class} & \nde{concept} & \\
\nde{hypothesis\_space} & \nde{concept} & \\
\nde{proper\_flag} & \nde{hypothesis\_space} & \\
\multicolumn{3}{c}{\footnotesize\itshape $\cdots$ 94 additional edges in \texttt{flt\_concept\_graph.json}} \\
\end{longtable}

% ----------------------------------------------------------
\section*{A.2\quad \rel{instance\_of} (25 edges)}
\label{sec:edges-instance-of}

\begin{longtable}{lll}
\caption{\rel{instance\_of} edges (representative sample; 25 total).}\label{tab:instance-of}\\
\toprule
\textbf{Source} & \textbf{Target} & \textbf{Note} \\
\midrule
\endfirsthead
\toprule
\textbf{Source} & \textbf{Target} & \textbf{Note} \\
\midrule
\endhead
\midrule
\multicolumn{3}{r}{\footnotesize\itshape continued on next page} \\
\endfoot
\bottomrule
\endlastfoot
\nde{text\_presentation} & \nde{data\_stream} & \\
\nde{informant\_presentation} & \nde{data\_stream} & \\
\nde{noisy\_input} & \nde{data\_stream} & \\
\nde{time\_index} & \nde{domain} & \\
\nde{iterative\_learner} & \nde{learner} & \\
\multicolumn{3}{c}{\footnotesize\itshape $\cdots$ 20 additional edges in \texttt{flt\_concept\_graph.json}} \\
\end{longtable}

% ----------------------------------------------------------
\section*{A.3\quad \rel{characterizes} (24 edges)}
\label{sec:edges-characterizes}

\begin{longtable}{lll}
\caption{\rel{characterizes} edges (representative sample; 24 total).}\label{tab:characterizes}\\
\toprule
\textbf{Source} & \textbf{Target} & \textbf{Citation} \\
\midrule
\endfirsthead
\toprule
\textbf{Source} & \textbf{Target} & \textbf{Citation} \\
\midrule
\endhead
\midrule
\multicolumn{3}{r}{\footnotesize\itshape continued on next page} \\
\endfoot
\bottomrule
\endlastfoot
\nde{vc\_characterization} & \nde{pac\_learning} & \cite{BlumerEhrenfeuchtHausslerWarmuth1989} \\
\nde{vc\_characterization} & \nde{vc\_dimension} & \cite{BlumerEhrenfeuchtHausslerWarmuth1989} \\
\nde{star\_number} & \nde{vc\_dimension} & \cite{Hanneke2024b} \\
\nde{trial\_and\_error} & \nde{ex\_learning} & \cite{Gold1965} \\
\nde{ds\_dimension} & \nde{pac\_learning} & \cite{BrukhimCarmonDinurMoranYehudayoff2022} \\
\multicolumn{3}{c}{\footnotesize\itshape $\cdots$ 19 additional edges in \texttt{flt\_concept\_graph.json}} \\
\end{longtable}

% ----------------------------------------------------------
\section*{A.4\quad \rel{analogy} (32 edges)}
\label{sec:edges-analogy}

\begin{longtable}{llll}
\caption{\rel{analogy} edges (representative sample; 32 total).}\label{tab:analogy}\\
\toprule
\textbf{Source} & \textbf{Target} & \textbf{Obstruction type} & \textbf{Citation} \\
\midrule
\endfirsthead
\toprule
\textbf{Source} & \textbf{Target} & \textbf{Obstruction type} & \textbf{Citation} \\
\midrule
\endhead
\midrule
\multicolumn{4}{r}{\footnotesize\itshape continued on next page} \\
\endfoot
\bottomrule
\endlastfoot
\nde{sample\_complexity} & \nde{vc\_dimension} & one\_way\_theorem\_only & \\
\nde{ordinal\_vc\_dim} & \nde{mind\_change\_ordinal} & proof\_method\_mismatch & \\
\nde{mistake\_bound} & \nde{littlestone\_dimension} & missing\_equiv\_witness & \\
\nde{label\_complexity} & \nde{sample\_complexity} & missing\_equiv\_witness & \\
\nde{concept\_drift} & \nde{online\_learning} & data\_model\_mismatch & \\
\multicolumn{4}{c}{\footnotesize\itshape $\cdots$ 27 additional edges in \texttt{flt\_concept\_graph.json}} \\
\end{longtable}

% ----------------------------------------------------------
\section*{A.5\quad \rel{measures} (22 edges)}
\label{sec:edges-measures}

\begin{longtable}{lll}
\caption{\rel{measures} edges (representative sample; 22 total).}\label{tab:measures}\\
\toprule
\textbf{Source} & \textbf{Target} & \textbf{Note} \\
\midrule
\endfirsthead
\toprule
\textbf{Source} & \textbf{Target} & \textbf{Note} \\
\midrule
\endhead
\midrule
\multicolumn{3}{r}{\footnotesize\itshape continued on next page} \\
\endfoot
\bottomrule
\endlastfoot
\nde{vc\_dimension} & \nde{concept\_class} & \\
\nde{shatters} & \nde{concept\_class} & \\
\nde{littlestone\_dimension} & \nde{concept\_class} & \\
\nde{star\_number} & \nde{version\_space} & \\
\nde{eluder\_dimension} & \nde{version\_space} & \\
\multicolumn{3}{c}{\footnotesize\itshape $\cdots$ 17 additional edges in \texttt{flt\_concept\_graph.json}} \\
\end{longtable}

% ----------------------------------------------------------
\section*{A.6\quad \rel{used\_in\_proof} (14 edges)}
\label{sec:edges-used-in-proof}

\begin{longtable}{lll}
\caption{\rel{used\_in\_proof} edges (representative sample; 14 total).}\label{tab:used-in-proof}\\
\toprule
\textbf{Source} & \textbf{Target} & \textbf{Citation} \\
\midrule
\endfirsthead
\toprule
\textbf{Source} & \textbf{Target} & \textbf{Citation} \\
\midrule
\endhead
\midrule
\multicolumn{3}{r}{\footnotesize\itshape continued on next page} \\
\endfoot
\bottomrule
\endlastfoot
\nde{gold\_theorem} & \nde{text\_presentation} & \cite{Gold1967} \\
\nde{gold\_theorem} & \nde{concept\_class} & \cite{Gold1967} \\
\nde{nfl\_theorem} & \nde{learner} & \cite{WolpertMacready1997} \\
\nde{nfl\_theorem} & \nde{inductive\_bias} & \cite{WolpertMacready1997} \\
\nde{pac\_lower\_bound} & \nde{vc\_dimension} & \cite{BlumerEhrenfeuchtHausslerWarmuth1989} \\
\multicolumn{3}{c}{\footnotesize\itshape $\cdots$ 9 additional edges in \texttt{flt\_concept\_graph.json}} \\
\end{longtable}

% ----------------------------------------------------------
\section*{A.7\quad \rel{does\_not\_imply} (9 edges)}
\label{sec:edges-does-not-imply}

\begin{longtable}{p{3cm}p{3cm}p{5.5cm}l}
\caption{\rel{does\_not\_imply} edges (all 9).}\label{tab:does-not-imply}\\
\toprule
\textbf{Source} & \textbf{Target} & \textbf{Witness} & \textbf{Citation} \\
\midrule
\endfirsthead
\toprule
\textbf{Source} & \textbf{Target} & \textbf{Witness} & \textbf{Citation} \\
\midrule
\endhead
\midrule
\multicolumn{4}{r}{\footnotesize\itshape continued on next page} \\
\endfoot
\bottomrule
\endlastfoot
\nde{pac\_learning} & \nde{mistake\_bounded} & Thresholds on $\R$: $\VC = 1$, $\Ldim = \infty$ & \cite{Littlestone1988} \\
\nde{ex\_learning} & \nde{pac\_learning} & All finite subsets of $\N$: identifiable from text, $\VC = \infty$ & \cite{Gold1967} \\
\nde{vc\_dimension} & \nde{computational\_hardness} & Poly-size circuits under DCRA: finite $\VC$ but not efficiently PAC-learnable & \cite{KearnsValiant1994} \\
\nde{vc\_dimension} & \nde{sq\_dimension} & Parities: $\VC = n$, $\SQdim = 2^n$ & \cite{BlumFurstJacksonKearnsMansourRudich1994} \\
\nde{proper\_improper\_separation} & \nde{pac\_learning} & 3-term DNF: NP-hard proper, poly-time improper & \cite{PittValiant1988} \\
\multicolumn{4}{c}{\footnotesize\itshape 4 additional edges in \texttt{flt\_concept\_graph.json}} \\
\end{longtable}

% ----------------------------------------------------------
\section*{A.8\quad \rel{upper\_bounds} (8 edges)}
\label{sec:edges-upper-bounds}

\begin{longtable}{lll}
\caption{\rel{upper\_bounds} edges (all 8).}\label{tab:upper-bounds}\\
\toprule
\textbf{Source} & \textbf{Target} & \textbf{Citation} \\
\midrule
\endfirsthead
\toprule
\textbf{Source} & \textbf{Target} & \textbf{Citation} \\
\midrule
\endhead
\bottomrule
\endlastfoot
\nde{littlestone\_dimension} & \nde{vc\_dimension} & \cite{Littlestone1988} \\
\nde{rademacher\_complexity} & \nde{generalization\_error} & \cite{BartlettMendelson2002} \\
\nde{pac\_bayes\_bound} & \nde{generalization\_error} & \cite{McAllester1999} \\
\nde{information\_theoretic\_bound} & \nde{generalization\_error} & \cite{XuRaginsky2017} \\
\nde{vc\_dimension} & \nde{rademacher\_complexity} & \cite{Sauer1972} \\
\nde{margin\_theory} & \nde{generalization\_error} & \cite{BartlettFosterTelgarsky2017} \\
\multicolumn{3}{c}{\footnotesize\itshape 2 additional edges in \texttt{flt\_concept\_graph.json}} \\
\end{longtable}

% ----------------------------------------------------------
\section*{A.9\quad \rel{restricts} (8 edges)}
\label{sec:edges-restricts}

\begin{longtable}{lll}
\caption{\rel{restricts} edges (representative sample; 8 total).}\label{tab:restricts}\\
\toprule
\textbf{Source} & \textbf{Target} & \textbf{Citation} \\
\midrule
\endfirsthead
\toprule
\textbf{Source} & \textbf{Target} & \textbf{Citation} \\
\midrule
\endhead
\bottomrule
\endlastfoot
\nde{bc\_learning} & \nde{ex\_learning} & \cite{CaseSmith1983} \\
\nde{agnostic\_pac\_learning} & \nde{pac\_learning} & \cite{Haussler1992} \\
\nde{universal\_learning} & \nde{pac\_learning} & \cite{BousquetHannekeMoranVanHandelYehudayoff2021} \\
\nde{meta\_pac\_bound} & \nde{pac\_learning} & \cite{Baxter2000} \\
\nde{anomalous\_learning} & \nde{ex\_learning} & \\
\multicolumn{3}{c}{\footnotesize\itshape 3 additional edges in \texttt{flt\_concept\_graph.json}} \\
\end{longtable}

% ----------------------------------------------------------
\section*{A.10\quad \rel{extends\_grammar} (8 edges)}
\label{sec:edges-extends-grammar}

\begin{longtable}{llll}
\caption{\rel{extends\_grammar} edges (representative sample; 8 total).}\label{tab:extends-grammar}\\
\toprule
\textbf{Source} & \textbf{Target} & \textbf{Generalization type} & \textbf{Citation} \\
\midrule
\endfirsthead
\toprule
\textbf{Source} & \textbf{Target} & \textbf{Generalization type} & \textbf{Citation} \\
\midrule
\endhead
\bottomrule
\endlastfoot
\nde{ordinal\_vc\_dim} & \nde{vc\_dimension} & new\_primitives\_required & \\
\nde{pseudodimension} & \nde{vc\_dimension} & new\_primitives\_required & \\
\nde{natarajan\_dimension} & \nde{vc\_dimension} & new\_primitives\_required & \\
\nde{mind\_change\_ordinal} & \nde{mind\_change\_count} & new\_primitives\_required & \\
\nde{pushdown\_automaton} & \nde{dfa} & new\_primitives\_required & \\
\multicolumn{4}{c}{\footnotesize\itshape 3 additional edges in \texttt{flt\_concept\_graph.json}} \\
\end{longtable}

% ----------------------------------------------------------
\section*{A.11\quad \rel{lower\_bounds} (5 edges)}
\label{sec:edges-lower-bounds}

\begin{longtable}{lll}
\caption{\rel{lower\_bounds} edges (all 5).}\label{tab:lower-bounds}\\
\toprule
\textbf{Source} & \textbf{Target} & \textbf{Citation} \\
\midrule
\endfirsthead
\toprule
\textbf{Source} & \textbf{Target} & \textbf{Citation} \\
\midrule
\endhead
\bottomrule
\endlastfoot
\nde{gold\_theorem} & \nde{ex\_learning} & \cite{Gold1967} \\
\nde{pac\_lower\_bound} & \nde{pac\_learning} & \cite{BlumerEhrenfeuchtHausslerWarmuth1989} \\
\nde{pac\_lower\_bound} & \nde{sample\_complexity} & \cite{BlumerEhrenfeuchtHausslerWarmuth1989} \\
\nde{proper\_improper\_separation} & \nde{pac\_learning} & \cite{PittValiant1988} \\
\nde{computational\_hardness} & \nde{pac\_learning} & \cite{KearnsValiant1994} \\
\end{longtable}

% ----------------------------------------------------------
\section*{A.12\quad \rel{strictly\_stronger} (4 edges)}
\label{sec:edges-strictly-stronger}

\begin{longtable}{p{2.8cm}p{2.8cm}p{5.5cm}l}
\caption{\rel{strictly\_stronger} edges (all 4).}\label{tab:strictly-stronger}\\
\toprule
\textbf{Source} & \textbf{Target} & \textbf{Witness} & \textbf{Citation} \\
\midrule
\endfirsthead
\toprule
\textbf{Source} & \textbf{Target} & \textbf{Witness} & \textbf{Citation} \\
\midrule
\endhead
\bottomrule
\endlastfoot
\nde{universal\_learning} & \nde{pac\_learning} & Classes with infinite $\VC$ but no infinite Littlestone tree & \cite{BousquetHannekeMoranVanHandelYehudayoff2021} \\
\nde{ex\_learning} & \nde{finite\_learning} & Classes requiring unbounded mind changes (e.g., pattern languages) & \cite{Gold1967} \\
\nde{online\_learning} & \nde{pac\_learning} & Finite $\Ldim \Rightarrow$ finite $\VC$; thresholds on finite domains: $\Ldim < \VC$ & \cite{Littlestone1988} \\
\nde{ds\_dimension} & \nde{natarajan\_dimension} & Hyperbolic pseudo-manifold construction & \cite{BrukhimCarmonDinurMoranYehudayoff2022} \\
\end{longtable}

% ----------------------------------------------------------
\section*{A.13\quad \rel{requires\_assumption} (2 edges)}
\label{sec:edges-requires-assumption}

\begin{longtable}{lll}
\caption{\rel{requires\_assumption} edges (all 2).}\label{tab:requires-assumption}\\
\toprule
\textbf{Source} & \textbf{Target} & \textbf{Citation} \\
\midrule
\bottomrule
\nde{computational\_hardness} & \nde{inductive\_bias} & \cite{ApplebaumBarakXiao2008} \\
\nde{posterior\_consistency} & \nde{bayesian\_inference} & \cite{Schwartz1965} \\
\end{longtable}

\bigskip
\noindent\textbf{Full inventory.}\quad
The complete 260-edge listing is available in \texttt{flt\_concept\_graph.json} (v1.0).
Programmatic access is recommended for any systematic traversal; see Appendix~\ref{app:validation} for the validation suite.

% ============================================================
% APPENDIX B: Graph Traversal Demonstrations
% ============================================================
\chapter{Graph Traversal Demonstrations}
\label{app:traversals}

These 15 tasks demonstrate the graph's utility for cross-paradigm reasoning.
Each traversal states a query, lists the sequence of edges traversed, and assembles an answer from the path.
Edge notation: $\edge{source}{relation}{target}$.

% ----------------------------------------------------------
% T01
% ----------------------------------------------------------
\begin{traversal}
\textbf{T01: Prerequisite retrieval for VC dimension.}\par\medskip
\emph{Query:} What concepts must be defined before $\VC$ can be stated?\par
\emph{Path:}\par
\begin{enumerate}[leftmargin=2em, nosep]
  \item $\edge{vc\_dimension}{defined\_using}{shatters}$
  \item $\edge{shatters}{defined\_using}{concept\_class}$ (via \nde{shatters} prerequisites)
  \item $\edge{concept\_class}{defined\_using}{concept}$
  \item $\edge{concept}{defined\_using}{domain}$
  \item $\edge{concept}{defined\_using}{label}$
\end{enumerate}
\emph{Answer:} The prerequisite chain is: \nde{domain}, \nde{label} $\to$ \nde{concept} $\to$ \nde{concept\_class} $\to$ \nde{shatters} $\to$ \nde{vc\_dimension}. In the graph, the single direct \rel{defined\_using} edge from \nde{vc\_dimension} targets \nde{shatters}.
\end{traversal}

% ----------------------------------------------------------
% T02
% ----------------------------------------------------------
\begin{traversal}
\textbf{T02: What characterizes PAC learnability?}\par\medskip
\emph{Query:} What are all known characterizations of PAC learnability?\par
\emph{Path:} Collect all nodes $X$ where $\edge{X}{characterizes}{pac\_learning}$.\par
\begin{enumerate}[leftmargin=2em, nosep]
  \item $\edge{vc\_characterization}{characterizes}{pac\_learning}$
  \item $\edge{vc\_dimension}{characterizes}{pac\_learning}$ (via \nde{vc\_characterization})
  \item $\edge{ds\_dimension}{characterizes}{pac\_learning}$
  \item $\edge{fundamental\_theorem}{characterizes}{pac\_learning}$
  \item $\edge{occam\_algorithm}{characterizes}{pac\_learning}$
\end{enumerate}
\emph{Answer:} Five nodes characterize PAC learnability: \nde{vc\_dimension}, \nde{vc\_characterization}, \nde{ds\_dimension}, \nde{fundamental\_theorem}, and \nde{occam\_algorithm}.
\end{traversal}

% ----------------------------------------------------------
% T03
% ----------------------------------------------------------
\begin{traversal}
\textbf{T03: Is a PAC-learnable class online-learnable?}\par\medskip
\emph{Query:} Is a PAC-learnable class necessarily online-learnable with finite mistakes?\par
\emph{Path:}\par
\begin{enumerate}[leftmargin=2em, nosep]
  \item $\edge{pac\_learning}{does\_not\_imply}{mistake\_bounded}$
\end{enumerate}
\emph{Answer:} No. The edge carries witness: thresholds on $\R$ have $\VC = 1$ but $\Ldim = \infty$ \cite{Littlestone1988}.
\end{traversal}

% ----------------------------------------------------------
% T04
% ----------------------------------------------------------
\begin{traversal}
\textbf{T04: Upper bounds on generalization error.}\par\medskip
\emph{Query:} What generalization bounds are available and what do they each depend on?\par
\emph{Path:} Find all $X$ with $\edge{X}{upper\_bounds}{generalization\_error}$, then follow \rel{defined\_using} from each.\par
\begin{enumerate}[leftmargin=2em, nosep]
  \item $\edge{rademacher\_complexity}{upper\_bounds}{generalization\_error}$ \cite{BartlettMendelson2002}
  \item $\edge{pac\_bayes\_bound}{upper\_bounds}{generalization\_error}$ \cite{McAllester1999}
  \item $\edge{information\_theoretic\_bound}{upper\_bounds}{generalization\_error}$ \cite{XuRaginsky2017}
  \item $\edge{margin\_theory}{upper\_bounds}{generalization\_error}$ \cite{BartlettFosterTelgarsky2017}
\end{enumerate}
\emph{Answer:} Four bound families provide upper bounds on generalization error: Rademacher complexity, PAC-Bayes, information-theoretic, and margin-based bounds. Each has distinct \rel{defined\_using} dependencies (e.g., Rademacher complexity depends on the growth function and hypothesis space).
\end{traversal}

% ----------------------------------------------------------
% T05
% ----------------------------------------------------------
\begin{traversal}
\textbf{T05: Cross-paradigm Bayesian--multiclass reasoning.}\par\medskip
\emph{Query:} Can a Bayesian learner be PAC-Bayes analyzed for a multiclass problem where $\Ndim < \infty$ but the class is not learnable?\par
\emph{Path:}\par
\begin{enumerate}[leftmargin=2em, nosep]
  \item $\edge{bayesian\_learner}{instance\_of}{learner}$
  \item $\edge{pac\_bayes\_bound}{upper\_bounds}{generalization\_error}$
  \item $\edge{natarajan\_dimension}{does\_not\_imply}{pac\_learning}$
  \item $\edge{ds\_dimension}{characterizes}{pac\_learning}$
\end{enumerate}
\emph{Answer:} PAC-Bayes bounds apply to any posterior $Q$ (edge~2), so a Bayesian learner (edge~1) can be analyzed.  However, finite Natarajan dimension does not imply PAC learnability (edge~3); the DS dimension is the correct multiclass characterization (edge~4).  If $\DSdim = \infty$, the bounds may be vacuous despite $\Ndim < \infty$.
\end{traversal}

% ----------------------------------------------------------
% T06
% ----------------------------------------------------------
\begin{traversal}
\textbf{T06: $\BC \supsetneq \Ex \supsetneq \FIN$ hierarchy.}\par\medskip
\emph{Query:} What is the strict power hierarchy among success criteria in Gold's model?\par
\emph{Path:}\par
\begin{enumerate}[leftmargin=2em, nosep]
  \item $\edge{bc\_learning}{restricts}{ex\_learning}$ \cite{CaseSmith1983}
  \item $\edge{ex\_learning}{strictly\_stronger}{finite\_learning}$ \cite{Gold1967}
\end{enumerate}
\emph{Answer:} $\BC \supsetneq \Ex \supsetneq \FIN$.  BC generalizes Ex by removing the syntactic convergence requirement (edge~1). Ex is strictly stronger than FIN because some classes require unbounded mind changes (edge~2).
\end{traversal}

% ----------------------------------------------------------
% T07
% ----------------------------------------------------------
\begin{traversal}
\textbf{T07: Why concept drift $\neq$ online learning (obstruction).}\par\medskip
\emph{Query:} Why is the analogy between concept drift and online learning not a formal theorem?\par
\emph{Path:}\par
\begin{enumerate}[leftmargin=2em, nosep]
  \item $\edge{concept\_drift}{analogy}{online\_learning}$, obstruction type: \texttt{data\_model\_mismatch}
\end{enumerate}
\emph{Answer:} Data model mismatch: drift uses i.i.d.\ draws per step from a slowly changing distribution; online learning uses adversarially chosen instances from a fixed class.
\end{traversal}

% ----------------------------------------------------------
% T08
% ----------------------------------------------------------
\begin{traversal}
\textbf{T08: Which generalizations of VC dimension required new primitives?}\par\medskip
\emph{Query:} Which generalizations of $\VC$ required inventing new formal primitives?\par
\emph{Path:} Find all $X$ with $\edge{X}{extends\_grammar}{vc\_dimension}$.\par
\begin{enumerate}[leftmargin=2em, nosep]
  \item $\edge{pseudodimension}{extends\_grammar}{vc\_dimension}$
  \item $\edge{natarajan\_dimension}{extends\_grammar}{vc\_dimension}$
  \item $\edge{ordinal\_vc\_dim}{extends\_grammar}{vc\_dimension}$
\end{enumerate}
\emph{Answer:} Three nodes extend VC grammar: pseudodimension (witness thresholds for real-valued functions), Natarajan dimension (two-coloring for multiclass), and ordinal VC dimension (transfinite generalization).
\end{traversal}

% ----------------------------------------------------------
% T09
% ----------------------------------------------------------
\begin{traversal}
\textbf{T09: What prevents efficient PAC learning even when $\VC < \infty$?}\par\medskip
\emph{Query:} What prevents efficient PAC learning even when the VC dimension is finite?\par
\emph{Path:}\par
\begin{enumerate}[leftmargin=2em, nosep]
  \item $\edge{computational\_hardness}{lower\_bounds}{pac\_learning}$ \cite{KearnsValiant1994}
  \item $\edge{computational\_hardness}{used\_in\_proof}{vc\_dimension}$
  \item $\edge{computational\_hardness}{requires\_assumption}{inductive\_bias}$ \cite{ApplebaumBarakXiao2008}
  \item $\edge{vc\_dimension}{does\_not\_imply}{computational\_hardness}$ \cite{KearnsValiant1994}
\end{enumerate}
\emph{Answer:} Computational hardness results (edge~1) show that under cryptographic assumptions (edge~3), polynomial-size circuits have finite $\VC$ but are not efficiently PAC-learnable. Finite $\VC$ alone does not imply tractability (edge~4).
\end{traversal}

% ----------------------------------------------------------
% T10
% ----------------------------------------------------------
\begin{traversal}
\textbf{T10: Curriculum ordering for PAC learning.}\par\medskip
\emph{Query:} Generate a reading order for learning PAC theory from scratch.\par
\emph{Path:} Topological sort of the \rel{defined\_using} subgraph restricted to PAC-related nodes.\par
\emph{Answer:} One valid ordering:\par
\begin{enumerate}[leftmargin=2em, nosep]
  \item \nde{domain}, \nde{label}
  \item \nde{concept} (depends on domain, label)
  \item \nde{shatters}, \nde{iid\_sample}
  \item \nde{concept\_class}, \nde{hypothesis\_space}
  \item \nde{pac\_learning}, \nde{generalization\_error}
  \item \nde{sample\_complexity}
  \item \nde{vc\_dimension}
  \item \nde{vc\_characterization}
  \item \nde{fundamental\_theorem}
\end{enumerate}
Multiple valid topological orderings exist; the key constraint is that each node appears after all its \rel{defined\_using} targets.
\end{traversal}

% ----------------------------------------------------------
% T11
% ----------------------------------------------------------
\begin{traversal}
\textbf{T11: Is reinforcement learning covered in this graph?}\par\medskip
\emph{Query:} Is reinforcement learning covered in this graph?\par
\emph{Path:} Check for \nde{scope\_boundary\_rl} node.\par
\emph{Answer:} Explicitly excluded. The node \nde{scope\_boundary\_rl} states that RL extends bandits to sequential decision-making with a different problem structure (MDP, policy, value function). No edges connect to kernel nodes.
\end{traversal}

% ----------------------------------------------------------
% T12
% ----------------------------------------------------------
\begin{traversal}
\textbf{T12: Count separation results.}\par\medskip
\emph{Query:} List all known separations between learning paradigms with their witnesses.\par
\emph{Path:} Collect all \rel{does\_not\_imply} and \rel{strictly\_stronger} edges.\par
\emph{Answer:} 9 \rel{does\_not\_imply} edges $+$ 4 \rel{strictly\_stronger} edges $= 13$ total separation results.  Each carries a \texttt{witness} field and a \texttt{citation}. See Tables~\ref{tab:does-not-imply} and~\ref{tab:strictly-stronger} in Appendix~A.
\end{traversal}

% ----------------------------------------------------------
% T13
% ----------------------------------------------------------
\begin{traversal}
\textbf{T13: How do $\VC$, $\Ldim$, $\SQdim$, and star number relate?}\par\medskip
\emph{Query:} How do VCdim, Littlestone dim, SQ dim, and star number relate to each other?\par
\emph{Path:}\par
\begin{enumerate}[leftmargin=2em, nosep]
  \item $\edge{star\_number}{characterizes}{vc\_dimension}$ \cite{Hanneke2024b}
  \item $\edge{littlestone\_dimension}{upper\_bounds}{vc\_dimension}$ \cite{Littlestone1988}
  \item $\edge{sq\_dimension}{analogy}{vc\_dimension}$ (exponential gap possible: parities)
  \item $\edge{sq\_dimension}{does\_not\_imply}{vc\_dimension}$ (parities: $\VC = n$, $\SQdim = 2^n$)
\end{enumerate}
\emph{Answer:} Star number is equivalent to $\VC$ (edge~1). Finite $\Ldim$ implies finite $\VC$ but the gap can be infinite (edge~2). SQ dimension is structurally analogous but formally independent, with exponential separations (edges~3--4).
\end{traversal}

% ----------------------------------------------------------
% T14
% ----------------------------------------------------------
\begin{traversal}
\textbf{T14: Real-valued extension of VC dimension.}\par\medskip
\emph{Query:} What changes when moving from binary classification to real-valued prediction?\par
\emph{Path:}\par
\begin{enumerate}[leftmargin=2em, nosep]
  \item $\edge{pseudodimension}{extends\_grammar}{vc\_dimension}$
  \item $\edge{pseudodimension}{restricts}{fat\_shattering\_dimension}$
  \item $\edge{fat\_shattering\_dimension}{characterizes}{real\_valued\_pac}$ (via appropriate node)
\end{enumerate}
\emph{Answer:} Pseudodimension extends VC grammar by introducing witness thresholds for real-valued functions (edge~1). It removes the scale-sensitivity from fat-shattering dimension (edge~2, constraint removal). Fat-shattering adds a margin parameter $\gamma$ and characterizes real-valued PAC learnability.  Three nodes in total extend the VC grammar: \nde{pseudodimension}, \nde{natarajan\_dimension}, \nde{ordinal\_vc\_dim}.
\end{traversal}

% ----------------------------------------------------------
% T15
% ----------------------------------------------------------
\begin{traversal}
\textbf{T15: Does the fundamental theorem extend to multiclass?}\par\medskip
\emph{Query:} Does the fundamental theorem of statistical learning extend to multiclass?\par
\emph{Path:}\par
\begin{enumerate}[leftmargin=2em, nosep]
  \item Check: \nde{fundamental\_theorem} has no direct edges to multiclass-specific nodes.
  \item $\edge{natarajan\_dimension}{does\_not\_imply}{pac\_learning}$ --- finite $\Ndim$ insufficient when $|Y| = \infty$.
  \item $\edge{ds\_dimension}{characterizes}{pac\_learning}$ \cite{BrukhimCarmonDinurMoranYehudayoff2022}
  \item $\edge{natarajan\_dimension}{extends\_grammar}{vc\_dimension}$
  \item $\edge{ds\_dimension}{strictly\_stronger}{natarajan\_dimension}$ \cite{BrukhimCarmonDinurMoranYehudayoff2022}
\end{enumerate}
\emph{Answer:} The fundamental theorem ($\VC \leftrightarrow$ PAC) is specific to binary classification (no multiclass edges from edge~1).  For multiclass: Natarajan dimension gives bounds with a $\log k$ gap but does not characterize learnability when $|Y| = \infty$ (edge~2).  The DS dimension is the correct multiclass characterization (edge~3), and it is strictly stronger than Natarajan dimension (edge~5).
\end{traversal}

% ============================================================
% APPENDIX C: Graph Validation
% ============================================================
\chapter{Graph Validation}
\label{app:validation}

This appendix documents the validation infrastructure that enforces structural integrity of \texttt{flt\_concept\_graph.json}.  Two scripts, both located in \texttt{scripts/}, perform complementary checks: \texttt{validate\_graph.py} enforces the embedded JSON schema and all relational constraints, while \texttt{validate\_bibliography\_links.py} cross-references every citation field against the BibTeX file.  Together they ensure that the graph remains self-consistent, fully cited, and release-ready.

% ----------------------------------------------------------
\section{Validation Scripts}
\label{sec:validation-scripts}

The concept graph embeds its own JSON schema in the \texttt{meta.json\_schema} object.  Both validators read this schema at runtime, so adding a new relation type, required field, or status value automatically propagates to the checks without modifying validator code.

\begin{table}[ht]
\centering\small
\caption{Companion files relevant to validation.}\label{tab:validation-files}
\begin{tabular}{lp{7.5cm}}
\toprule
\textbf{File} & \textbf{Purpose} \\
\midrule
\texttt{flt\_concept\_graph.json} & The knowledge graph (142 nodes, 260 edges, 13 relation types). \\
\texttt{flt\_bibliography.bib} & Complete BibTeX bibliography (${\sim}120$ entries). \\
\texttt{scripts/validate\_graph.py} & Schema and constraint validator (13 checks). \\
\texttt{scripts/validate\_bibliography\_links.py} & Bibliography cross-reference validator (3 checks). \\
\texttt{scripts/run\_reference\_tasks.py} & Executes the 15 benchmark tasks against the graph. \\
\texttt{scripts/evaluate\_reference\_tasks.py} & Scores task outputs against gold answers. \\
\bottomrule
\end{tabular}
\end{table}

\medskip
\noindent\textbf{Invocation.}  Both validators accept \texttt{-{}-graph} and \texttt{-{}-bib} flags with sensible defaults (the repository root).

\begin{verbatim}
  python3 scripts/validate_graph.py \
      --graph flt_concept_graph.json \
      --bib   flt_bibliography.bib

  python3 scripts/validate_bibliography_links.py \
      --graph flt_concept_graph.json \
      --bib   flt_bibliography.bib
\end{verbatim}

\noindent Both scripts exit with code~0 on success and code~1 on failure, making them suitable for continuous-integration pipelines.  A trailing-comma tolerance layer (regex-based) strips syntactic noise before JSON parsing, so hand-edited graph files with accidental trailing commas do not cause spurious parse failures.

% ----------------------------------------------------------
\section{The Thirteen Validation Checks}
\label{sec:thirteen-checks}

\texttt{validate\_graph.py} executes the following checks in order.  A single failure in any check causes the overall result to be \texttt{FAIL}.

\begin{longtable}{r p{3.8cm} p{7.2cm}}
\caption{Validation checks performed by \texttt{validate\_graph.py}.}\label{tab:validation-checks}\\
\toprule
\textbf{\#} & \textbf{Check name} & \textbf{Description} \\
\midrule
\endfirsthead
\toprule
\textbf{\#} & \textbf{Check name} & \textbf{Description} \\
\midrule
\endhead
\midrule
\multicolumn{3}{r}{\footnotesize\itshape continued on next page} \\
\endfoot
\bottomrule
\endlastfoot
1 & JSON parses correctly & The graph file is valid JSON after trailing-comma stripping. Executed before the \texttt{Validator} object is constructed; a parse failure terminates immediately. \\
2 & Unique node IDs & Every \texttt{id} field across the 142-element node array is unique.  Collected IDs are cached for use by subsequent checks. \\
3 & Edge endpoints exist & Every edge \texttt{source} and \texttt{target} value appears in the set of declared node IDs from Check~2. \\
4 & Relation values in enum & Every edge \texttt{relation} value is one of the 13 strings declared in \texttt{meta.json\_schema.relation\_enum}. \\
5 & Required node fields & Every node contains all 8 required fields: \texttt{id}, \texttt{name}, \texttt{category}, \texttt{layer}, \texttt{status}, \texttt{claim\_type}, \texttt{description}, \texttt{formal\_definition}. \\
6 & Required edge fields & Every edge contains all 3 universal required fields: \texttt{source}, \texttt{target}, \texttt{relation}. \\
7 & Edge constraints & Relation-specific required fields are present and non-empty.  See \S\ref{sec:edge-constraints} for the full constraint table. \\
8 & No duplicate edges & No two edges share the same (\texttt{source}, \texttt{relation}, \texttt{target}) triple. \\
9 & Bib keys resolve & All values in node \texttt{bib\_keys} arrays, \texttt{provenance.introduced\_by}, \texttt{provenance.proved\_by}, and edge \texttt{citation} fields that match the BibTeX key pattern \texttt{[A-Za-z]+[0-9]\{4\}[a-z]?} resolve against entries in the bibliography file. \\
10 & Node count matches meta & The actual length of the \texttt{nodes} array equals \texttt{meta.node\_count} (currently 142). \\
11 & Edge count matches meta & The actual length of the \texttt{edges} array equals \texttt{meta.edge\_count} (currently 260). \\
12 & Status values in enum & Every node \texttt{status} value is one of: \texttt{defined}, \texttt{proved}, \texttt{deferred}, \texttt{scope\_note}. \\
13 & No \texttt{generalizes\_unclassified} & No edges carry the legacy relation \texttt{generalizes\_unclassified}.  All such edges must be reclassified as either \rel{restricts} (constraint removal) or \rel{extends\_grammar} (new primitives) before release. \\
\end{longtable}

% ----------------------------------------------------------
\section{Edge Constraints}
\label{sec:edge-constraints}

Not all relation types carry the same metadata.  The schema declares relation-specific \emph{required fields} beyond the universal triple (\texttt{source}, \texttt{target}, \texttt{relation}).  Check~7 enforces these constraints.

\begin{longtable}{p{3.4cm} p{3.2cm} p{5.8cm}}
\caption{Relation-specific required fields enforced by Check~7.}\label{tab:edge-constraints}\\
\toprule
\textbf{Relation type} & \textbf{Required fields} & \textbf{Rationale} \\
\midrule
\endfirsthead
\toprule
\textbf{Relation type} & \textbf{Required fields} & \textbf{Rationale} \\
\midrule
\endhead
\bottomrule
\endlastfoot
\rel{strictly\_stronger}     & \texttt{witness}, \texttt{citation} & Strict hierarchy requires a separating construction and a proof reference. \\
\rel{does\_not\_imply}        & \texttt{witness}, \texttt{citation} & Negative result requires an explicit counterexample and a proof reference. \\
\rel{characterizes}           & \texttt{citation} & Equivalence theorem requires a traceable reference. \\
\rel{upper\_bounds}           & \texttt{citation} & Bound statement requires a proof reference. \\
\rel{lower\_bounds}           & \texttt{citation} & Bound statement requires a proof reference. \\
\rel{requires\_assumption}    & \texttt{citation} & Conditional result requires a reference. \\
\rel{used\_in\_proof}         & \texttt{citation} & Proof dependency requires a reference to the theorem that uses the concept. \\
\midrule
\rel{defined\_using}          & \emph{none beyond triple} & Definitional edges are self-documenting from node definitions. \\
\rel{instance\_of}            & \emph{none beyond triple} & Taxonomic specializations are self-documenting. \\
\rel{measures}                & \emph{none beyond triple} & Measurement relationships are self-documenting. \\
\rel{analogy}                 & \emph{none beyond triple} & Carries optional \texttt{obstruction\_type} and \texttt{obstruction} fields. \\
\rel{restricts}               & \emph{none beyond triple} & May carry optional \texttt{note}. \\
\rel{extends\_grammar}        & \emph{none beyond triple} & May carry optional \texttt{generalization\_type}. \\
\end{longtable}

\noindent\textbf{Design principle.}  The seven relations in the upper block encode mathematical claims---equivalences, bounds, separations, conditional validity, proof dependencies---and therefore require traceable citations.  The two separation types (\rel{strictly\_stronger} and \rel{does\_not\_imply}) additionally require a \texttt{witness} field: a concrete concept class, construction, or example that demonstrates the separation.  The six relations in the lower block encode structural or definitional relationships that do not assert a mathematical result and therefore need no external citation.

% ----------------------------------------------------------
\section{The Thirteen Relation Types}
\label{sec:relation-semantics}

Each relation type has a forward reading, an inverse reading, and specific semantics.  Table~\ref{tab:relation-semantics} documents all 13, ordered by frequency in the graph.

\begin{longtable}{p{3.2cm} r p{6.8cm}}
\caption{Relation types with semantics and edge counts.}\label{tab:relation-semantics}\\
\toprule
\textbf{Relation} & \textbf{Edges} & \textbf{Semantics (forward $\to$ inverse)} \\
\midrule
\endfirsthead
\toprule
\textbf{Relation} & \textbf{Edges} & \textbf{Semantics (forward $\to$ inverse)} \\
\midrule
\endhead
\midrule
\multicolumn{3}{r}{\footnotesize\itshape continued on next page} \\
\endfoot
\bottomrule
\endlastfoot
\rel{defined\_using} & 99 & $X$'s formal definition directly invokes $Y$.  Forms the definitional dependency DAG. \newline Inverse: \rel{definition\_of}. \\[3pt]
\rel{analogy} & 32 & Structural parallel that is formally independent.  Symmetric: $A \mathbin{\rel{analogy}} B$ iff $B \mathbin{\rel{analogy}} A$.  Each edge carries an optional \texttt{obstruction\_type} field (one of: \texttt{data\_model\_mismatch}, \texttt{missing\_equiv\_witness}, \texttt{one\_way\_theorem\_only}, \texttt{proof\_method\_mismatch}) classifying why the analogy does not upgrade to a theorem. \\[3pt]
\rel{instance\_of} & 25 & $X$ is a specific instance or specialization of $Y$. \newline Inverse: \rel{has\_instance}. \\[3pt]
\rel{characterizes} & 24 & $X \leftrightarrow Y$: equivalence under the stated conditions (the fundamental ``iff'' edges). \newline Inverse: \rel{characterized\_by}.  Requires \texttt{citation}. \\[3pt]
\rel{measures} & 22 & Complexity measure $X$ quantifies a property of object $Y$. \newline Inverse: \rel{measured\_by}. \\[3pt]
\rel{used\_in\_proof} & 14 & Theorem $X$ invokes concept $Y$ in its statement or proof.  Distinct from \rel{defined\_using}: the latter is a definitional prerequisite, while \rel{used\_in\_proof} records proof-level dependencies. \newline Inverse: \rel{proves\_about}.  Requires \texttt{citation}. \\[3pt]
\rel{does\_not\_imply} & 9 & $X$ does \emph{not} imply $Y$, with a concrete separating witness.  These are the negative-layer edges that encode the field's ``what does not hold'' structure. \newline Inverse: \rel{not\_implied\_by}.  Requires \texttt{witness} and \texttt{citation}. \\[3pt]
\rel{restricts} & 8 & $X$ generalizes $Y$ by removing a constraint: $Y = X +$ additional restriction.  Conservative generalization---no new formal primitives are needed. \newline Inverse: \rel{restricted\_from}. \\[3pt]
\rel{extends\_grammar} & 8 & $X$ generalizes $Y$ but required inventing new concepts or primitives not present in $Y$.  Generative grammar expansion.  Each edge carries an optional \texttt{generalization\_type} field (typically \texttt{new\_primitives\_required}). \newline Inverse: \rel{grammar\_extended\_by}. \\[3pt]
\rel{upper\_bounds} & 8 & $X$ provides an upper bound on $Y$. \newline Inverse: \rel{bounded\_above\_by}.  Requires \texttt{citation}. \\[3pt]
\rel{lower\_bounds} & 5 & $X$ provides a lower bound on $Y$. \newline Inverse: \rel{bounded\_below\_by}.  Requires \texttt{citation}. \\[3pt]
\rel{strictly\_stronger} & 4 & $X \supsetneq Y$: $X$ implies $Y$ but not conversely, with a known separating witness.  These edges encode the strict hierarchies among learning criteria. \newline Inverse: \rel{strictly\_weaker}.  Requires \texttt{witness} and \texttt{citation}. \\[3pt]
\rel{requires\_assumption} & 2 & $X$ holds only if $Y$ is assumed (e.g., cryptographic assumptions for hardness results). \newline Inverse: \rel{assumed\_by}.  Requires \texttt{citation}. \\
\end{longtable}

% ----------------------------------------------------------
\section{Bibliography Link Validation}
\label{sec:bib-validation}

\texttt{validate\_bibliography\_links.py} performs three additional checks beyond those in the graph validator:

\begin{enumerate}[leftmargin=2em]
\item \textbf{Missing keys} (FAIL).  BibTeX keys referenced in node \texttt{bib\_keys}, \texttt{provenance.introduced\_by}, \texttt{provenance.proved\_by}, or edge \texttt{citation} fields that do not appear in \texttt{flt\_bibliography.bib}.  Each missing key is listed with all locations where it is referenced.

\item \textbf{Malformed bib keys} (FAIL).  Entries in the \texttt{.bib} file whose keys do not match the canonical pattern \texttt{[A-Za-z]+[0-9]\{4\}[a-z]?} (e.g., \texttt{BlumerEhrenfeuchtHausslerWarmuth1989}).

\item \textbf{Orphan keys} (WARN).  Entries in the \texttt{.bib} file that are never referenced by any node or edge.  These do not cause a failure but indicate potential dead references.
\end{enumerate}

\noindent The script reports the total number of BibTeX entries and unique keys referenced in the graph, then lists any missing, malformed, or orphaned keys with their locations.

% ----------------------------------------------------------
\section{Sample Validator Output}
\label{sec:sample-output}

When all checks pass, \texttt{validate\_graph.py} produces:

\begin{verbatim}
=================================================================
  FLT Concept Graph Validation
=================================================================

  [PASS]  JSON parses correctly

  [PASS]  Unique node IDs

  [PASS]  Edge endpoints exist

  [PASS]  Relation values in enum

  [PASS]  Required node fields

  [PASS]  Required edge fields

  [PASS]  Edge constraints

  [PASS]  No duplicate edges

  [PASS]  Bib keys resolve

  [PASS]  Node count matches meta

  [PASS]  Edge count matches meta

  [PASS]  Status values in enum

  [PASS]  No generalizes_unclassified edges

=================================================================
  13/13 checks passed, 0 failed.
  RESULT: PASS
=================================================================
\end{verbatim}

\noindent When a check fails, the validator prints up to 20 detail lines identifying specific violations.  For example, a dangling edge endpoint produces:

\begin{verbatim}
  [FAIL]  Edge endpoints exist
           Edge 47: source 'nonexistent_node' is not a node ID
\end{verbatim}

\noindent A missing witness on a separation edge produces:

\begin{verbatim}
  [FAIL]  Edge constraints
           Edge pac_learning-[does_not_imply]->mistake_bounded:
           constraint requires 'witness' but it is missing or empty
\end{verbatim}

% ----------------------------------------------------------
\section{Running the Benchmark Suite}
\label{sec:run-benchmark}

The benchmark consists of 15 tasks (see Appendix~\ref{app:traversals}).  Execution and evaluation are separate steps.

\paragraph{Step 1: Execute tasks.}
\begin{verbatim}
  python3 scripts/run_reference_tasks.py \
      flt_concept_graph.json flt_tasks.json \
      --output results.json
\end{verbatim}

\noindent This traverses the graph for each task and writes structured outputs to \texttt{results.json}.

\paragraph{Step 2: Evaluate against gold answers.}
\begin{verbatim}
  python3 scripts/evaluate_reference_tasks.py \
      results.json flt_task_answers.json
\end{verbatim}

\noindent Scoring modes include \texttt{set\_match} (exact set equality), \texttt{chain\_match} (ordered sequence), and \texttt{structured\_explanation} (required elements present in structured output).  The expected output on a correct implementation is 15/15 tasks passed.

% ============================================================
% APPENDIX D: Notation Index
% ============================================================
\chapter{Notation Index}
\label{app:notation}

This appendix collects all notation used in the text, organized by category.
The \emph{Introduced} column gives the chapter of first use.
Symbols defined by \LaTeX\ macros in \texttt{main.tex} are marked with the macro name in the \emph{Notes} column.

% ----------------------------------------------------------
\section*{D.1\quad Sets and Spaces}
\label{sec:notation-sets}

\begin{longtable}{lp{7.5cm}l}
\toprule
\textbf{Symbol} & \textbf{Description} & \textbf{Introduced} \\
\midrule
\endfirsthead
\toprule
\textbf{Symbol} & \textbf{Description} & \textbf{Introduced} \\
\midrule
\endhead
\midrule
\multicolumn{3}{r}{\footnotesize\itshape continued on next page} \\
\endfoot
\bottomrule
\endlastfoot
$X$ & Domain (input space).  Typically $\{0,1\}^n$, $\R^d$, or a countable set. & Ch.~1 \\
$Y$ & Label space (output space).  $\{0,1\}$ for binary, finite set for multiclass, $\R$ for regression. & Ch.~1 \\
$\Sigma$ & Alphabet (finite, nonempty set of symbols). & Ch.~3 \\
$\Sigma^*$ & Free monoid: all finite strings over $\Sigma$. & Ch.~3 \\
$\mathcal{C}$ & Concept class: a collection of concepts $c : X \to Y$. & Ch.~1 \\
$\mathcal{H}$ & Hypothesis space: the set from which the learner selects hypotheses. & Ch.~1 \\
$\mathcal{F}$ & Real-valued function class $\mathcal{F} \subseteq Y^X$ (typically $Y = \R$). & Ch.~10 \\
$S = \{(x_i, y_i)\}_{i=1}^m$ & Training sample of $m$ labeled examples. & Ch.~2 \\
$D$ & Distribution on $X \times Y$ (or on $X$ in the realizable setting). & Ch.~2 \\
$\N$ & Natural numbers $\{0,1,2,\ldots\}$. & Ch.~1 \\
$\Z$ & Integers. & Ch.~1 \\
$\R$ & Real numbers. & Ch.~1 \\
$\{0,1\}^n$ & Binary strings of length $n$. & Ch.~1 \\
$\mathcal{B}_\varepsilon(x)$ & Open ball of radius $\varepsilon$ centered at $x$. & Ch.~10 \\
\end{longtable}

% ----------------------------------------------------------
\section*{D.2\quad Complexity Measures}
\label{sec:notation-complexity}

\begin{longtable}{lp{7.5cm}l}
\toprule
\textbf{Symbol} & \textbf{Description} & \textbf{Introduced} \\
\midrule
\endfirsthead
\toprule
\textbf{Symbol} & \textbf{Description} & \textbf{Introduced} \\
\midrule
\endhead
\midrule
\multicolumn{3}{r}{\footnotesize\itshape continued on next page} \\
\endfoot
\bottomrule
\endlastfoot
$\VC(\mathcal{H})$ & Vapnik--Chervonenkis dimension: largest set shattered by $\mathcal{H}$. Macro: \verb|\VC|. & Ch.~5 \\
$\Ldim(\mathcal{H})$ & Littlestone dimension: depth of deepest mistake tree for $\mathcal{H}$. Macro: \verb|\Ldim|. & Ch.~6 \\
$\DSdim(\mathcal{H})$ & DS dimension (Daniely--Shalev-Shwartz): multiclass PAC characterization. Macro: \verb|\DSdim|. & Ch.~17 \\
$\Ndim(\mathcal{H})$ & Natarajan dimension: multiclass generalization of VC dimension via two-colorings. Macro: \verb|\Ndim|. & Ch.~17 \\
$\Pdim(\mathcal{F})$ & Pseudodimension: VC dimension of the subgraph class $\{(x,t) : f(x) \geq t\}$. Macro: \verb|\Pdim|. & Ch.~10 \\
$\fatshat_\gamma(\mathcal{F})$ & Fat-shattering dimension at scale $\gamma > 0$. Macro: \verb|\fatshat|. & Ch.~10 \\
$\SQdim(\mathcal{C})$ & Statistical query dimension. Macro: \verb|\SQdim|. & Ch.~10 \\
$\Rad_n(\mathcal{H})$ & Empirical Rademacher complexity: $\E_\sigma\bigl[\sup_{h \in \mathcal{H}} \frac{1}{n}\sum_{i=1}^n \sigma_i h(x_i)\bigr]$. Macro: \verb|\Rad|. & Ch.~12 \\
$\growth_\mathcal{H}(n)$ & Growth function (Sauer--Shelah): $\max_{|S|=n} |\{h|_S : h \in \mathcal{H}\}|$. Macro: \verb|\growth|. & Ch.~10 \\
$\KL(P \| Q)$ & Kullback--Leibler divergence between distributions $P$ and $Q$. Macro: \verb|\KL|. & Ch.~12 \\
$\mathcal{N}(\varepsilon, \mathcal{F}, d)$ & Covering number: minimum $\varepsilon$-net size under metric $d$. & Ch.~10 \\
$K(x)$ & Kolmogorov complexity of string $x$. & Ch.~13 \\
$\mathrm{dl}(h)$ & Description length of hypothesis $h$. & Ch.~11 \\
\end{longtable}

% ----------------------------------------------------------
\section*{D.3\quad Learning Parameters}
\label{sec:notation-params}

\begin{longtable}{lp{7.5cm}l}
\toprule
\textbf{Symbol} & \textbf{Description} & \textbf{Introduced} \\
\midrule
\endfirsthead
\toprule
\textbf{Symbol} & \textbf{Description} & \textbf{Introduced} \\
\midrule
\endhead
\bottomrule
\endlastfoot
$\varepsilon$ & Accuracy parameter: $\risk(h) \leq \varepsilon$ or $\risk(h) \leq \risk(h^*) + \varepsilon$. & Ch.~5 \\
$\delta$ & Confidence parameter: bounds hold with probability $\geq 1 - \delta$. & Ch.~5 \\
$m$ & Sample size (number of training examples). & Ch.~2 \\
$m(\varepsilon, \delta)$ & Sample complexity: minimum $m$ sufficient for $(\varepsilon, \delta)$-learning. & Ch.~5 \\
$T$ & Number of rounds in online learning or time horizon. & Ch.~6 \\
$\gamma$ & Margin parameter (for fat-shattering dimension and margin bounds). & Ch.~10 \\
$k$ & Number of labels $|Y|$ in multiclass settings. & Ch.~17 \\
$n$ & Input dimensionality $|x|$ or sample subset size, depending on context. & Ch.~1 \\
\end{longtable}

% ----------------------------------------------------------
\section*{D.4\quad Risk and Loss}
\label{sec:notation-risk}

\begin{longtable}{lp{7.5cm}l}
\toprule
\textbf{Symbol} & \textbf{Description} & \textbf{Introduced} \\
\midrule
\endfirsthead
\toprule
\textbf{Symbol} & \textbf{Description} & \textbf{Introduced} \\
\midrule
\endhead
\bottomrule
\endlastfoot
$\ell(h(x), y)$ & Loss function (0-1 loss unless otherwise stated). & Ch.~5 \\
$\risk(h)$ & True (population) risk: $\E_{(x,y) \sim D}[\ell(h(x),y)]$. Macro: \verb|\risk|. & Ch.~5 \\
$\emprisk(h)$ or $\emprisk_S(h)$ & Empirical risk: $\frac{1}{m}\sum_{i=1}^m \ell(h(x_i),y_i)$. Macro: \verb|\emprisk|. & Ch.~5 \\
$\risk(h) - \emprisk(h)$ & Generalization gap. & Ch.~12 \\
$R_T$ & Cumulative regret over $T$ rounds. & Ch.~6 \\
$M_T$ & Mistake count over $T$ rounds. & Ch.~6 \\
\end{longtable}

% ----------------------------------------------------------
\section*{D.5\quad Distributions and Probability}
\label{sec:notation-prob}

\begin{longtable}{lp{7.5cm}l}
\toprule
\textbf{Symbol} & \textbf{Description} & \textbf{Introduced} \\
\midrule
\endfirsthead
\toprule
\textbf{Symbol} & \textbf{Description} & \textbf{Introduced} \\
\midrule
\endhead
\bottomrule
\endlastfoot
$D$ & Data-generating distribution on $X \times Y$. & Ch.~2 \\
$D_X$ & Marginal distribution on $X$. & Ch.~5 \\
$\Pr(\cdot)$ or $\mathbb{P}(\cdot)$ & Probability measure. Macro: \verb|\Pr|. & Ch.~2 \\
$\E[\cdot]$ & Expectation. Macro: \verb|\E|. & Ch.~2 \\
$\ind_A$ or $\ind[\text{event}]$ & Indicator function / indicator random variable. Macro: \verb|\ind|. & Ch.~2 \\
$\sigma_i$ & Rademacher random variable: $\Pr(\sigma_i = +1) = \Pr(\sigma_i = -1) = \tfrac{1}{2}$. & Ch.~12 \\
$Q$ & Posterior distribution (PAC-Bayes) or generic distribution over $\mathcal{H}$. & Ch.~12 \\
$P$ & Prior distribution over $\mathcal{H}$ (PAC-Bayes). & Ch.~12 \\
$S \sim D^m$ & Sample drawn i.i.d.\ from $D$. & Ch.~2 \\
\end{longtable}

% ----------------------------------------------------------
\section*{D.6\quad Ordinals and Mind-Changes}
\label{sec:notation-ordinals}

\begin{longtable}{lp{7.5cm}l}
\toprule
\textbf{Symbol} & \textbf{Description} & \textbf{Introduced} \\
\midrule
\endfirsthead
\toprule
\textbf{Symbol} & \textbf{Description} & \textbf{Introduced} \\
\midrule
\endhead
\bottomrule
\endlastfoot
$\omega$ & First infinite ordinal; cardinality of $\N$. & Ch.~7 \\
$\omega^k$ & Ordinal exponentiation (iterated $\omega$). & Ch.~13 \\
$\omega^\omega$ & First epsilon-unreachable ordinal power (limit of $\omega, \omega^2, \omega^3, \ldots$). & Ch.~13 \\
$\alpha, \beta$ & Arbitrary ordinals. & Ch.~13 \\
$\mathrm{MC}(L, \mathcal{C})$ & Mind-change count: maximum number of hypothesis changes by learner $L$ on class $\mathcal{C}$. & Ch.~7 \\
$\mathrm{MC}_\alpha$ & Ordinal mind-change bound at level $\alpha$. & Ch.~13 \\
$\omega\text{-}\VC$ & Ordinal VC dimension (transfinite extension). & Ch.~13 \\
$\Delta^0_2$ & Arithmetic hierarchy class: limits of computable functions. & Ch.~7 \\
\end{longtable}

% ----------------------------------------------------------
\section*{D.7\quad Information-Theoretic Quantities}
\label{sec:notation-info}

\begin{longtable}{lp{7.5cm}l}
\toprule
\textbf{Symbol} & \textbf{Description} & \textbf{Introduced} \\
\midrule
\endfirsthead
\toprule
\textbf{Symbol} & \textbf{Description} & \textbf{Introduced} \\
\midrule
\endhead
\bottomrule
\endlastfoot
$H(X)$ & Shannon entropy of random variable $X$. & Ch.~12 \\
$H(X \mid Y)$ & Conditional Shannon entropy. & Ch.~12 \\
$I(X; Y)$ & Mutual information between $X$ and $Y$: $H(X) - H(X \mid Y)$. & Ch.~12 \\
$\KL(P \| Q)$ & Kullback--Leibler divergence: $\E_P[\log(P/Q)]$. Macro: \verb|\KL|. & Ch.~12 \\
$I(S; A(S))$ & Input--output mutual information of algorithm $A$ on sample $S$ (Xu--Raginsky framework). & Ch.~12 \\
\end{longtable}

% ----------------------------------------------------------
\section*{D.8\quad Learning Criteria and Algorithms}
\label{sec:notation-criteria}

\begin{longtable}{lp{7.5cm}l}
\toprule
\textbf{Symbol} & \textbf{Description} & \textbf{Introduced} \\
\midrule
\endfirsthead
\toprule
\textbf{Symbol} & \textbf{Description} & \textbf{Introduced} \\
\midrule
\endhead
\bottomrule
\endlastfoot
$\Ex$ & Explanatory learning (Gold-style identification in the limit). Macro: \verb|\Ex|. & Ch.~7 \\
$\BC$ & Behaviorally correct learning. Macro: \verb|\BC|. & Ch.~7 \\
$\FIN$ & Finite identification (the learner eventually stops and is correct). Macro: \verb|\FIN|. & Ch.~7 \\
$\ERM$ & Empirical risk minimization. Macro: \verb|\ERM|. & Ch.~5 \\
$\SOA$ & Standard optimal algorithm (online prediction). Macro: \verb|\SOA|. & Ch.~6 \\
$L^*$ & Angluin's $L^*$ algorithm for learning regular languages from queries. & Ch.~8 \\
CEGIS & Counterexample-guided inductive synthesis. & Ch.~8 \\
\end{longtable}

% ----------------------------------------------------------
\section*{D.9\quad Functions and Maps}
\label{sec:notation-functions}

\begin{longtable}{lp{7.5cm}l}
\toprule
\textbf{Symbol} & \textbf{Description} & \textbf{Introduced} \\
\midrule
\endfirsthead
\toprule
\textbf{Symbol} & \textbf{Description} & \textbf{Introduced} \\
\midrule
\endhead
\bottomrule
\endlastfoot
$c : X \to Y$ & Target concept (the function the learner aims to identify). & Ch.~1 \\
$h : X \to Y$ & Hypothesis produced by the learner. & Ch.~1 \\
$h^* \in \mathcal{H}$ & Best hypothesis in $\mathcal{H}$: $h^* = \arg\min_{h \in \mathcal{H}} \risk(h)$. & Ch.~5 \\
$A : (X \times Y)^m \to \mathcal{H}$ & Learning algorithm (maps a sample to a hypothesis). & Ch.~4 \\
$\varphi_e$ & Partial computable function with G\"odel number $e$. & Ch.~3 \\
$W_e$ & Domain of $\varphi_e$ (recursively enumerable set with index $e$). & Ch.~3 \\
$h|_S$ & Restriction of $h$ to the set $S$: the behavior pattern of $h$ on points in $S$. & Ch.~5 \\
\end{longtable}

% ----------------------------------------------------------
\section*{D.10\quad Graph-Theoretic Notation}
\label{sec:notation-graph}

The following notation is specific to the knowledge graph structure described in this book.

\begin{longtable}{lp{7.5cm}l}
\toprule
\textbf{Symbol / Macro} & \textbf{Description} & \textbf{Introduced} \\
\midrule
\endfirsthead
\toprule
\textbf{Symbol / Macro} & \textbf{Description} & \textbf{Introduced} \\
\midrule
\endhead
\bottomrule
\endlastfoot
\verb|\nde{id}| $\to$ \nde{id} & Graph node identifier (sans-serif). & Preface \\
\verb|\rel{R}| $\to$ \rel{R} & Relation type label (monospaced). & Preface \\
$\edge{A}{R}{B}$ & Typed directed edge from \nde{A} to \nde{B} via \rel{R}. Macro: \verb|\edge{A}{R}{B}|. & Preface \\
Layer $\ell \in \{0, \ldots, 7\}$ & Stratification of nodes: 0~=~formal objects, 1~=~base types, 2~=~data presentations, 3~=~learner types, 4~=~success criteria, 5~=~complexity measures, 6~=~characterizations/impossibilities, 7~=~processes/scope boundaries. & Ch.~1 \\
Kernel node & A node with status \texttt{defined} or \texttt{proved} (133 of 142). & Preface \\
Deferred node & A node with status \texttt{deferred} (6 of 142). & Preface \\
Scope boundary & A node with status \texttt{scope\_note} marking an explicit exclusion (3 of 142). & Preface \\
\end{longtable}

\bigskip
\noindent\textbf{Complete symbol list.}\quad
This appendix covers the most frequently used symbols.  Additional notation local to a single proof or example is defined at the point of use.  All \LaTeX\ macros are defined in the preamble of \texttt{main.tex}; see \S\ref{sec:notation-graph} above for the graph-specific macros \verb|\nde|, \verb|\rel|, and \verb|\edge|.


% ============================================================
\backmatter
% ============================================================

\bibliographystyle{alpha}
\bibliography{../flt_bibliography}

\end{document}
